% !TeX program = pdflatex
% Integracja audytu II: 26 jednostek po wydzieleniu teorii dystrybucji; zachowano historyczną numerację.
% Wersje samodzielne zachowano w katalogu wersje_samodzielne/.
\documentclass[11pt,a4paper]{article}
\usepackage[T1]{fontenc}
\usepackage[utf8]{inputenc}
\usepackage[polish]{babel}
\usepackage{lmodern}
\usepackage{amsmath,amssymb,amsthm,mathtools}
\usepackage{microtype}
\usepackage{geometry}
\usepackage{enumitem}
\usepackage{tocloft}
\usepackage{xcolor}
\usepackage{tikz}
\usetikzlibrary{arrows.meta}
\usepackage[unicode]{hyperref}
\geometry{margin=2.5cm}
\hypersetup{colorlinks=true,linkcolor=blue!45!black,urlcolor=blue!45!black,pdftitle={Analiza funkcjonalna - wersja scalona}}
\setlength{\parindent}{0pt}
\setlength{\parskip}{4pt}
\setlength{\emergencystretch}{2em}
\setlength{\cftsecnumwidth}{3em}
\setlength{\cftsubsecnumwidth}{4.5em}
% Moduły 15A/15B mają numery podpodsekcji w rodzaju 15B.12.8.
\setlength{\cftsubsubsecnumwidth}{6em}
\setlist[enumerate]{topsep=4pt,itemsep=2pt}
\setlist[itemize]{topsep=4pt,itemsep=2pt}
\numberwithin{equation}{section}
\theoremstyle{plain}
\newtheorem{twierdzenie}{Twierdzenie}[section]
\newtheorem{stwierdzenie}[twierdzenie]{Stwierdzenie}
\newtheorem{lemat}[twierdzenie]{Lemat}
\newtheorem{wniosek}[twierdzenie]{Wniosek}
\theoremstyle{definition}
\newtheorem{definicja}[twierdzenie]{Definicja}
\newtheorem{przyklad}[twierdzenie]{Przykład}
\theoremstyle{remark}
\newtheorem{uwaga}[twierdzenie]{Uwaga}
\newcommand{\R}{\mathbb R}
\newcommand{\C}{\mathbb C}
\newcommand{\N}{\mathbb N}
\newcommand{\K}{\mathbb K}
\newcommand{\cl}[1]{\overline{#1}}
\newcommand{\TVS}{\mathrm{TVS}}
\DeclareMathOperator{\lin}{lin}
\DeclareMathOperator{\supp}{supp}
\DeclareMathOperator{\id}{id}
\title{Analiza funkcjonalna\\\large Przestrzenie liniowo-topologiczne, geometria przestrzeni Banacha i teoria operatorów}
\author{}
\date{Wersja scalona po drugiej rundzie audytu -- 19 listopada 2026}

% BEGIN GLOBAL CODE LISTINGS (generated from latex/pl/code-listings.tex)
\usepackage{xcolor}
\usepackage{listings}
\definecolor{codebg}{RGB}{248,249,251}
\definecolor{codeborder}{RGB}{180,191,203}
\definecolor{codecomment}{RGB}{61,113,79}
\definecolor{codekeyword}{RGB}{41,78,132}
\lstdefinelanguage{Julia}{
  morekeywords={abstract,baremodule,begin,break,catch,const,continue,do,else,elseif,
    end,export,false,finally,for,function,global,if,import,in,isa,let,local,macro,
    module,mutable,primitive,quote,return,struct,true,try,type,using,where,while,
    missing,nothing,println,print},
  sensitive=true,morecomment=[l]{\#},morecomment=[n]{\#=}{=\#},
  morestring=[b]",
}
\lstdefinestyle{skrypt}{
  backgroundcolor=\color{codebg},frame=single,rulecolor=\color{codeborder},
  framesep=5pt,frameround=tttt,framerule=.5pt,
  basicstyle=\ttfamily\footnotesize,keywordstyle=\color{codekeyword}\bfseries,
  commentstyle=\color{codecomment}\itshape,stringstyle=\color{purple!65!black},
  numbers=left,numberstyle=\tiny\color{gray!75!black},numbersep=8pt,
  stepnumber=1,breaklines=true,breakatwhitespace=false,
  columns=fullflexible,keepspaces=true,showstringspaces=false,tabsize=4,
  linewidth=.90\linewidth,captionpos=b,aboveskip=11pt,belowskip=12pt,
  xleftmargin=12pt,xrightmargin=2pt,
}
\lstdefinestyle{pythonpde}{style=skrypt,language=Python}
\lstdefinestyle{python}{style=skrypt,language=Python}
\lstdefinestyle{julia}{style=skrypt,language=Julia}
\lstset{style=skrypt}
\renewcommand{\lstlistingname}{Kod}
\renewcommand{\lstlistlistingname}{Spis kodów}
% END GLOBAL CODE LISTINGS
\begin{document}
\maketitle
\tableofcontents
\bigskip
% Wersja źródłowa: audyt III, jednostka 01
\section{Przestrzenie liniowo-topologiczne: definicje i pierwsze przykłady}\label{sec:tvs}

\subsection{Cel i konwencje}

W analizie funkcjonalnej bada się przestrzenie liniowe, w których można
wykonywać przejścia graniczne, oraz odwzorowania liniowe zgodne z tymi
przejściami. Norma jest tylko jednym ze sposobów wprowadzenia topologii.
Na przykład topologia zbieżności punktowej funkcji czy naturalna topologia
funkcji gładkich nie muszą pochodzić od pojedynczej normy. Dlatego zaczniemy
od aksjomatów, które wiążą strukturę liniową z topologiczną.

W całym skrypcie $\K$ oznacza $\R$ albo $\C$ z ich zwykłą topologią.
Przestrzenie liniowe rozpatrujemy nad $\K$. Przez otoczenie punktu $x$
rozumiemy zbiór zawierający pewien zbiór otwarty zawierający $x$;
w szczególności otoczenie nie musi być otwarte. Przez $\mathcal N_X(x)$
oznaczamy rodzinę wszystkich otoczeń $x$ w przestrzeni topologicznej $X$.
Słowo „ciągły” zawsze odnosi się do rozważanych topologii, nie zaś
do niejawnie wybranej normy.

\begin{definicja}[Przestrzeń liniowo-topologiczna]\label{def:tvs}
Przestrzenią liniowo-topologiczną nad $\K$ nazywamy przestrzeń liniową
$X$ z topologią $\tau$, dla której odwzorowania
\[
 +\colon X\times X\longrightarrow X,\qquad (x,y)\longmapsto x+y,
 \qquad
 m\colon\K\times X\longrightarrow X,\qquad (\lambda,x)\longmapsto\lambda x
\]
są ciągłe, gdy w dziedzinach przyjmiemy topologie produktowe.
Nie wymagamy na razie warunku Hausdorffa.
\end{definicja}

Warunek ciągłości mnożenia przez skalar jest \emph{łączny} względem
$\lambda$ i $x$. Sama ciągłość każdej funkcji $x\mapsto\lambda x$,
dla ustalonego $\lambda$, jest warunkiem słabszym.

\begin{przyklad}[Topologia trywialna]
Niech $X$ będzie dowolną niezerową przestrzenią liniową.
Topologia $\{\varnothing,X\}$ czyni $X$ przestrzenią
liniowo-topologiczną: każde odwzorowanie o wartościach w $X$ jest
ciągłe. Nie jest ona jednak Hausdorffa. Przykład wyjaśnia,
dlaczego rozdzielamy aksjomaty przestrzeni liniowo-topologicznej
od aksjomatu rozdzielania punktów.
\end{przyklad}

\subsection{Operacje liniowe jako homeomorfizmy}

\begin{stwierdzenie}\label{prop:translation}
Niech $X$ będzie przestrzenią liniowo-topologiczną nad $\K$.
Dla każdego $a\in X$ przesunięcie $T_a(x)=x+a$ jest
homeomorfizmem $X$ na $X$. Dla każdego $\lambda\in\K\setminus\{0\}$
odwzorowanie $M_\lambda(x)=\lambda x$ jest homeomorfizmem $X$ na $X$.
\end{stwierdzenie}
\begin{proof}
Ciągłość $T_a$ wynika ze złożenia ciągłego odwzorowania
$x\mapsto(x,a)$ z dodawaniem. Jego odwrotnością jest ciągłe $T_{-a}$.
Analogicznie $M_\lambda$ jest ciągłe jako funkcja jednej zmiennej
otrzymana z mnożenia przez skalar, a jego odwrotnością jest
$M_{1/\lambda}$.
\end{proof}

\begin{wniosek}\label{cor:zero-neighborhood}
Dla dowolnej przestrzeni liniowo-topologicznej $X$, punktu $x\in X$
i podzbioru $U\subseteq X$ zachodzi równoważność:
$U\in\mathcal N_X(x)$ wtedy i tylko wtedy, gdy
$U-x\in\mathcal N_X(0)$. W szczególności topologia przestrzeni
liniowo-topologicznej jest całkowicie określona przez otoczenia zera.
\end{wniosek}
\begin{proof}
Korzystamy z tego, że $T_{-x}$ jest homeomorfizmem i
$T_{-x}(x)=0$.
\end{proof}

\begin{uwaga}
Jednorodność topologiczna nie oznacza, że każda przestrzeń
liniowo-topologiczna jest metryzowalna. Problem metryzowalności
będziemy rozstrzygać za pomocą baz otoczeń zera.
\end{uwaga}

\subsection{Otoczenia zera, zbiory zrównoważone i pochłaniające}

\begin{definicja}\label{def:balanced-absorbing}
Niech $X$ będzie przestrzenią liniową nad $\K$, a $A\subseteq X$.
Mówimy, że:
\begin{enumerate}[label=\textup{(\roman*)}]
 \item $A$ jest \emph{zrównoważony}, jeżeli $\lambda A\subseteq A$
       dla każdego $\lambda\in\K$ z $|\lambda|\leq1$;
 \item $A$ jest \emph{pochłaniający}, jeżeli dla każdego $x\in X$
       istnieje $t>0$ takie, że $x\in tA$;
 \item $A$ jest \emph{wypukły}, jeżeli $sx+(1-s)y\in A$
       dla wszystkich $x,y\in A$ i $s\in[0,1]$.
\end{enumerate}
W drugim punkcie $tA=\{ta:a\in A\}$, gdzie $t$ jest rzeczywistą
liczbą dodatnią, również gdy $\K=\C$.
\end{definicja}

\begin{uwaga}
Wszystkie trzy warunki są algebraiczne. Zrównoważenie nie
implikuje wypukłości. W analizie lokalnie wypukłej będziemy
konstruować bazy otoczeń zera złożone ze zbiorów jednocześnie
wypukłych i zrównoważonych.
\end{uwaga}

\begin{twierdzenie}[Podstawowe własności otoczeń zera]\label{thm:zero-neighborhoods}
Niech $X$ będzie przestrzenią liniowo-topologiczną nad $\K$.
Dla każdego $U\in\mathcal N_X(0)$ zachodzą następujące własności:
\begin{enumerate}[label=\textup{(\roman*)}]
 \item istnieje $V\in\mathcal N_X(0)$ takie, że $V+V\subseteq U$;
 \item istnieje otwarte, symetryczne otoczenie $V$ zera
       ($V=-V$) zawarte w $U$;
 \item istnieje otwarte, zrównoważone otoczenie $W$ zera
       zawarte w $U$;
 \item $U$ jest zbiorem pochłaniającym.
\end{enumerate}
\end{twierdzenie}
\begin{proof}
Wybierzmy otwarte otoczenie $U_0$ zera z $U_0\subseteq U$.
Ciągłość dodawania w punkcie $(0,0)$ daje otoczenia otwarte
$V_1,V_2$ zera takie, że $V_1+V_2\subseteq U_0$.
Dla $V=V_1\cap V_2$ otrzymujemy (i). Zbiór
$U_0\cap(-U_0)$ jest otwartym, symetrycznym otoczeniem zera;
to dowodzi (ii).

Dla (iii) oznaczmy przez $D=\{\lambda\in\K:|\lambda|\leq1\}$
domknięty dysk jednostkowy (dla $\K=\R$ przedział $[-1,1]$).
Zbiór $D$ jest zwarty. Dla każdego $\lambda\in D$, z ciągłości
mnożenia przez skalar w $(\lambda,0)$, istnieją otoczenia otwarte
$O_\lambda\subseteq\K$ punktu $\lambda$ oraz $V_\lambda\subseteq X$
zera takie, że $O_\lambda V_\lambda\subseteq U_0$.
Wybieramy skończone podpokrycie
$D\subseteq O_{\lambda_1}\cup\dots\cup O_{\lambda_m}$ i kładziemy
$V=\bigcap_{j=1}^m V_{\lambda_j}$. Wtedy $DV\subseteq U_0$.
Zbiór $A=DV$ jest zrównoważony, ponieważ $\mu D\subseteq D$
dla $|\mu|\leq1$. Ponieważ $1\in D$, zachodzi $V\subseteq A$,
a zatem $A$ jest otoczeniem zera. Wreszcie $W=\operatorname{int}A$
jest otwarte i zawiera zero. Jeśli $0<|\mu|\leq1$, to
$\mu W$ jest otwartym podzbiorem $\mu A\subseteq A$, wobec czego
$\mu W\subseteq\operatorname{int}A=W$. Dla $\mu=0$ mamy
$\mu W=\{0\}\subseteq W$. Tak więc $W$ jest zrównoważone.

Dla (iv) ustalmy $x\in X$. Funkcja
$\K\ni\lambda\mapsto\lambda x\in X$ jest ciągła i przyjmuje
wartość zero dla $\lambda=0$. Istnieje zatem rzeczywiste
$\varepsilon>0$ takie, że $\varepsilon x\in U$.
Wtedy $x\in\varepsilon^{-1}U$, czyli $U$ jest pochłaniający.
\end{proof}

\begin{przyklad}[Zrównoważenie bez wypukłości]
Zbiór $A=\{(x,y)\in\R^2:|x|\leq1,\ |y|\leq1\}\cup
\{(x,y)\in\R^2:|x|\leq2,\ |y|\leq\frac12\}$
jest zrównoważony, lecz nie jest wypukły:
$(1,1),(2,\tfrac12)\in A$, podczas gdy
$(\tfrac32,\tfrac34)\notin A$.
\end{przyklad}

\subsection{Warunek Hausdorffa i domknięcie zera}

\begin{twierdzenie}\label{thm:hausdorff-zero}
Niech $X$ będzie przestrzenią liniowo-topologiczną nad $\K$.
Następujące warunki są równoważne:
\begin{enumerate}[label=\textup{(\roman*)}]
 \item $X$ jest przestrzenią Hausdorffa;
 \item zbiór $\{0\}$ jest domknięty w $X$;
 \item $\bigcap_{U\in\mathcal N_X(0)}U=\{0\}$.
\end{enumerate}
Ponadto $N=\cl{\{0\}}$ jest domkniętą podprzestrzenią liniową
$X$ oraz
\[
 N=\bigcap_{U\in\mathcal N_X(0)}U.
\]
\end{twierdzenie}
\begin{proof}
Implikacja (i)$\Rightarrow$(ii) zachodzi w każdej przestrzeni
Hausdorffa. Aby wykazać wzór na $N$, ustalmy $x\in X$.
Jeżeli $x\in N$, to każde otoczenie $x$ zawiera zero.
Dla dowolnego symetrycznego otoczenia $V$ zera oznacza to,
że $0\in x+V$, czyli $-x\in V$, a z symetrii także $x\in V$.
Ponieważ symetryczne otoczenia tworzą bazę, $x$ należy do
wszystkich otoczeń zera.

Odwrotnie, jeśli $x$ należy do wszystkich otoczeń zera,
to dla każdego otoczenia $x+V$ punktu $x$, gdzie $V$ jest
symetrycznym otoczeniem zera, mamy $-x\in V$, skąd
$0\in x+V$. Oznacza to $x\in N$. Udowodniliśmy wzór.

Jeżeli (ii), to $N=\{0\}$ i powyższy wzór daje (iii).
Jeżeli (iii), to dla dowolnych $x\ne y$ istnieje otoczenie
$U$ zera, do którego nie należy $y-x$. Dobierzmy symetryczne
otoczenie $V$ zera takie, że $V+V\subseteq U$. Otoczenia
$x+V$ i $y+V$ są rozłączne: w przeciwnym razie
$y-x\in V-V=V+V\subseteq U$. To daje (i).

Pozostaje wykazać, że $N$ jest podprzestrzenią.
Jeśli $x,y\in N$ i $U$ jest otoczeniem zera, wybierzmy
otoczenie $V$ z $V+V\subseteq U$. Wtedy $x,y\in V$,
a więc $x+y\in U$. Ponieważ $U$ było dowolne, $x+y\in N$.
Dla $\lambda\ne0$ i $x\in N$, zbiór $\lambda^{-1}U$
jest otoczeniem zera, więc $x\in\lambda^{-1}U$,
czyli $\lambda x\in U$; dla $\lambda=0$ jest to oczywiste.
Zatem $N$ jest podprzestrzenią. Jest domknięta z definicji.
\end{proof}

\begin{uwaga}
Od tego miejsca \emph{przestrzeń liniowo-topologiczna} będzie,
o ile nie zaznaczymy inaczej, oznaczać przestrzeń liniowo-topologiczną
Hausdorffa. Warto jednak pamiętać, że wymaganie Hausdorffa jest
założeniem, a nie konsekwencją ciągłości działań liniowych.
\end{uwaga}

\subsection{Ciągłe odwzorowania liniowe}

\begin{twierdzenie}[Kryterium ciągłości w zerze]\label{thm:linear-continuity-zero}
Niech $X,Y$ będą przestrzeniami liniowo-topologicznymi nad tym samym
ciałem $\K$ i niech $T\colon X\to Y$ będzie odwzorowaniem liniowym.
Następujące warunki są równoważne:
\begin{enumerate}[label=\textup{(\roman*)}]
 \item $T$ jest ciągłe na $X$;
 \item $T$ jest ciągłe w punkcie $0$;
 \item dla każdego $V\in\mathcal N_Y(0)$ istnieje
       $U\in\mathcal N_X(0)$ takie, że $T(U)\subseteq V$.
\end{enumerate}
\end{twierdzenie}
\begin{proof}
Implikacja (i)$\Rightarrow$(ii) jest oczywista.
Warunek (iii) jest dokładnym zapisem definicji ciągłości w zerze,
ponieważ $T(0)=0$.
Jeśli (iii) oraz $x\in X$, to dla dowolnego otoczenia
$T(x)+V$ punktu $T(x)$ istnieje otoczenie $U$ zera takie,
że $T(U)\subseteq V$. Z liniowości
$T(x+U)=T(x)+T(U)\subseteq T(x)+V$.
Oznacza to ciągłość $T$ w punkcie $x$, a więc (i).
\end{proof}

\begin{wniosek}\label{cor:kernel-closed}
Niech $X$ będzie przestrzenią liniowo-topologiczną, $Y$
przestrzenią liniowo-topologiczną Hausdorffa, a $T\colon X\to Y$
ciągłym odwzorowaniem liniowym. Wtedy $\ker T$ jest domkniętą
podprzestrzenią liniową $X$.
\end{wniosek}
\begin{proof}
Jest podprzestrzenią z liniowości $T$. Ponieważ $Y$ jest
Hausdorffa, zbiór $\{0\}$ jest domknięty na mocy
twierdzenia~\ref{thm:hausdorff-zero}. Mamy
$\ker T=T^{-1}(\{0\})$, zatem jądro jest domknięte.
\end{proof}

\begin{przyklad}[Ciągła bijekcja liniowa bez ciągłej odwrotności]
Na $X=C([0,1],\R)$ rozważmy dwie normy:
$\|f\|_\infty=\max_{x\in[0,1]}|f(x)|$ oraz
$\|f\|_1=\int_{[0,1]}|f(x)|\,\mathrm d\lambda(x)$,
gdzie $\lambda$ jest miarą Lebesgue'a.
Identyczność $I\colon(X,\|\cdot\|_\infty)\to(X,\|\cdot\|_1)$
jest liniową ciągłą bijekcją, bo $\|f\|_1\leq\|f\|_\infty$.
Odwrotność nie jest ciągła. Istotnie, funkcje
$f_n(x)=\max\{1-nx,0\}$ spełniają $\|f_n\|_\infty=1$
i $\|f_n\|_1=1/(2n)$.
Nie przeczy to twierdzeniu o odwzorowaniu otwartym:
przestrzeń $(X,\|\cdot\|_1)$ nie jest zupełna. Rzeczywiście,
można dobrać ciąg funkcji ciągłych $g_n$ przyjmujących wartości
w $[0,1]$, równych zeru na $[0,\tfrac12-\tfrac1n]$ oraz jeden
na $[\tfrac12+\tfrac1n,1]$, liniowych pomiędzy tymi przedziałami
(dla $n\geq3$). Wtedy
$\|g_n-\chi_{[1/2,1]}\|_1\leq2/n\to0$.
W szczególności $(g_n)$ jest ciągiem Cauchy'ego w normie
$\|\cdot\|_1$. Gdyby zbiegał do $g\in C([0,1])$,
jednoznaczność granicy w $L^1$ dawałaby
$g=\chi_{[1/2,1]}$ prawie wszędzie. Z ciągłości $g$
wynikałoby wtedy $g=0$ na $[0,1/2)$ oraz $g=1$ na
$(1/2,1]$, co jest sprzeczne z ciągłością w $1/2$.
Zatem $(X,\|\cdot\|_1)$ nie jest zupełna.
\end{przyklad}

\subsection{Zbiory ograniczone bez użycia normy}

\begin{definicja}[Ograniczoność w sensie von Neumanna]\label{def:tvs-bounded}
Niech $X$ będzie przestrzenią liniowo-topologiczną nad $\K$.
Zbiór $B\subseteq X$ nazywamy \emph{ograniczonym}, jeżeli
dla każdego otoczenia $U$ zera istnieje $t>0$ takie,
że $B\subseteq tU$. W definicji nie wymaga się, aby $B$
był domknięty ani aby był otoczeniem zera.
\end{definicja}

\begin{lemat}\label{lem:bounded-test}
Niech $X$ będzie przestrzenią liniowo-topologiczną, a $B\subseteq X$.
Zbiór $B$ jest ograniczony wtedy i tylko wtedy, gdy dla każdego
zrównoważonego otoczenia $U$ zera istnieje $\varepsilon>0$ takie,
że $\lambda B\subseteq U$ dla wszystkich $\lambda\in\K$
z $|\lambda|\leq\varepsilon$.
\end{lemat}
\begin{proof}
Załóżmy, że $B$ jest ograniczony, i weźmy zrównoważone
otoczenie $U$ zera. Istnieje $t>0$ z $B\subseteq tU$.
Dla $\varepsilon=1/t$ i $|\lambda|\leq\varepsilon$
mamy $\lambda B\subseteq(\lambda t)U\subseteq U$.

Odwrotnie, niech $V$ będzie dowolnym otoczeniem zera.
Na mocy twierdzenia~\ref{thm:zero-neighborhoods} istnieje
zrównoważone otoczenie $U\subseteq V$. Z założenia
$\varepsilon B\subseteq U$ dla pewnego $\varepsilon>0$,
a zatem $B\subseteq\varepsilon^{-1}V$.
\end{proof}

\begin{twierdzenie}\label{thm:compact-bounded}
Niech $X$ będzie przestrzenią liniowo-topologiczną.
Wówczas:
\begin{enumerate}[label=\textup{(\roman*)}]
 \item każdy skończony podzbiór $X$ jest ograniczony;
 \item każdy zwarty podzbiór $X$ jest ograniczony;
 \item jeżeli $T\colon X\to Y$ jest ciągłym odwzorowaniem liniowym
       między przestrzeniami liniowo-topologicznymi, to obraz
       każdego ograniczonego zbioru $B\subseteq X$ jest ograniczony w $Y$.
\end{enumerate}
\end{twierdzenie}
\begin{proof}
Wystarczy udowodnić (ii) i (iii), ponieważ każdy zbiór
skończony jest zwarty. Ustalmy zwarty zbiór $K\subseteq X$
oraz otoczenie $U$ zera. Wybierzmy otwarte, zrównoważone
otoczenie $V\subseteq U$. Zbiór $V$ jest pochłaniający,
więc $X=\bigcup_{n\in\N,\,n\geq1}nV$: jeżeli $x\in tV$ dla
pewnego $t>0$, to po wybraniu liczby całkowitej $n\geq t$
mamy $x\in nV$, ponieważ $V$ jest zrównoważony.
Rodzina $(nV)_{n\geq1}$ jest wstępująca, a każdy $nV$
jest otwarty. Zwartość $K$ daje skończone podpokrycie,
a w konsekwencji pewne $n$ z $K\subseteq nV\subseteq nU$.
To dowodzi (ii).

Dla (iii) niech $W\in\mathcal N_Y(0)$.
Z ciągłości $T$ zbiór $T^{-1}(W)$ jest otoczeniem zera w $X$.
Istnieje $t>0$ z $B\subseteq tT^{-1}(W)$.
Z liniowości $T(B)\subseteq tW$, co kończy dowód.
\end{proof}

\begin{przyklad}[Ograniczoność normowa]
Niech $(X,\|\cdot\|)$ będzie przestrzenią unormowaną.
Zbiór $B\subseteq X$ jest ograniczony w sensie
definicji~\ref{def:tvs-bounded} wtedy i tylko wtedy, gdy
$\sup_{x\in B}\|x\|<\infty$ (przyjmujemy supremum zbioru
pustego równe zero). Istotnie, wystarczy sprawdzać
otoczenia postaci $\{x:\|x\|<r\}$, $r>0$.
\end{przyklad}

\subsection{Przykłady przestrzeni nienormowalnych}

\begin{przyklad}[Przestrzeń wszystkich ciągów]\label{ex:product-tvs}
Niech $X=\K^{\N}$ będzie przestrzenią wszystkich ciągów
$x=(x_1,x_2,\ldots)$, wyposażoną w topologię produktową.
Bazę otoczeń zera tworzą zbiory
\[
 U(F,\varepsilon)=\big\{x\in\K^{\N}:|x_j|<\varepsilon_j
          \text{ dla }j\in F\big\},
\]
gdzie $F\subseteq\N$ jest skończony, a $\varepsilon_j>0$ dla
$j\in F$. Dodawanie i mnożenie przez skalar są ciągłe,
ponieważ ich złożenia ze wszystkimi rzutami współrzędnymi
są ciągłe. Zatem $X$ jest przestrzenią liniowo-topologiczną Hausdorffa.

Nie jest ona normowalna: żadne otoczenie zera nie jest ograniczone.
Istotnie, niech $U(F,\varepsilon)$ będzie bazowym otoczeniem
zera i wybierzmy $j\notin F$. Weźmy otoczenie
$W=\{x:|x_j|<1\}$. Dla każdego $t>0$ ciąg
$x=(t+1)e_j$ należy do $U(F,\varepsilon)$, ale nie należy
do $tW$. Stąd $U(F,\varepsilon)$ nie jest ograniczone,
a więc żadne większe otoczenie zera nie jest ograniczone.
Gdyby topologia pochodziła od normy, otwarta kula jednostkowa
byłaby ograniczonym otoczeniem zera, co prowadzi do sprzeczności.

Warto zauważyć, że $e_n\to0$ w topologii produktowej:
dla każdego ustalonego $j$ współrzędna $(e_n)_j$
jest równa zero dla wszystkich dostatecznie dużych $n$.
\end{przyklad}

\begin{przyklad}[Topologia zbieżności punktowej funkcji]
Niech $S$ będzie zbiorem, a $X\subseteq\K^S$ podprzestrzenią
liniową, np. przestrzenią wszystkich funkcji ciągłych na
$S=[0,1]$. Topologia zbieżności punktowej na $X$
jest topologią podprzestrzeni odziedziczoną z produktu
$\K^S$. Ciąg (lub ogólniej sieć) funkcji zbiega w tej
topologii wtedy i tylko wtedy, gdy zbiega w każdym punkcie $S$.
W szczególności topologia ta nie wymaga normy supremum.
\end{przyklad}

\begin{przyklad}[Nieciągły funkcjonał na przestrzeni unormowanej]
Rozważmy $X=\R[t]$, przestrzeń wielomianów rzeczywistych
na $[0,1]$, z normą supremum. Funkcjonał liniowy
$L(p)=p'(0)$ nie jest ciągły. Dla $p_n(t)=(1-t)^n$
mamy $\|p_n\|_\infty=1$, ale $|L(p_n)|=n$.
Gdyby $L$ był ciągły w zerze, z liniowości otrzymalibyśmy
$|L(p)|\leq C\|p\|_\infty$ dla pewnego $C>0$ i wszystkich
$p\in X$, co przeczy powyższemu rachunkowi.
\end{przyklad}

\subsection{Przykład zasadniczy: \texorpdfstring{$L^p$}{Lp} dla \texorpdfstring{$0<p<1$}{0 < p < 1}}

W tym przykładzie zakładamy znajomość elementarnych pojęć miary
Lebesgue'a. Ma on pokazać, dlaczego lokalna wypukłość będzie
później odgrywać kluczową rolę.

\begin{twierdzenie}[Zupełna przestrzeń liniowo-topologiczna
bez niezerowych funkcjonałów ciągłych]\label{thm:lp-trivial-dual}
Niech $0<p<1$, a $X=L^p([0,1],\lambda)$ oznacza przestrzeń
klas funkcji mierzalnych $f\colon[0,1]\to\K$ takich, że
$\int_{[0,1]}|f(x)|^p\,\mathrm d\lambda(x)<\infty$,
gdzie utożsamiamy funkcje równe prawie wszędzie.
Wówczas wzór
\[
 d(f,g)=\int_{[0,1]}|f(x)-g(x)|^p\,\mathrm d\lambda(x)
\]
definiuje zupełną metrykę na $X$ zgodną ze strukturą
przestrzeni liniowo-topologicznej. Co więcej, każdy
ciągły funkcjonał $\K$-liniowy $\Lambda\colon X\to\K$
jest równy zeru.
\end{twierdzenie}
\begin{proof}
Dla liczb rzeczywistych $a,b\geq0$ i $0<p<1$ zachodzi
$(a+b)^p\leq a^p+b^p$. W połączeniu z nierównością
$|f-h|\leq|f-g|+|g-h|$ daje to nierówność trójkąta dla $d$.
Pozostałe aksjomaty metryki wynikają z własności całki.
Dodawanie jest ciągłe na mocy
$d(f_1+g_1,f_2+g_2)\leq d(f_1,f_2)+d(g_1,g_2)$.
Dla $\alpha,\beta\in\K$ oraz $f,g\in X$ mamy
\[
 d(\alpha f,\beta g)
 \leq |\alpha|^p d(f,g)+|\alpha-\beta|^p d(g,0).
\]
Jest to oszacowanie wystarczające do łącznej ciągłości
mnożenia przez skalar.

Wykażemy zupełność. Niech $(f_n)$ będzie ciągiem Cauchy'ego
w metryce $d$; wybierzmy mierzalnych reprezentantów jego wyrazów.
Możemy wybrać podciąg $(f_{n_k})$ tak, aby
$\sum_{k=1}^\infty d(f_{n_{k+1}},f_{n_k})<\infty$.
Z twierdzenia o monotonicznej zbieżności wynika, że
\[
 \sum_{k=1}^\infty|f_{n_{k+1}}(x)-f_{n_k}(x)|^p<\infty
 \quad\text{dla prawie każdego }x\in[0,1].
\]
Ponieważ $p<1$, zbieżność szeregu dodatnich liczb
$\sum_k a_k^p$ implikuje zbieżność $\sum_k a_k$
(wszystkie dostatecznie późne $a_k$ należą do $[0,1]$,
a wtedy $a_k\leq a_k^p$). Zatem szereg różnic
$\sum_k(f_{n_{k+1}}-f_{n_k})$ zbiega bezwzględnie
prawie wszędzie. Oznaczmy jego granicę powiększoną o
$f_{n_1}$ przez $f$. Z nierówności $(\sum a_k)^p\leq\sum a_k^p$
i lematu Fatou dostajemy
\[
 d(f_{n_k},f)
 \leq \sum_{j=k}^\infty d(f_{n_{j+1}},f_{n_j})\longrightarrow0.
\]
W szczególności $f\in L^p$ (bo $|f|^p\leq|f-f_{n_1}|^p+|f_{n_1}|^p$),
a ciąg Cauchy'ego $(f_n)$ także zbiega do $f$.

Pozostało zbadać przestrzeń ciągłych funkcjonałów liniowych. Niech
$\Lambda\colon X\to\K$
będzie ciągłym funkcjonałem liniowym. Ciągłość w zerze
daje $\delta>0$ takie, że
$d(f,0)<\delta^p$ implikuje $|\Lambda(f)|<1$.
Z jednorodności $d(af,0)=|a|^p d(f,0)$ wynika istnienie
$C>0$ takie, że
\[
 |\Lambda(f)|\leq C\Big(\int_{[0,1]}|f|^p\,\mathrm d\lambda\Big)^{1/p}
 \qquad(f\in X).
\]
Rzeczywiście, dla $f\ne0$ wystarczy zastosować poprzedni
warunek do $a f$ z $|a|=\delta/(2d(f,0)^{1/p})$.

Ustalmy dowolny zbiór mierzalny $A\subseteq[0,1]$.
Dla każdego $n\geq1$ można rozbić $A$ na rozłączne
zbiory mierzalne $A_1,\ldots,A_n$ o równej mierze
$\lambda(A_i)=\lambda(A)/n$; wynika to z bezatomowości
miary Lebesgue'a. Stąd
\[
 |\Lambda(\chi_A)|
 \leq\sum_{i=1}^n|\Lambda(\chi_{A_i})|
 \leq C\lambda(A)^{1/p}n^{1-1/p}\longrightarrow0,
\]
ponieważ $1-1/p<0$. Funkcjonał $\Lambda$ znika więc
na funkcjach charakterystycznych, a z liniowości również
na funkcjach prostych. Uzasadnijmy również używaną gęstość funkcji prostych.
Dla $f\in L^p([0,1],\lambda)$ obcięcie
$f_N=f\chi_{\{|f|\leq N\}}$ spełnia
$d(f_N,f)\to0$ na mocy zbieżności zmajoryzowanej.
Funkcja $f_N$ jest ograniczona; dzieląc kwadrat
$[-N,N]^2$ (dla $\K=\C$) albo przedział $[-N,N]$
(dla $\K=\R$) na dostatecznie małe komórki i zastępując
wartości $f_N$ reprezentantami komórek, otrzymujemy
mierzalne funkcje proste $s_k$ z
$|s_k-f_N|\leq1/k$ prawie wszędzie. Wtedy
$d(s_k,f_N)\leq k^{-p}\lambda([0,1])\to0$.
Funkcje proste są więc gęste w $L^p$; z ciągłości
$\Lambda=0$ na całym $X$.
\end{proof}

\begin{uwaga}\label{rem:local-convexity-purpose}
Na przestrzeni $L^p([0,1])$ z $0<p<1$ nie można oddzielić
niezerowego wektora od zera za pomocą ciągłego funkcjonału
liniowego: przestrzeń ciągłych funkcjonałów liniowych jest trywialna.
W dalszych
rozdziałach udowodnimy, że w każdej przestrzeni lokalnie
wypukłej Hausdorffa ciągłe funkcjonały liniowe oddzielają
punkty. Zatem powyższa przestrzeń \emph{nie jest lokalnie wypukła},
choć jest metryzowalna i zupełna.
\end{uwaga}

\subsection{Co zostało ustalone, a co wymaga dodatkowych założeń?}

W dowolnej przestrzeni liniowo-topologicznej możemy mówić o
ciągłości operatorów liniowych, zbiorach ograniczonych, otoczeniach
zera, podprzestrzeniach i zbieżności. Nie potrzebowaliśmy do tego
ani normy, ani iloczynu skalarnego, ani zupełności.

Nie możemy natomiast bez dalszych założeń zakładać, że istnieją
niezerowe ciągłe funkcjonały liniowe. Twierdzenie
\ref{thm:lp-trivial-dual} pokazuje, że jest to rzeczywista
przeszkoda. W kolejnych rozdziałach zbadamy podprzestrzenie, produkty
i ilorazy, następnie zbieżność sieci, ograniczoność oraz
metryzowalność. Później wprowadzimy seminormy i przestrzenie
lokalnie wypukłe.

\begin{uwaga}[Dalszy kierunek skryptu]
Pochodne Gâteaux i Frécheta nie wymagają teorii dystrybucji.
Wprowadzimy je w osobnym bloku dotyczącym przestrzeni Banacha
i Hilberta, a następnie wykorzystamy w twierdzeniu o funkcji
uwikłanej, zasadach wariacyjnych i rachunku na przestrzeniach
funkcyjnych. Dla przestrzeni zespolonych wyraźnie odróżnimy
różniczkowalność w sensie rzeczywistym od pochodnej
zespolono-liniowej.
\end{uwaga}

% Wersja źródłowa: audyt III, jednostka 02
\bigskip
\setcounter{section}{1}
\section{Podprzestrzenie, produkty, ilorazy i topologie uniwersalne}\label{sec:konstrukcje}

\subsection{Dlaczego potrzebujemy konstrukcji nowych przestrzeni?}

Z jednej przestrzeni liniowo-topologicznej (TVS) powstają naturalnie nowe:
zbiór funkcji spełniających liniowy warunek, rodzina funkcji zależnych
od wielu parametrów, przestrzeń funkcji rozpatrywanych z dokładnością
do pewnej relacji równoważności. Odpowiadają im odpowiednio
\emph{podprzestrzenie}, \emph{produkty} oraz \emph{ilorazy}.
Istotą konstrukcji jest nie tylko zdefiniowanie zbioru i działań,
lecz także dobranie topologii tak, aby właściwe odwzorowania były
ciągłe. W dalszych rozdziałach ta idea posłuży do budowy
przestrzeni funkcji testowych i dystrybucji.

Przyjmujemy konwencje \hyperref[sec:tvs]{rozdziału 1}: $\K\in\{\R,\C\}$,
a wszystkie przestrzenie liniowo-topologiczne są Hausdorffa,
chyba że wyraźnie zaznaczymy inaczej. Pojęcie \emph{topologii
indukowanej} i \emph{topologii produktowej} rozumiemy w zwykłym sensie
topologii ogólnej. W tej części przypomnimy ich precyzyjne własności.

\subsection{Podprzestrzenie liniowo-topologiczne}

\begin{twierdzenie}[Topologia podprzestrzeni]\label{thm:subspace-tvs}
Niech $X$ będzie przestrzenią liniowo-topologiczną nad $\K$,
a $M\subseteq X$ podprzestrzenią liniową. Z topologią indukowaną
$\tau_M=\{O\cap M:O\subseteq X\text{ jest otwarty}\}$
przestrzeń $M$ jest przestrzenią liniowo-topologiczną nad $\K$.
Jeżeli $X$ jest Hausdorffa, to również $M$ jest Hausdorffa.
Otoczenia zera w $M$ mają bazę postaci $U\cap M$, gdzie
$U$ przebiega otoczenia zera w $X$.
\end{twierdzenie}
\begin{proof}
Ograniczenie dodawania $+_X$ do $M\times M$ przyjmuje wartości
w $M$, a topologia na $M\times M$ jest topologią podprzestrzeni
$X\times X$. Jeżeli $O\cap M$ jest otwarty w $M$, to
\[
 (+_M)^{-1}(O\cap M)=(+_X)^{-1}(O)\cap(M\times M),
\]
więc dodawanie w $M$ jest ciągłe. Identyczny argument dla
$\K\times M\subseteq\K\times X$ daje ciągłość mnożenia
przez skalar. Dziedziczenie własności Hausdorffa jest własnością
każdej podprzestrzeni topologicznej. Ostatnie zdanie wynika
wprost z definicji topologii indukowanej i faktu, że $0\in M$.
\end{proof}

\begin{lemat}[Domknięcie podprzestrzeni]\label{lem:closure-subspace}
Niech $M$ będzie podprzestrzenią liniową przestrzeni
liniowo-topologicznej $X$, a $\mathcal N_X(0)$ rodziną
otoczeń zera w $X$. Wówczas $\overline M$ jest podprzestrzenią
liniową oraz
\[
 \overline M=\bigcap_{U\in\mathcal N_X(0)}(M+U),
 \qquad M+U=\{m+u:m\in M,\ u\in U\}.
\]
\end{lemat}
\begin{proof}
Jeżeli $x,y\in\overline M$, to dla dowolnego otoczenia $U$
zera wybieramy symetryczne otoczenie $V$ z $V+V\subseteq U$.
Istnieją $m_x,m_y\in M$ takie, że $x-m_x\in V$ oraz
$y-m_y\in V$. Stąd $(x+y)-(m_x+m_y)\in U$,
więc $x+y\in\overline M$. Dla $\lambda\ne0$ ciągłość
mnożenia przez skalar (a nawet fakt, że jest homeomorfizmem)
daje $\lambda\overline M=\overline{\lambda M}=\overline M$;
przypadek $\lambda=0$ jest oczywisty.

Jeżeli $x\in\overline M$, to dla każdego $U$ możemy wybrać
symetryczne $V\subseteq U$; przecięcie $(x+V)\cap M$
jest niepuste, więc $x=m-v\in M+U$ dla odpowiednich
$m\in M$, $v\in V$. Odwrotnie, niech $x\in\bigcap_U(M+U)$.
Dla otoczenia $x+V$ punktu $x$ wybieramy symetryczne $V$;
z założenia $x=m+v$ z $m\in M$, $v\in V$,
a więc $m=x-v\in x+V$. Każde otoczenie $x$ spotyka $M$,
czyli $x\in\overline M$.
\end{proof}

\begin{przyklad}[Dwie topologie na wielomianach]
Niech $P=\K[t]$, a $X=C([0,1],\K)$ z normą supremum.
Topologia indukowana na $P$ jest unormowana, lecz $P$
nie jest domknięte w $X$: wielomiany są gęste w $X$
(z twierdzenia Weierstrassa o aproksymacji). W szczególności
zupełność $X$ nie przechodzi na dowolną podprzestrzeń.
W przyszłości udowodnimy, że podprzestrzeń przestrzeni Banacha
jest zupełna w normie indukowanej wtedy i tylko wtedy, gdy
jest domknięta.
\end{przyklad}

\begin{stwierdzenie}[Uniwersalna własność podprzestrzeni]
Niech $M\subseteq X$ będzie podprzestrzenią liniową przestrzeni
liniowo-topologicznej $X$, $i:M\hookrightarrow X$ naturalnym
włożeniem, $Z$ dowolną przestrzenią topologiczną oraz
$f:Z\to M$ odwzorowaniem. Wówczas $f$ jest ciągłe
wtedy i tylko wtedy, gdy $i\circ f:Z\to X$ jest ciągłe.
\end{stwierdzenie}
\begin{proof}
Konieczność jest natychmiastowa, gdyż $i$ jest ciągłe.
Jeśli $i\circ f$ jest ciągłe, to dla otwartego $O\cap M$
w $M$ mamy $f^{-1}(O\cap M)=(i\circ f)^{-1}(O)$,
więc $f$ jest ciągłe.
\end{proof}

\subsection{Produkty dowolnych rodzin przestrzeni}

\begin{definicja}[Produkt przestrzeni liniowo-topologicznych]
Niech $I$ będzie zbiorem indeksów, a $(X_i)_{i\in I}$ rodziną
przestrzeni liniowo-topologicznych nad $\K$. Na przestrzeni
liniowej $X=\prod_{i\in I}X_i$ działań współrzędnych
wprowadzamy \emph{topologię produktową}: jej bazę tworzą zbiory
$\prod_{i\in I}O_i$, gdzie każde $O_i\subseteq X_i$ jest
otwarte i $O_i=X_i$ dla wszystkich poza skończenie wieloma
indeksami. Rzut na $i$-tą współrzędną oznaczamy przez
$\pi_i:X\to X_i$.
\end{definicja}

\begin{twierdzenie}[Produkt i jego własność uniwersalna]\label{thm:product-tvs}
Niech $(X_i)_{i\in I}$ będzie rodziną przestrzeni
liniowo-topologicznych Hausdorffa nad $\K$.
Wówczas $X=\prod_{i\in I}X_i$ z topologią produktową
jest przestrzenią liniowo-topologiczną Hausdorffa, a każdy
$\pi_i$ jest ciągłym, liniowym i otwartym odwzorowaniem na $X_i$.
Dla dowolnej przestrzeni topologicznej $Z$ odwzorowanie
$f:Z\to X$ jest ciągłe wtedy i tylko wtedy, gdy
$\pi_i\circ f:Z\to X_i$ jest ciągłe dla każdego $i\in I$.
\end{twierdzenie}
\begin{proof}
Rzuty są ciągłe z definicji topologii produktowej. Dla zbioru
bazowego $\prod_i O_i$ zachodzi $\pi_j(\prod_i O_i)=O_j$,
ponieważ wszystkie pozostałe $O_i$ są niepuste, o ile
rozpatrujemy niepusty zbiór bazowy; rzut zbioru pustego
jest pusty. Zatem rzuty są otwarte.

Dodawanie $a:X\times X\to X$ jest ciągłe, bo każde
$\pi_i\circ a:(x,y)\mapsto x_i+y_i$ jest ciągłe
(jest to złożenie rzutowań i dodawania w $X_i$).
To samo rozumowanie stosuje się do mnożenia skalarnego
$m:\K\times X\to X$. Przestrzeń $X$ jest Hausdorffa,
ponieważ dwa różne elementy różnią się na pewnej współrzędnej,
a tę można rozdzielić otoczeniami w odpowiednim $X_i$.

Jeśli wszystkie $f_i=\pi_i\circ f$ są ciągłe,
to dla podzbioru bazowego $\prod_iO_i$ mamy
$f^{-1}(\prod_iO_i)=\bigcap_{i\in F}f_i^{-1}(O_i)$,
gdzie $F=\{i:O_i\ne X_i\}$ jest skończony;
jest to zbiór otwarty. Implikacja odwrotna wynika
ze złożenia odwzorowań ciągłych.
\end{proof}

\begin{uwaga}[Dlaczego tylko skończenie wiele warunków?]
W produkcie nieskończonym wymaganie, aby wszystkie
współrzędne jednocześnie należały do wybranych małych otoczeń,
może zniszczyć ciągłość działań. Na przykład w $\K^{\N}$
zbiory $\{x:|x_n|<1\text{ dla każdego }n\}$ nie są otwarte
w topologii produktowej. Ograniczanie skończenie wielu
współrzędnych jest kluczowe dla uniwersalnej własności produktu.
\end{uwaga}

\begin{twierdzenie}[Ograniczoność w produkcie]\label{thm:product-bounded}
Niech $(X_i)_{i\in I}$ będzie rodziną przestrzeni
liniowo-topologicznych i niech $B\subseteq\prod_{i\in I}X_i$.
Wówczas $B$ jest ograniczony w sensie von Neumanna wtedy
i tylko wtedy, gdy $\pi_i(B)$ jest ograniczony w $X_i$
dla każdego $i\in I$.
\end{twierdzenie}
\begin{proof}
W jedną stronę stosujemy rezultat z \hyperref[sec:tvs]{rozdziału 1}: ciągły operator liniowy
przenosi zbiory ograniczone na zbiory ograniczone. Stosujemy go do ciągłych rzutów $\pi_i$.
Odwrotnie, niech $U$ będzie otoczeniem zera w produkcie.
Możemy wybrać bazowe otoczenie $V=\prod_i V_i\subseteq U$,
gdzie $V_i=X_i$ poza skończonym zbiorem $F$, a dla $i\in F$
przyjmijmy $V_i$ zrównoważone. Dla każdego $i\in F$ istnieje
$t_i>0$ takie, że $\pi_i(B)\subseteq t_iV_i$.
Jeśli $F$ jest niepusty, ustalmy $t\geq\max(1,\max_{i\in F}t_i)$.
Zrównoważenie $V_i$ daje $t_iV_i\subseteq tV_i$, więc
$B\subseteq tV\subseteq tU$. Jeżeli $F$ jest pusty,
wystarczy przyjąć $t=1$.
\end{proof}

\begin{przyklad}[Dwa rodzaje zbieżności funkcji]
Niech $X=C([0,1],\K)$. Norma supremum wyznacza topologię
zbieżności jednostajnej; topologię zbieżności punktowej
otrzymujemy przez zanurzenie $X\subseteq\K^{[0,1]}$.
Identyczność z topologii normowej do punktowej jest ciągła.
Nie jest jednak homeomorfizmem: funkcje
$f_n(t)=\max\{0,1-n^2|t-1/n|\}$, $n\geq2$, są ciągłe,
zbiegają punktowo do zera, ale $\|f_n\|_\infty=1$.
Istotnie, dla każdego ustalonego $t\in[0,1]$ nośnik $f_n$
nie zawiera $t$ dla wszystkich dostatecznie dużych $n$.
\end{przyklad}

\begin{twierdzenie}[Funkcjonały na produkcie przestrzeni skalarnych]\label{thm:product-dual}
Niech $I$ będzie dowolnym zbiorem, a $X=\K^I$ wyposażmy
w topologię produktową. Każdy ciągły funkcjonał liniowy
$L:X\to\K$ ma postać
\[
 L(x)=\sum_{i\in F}a_i x_i
\]
przy pewnym skończonym zbiorze $F\subseteq I$ oraz
skalarach $a_i\in\K$. Odwrotnie, każdy taki funkcjonał
jest ciągły.
\end{twierdzenie}
\begin{proof}
Z ciągłości $L$ w zerze istnieje bazowe otoczenie zera
$U=\{x:|x_i|<\varepsilon_i\text{ dla }i\in F\}$
takie, że $|L(x)|<1$ dla $x\in U$.
Jeśli $x_i=0$ dla każdego $i\in F$, to $\lambda x\in U$
dla dowolnego $\lambda\in\K$. Stąd
$|\lambda L(x)|<1$ dla wszystkich $\lambda$, więc $L(x)=0$.
Dla dowolnego $x\in X$ odejmijmy skończoną sumę
$\sum_{i\in F}x_i e_i$, gdzie $e_i$ oznacza funkcję
równą $1$ na $i$ i $0$ poza nim. Pozostały wektor
ma zera na $F$, dlatego $L(x)=\sum_{i\in F}L(e_i)x_i$.
Odwrotność wynika z ciągłości rzutów i działań skalarnych.
\end{proof}

\begin{uwaga}
Twierdzenie~\ref{thm:product-dual} jest pierwszą wyraźną
manifestacją tego, że \emph{przestrzeń ciągłych funkcjonałów liniowych}
może być znacznie mniejsza od przestrzeni funkcjonałów algebraicznych. Nie należy go
mylić z analogiczną własnością dowolnych produktów:
ciągły operator $\K^{\N}\to\K^{\N}$ może zależeć
od nieskończenie wielu współrzędnych (np. identyczność).
\end{uwaga}

\subsection{Przestrzeń ilorazowa i topologia ilorazowa}

\begin{definicja}[Iloraz przez podprzestrzeń]
Niech $X$ będzie przestrzenią liniowo-topologiczną nad $\K$,
a $M\subseteq X$ podprzestrzenią liniową. Na przestrzeni
wektorowej $X/M=\{x+M:x\in X\}$ wprowadzamy
\emph{topologię ilorazową}: zbiór $W\subseteq X/M$
jest otwarty wtedy i tylko wtedy, gdy $q^{-1}(W)$
jest otwarty w $X$, gdzie $q:X\to X/M$, $q(x)=x+M$.
\end{definicja}

\begin{twierdzenie}[Iloraz jako TVS]\label{thm:quotient-tvs}
Niech $X$ będzie dowolną przestrzenią liniowo-topologiczną,
niekoniecznie Hausdorffa, a $M\subseteq X$ podprzestrzenią.
Wówczas:
\begin{enumerate}[label=\textup{(\roman*)}]
\item odwzorowanie $q:X\to X/M$ jest liniowe, ciągłe,
      surjektywne i otwarte;
\item topologia ilorazowa czyni $X/M$ przestrzenią
      liniowo-topologiczną;
\item $X/M$ jest Hausdorffa wtedy i tylko wtedy,
      gdy $M$ jest domknięta w $X$.
\end{enumerate}
\end{twierdzenie}
\begin{proof}
Liniowość i surjektywność $q$ wynikają z konstrukcji
ilorazu, a ciągłość wynika z definicji topologii ilorazowej.
Jeżeli $O\subseteq X$ jest otwarty, to
\[
 q^{-1}(q(O))=O+M=\bigcup_{m\in M}(O+m)
\]
jest otwarty, ponieważ każde przesunięcie $O+m$
jest otwarte. Zatem $q(O)$ jest otwarty w $X/M$.

Wykażemy ciągłość dodawania w ilorazie. Odwzorowanie
$q\times q:X\times X\to(X/M)\times(X/M)$ jest ciągłą,
otwartą surjekcją: otwartość wynika z faktu, że obrazy
zbiorów bazowych $O_1\times O_2$ są produktami
otwartych $q(O_1)\times q(O_2)$. Każda otwarta surjekcja
ciągła jest odwzorowaniem ilorazowym. Mamy równość
$+_{X/M}\circ(q\times q)=q\circ+_X$, a prawa strona
jest ciągła. Z własności odwzorowania ilorazowego wynika
ciągłość $+_{X/M}$. Podobnie odwzorowanie
$\id_\K\times q:\K\times X\to\K\times(X/M)$ jest
otwartą surjekcją, więc z równości
$m_{X/M}\circ(\id_\K\times q)=q\circ m_X$
wynika ciągłość mnożenia skalarnego. Tym samym $X/M$
jest TVS, także gdy $M$ nie jest domknięta.

W dowolnej TVS singleton $\{0\}$ jest domknięty
wtedy i tylko wtedy, gdy przestrzeń jest Hausdorffa
(por. twierdzenie o domknięciu zera z \hyperref[sec:tvs]{rozdziału 1}).
Ponieważ $q^{-1}(\{0+M\})=M$, definicja topologii
ilorazowej pokazuje, że $\{0+M\}$ jest domknięty
w $X/M$ wtedy i tylko wtedy, gdy $M$ jest domknięta w $X$.
\end{proof}

\begin{uwaga}[Otoczenia w ilorazie]\label{rem:q-neighborhood}
Ponieważ $q$ jest otwarty, zbiory $q(U)$, gdzie
$U\in\mathcal N_X(0)$, tworzą bazę otoczeń zera
w $X/M$. Dokładniej: dla każdego otoczenia $W$
punktu $0+M$ zbiór $q^{-1}(W)$ jest otoczeniem zera,
a $q(q^{-1}(W))=W$; z drugiej strony każdy $q(U)$
jest otoczeniem zera, bo zawiera $q(O)$ dla pewnego
otwartego otoczenia $O\subseteq U$ zera.
\end{uwaga}

\begin{center}
\begin{tikzpicture}[x=0.95cm,y=0.82cm,>=Latex]
\draw[->] (-0.5,0)--(5.35,0) node[right] {$x$};
\draw[->] (0,-0.55)--(0,3.75) node[above] {$y$};
\foreach \yy in {0.8,1.7,2.65} {
  \draw[blue!55!black,thick] (0.2,\yy)--(4.7,\yy);
}
\filldraw[blue!55!black] (1.0,1.7) circle (1.7pt);
\filldraw[blue!55!black] (3.3,1.7) circle (1.7pt);
\node[anchor=west] at (4.75,1.7) {$x_0+M$};
\node[below] at (2.4,-0.10) {$M=\R\times\{0\}$};
\draw[->,very thick] (5.15,2.8) -- (6.4,2.8)
  node[midway,above] {$q$};
\draw[->] (7.2,-0.5)--(7.2,3.65);
\foreach \yy in {0.8,1.7,2.65} {
 \filldraw[blue!55!black] (7.2,\yy) circle (2.1pt);
}
\node[right] at (7.4,1.7) {$[x_0]$};
\node[below] at (7.2,-0.65) {$\R^2/M\simeq\R$};
\end{tikzpicture}

{\small Rysunek 2.1. Iloraz $\R^2/M$, gdzie $M$ jest osią poziomą.
Każda pozioma prosta jest jedną klasą równoważności
(prawa strona przedstawia zbiór klas).}
\end{center}

\begin{twierdzenie}[Własność uniwersalna ilorazu]\label{thm:quotient-universal}
Niech $X$ będzie TVS nad $\K$, $M\subseteq X$ podprzestrzenią,
a $q:X\to X/M$ rzutowaniem. Niech $Z$ będzie dowolną
przestrzenią topologiczną, a $g:X/M\to Z$ odwzorowaniem.
Wówczas $g$ jest ciągłe wtedy i tylko wtedy,
gdy $g\circ q$ jest ciągłe. Ponadto jeśli $Z$ jest TVS,
a $T:X\to Z$ jest liniowe i $M\subseteq\ker T$, to
istnieje dokładnie jedno liniowe $\widetilde T:X/M\to Z$
takie, że $T=\widetilde T\circ q$; jest ono ciągłe
dokładnie wtedy, gdy $T$ jest ciągłe.
\end{twierdzenie}
\begin{proof}
Pierwsza równoważność wynika z definicji topologii ilorazowej:
$g^{-1}(O)$ jest otwarty w $X/M$ dokładnie wtedy,
gdy $(g\circ q)^{-1}(O)$ jest otwarty w $X$, dla każdego
otwartego $O\subseteq Z$. Jeśli $M\subseteq\ker T$,
to wzór $\widetilde T(x+M)=T(x)$ jest dobrze określony:
z równości $x+M=y+M$ wynika $x-y\in M$, czyli $T(x)=T(y)$.
Wzór jest liniowy, a surjektywność $q$ zapewnia jednoznaczność.
Równoważność ciągłości stosujemy do $g=\widetilde T$.
\end{proof}

\begin{wniosek}[Pierwsze twierdzenie o izomorfizmie --- uwaga topologiczna]
Niech $T:X\to Y$ będzie ciągłym, liniowym odwzorowaniem
między przestrzeniami liniowo-topologicznymi Hausdorffa.
Wówczas odwzorowanie
$\widetilde T:X/\ker T\to T(X)$, $x+\ker T\mapsto T(x)$,
jest ciągłą bijekcją liniową, gdzie $T(X)$ ma topologię
indukowaną z $Y$. Jest izomorfizmem liniowo-topologicznym
wtedy i tylko wtedy, gdy $T:X\to T(X)$ jest odwzorowaniem
otwartym (przy topologii podprzestrzeni na obrazie).
\end{wniosek}
\begin{proof}
Jądro $T$ jest domknięte, ponieważ $T$ jest ciągłe,
a $Y$ jest Hausdorffa; dlatego iloraz jest Hausdorffa.
Istnienie i ciągłość $\widetilde T$ daje
\ref{thm:quotient-universal}; jest ona różnowartościowa,
bo $\ker T$ jest dokładnie jądrem $T$.
Jeśli $T$ jest otwarte na swój obraz, to
$\widetilde T(q(O))=T(O)$ jest otwarty dla każdego otwartego
$O\subseteq X$, a z otwartości $q$ i opisu zbiorów
otwartych w ilorazie wynika otwartość $\widetilde T$.
Odwrotnie, jeżeli $\widetilde T$ jest homeomorfizmem,
to $T=\widetilde T\circ q$ jest otwarte na obraz jako
złożenie dwóch odwzorowań otwartych.
\end{proof}

\begin{przyklad}[Dlaczego ciągła bijekcja nie wystarcza]
Niech $X=C([0,1],\R)$ z normą supremum, a $Y$ to ta sama
przestrzeń liniowa z normą $\|f\|_1=\int_{[0,1]}|f(t)|\,\mathrm d\lambda(t)$.
Identyczność $T:X\to Y$ jest ciągłą bijekcją liniową,
ale nie jest otwarta. Jej jądro jest zerowe, a mimo to
$X/\ker T\simeq X$ nie jest topologicznie izomorficzne z $Y$
poprzez odwzorowanie indukowane przez $T$. Dowód braku
ciągłości odwrotności podaliśmy w \hyperref[sec:tvs]{rozdziale 1}.
\end{przyklad}

\begin{przyklad}[Iloraz przez jądro wartościowania]
Niech $X=C([0,1],\K)$ z normą supremum, a
$E_0(f)=f(0)$. Jądro $M=\{f:f(0)=0\}$ jest domknięte.
Rzut $X/M\to\K$, $[f]\mapsto f(0)$, jest izomorfizmem
liniowo-topologicznym. Istotnie, dla $a\in\K$ stała
funkcja $s(a)(t)=a$ jest ciągłym liniowym prawym odwrotnym
do $E_0$, a odwrotne odwzorowanie na iloraz ma postać
$a\mapsto q(s(a))$. Jest więc ciągłe jako złożenie ciągłych.
\end{przyklad}

\subsection{Norma ilorazowa: ważny przypadek szczególny}

\begin{twierdzenie}[Norma w ilorazie]\label{thm:quotient-norm}
Niech $(X,\|\cdot\|)$ będzie przestrzenią unormowaną,
a $M\subseteq X$ domkniętą podprzestrzenią liniową. Wzór
\[
 \|x+M\|_{X/M}:=\inf_{m\in M}\|x-m\|=\operatorname{dist}(x,M)
\]
definiuje normę na $X/M$ generującą topologię ilorazową.
Dla każdego $r>0$ zachodzi
$q(B_X(0,r))=B_{X/M}(0,r)$.
Jeżeli dodatkowo $X$ jest przestrzenią Banacha,
to $X/M$ jest przestrzenią Banacha.
\end{twierdzenie}
\begin{proof}
Zastąpienie $x$ przez $x+m_0$, $m_0\in M$, nie zmienia
zbioru wartości $\|x+m_0-m\|$ dla $m\in M$, więc wzór
jest dobrze określony. Jednorodność wynika z liniowości $M$;
dla $\lambda=0$ obie strony są równe zeru. Dla dowolnych
$m_1,m_2\in M$ nierówność trójkąta w $X$ daje
$\|(x+y)-(m_1+m_2)\|\leq\|x-m_1\|+\|y-m_2\|$.
Biorąc infima, otrzymujemy nierówność trójkąta w ilorazie.
Jeśli $\operatorname{dist}(x,M)=0$, to $x\in\overline M=M$,
więc $x+M=0+M$. Zatem mamy normę.

Równość kul zachodzi, bo $\|x+M\|_{X/M}<r$ dokładnie
wtedy, gdy istnieje $m\in M$ z $\|x-m\|<r$.
Wobec otwartości $q$ wynika z tego, że obrazy kul
są otwarte. Dla dowolnego zbioru $W\subseteq X/M$
otwartego w topologii normowej jego przeciwobraz
$q^{-1}(W)$ jest otwarty, gdyż
$\|q(x)-q(y)\|_{X/M}\leq\|x-y\|$.
W drugą stronę, jeśli $W$ jest otwarty w topologii ilorazowej,
to dla każdego $[x]\in W$ istnieje $r>0$ takie, że
$B_X(x,r)\subseteq q^{-1}(W)$; wtedy
$B_{X/M}([x],r)=q(B_X(x,r))\subseteq W$.
Topologie są więc identyczne.

Niech teraz $X$ będzie zupełne, a $(u_n)$ ciągiem Cauchy'ego
w $X/M$. Wybierzmy podciąg $(u_{n_k})$ taki, że
$\|u_{n_{k+1}}-u_{n_k}\|_{X/M}<2^{-k-1}$.
Ustalmy przedstawiciela $x_1$ klasy $u_{n_1}$.
Indukcyjnie wybieramy przedstawiciela $x_{k+1}$ klasy
$u_{n_{k+1}}$ z $\|x_{k+1}-x_k\|<2^{-k}$.
Jest to możliwe, gdyż definicja normy ilorazowej
pozwala, po ustaleniu $x_k$, dobrać przedstawiciela
kolejnej klasy tak, aby odległość była dowolnie bliska
normie różnicy klas. Szereg odległości jest zbieżny,
więc $(x_k)$ jest ciągiem Cauchy'ego w $X$ i zbiega
do pewnego $x\in X$. Z ciągłości $q$ mamy
$u_{n_k}=q(x_k)\to q(x)$. Ponieważ cały $(u_n)$
jest ciągiem Cauchy'ego, również zbiega do $q(x)$.
\end{proof}

\begin{uwaga}
Gdy $M$ nie jest domknięta, ten sam wzór definiuje
jedynie seminormę: $\|x+M\|_{X/M}=0$ może zachodzić
dla niezerowej klasy. W szczególności taki wzór nie
może generować topologii Hausdorffa na $X/M$.
\end{uwaga}

\subsection{Topologie początkowe: najoszczędniejszy wybór topologii}

\begin{definicja}[Topologia początkowa]
Niech $S$ będzie zbiorem, a $f_i:S\to Y_i$, $i\in I$,
rodziną odwzorowań do przestrzeni topologicznych $Y_i$.
\emph{Topologią początkową} na $S$ względem $(f_i)$
nazywamy najmniejszą topologię, względem której wszystkie
$f_i$ są ciągłe. Jest ona generowana przez podbazę
$\{f_i^{-1}(O):i\in I,\ O\subseteq Y_i\text{ otwarty}\}$.
\end{definicja}

\begin{twierdzenie}[Początkowa topologia liniowa]\label{thm:initial-tvs}
Niech $X$ będzie przestrzenią liniową nad $\K$, $(Y_i)_{i\in I}$
rodziną przestrzeni liniowo-topologicznych nad $\K$,
a $T_i:X\to Y_i$ odwzorowaniami liniowymi.
Topologia początkowa względem $(T_i)$ czyni $X$
przestrzenią liniowo-topologiczną, niekoniecznie Hausdorffa.
Dla dowolnej przestrzeni topologicznej $Z$ odwzorowanie
$f:Z\to X$ jest ciągłe wtedy i tylko wtedy, gdy
$T_i\circ f$ jest ciągłe dla każdego $i\in I$.
Jeśli wszystkie $Y_i$ są Hausdorffa, to $X$ jest Hausdorffa
dokładnie wtedy, gdy $\bigcap_{i\in I}\ker T_i=\{0\}$.
\end{twierdzenie}
\begin{proof}
Ciągłość odwzorowań $T_i$ wynika z definicji topologii.
Ponieważ $T_i\circ+_X=+_{Y_i}\circ(T_i\times T_i)$,
a odwzorowania $T_i\times T_i$ są ciągłe względem
topologii początkowej na $X$, wszystkie złożenia
$T_i\circ+_X$ są ciągłe. Z uniwersalnej własności
topologii początkowej wynika ciągłość dodawania na $X$.
Tak samo $T_i\circ m_X=m_{Y_i}\circ(\id_\K\times T_i)$
daje ciągłość mnożenia skalarnego. Równoważność dotycząca
$f$ wynika wprost z preobrazów podbazowych zbiorów otwartych.

Jeśli $x\ne y$ i istnieje $i$ z $T_i(x)\ne T_i(y)$,
to dzięki własności Hausdorffa $Y_i$ istnieją rozłączne
otoczenia tych obrazów, których przeciwobrazy rozdzielają
$x$ i $y$. Jeśli natomiast $0\ne z\in\bigcap_i\ker T_i$,
to każde bazowe otoczenie zera zawiera $z$, a więc
$X$ nie jest Hausdorffa na mocy kryterium Hausdorffa
z \hyperref[sec:tvs]{rozdziału 1}.
\end{proof}

\begin{przyklad}[Produkty i podprzestrzenie jako przypadki topologii początkowej]
Topologia produktowa na $\prod_iX_i$ jest topologią
początkową względem rzutów $\pi_i$.
Topologia podprzestrzeni na $M\subseteq X$ jest
początkową względem włożenia $M\hookrightarrow X$.
W szczególności naturalna topologia na $C([0,1])$
zbieżności punktowej pochodzi od rodziny funkcjonałów
wartościowania $f\mapsto f(t)$, $t\in[0,1]$.
\end{przyklad}

\subsection{Topologie końcowe: przypadek delikatniejszy}

\begin{definicja}[Topologia końcowa zbiorów]
Niech $S$ będzie zbiorem, a $f_i:Y_i\to S$, $i\in I$,
rodziną odwzorowań z przestrzeni topologicznych.
\emph{Topologią końcową} na $S$ nazywamy najsilniejszą
topologię, przy której wszystkie $f_i$ są ciągłe:
\[
 \tau_{\mathrm{fin}}=\{O\subseteq S:
    f_i^{-1}(O)\text{ jest otwarty w }Y_i
    \text{ dla wszystkich }i\in I\}.
\]
Dla przestrzeni ilorazowej $X/M$ otrzymujemy właśnie
tę topologię przy jednym odwzorowaniu $q:X\to X/M$.
\end{definicja}

\begin{uwaga}[Niebezpieczeństwo: topologia końcowa nie musi być liniowa]
Jeżeli $S$ jest przestrzenią liniową, a $f_i$ są liniowe,
topologia końcowa w kategorii \emph{przestrzeni topologicznych}
nie musi czynić $S$ przestrzenią liniowo-topologiczną.
Przykład poniżej pokazuje, że nie można bez dowodu
przenosić własności ilorazu na dowolne rodziny odwzorowań.
\end{uwaga}

\begin{przyklad}[Dwie osie i błędna topologia końcowa]\label{ex:final-cross}
Niech $S=\R^2$, $Y_1=Y_2=\R$ z topologiami zwykłymi,
a $f_1(t)=(t,0)$ i $f_2(t)=(0,t)$. W topologii końcowej
zbiór jednoelementowy $\{(1,1)\}$ jest otwarty,
gdyż oba jego przeciwobrazy są puste. Jednak $\{(0,0)\}$
nie jest otwarty: jego przeciwobraz przez $f_1$ jest
zbiorem $\{0\}\subseteq\R$. Przesunięcie o $(-1,-1)$
przenosi pierwszy zbiór na drugi, więc nie jest
homeomorfizmem. Tymczasem w każdej TVS wszystkie
przesunięcia są homeomorfizmami. A zatem ta topologia
końcowa nie jest topologią przestrzeni liniowo-topologicznej.
\end{przyklad}

\begin{definicja}[Końcowa topologia liniowo-topologiczna]
Niech $X$ będzie przestrzenią liniową nad $\K$, a
$T_i:X_i\to X$, $i\in I$, rodziną odwzorowań
liniowych z przestrzeni liniowo-topologicznych.
\emph{Końcową topologią liniowo-topologiczną}
nazywamy najsilniejszą topologię na $X$, która czyni $X$
przestrzenią liniowo-topologiczną i wszystkie $T_i$
odwzorowaniami ciągłymi. W tej definicji nie zakładamy
warunku Hausdorffa.
\end{definicja}

\begin{twierdzenie}[Istnienie końcowej topologii liniowo-topologicznej]
\label{thm:final-tvs}
Dla każdej przestrzeni liniowej $X$ nad $\K$ i rodziny
liniowych odwzorowań $T_i:X_i\to X$ z TVS istnieje
końcowa topologia liniowo-topologiczna. Ma ona następującą
własność: jeżeli $Z$ jest TVS, a $S:X\to Z$ odwzorowaniem
liniowym, to $S$ jest ciągłe dokładnie wtedy, gdy
$S\circ T_i$ jest ciągłe dla wszystkich $i\in I$.
\end{twierdzenie}
\begin{proof}
Rozpatrzmy rodzinę $\mathfrak T$ wszystkich topologii
liniowo-topologicznych na $X$, względem których wszystkie
$T_i$ są ciągłe. Nie jest ona pusta: zawiera topologię
trywialną. Niech $\tau$ będzie topologią generowaną
przez sumę wszystkich topologii z $\mathfrak T$.
Aby nie pomylić topologii na dziedzinach działań,
ustalmy otwarty w pewnej $\sigma\in\mathfrak T$ zbiór
$O\subseteq X$. Wówczas
$+_X^{-1}(O)$ jest otwarty w $(X,\sigma)\times(X,\sigma)$,
a więc również w $(X,\tau)\times(X,\tau)$.
Analogicznie $m_X^{-1}(O)$ jest otwarty w
$\K\times(X,\tau)$. Ponieważ takie zbiory $O$
tworzą podbazę $\tau$, oba działania są łącznie ciągłe.
Dla każdego $i$ przeciwobrazy zbiorów podbazowych przez
$T_i$ są otwarte w $X_i$; ich skończone przecięcia
i dowolne sumy również. Tak więc $T_i:X_i\to(X,\tau)$
są ciągłe. Każda dopuszczalna topologia jest zawarta
w $\tau$, co dowodzi maksymalności. W szczególności
w tej konstrukcji nie trzeba zakładać, że obrazy $T_i$
generują algebraicznie $X$.

Jeśli $S:X\to Z$ jest liniowe, a wszystkie $S\circ T_i$
są ciągłe, to topologia początkowa $\sigma_S$ na $X$
indukowana przez $S$ jest topologią liniowo-topologiczną
(twierdzenie~\ref{thm:initial-tvs}); wszystkie $T_i$
są ciągłe do $(X,\sigma_S)$, a zatem $\sigma_S\subseteq\tau$.
Z definicji $\sigma_S$ wynika, że $S:(X,\tau)\to Z$
jest ciągłe. W drugą stronę wystarczy składać odwzorowania
ciągłe.
\end{proof}

\begin{uwaga}[Wersja lokalnie wypukła]
Jeżeli wszystkie $X_i$ są lokalnie wypukłe, można
analogicznie rozważać \emph{najsilniejszą lokalnie wypukłą}
topologię na $X$, przy której $T_i$ są ciągłe.
Będzie to szczególnie ważne dla przestrzeni funkcji
testowych $\mathcal D(\Omega)$; tam użyjemy granicy
induktywnej w kategorii przestrzeni lokalnie wypukłych.
Nie należy bez dodatkowych argumentów utożsamiać
tej topologii z końcową topologią zwykłych przestrzeni
topologicznych.
\end{uwaga}

\subsection{Sumy proste: produkt skończony i nieskończony}

\begin{definicja}[Algebraiczna suma prosta]
Dla rodziny przestrzeni liniowych $(X_i)_{i\in I}$
przez $\bigoplus_{i\in I}X_i$ oznaczamy przestrzeń
rodzin $(x_i)\in\prod_iX_i$, w których tylko skończenie
wiele współrzędnych jest niezerowych. Naturalne włożenia
$j_i:X_i\to\bigoplus_iX_i$ umieszczają element w $i$-tej
współrzędnej, a na pozostałych wpisują zero.
\end{definicja}

\begin{twierdzenie}[Skończone sumy proste]
Niech $X_1,\dots,X_n$ będą przestrzeniami liniowo-topologicznymi.
Na $X_1\oplus\dots\oplus X_n$ utożsamionym algebraicznie
z $X_1\times\dots\times X_n$ topologia produktowa jest
końcową topologią liniowo-topologiczną względem włożeń $j_i$.
\end{twierdzenie}
\begin{proof}
Włożenia $j_i$ są ciągłe w topologii produktowej.
Jeśli $Z$ jest TVS i $S:X_1\oplus\dots\oplus X_n\to Z$
jest liniowe, a każde $S\circ j_i$ jest ciągłe,
to dla $(x_1,\dots,x_n)$ mamy
$S(x_1,\dots,x_n)=\sum_{i=1}^n(S\circ j_i)(x_i)$.
Skończona suma ciągłych odwzorowań jest ciągła.
Z uniwersalnej własności końcowej topologii
liniowo-topologicznej (twierdzenie~\ref{thm:final-tvs})
wynika teza.
\end{proof}

\begin{przyklad}[Przestrzeń ciągów o skończonym nośniku]\label{ex:direct-sum}
Niech $\varphi=\bigoplus_{n\geq1}\K e_n$ oznacza
przestrzeń ciągów o skończonym nośniku. Topologia
indukowana z produktu $\K^{\N}$ i \emph{topologia
lokalnie wypukłej sumy prostej} są różne.
Dla jasności: \emph{seminorma} to funkcja $p:X\to[0,\infty)$
spełniająca $p(\lambda x)=|\lambda|p(x)$ i
$p(x+y)\leq p(x)+p(y)$; może przyjmować wartość zero
na niezerowych wektorach. Rodzina seminorm definiuje topologię,
w której bazę otoczeń zera stanowią skończone przecięcia
zbiorów $\{x:p(x)<\varepsilon\}$, $\varepsilon>0$.
Jest to topologia \emph{lokalnie wypukła}: zero posiada bazę
wypukłych otoczeń. Teorię seminorm rozwiniemy w \hyperref[sec:lokalnie-wypukle]{rozdziale 5}.

Na $\varphi$ rozważmy rodzinę seminorm
\[
 p_a(x)=\sum_{n\geq1}a_n|x_n|,
 \qquad a=(a_n)_{n\geq1}\in[0,\infty)^{\N}.
\]
Każda suma jest skończona dla ustalonego $x\in\varphi$.
Wszystkie $p_a$ są seminormami, a ich rodzina
rozdziela punkty, więc określają lokalnie wypukłą
topologię Hausdorffa. Niech $Z$ będzie przestrzenią
liniowo-topologiczną z topologią zadaną przez rodzinę
seminorm. Dla liniowego $S:\varphi\to Z$ oraz
dowolnej seminormy $p$ z tej rodziny zachodzi
\[
 p(Sx)\leq\sum_{n\geq1}|x_n|\,p(S e_n)
         =p_a(x),\qquad a_n=p(S e_n).
\]
Wynika stąd ciągłość każdego liniowego $S:\varphi\to Z$:
dla dowolnego $p$ i $\varepsilon>0$ otwarte otoczenie zera
$\{x\in\varphi:p_a(x)<\varepsilon\}$ zawiera się w
$S^{-1}(\{z\in Z:p(z)<\varepsilon\})$.
Ciągłość $S$ w zerze, a następnie w każdym punkcie,
wynika z opisu topologii $Z$ przez seminormy i liniowości $S$.
W \hyperref[sec:lokalnie-wypukle]{rozdziale 5} udowodnimy, że topologie wszystkich
przestrzeni lokalnie wypukłych można opisać seminormami;
wtedy powyższy argument da pełną własność uniwersalną
sumy prostej \emph{w kategorii przestrzeni lokalnie wypukłych}.
Nie dowodzi on analogicznej własności względem wszystkich
przestrzeni liniowo-topologicznych bez lokalnej wypukłości.

W topologii produktowej odziedziczonej z $\K^{\N}$
ciąg $e_n$ zbiega do zera. W topologii sumy prostej
nie zbiega: $p_{(1,1,\dots)}(e_n)=1$ dla każdego $n$.
To kolejny przykład, w którym ten sam zbiór liniowy
niesie dwie naturalne, ale istotnie różne topologie.
\end{przyklad}

\subsection{Iloraz Hausdorffa dowolnej TVS}

\begin{twierdzenie}[Hausdorffyzacja]\label{thm:hausdorffization}
Niech $X$ będzie przestrzenią liniowo-topologiczną
(niekoniecznie Hausdorffa), a $N=\overline{\{0\}}$.
Wówczas $N$ jest domkniętą podprzestrzenią, $X/N$
jest przestrzenią liniowo-topologiczną Hausdorffa
oraz dla każdego ciągłego liniowego $T:X\to Y$,
gdzie $Y$ jest TVS Hausdorffa, istnieje dokładnie
jedno ciągłe liniowe $\widetilde T:X/N\to Y$
spełniające $T=\widetilde T\circ q$.
\end{twierdzenie}
\begin{proof}
Domkniętość i liniowość $N$ wykazano w \hyperref[sec:tvs]{rozdziale 1}. Wobec
\ref{thm:quotient-tvs} iloraz $X/N$ jest Hausdorffa.
Ponieważ $T$ jest ciągłe i $Y$ jest Hausdorffa,
$\ker T=T^{-1}(\{0\})$ jest domknięte i zawiera zero,
a więc $N\subseteq\ker T$. Z własności uniwersalnej
ilorazu (twierdzenie~\ref{thm:quotient-universal})
istnieje jednoznaczne ciągłe $\widetilde T$.
\end{proof}

\subsection{Podsumowanie: kiedy topologia jest właściwa?}

W konstrukcjach tego rozdziału występują cztery pojęcia:
\begin{enumerate}[label=\textup{(\roman*)}]
\item topologia podprzestrzeni jest \emph{najsłabszą} topologią,
      przy której włożenie do większej przestrzeni jest ciągłe;
\item topologia produktowa jest \emph{najsłabszą} topologią,
      przy której wszystkie rzuty współrzędne są ciągłe;
\item topologia ilorazowa jest \emph{najsilniejszą} topologią,
      przy której rzutowanie ilorazowe jest ciągłe; w tym
      przypadku pozostaje topologią przestrzeni liniowo-topologicznej;
\item ogólna topologia końcowa dla rodziny odwzorowań liniowych
      nie musi być zgodna z działaniami liniowymi; trzeba
      rozróżniać wersję topologiczną, liniowo-topologiczną
      i lokalnie wypukłą.
\end{enumerate}

Warto zapamiętać dwie przeszkody: iloraz przez podprzestrzeń
niedomkniętą nie jest Hausdorffa, a ciągła liniowa bijekcja
nie musi być izomorfizmem topologicznym.
W następnym rozdziale zajmiemy się sieciami, strukturą jednostajną
przestrzeni liniowo-topologicznych, zupełnością oraz uzupełnieniami.

% Wersja źródłowa: audyt III, jednostka 03
\bigskip
\setcounter{section}{2}
\section{Sieci, struktura jednostajna i zupełność przestrzeni liniowo-topologicznych}\label{sec:zupelnosc}

\subsection{Dlaczego ciągi nie wystarczają?}

W przestrzeni metrycznej ciągi wystarczają do charakteryzowania domknięcia,
ciągłości i zupełności. W dowolnej przestrzeni topologicznej jest inaczej:
zbiór może zawierać granice wszystkich swoich zbieżnych ciągów, a mimo to
nie być domknięty. Dlatego zamiast samych ciągów będziemy posługiwać się
\emph{sieciami}. Później pojęcie sieci Cauchy'ego pozwoli określić
zupełność bez wybierania normy czy metryki.

Dla wygody przypomnijmy: $U\subseteq X$ jest otoczeniem
zera, jeżeli zawiera otwarty zbiór zawierający zero;
jest zrównoważony, gdy $\lambda U\subseteq U$
dla $|\lambda|\leq1$.
W tym rozdziale $\K\in\{\R,\C\}$, a TVS jest, jeśli nie zaznaczono
inaczej, przestrzenią liniowo-topologiczną Hausdorffa. Otoczenie punktu
nie musi być otwarte. Gdy piszemy $U\in\mathcal N_X(0)$, mamy na myśli
otoczenie zera w $X$.

\begin{definicja}[Zbiór skierowany i sieć]\label{def:net}
\emph{Zbiorem skierowanym} nazywamy niepusty zbiór $D$ z relacją
refleksyjną i przechodnią $\preceq$, w którym dla dowolnych $d,e\in D$
istnieje $f\in D$ takie, że $d\preceq f$ i $e\preceq f$.
\emph{Sieć} w zbiorze $X$ to odwzorowanie $D\to X$,
$d\mapsto x_d$, gdzie $D$ jest zbiorem skierowanym.
Jeżeli $X$ jest przestrzenią topologiczną, mówimy, że $x_d\to x\in X$,
gdy dla każdego otoczenia $U$ punktu $x$ istnieje $d_0\in D$ takie,
że $x_d\in U$ dla wszystkich $d\succeq d_0$.
\end{definicja}

Każdy ciąg jest siecią, w której $D=\N$ z naturalnym porządkiem.
Innym przykładem jest rodzina $\mathcal F(I)$ wszystkich skończonych
podzbiorów ustalonego zbioru $I$, skierowana przez inkluzję.
Jeśli $X$ jest przestrzenią topologiczną i $x\in X$, to rodzina
otwartych otoczeń $x$, skierowana przez \emph{odwrotną} inkluzję,
również jest zbiorem skierowanym: wspólnym późniejszym indeksem
jest przecięcie dwóch otoczeń.

\begin{twierdzenie}[Domknięcie i ciągłość przez sieci]\label{thm:nets-topology}
Niech $X,Y$ będą przestrzeniami topologicznymi, $A\subseteq X$,
$x\in X$ oraz $f:X\to Y$. Wtedy:
\begin{enumerate}[label=\textup{(\roman*)}]
\item $x\in\overline A$ dokładnie wtedy, gdy istnieje sieć
      $(a_d)_{d\in D}$ elementów $A$ zbieżna do $x$;
\item $f$ jest ciągłe w $x$ dokładnie wtedy, gdy dla każdej sieci
      $x_d\to x$ zachodzi $f(x_d)\to f(x)$.
\end{enumerate}
\end{twierdzenie}
\begin{proof}
Jeśli $a_d\in A$ i $a_d\to x$, każde otoczenie $x$ zawiera element $A$,
zatem $x\in\overline A$. Odwrotnie, niech $x\in\overline A$.
Skierujmy rodzinę otwartych otoczeń $x$ przez odwrotną inkluzję,
a dla każdego takiego otoczenia $U$ wybierzmy $a_U\in A\cap U$.
Dla danego otoczenia $V$ wszystkie indeksy $U\subseteq V$
spełniają $a_U\in V$. Stąd $a_U\to x$.

Jeśli $f$ jest ciągłe w $x$, zachowuje granice sieci zbieżnych do $x$.
Gdyby $f$ nie było ciągłe w $x$, istniałoby otwarte otoczenie $W$
punktu $f(x)$, dla którego w każdym otwartym otoczeniu $U$ punktu $x$
można wybrać $x_U\in U$ z $f(x_U)\notin W$. Wówczas $x_U\to x$,
ale $f(x_U)$ nie zbiega do $f(x)$, co daje sprzeczność.
\end{proof}

\begin{przyklad}[Domknięcie niewykrywalne przez ciągi]\label{ex:sigma-net}
Niech $I$ będzie zbiorem nieprzeliczalnym, a
$X=\K^I$ z topologią produktową. Oznaczmy przez
\[
 \Sigma(I)=\{x\in\K^I:\supp x=\{i\in I:x_i\ne0\}
                 \text{ jest przeliczalny}\}.
\]
Jest to podprzestrzeń liniowa, bo suma dwóch zbiorów przeliczalnych
jest przeliczalna. Jest gęsta w $X$: dla $x\in X$ sieć
$x_F=\sum_{i\in F}x_i e_i$, gdzie $F$ przebiega skończone podzbiory
$I$ uporządkowane przez inkluzję, zbiega współrzędnie do $x$.
Nie jest jednak równa $X$: wektor $\mathbf1=(1)_{i\in I}$
ma nieprzeliczalny nośnik.

Co istotne, $\Sigma(I)$ jest \emph{ciągowo domknięta} w $X$.
Jeśli $x_n\in\Sigma(I)$, to zbiór $J=\bigcup_n\supp x_n$
jest przeliczalny. Jeżeli $x_n\to x$ współrzędnie, to $x_i=0$
dla $i\notin J$, a więc $x\in\Sigma(I)$. Zatem żaden ciąg
z $\Sigma(I)$ nie może zbiegać do $\mathbf1$, choć pewna sieć może.
\end{przyklad}

\subsection{Filtry i ich zbieżność}

Filtry są alternatywnym językiem opisu zbieżności, szczególnie
wygodnym w konstrukcji uzupełnień.

\begin{definicja}[Filtr]\label{def:filter}
\emph{Filtrem} na niepustym zbiorze $X$ nazywamy niepustą rodzinę
$\mathcal F$ podzbiorów $X$ taką, że:
\begin{enumerate}[label=\textup{(\roman*)}]
\item $\varnothing\notin\mathcal F$;
\item $A,B\in\mathcal F$ pociąga $A\cap B\in\mathcal F$;
\item $A\in\mathcal F$ i $A\subseteq B\subseteq X$ pociąga
      $B\in\mathcal F$.
\end{enumerate}
Filtr $\mathcal F$ \emph{zbiega} do $x$ w przestrzeni topologicznej $X$,
gdy $\mathcal N_X(x)\subseteq\mathcal F$.
\end{definicja}

Z każdej sieci $(x_d)_{d\in D}$ powstaje filtr jej \emph{ogonów}:
jest to rodzina wszystkich zbiorów zawierających pewien ogon
$\{x_d:d\succeq d_0\}$. Sieć zbiega do $x$ dokładnie wtedy,
gdy filtr ogonów zbiega do $x$. Jeśli filtr $\mathcal F$ jest zadany,
można w drugą stronę utworzyć sieć: indeksami są pary $(A,a)$,
gdzie $A\in\mathcal F$ i $a\in A$, a
$(A,a)\preceq(B,b)$, gdy $B\subseteq A$.
Wartością sieci w $(A,a)$ jest $a$. Przecięcie dwóch elementów filtra
jest niepuste, więc podana relacja jest skierowana. Ta sieć zbiega
do $x$ dokładnie wtedy, gdy $\mathcal F$ zbiega do $x$.

\subsection{Kanoniczna struktura jednostajna TVS}

W przestrzeni metrycznej stwierdzenie, że $x$ i $y$ są blisko,
wyrażamy nierównością $d(x,y)<\varepsilon$. W TVS rolę tej nierówności
pełni warunek $x-y\in U$, gdzie $U$ jest otoczeniem zera.

\begin{definicja}[Otoczenie przekątnej i struktura jednostajna]
Niech $X$ będzie przestrzenią liniowo-topologiczną. Dla otoczenia
$U$ zera określamy
\[
 E_U=\{(x,y)\in X\times X:x-y\in U\}.
\]
Rodzinę zbiorów $W\subseteq X\times X$ zawierających przynajmniej
jeden $E_U$ nazywamy \emph{kanoniczną strukturą jednostajną} $X$.
Jej elementy to \emph{otoczenia przekątnej} (ang. \emph{entourages}).
\end{definicja}

\begin{twierdzenie}[Własności kanonicznej struktury jednostajnej]
\label{thm:uniform-tvs}
Niech $X$ będzie TVS. Rodzina $\{E_U:U\in\mathcal N_X(0)\}$
spełnia następujące warunki:
\begin{enumerate}[label=\textup{(\roman*)}]
\item każde $E_U$ zawiera przekątną $\Delta_X=\{(x,x):x\in X\}$;
\item dla otoczeń $U,V$ zachodzi $E_{U\cap V}=E_U\cap E_V$;
\item dla każdego $U$ istnieje $V$ takie, że $E_V^{-1}\subseteq E_U$;
\item dla każdego $U$ istnieje $V$ takie, że
      $E_V\circ E_V\subseteq E_U$, gdzie
      $E\circ F=\{(x,z):\exists y\ ((x,y)\in E,(y,z)\in F)\}$.
\end{enumerate}
Topologia wyznaczona przez tę strukturę jednostajną jest oryginalną
topologią $X$: przekrój $E_U[x]=\{y:(x,y)\in E_U\}$
tworzy wraz z $U$ bazę otoczeń $x$.
\end{twierdzenie}
\begin{proof}
Własności (i) i (ii) są bezpośrednie. Do (iii) wybieramy symetryczne
otoczenie $V\subseteq U$; wtedy $(x,y)\in E_V^{-1}$ oznacza
$y-x\in V$, zatem $x-y\in-V=V\subseteq U$.
Dla (iv) wybieramy $V$ z $V+V\subseteq U$. Jeżeli
$x-y\in V$ i $y-z\in V$, to $x-z\in U$.
Wreszcie $E_U[x]=x-U$, a symetryczne otoczenia zera tworzą bazę,
więc otrzymujemy dokładnie topologię $X$.
\end{proof}

\begin{uwaga}
Struktura jednostajna pozwala definiować \emph{jednostajną
ciągłość} również bez metryki: odwzorowanie $f:X\to Y$
między przestrzeniami z ustalonymi strukturami jednostajnymi
jest jednostajnie ciągłe, gdy dla każdego otoczenia przekątnej
$W$ w $Y$ istnieje otoczenie przekątnej $V$ w $X$ takie,
że $(x,y)\in V$ pociąga $(f(x),f(y))\in W$.
Przestrzenie liniowo-topologiczne mają kanoniczne struktury
jednostajne; dowolna przestrzeń topologiczna nie musi ich posiadać.
\end{uwaga}

\begin{twierdzenie}[Ciągłość liniowa implikuje jednostajną]
\label{thm:linear-uniform}
Jeśli $X,Y$ są TVS, a $T:X\to Y$ jest liniowe i ciągłe,
to $T$ jest jednostajnie ciągłe względem kanonicznych
struktur jednostajnych.
\end{twierdzenie}
\begin{proof}
Niech $W$ będzie otoczeniem zera w $Y$. Z ciągłości $T$ w zerze
wybieramy otoczenie $U$ zera w $X$ takie, że $T(U)\subseteq W$.
Jeśli $x-y\in U$, to $T(x)-T(y)=T(x-y)\in W$, co jest
żądaną jednostajną ciągłością.
\end{proof}

\subsection{Sieci Cauchy'ego i zupełność}

\begin{definicja}[Sieć Cauchy'ego, filtr Cauchy'ego i zupełność]
\label{def:tvs-complete}
Niech $X$ będzie TVS. Sieć $(x_d)_{d\in D}$ w $X$ nazywamy
\emph{siecią Cauchy'ego}, jeżeli dla każdego $U\in\mathcal N_X(0)$
istnieje $d_0\in D$ takie, że
$ x_d-x_e\in U $ dla wszystkich $d,e\succeq d_0$.
Filtr $\mathcal F$ na $X$ nazywamy \emph{filtrem Cauchy'ego},
jeżeli dla każdego $U\in\mathcal N_X(0)$ istnieje
$A\in\mathcal F$ takie, że $A-A\subseteq U$.
Przestrzeń $X$ jest \emph{zupełna}, jeżeli każda sieć Cauchy'ego
w $X$ ma granicę w $X$. Jest \emph{ciągowo zupełna},
jeżeli każdy jej ciąg Cauchy'ego ma granicę w $X$.
\end{definicja}

\begin{twierdzenie}[Równoważność opisu sieciowego i filtrowego]
\label{thm:filter-cauchy}
Niech $X$ będzie TVS. Sieć jest siecią Cauchy'ego dokładnie wtedy,
gdy jej filtr ogonów jest filtrem Cauchy'ego. Ponadto $X$ jest
zupełna dokładnie wtedy, gdy każdy filtr Cauchy'ego na $X$
jest zbieżny.
\end{twierdzenie}
\begin{proof}
Jeśli sieć jest Cauchy'ego, to dla otoczenia $U$ wybieramy
ogon $A$, w którym różnice dowolnych dwóch wyrazów należą do $U$;
wtedy $A-A\subseteq U$ i $A$ należy do filtra ogonów.
W drugą stronę, jeśli filtr ogonów jest Cauchy'ego, wybieramy
$A$ z $A-A\subseteq U$ oraz ogon zawarty w $A$. Różnice
wyrazów tego ogona należą do $U$.

Załóżmy, że $X$ jest zupełna, i niech $\mathcal F$ będzie filtrem
Cauchy'ego. Rozważmy związaną z nim sieć indeksowaną parami
$(A,a)$, opisaną po definicji~\ref{def:filter}. Dla otoczenia
$U$ wybieramy $A\in\mathcal F$ z $A-A\subseteq U$.
Wszystkie wyrazy sieci o indeksach późniejszych niż $(A,a)$
należą do $A$, więc sieć jest Cauchy'ego. Niech $x$ będzie jej
granicą. Dla dowolnego otoczenia $V$ punktu $x$ istnieje
$(A,a)$ takie, że wszystkie późniejsze wyrazy należą do $V$.
W szczególności wszystkie punkty $b\in A$ (otrzymane
z indeksów $(A,b)$) należą do $V$. Stąd $A\subseteq V$,
a więc $V\in\mathcal F$, co oznacza zbieżność filtra.
Odwrotna implikacja wynika z zastosowania założenia o filtrach
do filtra ogonów dowolnej sieci Cauchy'ego.
\end{proof}

\begin{stwierdzenie}[Granice i jednoznaczność]
\label{prop:complete-basic}
Niech $X$ będzie TVS. Każda sieć zbieżna w $X$ jest Cauchy'ego.
Każda sieć Cauchy'ego ma co najwyżej jedną granicę.
Jeżeli sieć Cauchy'ego ma podsieć zbieżną do $x\in X$, to sama
zbiega do $x$.
\end{stwierdzenie}
\begin{proof}
Jeśli $x_d\to x$ i $U$ jest otoczeniem zera, wybieramy
symetryczne otoczenie $V$ z $V+V\subseteq U$.
Dla późnych indeksów $x_d-x\in V$, a zatem
$x_d-x_e\in V-V\subseteq U$. Jednoznaczność granicy
wynika z własności Hausdorffa.

Załóżmy, że sieć jest Cauchy'ego, a jej podsieć zbiega do $x$.
Dla otoczenia $U$ wybieramy symetryczne $V$ z $V+V\subseteq U$.
Istnieje $d_0$ takie, że $x_d-x_e\in V$ dla $d,e\succeq d_0$.
Zbieżna podsieć ma wyraz $x_{e}$ o indeksie $e\succeq d_0$
i z $x_e-x\in V$. Wówczas dla każdego $d\succeq d_0$
mamy $x_d-x\in V+V\subseteq U$.
\end{proof}

\begin{uwaga}[Podsieć]
Podsieć sieci $(x_d)_{d\in D}$ rozumiemy jako sieć
$(x_{\phi(e)})_{e\in E}$, gdzie $E$ jest skierowany,
a $\phi:E\to D$ ma własność: dla każdego $d_0\in D$
istnieje $e_0\in E$ takie, że $\phi(e)\succeq d_0$
dla wszystkich $e\succeq e_0$. Ta własność, zwana
\emph{współkońcowością}, zapewnia przenoszenie warunków
prawdziwych od pewnego indeksu na podsieć.
\end{uwaga}

\begin{twierdzenie}[Zupełność ciągowa w przestrzeniach metryzowalnych]
\label{thm:metric-complete}
Niech $X$ będzie metryzowalną przestrzenią liniowo-topologiczną
Hausdorffa. Wtedy $X$ jest zupełna dokładnie wtedy,
gdy jest ciągowo zupełna.
\end{twierdzenie}
\begin{proof}
Jedna implikacja jest oczywista. W drugą stronę załóżmy,
że wszystkie ciągi Cauchy'ego są zbieżne. Niech $(x_d)$
będzie dowolną siecią Cauchy'ego. Wybierzmy malejącą
przeliczalną bazę symetrycznych otoczeń zera $(U_n)$
taką, że $U_{n+1}+U_{n+1}\subseteq U_n$.
Dla każdego $n$ istnieje $d_n$ takie, że
$x_d-x_e\in U_n$ dla $d,e\succeq d_n$.
Rekurencyjnie wybierzmy indeksy $e_n$ spełniające
$e_n\succeq d_1,\ldots,d_n,e_{n-1}$.
Wówczas $(x_{e_n})$ jest ciągiem Cauchy'ego: dla
$m,n\geq N$ oba indeksy są późniejsze od $d_N$,
więc $x_{e_n}-x_{e_m}\in U_N$.
Niech $x_{e_n}\to x$. Dla zadanego otoczenia $W$ zera
wybierzmy $N$ tak duże, że $U_N+U_N\subseteq W$.
Dla każdego $d\succeq d_N$ i $n\geq N$ mamy
$x_d-x_{e_n}\in U_N$. Wybierając $n$ dostatecznie duże,
aby $x_{e_n}-x\in U_N$, otrzymujemy $x_d-x\in W$.
Zatem oryginalna sieć zbiega do $x$.
\end{proof}

\subsection{Operacje na przestrzeniach zupełnych}

\begin{twierdzenie}[Podprzestrzenie i produkty zupełne]
\label{thm:closed-complete}
Niech $X$ będzie TVS Hausdorffa.
\begin{enumerate}[label=\textup{(\roman*)}]
\item Jeśli $X$ jest zupełna, to każda domknięta
      podprzestrzeń liniowa $M\subseteq X$ jest zupełna.
\item Jeśli $M\subseteq X$ jest podprzestrzenią liniową zupełną
      w topologii indukowanej, to $M$ jest domknięta w $X$.
\item Dla rodziny TVS Hausdorffa $(X_i)_{i\in I}$ produkt
      $\prod_{i\in I}X_i$ jest zupełny wtedy i tylko wtedy,
      gdy wszystkie $X_i$ są zupełne.
\end{enumerate}
\end{twierdzenie}
\begin{proof}
Sieć Cauchy'ego w $M$ jest Cauchy'ego w $X$.
Jeżeli $X$ jest zupełna, ma ona granicę w $X$; domkniętość
$M$ zapewnia, że granica należy do $M$. To daje (i).
Jeśli $x\in\overline M$, twierdzenie~\ref{thm:nets-topology}
daje sieć w $M$ zbieżną do $x$ w $X$. Jest ona Cauchy'ego
w $M$, więc z założenia ma granicę $m\in M$.
Jednoznaczność granicy w $X$ daje $x=m$, co dowodzi (ii).

Niech teraz wszystkie $X_i$ będą zupełne, a $(x_d)$
siecią Cauchy'ego w produkcie. Dla każdego $i$ sieć
$(\pi_i x_d)$ jest Cauchy'ego w $X_i$, gdyż rzuty
są ciągłe i liniowe, zatem jednostajnie ciągłe.
Istnieją więc $x_i\in X_i$ takie, że $\pi_i x_d\to x_i$.
Z uniwersalnej własności produktu $x_d\to(x_i)_{i\in I}$.
W drugą stronę, dla każdego $j\in I$ kopia $X_j$
złożona z wektorów o zerowych współrzędnych poza $j$
jest domkniętą podprzestrzenią produktu (pozostałe czynniki
są Hausdorffa). Jeśli produkt jest zupełny, część (i)
pokazuje zupełność $X_j$.
\end{proof}

\begin{przyklad}[Przestrzeń ciągowo zupełna, ale niezupełna]
\label{ex:seq-not-complete}
Niech $I$ będzie nieprzeliczalny, $X=\K^I$ z topologią
produktową, a $M=\Sigma(I)$ z przykładu~\ref{ex:sigma-net}.
Produkt $X$ jest zupełny, ponieważ $\K$ jest zupełne.
Każdy ciąg Cauchy'ego $(x_n)$ w $M$ jest Cauchy'ego w $X$,
więc zbiega tam do pewnego $x\in X$. Nośnik $x$ zawiera
się w przeliczalnej sumie $\bigcup_n\supp x_n$, toteż $x\in M$.
Zatem $M$ jest ciągowo zupełna. Nie jest zupełna,
ponieważ jest gęstą właściwą podprzestrzenią zupełnej
przestrzeni $X$ i nie jest w niej domknięta; do
niezupełności wystarcza też część (ii) twierdzenia
\ref{thm:closed-complete}.

Konkretnym świadkiem jest sieć $x_F=\sum_{i\in F}e_i$,
$F\in\mathcal F(I)$: jest Cauchy'ego w $M$, gdyż
zbiega do $\mathbf1$ w $X$, lecz nie zbiega w $M$.
\end{przyklad}

\subsection{Uzupełnienie TVS: konstrukcja bez normy}

Chcemy skonstruować przestrzeń zupełną $\widehat X$
i ciągłe liniowe zanurzenie $j:X\to\widehat X$ o gęstym obrazie,
bez zakładania lokalnej wypukłości czy metryzowalności.
Poniższy dowód jest konstrukcją opartą na sieciach Cauchy'ego.

\begin{twierdzenie}[Istnienie i własność uniwersalna uzupełnienia]
\label{thm:tvs-completion}
Niech $X$ będzie przestrzenią liniowo-topologiczną Hausdorffa
nad $\K$. Istnieją zupełna przestrzeń liniowo-topologiczna
Hausdorffa $\widehat X$ oraz liniowe zanurzenie topologiczne
$j:X\to\widehat X$ o gęstym obrazie. Dla każdej zupełnej
przestrzeni liniowo-topologicznej Hausdorffa $Y$ oraz każdego
ciągłego liniowego $T:X\to Y$ istnieje dokładnie jedno
ciągłe liniowe $\widehat T:\widehat X\to Y$ takie,
że $\widehat T\circ j=T$. Para $(\widehat X,j)$ jest
jednoznaczna z dokładnością do jedynego izomorfizmu
liniowo-topologicznego zachowującego osadzenie $X$.
\end{twierdzenie}
\begin{proof}
\emph{Krok 1: klasy sieci Cauchy'ego.}
Aby otrzymać rzeczywisty zbiór klas, a nie klasę właściwą,
ustalmy rodzinę $\mathcal U$ wszystkich otwartych zrównoważonych
otoczeń zera w $X$ oraz zbiór skierowany
$D_*=\mathcal F(\mathcal U)$ jej skończonych podzbiorów,
uporządkowany przez inkluzję. Każda sieć Cauchy'ego $(x_d)_{d\in D}$
ma równoważną sieć indeksowaną przez $D_*$: dla każdego
$U\in\mathcal U$ wybierzmy $d_U$ takie, że
$x_d-x_e\in U$ dla $d,e\succeq d_U$; dla skończonego
$F\subseteq\mathcal U$ wybierzmy $d_F$ późniejszy od wszystkich
$d_U$ z $U\in F$. Pokażemy, że sieć $(x_{d_F})_{F\in D_*}$
jest Cauchy'ego i równoważna pierwotnej. Dla otoczenia
$W$ zera wybierzmy $U\in\mathcal U$ z $U\subseteq W$.
Gdy $F,G\supseteq\{U\}$, mamy $d_F,d_G\succeq d_U$,
a więc $x_{d_F}-x_{d_G}\in U\subseteq W$.
Podobnie dla $F\supseteq\{U\}$ i $d\succeq d_U$
zachodzi $x_{d_F}-x_d\in W$; jest to dokładnie
wymagana równoważność dwóch sieci. Nie wymagamy,
aby wybór $F\mapsto d_F$ był monotoniczny: podane
oszacowanie zapewnia wszystkie potrzebne własności.
W konstrukcji $\widehat X$ wystarczy zatem uwzględniać
sieci o ustalonym zbiorze indeksów $D_*$. Tworzą one
zbiór, co usuwa problem klas właściwych.

Dwie sieci Cauchy'ego $(x_d)_{d\in D}$ i $(y_e)_{e\in E}$
nazwiemy równoważnymi, gdy dla każdego otoczenia $U$ zera
istnieją $d_0,e_0$ takie, że
$x_d-y_e\in U$ dla $d\succeq d_0$ i $e\succeq e_0$.
Relacja ta jest zwrotna, ponieważ sieci są Cauchy'ego,
symetryczna dzięki symetrycznym otoczeniom i przechodnia:
dla danego $U$ wybieramy symetryczne $V$ z $V+V\subseteq U$
i łączymy dwa oszacowania przez wspólny późny wyraz
sieci pośredniej. Relację ograniczamy formalnie do zbioru sieci indeksowanych
przez $D_*$, a przez $[x_d]$ oznaczamy klasę
równoważności; $\widehat X$ jest zbiorem tych klas.
Dzięki powyższemu przeindeksowaniu zapis $[x_d]$
pozostaje jednoznaczny również dla sieci o innym zbiorze
indeksów. Algebraicznie sumujemy klasy, przechodząc do sieci
$(x_d+y_e)_{(d,e)\in D\times E}$, a następnie w razie potrzeby
do równoważnego przedstawiciela na $D_*$;
mnożenie przez ustalony skalar określamy przez
$\lambda[x_d]=[\lambda x_d]$. Liniowość działań w $X$
i ich ciągłość w zerze zapewniają, że są to działania
dobrze określone na klasach sieci Cauchy'ego.
Definiujemy $j(x)$ jako klasę stałej sieci o wartości $x$.

\emph{Krok 2: topologia.}
Wybierzmy otwarte zrównoważone otoczenie $U$ zera w $X$.
Niech $\widehat U$ oznacza zbiór klas $[x_d]$, dla których
istnieje otwarte otoczenie $V$ zera oraz indeks $d_0$ takie, że
\[
       x_d+V\subseteq U\qquad(d\succeq d_0).
\]
Definicja nie zależy od przedstawiciela: jeśli $(y_e)$
jest równoważna $(x_d)$, wybieramy otoczenie $W$ z
$W+W\subseteq V$. Dla dostatecznie późnych $d,e$
mamy $y_e-x_d\in W$, a więc
$y_e+W\subseteq x_d+(W+W)\subseteq x_d+V\subseteq U$.

Rodzina $\{\widehat U\}$ jest bazą otoczeń zera w
$\widehat X$. Istotnie, dla $V+V\subseteq U$ zachodzi
$\widehat V+\widehat V\subseteq\widehat U$,
a dla $U_1,U_2$ można wybrać otwarte zrównoważone
$V\subseteq U_1\cap U_2$. Wyjaśnijmy dokładnie, dlaczego przesunięcia tych zbiorów
określają topologię grupy addytywnej. Jeśli
$z=[x_d]\in\widehat U$ i $x_d+W\subseteq U$ dla późnych $d$,
wybierzmy otwarte zrównoważone $V\subseteq W$.
Dla $w=[y_e]\in\widehat V$ mamy
$y_e+W_0\subseteq V$ przy dostatecznie późnych $e$
i pewnym otwartym otoczeniu $W_0$ zera (niezależnym od $e$).
Wówczas dla późnych $d,e$ zachodzi
$x_d+y_e+W_0\subseteq x_d+V\subseteq U$,
a zatem $z+\widehat V\subseteq\widehat U$.
Tak więc każda klasa należąca do $\widehat U$ ma
otoczenie będące przesunięciem pewnego $\widehat V$
i mieszczące się w $\widehat U$; w połączeniu z
$\widehat V+\widehat V\subseteq\widehat U$ oraz
stabilnością na skończone przecięcia uzasadnia to
bazę otoczeń i ciągłość dodawania.
Zbiory $\widehat U$ są zrównoważone, gdyż dla
$0<|\lambda|\leq1$ mnożenie przez $\lambda$
przenosi $x_d+V\subseteq U$ w
$\lambda x_d+\lambda V\subseteq\lambda U\subseteq U$,
a przypadek $\lambda=0$ jest oczywisty.
Dlatego dodawanie i zmiana znaku są ciągłe. Ciągłość
mnożenia przez skalar uzasadnimy po wykazaniu gęstości $j(X)$.

Odwzorowanie $j$ jest liniowe, a z definicji
$j^{-1}(\widehat U)=U$: jeśli $x\in U$, to
$x+V\subseteq U$ dla pewnego otwartego otoczenia $V$ zera,
a implikacja odwrotna jest natychmiastowa.
Stąd $j$ jest zanurzeniem topologicznym. Co więcej,
$\widehat X$ jest Hausdorffa. Jeśli
$z=[x_d]\in\bigcap_U\widehat U$, to dla każdego otoczenia
$U$ można wybrać otwarte zrównoważone $V\subseteq U$
i $z\in\widehat V$, skąd ostatecznie $x_d\in V\subseteq U$.
Zatem $x_d\to0$ i $z=0$.

\emph{Krok 3: gęstość oraz granice sieci z $X$.}
Jeżeli $z=[x_d]$ i $U$ jest otoczeniem zera, wybieramy
otwarte zrównoważone $V,W$ z $V+W\subseteq U$.
Z warunku Cauchy'ego dla $(x_d)$ istnieje $d_0$ takie, że
$x_d-x_e\in V$ dla $d,e\succeq d_0$.
Dla każdego $e\succeq d_0$ sieć $(x_d-x_e)_d$ jest
ostatecznie zawarta w $V$, a wraz z $W$ mieści się w $U$.
Stąd $z-j(x_e)\in\widehat U$. Dowodzi to zarówno
gęstości $j(X)$, jak i zbieżności $j(x_d)\to z$.

Teraz domykamy dowód ciągłości mnożenia przez skalar.
Dla ustalonego $z_0\in\widehat X$ i zrównoważonego
$\widehat U$ wybierzmy zrównoważone $\widehat V$
takie, że $\widehat V+\widehat V\subseteq\widehat U$.
Z gęstości istnieje $x\in X$ takie, że
$z_0-j(x)\in\widehat V$. Dla $|\mu|\leq1$ mamy
$\mu(z_0-j(x))\in\widehat V$, natomiast
$\mu j(x)=j(\mu x)\to0$ przy $\mu\to0$.
Zatem $\mu z_0\to0$. W dowolnym punkcie
$(\lambda_0,z_0)$ korzystamy z tożsamości
$\lambda z-\lambda_0z_0=\lambda(z-z_0)+
(\lambda-\lambda_0)z_0$. Dla $|\lambda|\leq M$
i otwartego zrównoważonego $\widehat W$ z
$M\widehat W\subseteq\widehat V$ pierwszy składnik
należy do $\widehat V$, gdy $z-z_0\in\widehat W$,
a drugi przy $\lambda\to\lambda_0$ również należy
do $\widehat V$. To daje ciągłość łączną.

\emph{Krok 4: zupełność $\widehat X$.}
Niech $(z_\alpha)_{\alpha\in A}$ będzie siecią Cauchy'ego
w $\widehat X$. Niech $\mathcal U$ będzie rodziną otwartych,
zrównoważonych otoczeń zera w $X$, skierowaną przez odwrotną
inkluzję. Z gęstości dla każdej pary $(\alpha,U)$ wybieramy
$x_{\alpha,U}\in X$ tak, aby
$z_\alpha-j(x_{\alpha,U})\in\widehat U$.
Indeksami nowej sieci są pary $(\alpha,U)$, uporządkowane
wzrostem $\alpha$ i malejącym $U$.
Ta sieć w $X$ jest Cauchy'ego. Dla danego otwartego
zrównoważonego $V$ wybieramy otwarte zrównoważone $W$
z $W+W+W\subseteq V$. Dla dostatecznie późnych
$\alpha,\beta$ mamy $z_\alpha-z_\beta\in\widehat W$,
a dla $U_1,U_2\subseteq W$ zachodzi
\[
 j(x_{\alpha,U_1}-x_{\beta,U_2})
 \in\widehat W+\widehat W+\widehat W
 \subseteq\widehat V.
\]
Ponieważ $j^{-1}(\widehat V)=V$, różnica wyrazów należy do $V$.
Oznaczmy klasę tej sieci przez $z\in\widehat X$.
Na mocy kroku 3 zachodzi $j(x_{\alpha,U})\to z$.
Pokażemy, że $z_\alpha\to z$. Ustalmy otoczenie
$V$ zera w $X$ i wybierzmy $W$ z $W+W+W\subseteq V$.
Dla $\alpha,\beta$ dostatecznie późnych
$z_\alpha-z_\beta\in\widehat W$.
Wybieramy parę $(\beta,U)$ z późnym $\beta$ i $U\subseteq W$
na tyle późną, aby ponadto
$j(x_{\beta,U})-z\in\widehat W$. Ponieważ
$z_\beta-j(x_{\beta,U})\in\widehat U\subseteq\widehat W$,
dla wszystkich odpowiednio późnych $\alpha$ zachodzi
\[
 z_\alpha-z=(z_\alpha-z_\beta)
 +(z_\beta-j(x_{\beta,U}))+(j(x_{\beta,U})-z)
 \in\widehat V.
\]
Przestrzeń $\widehat X$ jest więc zupełna.

\emph{Krok 5: przedłużanie i jednoznaczność.}
Niech $T:X\to Y$ będzie jak w tezie. Jednostajna ciągłość
$T$ zapewnia, że dla każdej sieci Cauchy'ego $(x_d)$
sieć $(Tx_d)$ jest Cauchy'ego w $Y$, a więc zbiega.
Definiujemy
$\widehat T([x_d])=\lim_d Tx_d$.
Równoważne sieci mają tę samą granicę, gdyż ich różnice
na produkcie indeksów zbiegają do zera; liniowość jest
natychmiastowa. Aby wykazać ciągłość, niech $W$ będzie
otwartym otoczeniem zera w $Y$. Wybierzmy otwarte
symetryczne $V$ z $V+V\subseteq W$. Wówczas
$\overline V\subseteq V+V\subseteq W$: dla $y\in\overline V$
zbiór $y+V$ spotyka $V$, więc $y\in V-V$.
Z ciągłości $T$ istnieje otwarte zrównoważone otoczenie
$U$ zera w $X$ takie, że $T(U)\subseteq V$.
Dla $[x_d]\in\widehat U$ wyrazy $x_d$ o późnych indeksach
należą do $U$, a więc
$\widehat T([x_d])=\lim_d Tx_d\in\overline V\subseteq W$.
Otrzymujemy ciągłość $\widehat T$ w zerze, a z liniowości
wszędzie. Ponieważ $j(X)$ jest gęsta, przedłużenie jest
jedyne: dwa ciągłe odwzorowania do przestrzeni Hausdorffa
zgodne na zbiorze gęstym są równe. Jeśli
$(\widehat X',j')$ jest drugą parą o podanych własnościach,
przedłużamy $j'$ do $\widehat X\to\widehat X'$ oraz $j$
do $\widehat X'\to\widehat X$. Z jednoznaczności ich
złożenia są identycznościami, a więc otrzymujemy żądany
izomorfizm.
\end{proof}

\begin{center}
\begin{tikzpicture}[>=Latex, every node/.style={font=\small}]
  \node (X) at (0,1.3) {$X$};
  \node (H) at (5.0,1.3) {$\widehat X$};
  \node (Y) at (5.0,-0.45) {$Y$};
  \draw[->,thick] (X)--node[above] {$j$ (gęsty obraz)} (H);
  \draw[->,thick] (X)--node[below left] {$T$} (Y);
  \draw[->,thick] (H)--node[right] {$\widehat T$} (Y);
\end{tikzpicture}

{\small Rysunek 3.1. Każdy ciągły operator liniowy $T$ z $X$
do zupełnej TVS Hausdorffa $Y$ przedłuża się jednoznacznie
na uzupełnienie $\widehat X$.}
\end{center}

\begin{uwaga}[Dlaczego w konstrukcji pojawia się dodatkowe $V$?]
Nie wystarczyłoby zdefiniować $\widehat U$ przez wymaganie,
aby reprezentująca sieć \emph{ostatecznie należała} do $U$.
W $\R$ sieci $1-1/n$ i $1+1/n$ reprezentują tę samą klasę,
ale tylko pierwsza jest ostatecznie zawarta w $(-1,1)$.
Warunek $x_d+V\subseteq U$ wymaga odpowiedniego zapasu
do brzegu $U$ i nie zależy od reprezentanta.
\end{uwaga}

\subsection{Jak rozpoznawać uzupełnienie w praktyce?}

\begin{stwierdzenie}[Kryterium uzupełnienia]\label{prop:completion-criterion}
Niech $X$ będzie TVS Hausdorffa, $Y$ zupełną TVS Hausdorffa,
a $T:X\to Y$ liniowym zanurzeniem topologicznym o gęstym obrazie.
Wtedy $Y$, razem z $T$, jest uzupełnieniem $X$ w sensie
twierdzenia~\ref{thm:tvs-completion}.
\end{stwierdzenie}
\begin{proof}
Niech $(\widehat X,j)$ będzie uzupełnieniem z poprzedniego
twierdzenia. Z jego własności uniwersalnej otrzymujemy ciągłe
liniowe $S:\widehat X\to Y$ takie, że $S\circ j=T$.
Zbudujemy odwrotność. Odwzorowanie
$R_0:T(X)\to\widehat X$, $R_0(Tx)=j(x)$, jest liniowe,
ciągłe i jednostajnie ciągłe, ponieważ $T$ jest zanurzeniem
topologicznym. Dla $y\in Y$ wybierzmy sieć
$(y_d)$ punktów $T(X)$ zbieżną do $y$ (istnieje dzięki
gęstości). Jest to sieć Cauchy'ego, zatem $(R_0y_d)$
jest Cauchy'ego w zupełnej $\widehat X$. Definiujemy
$R(y)=\lim_d R_0(y_d)$.
Jeżeli dwie sieci z $T(X)$ zbiegają do tego samego $y$,
to ich różnice na produkcie indeksów zbiegają do zera.
Z jednostajnej ciągłości $R_0$ wynika równość granic;
zatem definicja nie zależy od wyboru sieci. Z tego samego
powodu $R$ jest liniowe.

Sprawdzimy ciągłość $R$ w zerze. Niech $U$ będzie otwartym
otoczeniem zera w $\widehat X$. Wybierzmy otwarte
symetryczne otoczenie $V$ z $V+V\subseteq U$;
wówczas $\overline V\subseteq U$, co wynika z argumentu
użytego w kroku 5 poprzedniego dowodu. Z ciągłości $R_0$
istnieje otwarte otoczenie $W$ zera w $Y$ takie, że
$R_0(W\cap T(X))\subseteq V$.
Dla $y\in W$ możemy dobrać sieć punktów z $W\cap T(X)$
zbieżną do $y$ (otwartość $W$ i gęstość $T(X)$).
Wtedy $R(y)\in\overline V\subseteq U$, zatem $R$
jest ciągłe. Na gęstym zbiorze $T(X)$ zachodzi
$S\circ R=\mathrm{id}_Y$, a na gęstym $j(X)$ zachodzi
$R\circ S=\mathrm{id}_{\widehat X}$. Z ciągłości obie
równości rozciągają się na całe dziedziny. Wobec tego
$S$ jest izomorfizmem liniowo-topologicznym.
\end{proof}

\begin{przyklad}[Trzy klasyczne uzupełnienia]
\begin{enumerate}[label=\textup{(\roman*)}]
\item Przestrzeń $c_{00}$ z normą $\ell^p$, $1\leq p<\infty$,
      ma uzupełnienie $\ell^p$: obcięcia ciągów do pierwszych
      $n$ współrzędnych są zbieżne w normie $\ell^p$.
\item Przestrzeń $C([0,1],\K)$ z normą
      $\|f\|_1=\int_{[0,1]}|f(t)|\,\mathrm d\lambda(t)$
      ma uzupełnienie $L^1([0,1],\lambda)$: funkcje ciągłe
      są gęste w $L^1$, a $L^1$ jest zupełna.
\item Przestrzeń $c_{00}$ z topologią odziedziczoną
      z produktu $\K^{\N}$ ma uzupełnienie $\K^{\N}$.
      W tym przypadku nie używa się jednej normy.
\end{enumerate}
\end{przyklad}

\subsection{Co wykorzystamy w dalszej części skryptu?}

Zupełność nie jest własnością wyłącznie przestrzeni Banacha.
Przykładowo $\K^{\N}$ z topologią produktową jest zupełna,
ale nie ma ograniczonego otoczenia zera i nie jest normowalna.
Następnie przejdziemy do własności rodzin operatorów, warunków
metryzowalności i lokalnej wypukłości. To przygotuje narzędzia
do badania przestrzeni $C^\infty(\Omega)$ oraz $\mathcal D(\Omega)$,
a docelowo ich dualnych przestrzeni dystrybucji.

% Wersja źródłowa: audyt II, jednostka 04
\setcounter{section}{3}
\section{Zbiory ograniczone, rodziny operatorów i metryzowalność TVS}\label{sec:rodziny-metryzowalnosc}

\subsection{Motywacja i przyjęte założenia}

W przestrzeni unormowanej liczba $\sup_{x\in B}\|x\|$ mierzy
rozmiar zbioru $B$. W ogólnej przestrzeni liniowo-topologicznej
nie mamy wyróżnionej normy. Możemy jednak określić, czy wszystkie
punkty zbioru $B$ mieszczą się w odpowiednim wielokrotnym
powiększeniu dowolnego otoczenia zera. W \hyperref[sec:tvs]{rozdziale 1} przyjęliśmy
następującą definicję: $B\subseteq X$ jest \emph{ograniczony}, gdy
dla każdego otoczenia $U$ zera istnieje $t>0$ takie, że
$B\subseteq tU$. Jest to ograniczoność w sensie von Neumanna.

W całym rozdziale $\K\in\{\R,\C\}$, a $X,Y$ oznaczają
przestrzenie liniowo-topologiczne nad $\K$; jeśli nie wskazano
inaczej, są one Hausdorffa. Otoczenia zera nie muszą być otwarte.
Przez $\mathcal L(X,Y)$ rozumiemy \emph{zbiór ciągłych odwzorowań
liniowych} $X\to Y$. Dla rodziny odwzorowań $\mathcal T$
zapis $\mathcal T(B)$ oznacza zbiór
$\{T(x):T\in\mathcal T,\ x\in B\}$.

\subsection{Sekwencyjne kryterium ograniczoności}

\begin{twierdzenie}[Kryterium ograniczoności przez skalarne ciągi]
\label{thm:bounded-sequential-test}
Niech $X$ będzie przestrzenią liniowo-topologiczną, a
$B\subseteq X$. Następujące warunki są równoważne:
\begin{enumerate}[label=\textup{(\roman*)}]
\item $B$ jest ograniczony w sensie von Neumanna;
\item dla każdego ciągu $(x_n)$ w $B$ i każdego ciągu
      $(\lambda_n)$ w $\K$ zbieżnego do zera zachodzi
      $\lambda_nx_n\to0$ w $X$.
\end{enumerate}
\end{twierdzenie}
\begin{proof}
Załóżmy (i). Niech $U$ będzie otoczeniem zera; wybierzmy
zrównoważone otoczenie $V\subseteq U$.
Istnieje $t>0$ takie, że $B\subseteq tV$.
Gdy $|\lambda_n|t\le1$, dostajemy
$\lambda_nx_n\in\lambda_ntV\subseteq V\subseteq U$.
Zatem $\lambda_nx_n\to0$.

Jeżeli (i) nie zachodzi, istnieje zrównoważone otoczenie
$U$ zera takie, że $B\not\subseteq tU$ dla żadnego $t>0$.
Dla każdego $n\ge1$ wybierzmy $x_n\in B\setminus nU$
i połóżmy $\lambda_n=1/n$. Wtedy $\lambda_n\to0$,
ale $\lambda_nx_n\notin U$ dla żadnego $n$, co przeczy (ii).
\end{proof}

\begin{stwierdzenie}[Podstawowe działania na zbiorach ograniczonych]
\label{prop:bounded-operations}
Niech $X$ będzie TVS, a $B,C\subseteq X$ zbiorami ograniczonymi.
Wówczas $B\cup C$, $B+C$, $\lambda B$ dla $\lambda\in\K$
oraz domknięcie $\overline B$ są ograniczone. Każdy podzbiór
zbioru ograniczonego jest ograniczony.
\end{stwierdzenie}
\begin{proof}
Wybierzmy zrównoważone otoczenie $V$ zera z
$V+V\subseteq U$, gdzie $U$ jest zadanym otoczeniem zera.
Istnieją liczby $s,t>0$ z $B\subseteq sV$ i $C\subseteq tV$.
Jeżeli $r=\max(s,t)$, to $B\cup C\subseteq rV$, ponieważ
$V$ jest zrównoważony. Ponadto
$B+C\subseteq sV+tV\subseteq (s+t)(V+V)\subseteq(s+t)U$;
wystarcza też użyć $r$ i sumy $V+V$. Ograniczoność
$\lambda B$ i podzbiorów wynika bezpośrednio z definicji.

Dla domknięcia wybierzmy otwarte zrównoważone otoczenie
$V$ z $V+V\subseteq U$. Zachodzi $\overline V\subseteq V+V$:
gdy $x\in\overline V$, otwarty zbiór $x+V$ przecina $V$,
zatem $x\in V-V=V+V$. Skoro $B\subseteq tV$ dla pewnego
$t>0$, mamy
$\overline B\subseteq t\overline V\subseteq tU$.
\end{proof}

\begin{uwaga}[Ważna pułapka: otoczka wypukła]
Dla $B\subseteq X$ przez $\operatorname{conv}(B)$ rozumiemy
zbiór wszystkich skończonych kombinacji wypukłych punktów $B$
(przyjmując $\operatorname{conv}(\varnothing)=\varnothing$).
W \emph{dowolnej} TVS otoczka wypukła zbioru ograniczonego
nie musi być ograniczona. Własność ta stanie się prawdziwa
po założeniu lokalnej wypukłości: dla dowolnego otoczenia
$U$ zera wybieramy wypukłe, zrównoważone otoczenie
$V\subseteq U$. Jeżeli $B\subseteq tV$, to
$\operatorname{conv}(B)\subseteq tV\subseteq tU$, więc
otoczka wypukła jest ograniczona.
\end{uwaga}

\begin{przyklad}[Nieograniczona otoczka wypukła kuli jednostkowej]
\label{ex:quasinorm-convex-hull}
Niech $0<p<1$ oraz $X=L^p([0,1])$ z metryką
$d(f,g)=\int_{[0,1]}|f(t)-g(t)|^p\,\mathrm d\lambda(t)$.
Zbiór $B=\{f:\int_{[0,1]}|f(t)|^p\,\mathrm d\lambda(t)\le1\}$ jest
ograniczony: dla dowolnego $\varepsilon>0$ jego odpowiednio
mała homotetia mieści się w
$\{f:\int_{[0,1]}|f(t)|^p\,\mathrm d\lambda(t)<\varepsilon\}$.
Dla $n\ge1$ podzielmy $[0,1]$ na przedziały $I_{n,j}$,
$j=1,\ldots,n$, o długości $1/n$ i połóżmy
$f_{n,j}=n^{1/p}\mathbf1_{I_{n,j}}$. Wtedy
$\int_{[0,1]}|f_{n,j}(t)|^p\,\mathrm d\lambda(t)=1$,
natomiast ich średnia arytmetyczna
$g_n=n^{-1}\sum_{j=1}^n f_{n,j}=n^{1/p-1}\mathbf1_{[0,1]}$
należy do $\operatorname{conv}(B)$ i spełnia
\[
 \int_{[0,1]}|g_n(t)|^p\,\mathrm d\lambda(t)=n^{1-p}\longrightarrow\infty.
\]
Otoczka wypukła $B$ nie jest więc ograniczona. Dowodzi to
również, że $X$ nie ma bazy otoczeń zera złożonej ze
zbiorów wypukłych: w przeciwnym razie otoczka wypukła
każdego zbioru ograniczonego byłaby ograniczona.
\end{przyklad}

\subsection{Lokalna ograniczoność i jej granice}

\begin{definicja}[Przestrzeń lokalnie ograniczona]
Przestrzeń liniowo-topologiczną $X$ nazywamy
\emph{lokalnie ograniczoną}, jeżeli istnieje ograniczone
otoczenie zera w $X$.
\end{definicja}

\begin{stwierdzenie}[Kryterium lokalnej ograniczoności]
\label{prop:locally-bounded}
Dla TVS $X$ następujące warunki są równoważne:
\begin{enumerate}[label=\textup{(\roman*)}]
\item $X$ ma ograniczone otoczenie zera;
\item istnieje otoczenie zera $U$ takie, że dla każdego
otoczenia zera $V$ istnieje $t>0$ z $U\subseteq tV$;
\item istnieje baza otoczeń zera postaci $\{tU:t>0\}$
dla pewnego zrównoważonego otoczenia $U$.
\end{enumerate}
\end{stwierdzenie}
\begin{proof}
Równoważność (i) i (ii) jest treścią definicji.
Załóżmy (i) i wybierzmy zrównoważone otoczenie
$U$ zawarte w ograniczonym otoczeniu. Jest ono również
ograniczone. Każdy zbiór $tU$, $t>0$, jest otoczeniem
zera, a ograniczoność $U$ daje: dla dowolnego otoczenia
$V$ istnieje $s>0$ takie, że $U\subseteq sV$, czyli
$s^{-1}U\subseteq V$. Dowodzi to (iii).
Implikacja (iii)$\Rightarrow$(i) wynika z warunku, że
dla danego otoczenia $V$ pewna homotetia $tU$ jest w $V$.
\end{proof}

\begin{przyklad}[Dwa odmienne przypadki]
Przestrzeń $\K^{\N}$ z topologią produktową nie jest
lokalnie ograniczona. Każde jej otoczenie zera zawiera
zbiór, w którym tylko skończenie wiele współrzędnych jest
ograniczonych; poza tym zbiorem współrzędnych można wybrać
współrzędną $j$ całkowicie swobodną. Takie otoczenie nie
mieści się w żadnym wielokrotnym powiększeniu zbioru
$\{x:|x_j|<1\}$.

Z kolei $L^p([0,1])$, $0<p<1$, jest lokalnie ograniczona:
metryczna kula zera określona przez
$\int_{[0,1]}|f(t)|^p\,\mathrm d\lambda(t)<1$ jest ograniczona.
Nie jest jednak normowalna, ponieważ nie jest lokalnie
wypukła (przykład~\ref{ex:quasinorm-convex-hull}).
Tak więc \emph{lokalna ograniczoność nie oznacza normowalności}.
\end{przyklad}

\subsection{Rodziny operatorów: ograniczoność punktowa i równociągłość}

\begin{definicja}[Ograniczoność punktowa i równociągłość]
\label{def:operator-family}
Niech $X,Y$ będą TVS i niech
$\mathcal T\subseteq\mathcal L(X,Y)$.
Rodzinę $\mathcal T$ nazywamy \emph{punktowo ograniczoną}, jeśli
dla każdego $x\in X$ zbiór $\{Tx:T\in\mathcal T\}$
jest ograniczony w $Y$. Rodzinę nazywamy
\emph{równociągłą} (w zerze), jeśli dla każdego otoczenia
$V$ zera w $Y$ istnieje otoczenie $U$ zera w $X$ takie,
że $T(U)\subseteq V$ dla wszystkich $T\in\mathcal T$.
\end{definicja}

\begin{twierdzenie}[Równociągłość i zbiory ograniczone]
\label{thm:equi-bounded}
Niech $X,Y$ będą TVS i $\mathcal T\subseteq\mathcal L(X,Y)$.
Jeśli $\mathcal T$ jest równociągła, to dla każdego
ograniczonego $B\subseteq X$ zbiór $\mathcal T(B)$
jest ograniczony w $Y$. W szczególności każda rodzina
równociągła jest punktowo ograniczona.
\end{twierdzenie}
\begin{proof}
Ustalmy ograniczony $B\subseteq X$ i otoczenie $V$ zera
w $Y$. Z równociągłości istnieje otoczenie $U$ zera
w $X$ z $T(U)\subseteq V$ dla wszystkich $T\in\mathcal T$.
Ponieważ $B\subseteq tU$ dla pewnego $t>0$, mamy
$T(B)\subseteq tT(U)\subseteq tV$ dla każdego $T$.
Zatem $\mathcal T(B)\subseteq tV$.
\end{proof}

\begin{przyklad}[Punktowa ograniczoność bez równociągłości]
\label{ex:c00-uniform-boundedness-failure}
Niech $X=c_{00}$ będzie przestrzenią ciągów o skończonym
nośniku z normą $\|x\|_\infty=\sup_{k\ge1}|x_k|$.
Dla $n\ge1$ określmy ciągły funkcjonał
$T_n:X\to\K$, $T_nx=nx_n$. Każdy $x\in c_{00}$
ma skończony nośnik, zatem $T_nx=0$ dla wszystkich
dostatecznie dużych $n$. Rodzina $\{T_n:n\ge1\}$
jest więc punktowo ograniczona (a nawet $T_nx\to0$
dla każdego $x$). Nie jest równociągła: zbiór
$B=\{e_n:n\ge1\}$ (gdzie $e_n$ jest $n$-tym wektorem
standardowej bazy współrzędnościowej) jest ograniczony w $X$, lecz
$T_n(e_n)=n$ i $\mathcal T(B)$ jest nieograniczony w $\K$.

Ten przykład pokazuje, że punktowa ograniczoność rodzin
operatorów \emph{nie wystarcza} do równociągłości w dowolnej
przestrzeni normowanej. Zupełność dziedziny będzie kluczowa
w późniejszym twierdzeniu Banacha--Steinhausa.
\end{przyklad}

\begin{stwierdzenie}[Rodziny na przestrzeniach unormowanych]
\label{prop:equi-normed}
Niech $X,Y$ będą przestrzeniami unormowanymi, a
$\mathcal T\subseteq\mathcal L(X,Y)$. Wtedy $\mathcal T$
jest równociągła dokładnie wtedy, gdy
$\sup_{T\in\mathcal T}\|T\|<\infty$,
gdzie $\|T\|=\sup_{\|x\|\le1}\|Tx\|$.
\end{stwierdzenie}
\begin{proof}
Jeśli rodzina jest równociągła, to dla otoczenia
$V=\{y:\|y\|<1\}$ istnieje $r>0$ takie, że
$\|x\|<r$ pociąga $\|Tx\|<1$ dla wszystkich $T$.
Dla $\|x\|\le1$ zastosujmy to do $(r/2)x$:
$\|Tx\|<2/r$. Stąd supremum norm operatorowych jest
skończone.
Odwrotnie, jeśli $M=\sup_T\|T\|<\infty$, to dla
$\varepsilon>0$ wystarczy wybrać
$U=\{x:\|x\|<\varepsilon/(M+1)\}$.
\end{proof}

\subsection{Topologie zbieżności operatorów}

\begin{definicja}[Zbieżność punktowa i jednostajna na ograniczonych]
\label{def:operator-topologies}
Niech $X,Y$ będą TVS Hausdorffa. Na przestrzeni
$\mathcal L(X,Y)$ rozważamy dwie topologie:
\begin{enumerate}[label=\textup{(\roman*)}]
\item topologię $\tau_s$ \emph{zbieżności punktowej},
wyznaczoną odwzorowaniami $T\mapsto Tx$, $x\in X$;
\item topologię $\tau_b$ \emph{zbieżności jednostajnej na
zbiorach ograniczonych}, której bazę otoczeń zera
stanowią skończone przecięcia zbiorów
\[
 \mathcal N(B,V)=\{T\in\mathcal L(X,Y):T(B)\subseteq V\},
\]
gdzie $B\subseteq X$ jest ograniczony, a $V$ jest
otoczeniem zera w $Y$.
\end{enumerate}
\end{definicja}

\begin{twierdzenie}[Struktura liniowo-topologiczna przestrzeni operatorów]
\label{thm:op-topologies}
Niech $X,Y$ będą TVS Hausdorffa. Topologie $\tau_s$ i
$\tau_b$ czynią $\mathcal L(X,Y)$ przestrzenią
liniowo-topologiczną Hausdorffa, a $\tau_b$ jest silniejsza
od $\tau_s$. Dla sieci $(T_d)$ i $T\in\mathcal L(X,Y)$
zachodzi
\begin{align*}
 T_d\xrightarrow{\tau_s}T
 &\iff T_dx\to Tx\ \text{dla każdego }x\in X,\\
 T_d\xrightarrow{\tau_b}T
 &\iff \text{dla każdego ograniczonego }B\subseteq X
 \text{ i }V\in\mathcal N_Y(0)\\
 &\hspace{18mm}\text{ostatecznie }(T_d-T)(B)\subseteq V.
\end{align*}
\end{twierdzenie}
\begin{proof}
Topologia $\tau_s$ jest topologią podprzestrzeni
otrzymaną przez zanurzenie $T\mapsto(Tx)_{x\in X}$
w produkt $Y^X$, więc jest Hausdorffa i zgodna z
działaniami liniowymi.

Dla $\tau_b$ zbiory $\mathcal N(B,V)$ zawierają zero.
Wystarczy rozważać otwarte, zrównoważone otoczenia $V$,
ponieważ tworzą one bazę otoczeń zera w $Y$.
Dla takiego $V$ można dobrać otwarte, zrównoważone
$W$ z $W+W\subseteq V$;
wtedy
$\mathcal N(B,W)+\mathcal N(B,W)\subseteq\mathcal N(B,V)$.
Przecięcia wyraża się przez skończone układy warunków.
Zbiory te są zrównoważone, jeżeli $V$ jest zrównoważony.
Dla każdego $S\in\mathcal N(B,W)$ zachodzi
$S+\mathcal N(B,W)\subseteq\mathcal N(B,V)$.
Warunki te zapewniają istnienie topologii grupowej
z podaną bazą \emph{otoczeń} zera. Nie twierdzimy, że każdy
zbiór $\mathcal N(B,V)$ jest otwarty; może nim nie być.

Sprawdźmy łączną ciągłość mnożenia przez skalar.
Ustalmy $\lambda_0\in\K$, $T_0\in\mathcal L(X,Y)$,
ograniczony $B\subseteq X$ i otoczenie $V$ zera w $Y$.
Wybierzmy zrównoważone otoczenie $W$ z $W+W\subseteq V$.
Ponieważ $T_0(B)$ jest ograniczony, dla $\lambda$
dostatecznie bliskich $\lambda_0$ mamy
$(\lambda-\lambda_0)T_0(B)\subseteq W$.
W otoczeniu $\lambda_0$ liczby $|\lambda|$ są ograniczone
przez pewne $M>0$. Dobierzmy zrównoważone otoczenie
$U$ takie, że $\lambda U\subseteq W$ dla $|\lambda|\le M$.
Jeżeli $T-T_0\in\mathcal N(B,U)$, to
\[
 (\lambda T-\lambda_0T_0)(B)
 \subseteq\lambda(T-T_0)(B)+(\lambda-\lambda_0)T_0(B)
 \subseteq W+W\subseteq V.
\]
To dowodzi ciągłości łącznej. Ponieważ każdy singleton
$\{x\}\subseteq X$ jest ograniczony, otoczenia
$\mathcal N(\{x\},V)$ należą do bazy $\tau_b$,
a zatem $\tau_b\supseteq\tau_s$. Warunek Hausdorffa
wynika z faktu, że jeżeli $T\ne0$, to dla pewnego $x$
mamy $Tx\ne0$ i można oddzielić $Tx$ od zera w $Y$.
Opisy zbieżności wynikają wprost z definicji.
\end{proof}

\begin{stwierdzenie}[Ograniczoność rodzin w topologiach operatorów]
\label{prop:op-family-bounded}
Niech $X,Y$ będą TVS Hausdorffa i
$\mathcal T\subseteq\mathcal L(X,Y)$. Wówczas:
\begin{enumerate}[label=\textup{(\roman*)}]
\item $\mathcal T$ jest ograniczona w $(\mathcal L(X,Y),\tau_s)$
dokładnie wtedy, gdy jest punktowo ograniczona;
\item $\mathcal T$ jest ograniczona w $(\mathcal L(X,Y),\tau_b)$
dokładnie wtedy, gdy dla każdego ograniczonego $B\subseteq X$
zbiór $\mathcal T(B)$ jest ograniczony w $Y$.
\end{enumerate}
\end{stwierdzenie}
\begin{proof}
W (i) sprawdzamy ograniczoność we wszystkich ewaluacjach
$T\mapsto Tx$ i korzystamy z kryterium ograniczoności
w produkcie. W (ii) z definicji ograniczoności w
$\tau_b$ dla danego $B,V$ istnieje $t>0$ takie,
że $\mathcal T\subseteq t\mathcal N(B,V)$; jest to
równoważne zawieraniu $\mathcal T(B)\subseteq tV$.
Skończone przecięcia warunków nie zmieniają tego kryterium,
ponieważ bierzemy największą z odpowiednich stałych
skalowania po uprzednim zastąpieniu otoczeń zrównoważonymi.
\end{proof}

\begin{przyklad}[Dwie różne zbieżności operatorów]
Niech $X=c_{00}$ z normą $\ell^\infty$, a $Y=\K$.
Funkcjonały $T_nx=nx_n$ spełniają
$T_n\xrightarrow{\tau_s}0$, ale nie
$T_n\xrightarrow{\tau_b}0$, ponieważ
$\{e_n:n\ge1\}$ jest ograniczony i
$\sup_{x\in\{e_n:n\ge1\}}|T_nx|=n$.
\end{przyklad}

\begin{twierdzenie}[Przypadek przestrzeni unormowanych]
\label{thm:op-b-norm}
Jeżeli $X,Y$ są przestrzeniami unormowanymi, to topologia
$\tau_b$ na $\mathcal L(X,Y)$ jest topologią normy
operatorowej. Jeżeli $X$ ma dodatkowo wymiar skończony,
również $\tau_s=\tau_b$.
\end{twierdzenie}
\begin{proof}
Kula jednostkowa $B_X=\{x:\|x\|\le1\}$ jest ograniczona.
Warunek $(T_d-T)(B_X)\subseteq\{y:\|y\|<\varepsilon\}$
pociąga $\|T_d-T\|\le\varepsilon$. Wybierając $\varepsilon/2$ zamiast $\varepsilon$ widzimy, że $\tau_b$ jest co najmniej tak silna jak topologia normy. Odwrotnie, jeśli
$B\subseteq X$ jest ograniczony, to
$B\subseteq RB_X$ dla pewnego $R>0$, więc
$\|S\|<\varepsilon/R$ implikuje $S(B)\subseteq
\{y:\|y\|<\varepsilon\}$. W ten sposób otrzymujemy
tę samą topologię.

Jeśli $\dim X=0$, to $\mathcal L(X,Y)=\{0\}$
i obie topologie są jedyną topologią tej przestrzeni.
Załóżmy więc $\dim X=m\geq1$ i wybierzmy bazę
$e_1,\ldots,e_m$. Dla każdego $x=\sum_i a_ie_i$
zachodzą oszacowania
$|a_i|\le C\|x\|$ dla pewnej wspólnej stałej $C$:
aby to uzasadnić, rozważmy zbiór
$S=\{(a_1,\ldots,a_m):\max_i|a_i|=1\}$.
Jest on zwarty w $\K^m$, a ciągła funkcja
$(a_i)\mapsto\|\sum_i a_ie_i\|$ ma na nim dodatnie
minimum $c$. Skalowanie daje
$\|\sum_i a_ie_i\|\ge c\max_i|a_i|$, co
wystarcza przy $C=1/c$.
Jeżeli $\|S e_i\|<\delta$ dla każdego $i$, to
$\|Sx\|\le mC\delta\|x\|$, a zatem ograniczenie
operatora kontroluje skończenie wiele ewaluacji.
Wynika stąd $\tau_s=\tau_b$.
\end{proof}

\subsection{Twierdzenie o metryzowalności przestrzeni liniowo-topologicznych}

\begin{definicja}[Pierwszy aksjomat przeliczalności]
Przestrzeń topologiczną $X$ nazywamy
\emph{pierwszoprzeliczalną}, jeżeli każdy punkt ma
przeliczalną bazę otoczeń. W TVS wystarcza, aby
przeliczalna baza otoczeń istniała w zerze.
\end{definicja}

\begin{twierdzenie}[Kryterium metryzowalności TVS]
\label{thm:tvs-metrizable}
Niech $X$ będzie przestrzenią liniowo-topologiczną
Hausdorffa nad $\K$. Następujące warunki są równoważne:
\begin{enumerate}[label=\textup{(\roman*)}]
\item $X$ jest metryzowalna;
\item $X$ ma przeliczalną bazę otoczeń zera;
\item topologia $X$ jest wyznaczona przez pewną
      metrykę niezmienniczą względem przesunięć,
      tj. $d(x+a,y+a)=d(x,y)$ dla wszystkich $x,y,a\in X$.
\end{enumerate}
Co więcej, metryka z (iii) wyznacza kanoniczną strukturę
jednostajną $X$, a więc zupełność tej metryki jest
równoważna zupełności TVS.
\end{twierdzenie}
\begin{proof}
Implikacje $(\mathrm{iii})\Rightarrow(\mathrm{i})\Rightarrow(\mathrm{ii})$ są oczywiste:
w przestrzeni metrycznej kule o promieniach $1/n$
tworzą przeliczalną bazę otoczeń zera.

Załóżmy (ii). Wybierzmy malejącą bazę otwartych otoczeń
zera $(W_n)_{n\ge1}$ i skonstruujmy rekurencyjnie
otwarte, symetryczne i zrównoważone otoczenia $U_n$,
$n\ge0$, takie, że
$U_0=X$, $U_n\subseteq W_n$ dla $n\ge1$ oraz
$U_{n+1}+U_{n+1}+U_{n+1}\subseteq U_n$.
Jest to możliwe dzięki ciągłości dodawania i istnieniu
zrównoważonych otoczeń zera. Mamy $\bigcap_nU_n=\{0\}$:
przeciwne założenie przeczyłoby własności Hausdorffa.

Dla $x,y\in X$ połóżmy
\[
 d(x,y)=\inf\Big\{\sum_{j=1}^{m}2^{-n_j}:\
 m\ge1,\ x=x_0,\ldots,x_m=y,\
 x_{j-1}-x_j\in U_{n_j},\ n_j\ge0\Big\}.
\]
Infimum dotyczy wszystkich skończonych łańcuchów.
Ponieważ $U_0=X$, zbiór, z którego bierzemy infimum,
jest niepusty. Symetria zbiorów $U_n$ daje
$d(x,y)=d(y,x)$, a konkatenacja łańcuchów nierówność
trójkąta. Wzór zależy tylko od różnic $x-y$, więc
$d(x+a,y+a)=d(x,y)$.

Sprawdzimy, że $d$ jest dodatnio określona i generuje
topologię $X$. Potrzebujemy następującego rachunku:
\emph{jeżeli łańcuch ma łączną wagę mniejszą niż $2^{-n}$,
$n\ge1$, to różnica jego końców należy do $U_{n-1}$.}
Dowód prowadzimy indukcyjnie względem liczby jego
odcinków. Dla jednego odcinka warunek
$2^{-n_1}<2^{-n}$ daje $n_1>n$, czyli różnica należy do
$U_{n+1}\subseteq U_{n-1}$. Dla łańcucha dłuższego
wybierzmy odcinek przecinający połowę jego łącznej
wagi. Suma wag odcinków poprzedzających ten odcinek
oraz suma wag odcinków następujących po nim są
obie mniejsze niż $2^{-(n+1)}$. Z założenia
indukcyjnego ich całkowite różnice należą do $U_n$.
Waga wyróżnionego odcinka jest potęgą $2^{-n_j}$
i jest mniejsza niż $2^{-n}$, skąd $n_j\geq n+1$;
jego różnica należy więc do $U_{n_j}\subseteq U_n$.
Różnica końców całego łańcucha leży zatem w
$U_n+U_n+U_n\subseteq U_{n-1}$.

Dla każdego $n\ge1$ wynika stąd
\[
 U_{n+1}\subseteq\{x:d(x,0)<2^{-n}\}\subseteq U_{n-1}.
\]
Rzeczywiście pierwsza inkluzja odpowiada łańcuchowi
jednoodcinkowemu, a druga -- powyższemu rachunkowi.
Jeżeli $d(x,0)=0$, to $x$ należy do każdego $U_n$,
więc $x=0$. Metryka jest zatem dodatnio określona.
Dwie inkluzje dowodzą zarazem, że jej kule stanowią
tę samą bazę otoczeń zera co $U_n$; niezmienniczość
względem przesunięć daje zgodność topologii w każdym punkcie.

Otoczenia przekątnej dla metryki $d$ są postaci
$\{(x,y):d(x,y)<r\}$. Powyższe inkluzje pokazują,
że tworzą one dokładnie tę samą strukturę jednostajną,
co otoczenia $E_U=\{(x,y):x-y\in U\}$.
Sieci Cauchy'ego obu struktur są więc identyczne.
To dowodzi ostatniej części tezy.
\end{proof}

\begin{przyklad}[Metryzowalna, zupełna, ale nienormowalna]
\label{ex:metrizable-not-normable}
Na $X=\K^{\N}$ z topologią produktową określmy
\[
 d(x,y)=\sum_{n=1}^{\infty}2^{-n}\min\{1,|x_n-y_n|\}.
\]
Szereg jest zbieżny, $d$ jest metryką niezmienniczą
względem przesunięć i opisuje zbieżność współrzędną.
Przestrzeń $X$ jest zupełna jako produkt zupełnych
przestrzeni $\K$ (co wynika z zupełności produktów),
lecz nie jest normowalna, ponieważ nie ma ograniczonego
otoczenia zera. To pokazuje, że metryzowalność i
zupełność nie zastępują normowalności.
\end{przyklad}

\subsection{Jakie założenia będą potrzebne w następnych rozdziałach?}

Różnicę między ograniczonością punktową i równociągłością
usuwa w ważnych sytuacjach zasada jednostajnej ograniczoności;
jednym z jej naturalnych kontekstów są przestrzenie beczkowe.
Twierdzenie to wymaga dodatkowej struktury i zostanie
udowodnione później, a nie założone milcząco.

W kolejnym rozdziale wprowadzimy lokalną wypukłość,
seminormy i funkcjonał Minkowskiego. Dowiemy się między
innymi, dlaczego wypukłość otoczeń pozwala zamieniać
ograniczoność na oszacowania przez seminormy oraz
kiedy lokalna ograniczoność prowadzi już do normowalności.

% Wersja źródłowa: audyt III, jednostka 05
\setcounter{section}{4}
\section{Przestrzenie lokalnie wypukłe i seminormy}\label{sec:lokalnie-wypukle}

\subsection{Motywacja: od topologii do nierówności}

W poprzednich rozdziałach badaliśmy przestrzenie liniowo-topologiczne
bez wymagania, aby ich otoczenia były wypukłe. Okazało się, że nawet
zupełna i metryzowalna TVS może zachowywać się odmiennie od przestrzeni
unormowanej: w $L^p([0,1])$, $0<p<1$, otoczka wypukła ograniczonej
kuli nie jest ograniczona, a każdy ciągły funkcjonał liniowy jest zerowy.
Ten drugi fakt wynika z następującego argumentu. Niech
$L:L^p([0,1])\to\K$ będzie ciągłym funkcjonałem liniowym.
Dla $f\in L^p$ piszemy
$\|f\|_p=(\int_{[0,1]}|f|^p\,\mathrm d\lambda)^{1/p}$;
przy $p<1$ nie jest to norma. Z ciągłości $L$ w zerze istnieje $\delta>0$ takie, że
$\|g\|_p<\delta$ implikuje $|L(g)|<1$.
Dla $g\ne0$ podstawiamy $\delta g/(2\|g\|_p)$,
a z jednorodności dostajemy
$|L(g)|\le 2\|g\|_p/\delta$; nierówność ta
jest oczywista także dla $g=0$.
Przyjmujemy zatem $C=2/\delta$.
Jeśli $f\ne0$, to nieatomową miarę
$\nu_f(E)=\int_E|f|^p\,\mathrm d\lambda$ można dla
każdego $n\ge1$ podzielić na $n$ rozłącznych
zbiorów mierzalnych $E_1,\ldots,E_n$ o jednakowej
mierze $\nu_f(E_j)=\nu_f([0,1])/n$. Zatem
$f=\sum_{j=1}^n f\chi_{E_j}$ prawie wszędzie, a
\[
 |L(f)|\le\sum_{j=1}^n|L(f\chi_{E_j})|
 \le C n^{1-1/p}\|f\|_p\longrightarrow0.
\]
Ponieważ $1-1/p<0$, mamy $L(f)=0$ dla każdego $f$.
Wprowadzenie \emph{lokalnej wypukłości} pozwoli nam powiązać topologię
z oszacowaniami przez seminormy, a następnie z teorią dualności.

W całym rozdziale $\K\in\{\R,\C\}$. Jeżeli nie zaznaczono inaczej,
$X$ i $Y$ oznaczają przestrzenie liniowo-topologiczne nad $\K$;
nie zakładamy milcząco ani normowalności, ani metryzowalności.
Przez $\mathcal N_X(0)$ oznaczamy rodzinę otoczeń zera w $X$.
Zbiór $A\subseteq X$ nazywamy \emph{pochłaniającym}, jeśli
$\forall x\in X\ \exists t>0:\ x\in tA$.
W TVS każde otoczenie zera jest pochłaniające, ponieważ
$\lambda x\to0$ dla $\lambda\to0$.

\begin{definicja}[Wypukłość i zrównoważenie]
Niech $X$ będzie przestrzenią liniową nad $\K$.
Zbiór $A\subseteq X$ nazywamy \emph{wypukłym}, gdy
$(1-t)x+ty\in A$ dla wszystkich $x,y\in A$, $t\in[0,1]$.
Nazywamy go \emph{zrównoważonym}, gdy $\lambda A\subseteq A$
dla każdej $\lambda\in\K$ z $|\lambda|\le1$.
Zbiór jednocześnie wypukły i zrównoważony nazywamy
\emph{absolutnie wypukłym}.
\end{definicja}

\begin{stwierdzenie}[Wypukłość absolutna jako nierówność na sumę]
\label{prop:absolutna}
Niech $X$ będzie przestrzenią liniową nad $\K$ i
$A\subseteq X$ zbiorem niepustym. Jest on absolutnie wypukły
wtedy i tylko wtedy, gdy dla wszystkich $x,y\in A$ oraz
$\alpha,\beta\in\K$ spełniających $|\alpha|+|\beta|\le1$
zachodzi $\alpha x+\beta y\in A$.
\end{stwierdzenie}
\begin{proof}
Jeśli $A$ jest absolutnie wypukły, to
$\alpha x+\beta y=(|\alpha|+|\beta|)z$ przy
$z=\frac{|\alpha|}{|\alpha|+|\beta|}\frac{\alpha}{|\alpha|}x+
\frac{|\beta|}{|\alpha|+|\beta|}\frac{\beta}{|\beta|}y$
(wyrazy z zerowym współczynnikiem pomijamy).
Wypukłość i zrównoważenie dają $z\in A$ oraz
$\alpha x+\beta y\in A$. Gdy oba współczynniki są zerowe,
$0\in A$ wynika ze zrównoważenia. Odwrotnie, kładąc
$\beta=0$ otrzymujemy zrównoważenie, a kładąc
$\alpha=t$, $\beta=1-t$ dla $t\in[0,1]$, wypukłość.
\end{proof}

\begin{definicja}[Przestrzeń lokalnie wypukła]
\label{def:lcs}
Przestrzeń liniowo-topologiczną $X$ nazywamy
\emph{lokalnie wypukłą} (LCS), jeśli zero ma bazę otoczeń
złożoną ze zbiorów wypukłych. Przez \emph{lokalnie wypukłą
przestrzeń Hausdorffa} rozumiemy LCS, której topologia jest
Hausdorffa; rozróżnienie to jest istotne w twierdzeniach o normowalności.
\end{definicja}

\begin{stwierdzenie}[Absolutnie wypukła baza otoczeń]
\label{prop:absconv-base}
Niech $X$ będzie TVS. Przestrzeń $X$ jest lokalnie wypukła
wtedy i tylko wtedy, gdy zero ma bazę otoczeń absolutnie wypukłych.
Bazę można wybrać ze zbiorów otwartych.
\end{stwierdzenie}
\begin{proof}
Załóżmy lokalną wypukłość i ustalmy otoczenie $U$ zera.
Wybierzmy wypukłe otoczenie $V\subseteq U$, a następnie
otwarte zrównoważone otoczenie $O\subseteq V$.
Zbiór $\operatorname{conv}(O)$ jest zawarty w $V$,
bo $V$ jest wypukły; jest też zrównoważony,
ponieważ dla $|\lambda|\le1$ zachodzi
$\lambda\operatorname{conv}(O)=\operatorname{conv}(\lambda O)
\subseteq\operatorname{conv}(O)$.
Zbiór $\operatorname{conv}(O)$ jest absolutnie wypukły.
Ponadto jest otwarty: jeżeli $x\in\operatorname{conv}(O)$,
to $x+O\subseteq\operatorname{conv}(O)+O\subseteq
2\operatorname{conv}(O)$ nie wystarcza do żądanej tezy;
można natomiast zapisać $x=\sum_{j=1}^m t_jx_j$
z $t_j\ge0$, $\sum t_j=1$, i wybrać $j$ z $t_j>0$.
Ponieważ $O$ jest otwarty, dla dostatecznie małego
otoczenia $N$ zera mamy $x_j+t_j^{-1}N\subseteq O$,
stąd $x+N\subseteq\operatorname{conv}(O)$.
Druga implikacja wynika bezpośrednio z definicji.
\end{proof}

\begin{uwaga}[Dlaczego lokalna wypukłość jest istotna?]
W przestrzeni $L^p([0,1])$ dla $0<p<1$ nie ma bazy
wypukłych otoczeń zera; dowód podano w
przykładzie z \hyperref[sec:rodziny-metryzowalnosc]{rozdziału 4}.
Tak więc zupełność i metryzowalność nie wymuszają lokalnej wypukłości.
\end{uwaga}

\subsection{Seminormy i topologie zadane przez rodziny seminorm}

\begin{definicja}[Seminorma]
Niech $X$ będzie przestrzenią liniową nad $\K$.
\emph{Seminormą} na $X$ nazywamy funkcję
$p:X\to[0,\infty)$ spełniającą
$p(\lambda x)=|\lambda|p(x)$ dla $\lambda\in\K$, $x\in X$,
oraz $p(x+y)\le p(x)+p(y)$ dla $x,y\in X$.
Jeżeli dodatkowo $p(x)=0$ implikuje $x=0$, to $p$ jest normą.
\end{definicja}

\begin{przyklad}[Pierwsze seminormy]
Na $\K^{\N}$ funkcje $p_n(x)=|x_n|$ są seminormami, ale
nie normami. Na przestrzeni $C^\infty([0,1])$ seminormami są
$p_m(f)=\max_{0\le j\le m}\sup_{t\in[0,1]}|f^{(j)}(t)|$,
$m\in\N\cup\{0\}$. Na $C(\Omega)$ dla otwartego
$\Omega\subseteq\R^d$ seminormą jest
$p_K(f)=\sup_{x\in K}|f(x)|$ dla dowolnego zwartego $K\subseteq\Omega$.
\end{przyklad}

\begin{twierdzenie}[Topologia generowana przez seminormy]
\label{thm:seminorm-topology}
Niech $X$ będzie przestrzenią liniową nad $\K$, a
$\mathcal P$ niepustą rodziną seminorm na $X$. Istnieje
najmniejsza topologia $\tau_{\mathcal P}$ \emph{zgodna ze
strukturą przestrzeni liniowo-topologicznej}, w której
wszystkie seminormy $p\in\mathcal P$ są ciągłe.
Topologia ta czyni $X$ przestrzenią lokalnie wypukłą,
a bazą otoczeń zera są zbiory
\[
 U(p_1,\ldots,p_m;\varepsilon)
 =\{x\in X:\max_{1\le j\le m}p_j(x)<\varepsilon\},
 \quad m\ge1,\quad p_j\in\mathcal P,\quad\varepsilon>0.
\]
Jest ona Hausdorffa wtedy i tylko wtedy, gdy
$\bigcap_{p\in\mathcal P}\ker p=\{0\}$.
Ponadto sieć $(x_d)$ zbiega do $x$ w $\tau_{\mathcal P}$
dokładnie wtedy, gdy $p(x_d-x)\to0$ dla każdego $p\in\mathcal P$.
\end{twierdzenie}
\begin{proof}
Zbiory z tezy są absolutnie wypukłe, pochłaniające i
zamknięte na skończone przecięcia (po ewentualnym
zmniejszeniu $\varepsilon$). Dla zbioru z promieniem
$\varepsilon$ suma dwóch analogicznych zbiorów z promieniem
$\varepsilon/2$ mieści się w wyjściowym zbiorze.
Aby uzasadnić ciągłość mnożenia przez skalar w punkcie
$(\lambda_0,x_0)$, dla każdego z ustalonych $p_j$
używamy oszacowania
$p_j(\lambda x-\lambda_0x_0)\le
|\lambda|p_j(x-x_0)+|\lambda-\lambda_0|p_j(x_0)$.
Wystarczy ograniczyć $|\lambda|$ w małym otoczeniu
$\lambda_0$, a następnie wybrać odpowiednio małe
$|\lambda-\lambda_0|$ i $p_j(x-x_0)$ dla wszystkich
skończenie wielu indeksów $j$.
W ten sposób bazy przesuniętych zbiorów definiują
topologię TVS.
Każde $p\in\mathcal P$ jest ciągłe w zerze, a z nierówności
$|p(x)-p(y)|\le p(x-y)$ wynika jej ciągłość wszędzie.
Każda inna topologia TVS, w której seminormy są ciągłe,
zawiera wszystkie powyższe otoczenia, zatem jest silniejsza.

Jeżeli $x\ne0$ i $p(x)>0$ dla pewnego $p$, otoczenie
$\{y:p(y)<p(x)/2\}$ oddziela $0$ od $x$.
Jeśli natomiast $p(x)=0$ dla wszystkich $p$, to $x$
należy do każdego otoczenia zera. W konsekwencji
$\tau_{\mathcal P}$ jest Hausdorffa dokładnie przy
podanym warunku. Opis zbieżności wynika z bazy otoczeń.
\end{proof}

\subsection{Funkcjonał Minkowskiego}

\begin{definicja}[Funkcjonał Minkowskiego]
\label{def:gauge}
Niech $X$ będzie przestrzenią liniową, a $A\subseteq X$
niepustym zbiorem pochłaniającym. Funkcjonałem Minkowskiego
(\emph{funkcjonałem cechującym}) zbioru $A$ nazywamy
\[
 p_A(x)=\inf\{t>0:x\in tA\},\qquad x\in X.
\]
Pochłanianie zapewnia, że infimum jest skończone.
\end{definicja}

\begin{twierdzenie}[Od zbioru absolutnie wypukłego do seminormy]
\label{thm:minkowski}
Niech $X$ będzie TVS nad $\K$, a $A\subseteq X$
pochłaniającym zbiorem absolutnie wypukłym.
Wówczas $p_A$ jest seminormą i
$\{x:p_A(x)<1\}\subseteq A\subseteq\{x:p_A(x)\le1\}$.
Jeśli dodatkowo $A$ jest \emph{otwarty}, to
$A=\{x:p_A(x)<1\}$ i $p_A$ jest ciągła.
\end{twierdzenie}
\begin{proof}
Pochłanianie oraz $0\in A$ dają $p_A(x)<\infty$ i
$p_A(0)=0$. Dla $\lambda\ne0$ mamy
$x\in tA$ dokładnie wtedy, gdy $\lambda x\in|\lambda|tA$:
przy zespolonym $\lambda$ korzystamy ze zrównoważenia
$A$ i z równości $\frac{\lambda}{|\lambda|}A=A$.
Otrzymujemy $p_A(\lambda x)=|\lambda|p_A(x)$;
przypadek $\lambda=0$ jest osobny i natychmiastowy.

Niech $s>p_A(x)$ oraz $t>p_A(y)$. Z definicji infimum
i z własności $rA\subseteq sA$ dla $0<r\le s$
wynika $x\in sA$ i $y\in tA$. Wypukłość daje
$sA+tA\subseteq(s+t)A$, zatem
$p_A(x+y)\le s+t$. Przechodząc z $s$ i $t$ do
odpowiednich infimów otrzymujemy podaddytywność.
Jeżeli $p_A(x)<1$, to $x\in tA$ dla pewnego $t<1$,
więc $x\in A$ dzięki zrównoważeniu.
Jeżeli $x\in A$, to $p_A(x)\le1$.

Niech teraz $A$ będzie otwarty i niech $x\in A$.
Ciągłość $\lambda\mapsto\lambda x$ w $\lambda=1$
daje $\lambda>1$ dostatecznie bliskie $1$ z
$\lambda x\in A$; stąd $p_A(x)\le1/\lambda<1$.
Zatem $A=\{p_A<1\}$.
Ponieważ $\{p_A<\varepsilon\}=\varepsilon A$
dla każdego $\varepsilon>0$, seminorma jest ciągła
w zerze, a nierówność
$|p_A(x)-p_A(y)|\le p_A(x-y)$ zapewnia ciągłość wszędzie.
\end{proof}

\begin{center}
\begin{tikzpicture}[scale=.85,>=Latex]
 \draw[fill=blue!5,draw=blue!50!black,thick] (0,0) ellipse (2.3 and 1.35);
 \draw[->,gray] (-2.8,0)--(4.35,0) node[right] {$x_1$};
 \draw[->,gray] (0,-1.65)--(0,1.9) node[above] {$x_2$};
 \fill (0,0) circle (1.8pt) node[below left] {$0$};
 \draw[->,thick] (0,0)--(3.45,1.12) node[above right] {$x$};
 \fill[blue!70!black] (2.012,.653) circle (2.3pt);
 \node[blue!60!black,above] at (1.50,.84) {$x/p_A(x)$};
 \node[blue!60!black] at (-1.45,.7) {$A$};
\end{tikzpicture}

\small Rysunek: dla absolutnie wypukłego otwartego otoczenia $A$
dla punktu $x$ z $p_A(x)>0$ wektor $x/p_A(x)$ leży
na brzegu $A$; wartość $p_A(x)$ określa współczynnik skalowania.
\end{center}

\begin{przyklad}[Seminorma nie musi być normą]
Dla $X=\R^2$ i $A=\{(x,y):|x|<1\}$ zachodzi
$p_A(x,y)=|x|$. Zbiór $A$ jest otwarty, absolutnie
wypukły i pochłaniający, lecz nie jest ograniczony,
a $p_A(0,1)=0$. Takie zdegenerowanie nie przeczy
Twierdzeniu~\ref{thm:minkowski}.
\end{przyklad}

\begin{przyklad}[Pochłanianie nie gwarantuje ciągłości]
Niech $X=c_{00}$ z normą $\|x\|_\infty$ oraz
$A=\{x\in c_{00}:\sup_{n\ge1}n|x_n|\le1\}$.
Każdy $x\in c_{00}$ ma skończony nośnik, więc $A$
jest pochłaniający i absolutnie wypukły. Jego funkcjonał
to $p_A(x)=\sup_{n\ge1}n|x_n|$. Nie jest on ciągły
w topologii normy supremum: $x^{(n)}=e_n/n\to0$,
ale $p_A(x^{(n)})=1$. W szczególności $A$ nie jest
otoczeniem zera w wyjściowej topologii.
\end{przyklad}

\subsection{Charakteryzacja lokalnej wypukłości przez seminormy}

\begin{twierdzenie}[Charakteryzacja lokalnej wypukłości]
\label{thm:lcs-seminorms}
Dla dowolnej TVS $X$ nad $\K$ następujące warunki są
równoważne: (i) $X$ jest lokalnie wypukła;
(ii) topologia $X$ jest generowana przez rodzinę
ciągłych seminorm. W przypadku Hausdorffa rodzinę
można wybrać rozdzielającą punkty.
\end{twierdzenie}
\begin{proof}
Implikację (ii)$\Rightarrow$(i) daje
Twierdzenie~\ref{thm:seminorm-topology}.
Załóżmy (i) i weźmy rodzinę $\mathcal U$ wszystkich
otwartych absolutnie wypukłych otoczeń zera.
Dla każdego $U\in\mathcal U$ określmy funkcjonał
Minkowskiego $p_U$ z poprzedniej podsekcji. Dowód
Twierdzenia~\ref{thm:minkowski} pokazuje, że $p_U$
jest ciągłą seminormą i $U=\{p_U<1\}$. Zatem
rodzina $(p_U)_{U\in\mathcal U}$ generuje wyjściową
topologię. Jeśli $X$ jest Hausdorffa, dla $x\ne0$
możemy wybrać $U\in\mathcal U$ niezawierające $x$,
a wówczas $p_U(x)\ge1$. Żaden niezerowy wektor
nie należy do wszystkich jąder.
\end{proof}


\subsection{Ciągłość odwzorowań liniowych i ograniczoność}

\begin{twierdzenie}[Kryterium ciągłości przez seminormy]
\label{thm:continuous-sem}
Niech $X,Y$ będą przestrzeniami lokalnie wypukłymi,
$\mathcal P$ i $\mathcal Q$ rodzinami seminorm generującymi
odpowiednio ich topologie, a $T:X\to Y$ odwzorowaniem liniowym.
Następujące warunki są równoważne:
\begin{enumerate}[label=\textup{(\roman*)}]
\item $T$ jest ciągłe;
\item dla każdej $q\in\mathcal Q$ istnieją $m\ge1$,
$p_1,\ldots,p_m\in\mathcal P$ i $C>0$ takie, że
$q(Tx)\le C\max_{1\le j\le m}p_j(x)$ dla wszystkich $x\in X$.
\end{enumerate}
\end{twierdzenie}
\begin{proof}
Załóżmy ciągłość $T$. Dla $q\in\mathcal Q$
znajdźmy $p_1,\ldots,p_m$ i $\delta>0$ takie, że
$\max_jp_j(x)<\delta$ implikuje $q(Tx)<1$.
Oznaczmy $r(x)=\max_jp_j(x)$. Jeżeli $r(x)=0$,
to $r(tx)=0$ dla każdego $t>0$, a więc
$tq(Tx)<1$ dla każdego $t>0$; stąd $q(Tx)=0$.
Jeżeli $r(x)>0$, to dla
$t=\delta/(2r(x))$ mamy $r(tx)=\delta/2$,
a zatem $q(Tx)<2r(x)/\delta$. Uzyskujemy (ii)
z $C=2/\delta$. Odwrotnie, każda z nierówności
z (ii) zapewnia ciągłość $q\circ T$ w zerze;
opis bazy otoczeń z Twierdzenia~\ref{thm:seminorm-topology}
w $Y$ daje ciągłość $T$.
\end{proof}

\begin{twierdzenie}[Ograniczoność przez seminormy]
\label{thm:bounded-sem}
Niech $X$ będzie przestrzenią lokalnie wypukłą,
a $\mathcal P$ rodziną seminorm generującą topologię $X$.
Zbiór $B\subseteq X$ jest ograniczony w sensie von Neumanna
wtedy i tylko wtedy, gdy
$\sup_{x\in B}p(x)<\infty$ dla każdego $p\in\mathcal P$
(przyjmujemy $\sup\varnothing=0$).
\end{twierdzenie}
\begin{proof}
Jeśli $B$ jest ograniczony, to dla danego $p$ i
$V=\{x:p(x)<1\}$ istnieje $t>0$ z $B\subseteq tV$,
więc $\sup_Bp\le t$.
Odwrotnie, niech $U$ będzie otoczeniem zera.
Wybierzmy $p_1,\ldots,p_m\in\mathcal P$ i
$\varepsilon>0$ z
$\{x:\max_jp_j(x)<\varepsilon\}\subseteq U$.
Niech $M=\max_j\sup_{x\in B}p_j(x)<\infty$.
Dla $t>(M/\varepsilon)$ mamy
$B\subseteq t\{x:\max_jp_j(x)<\varepsilon\}\subseteq tU$.
\end{proof}

\begin{wniosek}[Otoczki ograniczone]
\label{cor:absconv-bounded}
Jeżeli $X$ jest przestrzenią lokalnie wypukłą,
a $B\subseteq X$ jest ograniczony, to jego otoczka wypukła
i jego otoczka absolutnie wypukła są ograniczone.
\end{wniosek}
\begin{proof}
Dla każdej ciągłej seminormy $p$ i dla kombinacji
$z=\sum_{j=1}^m\lambda_jx_j$ z $x_j\in B$ oraz
$\sum_j|\lambda_j|\le1$ zachodzi
$p(z)\le\sup_{x\in B}p(x)$. Takie kombinacje
wypełniają otoczkę absolutnie wypukłą $B$.
Stosujemy Twierdzenie~\ref{thm:bounded-sem}.
\end{proof}

\subsection{Twierdzenie Kolmogorowa o normowalności}

\begin{twierdzenie}[Kryterium normowalności]
\label{thm:kolmogorov-normability}
Niech $X$ będzie lokalnie wypukłą przestrzenią
liniowo-topologiczną Hausdorffa nad $\K$.
Następujące warunki są równoważne:
\begin{enumerate}[label=\textup{(\roman*)}]
\item topologia $X$ pochodzi od pewnej normy na $X$;
\item $X$ ma ograniczone otoczenie zera;
\item $X$ ma ograniczone, otwarte i absolutnie wypukłe otoczenie zera.
\end{enumerate}
\end{twierdzenie}
\begin{proof}
Jeżeli $X$ jest unormowana, jej otwarta kula jednostkowa
jest otoczeniem ograniczonym. To daje (i)$\Rightarrow$(ii).
Dzięki lokalnej wypukłości wewnątrz ograniczonego
otoczenia możemy wybrać otwarte absolutnie wypukłe
otoczenie $U$, które jest ograniczone jako podzbiór
zbioru ograniczonego. Mamy (ii)$\Rightarrow$(iii).

Załóżmy (iii). Funkcjonał Minkowskiego $p_U$ jest
ciągłą seminormą i $U=\{p_U<1\}$.
Pokażemy, że jest normą. Jeśli $p_U(x)=0$,
to dla każdego $s>0$ zachodzi $x\in sU$.
Niech $V$ będzie dowolnym otoczeniem zera w $X$.
Ograniczoność $U$ daje $t>0$ z $U\subseteq tV$,
zatem $x\in t^{-1}U\subseteq V$.
Ponieważ $X$ jest Hausdorffa, przecięcie wszystkich
otoczeń zera jest równe $\{0\}$, czyli $x=0$.

Topologia normy $p_U$ nie jest silniejsza od wyjściowej,
ponieważ $p_U$ jest ciągła. Odwrotnie, dla dowolnego
otoczenia $V$ ograniczoność $U$ daje
$t^{-1}U\subseteq V$, a zbiór ten jest kulą
$\{p_U<1/t\}$. Zatem obie topologie są równe.
\end{proof}

\begin{uwaga}[Gdzie konieczna jest własność Hausdorffa?]
Jeżeli lokalna wypukłość jest spełniona, ale przestrzeń
nie jest Hausdorffa, funkcjonał Minkowskiego ograniczonego
otoczenia może być tylko seminormą z niezerowym jądrem.
Topologia generowana przez seminormę nieoddzielającą punktów
nie może pochodzić od normy.
\end{uwaga}

\subsection{Metryzowalność lokalnie wypukła i przestrzenie Frécheta}

\begin{twierdzenie}[Przeliczalna rodzina seminorm]
\label{thm:countable-seminorms}
Niech $X$ będzie lokalnie wypukłą TVS Hausdorffa.
Jest ona metryzowalna wtedy i tylko wtedy, gdy jej topologię
generuje przeliczalna rodzina seminorm $(p_n)_{n\ge1}$.
W takim przypadku zgodną, niezmienniczą względem przesunięć
metryką jest
\[
 d(x,y)=\sum_{n=1}^{\infty}2^{-n}\min\{1,p_n(x-y)\}.
\]
\end{twierdzenie}
\begin{proof}
Jeżeli istnieje przeliczalna rodzina $(p_n)$,
to każda funkcja $\min\{1,p_n(x-y)\}$ spełnia
nierówność trójkąta, a szereg jest jednostajnie zbieżny
(bo jego wyrazy nie przekraczają $2^{-n}$).
Ponieważ seminormy rozdzielają punkty, $d$ jest metryką.
Jeśli $x_k\to x$ w topologii seminorm, to dla ustalonego
$N$ pierwszych $N$ składników sumy dążą do zera,
a ogon jest jednostajnie mniejszy od $2^{-N}$;
stąd $d(x_k,x)\to0$. W drugą stronę z
$d(x_k,x)\ge2^{-n}\min\{1,p_n(x_k-x)\}$
wynika zbieżność każdej seminormy. Ten sam argument
zamiast ciągów stosuje się do otoczeń i dowodzi
równości topologii.

Jeżeli $X$ jest metryzowalna, posiada przeliczalną bazę
otoczeń zera (kryterium metryzowalności TVS z \hyperref[sec:rodziny-metryzowalnosc]{rozdziału 4}).
Do każdego jej elementu wybierzmy, korzystając z
Twierdzenia~\ref{thm:lcs-seminorms}, skończoną rodzinę
ciągłych seminorm i wspólny promień tak, aby odpowiadający
zbiór był zawarty w danym otoczeniu. Suma tych skończonych
rodzin jest przeliczalna i generuje topologię $X$.
\end{proof}

\begin{definicja}[Przestrzeń Frécheta]
\label{def:frechet-space}
\emph{Przestrzenią Frécheta} nazywamy przestrzeń
liniowo-topologiczną, która jest lokalnie wypukła,
metryzowalna, zupełna i Hausdorffa. Nie wymagamy, by
metryka generująca topologię była zadana normą.
\end{definicja}

\begin{przyklad}[Produkt, który nie jest normowalny]
\label{ex:product-lcs}
Niech $X=\K^{\N}$ z topologią produktową.
Generują ją seminormy $p_n(x)=|x_n|$.
Topologia jest Hausdorffa, metryzowalna, a przestrzeń
jest zupełna jako produkt zupełnych przestrzeni $\K$.
Jest zatem przestrzenią Frécheta. Nie jest natomiast
normowalna: żadne otoczenie zera nie jest ograniczone,
co pokazaliśmy w \hyperref[sec:rodziny-metryzowalnosc]{rozdziale 4}.
\end{przyklad}

\begin{przyklad}[Gładkie funkcje i nieskończenie wiele pochodnych]
\label{ex:smooth-frechet-preview}
Niech $X=C^\infty([0,1])$ z topologią generowaną przez
$p_m(f)=\max_{0\le j\le m}\sup_{t\in[0,1]}|f^{(j)}(t)|$,
$m=0,1,\ldots$. Seminormy te rozdzielają punkty,
a topologia jest metryzowalna i lokalnie wypukła.
Zupełność wykażemy w następnym stwierdzeniu. Już teraz
można dowieść \emph{nienormowalności}. Gdyby $U$ było ograniczonym
otoczeniem zera, zawierałoby
$V=\{f:p_m(f)<\delta\}$ dla pewnych $m$ i $\delta>0$;
zbiór $V$ byłby więc ograniczony. Tymczasem dla
$N\ge2\pi$ funkcje
$f_N(t)=\frac{\delta}{2N^m}\sin(Nt)$
spełniają $p_m(f_N)\le\delta/2$, natomiast
$p_{m+1}(f_N)\ge\delta N/2\to\infty$.
To przeczy Twierdzeniu~\ref{thm:bounded-sem}.
Przestrzeń $C^\infty([0,1])$ jest więc lokalnie wypukła
oraz metryzowalna, ale nie normowalna.
\end{przyklad}


\subsection{Sieci Cauchy'ego i zupełność w języku seminorm}

\begin{twierdzenie}[Kryterium Cauchy'ego w przestrzeniach lokalnie wypukłych]
\label{thm:lcs-cauchy}
Niech $X$ będzie przestrzenią lokalnie wypukłą, a
$\mathcal P$ rodziną seminorm generującą topologię $X$.
Sieć $(x_d)_{d\in D}$ w $X$ jest siecią Cauchy'ego
w kanonicznej strukturze jednostajnej TVS dokładnie wtedy,
gdy dla każdej $p\in\mathcal P$ i każdego $\varepsilon>0$
istnieje $d_0\in D$ takie, że $p(x_d-x_e)<\varepsilon$
dla wszystkich $d,e\succeq d_0$.
W szczególności $X$ jest zupełna dokładnie wtedy, gdy każda
taka sieć jest zbieżna w $X$.
\end{twierdzenie}
\begin{proof}
Jeżeli sieć jest Cauchy'ego, zastosujmy definicję z
\hyperref[sec:zupelnosc]{rozdziału 3} do otoczenia $\{x:p(x)<\varepsilon\}$.
Otrzymujemy żądany warunek dla każdej seminormy.
Na odwrót, niech $U$ będzie dowolnym otoczeniem zera.
Znajdźmy seminormy $p_1,\ldots,p_m$ i $\varepsilon>0$
takie, że $\{x:\max_jp_j(x)<\varepsilon\}\subseteq U$.
Dla każdego $j$ mamy indeks $d_j$, po którym
$p_j(x_d-x_e)<\varepsilon$; skierowanie $D$ daje wspólny
indeks $d_0\succeq d_j$, $j=1,\ldots,m$.
Wtedy $x_d-x_e\in U$ dla $d,e\succeq d_0$.
Ostatnie zdanie jest dokładnie definicją zupełności TVS.
\end{proof}

\begin{twierdzenie}[Zupełność przestrzeni funkcji gładkich]
\label{thm:smooth-complete}
Niech $X=C^\infty([0,1])$ nad $\K$ z topologią generowaną
przez seminormy $p_m(f)=\max_{0\le j\le m}
\sup_{t\in[0,1]}|f^{(j)}(t)|$, $m=0,1,\ldots$.
Wtedy $X$ jest przestrzenią Frécheta, ale nie jest normowalna.
\end{twierdzenie}
\begin{proof}
Lokalna wypukłość, warunek Hausdorffa, metryzowalność
i brak normowalności zostały wykazane powyżej.
Dla zupełności, na mocy twierdzenia o metryzowalnych
TVS z \hyperref[sec:zupelnosc]{rozdziału 3}, wystarczy wykazać zupełność ciągową.

Niech $(f_n)$ będzie ciągiem Cauchy'ego względem wszystkich
$p_m$. Dla każdego $j\ge0$ ciąg funkcji ciągłych
$(f_n^{(j)})$ jest jednostajnie Cauchy'ego na $[0,1]$.
Zupełność $C([0,1])$ w normie supremum daje
funkcję ciągłą $g_j$ taką, że $f_n^{(j)}\to g_j$
jednostajnie. Korzystając z podstawowego twierdzenia
rachunku całkowego, dla $t\in[0,1]$ mamy
\[
 f_n^{(j)}(t)-f_n^{(j)}(0)
 =\int_{[0,t]}f_n^{(j+1)}(s)\,\mathrm d\lambda(s).
\]
Jednostajna zbieżność całkowanych funkcji i wartości
w zerze pozwala przejść do granicy:
$g_j(t)-g_j(0)=\int_{[0,t]}g_{j+1}(s)\,\mathrm d\lambda(s)$.
Ponieważ $g_{j+1}$ jest ciągła, wynika stąd
$g_j'=g_{j+1}$. Indukcyjnie $g_0$ jest gładka i
$g_0^{(j)}=g_j$ dla wszystkich $j\ge0$.
Ponadto $p_m(f_n-g_0)\to0$ dla każdego $m$.
Oznacza to, że ciąg Cauchy'ego ma granicę w $X$,
a więc $X$ jest zupełna.
\end{proof}

\begin{uwaga}[Rzut na skończenie wiele pochodnych]
Z punktu widzenia konstrukcji przestrzeni funkcyjnych
$C^\infty([0,1])$ można uważać za zgodną rodzinę
funkcji należących do przestrzeni Banacha $C^m([0,1])$,
$m\ge0$, z normami $p_m$. Każdy operator
$f\mapsto f^{(j)}(t_0)$, $t_0\in[0,1]$, jest ciągłym
funkcjonałem liniowym na tej przestrzeni, gdyż
$|f^{(j)}(t_0)|\le p_j(f)$. Tę strukturę będziemy
wykorzystywać przy definiowaniu dystrybucji.
\end{uwaga}

\subsection{Podprzestrzenie, produkty i ilorazy}

\begin{stwierdzenie}[Trwałość lokalnej wypukłości]
\label{prop:lcs-constructions}
Niech $(X_i)_{i\in I}$ będzie rodziną przestrzeni lokalnie
wypukłych nad $\K$. Wówczas:
\begin{enumerate}[label=\textup{(\roman*)}]
\item każda podprzestrzeń liniowa $M\subseteq X_i$ z topologią
podprzestrzeni jest lokalnie wypukła;
\item produkt $\prod_{i\in I}X_i$ z topologią produktową
jest lokalnie wypukły;
\item dla dowolnej podprzestrzeni liniowej $M\subseteq X_i$
iloraz $X_i/M$ z topologią ilorazową jest lokalnie wypukły.
Iloraz jest Hausdorffa wtedy i tylko wtedy, gdy $M$ jest domknięta.
\end{enumerate}
\end{stwierdzenie}
\begin{proof}
W (i) przecinamy wypukłe otoczenia zera w $X_i$
z $M$. W (ii) podstawowe otoczenia zera są produktami,
w których jedynie skończenie wiele współrzędnych
jest ograniczonych do wypukłych otoczeń; takie produkty
są wypukłe. W (iii) naturalne rzutowanie
$q:X_i\to X_i/M$ jest otwarte, a obrazy wypukłych
otoczeń zera są wypukłymi otoczeniami zera w ilorazie.
Kryterium Hausdorffa udowodniono w
\hyperref[sec:konstrukcje]{rozdziale 2}.
\end{proof}

\begin{uwaga}[Dalszy kierunek]
Związek topologii z seminormami jest punktem wyjścia
teorii dualności. W następnym rozdziale sformułujemy i
udowodnimy twierdzenie Hahna--Banacha, które pozwoli
przejść od seminorm i rozdzielania punktów do istnienia
nietrywialnych ciągłych funkcjonałów liniowych.
Później wprowadzimy topologie słabe, przestrzenie Frécheta,
przestrzenie funkcji testowych $\mathcal D(\Omega)$ oraz ich
przestrzeń sprzężoną $\mathcal D'(\Omega)$.
\end{uwaga}

% Wersja źródłowa: audyt II, jednostka 06
\setcounter{section}{5}
\section{Twierdzenie Hahna--Banacha i ciągły dual}\label{sec:hahn-banach}

\subsection{Dlaczego potrzebujemy twierdzenia o przedłużaniu?}

W skończonym wymiarze każdy funkcjonał liniowy jest ciągły, a jego
przedłużenie z podprzestrzeni można skonstruować, uzupełniając bazę.
W ogólnej przestrzeni liniowo-topologicznej postępowanie takie daje
jedynie funkcjonał \emph{algebraiczny}: jego ciągłość nie jest
zagwarantowana. Twierdzenie Hahna--Banacha pozwala zachować
oszacowanie wyrażone przez funkcję podliniową, a w przestrzeniach
lokalnie wypukłych --- także ciągłość funkcjonału.

W całym rozdziale $\K\in\{\R,\C\}$. Przestrzeń liniowa nad $\K$
nie musi być wyposażona w topologię, dopóki nie zaznaczymy tego
jawnie. Nie przyjmujemy bez dodatkowego założenia, że przestrzeń
liniowo-topologiczna jest Hausdorffa.

\begin{definicja}[Dual algebraiczny i topologiczny]\label{def:dual-topological}
Niech $X$ będzie przestrzenią liniową nad $\K$. Przez $X^*$
oznaczamy przestrzeń wszystkich $\K$-liniowych funkcjonałów
$f:X\to\K$ (\emph{dual algebraiczny}). Jeśli $X$ jest TVS,
przez $X'$ oznaczamy przestrzeń funkcjonałów liniowych ciągłych
(\emph{dual topologiczny}). Wtedy $X'\subseteq X^*$.
\end{definicja}

\begin{stwierdzenie}[Elementarne kryterium ciągłości]
\label{prop:dual-continuity}
Niech $X$ będzie TVS nad $\K$ i $f:X\to\K$ funkcjonałem
liniowym. Wtedy $f$ jest ciągły wtedy i tylko wtedy, gdy
istnieje otoczenie $U$ zera w $X$ takie, że $|f(x)|<1$
dla każdego $x\in U$.
\end{stwierdzenie}
\begin{proof}
Ciągłość w zerze zapewnia istnienie otoczenia
$U=f^{-1}(\{z\in\K:|z|<1\})$ z wymaganą własnością.
W drugą stronę, dla $\varepsilon>0$ zbiór
$\varepsilon U$ jest otoczeniem zera i dla $x\in\varepsilon U$
mamy $|f(x)|<\varepsilon$. To dowodzi ciągłości w zerze.
Liniowość daje następnie ciągłość w każdym punkcie.
\end{proof}

\begin{przyklad}[Algebraiczny dual może być znacznie większy]
\label{ex:discontinuous-functional}
Niech $X$ będzie nieskończenie wymiarową przestrzenią
unormowaną. Wybierzmy liniowo niezależne wektory $(e_n)_{n\ge1}$
z $\|e_n\|=1$. Na ich rozpiętości określmy
$f(\sum_{j=1}^m a_je_j)=\sum_{j=1}^m j a_j$.
Po uzupełnieniu tej rodziny do bazy Hamela przestrzeni $X$
przedłużamy $f$ liniowo na całe $X$. Wtedy $f\in X^*$,
ale $f\notin X'$, ponieważ $|f(e_n)|=n\to\infty$,
podczas gdy $\|e_n\|=1$. Konstrukcja bazy Hamela
wykorzystuje aksjomat wyboru.
\end{przyklad}

\subsection{Wersja rzeczywista: funkcje podliniowe}

\begin{definicja}[Funkcja podliniowa]\label{def:sublinear}
Niech $E$ będzie przestrzenią liniową nad $\R$.
Funkcję $p:E\to\R$ nazywamy \emph{podliniową}, jeśli
$p(x+y)\le p(x)+p(y)$ dla wszystkich $x,y\in E$ oraz
$p(tx)=tp(x)$ dla wszystkich $x\in E$ i $t\ge0$.
Nie wymagamy symetrii $p(-x)=p(x)$ ani nieujemności $p(x)$.
Z definicji wynika $p(0)=0$.
\end{definicja}

\begin{lemat}[Przedłużenie o jeden wymiar]
\label{lem:hb-one-step}
Niech $E$ będzie rzeczywistą przestrzenią liniową,
$p:E\to\R$ funkcją podliniową, $M\subseteq E$ podprzestrzenią,
a $f:M\to\R$ funkcjonałem liniowym spełniającym
$f(m)\le p(m)$ dla każdego $m\in M$. Dla każdego
$x_0\in E\setminus M$ istnieje funkcjonał liniowy
$g:M+\R x_0\to\R$, który przedłuża $f$ i spełnia
$g(x)\le p(x)$ dla każdego $x\in M+\R x_0$.
\end{lemat}
\begin{proof}
Każdy $x\in M+\R x_0$ ma jednoznaczny zapis $x=m+tx_0$.
Szukamy liczby $a\in\R$ takiej, by wzór
$g(m+tx_0)=f(m)+ta$ spełniał $g\le p$.
Dla $t>0$ warunek ten, po podzieleniu przez $t$,
jest równoważny $a\le p(z+x_0)-f(z)$ dla każdego $z\in M$.
Dla $t<0$ analogicznie otrzymujemy
$a\ge f(y)-p(y-x_0)$ dla każdego $y\in M$.
Zdefiniujmy więc
\[
 L=\sup_{y\in M}\bigl(f(y)-p(y-x_0)\bigr),\qquad
 R=\inf_{z\in M}\bigl(p(z+x_0)-f(z)\bigr).
\]
Dla dowolnych $y,z\in M$, z liniowości $f$ oraz
podaddytywności $p$ mamy
\[
 f(y)+f(z)=f(y+z)\le p(y+z)
 \le p(y-x_0)+p(z+x_0).
\]
Stąd $L\le R$. Ponieważ $0\in M$, mamy też
$L\ge-p(-x_0)$ i $R\le p(x_0)$, zatem oba krańce
są skończone. Wybierzmy dowolne $a\in[L,R]$.
Wówczas wyprowadzone nierówności dają $g(m+tx_0)\le
p(m+tx_0)$, odpowiednio dla $t>0$ i $t<0$.
Dla $t=0$ korzystamy z założenia $f\le p$.
\end{proof}

\begin{twierdzenie}[Hahn--Banach, wersja rzeczywista]
\label{thm:hb-real}
Niech $E$ będzie przestrzenią liniową nad $\R$,
$p:E\to\R$ funkcją podliniową, $M\subseteq E$ podprzestrzenią,
a $f:M\to\R$ funkcjonałem liniowym takim, że
$f(m)\le p(m)$ dla każdego $m\in M$.
Wtedy istnieje liniowy $F:E\to\R$ spełniający
$F|_M=f$ oraz $F(x)\le p(x)$ dla każdego $x\in E$.
\end{twierdzenie}
\begin{proof}
Rozważmy zbiór par $(N,g)$, gdzie $N$ jest podprzestrzenią
$E$ zawierającą $M$, a $g:N\to\R$ jest liniowy, przedłuża
$f$ i spełnia $g\le p$ na $N$. Jest on niepusty, ponieważ
zawiera $(M,f)$. Porządkujemy pary przez przedłużanie:
$(N_1,g_1)\preceq(N_2,g_2)$, gdy $N_1\subseteq N_2$
i $g_2|_{N_1}=g_1$.

Jeżeli $\mathcal C$ jest łańcuchem, to $N_\mathcal C=
\bigcup_{(N,g)\in\mathcal C}N$ jest podprzestrzenią:
dwa dowolne jej elementy należą do jednego z dwóch
porównywalnych członów łańcucha. Funkcjonały z łańcucha
zgadzają się na przecięciach, więc określają wspólny
funkcjonał $g_\mathcal C$ na $N_\mathcal C$.
Jest on liniowy, przedłuża $f$ i nadal spełnia
$g_\mathcal C\le p$. Zatem każdy łańcuch ma ograniczenie
górne. Lemat Kuratowskiego--Zorna daje element maksymalny
$(N_0,g_0)$. Gdyby $N_0\ne E$, lemat
\ref{lem:hb-one-step} przedłużyłby $g_0$ do
$N_0+\R x_0$ dla $x_0\notin N_0$, co przeczyłoby
maksymalności. Wobec tego $N_0=E$ i $F=g_0$.
\end{proof}

\begin{uwaga}[O roli aksjomatu wyboru]
Dowód odwołuje się do lematu Kuratowskiego--Zorna,
który w zwykłej teorii mnogości ZFC wynika z aksjomatu
wyboru. Przy konstrukcji przedłużeń w skończonym wymiarze
wystarcza skończona iteracja Lematu~\ref{lem:hb-one-step}.
\end{uwaga}

\subsection{Wersja zespolona i dominacja przez seminormę}

W przypadku zespolonym nie porównujemy liczb zespolonych
nierównością $F(x)\le p(x)$. Zamiast tego wymagamy
$|F(x)|\le p(x)$, gdzie $p$ jest seminormą.

\begin{lemat}[Odtworzenie funkcjonału zespolonego]
\label{lem:complex-real-part}
Niech $E$ będzie przestrzenią zespoloną traktowaną także
jako przestrzeń rzeczywista. Każdemu funkcjonałowi
rzeczywiście liniowemu $u:E\to\R$ odpowiada dokładnie
jeden zespolono-liniowy funkcjonał $G:E\to\C$ spełniający
$\operatorname{Re}G=u$. Jest on dany wzorem
$G(x)=u(x)-i\,u(ix)$.
\end{lemat}
\begin{proof}
Z liniowości rzeczywistej $u$ wynika
$G(ix)=u(ix)-i u(-x)=iG(x)$ oraz
$G(ax+by)=aG(x)+bG(y)$ dla $a,b\in\R$.
Każdy skalar zespolony ma postać $a+ib$, więc $G$ jest
zespolono-liniowy. Ponadto $\operatorname{Re}G(x)=u(x)$.
Jeśli $H$ jest innym funkcjonałem zespolono-liniowym
o tej samej części rzeczywistej, to
$\operatorname{Im}H(x)=-\operatorname{Re}H(ix)=-u(ix)$;
stąd $H(x)=G(x)$.
\end{proof}

\begin{twierdzenie}[Hahn--Banach, dominacja przez seminormę]
\label{thm:hb-complex-seminorm}
Niech $E$ będzie przestrzenią liniową nad
$\K\in\{\R,\C\}$, $p:E\to[0,\infty)$ seminormą,
$M\subseteq E$ podprzestrzenią liniową, a
$f:M\to\K$ funkcjonałem $\K$-liniowym spełniającym
$|f(m)|\le p(m)$ dla każdego $m\in M$.
Wtedy istnieje $\K$-liniowy $F:E\to\K$ taki, że
$F|_M=f$ i $|F(x)|\le p(x)$ dla każdego $x\in E$.
\end{twierdzenie}
\begin{proof}
Dla $\K=\R$ stosujemy Twierdzenie~\ref{thm:hb-real}
do $p$ i $f$; dominacja $F(x)\le p(x)$ zastosowana
również do $-x$ daje $|F(x)|\le p(x)$.

Niech $\K=\C$. Na przestrzeni rzeczywistej
$E_{\R}$ funkcja $u=\operatorname{Re}f$ jest liniowa
na $M_{\R}$ i spełnia $u(m)\le p(m)$.
Twierdzenie~\ref{thm:hb-real} zapewnia przedłużenie
rzeczywiście liniowe $v:E_{\R}\to\R$ z $v\le p$.
Z Lematu~\ref{lem:complex-real-part} otrzymujemy
$F(x)=v(x)-iv(ix)$. Na $M$ funkcjonały $F$ i $f$
mają tę samą część rzeczywistą, zatem są równe.

Pozostaje sprawdzić \emph{moduł} oszacowania. Dla
$x\in E$ z $F(x)\ne0$ wybierzmy $\alpha\in\C$,
$|\alpha|=1$, w taki sposób, aby
$\alpha F(x)=|F(x)|$. Wówczas
\[
 |F(x)|=\operatorname{Re}F(\alpha x)
 =v(\alpha x)\le p(\alpha x)=p(x).
\]
Gdy $F(x)=0$, nierówność jest oczywista.
\end{proof}

\begin{uwaga}[Co zostaje zachowane przy przedłużaniu?]
Przedłużenie zachowuje dominację przez \emph{tę samą}
seminormę, bez dodatkowego współczynnika. Samo twierdzenie
jest algebraiczne: $E$ nie musi być przestrzenią topologiczną.
Dopiero ciągłość $p$ daje automatycznie ciągłość $F$.
\end{uwaga}

\subsection{Hahn--Banach w przestrzeniach lokalnie wypukłych}

\begin{stwierdzenie}[Ciągłość funkcjonału a seminorma]
\label{prop:dual-seminorm-criterion}
Niech $X$ będzie lokalnie wypukłą TVS nad $\K$,
a $f:X\to\K$ funkcjonałem liniowym. Wtedy
$f\in X'$ dokładnie wtedy, gdy istnieje ciągła
seminorma $p$ na $X$ taka, że
$|f(x)|\le p(x)$ dla każdego $x\in X$.
\end{stwierdzenie}
\begin{proof}
Jeśli $f$ jest ciągły, to $p(x)=|f(x)|$ jest
ciągłą seminormą i $|f|\le p$. Odwrotnie, gdy
$|f|\le p$ i $p$ jest ciągła, to
$f^{-1}(\{z:|z|<\varepsilon\})$ zawiera otwarte
otoczenie $\{x:p(x)<\varepsilon\}$ dla każdego
$\varepsilon>0$. Wynika stąd ciągłość $f$.
\end{proof}

\begin{twierdzenie}[Topologiczny Hahn--Banach]\label{thm:hb-topological}
Niech $X$ będzie lokalnie wypukłą przestrzenią
liniowo-topologiczną nad $\K\in\{\R,\C\}$,
$M\subseteq X$ podprzestrzenią liniową z topologią
indukowaną z $X$, a $f:M\to\K$ ciągłym funkcjonałem
liniowym. Istnieje $F\in X'$ takie, że $F|_M=f$.
Nie trzeba zakładać, że $M$ jest domknięta, a $X$
jest Hausdorffa lub zupełna.
\end{twierdzenie}
\begin{proof}
Ciągłość $f$ w zerze oznacza, że istnieje otoczenie
$W$ zera w $X$ takie, iż $|f(m)|<1$ dla wszystkich
$m\in M\cap W$. Korzystając z lokalnej wypukłości,
wybieramy otwarte absolutnie wypukłe otoczenie
$U\subseteq W$. Jego funkcjonał Minkowskiego $p_U$
jest ciągłą seminormą na $X$ na mocy
twierdzenia o funkcjonale Minkowskiego z \hyperref[sec:lokalnie-wypukle]{rozdziału 5}.

Jeżeli $m\in M$ i $t>p_U(m)$, to $m\in tU$,
zatem $m/t\in M\cap U$ i $|f(m)|<t$.
Ponieważ jest to prawdą dla każdego $t>p_U(m)$,
uzyskujemy $|f(m)|\le p_U(m)$.
Stosując Twierdzenie~\ref{thm:hb-complex-seminorm}
do $p_U$ i podprzestrzeni $M$, otrzymujemy funkcjonał
$F:X\to\K$ z $F|_M=f$ i $|F|\le p_U$.
Ciągłość $p_U$ implikuje ciągłość $F$.
\end{proof}

\begin{wniosek}[Przedłużanie z podprzestrzeni gęstej]
\label{cor:dense-extension}
Jeśli w założeniach Twierdzenia~\ref{thm:hb-topological}
podprzestrzeń $M$ jest gęsta w $X$, to ciągłe
przedłużenie $f$ jest jednoznaczne.
\end{wniosek}
\begin{proof}
Istnienie wynika z twierdzenia. Jeśli $F_1,F_2\in X'$
przedłużają $f$, to ciągły funkcjonał $F_1-F_2$
zeruje się na gęstym zbiorze $M$. Jego zbiór zerowy
jest domknięty (bo $\K$ jest Hausdorffa), więc
$F_1=F_2$ na całym $X$.
\end{proof}

\begin{przyklad}[Sumowanie w $\ell^1$ i jednoznaczność]
Niech $X=\ell^1(\N)$ z normą
$\|x\|_1=\sum_{n\ge1}|x_n|$, a $M=c_{00}\subseteq X$
niech oznacza przestrzeń ciągów o skończonym nośniku.
Przez $e_n$ oznaczamy ciąg mający $1$ na miejscu $n$
i zera na pozostałych miejscach. Funkcjonał $f(x)=\sum_{n\ge1}x_n$ jest dobrze określony
na $M$ i ma normę $1$, bo $|f(x)|\le\|x\|_1$
i $f(e_1)=\|e_1\|_1=1$. Podprzestrzeń $M$ jest gęsta
w $X$: obcięcia $x^{(N)}=(x_1,\ldots,x_N,0,\ldots)$
spełniają $\|x^{(N)}-x\|_1\to0$. Jedynym ciągłym
przedłużeniem $f$ jest więc
$F(x)=\sum_{n=1}^\infty x_n$ na $\ell^1$.
Dla porównania, ten sam wzór na $c_{00}$ z normą
supremum nie jest ciągły: dla
$x^{(N)}=\sum_{n=1}^N e_n$ zachodzi
$\|x^{(N)}\|_\infty=1$ oraz $f(x^{(N)})=N$.
\end{przyklad}

\begin{przyklad}[Funkcjonały na przestrzeni funkcji gładkich]
Niech $X=C^\infty([0,1])$ z topologią seminorm
$p_m$ z \hyperref[sec:lokalnie-wypukle]{rozdziału 5}. Dla $t_0\in[0,1]$ i
$j\in\N\cup\{0\}$ funkcjonał
$f\mapsto f^{(j)}(t_0)$ jest ciągły, gdyż
$|f^{(j)}(t_0)|\le p_j(f)$. Każdy taki funkcjonał
będzie później przykładem dystrybucji na odpowiedniej
przestrzeni funkcji testowych.
\end{przyklad}

\subsection{Przestrzenie unormowane: zachowanie normy}

\begin{twierdzenie}[Hahn--Banach dla przestrzeni unormowanych]
\label{thm:hb-normed}
Niech $X$ będzie przestrzenią unormowaną nad $\K$,
$M\subseteq X$ podprzestrzenią liniową, a
$f:M\to\K$ ograniczonym funkcjonałem liniowym
(względem normy odziedziczonej z $X$).
Istnieje funkcjonał ograniczony $F:X\to\K$ taki, że
$F|_M=f$ i $\|F\|=\|f\|$, gdzie
$\|f\|=\sup\{|f(m)|:\|m\|\le1,\ m\in M\}$.
\end{twierdzenie}
\begin{proof}
W Twierdzeniu~\ref{thm:hb-complex-seminorm}
weźmy $p(x)=\|f\|\,\|x\|$. Otrzymujemy $F$
z oszacowaniem $|F(x)|\le\|f\|\,\|x\|$,
czyli $\|F\|\le\|f\|$. Ponieważ $F|_M=f$,
mamy także $\|F\|\ge\|f\|$. Zatem normy są równe.
Nie wymagamy zupełności $X$ ani domkniętości $M$.
\end{proof}

\begin{wniosek}[Funkcjonał normujący wektor]\label{cor:norming}
Niech $X$ będzie przestrzenią unormowaną nad $\K$,
a $x_0\in X\setminus\{0\}$. Istnieje $f\in X'$
taki, że $\|f\|=1$ i $f(x_0)=\|x_0\|$.
W szczególności, dla każdego $x\in X$ zachodzi
\[
 \|x\|=\sup\bigl\{|f(x)|:f\in X',\ \|f\|\le1\bigr\}.
\]
\end{wniosek}
\begin{proof}
Na $M=\K x_0$ określmy
$f_0(\lambda x_0)=\lambda\|x_0\|$.
Wówczas $\|f_0\|=1$ i z
Twierdzenia~\ref{thm:hb-normed} otrzymujemy przedłużenie
$f$ o normie $1$ i wartości $f(x_0)=\|x_0\|$.
Dla dowolnego $x$ nierówność $|g(x)|\le\|x\|$
dla $\|g\|\le1$ wynika z definicji normy operatorowej,
a dla $x\ne0$ równość zostaje osiągnięta przez
skonstruowany funkcjonał. Dla $x=0$ obie strony są zerowe.
\end{proof}

\begin{przyklad}[Przedłużenie nie musi być jednoznaczne]
Niech $X=\R^2$ z normą $\|(x,y)\|_\infty
=\max(|x|,|y|)$, $M=\{(t,t):t\in\R\}$, a
$f(t,t)=t$. Wtedy $\|f\|=1$. Dla każdego
$a\in[0,1]$ wzór
$F_a(x,y)=a x+(1-a)y$ określa przedłużenie $f$
spełniające $\|F_a\|=|a|+|1-a|=1$.
Twierdzenie Hahna--Banacha nie daje zatem
jednoznaczności bez dodatkowych założeń.
\end{przyklad}

\subsection{Separacja punktów, podprzestrzeni i rola lokalnej wypukłości}

\begin{twierdzenie}[Seminorma jako supremum funkcjonałów]
\label{thm:seminorm-as-sup}
Niech $X$ będzie lokalnie wypukłą TVS nad $\K$,
$p$ ciągłą seminormą na $X$, a $x_0\in X$.
Wtedy
\[
 p(x_0)=\sup\bigl\{|f(x_0)|:f\in X',\quad
 |f(x)|\le p(x)\ \text{dla wszystkich }x\in X\bigr\}.
\]
Ponadto istnieje $f\in X'$ spełniający $|f|\le p$
i $f(x_0)=p(x_0)$.
\end{twierdzenie}
\begin{proof}
Jeśli $p(x_0)=0$, wybieramy funkcjonał zerowy.
Gdy $p(x_0)>0$, na $M=\K x_0$ określamy
$f_0(\lambda x_0)=\lambda p(x_0)$.
Z jednorodności seminormy wynika
$|f_0(\lambda x_0)|=p(\lambda x_0)$.
Twierdzenie~\ref{thm:hb-complex-seminorm} daje
przedłużenie $f$ z $|f|\le p$, a ciągłość $p$
gwarantuje $f\in X'$. Żaden funkcjonał w podanym
supremum nie może mieć wartości większej niż $p(x_0)$,
a skonstruowany osiąga tę wartość.
\end{proof}

\begin{wniosek}[Dual rozdziela punkty]\label{cor:dual-separates}
Niech $X$ będzie lokalnie wypukłą TVS Hausdorffa
nad $\K$. Dla każdych różnych $x,y\in X$ istnieje
$f\in X'$ takie, że $f(x)\ne f(y)$.
\end{wniosek}
\begin{proof}
Dla $z=x-y\ne0$ istnieje ciągła seminorma $p$
z $p(z)>0$, ponieważ seminormy generujące topologię
rozdzielają punkty (twierdzenie o topologii seminorm z \hyperref[sec:lokalnie-wypukle]{rozdziału 5}).
Twierdzenie~\ref{thm:seminorm-as-sup} daje $f\in X'$
z $f(z)=p(z)>0$, co kończy dowód.
\end{proof}

\begin{twierdzenie}[Oddzielanie punktu od domkniętej podprzestrzeni]
\label{thm:point-closed-subspace}
Niech $X$ będzie lokalnie wypukłą TVS nad $\K$,
$M\subseteq X$ domkniętą podprzestrzenią liniową,
a $x_0\in X\setminus M$. Istnieje $f\in X'$
takie, że $f|_M=0$ i $f(x_0)=1$.
Nie jest potrzebne założenie Hausdorffa całej przestrzeni $X$.
\end{twierdzenie}
\begin{proof}
Ponieważ $M$ jest domknięta, istnieje otoczenie
$W$ zera takie, że $(x_0+W)\cap M=\varnothing$.
Wybierzmy otwarte absolutnie wypukłe otoczenie
$U\subseteq W$ i oznaczmy jego ciągłą seminormę
Minkowskiego przez $p=p_U$.
Dla każdego $m\in M$ mamy $x_0-m\notin U$,
a więc $p(x_0-m)\ge1$, bo $U=\{x:p(x)<1\}$.
Na podprzestrzeni $M+\K x_0$ określmy
$g(m+\lambda x_0)=\lambda$. Dla $\lambda\ne0$,
jednorodność $p$ oraz $-m/\lambda\in M$ dają
\[
 p(m+\lambda x_0)=|\lambda|\,
 p\bigl(x_0+m/\lambda\bigr)\ge|\lambda|
 =|g(m+\lambda x_0)|.
\]
Dla $\lambda=0$ nierówność też zachodzi.
Z Twierdzenia~\ref{thm:hb-complex-seminorm} otrzymujemy
przedłużenie $f$ z $|f|\le p$. Jest ono ciągłe,
zeruje się na $M$ i spełnia $f(x_0)=1$.
\end{proof}

\begin{wniosek}[Charakteryzacja domkniętych podprzestrzeni]
\label{cor:closed-as-annihilator}
Niech $X$ będzie lokalnie wypukłą TVS nad $\K$,
a $M\subseteq X$ podprzestrzenią liniową. Wtedy
\[
 \overline M=\bigcap\{\ker f:f\in X',\ f|_M=0\}.
\]
W szczególności $M$ jest domknięta dokładnie wtedy,
gdy stanowi przecięcie jąder pewnej rodziny ciągłych
funkcjonałów liniowych.
\end{wniosek}
\begin{proof}
Dla każdego $f\in X'$ zerującego się na $M$
jego domknięte jądro zawiera $\overline M$.
Jeśli $x_0\notin\overline M$, zastosujmy
Twierdzenie~\ref{thm:point-closed-subspace}
do domkniętej podprzestrzeni $\overline M$.
Otrzymujemy funkcjonał z $f(x_0)=1$, który
wyklucza $x_0$ z podanego przecięcia.
\end{proof}

\begin{przyklad}[Dlaczego nie możemy pominąć lokalnej wypukłości?]
Niech $X=L^p([0,1])$, $0<p<1$, z topologią metryki
$d(f,g)=\int_{[0,1]}|f(t)-g(t)|^p\,\mathrm d\lambda(t)$.
Z \hyperref[sec:tvs]{rozdziału 1} wiemy, że $X'=\{0\}$.
Tymczasem $M=\K\mathbf 1$ jest jednowymiarową
podprzestrzenią, a funkcjonał $f(c\mathbf1)=c$
jest na niej ciągły: $d(c\mathbf1,0)=|c|^p$.
Nie istnieje ciągłe przedłużenie $f$ na $X$.
To pokazuje, że topologiczna wersja Hahna--Banacha
wymaga dodatkowych założeń; lokalna wypukłość jest
wystarczającym, lecz nie zawsze koniecznym warunkiem.
\end{przyklad}

\begin{stwierdzenie}[Dual topologiczny przestrzeni produktowej]
\label{prop:product-dual}
Niech $I$ będzie dowolnym zbiorem, a $X=\K^I$
przestrzenią liniowo-topologiczną z topologią produktową.
Każdy funkcjonał $f\in X'$ ma postać
$f(x)=\sum_{i\in I_0}a_i x_i$, gdzie $I_0\subseteq I$
jest skończony, a $a_i\in\K$. Odwrotnie, każdy taki
funkcjonał jest ciągły.
\end{stwierdzenie}
\begin{proof}
Niech $f\in X'$. Z ciągłości w zerze istnieje
podstawowe otoczenie produktowe $U$ zera, na którym
$|f(x)|<1$. Możemy wybrać je tak, aby ograniczało
współrzędne tylko dla indeksów z pewnego skończonego
$I_0\subseteq I$. Jeśli $x_i=0$ dla $i\in I_0$,
to $tx\in U$ dla każdego $t\in\K$, zatem
$|t f(x)|<1$ dla wszystkich $t$; koniecznie $f(x)=0$.
Niech $e_i\in\K^I$ oznacza wektor o współrzędnej
$1$ na pozycji $i$ i zerach na pozostałych pozycjach.
Każde $x\in X$ jest sumą wektora
$\sum_{i\in I_0}x_i e_i$ i wektora o zerowych
współrzędnych w $I_0$. Stąd
$f(x)=\sum_{i\in I_0}x_i f(e_i)$.
W drugą stronę, rzuty współrzędnych $x\mapsto x_i$
są ciągłe z definicji topologii produktowej;
skończona kombinacja liniowa też jest ciągła.
\end{proof}

\begin{uwaga}[Jak silne jest to ograniczenie?]
Gdy $I=\N$, na przestrzeni wszystkich ciągów
$\K^\N$ z topologią produktową żaden ciągły funkcjonał
nie może zależeć istotnie od nieskończenie wielu wyrazów.
Inaczej jest na $\ell^1$: funkcjonał sumowania z
poprzedniego przykładu jest ciągły.
Różnica pochodzi z topologii, nie z samej liniowości.
\end{uwaga}

\subsection{Norma odległości i anihilatory}

\begin{definicja}[Anihilator]\label{def:annihilator}
Niech $X$ będzie TVS nad $\K$, a $M\subseteq X$
podprzestrzenią liniową. Jej \emph{anihilatorem}
w dualu topologicznym nazywamy
$M^\perp=\{f\in X':f(m)=0\ \text{dla wszystkich }m\in M\}$.
Zapis $M^\perp$ nie wymaga iloczynu skalarnego.
\end{definicja}

\begin{stwierdzenie}[Dual ilorazu]\label{prop:quotient-dual}
Niech $X$ będzie TVS nad $\K$, a $M\subseteq X$
podprzestrzenią liniową. Niech $q:X\to X/M$
będzie rzutowaniem z topologią ilorazową.
Odwzorowanie $q^*:(X/M)'\to X'$ określone przez
$q^*(\varphi)=\varphi\circ q$ jest liniową bijekcją
z $(X/M)'$ na $M^\perp$.
\end{stwierdzenie}
\begin{proof}
Jeśli $\varphi$ jest ciągłym funkcjonałem na ilorazie,
to $\varphi\circ q$ jest ciągły i zeruje się na $M$.
Odwzorowanie $q^*$ jest różnowartościowe, bo $q$
jest surjekcją. Jeżeli $f\in M^\perp$, to wzór
$\varphi(x+M)=f(x)$ jest dobrze określony, liniowy
i spełnia $\varphi\circ q=f$.
Z własności topologii ilorazowej (\hyperref[sec:konstrukcje]{rozdział 2})
wynika ciągłość $\varphi$, gdyż $f$ jest ciągły.
\end{proof}

\begin{twierdzenie}[Odległość od podprzestrzeni i dual]\label{thm:distance-dual}
Niech $X$ będzie przestrzenią unormowaną nad $\K$,
$M\subseteq X$ domkniętą podprzestrzenią liniową,
a $x_0\in X$. Wtedy
\[
 \operatorname{dist}(x_0,M):=\inf_{m\in M}\|x_0-m\|
 =\sup\{|f(x_0)|:f\in M^\perp,\ \|f\|\le1\}.
\]
Jeśli $x_0\notin M$, istnieje $f\in M^\perp$
o normie $1$, dla którego $f(x_0)=\operatorname{dist}(x_0,M)$.
\end{twierdzenie}
\begin{proof}
Niech $d=\operatorname{dist}(x_0,M)$.
Dla $f\in M^\perp$ z $\|f\|\le1$ i $m\in M$
mamy $|f(x_0)|=|f(x_0-m)|\le\|x_0-m\|$.
Po wzięciu infimum po $m$ otrzymujemy $|f(x_0)|\le d$.
Jeśli $x_0\in M$, to $d=0$ i funkcjonał zerowy
wystarcza do równości.

Załóżmy $x_0\notin M$. Ponieważ $M$ jest domknięta,
$d>0$. Na $M+\K x_0$ określmy
$g(m+\lambda x_0)=\lambda d$.
Dla $\lambda\ne0$ mamy
$\|m+\lambda x_0\|=|\lambda|\,
\|x_0+m/\lambda\|\ge|\lambda|d=|g(m+\lambda x_0)|$;
dla $\lambda=0$ jest to oczywiste.
Z Twierdzenia~\ref{thm:hb-normed} otrzymujemy
przedłużenie $f$ z $\|f\|\le1$, $f|_M=0$ i
$f(x_0)=d$. Równość $\|f\|=1$ wynika z faktu, że
$\|g\|=1$: dla $\lambda=1$ i $m\in M$ ilorazy
$|g(x_0-m)|/\|x_0-m\|=d/\|x_0-m\|$
mają supremum $1$ z definicji $d$.
\end{proof}

\begin{uwaga}[Następne rozdziały]
Rozdzielanie \emph{ogólnych zbiorów wypukłych}, twierdzenie
bipolarne i topologie słabe wymagają dodatkowych konstrukcji.
Omówimy je osobno, opierając się na udowodnionych
wersjach twierdzenia Hahna--Banacha. W części o przestrzeniach
dystrybucji ciągły dual $\mathcal D(\Omega)'$ otrzyma
kanoniczną interpretację jako przestrzeń dystrybucji
$\mathcal D'(\Omega)$.
\end{uwaga}

% Wersja źródłowa: audyt IV, jednostka 07
\setcounter{section}{6}
\section{Oddzielanie zbiorów wypukłych, polary i twierdzenie bipolarne}\label{sec:separacja-bipolarne}

\subsection{Przejście do rzeczywistych funkcjonałów}

W \hyperref[sec:hahn-banach]{rozdziale 6} udowodniliśmy twierdzenie Hahna--Banacha dla funkcjonałów liniowych.
Teraz wykorzystamy je do opisu geometrii zbiorów wypukłych i topologii,
którą można odtworzyć za pomocą przestrzeni sprzężonej. W szczególności odróżnimy
\emph{nierówności ścisłe dla każdego punktu} od \emph{oddzielenia z jednolitą luką}.

Przez cały rozdział $\K\in\{\R,\C\}$, a $X$ jest przestrzenią liniowo-topologiczną
nad $\K$. Gdy jest to wyraźnie zaznaczone, zakładamy lokalną wypukłość
oraz własność Hausdorffa. W topologii zespolonej oddzielanie zbiorów
odbywa się za pomocą \emph{rzeczywistych części} funkcjonałów z $X'$.

\begin{definicja}[Rzeczywiste funkcjonały i rodzaje separacji]\label{def:separacja}
Niech $A,B\subseteq X$ będą niepustymi zbiorami, a
$\ell:X\to\R$ niezerowym, ciągłym, rzeczywiście liniowym funkcjonałem
(na przestrzeni $X$ rozpatrywanej nad $\R$). Mówimy, że
\emph{oddziela} $A$ i $B$, gdy
$\sup_{a\in A}\ell(a)\le\inf_{b\in B}\ell(b)$.
Oddziela je \emph{ściśle punktowo}, gdy $\ell(a)<\ell(b)$
dla wszystkich $a\in A$, $b\in B$, a \emph{silnie}, gdy istnieje
$\delta>0$ takie, że $\sup_A\ell+\delta\le\inf_B\ell$.
Hiperpłaszczyzną afiniczną funkcjonału $\ell$ na poziomie $c\in\R$
nazywamy $H_{\ell,c}=\{x\in X:\ell(x)=c\}$.
\end{definicja}

\begin{lemat}[Rzeczywista część funkcjonału zespolonego]\label{lem:realification}
Niech $X$ będzie zespoloną TVS. Każdy ciągły, rzeczywiście liniowy
$\ell:X\to\R$ ma dokładnie jedno przedstawienie
$\ell=\operatorname{Re}F$, gdzie $F\in X'$. Mianowicie
$F(x)=\ell(x)-i\ell(ix)$.
\end{lemat}
\begin{proof}
Liniowość zespolona i jednoznaczność zostały wykazane w
lemacie o odtworzeniu funkcjonału zespolonego z \hyperref[sec:hahn-banach]{rozdziału 6}. Ponieważ $x\mapsto ix$
jest ciągły, także funkcjonał $F$ jest ciągły.
\end{proof}

\subsection{Oddzielanie zbioru otwartego od punktu}

\begin{lemat}[Funkcjonał Minkowskiego dla zbioru wypukłego]\label{lem:open-gauge}
Niech $E$ będzie rzeczywistą TVS, a $V\subseteq E$ otwartym
wypukłym otoczeniem zera (niekoniecznie symetrycznym).
Wówczas $V$ jest pochłaniający, a jego funkcjonał
$p_V(x)=\inf\{t>0:x\in tV\}$ jest skończony, podliniowy i ciągły.
Ponadto $V=\{x:p_V(x)<1\}$.
\end{lemat}
\begin{proof}
Pochłanianie wynika z ciągłości $t\mapsto tx$ w zerze.
Dla dowolnego $x\in E$ i $r>p_V(x)$ zachodzi $x\in rV$:
istnieje bowiem $0<s<r$ z $x\in sV$, a ponieważ $0\in V$
i $V$ jest wypukły, mamy $(s/r)V\subseteq V$.
Dla $t>0$ bezpośrednio z definicji otrzymujemy
$p_V(tx)=t p_V(x)$, a $p_V(0)=0$. Jeżeli
$r>p_V(x)$ i $s>p_V(y)$, to możemy napisać
$x=ru$, $y=sv$ z $u,v\in V$. Wypukłość daje
$(x+y)/(r+s)=(r/(r+s))u+(s/(r+s))v\in V$,
czyli $p_V(x+y)\le r+s$. Przechodząc z $r$ i $s$
do odpowiednich infimów, otrzymujemy
$p_V(x+y)\le p_V(x)+p_V(y)$. Zatem $p_V$ jest
podliniowy; w tym dowodzie nie korzystaliśmy z symetrii $V$.
Z wypukłości i $0\in V$ wynika
$tV\subseteq V$ dla $0<t\le1$. Jeżeli $x\in V$, to
z otwartości $V$ i ciągłości $s\mapsto sx$ w $s=1$
można wybrać $s>1$ takie, że $sx\in V$, czyli
$p_V(x)\le1/s<1$. Odwrotnie, jeżeli $p_V(x)<1$,
można znaleźć $t<1$ z $x\in tV\subseteq V$.
Stąd $V=\{p_V<1\}$.

Dla $\varepsilon>0$ zbiór $\{p_V<\varepsilon\}=\varepsilon V$
jest otwarty. Zatem $p_V$ jest górnie półciągły. Aby
wykazać ciągłość, dla każdego $x,y\in E$ korzystamy z
$\,p_V(x)\le p_V(y)+p_V(x-y)$ i
$p_V(y)\le p_V(x)+p_V(y-x)$.
Jeżeli $x-y\in\varepsilon V\cap(-\varepsilon V)$,
to obie różnice są mniejsze niż $\varepsilon$.
Takie zbiory tworzą otoczenia zera, co daje ciągłość.
\end{proof}

\begin{twierdzenie}[Oddzielanie punktu od otwartego zbioru wypukłego]
\label{thm:open-point-sep}
Niech $X$ będzie lokalnie wypukłą TVS nad $\K$,
$C\subseteq X$ niepustym otwartym zbiorem wypukłym
(w sensie rzeczywistym), a $x_0\in X\setminus C$.
Wówczas istnieje niezerowy ciągły, rzeczywiście liniowy
$\ell:X\to\R$ taki, że $\ell(c)<\ell(x_0)$
dla każdego $c\in C$. Jeżeli $\K=\C$, można przyjąć
$\ell=\operatorname{Re}F$ z $F\in X'$.
Nie twierdzimy ogólnie, że $\sup_C\ell<\ell(x_0)$.
\end{twierdzenie}
\begin{proof}
Traktujemy $X$ jako rzeczywistą TVS; lokalna wypukłość
jest zachowana. Zbiór $D=C-x_0$ jest niepusty, otwarty
i wypukły, a $0\notin D$. Wybierzmy $a\in D$ i połóżmy
$V=D-a$. Zbiór $V$ jest otwartym, wypukłym otoczeniem zera,
a $-a\notin V$, stąd $p_V(-a)\ge1$ na mocy
Lematu~\ref{lem:open-gauge}. Na prostej $\R(-a)$
określmy $g(t(-a))=t$. Dla $t\ge0$ mamy
$g(t(-a))=t\le p_V(t(-a))$. Dla $t<0$ nierówność
także zachodzi, bo $p_V\ge0$. Rzeczywiste twierdzenie
Hahna--Banacha daje liniowe $\ell:X\to\R$ spełniające
$\ell\le p_V$ i $\ell(-a)=1$. Z ciągłości $p_V$
i nierówności $-p_V(-x)\le\ell(x)\le p_V(x)$
wynika ciągłość $\ell$.

Dla dowolnego $d\in D$ mamy $d-a\in V$, więc
$\ell(d-a)\le p_V(d-a)<1$. Ponieważ $\ell(a)=-1$,
otrzymujemy $\ell(d)<0$, czyli
$\ell(c)<\ell(x_0)$ dla $c\in C$. Równość
$\ell(-a)=1$ dowodzi, że $\ell\ne0$. W przypadku
zespolonym stosujemy Lemat~\ref{lem:realification}.
\end{proof}

\begin{twierdzenie}[Oddzielanie dwóch zbiorów wypukłych]
\label{thm:open-two-sep}
Niech $X$ będzie lokalnie wypukłą TVS nad $\K$,
a $A,B\subseteq X$ niepustymi, rozłącznymi zbiorami wypukłymi.
Jeśli $A$ jest otwarty, istnieje niezerowy, ciągły, rzeczywiście
liniowy funkcjonał $\ell$ spełniający
$\ell(a)<\ell(b)$ dla wszystkich $a\in A$, $b\in B$.
W szczególności istnieje $c\in\R$ takie, że
$\sup_A\ell\le c\le\inf_B\ell$.
\end{twierdzenie}
\begin{proof}
Zbiór $D=A-B=\{a-b:a\in A,\ b\in B\}$ jest wypukły
i otwarty (jest sumą zbioru otwartego i dowolnego),
a $0\notin D$. Z Twierdzenia~\ref{thm:open-point-sep}
dostajemy $\ell(d)<0$ dla wszystkich $d\in D$,
co daje nierówność punktową. Dla wybranego $b_0\in B$
mamy $\sup_A\ell\le\ell(b_0)<\infty$; analogicznie,
dla $a_0\in A$ mamy $\inf_B\ell\ge\ell(a_0)>-\infty$.
Stąd obie liczby są skończone i możemy wybrać $c$
pomiędzy nimi.
\end{proof}

\begin{przyklad}[Dlaczego otwartość nie daje jednolitej luki?]
Dla $X=\R$, $A=(-1,0)$, $B=\{0\}$ funkcjonał
$\ell(x)=x$ spełnia $\ell(a)<\ell(b)$ dla wszystkich
$a\in A$, $b\in B$, lecz $\sup_A\ell=\inf_B\ell=0$.
Nie istnieje silna separacja tych zbiorów, ponieważ
$0\in\overline A\cap B$.
\end{przyklad}

\subsection{Oddzielanie silne i hiperpłaszczyzny podpierające}

\begin{lemat}[Wspólne otoczenie dla zbioru zwartego]
\label{lem:compact-neighborhood}
Niech $X$ będzie TVS, $K\subseteq X$ zbiorem zwartym,
a $F\subseteq X$ zbiorem domkniętym, przy czym $K\cap F=\varnothing$.
Wtedy istnieje otoczenie $U$ zera takie, że
$(K+U)\cap F=\varnothing$. Jeżeli $X$ jest lokalnie wypukła,
można przyjąć $U$ otwarte i wypukłe.
\end{lemat}
\begin{proof}
Dla każdego $k\in K$ istnieje otoczenie $V_k$ zera
z $(k+V_k)\cap F=\varnothing$. Wybieramy otoczenie
$W_k$ zera takie, że $W_k+W_k\subseteq V_k$,
a w przypadku lokalnie wypukłym możemy żądać
jego wypukłości i otwartości. Zwartość $K$ daje
pokrycie skończone $K\subseteq\bigcup_{j=1}^m(k_j+W_{k_j})$.
Wybierzmy otoczenie $U\subseteq\bigcap_j W_{k_j}$.
Jeżeli $x\in K$ należy do $k_j+W_{k_j}$, to
$x+U\subseteq k_j+V_{k_j}$, zatem $(K+U)\cap F=\varnothing$.
\end{proof}

\begin{twierdzenie}[Silna separacja zbioru zwartego i domkniętego]
\label{thm:strong-separation}
Niech $X$ będzie lokalnie wypukłą TVS nad $\K$,
$A,B\subseteq X$ niepustymi, rozłącznymi zbiorami wypukłymi.
Jeżeli $A$ jest zwarty, a $B$ domknięty, to istnieje
niezerowy ciągły, rzeczywiście liniowy funkcjonał $\ell$
oraz $\delta>0$ takie, że
\[
\sup_{a\in A}\ell(a)+\delta\le\inf_{b\in B}\ell(b).
\]
W przypadku zespolonym $\ell=\operatorname{Re}F$ dla pewnego $F\in X'$.
\end{twierdzenie}
\begin{proof}
Na mocy Lematu~\ref{lem:compact-neighborhood} istnieje
otwarte wypukłe otoczenie $U$ zera takie, że
$(A+U)\cap B=\varnothing$. Zbiór $A+U$ jest otwarty
i wypukły. Z Twierdzenia~\ref{thm:open-two-sep}
otrzymujemy niezerowy $\ell$ z
$\ell(a+u)<\ell(b)$ dla wszystkich
$a\in A$, $u\in U$, $b\in B$. Ponieważ $\ell\ne0$
i $U$ jest otoczeniem zera, istnieje $u_0\in U$
z $\ell(u_0)>0$. Kładąc $\delta=\ell(u_0)$, uzyskujemy
$\ell(a)+\delta<\ell(b)$ dla wszystkich $a,b$,
a zatem $\sup_A\ell+\delta\le\inf_B\ell$.
\end{proof}

\begin{wniosek}[Separacja punktu i domkniętego zbioru wypukłego]
\label{cor:point-closed-convex}
Niech $X$ będzie lokalnie wypukłą TVS, $C\subseteq X$
niepustym, domkniętym i wypukłym zbiorem oraz
$x_0\in X\setminus C$. Wówczas istnieją
ciągły rzeczywiście liniowy $\ell\ne0$ i $\delta>0$
takie, że $\ell(x_0)+\delta\le\inf_{c\in C}\ell(c)$.
\end{wniosek}
\begin{proof}
Stosujemy Twierdzenie~\ref{thm:strong-separation}
do $A=\{x_0\}$ i $B=C$.
\end{proof}

\begin{twierdzenie}[Hiperpłaszczyzna podpierająca]
\label{thm:supporting-hyperplane}
Niech $X$ będzie lokalnie wypukłą TVS nad $\K$,
a $C\subseteq X$ niepustym, domkniętym i wypukłym
zbiorem o niepustym wnętrzu. Dla każdego $x_0\in\partial C$
istnieje niezerowy, ciągły, rzeczywiście liniowy
$\ell:X\to\R$ taki, że $\ell(x)\le\ell(x_0)$
dla wszystkich $x\in C$.
\end{twierdzenie}
\begin{proof}
Wnętrze $G=\operatorname{int} C$ jest wypukłe.
Ponieważ $x_0\in\partial C$, $x_0\notin G$.
Twierdzenie~\ref{thm:open-point-sep} daje $\ell\ne0$
z $\ell(g)<\ell(x_0)$ dla $g\in G$.
Dla $x\in C$ i ustalonego $g_0\in G$ punkty
$x_t=(1-t)x+tg_0$ należą do $G$ dla $0<t\le1$:
małe otoczenie zera $V$ z $g_0+V\subseteq C$
spełnia $x_t+tV\subseteq C$ na mocy wypukłości.
Ponieważ $x_t\to x$ przy $t\downarrow0$, ciągłość
$\ell$ daje $\ell(x)\le\ell(x_0)$.
\end{proof}

\begin{uwaga}[Czego nie należy wnioskować]
Dla zbioru domkniętego wypukłego o pustym wnętrzu
funkcjonał podpierający w zadanym punkcie nie musi
istnieć bez dodatkowych hipotez. Podobnie z istnienia
hiperpłaszczyzny podpierającej nie wynika, że jest ona
jednoznaczna. W przestrzeni zespolonej pojęcie
hiperpłaszczyzny podpierającej odnosi się do równania
$\operatorname{Re}F(x)=c$, a nie do $F(x)=c$.
\end{uwaga}

\subsection{Zbiory absolutnie wypukłe i polary}

\begin{definicja}[Otoczka absolutnie wypukła]
\label{def:absolutely-convex-hull}
Niech $X$ będzie przestrzenią liniową nad $\K$, a $A\subseteq X$.
Definiujemy
\[
\operatorname{aco}(A)=\left\{\sum_{j=1}^n\lambda_j a_j:
 n\in\N,\ a_j\in A,\ \lambda_j\in\K,
 \sum_{j=1}^n|\lambda_j|\le1\right\}\cup\{0\}.
\]
Jest to najmniejszy zbiór wypukły i zrównoważony zawierający $A$.
Przyjmujemy $\operatorname{aco}(\varnothing)=\{0\}$.
\end{definicja}

\begin{definicja}[Polary w parze $(X,X')$]\label{def:polar}
Niech $X$ będzie TVS nad $\K$ i $X'$ przestrzenią sprzężoną.
Dla $A\subseteq X$ i $B\subseteq X'$ definiujemy
\[
 \begin{aligned}
 A^\circ&=\{f\in X':|f(a)|\le1\ \text{dla każdego }a\in A\},\\
 B^\circ&=\{x\in X:|f(x)|\le1\ \text{dla każdego }f\in B\}.
 \end{aligned}
\]
W szczególności $A^{\circ\circ}=(A^\circ)^\circ\subseteq X$.
Powyższa biegunowość jest \emph{absolutna}: w przypadku
zespolonym również używamy modułu, nie samej części rzeczywistej.
\end{definicja}

\begin{stwierdzenie}[Podstawowe własności polary]\label{prop:polar-properties}
Niech $X$ będzie TVS nad $\K$ i $A,A_1,A_2\subseteq X$.
Wówczas:
\begin{enumerate}
\item $A_1\subseteq A_2$ implikuje $A_2^\circ\subseteq A_1^\circ$;
\item $A^\circ=(\operatorname{aco}A)^\circ$ oraz
      $A\subseteq A^{\circ\circ}$;
\item $A^\circ$ jest wypukły, zrównoważony i domknięty
      w topologii punktowej zbieżności na $X'$, tj. $\sigma(X',X)$;
\item $A^{\circ\circ}$ jest wypukły, zrównoważony i domknięty
      w oryginalnej topologii $X$.
\end{enumerate}
\end{stwierdzenie}
\begin{proof}
Pierwsza własność wynika bezpośrednio z kwantyfikatora
w definicji. Dla $f\in A^\circ$ i
$x=\sum_j\lambda_j a_j$ z $\sum_j|\lambda_j|\le1$
mamy $|f(x)|\le\sum_j|\lambda_j||f(a_j)|\le1$,
co dowodzi równości polar. Każdy $a\in A$ spełnia
$|f(a)|\le1$ dla $f\in A^\circ$, zatem $a\in A^{\circ\circ}$.
Zbiory $\{f:|f(a)|\le1\}$ są domknięte, wypukłe i
zrównoważone w topologii punktowej; ich przecięcie
jest zbiorem $A^\circ$. Analogicznie przecięcie
$\{x:|f(x)|\le1\}$ po $f\in A^\circ$ daje własność (4).
\end{proof}

\begin{twierdzenie}[Twierdzenie bipolarne]\label{thm:bipolar}
Niech $X$ będzie lokalnie wypukłą TVS Hausdorffa nad
$\K\in\{\R,\C\}$, a $A\subseteq X$ dowolnym zbiorem.
Wówczas
\[
 A^{\circ\circ}=\overline{\operatorname{aco}(A)},
\]
gdzie domknięcie jest brane w oryginalnej topologii $X$.
\end{twierdzenie}
\begin{proof}
Z poprzedniego stwierdzenia wynika
$\overline{\operatorname{aco}(A)}\subseteq A^{\circ\circ}$.
Oznaczmy $C=\overline{\operatorname{aco}(A)}$.
Jest to zbiór niepusty, domknięty, wypukły i zrównoważony.
Jeśli $x_0\notin C$, Wniosek~\ref{cor:point-closed-convex}
zastosowany do $C$ daje ciągły rzeczywiście liniowy
$\ell$ z $\ell(x_0)>\sup_{c\in C}\ell(c)$
(po ewentualnej zmianie znaku).
Niech $F=\ell$ w przypadku rzeczywistym, a
$F(x)=\ell(x)-i\ell(ix)$ w zespolonym.
Wówczas $F\in X'$, $\operatorname{Re}F=\ell$.
Ze zrównoważenia $C$ wynika
\[
 m:=\sup_{c\in C}|F(c)|=\sup_{c\in C}\operatorname{Re}F(c)
 =\sup_{c\in C}\ell(c)<\ell(x_0)\le|F(x_0)|.
\]
W szczególności $m$ jest skończone i nieujemne.
Jeżeli $m>0$, to $G=F/m$ spełnia $|G(a)|\le1$
dla wszystkich $a\in A$, lecz $|G(x_0)|>1$.
Jeżeli $m=0$, wybieramy dodatnie $t$ z
$t|F(x_0)|>1$ i kładziemy $G=tF$.
W obu przypadkach $G\in A^\circ$, a $x_0\notin A^{\circ\circ}$.
Otrzymujemy więc odwrotną inkluzję.
\end{proof}

\begin{wniosek}[Domknięcie wypukłe nie zmienia się po osłabieniu topologii]
\label{cor:weak-convex-closure}
Niech $X$ będzie lokalnie wypukłą TVS Hausdorffa,
$C\subseteq X$ zbiorem wypukłym, a $\sigma(X,X')$
najsłabszą topologią, w której wszystkie $f\in X'$
są ciągłe. Wtedy domknięcie $C$ w $\sigma(X,X')$
jest równe domknięciu $C$ w oryginalnej topologii $X$.
\end{wniosek}
\begin{proof}
Jeśli $C=\varnothing$, teza jest oczywista. W dalszej części
zakładamy $C\ne\varnothing$. Topologia $\sigma(X,X')$ jest słabsza od wyjściowej,
zatem domknięcie słabe zawiera domknięcie wyjściowe.
Niech $D=\overline C$ oznacza domknięcie w topologii
wyjściowej. Zbiór $D$ jest niepusty, domknięty i wypukły.
Jeżeli $x_0\notin D$, stosujemy
Wniosek~\ref{cor:point-closed-convex} do $D$, a nie do
niekoniecznie domkniętego $C$. Otrzymujemy ciągły
rzeczywiście liniowy $\ell$ i $\alpha\in\R$ takie, że
$\ell(x_0)<\alpha<\inf_{d\in D}\ell(d)
\le\inf_{c\in C}\ell(c)$.
Jego reprezentacja przez część rzeczywistą funkcjonału
$F\in X'$ pokazuje, że zbiór
$\{x:\operatorname{Re}F(x)<\alpha\}$ jest słabo otwartym
otoczeniem $x_0$ rozłącznym z $C$.
\end{proof}

\subsection{Przykłady i pułapki}

\begin{przyklad}[Polary kul jednostkowych w $\R^2$]\label{ex:polar-balls}
W parze $X=X'=\R^2$ z parowaniem
$\langle a,x\rangle=a_1x_1+a_2x_2$ polara
$A=\{x:|x_1|+|x_2|\le1\}$ jest zbiorem
$A^\circ=\{a:\max(|a_1|,|a_2|)\le1\}$.
Jest to konsekwencja tożsamości
$\sup_{\|x\|_1\le1}|\langle a,x\rangle|=\|a\|_\infty$.
\begin{center}
\begin{tikzpicture}[scale=1.07,>=Stealth]
\begin{scope}[shift={(-2.65,0)}]
\draw[->,gray] (-1.55,0)--(1.55,0) node[right] {$x_1$};
\draw[->,gray] (0,-1.48)--(0,1.5) node[above] {$x_2$};
\fill[blue!12] (1,0)--(0,1)--(-1,0)--(0,-1)--cycle;
\draw[blue!65!black,very thick] (1,0)--(0,1)--(-1,0)--(0,-1)--cycle;
\node at (0,-1.85) {$A:\ \|x\|_1\le1$};
\end{scope}
\draw[->,thick] (-.72,0)--(.75,0) node[midway,above] {$\circ$};
\begin{scope}[shift={(2.65,0)}]
\draw[->,gray] (-1.55,0)--(1.55,0) node[right] {$a_1$};
\draw[->,gray] (0,-1.48)--(0,1.5) node[above] {$a_2$};
\fill[blue!12] (-1,-1) rectangle (1,1);
\draw[blue!65!black,very thick] (-1,-1) rectangle (1,1);
\node at (0,-1.85) {$A^\circ:\ \|a\|_\infty\le1$};
\end{scope}
\end{tikzpicture}
\end{center}
\end{przyklad}

\begin{przyklad}[Polara w produkcie $\K^\N$]\label{ex:polar-product}
Niech $X=\K^\N$ z topologią produktową, a
$A=\{e_n:n\in\N\}$, gdzie $e_n$ ma jedynkę na
$n$-tej współrzędnej. Z \hyperref[sec:hahn-banach]{rozdziału 6} wiemy, że $X'$
składa się z funkcjonałów $f_a(x)=\sum_n a_nx_n$
o skończonym nośniku. Zatem
$A^\circ=\{f_a:\ a\in c_{00},\ |a_n|\le1\ \text{dla każdego }n\}$,
a więc
\[
 A^{\circ\circ}=
 \left\{x\in\K^\N:\sum_{n=1}^\infty|x_n|\le1\right\}.
\]
Istotnie, supremum $|\sum_n a_nx_n|$ po
skończenie wspartych ciągach $(a_n)$ z $|a_n|\le1$
wynosi $\sup_{m\in\N}\sum_{n=1}^m|x_n|$. Z drugiej strony $\operatorname{aco}(A)$ składa się dokładnie
z ciągów o skończonym nośniku spełniających
$\sum_n|x_n|\le1$. Dla każdego $m$ funkcja
$x\mapsto\sum_{n=1}^m|x_n|$ jest ciągła w topologii
produktowej i nie przekracza $1$ na $\operatorname{aco}(A)$.
Dlatego nie przekracza też $1$ na jego domknięciu,
a po wzięciu supremum po $m$ dostajemy
$\sum_{n=1}^\infty|x_n|\le1$. Odwrotnie, jeśli
$x$ spełnia tę nierówność, jego obcięcia
$x^{(m)}=(x_1,\ldots,x_m,0,\ldots)$ należą do
$\operatorname{aco}(A)$ i zbiegają do $x$ w topologii
produktowej. Dowodzi to bezpośrednio opisanego
w twierdzeniu bipolarnym domknięcia.
\end{przyklad}

\begin{przyklad}[Dlaczego lokalna wypukłość ma znaczenie?]
Niech $X=L^p([0,1])$, $0<p<1$, z naturalną topologią
metryczną. Jak wiemy, $X'=\{0\}$. Dlatego dla
$A=\{0\}$ zachodzi $A^\circ=\{0\}$ oraz
$A^{\circ\circ}=X$, mimo że
$\overline{\operatorname{aco}(A)}=\{0\}$.
Twierdzenie bipolarne nie zachodzi więc w dowolnych TVS.
\end{przyklad}

\begin{uwaga}[Co dalej?]
W następnym rozdziale szczegółowo zbadamy pary
przestrzeni sprzężonych i topologie $\sigma(X,X')$ oraz
$\sigma(X',X)$. Dzięki temu pojęcia polary i
zamknięcia słabego otrzymają naturalną interpretację
jako narzędzia opisu przestrzeni funkcjonałów ciągłych,
a później także dystrybucji.
\end{uwaga}

% Wersja źródłowa: audyt II, jednostka 08
\setcounter{section}{7}
\section{Pary dualne i topologie słabe}\label{sec:topologie-slabe}

\subsection{Dlaczego potrzebujemy kolejnej topologii?}

Dotąd topologia przestrzeni liniowo-topologicznej $X$ była dana z góry,
a przestrzeń dualna topologiczna $X'$ składała się z tych
funkcjonałów liniowych, które są względem niej ciągłe.
Można postąpić odwrotnie: wybrać rodzinę funkcjonałów i skonstruować
\emph{najsłabszą} topologię, w której są one ciągłe.
Z tego pomysłu powstają topologie słabe, kluczowe w badaniu zbieżności
funkcji, operatorów i dystrybucji.

W rozdziale $\K\in\{\R,\C\}$. Wszystkie przestrzenie liniowe
są nad $\K$. Przez $E^*$ oznaczamy przestrzeń dualną algebraiczną,
a przez $E'$ przestrzeń dualną topologiczną, gdy $E$ jest TVS.

\begin{definicja}[Para dualna]\label{def:para-dualna}
Parą dualną nazywamy dwie przestrzenie liniowe $E,F$ wraz
z dwuliniowym parowaniem $\langle\cdot,\cdot\rangle:E\times F\to\K$
takim, że
\begin{enumerate}
\item dla każdego $x\in E\setminus\{0\}$ istnieje $f\in F$
      spełniający $\langle x,f\rangle\ne0$;
\item dla każdego $f\in F\setminus\{0\}$ istnieje $x\in E$
      spełniający $\langle x,f\rangle\ne0$.
\end{enumerate}
Mówimy, że parowanie jest \emph{niezdegenerowane} po obu stronach.
Nie zakładamy istnienia topologii na $E$ ani $F$.
\end{definicja}

\begin{przyklad}[Para kanoniczna]
Jeżeli $E$ jest lokalnie wypukłą TVS Hausdorffa, to
$(E,E')$, z parowaniem $\langle x,f\rangle=f(x)$,
jest parą dualną. Pierwsza niezdegenerowaność wynika
z konsekwencji twierdzenia Hahna--Banacha (\hyperref[sec:hahn-banach]{rozdział 6}),
a druga z faktu, że niezerowy funkcjonał ma niezerową wartość
na pewnym wektorze.
\end{przyklad}

\begin{przyklad}[Para ciągów]
Niech $E=c_0$, a $F=\ell^1$. Parowanie
$\langle x,a\rangle=\sum_{n=1}^{\infty}x_na_n$
jest poprawnie określone, gdyż
$\sum_n|x_na_n|\le\|x\|_\infty\|a\|_1$.
Jest niezdegenerowane: wektory współrzędnych $e_n$ należą
zarówno do $c_0$, jak i do $\ell^1$, i wykrywają odpowiednie
współrzędne.
\end{przyklad}

\subsection{Topologie słabe jako topologie początkowe}

\begin{definicja}[Topologia słaba pary dualnej]\label{def:slaba-para}
Niech $(E,F,\langle\cdot,\cdot\rangle)$ będzie parą dualną.
Topologią $\sigma(E,F)$ nazywamy najsłabszą topologię na $E$,
w której każdy funkcjonał $x\mapsto\langle x,f\rangle$,
$f\in F$, jest ciągły. Analogicznie definiujemy
$\sigma(F,E)$, wykorzystując funkcjonały
$f\mapsto\langle x,f\rangle$ dla $x\in E$.
\end{definicja}

\begin{twierdzenie}[Seminormy i baza otoczeń]\label{thm:slaba-baza}
Niech $(E,F)$ będzie parą dualną. Wtedy $\sigma(E,F)$
jest lokalnie wypukłą topologią Hausdorffa, generowaną przez
seminormy $p_f(x)=|\langle x,f\rangle|$, $f\in F$.
Bazę otoczeń zera stanowią zbiory
\[
 U(f_1,\ldots,f_m;\varepsilon)
 =\{x\in E:|\langle x,f_j\rangle|<\varepsilon,
 \ j=1,\ldots,m\},
\]
gdzie $m\ge1$, $f_1,\ldots,f_m\in F$ i $\varepsilon>0$.
W szczególności sieć $(x_i)_{i\in I}$ zbiega do $x\in E$
w $\sigma(E,F)$ wtedy i tylko wtedy, gdy
$\langle x_i,f\rangle\to\langle x,f\rangle$ dla każdego $f\in F$.
Analogiczne fakty zachodzą dla $\sigma(F,E)$.
\end{twierdzenie}
\begin{proof}
Z definicji topologii początkowej względem rodziny
odwzorowań do $\K$ otrzymujemy zbiory opisane skończenie
wieloma nierównościami. Każda funkcja $p_f$ jest seminormą;
przecięcia jej podpoziomic są absolutnie wypukłe, więc
otrzymana topologia jest lokalnie wypukła. Jeżeli $x\ne0$,
niezdegenerowaność parowania daje $f$ z $p_f(x)>0$.
Można zatem oddzielić $x$ od zera otoczeniem zera,
co dowodzi własności Hausdorffa. Kryterium zbieżności
jest własnością definiującą topologię początkową.
\end{proof}

\begin{stwierdzenie}[Słaba ograniczoność]\label{prop:slaba-ograniczonosc}
Niech $(E,F)$ będzie parą dualną, a $B\subseteq E$.
Zbiór $B$ jest ograniczony w TVS $(E,\sigma(E,F))$
wtedy i tylko wtedy, gdy dla każdego $f\in F$
zachodzi $\sup_{x\in B}|\langle x,f\rangle|<\infty$.
\end{stwierdzenie}
\begin{proof}
Jeżeli $B$ jest ograniczony, jego obraz przez każdy ciągły
funkcjonał jest ograniczonym zbiorem w $\K$.
Odwrotnie, niech $U=U(f_1,\ldots,f_m;\varepsilon)$
będzie podstawowym otoczeniem zera. Z założeń wynika, że
$M=\max_{1\le j\le m}\sup_{x\in B}|\langle x,f_j\rangle|$
jest skończone. Dla każdego $t>\max\{1,M/\varepsilon\}$
mamy $B\subseteq tU$, co jest definicją ograniczoności.
\end{proof}

\begin{twierdzenie}[Przestrzeń dualna topologii słabej]\label{thm:slaby-dual}
Niech $(E,F)$ będzie parą dualną. Odwzorowanie
$F\to(E,\sigma(E,F))'$, przypisujące $f$ funkcjonał
$x\mapsto\langle x,f\rangle$, jest izomorfizmem
algebraicznym. W szczególności topologia słaba nie wytwarza
nowych ciągłych funkcjonałów liniowych poza elementami $F$.
\end{twierdzenie}
\begin{proof}
Niezdegenerowaność zapewnia różnowartościowość, a
ciągłość funkcjonałów $f\in F$ wynika z definicji.
Niech teraz $g:E\to\K$ będzie liniowy i ciągły w
$\sigma(E,F)$. Istnieją $f_1,\ldots,f_m\in F$
i $\varepsilon>0$ takie, że
$|g(x)|<1$ dla $x\in U(f_1,\ldots,f_m;\varepsilon)$.
Połóżmy $T(x)=(\langle x,f_1\rangle,\ldots,
\langle x,f_m\rangle)\in\K^m$. Jeżeli $x\in\ker T$,
to $tx\in U$ dla każdego $t\in\K$; nierówność
$|t g(x)|<1$ dla wszystkich $t$ oznacza $g(x)=0$.
Zatem $g$ faktoryzuje się przez liniową mapę
$\widetilde g:T(E)\to\K$. Przedłużamy ją algebraicznie
z podprzestrzeni $T(E)$ na $\K^m$. Każdy funkcjonał
liniowy na $\K^m$ jest kombinacją współrzędnych,
więc $g(x)=\sum_{j=1}^m a_j\langle x,f_j\rangle
=\langle x,\sum_{j=1}^m a_j f_j\rangle$.
To kończy dowód surjektywności.
\end{proof}

\begin{wniosek}[Topologia wyjściowa i topologia słaba]\label{cor:slaba-wyjsciowa}
Jeżeli $E$ jest lokalnie wypukłą TVS Hausdorffa z topologią
$\tau$ i przestrzenią dualną topologiczną $E'$, to
$\sigma(E,E')\subseteq\tau$ (inkluzja topologii),
a przestrzenie dualne topologiczne $(E,\sigma(E,E'))'$ i
$(E,\tau)'$ są algebraicznie identyczne.
\end{wniosek}
\begin{proof}
Wszystkie elementy $E'$ są ciągłe względem $\tau$,
zatem minimalność topologii początkowej daje inkluzję.
Druga teza wynika z Twierdzenia~\ref{thm:slaby-dual}.
\end{proof}

\subsection{Słabe domknięcia i zjawiska nieskończenie wymiarowe}

\begin{twierdzenie}[Domknięcie podprzestrzeni przez anihilator]
\label{thm:anihilator-domkniecie}
Niech $E$ będzie lokalnie wypukłą TVS Hausdorffa,
$M\subseteq E$ podprzestrzenią liniową, a
$M^\perp=\{f\in E': f(m)=0\ \text{dla wszystkich }m\in M\}$.
Wtedy
\[
\overline M^{\,\tau}=\overline M^{\,\sigma(E,E')}
=\{x\in E:f(x)=0\ \text{dla wszystkich }f\in M^\perp\}.
\]
\end{twierdzenie}
\begin{proof}
Każdy funkcjonał z $M^\perp$ znika na obu domknięciach,
co daje jedną inkluzję. Jeżeli $x_0\notin\overline M^{\,\tau}$,
to w \hyperref[sec:hahn-banach]{rozdziale 6} oddzieliliśmy punkt od domkniętej
podprzestrzeni za pomocą $f\in E'$ spełniającego
$f|_M=0$ i $f(x_0)=1$. Tak więc $x_0$ nie leży w
zbiorze po prawej stronie. Ponieważ domknięcie słabe
zawiera domknięcie w $\tau$, otrzymujemy wszystkie równości.
\end{proof}

\begin{stwierdzenie}[Otoczenia słabe w nieskończonym wymiarze]
\label{prop:slabe-otoczenia}
Niech $(E,F)$ będzie parą dualną, przy czym
$\dim E=\infty$. Wtedy każde otoczenie zera w
$\sigma(E,F)$ zawiera niezerową podprzestrzeń liniową,
a więc nie jest ograniczone w tej topologii.
W konsekwencji $(E,\sigma(E,F))$ nie jest normowalna.
\end{stwierdzenie}
\begin{proof}
Dane otoczenie zawiera $U(f_1,\ldots,f_m;\varepsilon)$,
a ono zawiera $M=\bigcap_{j=1}^m\ker f_j$.
Podprzestrzeń $M$ ma kowymiar nie większy niż $m$,
więc jest niezerowa. Jeżeli $v\in M\setminus\{0\}$,
to $\K v\subseteq U$, lecz zbiór $\K v$ nie jest
słabo ograniczony: wybierając $f\in F$ z
$\langle v,f\rangle\ne0$, otrzymujemy nieograniczone
wartości $|\langle tv,f\rangle|$.
Przestrzeń normowalna ma ograniczone otoczenie zera,
co przeczy powyższemu faktowi.
\end{proof}

\begin{przyklad}[Słaba zbieżność bez zbieżności w normie]
Niech $E=\ell^2$ nad $\R$ albo $\C$ z topologią normową,
a $e_n$ oznacza $n$-ty wektor jednostkowy. Wtedy
$e_n\to0$ słabo, ponieważ każdy funkcjonał ciągły na
$\ell^2$ ma postać $x\mapsto\sum_k x_k a_k$
z pewnym $a\in\ell^2$ (w przypadku zespolonym
odpowiednio wybieramy współrzędne liniowe), a $a_n\to0$.
Jednocześnie $\|e_n\|_2=1$ dla każdego $n$.
\end{przyklad}

\begin{twierdzenie}[Słabe domknięcie sfery jednostkowej]
\label{thm:sfera-slaba}
Niech $E$ będzie nieskończenie wymiarową przestrzenią
unormowaną nad $\K$, $S=\{x\in E:\|x\|=1\}$,
$B=\{x\in E:\|x\|\le1\}$. Wtedy
$\overline S^{\,\sigma(E,E')}=B$.
\end{twierdzenie}
\begin{proof}
Kula $B$ jest słabo domknięta. Istotnie, jeżeli
$\|x_0\|>1$, twierdzenie Hahna--Banacha daje
$f\in E'$ z $\|f\|=1$ i $f(x_0)=\|x_0\|$,
a zbiór $\{x:\operatorname{Re}f(x)>1\}$ jest słabo
otwarty, zawiera $x_0$ i nie przecina $B$.
Stąd $\overline S^{\,\sigma(E,E')}\subseteq B$.

Niech $\|x\|<1$ i $U$ będzie podstawowym słabym
otoczeniem $x$, określonym przez skończenie wiele
funkcjonałów $f_1,\ldots,f_m$. Z nieskończonego wymiaru
wynika istnienie $v\ne0$ we wspólnym jądrze tych
funkcjonałów. Funkcja $t\mapsto\|x+tv\|$, $t\ge0$,
jest ciągła, w zerze mniejsza od $1$, a dla $t\to\infty$
dąży do nieskończoności. Istnieje więc $t_0>0$
z $\|x+t_0v\|=1$. Punkt $x+t_0v$ należy do
$S\cap U$, gdyż funkcjonały definiujące $U$ mają
na nim te same wartości co na $x$. To dowodzi
$B\subseteq\overline S^{\,\sigma(E,E')}$.
\end{proof}

\begin{uwaga}[Co stało się z intuicją geometryczną?]
W przestrzeni skończenie wymiarowej sfera jest domknięta.
W nieskończenie wymiarowej przestrzeni unormowanej topologia
słaba nie rozróżnia, za pomocą skończenie wielu funkcjonałów,
punktów sfery od punktów wewnątrz kuli. To jednak nie znaczy,
że kula jest słabo zwarta: zwartość omówimy osobno.
\end{uwaga}

\subsection{Topologia słaba-* na przestrzeni dualnej}

\begin{definicja}[Topologia słaba-*]\label{def:slaba-gwiazdka}
Niech $E$ będzie lokalnie wypukłą TVS Hausdorffa.
Na przestrzeni dualnej topologicznej $E'$ topologią
\emph{słabą-*} nazywamy $\sigma(E',E)$, tj. topologię
punktowej zbieżności funkcjonałów na $E$. Dla sieci
$(f_i)_{i\in I}\subseteq E'$ zapis
$f_i\overset{*}{\rightharpoonup}f$ oznacza, że
$f_i(x)\to f(x)$ dla każdego $x\in E$.
\end{definicja}

\begin{stwierdzenie}[Słaba i słaba-* topologia na przestrzeni dualnej]
\label{prop:slaba-weakstar}
Niech $E$ będzie przestrzenią Banacha. Na przestrzeni
Banacha $E'$ rozważmy topologię słabą
$\sigma(E',E'')$ i słabą-* $\sigma(E',E)$,
gdzie $E$ utożsamiamy z obrazem kanonicznego zanurzenia
$J:E\to E''$, $J(x)(f)=f(x)$.
Wtedy $\sigma(E',E)\subseteq\sigma(E',E'')$,
a równość zachodzi wtedy i tylko wtedy, gdy $E$ jest
refleksywna (tj. $J$ jest surjekcją).
\end{stwierdzenie}
\begin{proof}
Pierwsza inkluzja wynika z $J(E)\subseteq E''$.
Jeżeli $J(E)=E''$, rodziny definiujących funkcjonałów
są jednakowe, więc topologie są równe.
Odwrotnie, jeżeli topologie są równe, ich przestrzenie
dualne topologiczne są równe. Z Twierdzenia
\ref{thm:slaby-dual}, zastosowanego do par
$(E',J(E))$ i $(E',E'')$, wynika odpowiednio,
że te przestrzenie dualne są algebraicznie równe
$J(E)$ i $E''$. Zatem $J(E)=E''$.
\end{proof}

\begin{przyklad}[Zbieżność słaba-* nie musi być słaba]
Niech $E=c_0$, $E'=\ell^1$ z naturalnym parowaniem.
Ciąg wektorów jednostkowych $e_n\in\ell^1$ spełnia
$e_n\to0$ w $\sigma(\ell^1,c_0)$, ponieważ dla każdego
$x\in c_0$ mamy $\langle x,e_n\rangle=x_n\to0$.
Nie jest to jednak zbieżność słaba w $\ell^1$:
funkcjonał z $\ell^\infty=(\ell^1)'$ odpowiadający
ciągowi $(1,1,\ldots)$ ma na każdym $e_n$ wartość $1$.
\end{przyklad}

\begin{przyklad}[Topologia produktowa jako topologia słaba]
Niech $E=\K^{\N}$ z topologią produktową i $F=c_{00}$
z parowaniem $\langle x,a\rangle=\sum_n x_na_n$
(suma jest skończona). Ponieważ współrzędne
$x\mapsto x_n$ odpowiadają wektorom $e_n\in F$,
topologia produktowa zawiera $\sigma(E,F)$ i odwrotnie:
każdy funkcjonał z $F$ jest skończoną kombinacją
odwzorowań współrzędnych. Zatem topologia produktowa
jest dokładnie $\sigma(E,F)$.
\end{przyklad}

\subsection{Zupełność topologii słabej i przestrzeń dualna algebraiczna}

\begin{twierdzenie}[Uzupełnienie przestrzeni ze słabą topologią]
\label{thm:slabe-uzupelnienie}
Niech $(E,F)$ będzie parą dualną. Niech $F^*$ oznacza
przestrzeń dualną algebraiczną przestrzeni liniowej $F$,
wyposażoną w topologię punktowej zbieżności na $F$.
Wówczas odwzorowanie $J:E\to F^*$, określone wzorem
$J(x)(f)=\langle x,f\rangle$, jest liniowym zanurzeniem
topologicznym przestrzeni $(E,\sigma(E,F))$ w $F^*$.
Obraz $J(E)$ jest gęsty, a $F^*$ jest zupełną TVS
Hausdorffa. Tym samym $F^*$ jest uzupełnieniem
$(E,\sigma(E,F))$ w sensie \hyperref[sec:zupelnosc]{rozdziału 3}.
\end{twierdzenie}
\begin{proof}
Traktujemy $F^*$ jako podzbiór produktu $\K^F$.
Jest to zbiór tych funkcji $\Phi:F\to\K$, które spełniają
$\Phi(f+g)=\Phi(f)+\Phi(g)$ i
$\Phi(af)=a\Phi(f)$ dla $f,g\in F$, $a\in\K$.
Każdy taki warunek definiuje zbiór domknięty w topologii
produktowej, zatem $F^*$ jest domkniętą podprzestrzenią
$\K^F$. Produkt zupełnych przestrzeni liniowo-topologicznych
jest zupełny, a podprzestrzeń domknięta przestrzeni
zupełnej jest zupełna. Topologia jest Hausdorffa.

Niezdegenerowaność parowania daje różnowartościowość $J$,
a definicje topologii początkowych oznaczają, że jest to
zanurzenie topologiczne. Aby wykazać gęstość, weźmy
$\Phi\in F^*$ i podstawowe otoczenie punktu $\Phi$,
wyznaczone przez $f_1,\ldots,f_m\in F$ oraz $\varepsilon>0$.
Niech $G=\operatorname{lin}\{f_1,\ldots,f_m\}\subseteq F$.
Odwzorowanie $T:E\to G^*$, $T(x)(f)=\langle x,f\rangle$,
jest surjekcją. Istotnie, gdyby jego obraz był właściwą
podprzestrzenią skończenie wymiarowej przestrzeni $G^*$,
istniałby niezerowy element $g\in G$ z
$\langle x,g\rangle=0$ dla każdego $x\in E$,
co przeczy niezdegenerowaności parowania.
Wybieramy zatem $x\in E$ takie, że
$T(x)=\Phi|_G$. Wtedy $J(x)(f_j)=\Phi(f_j)$
dla wszystkich $j$, a więc $J(x)$ należy do rozważanego
otoczenia. To dowodzi gęstości.
\end{proof}

\begin{wniosek}[Zaskakująca niezupełność topologii słabej]
\label{cor:slabe-niezupełne}
Jeśli $E$ jest nieskończenie wymiarową przestrzenią
unormowaną, to przestrzeń $(E,\sigma(E,E'))$ nie jest
zupełna, nawet gdy $E$ jest przestrzenią Hilberta.
\end{wniosek}
\begin{proof}
Przestrzeń $E'$ jest nieskończenie wymiarową przestrzenią
Banacha. Istnieje na niej nieciągły funkcjonał liniowy
$\Phi\in(E')^*$ (przykład z \hyperref[sec:hahn-banach]{rozdziału 6}).
Każdy $J(x)$, $x\in E$, jest jednak ciągły względem normy
operatorowej na $E'$, ponieważ
$|J(x)(f)|=|f(x)|\le\|x\|\|f\|$.
Stąd $\Phi\notin J(E)$. Twierdzenie
\ref{thm:slabe-uzupelnienie} mówi, że $J(E)$ jest
gęstą, właściwą podprzestrzenią zupełnej TVS Hausdorffa
$F^*=(E')^*$. Zatem $J(E)$ nie jest zupełne.
\end{proof}

\begin{uwaga}
Powyższy wniosek nie przeczy zupełności $E$ w normie:
zupełność zależy od topologii i jej struktury jednostajnej.
Nie należy też mylić przestrzeni dualnej algebraicznej
$(E')^*$ z bidualem topologicznym $E''=(E')'$.
\end{uwaga}

\subsection{Topologia silna i zapowiedź topologii Mackeya}

\begin{definicja}[Topologia silna na przestrzeni dualnej]
\label{def:silna-dualna}
Niech $E$ będzie lokalnie wypukłą TVS Hausdorffa.
Dla każdego niepustego ograniczonego zbioru $B\subseteq E$ i $f\in E'$
kładziemy $p_B(f)=\sup_{x\in B}|f(x)|$.
Każdy $p_B$ jest skończony, ponieważ ciągły funkcjonał
przekształca zbiory ograniczone na ograniczone zbiory skalarów.
Topologię na $E'$ generowaną przez rodzinę takich seminorm
nazywamy topologią silną i oznaczamy $\beta(E',E)$.
\end{definicja}
\begin{stwierdzenie}[Przypadek przestrzeni unormowanych]
\label{prop:silna-unormowana}
Jeśli $E$ jest przestrzenią unormowaną, to topologia
$\beta(E',E)$ jest topologią normy operatorowej
$\|f\|=\sup_{\|x\|\le1}|f(x)|$ na $E'$.
\end{stwierdzenie}
\begin{proof}
Dla ograniczonego $B$ mamy
$\sup_{x\in B}\|x\|=M<\infty$, więc
$p_B(f)\le M\|f\|$. Z drugiej strony kula jednostkowa
$B_E$ jest ograniczona, a $p_{B_E}(f)=\|f\|$.
Obie rodziny seminorm definiują zatem tę samą topologię.
\end{proof}

\begin{definicja}[Topologia Mackeya]\label{def:mackey}
Niech $(E,F)$ będzie parą dualną. Dla każdego absolutnie
wypukłego, niepustego i zwartego w $\sigma(F,E)$ zbioru $K\subseteq F$
określamy $p_K(x)=\sup_{f\in K}|\langle x,f\rangle|$.
Topologię lokalnie wypukłą generowaną przez wszystkie takie
seminormy nazywamy \emph{topologią Mackeya} $\tau(E,F)$.
\end{definicja}

\begin{stwierdzenie}[Poprawność definicji Mackeya]\label{prop:mackey-basic}
Dla dowolnej pary dualnej $(E,F)$ funkcje $p_K$ są
skończonymi seminormami, a topologia $\tau(E,F)$
jest lokalnie wypukła, Hausdorffa i silniejsza od
$\sigma(E,F)$.
\end{stwierdzenie}
\begin{proof}
Dla ustalonego $x\in E$ funkcja
$f\mapsto|\langle x,f\rangle|$ jest ciągła względem
$\sigma(F,E)$, a zatem ograniczona na każdym zwartym $K$.
Nierówność trójkąta i jednorodność dają własności seminormy.
Dla każdego $f\in F$ zbiór
$K_f=\{af:|a|\le1\}$ jest absolutnie wypukły i zwarty
jako obraz zwartego domkniętego koła w $\K$.
Mamy $p_{K_f}(x)=|\langle x,f\rangle|$.
Topologia Mackeya zawiera więc topologię słabą,
a ponieważ ta ostatnia jest Hausdorffa, taką własność
ma również $\tau(E,F)$.
\end{proof}

\begin{uwaga}[Twierdzenie Mackeya--Arensa: zapowiedź]
Twierdzenie Mackeya--Arensa orzeka, że $\tau(E,F)$ jest
najsilniejszą topologią lokalnie wypukłą na $E$, dla której
przestrzeń dualna topologiczna jest algebraicznie równa $F$
(przy podanym parowaniu). Wymaga to jednak twierdzeń
o zwartości i zbiorach polarnych, które rozwiniemy w następnym
rozdziale. Nie traktujemy tej charakterystyki jako
udowodnionej w bieżącym miejscu.
\end{uwaga}

\begin{uwaga}[W stronę dystrybucji]
Dla przestrzeni funkcji testowych $\mathcal D(\Omega)$
topologia słaba-* $\sigma(\mathcal D'(\Omega),\mathcal D(\Omega))$
opisuje zbieżność dystrybucji na każdej ustalonej funkcji
testowej. Topologia silna przestrzeni dualnej opisuje
natomiast zbieżność jednostajną na wszystkich zbiorach
ograniczonych funkcji testowych. Ich porównanie będzie
istotne w dalszych rozdziałach.
\end{uwaga}

% Wersja źródłowa: audyt II, jednostka 09
\setcounter{section}{8}
\section{Banach--Alaoglu, Krein--Milman i topologia Mackeya}\label{sec:alaoglu-mackey}

\subsection{Zwartość słaba-* i równociągłość}

W przestrzeni nieskończenie wymiarowej kula jednostkowa nie jest zwarta
w topologii normowej. Zmiana topologii na słabą-* pozwala odzyskać zwartość,
ale zasadniczą rolę odgrywają wtedy \emph{zbiory polarne otoczeń zera}.
Najpierw dowiedziemy twierdzenia w klasie dowolnych przestrzeni
liniowo-topologicznych, bez zakładania lokalnej wypukłości i zupełności.

W całym rozdziale $\K\in\{\R,\C\}$. Dla TVS $E$ symbol $E'$ oznacza
przestrzeń sprzężoną, a $E^*$ przestrzeń algebraicznie sprzężoną.
Zbiorem polarnym $A\subseteq E$ nazywamy
$A^\circ=\{f\in E':|f(x)|\le1\ \text{dla każdego }x\in A\}$.
Na $E'$ rozpatrujemy topologię $\sigma(E',E)$ punktowej zbieżności na $E$.

\begin{definicja}[Równociągła rodzina funkcjonałów]\label{def:equicont-family}
Niech $E$ będzie TVS i $H\subseteq E'$. Rodzinę $H$ nazywamy
\emph{równociągłą}, jeżeli dla każdego $\varepsilon>0$ istnieje otoczenie
$U$ zera w $E$ takie, że $|f(u)|<\varepsilon$ dla wszystkich
$f\in H$ i $u\in U$.
\end{definicja}

\begin{stwierdzenie}[Równociągłość i zbiory polarne]\label{prop:equicont-polar}
Niech $E$ będzie TVS, a $H\subseteq E'$. Rodzina $H$ jest równociągła
wtedy i tylko wtedy, gdy istnieje otoczenie $U$ zera w $E$ takie,
że $H\subseteq U^\circ$. Każda rodzina równociągła jest punktowo
ograniczona, tzn. $\sup_{f\in H}|f(x)|<\infty$ dla każdego $x\in E$.
\end{stwierdzenie}
\begin{proof}
Jeżeli $H$ jest równociągła, bierzemy w definicji $\varepsilon=1$;
odpowiednie $U$ spełnia $H\subseteq U^\circ$. Odwrotnie,
jeżeli $H\subseteq U^\circ$, to dla ustalonego $\varepsilon>0$
otoczenie $(\varepsilon/2)U$ daje
$|f(v)|\le\varepsilon/2<\varepsilon$ dla $v\in(\varepsilon/2)U$
i $f\in H$. Dla $x\in E$ pochłanianie przez $U$ zapewnia
$x\in tU$ dla pewnego $t>0$, zatem $|f(x)|\le t$ dla $f\in H$.
\end{proof}

\begin{uwaga}[Wykorzystywane twierdzenie topologiczne]
W dowodzie twierdzenia Banacha--Alaoglu zastosujemy \emph{twierdzenie
Tychonowa}: iloczyn dowolnej rodziny zwartych przestrzeni topologicznych,
wyposażony w topologię produktową, jest zwarty. Jest to klasyczne
wynikanie z aksjomatu wyboru; nie będziemy tu ponownie dowodzić
tego ogólnotopologicznego rezultatu. Stosujemy go jedynie do iloczynów
zwartych przedziałów lub domkniętych kół w $\K$.
\end{uwaga}

\begin{twierdzenie}[Banach--Alaoglu--Bourbaki]\label{thm:banach-alaoglu-general}
Niech $E$ będzie przestrzenią liniowo-topologiczną nad $\K$, a
$U\subseteq E$ otoczeniem zera. Wówczas zbiór polarny $U^\circ$
jest zwarty w topologii $\sigma(E',E)$.
Nie zakładamy, że $E$ jest Hausdorffa, lokalnie wypukła, metryzowalna
lub zupełna.
\end{twierdzenie}
\begin{proof}
Umieszczamy $E^*$ w $\K^E$, utożsamiając funkcjonał z rodziną
jego wartości $(f(x))_{x\in E}$. Topologia punktowej zbieżności
jest topologią indukowaną przez topologię produktową.
Ponieważ $U$ jest pochłaniające, dla każdego $x\in E$ można wybrać
$t_x>0$ takie, że $x\in t_xU$. Każdy $f\in U^\circ$ spełnia wtedy
$|f(x)|\le t_x$. Niech $D_x=\{z\in\K:|z|\le t_x\}$ i
$P=\prod_{x\in E}D_x$. Zbiory $D_x$ są zwarte, więc $P$ jest zwarty
na mocy twierdzenia Tychonowa.

Zbiór $U^\circ$ jest domknięty w $P$. Rzeczywiście, warunki
$f(x+y)=f(x)+f(y)$, $f(ax)=af(x)$ oraz $|f(u)|\le1$ dla $u\in U$
opisują przecięcia zbiorów domkniętych w topologii produktowej.
Każda funkcja $f:E\to\K$ spełniająca te warunki jest liniowa,
a przy tym ciągła: dla $\varepsilon>0$ na otoczeniu
$(\varepsilon/2)U$ zachodzi $|f(v)|\le\varepsilon/2$.
Zatem otrzymany zbiór jest dokładnie $U^\circ\subseteq E'$.
Jako domknięty podzbiór zwartego $P$ jest on zwarty. Topologia
indukowana z $P$ na $U^\circ$ to $\sigma(E',E)$.
\end{proof}

\begin{wniosek}[Kula przestrzeni sprzężonej]\label{cor:alaoglu-norm}
Niech $E$ będzie przestrzenią unormowaną, niekoniecznie zupełną.
Dla każdego $R\ge0$ kula
$B_R'=\{f\in E':\|f\|\le R\}$ jest zwarta w $\sigma(E',E)$.
W szczególności normowa kula jednostkowa przestrzeni sprzężonej
każdej przestrzeni Banacha jest słabo-* zwarta.
\end{wniosek}
\begin{proof}
Dla $R>0$ niech $U=\{x\in E:\|x\|<1/R\}$.
Wtedy $U^\circ=B_R'$: warunek $|f(u)|\le1$ dla $\|u\|<1/R$
jest równoważny $\|f\|\le R$. Zastosowanie
Twierdzenia~\ref{thm:banach-alaoglu-general} kończy dowód.
Dla $R=0$ zbiór $B_0'=\{0\}$ jest zwarty.
\end{proof}

\begin{wniosek}[Domknięte podzbiory kuli]\label{cor:alaoglu-closed}
Niech $E$ będzie przestrzenią unormowaną, $R\ge0$, a
$A\subseteq B_R'$ będzie domknięty w topologii
$\sigma(E',E)$ indukowanej na $B_R'$. Wtedy $A$ jest
słabo-* zwarty.
\end{wniosek}
\begin{proof}
Domknięty podzbiór zwartej przestrzeni jest zwarty.
\end{proof}

\begin{uwaga}[Czego twierdzenie nie mówi?]
Twierdzenie Banacha--Alaoglu nie orzeka, że cała przestrzeń $E'$
jest słabo-* zwarta ani że kula jest zwarta w normie.
Zwartość nie daje też automatycznie możliwości zastąpienia
sieci przez ciągi: do tego potrzebna jest dodatkowa metryzowalność.
\end{uwaga}

\subsection{Metryzowalność i kryterium ośrodkowości}

\begin{twierdzenie}[Metryzowalność na rodzinie równociągłej]
\label{thm:weakstar-metrizable-equicont}
Niech $E$ będzie ośrodkową TVS Hausdorffa, a $H\subseteq E'$
rodziną równociągłą. Wówczas topologia indukowana na $H$
przez $\sigma(E',E)$ jest metryzowalna.
\end{twierdzenie}
\begin{proof}
Niech $D=\{x_n:n\in\N\}$ będzie zbiorem gęstym w $E$.
Określamy na $H$ funkcję
\[
 d(f,g)=\sum_{n=1}^\infty 2^{-n}\min\{1,|f(x_n)-g(x_n)|\}.
\]
Jest to metryka: nierówność trójkąta wynika z
$\min\{1,a+b\}\le\min\{1,a\}+\min\{1,b\}$,
a jeśli $d(f,g)=0$, to $f$ i $g$ pokrywają się na gęstym
$D$, zatem przez ciągłość na całym $E$.
Topologia metryki jest topologią punktowej zbieżności na $D$,
więc jest nie silniejsza od $\sigma(E',E)|_H$.

Pokażemy odwrotną inkluzję. Ustalmy $x\in E$, $f_0\in H$
i $\varepsilon>0$. Ponieważ $H$ jest równociągła,
istnieje zrównoważone otoczenie $V$ zera takie, że
$|h(v)|<\varepsilon/3$ dla wszystkich $h\in H$, $v\in V$.
Wybieramy $x_n\in D\cap(x+V)$. Dla $f\in H$ mamy
\[
 |(f-f_0)(x)|\le |f(x-x_n)|+|(f-f_0)(x_n)|
                      +|f_0(x_n-x)|
 <2\varepsilon/3+|(f-f_0)(x_n)|.
\]
Warunek $|(f-f_0)(x_n)|<\varepsilon/3$ jest otwarty
w metryce $d$. Zatem wartościowania w każdym $x\in E$
są ciągłe na $(H,d)$, co daje odwrotną inkluzję topologii.
\end{proof}

\begin{twierdzenie}[Kryterium dla kuli przestrzeni sprzężonej]
\label{thm:banach-ball-metrizable-iff-separable}
Niech $E$ będzie przestrzenią unormowaną nad $\K$,
a $B'=\{f\in E':\|f\|\le1\}$. Wówczas następujące
warunki są równoważne:
\begin{enumerate}
\item $E$ jest ośrodkowa w topologii normowej;
\item $B'$ z topologią $\sigma(E',E)$ jest metryzowalna.
\end{enumerate}
Jeśli warunki zachodzą, to $B'$ jest zwarta, a każdy jej ciąg ma podciąg zbieżny.
\end{twierdzenie}
\begin{proof}
Kula $B'$ jest równociągła, bo $|f(x)|\le\|x\|$ dla $f\in B'$.
Jeżeli $E$ jest ośrodkowa, stosujemy
Twierdzenie~\ref{thm:weakstar-metrizable-equicont};
zwartość wynika z Twierdzenia~\ref{thm:banach-alaoglu-general}.
Każda zwarta przestrzeń metryczna ma własność zwartości ciągowej.

Odwrotnie, załóżmy, że $B'$ jest metryzowalna. Istnieje
przeliczalna baza otoczeń zera w $B'$, a więc można wybrać
ciąg podstawowych otoczeń słabych-* zera
$V_n=\{f\in B':|f(x)|<\varepsilon_n\ \text{dla }x\in F_n\}$,
gdzie $F_n\subseteq E$ jest skończony, tak aby $\{V_n\}$
była bazą otoczeń zera w $B'$. Niech
$M=\operatorname{lin}(\bigcup_{n=1}^\infty F_n)$.
Jeśli $\overline M^{\|\cdot\|}\ne E$, twierdzenie Hahna--Banacha
daje niezerowy $f\in B'$ zerujący się na $M$.
Wtedy $f\in V_n$ dla każdego $n$. Ponieważ $B'$ jest Hausdorffa,
przecięcie wszystkich otoczeń zera w $B'$ jest równe $\{0\}$;
sprzeczność. Zatem $M$ jest gęsta w $E$, a jako przestrzeń
rozpięta nad $\K$ przez zbiór przeliczalny ma przeliczalny
gęsty podzbiór (kombinacje wymierne, a nad $\C$ -- z
wymiernymi częściami rzeczywistymi i urojonymi współczynników).
Stąd $E$ jest ośrodkowa.
\end{proof}

\begin{przyklad}[Kula przestrzeni dualnej $\ell^1$]
Dla $E=c_0$ mamy $E'=\ell^1$. Ponieważ $c_0$ jest ośrodkowa,
kula jednostkowa w $\ell^1$ jest zwarta i metryzowalna
w $\sigma(\ell^1,c_0)$. Nie jest natomiast zwarta w normie:
wektory $e_n$ należą do kuli, lecz $\|e_n-e_m\|_1=2$
dla $n\ne m$. Przy tym $e_n\to0$ słabo-*.
\end{przyklad}

\subsection{Dlaczego sieci są niezbędne?}

\begin{przyklad}[Zwartość bez metryzowalności]\label{ex:uncountable-product}
Niech $I$ będzie zbiorem nieprzeliczalnym,
$E=\ell^1(I)$, a $E'=\ell^\infty(I)$ z kanonicznym parowaniem
$\langle x,f\rangle=\sum_{i\in I}x_i f_i$.
Kula jednostkowa $B'$ przestrzeni $\ell^\infty(I)$
wyposażona w $\sigma(\ell^\infty(I),\ell^1(I))$
jest homeomorficzna z produktem
$\prod_{i\in I}\{z\in\K:|z|\le1\}$.
Istotnie, zbieżność słaba-* implikuje współrzędną przez
podstawienie $x=e_i$. Odwrotnie, jeśli sieć $f_\alpha\in B'$
zbiega współrzędnie do $f$, to dla $x\in\ell^1(I)$
wybieramy skończony $F\subseteq I$ z
$\sum_{i\notin F}|x_i|<\varepsilon/4$ i kontrolujemy
skończoną sumę $\sum_{i\in F}x_i(f_{\alpha,i}-f_i)$;
reszta jest mniejsza od $2\sum_{i\notin F}|x_i|<\varepsilon/2$.
Produkt jest zwarty, lecz nie jest metryzowalny, ponieważ
$\ell^1(I)$ nie jest ośrodkowa.

Wewnątrz $B'$ rozpatrzmy zbiór
$A=\{\chi_F:F\subseteq I\ \text{jest skończony}\}$,
gdzie $\chi_F$ jest funkcją charakterystyczną $F$.
Sieć $(\chi_F)$, indeksowana skończonymi podzbiorami $I$
i uporządkowana inkluzją, zbiega słabo-* do funkcji
$\chi_I$. Żaden \emph{ciąg} elementów $A$ nie zbiega
jednak słabo-* do $\chi_I$: suma jego przeliczalnie wielu
skończonych nośników jest przeliczalna, a więc istnieje
$i\in I$, na którym wszystkie wyrazy ciągu mają wartość zero.
\end{przyklad}

\begin{przyklad}[Słabo-* zwarta rodzina nierównociągła]
\label{ex:compact-not-equicont}
Niech $E=c_{00}$ będzie przestrzenią ciągów o skończonym nośniku,
z normą $\|x\|_\infty=\sup_n|x_n|$. Przestrzeń $E$ nie jest
zupełna. Dla $n\in\N$ określamy $f_n\in E'$ wzorem
$f_n(x)=n x_n$. Ponieważ każdy $x\in c_{00}$ ma skończony
nośnik, $f_n(x)=0$ dla wszystkich dostatecznie dużych $n$.
Wobec tego $f_n\to0$ w $\sigma(E',E)$, a zbiór
$K=\{0\}\cup\{f_n:n\in\N\}$ jest słabo-* zwarty:
każdy otwarty zbiór zawierający zero zawiera prawie wszystkie $f_n$.
Jednak $\|f_n\|=n$, zatem $K$ nie jest równociągły.
To pokazuje, że słaba-* zwartość nie oznacza równociągłości
bez dodatkowych założeń. W rozdziale o twierdzeniu
Banacha--Steinhausa omówimy warunki, pod którymi implikacja
staje się prawdziwa.
\end{przyklad}

\subsection{Punkty ekstremalne i twierdzenie Kreina--Milmana}
\label{subsec:krein-milman}

Twierdzenie Banacha--Alaoglu dostarcza zwartych zbiorów wypukłych
w topologiach słabych-*. Kolejne twierdzenie opisuje ich geometrię:
cały zbiór można odzyskać z jego punktów ekstremalnych. W tej
części, jeśli nie zaznaczono inaczej, $E$ oznacza przestrzeń
lokalnie wypukłą Hausdorffa nad $\K\in\{\R,\C\}$. W przypadku
zespolonym wypukłość rozumiemy w sensie rzeczywistym.

\begin{definicja}[Punkt ekstremalny i ściana]
\label{def:extreme-face9}
Niech $K\subseteq E$ będzie niepustym zbiorem wypukłym.
Punkt $e\in K$ jest \emph{ekstremalny}, jeżeli z
$e=tx+(1-t)y$, $x,y\in K$, $0<t<1$ wynika $x=y=e$.
Zbiór wszystkich takich punktów oznaczamy $\operatorname{Ext}(K)$.
Niepusty podzbiór $F\subseteq K$ jest \emph{ścianą} zbioru $K$,
jeżeli jest wypukły i dla $x,y\in K$ oraz $0<t<1$ z
$tx+(1-t)y\in F$ wynika $x,y\in F$.
Przez $\operatorname{conv}(A)$ oznaczamy otoczkę wypukłą $A$,
a przez $\overline A$ jej domknięcie w $E$.
\end{definicja}

\begin{przyklad}[Wielokąt i kula w przestrzeni Hilberta]
Wierzchołki wypukłego wielokąta o niepustym wnętrzu w $\R^2$ są dokładnie jego
punktami ekstremalnymi. Jeżeli $H$ jest rzeczywistą przestrzenią
Hilberta, to punktami ekstremalnymi domkniętej kuli jednostkowej
$B_H$ są dokładnie punkty sfery $\{x:\|x\|=1\}$.
Istotnie, dla $\|x\|<1$ i dostatecznie małego $h\ne0$
mamy $x=\tfrac12(x+h)+\tfrac12(x-h)$, a oba punkty należą do kuli.
Natomiast jeśli $\|x\|=1$ i $x=(u+v)/2$ przy $u,v\in B_H$,
to tożsamość równoległoboku daje
$\|u-v\|^2=2\|u\|^2+2\|v\|^2-4\|x\|^2\le0$, więc $u=v=x$.
\end{przyklad}

\begin{center}
\begin{tikzpicture}[scale=0.90, every node/.style={font=\small}]
\filldraw[fill=blue!7,draw=blue!55!black,very thick]
(0.1,0.6)--(1.5,0.15)--(3.9,0.65)--(4.35,2.05)--(2.75,2.85)--(0.4,2.3)--cycle;
\foreach \x/\y in {0.1/0.6,1.5/0.15,3.9/0.65,4.35/2.05,2.75/2.85,0.4/2.3}
 {\fill[blue!65!black] (\x,\y) circle (2.7pt);}
\draw[orange!85!black,densely dashed,thick] (-0.1,2.85)--(4.75,2.85);
\fill[orange!80!black] (2.75,2.85) circle (3.5pt);
\node[anchor=west,align=left] at (4.83,2.9) {$\ell=\max_K\ell$};
\node at (2.2,1.5) {$K$};
\node[anchor=north] at (2.2,-0.05) {Wierzchołki są punktami ekstremalnymi.};
\end{tikzpicture}
\end{center}
\noindent
\emph{Ilustracja.} Linia przerywana jest poziomicą funkcjonału
liniowego podpierającą $K$; odpowiednia ściana składa się tutaj
z jednego punktu ekstremalnego.

\begin{lemat}[Ściana poziomicy maksymalnej]
\label{lem:max-face9}
Niech $K\subseteq E$ będzie niepusty, zwarty i wypukły,
a $\ell:E\to\R$ niech będzie ciągłym rzeczywistym
funkcjonałem liniowym. Wówczas zbiór
$F_\ell=\{x\in K:\ell(x)=\max_{y\in K}\ell(y)\}$
jest niepustą zwartą ścianą $K$.
Jeżeli $F$ jest ścianą $K$, a $G$ ścianą $F$,
to $G$ jest również ścianą $K$.
\end{lemat}
\begin{proof}
Ciągłość $\ell$ i zwartość $K$ zapewniają istnienie maksimum;
$F_\ell$ jest domkniętym, a więc zwartym podzbiorem $K$.
Jest wypukły z liniowości $\ell$. Jeśli
$tx+(1-t)y\in F_\ell$, gdzie $x,y\in K$ i $0<t<1$, to
$\max_K\ell=t\ell(x)+(1-t)\ell(y)$, podczas gdy
$\ell(x),\ell(y)\le\max_K\ell$. Stąd obie wartości
są maksymalne i $x,y\in F_\ell$. Druga część wynika
z dwukrotnego zastosowania definicji ściany.
\end{proof}

\begin{lemat}[Istnienie punktu ekstremalnego]
\label{lem:exist-extreme9}
Każdy niepusty, zwarty i wypukły podzbiór przestrzeni lokalnie
wypukłej Hausdorffa ma co najmniej jeden punkt ekstremalny.
Co więcej, każda jego niepusta, domknięta ściana zawiera
punkt ekstremalny całego zbioru.
\end{lemat}
\begin{proof}
Niech $K\subseteq E$ spełnia założenia. Rozważmy rodzinę
$\mathcal F$ wszystkich niepustych \emph{domkniętych} ścian $K$,
uporządkowaną przez inkluzję odwrotną.
Każdy łańcuch w $\mathcal F$ ma górne ograniczenie:
jego elementy są zagnieżdżonymi niepustymi zbiorami
domkniętymi w zwartym $K$, więc mają niepuste przecięcie.
Przecięcie jest domkniętą ścianą, ponieważ warunek definiujący
ścianę jest zachowywany przy przecięciach.
Lemat Kuratowskiego--Zorna zapewnia istnienie minimalnej
niepustej domkniętej ściany $F$.

Gdyby $F$ zawierała różne punkty $x,y$, lokalna wypukłość
i warunek Hausdorffa zapewniałyby ciągły rzeczywisty
funkcjonał liniowy $\ell$ rozróżniający $x$ i $y$.
Dla $\K=\C$ możemy wziąć część rzeczywistą ciągłego
zespolonego funkcjonału. Z Lematu~\ref{lem:max-face9}
wynika, że zbiór punktów, w których $\ell$ osiąga
maksimum na $F$, jest niepustą, domkniętą ścianą $K$,
prawdziwie zawartą w $F$. To przeczy minimalności.
Zatem $F=\{e\}$, a warunek bycia ścianą oznacza,
że $e\in\operatorname{Ext}(K)$.
Jeśli zaczniemy od dowolnej niepustej domkniętej ściany
$F_0$ zamiast od $K$, ten sam argument da jej punkt
ekstremalny, ekstremalny również w $K$.
\end{proof}

\begin{twierdzenie}[Krein--Milman]
\label{thm:krein-milman}
Niech $E$ będzie rzeczywistą lub zespoloną przestrzenią
lokalnie wypukłą Hausdorffa, a $K\subseteq E$ niepustym,
zwartym i wypukłym zbiorem. Wtedy
\[
 K=\overline{\operatorname{conv}(\operatorname{Ext}(K))}.
\]
Domknięcie rozumiemy w wyjściowej topologii $E$.
\end{twierdzenie}
\begin{proof}
Z Lematu~\ref{lem:exist-extreme9} zbiór
$\operatorname{Ext}(K)$ jest niepusty. Ponieważ $K$ jest
wypukły i domknięty, zbiór
$C=\overline{\operatorname{conv}(\operatorname{Ext}(K))}$
jest zawarty w $K$ i jest zwarty jako domknięty podzbiór $K$.

Załóżmy, że istnieje $x_0\in K\setminus C$.
Twierdzenie o ścisłym oddzielaniu punktu od domkniętego
zbioru wypukłego w przestrzeni lokalnie wypukłej
(twierdzenie Hahna--Banacha o oddzielaniu zbiorów wypukłych) daje ciągły
rzeczywisty funkcjonał liniowy $\ell:E\to\R$ spełniający
$\ell(x_0)>\sup_{c\in C}\ell(c)$. W przypadku zespolonym
$\ell$ można otrzymać jako część rzeczywistą funkcjonału
zespolonego.

Niech $M=\max_{x\in K}\ell(x)$.
Z Lematu~\ref{lem:max-face9} zbiór
$F_\ell=\{x\in K:\ell(x)=M\}$ jest niepustą zwartą
ścianą $K$. Z Lematu~\ref{lem:exist-extreme9}
ma on punkt $e\in\operatorname{Ext}(K)$.
Wtedy $e\in C$, a równocześnie
$\ell(e)=M\ge\ell(x_0)>\sup_C\ell$, co jest sprzecznością.
Stąd $C=K$.
\end{proof}

\begin{uwaga}[Dlaczego hipotezy są istotne?]
Zwartość nie może być pominięta: cała przestrzeń $E=\R$
ma niepusty, domknięty i wypukły zbiór $K=\R$,
ale nie ma punktów ekstremalnych. Lokalna wypukłość
gwarantuje istnienie dostatecznie wielu ciągłych
funkcjonałów do oddzielania punktów. W ogólnych TVS
bez lokalnej wypukłości powyższy dowód nie działa.
Twierdzenie nie mówi także, że każdy punkt $K$ jest
\emph{skończoną} kombinacją wypukłą punktów ekstremalnych;
w ogólności konieczne jest domknięcie.
\end{uwaga}

\begin{wniosek}[Zasada maksimum Bauera]
\label{cor:bauer9}
Niech $K$ będzie niepustym zwartym wypukłym podzbiorem
przestrzeni lokalnie wypukłej Hausdorffa $E$ nad $\K$.
Jeżeli $f:K\to\R$ jest wypukła i półciągła z góry,
to istnieje $e\in\operatorname{Ext}(K)$ takie,
że $f(e)=\max_{x\in K}f(x)$.
\end{wniosek}
\begin{proof}
Niech $M=\max_K f$; maksimum istnieje dzięki zwartości $K$
i półciągłości $f$ z góry. Zbiór $A=\{x\in K:f(x)=M\}$
jest niepusty i domknięty. Nie musi być wypukły, ale ma
\emph{własność ekstremalności}: jeśli $tx+(1-t)y\in A$
przy $x,y\in K$ i $0<t<1$, to $x,y\in A$. Rzeczywiście,
$M=f(tx+(1-t)y)\le tf(x)+(1-t)f(y)\le M$,
więc obie nierówności są równościami.

Wybierzmy, korzystając z lematu Kuratowskiego--Zorna,
minimalny niepusty, domknięty podzbiór $B\subseteq A$
mający własność ekstremalności względem $K$.
Argument z przecięciem łańcucha jest taki sam jak
w Lemacie~\ref{lem:exist-extreme9}; nie używa wypukłości $B$.
Gdyby $B$ zawierał różne punkty, istniałby ciągły rzeczywisty
funkcjonał $\ell$ niebędący na nim stałym. Niepusty zbiór
$B_\ell=\{x\in B:\ell(x)=\max_{y\in B}\ell(y)\}$
byłby zwartym, właściwym podzbiorem $B$.
Jeśli $z=tx+(1-t)y\in B_\ell$ dla $x,y\in K$, to
$x,y\in B$, gdyż $B$ jest ekstremalny, a z liniowości
$\ell$ wynika $x,y\in B_\ell$. Jest to sprzeczność
z minimalnością. Zatem $B=\{e\}$; własność ekstremalności
tej jednoelementowej rodziny mówi dokładnie,
że $e\in\operatorname{Ext}(K)$, a $e\in A$ kończy dowód.
\end{proof}

\begin{przyklad}[Produkt przedziałów]
\label{ex:krein-cube9}
Niech $I$ będzie dowolnym niepustym zbiorem,
a $K=[-1,1]^I$ z topologią produktową, uważany za
podzbiór $\R^I$. Jest to zwarty zbiór wypukły na mocy
twierdzenia Tychonowa. Jego punkty ekstremalne to dokładnie
$x\in K$ takie, że $x_i\in\{-1,1\}$ dla każdego $i\in I$.
Istotnie, jeżeli $|x_i|<1$, można zmienić tylko współrzędną
$i$ o małe wartości $\pm\varepsilon$, aby przedstawić
$x$ jako środek dwóch różnych punktów $K$.
Jeśli wszystkie współrzędne są równe $\pm1$, to każdy
rozkład $x=ty+(1-t)z$ w $K$ musi być równy $y=z=x$
w każdej współrzędnej.
Dla dowolnego nieprzeliczalnego $I$ argument zwartości
opiera się na twierdzeniu Tychonowa, a nie na kryterium
sprowadzającym zwartość wyłącznie do zachowania ciągów.
\end{przyklad}

\begin{przyklad}[Słabo-* zwarta kula $\ell^1$]
\label{ex:krein-l1-weakstar9}
Na przestrzeni $\ell^1(\N;\R)=c_0(\N;\R)'$ rozważmy
słabą-* topologię $\sigma(\ell^1,c_0)$.
Kula $B=\{x\in\ell^1:\|x\|_1\le1\}$ jest zwarta
w tej topologii na mocy twierdzenia Banacha--Alaoglu.
Jej punktami ekstremalnymi są dokładnie $\pm e_n$.
Dla $\|x\|_1<1$ można rozdzielić $x$ jako średnią
$x\pm\varepsilon e_1$. Jeśli $\|x\|_1=1$ i $x$ ma
co najmniej dwie niezerowe współrzędne, odpowiednio małe
przeniesienie masy między nimi pozwala zapisać
$x=(y+z)/2$ dla różnych $y,z\in B$.
Z kolei $\pm e_n$ są ekstremalne, gdyż równość
$\pm1=(y_n+z_n)/2$ dla $y,z\in B$ wymusza
$y_n=z_n=\pm1$ i zerowanie się innych współrzędnych.
Wobec twierdzenia Kreina--Milmana kula $B$ jest słabo-*
domkniętą otoczką wypukłą tych punktów.
\end{przyklad}


\subsection{Twierdzenie Mackeya--Arensa}

Przypomnijmy, że dla pary dualnej $(E,F)$ topologia
$\tau(E,F)$ jest generowana przez seminormy
$p_K(x)=\sup_{f\in K}|\langle x,f\rangle|$, gdzie $K$
przebiega niepuste, absolutnie wypukłe i zwarte zbiory
w topologii $\sigma(F,E)$.

\begin{definicja}[Topologia zgodna z parą dualną]\label{def:compatible-duality}
Niech $(E,F,\langle\cdot,\cdot\rangle)$ będzie parą dualną.
Lokalnie wypukłą topologię Hausdorffa $\rho$ na $E$ nazywamy
\emph{zgodną z parą dualną $(E,F)$}, jeżeli odwzorowanie
$F\to(E,\rho)'$, $f\mapsto(x\mapsto\langle x,f\rangle)$,
jest izomorfizmem algebraicznym.
\end{definicja}

\begin{lemat}[Zwartość i dominacja funkcjonału]\label{lem:compact-domination}
Niech $(E,F)$ będzie parą dualną, a $K\subseteq F$ niepustym,
absolutnie wypukłym i zwartym zbiorem w $\sigma(F,E)$.
Jeżeli $g\in E^*$ spełnia
$|g(x)|\le p_K(x)$ dla każdego $x\in E$,
to istnieje $f\in K$ takie, że
$g(x)=\langle x,f\rangle$ dla wszystkich $x\in E$.
\end{lemat}
\begin{proof}
Dla każdego skończonego zbioru $S=\{x_1,\ldots,x_m\}\subseteq E$
rozważamy odwzorowanie
$T_S:K\to\K^m$, $T_S(f)=(\langle x_1,f\rangle,\ldots,
\langle x_m,f\rangle)$. Jego obraz jest zwarty i absolutnie wypukły.
Pokażemy, że $y=(g(x_1),\ldots,g(x_m))\in T_S(K)$.
Gdyby tak nie było, twierdzenie o silnym oddzielaniu punktu
od zwartego wypukłego zbioru w skończonym wymiarze dałoby
$a_1,\ldots,a_m\in\K$ takie, że
\[
 \operatorname{Re}\sum_{j=1}^m a_jg(x_j)
 >\sup_{f\in K}\operatorname{Re}
       \sum_{j=1}^m a_j\langle x_j,f\rangle.
\]
Połóżmy $x=\sum_{j=1}^m a_jx_j$.
Absolutna wypukłość $K$ oznacza, że zbiór wartości
$\{\langle x,f\rangle:f\in K\}$ jest zrównoważony,
a więc jego supremum części rzeczywistych jest równe
supremum modułów. Otrzymalibyśmy
$\operatorname{Re}g(x)>p_K(x)$, sprzecznie z założeniem.

Zatem dla każdego skończonego $S$ zbiór
$K_S=\{f\in K:\langle x,f\rangle=g(x)\ \text{dla }x\in S\}$
jest niepusty i domknięty w $K$. Rodzina $(K_S)$ ma
własność skończonych przecięć, ponieważ
$K_{S_1}\cap\cdots\cap K_{S_r}=K_{S_1\cup\cdots\cup S_r}$.
Zwartość $K$ daje punkt we wspólnym przecięciu,
co kończy dowód.
\end{proof}

\begin{twierdzenie}[Mackey--Arens]\label{thm:mackey-arens}
Niech $(E,F,\langle\cdot,\cdot\rangle)$ będzie parą dualną
przestrzeni liniowych nad $\K$. Wówczas topologia
$\tau(E,F)$ jest \emph{najsilniejszą} lokalnie wypukłą
topologią Hausdorffa na $E$ zgodną z tą parą dualną.
Dokładniej:
\begin{enumerate}
\item przestrzeń sprzężona $(E,\tau(E,F))'$ jest
      algebraicznie równa $F$ przez dane parowanie;
\item jeżeli $\rho$ jest dowolną inną lokalnie wypukłą topologią
      Hausdorffa na $E$ zgodną z $(E,F)$, to $\rho\subseteq\tau(E,F)$.
\end{enumerate}
\end{twierdzenie}
\begin{proof}
Z definicji $\tau(E,F)$ każdy element $F$ wyznacza
funkcjonał ciągły na $E$; istotnie, zbiór
$\{af:|a|\le1\}$ jest absolutnie wypukły i zwarty
w $\sigma(F,E)$.

Niech $g\in(E,\tau(E,F))'$. Ciągłość w zerze daje
niepuste absolutnie wypukłe zbiory zwarte
$K_1,\ldots,K_m\subseteq F$ i stałą $C>0$ takie, że
$|g(x)|\le C\max_j p_{K_j}(x)$ dla wszystkich $x\in E$.
Aby uzyskać to oszacowanie, bierzemy bazowe otoczenie
$\{x:\max_j p_{K_j}(x)<\varepsilon\}$ zawarte w
$\{x:|g(x)|<1\}$ i skalujemy $x$; można przyjąć
$C=2/\varepsilon$. Niech $K=K_1+\cdots+K_m$.
Jest to niepusty, absolutnie wypukły zbiór zwarty
w $\sigma(F,E)$ jako obraz zwartego produktu przez dodawanie.
Ponieważ $0\in K_j$, mamy $p_K\ge\max_j p_{K_j}$.
Zatem $|g(x)|\le C p_K(x)=p_{CK}(x)$.
Lemat~\ref{lem:compact-domination}, zastosowany do zwartego
absolutnie wypukłego $CK$, zapewnia $g\in F$.
To dowodzi punktu (1).

Niech teraz $\rho$ będzie zgodną lokalnie wypukłą topologią.
Wystarczy sprawdzić, że każda $\rho$-ciągła seminorma
$p:E\to[0,\infty)$ jest $\tau(E,F)$-ciągła.
Określmy $U=\{x:p(x)\le1\}$ i
\[
 K_p=\{f\in F:|\langle x,f\rangle|\le p(x)
                 \ \text{dla wszystkich }x\in E\}.
\]
Zbiór $U$ jest otoczeniem zera dla $\rho$, a
$K_p=U^\circ$ przy danym utożsamieniu $F=(E,\rho)'$.
Istotnie, jeśli $|\langle x,f\rangle|\le1$ dla $p(x)\le1$,
to dla $p(x)>0$ stosujemy $x/p(x)$, a przy $p(x)=0$
skalujemy $nx$ i przechodzimy z $n\to\infty$.
Na mocy Twierdzenia~\ref{thm:banach-alaoglu-general}
zbiór $K_p$ jest zwarty w $\sigma(F,E)$;
jest również niepusty i absolutnie wypukły.

Pokażemy, że $p=p_{K_p}$. Nierówność $p_{K_p}\le p$
wynika z definicji. Dla ustalonego $x\ne0$ funkcjonał
na $\K x$, $\lambda x\mapsto\lambda p(x)$,
jest dominowany przez seminormę $p$ co do modułu.
Twierdzenie Hahna--Banacha (w wersji dla seminormy)
przedłuża go do funkcjonału $f\in E^*$ z
$|f(y)|\le p(y)$ dla każdego $y\in E$.
Funkcjonał ten jest $\rho$-ciągły, należy więc do $F$,
a w konsekwencji $f\in K_p$ i $|f(x)|=p(x)$.
Zachodzi zatem $p=p_{K_p}$, czyli $p$ jest
$\tau(E,F)$-ciągła. Ponieważ topologia $\rho$ jest
generowana przez swoje ciągłe seminormy, otrzymujemy
$\rho\subseteq\tau(E,F)$.
\end{proof}

\begin{wniosek}[Najsłabsza i najsilniejsza topologia zgodna]\label{cor:mackey-extremes}
Niech $(E,F)$ będzie parą dualną. Topologia słaba
$\sigma(E,F)$ jest najsłabszą, zaś topologia Mackeya
$\tau(E,F)$ najsilniejszą lokalnie wypukłą topologią
Hausdorffa zgodną z tą parą. W szczególności dla każdej
takiej topologii $\rho$ zachodzi
$\sigma(E,F)\subseteq\rho\subseteq\tau(E,F)$.
\end{wniosek}
\begin{proof}
Jeżeli $\rho$ jest zgodna, wszystkie funkcjonały z $F$
są $\rho$-ciągłe, więc minimalność topologii początkowej
daje $\sigma(E,F)\subseteq\rho$. Druga inkluzja wynika
z Twierdzenia~\ref{thm:mackey-arens}.
\end{proof}

\begin{przyklad}[Przestrzeń słaba i topologia Mackeya nie muszą być równe]
Niech $E$ będzie nieskończenie wymiarową przestrzenią Banacha,
a $F=E'$. Topologia $\sigma(E,E')$ nie jest normowalna
(co udowodniono w \hyperref[sec:topologie-slabe]{rozdziale 8}), a topologia normowa
jest zgodna z parą $(E,E')$. Z wniosku
\ref{cor:mackey-extremes} wynika więc, że
$\sigma(E,E')\subsetneq\tau(E,E')$.
Nie utożsamiamy tu jednak automatycznie $\tau(E,E')$
z topologią normową: odpowiednie twierdzenie dla przestrzeni
Banacha powiążemy w dalszej części skryptu z zasadą
jednostajnej ograniczoności.
\end{przyklad}

\begin{uwaga}[Dalsza perspektywa: przestrzenie dystrybucji]
W teorii dystrybucji zbiory polarne otoczeń zera w przestrzeni
funkcji testowych są słabo-* zwarte, natomiast topologia silna
na przestrzeni dystrybucji oznacza zbieżność jednostajną
na zbiorach ograniczonych funkcji testowych. Różnica między
zbiorem ograniczonym, zwartym i równociągłym będzie podstawowa
w opisie topologii $\mathcal D'(\Omega)$.
\end{uwaga}

\subsection{Zastosowanie zwartości: nieskończenie wiele warunków liniowych}

\begin{twierdzenie}[Kryterium przedłużenia z wieloma warunkami]
\label{thm:helly-interpolation}
Niech $E$ będzie przestrzenią unormowaną nad $\K$,
$R\ge0$, $I$ dowolnym zbiorem, $(x_i)_{i\in I}$ rodziną
punktów $E$, a $(b_i)_{i\in I}$ rodziną skalarów.
Istnieje funkcjonał $f\in E'$ spełniający
$\|f\|\le R$ i $f(x_i)=b_i$ dla każdego $i\in I$
wtedy i tylko wtedy, gdy dla dowolnego $m\in\N$,
$i_1,\ldots,i_m\in I$ oraz $c_1,\ldots,c_m\in\K$
zachodzi
\[
 \left|\sum_{j=1}^{m}c_jb_{i_j}\right|
 \le R\left\|\sum_{j=1}^{m}c_jx_{i_j}\right\|.
\]
\end{twierdzenie}
\begin{proof}
Konieczność wynika z liniowości funkcjonału i oszacowania
$|f(x)|\le\|f\|\,\|x\|$.

Udowodnijmy dostateczność w sposób wykorzystujący zwartość.
Niech $B_R'=\{f\in E':\|f\|\le R\}$; jest to słabo-* zwarty
zbiór na mocy Wniosku~\ref{cor:alaoglu-norm}. Dla $i\in I$
połóżmy $C_i=\{f\in B_R':f(x_i)=b_i\}$.
Każdy $C_i$ jest słabo-* domknięty w $B_R'$.
Dla ustalonego skończonego $J\subseteq I$ definiujemy
na $M_J=\operatorname{lin}\{x_j:j\in J\}$ funkcjonał
$g_J(\sum_{j\in J}c_jx_j)=\sum_{j\in J}c_jb_j$.
Warunek z tezy zapewnia, że definicja nie zależy od
przedstawienia argumentu, a ponadto $|g_J(x)|\le R\|x\|$.
Twierdzenie Hahna--Banacha daje przedłużenie
$f_J\in E'$ z $\|f_J\|\le R$; stąd
$f_J\in\bigcap_{j\in J} C_j$.
Rodzina $(C_i)_{i\in I}$ ma więc własność skończonych przecięć.
Zwartość $B_R'$ daje $f\in\bigcap_{i\in I}C_i$,
który spełnia wszystkie żądane warunki.
\end{proof}

\begin{przyklad}[Przedłużanie wartości na wektorach współrzędnych]
Niech $E=\ell^1(\N)$, $x_n=e_n$ i $R>0$.
Dla dowolnej rodziny liczb $b_n$ ograniczonej przez
$|b_n|\le R$ otrzymujemy funkcjonał $f\in E'$
spełniający $f(e_n)=b_n$ i $\|f\|\le R$:
\[
 \left|\sum_{j=1}^m c_jb_{i_j}\right|
 \le R\sum_{n=1}^\infty
       \left|\sum_{j:i_j=n}c_j\right|
 =R\left\|\sum_{j=1}^m c_je_{i_j}\right\|_1.
\]
Jest to zgodne z identyfikacją $(\ell^1)'=\ell^\infty$.
\end{przyklad}

% Wersja źródłowa: audyt II, jednostka 10
\setcounter{section}{9}
\section{Przestrzenie beczkowe i bornologiczne. Twierdzenie Banacha--Steinhausa}\label{sec:beczkowe-bornologiczne}

\subsection{Od ograniczoności punktowej do równociągłości}

W rozdziale o twierdzeniu Banacha--Alaoglu pokazaliśmy, że rodzina
$H\subseteq E'$ jest równociągła dokładnie wtedy, gdy mieści się
w zbiorze polarnym pewnego otoczenia zera. Równociągłość implikuje
ograniczoność punktową, lecz przykład z funkcjonałami $f_n(x)=nx_n$ na $c_{00}$ pokazuje, że implikacji nie można odwrócić w każdej
przestrzeni liniowo-topologicznej. Zbadamy teraz, jakie dodatkowe
własności przestrzeni dziedziny gwarantują tę implikację.

W całym rozdziale $\K\in\{\R,\C\}$. Przestrzenie $E$ i $F$
są lokalnie wypukłymi przestrzeniami liniowo-topologicznymi Hausdorffa,
chyba że wyraźnie zaznaczono inaczej. Symbol $E'$ oznacza przestrzeń
dualną topologiczną, a $\mathcal L(E,F)$ zbiór ciągłych operatorów
liniowych z $E$ do $F$. Przez zbiór ograniczony rozumiemy zbiór
ograniczony w sensie von Neumanna. Zbiór $C$ jest \emph{absolutnie
wypukły}, jeżeli jest wypukły i zrównoważony.

\begin{definicja}[Beczka i przestrzeń beczkowa]\label{def:beczka}
Niech $E$ będzie lokalnie wypukłą TVS. Zbiór $B\subseteq E$
nazywamy \emph{beczką}, jeżeli jest domknięty, absolutnie wypukły
i pochłaniający, to znaczy dla każdego $x\in E$ istnieje $t>0$
takie, że $x\in tB$. Przestrzeń $E$ jest \emph{beczkowa}, jeśli
każda beczka w $E$ jest otoczeniem zera.
\end{definicja}

Warunek pochłaniania dotyczy pojedynczych punktów, a nie całych
zbiorów ograniczonych. To rozróżnienie będzie istotne przy
wprowadzaniu przestrzeni bornologicznych.

\begin{lemat}[Zbiór polarny i ograniczoność punktowa]\label{lem:polar-barrel}
Niech $H\subseteq E'$ i niech
$H^\circ=\{x\in E:|f(x)|\le1\ \text{dla każdego }f\in H\}$.
Wówczas $H^\circ$ jest domknięty i absolutnie wypukły.
Ponadto $H^\circ$ jest pochłaniający wtedy i tylko wtedy, gdy
$H$ jest punktowo ograniczony, czyli
$\sup_{f\in H}|f(x)|<\infty$ dla każdego $x\in E$.
\end{lemat}
\begin{proof}
Każdy zbiór $\{x:|f(x)|\le1\}$ jest domknięty i absolutnie
wypukły, a przecięcie dowolnej rodziny takich zbiorów
zachowuje obie własności. Jeśli $H$ jest punktowo ograniczony,
to dla ustalonego $x$ liczba $M_x=\sup_{f\in H}|f(x)|$
jest skończona; przyjmujemy $M_x=0$, gdy $H=\varnothing$.
Dla $t>\max\{1,M_x\}$ mamy $x/t\in H^\circ$, czyli
$x\in tH^\circ$. Odwrotnie, jeśli $x\in tH^\circ$, to
$|f(x)|\le t$ dla wszystkich $f\in H$. W obu kierunkach
dopuszczamy również $x=0$ i rodzinę pustą.
\end{proof}

\begin{twierdzenie}[Charakteryzacja beczkowości]\label{thm:barrelled-char}
Niech $E$ będzie lokalnie wypukłą TVS Hausdorffa.
Następujące warunki są równoważne:
\begin{enumerate}[label=\textup{(\roman*)}]
\item $E$ jest beczkowa;
\item każda punktowo ograniczona rodzina $H\subseteq E'$
      jest równociągła;
\item każda rodzina $H\subseteq E'$ ograniczona w
      $\sigma(E',E)$ jest równociągła.
\end{enumerate}
\end{twierdzenie}
\begin{proof}
Warunki (ii) i (iii) są równoważne, gdyż topologia
$\sigma(E',E)$ jest topologią punktowej zbieżności,
a ograniczoność zbioru w produkcie jest równoważna
ograniczoności każdej jego projekcji na $\K$.

Załóżmy (i) i niech $H$ będzie punktowo ograniczona.
Z lematu~\ref{lem:polar-barrel} wynika, że $H^\circ$
jest beczką, zatem jest otoczeniem zera. Dla
$u\in H^\circ$ zachodzi $|f(u)|\le1$ dla wszystkich
$f\in H$, czyli $H$ jest równociągła.

Załóżmy (ii) i niech $B\subseteq E$ będzie beczką.
Zbiór polarny $B^\circ\subseteq E'$ jest punktowo ograniczony:
dla każdego $x$ istnieje $t>0$ z $x\in tB$, zatem
$|f(x)|\le t$ dla wszystkich $f\in B^\circ$.
Z (ii) otrzymujemy otoczenie $U$ takie, że
$B^\circ\subseteq U^\circ$. Wskutek odwracania inkluzji
przez operację przejścia do zbioru polarnego i twierdzenia
bipolarnego dla lokalnie wypukłej przestrzeni Hausdorffa
zachodzi
$U\subseteq U^{\circ\circ}\subseteq B^{\circ\circ}=B$.
Zatem $B$ jest otoczeniem zera.
\end{proof}

\begin{uwaga}[Dlaczego konieczna jest lokalna wypukłość?]
W dowodzie odwrotnej implikacji użyliśmy równości
$B^{\circ\circ}=B$ dla domkniętego zbioru absolutnie
wypukłego. Opiera się ona na twierdzeniu Hahna--Banacha
i może zawieść w ogólnej TVS. Charakteryzację
\ref{thm:barrelled-char} formułujemy więc w klasie
przestrzeni lokalnie wypukłych Hausdorffa.
\end{uwaga}

\subsection{Własność Baire'a i przestrzenie Frécheta}

\begin{definicja}[Przestrzeń Baire'a]
Przestrzeń topologiczną $X$ nazywamy \emph{przestrzenią Baire'a},
jeżeli każde przeliczalne przecięcie otwartych zbiorów gęstych
w $X$ jest gęste. Równoważnie: przeliczalna suma zbiorów
domkniętych o pustym wnętrzu ma puste wnętrze.
\end{definicja}

\begin{twierdzenie}[Baire: zupełność metryczna wystarcza]
\label{thm:baire-complete}
Każda niepusta zupełna przestrzeń metryczna $(X,d)$
jest przestrzenią Baire'a.
\end{twierdzenie}
\begin{proof}
Niech $G_1,G_2,\ldots$ będą otwartymi podzbiorami
gęstymi w $X$. Pokażemy, że ich przecięcie przecina
każdy niepusty zbiór otwarty $O\subseteq X$.
Wybieramy indukcyjnie punkty $x_n\in X$ oraz
liczby $0<r_n<2^{-n}$ tak, by
\[
 \overline{B}(x_1,r_1)\subseteq O\cap G_1,\qquad
 \overline{B}(x_{n+1},r_{n+1})\subseteq
 B(x_n,r_n)\cap G_{n+1}.
\]
W pierwszym kroku $O\cap G_1$ jest niepusty i otwarty,
więc zawiera domkniętą kulę o dodatnim promieniu.
Analogicznie $B(x_n,r_n)\cap G_{n+1}$ jest niepusty
i otwarty, ponieważ $G_{n+1}$ jest gęsty. Możemy
zatem kontynuować konstrukcję.

Dla $m>n$ mamy $x_m\in B(x_n,r_n)$, więc
$d(x_m,x_n)<r_n<2^{-n}$. Ciąg $(x_n)$ jest ciągiem
Cauchy'ego. Z zupełności otrzymujemy punkt $x\in X$
z $x_n\to x$. Dla ustalonego $n$ wszystkie punkty
$x_m$ o $m>n$ leżą w $\overline B(x_n,r_n)$,
a domkniętość tej kuli daje
$x\in\overline B(x_n,r_n)\subseteq G_n$.
Ponadto $x\in\overline B(x_1,r_1)\subseteq O$.
Wobec dowolności $O$ przecięcie $\bigcap_nG_n$
jest gęste w $X$.
\end{proof}

\begin{twierdzenie}[Przestrzenie Baire'a są beczkowe]\label{thm:baire-barrelled}
Niech $E$ będzie lokalnie wypukłą TVS Hausdorffa, której
przestrzeń topologiczna jest przestrzenią Baire'a.
Wówczas $E$ jest beczkowa.
\end{twierdzenie}
\begin{proof}
Niech $B$ będzie beczką w $E$. Ponieważ jest pochłaniająca
i zrównoważona, $E=\bigcup_{n=1}^\infty nB$. Każdy zbiór
$nB$ jest domknięty, bo mnożenie przez niezerowy skalar
jest homeomorfizmem. Z własności Baire'a otrzymujemy
$\operatorname{int}(n_0B)\ne\varnothing$ dla pewnego
$n_0\ge1$, a zatem $\operatorname{int}B\ne\varnothing$.

Wybierzmy $y\in\operatorname{int}B$. Zrównoważenie daje
$-y\in\operatorname{int}B$. Istnieje zrównoważone
otwarte otoczenie $V$ zera takie, że
$y+V\subseteq B$ i $-y+V\subseteq B$. Dla każdego
$v\in V$ wypukłość $B$ daje
$v=\frac12(y+v)+\frac12(-y+v)\in B$.
Zatem $V\subseteq B$ i $B$ jest otoczeniem zera.
\end{proof}

\begin{wniosek}[Przestrzenie Frécheta i Banacha]\label{cor:frechet-barrelled}
Każda przestrzeń Frécheta i każda przestrzeń Banacha
jest beczkowa. Również każda lokalnie wypukła TVS
Hausdorffa, której przestrzeń topologiczna jest Baire'a,
jest beczkowa.
\end{wniosek}
\begin{proof}
Z Twierdzenia~\ref{thm:baire-complete} wynika, że
przestrzeń Frécheta jest Baire'a.
Przestrzeń Banacha jest szczególnym przypadkiem przestrzeni
Frécheta. Pozostała część jest
Twierdzeniem~\ref{thm:baire-barrelled}.
\end{proof}

\begin{uwaga}[Dwa twierdzenia o nazwie Baire'a]
Należy odróżnić własność Baire'a przestrzeni topologicznej
od twierdzenia Baire'a, które pozwala ją uzyskać z
zupełności metrycznej. Dowód twierdzenia Baire'a podaliśmy powyżej;
w następnym rozdziale zastosujemy go do twierdzeń
o odwzorowaniu otwartym i wykresie domkniętym.
\end{uwaga}

\subsection{Twierdzenie Banacha--Steinhausa w przestrzeniach lokalnie wypukłych}

\begin{definicja}[Ograniczoność punktowa i równociągłość operatorów]
Niech $\mathcal T\subseteq\mathcal L(E,F)$.
Rodzina $\mathcal T$ jest \emph{punktowo ograniczona}, gdy
$\{Tx:T\in\mathcal T\}$ jest ograniczony w $F$ dla każdego
$x\in E$. Jest \emph{równociągła}, gdy dla każdego otoczenia
$V$ zera w $F$ istnieje otoczenie $U$ zera w $E$
takie, że $T(U)\subseteq V$ dla każdego $T\in\mathcal T$.
\end{definicja}

\begin{twierdzenie}[Banach--Steinhaus dla przestrzeni beczkowych]
\label{thm:banach-steinhaus-lcs}
Niech $E$ będzie beczkową lokalnie wypukłą TVS Hausdorffa,
$F$ lokalnie wypukłą TVS Hausdorffa, a
$\mathcal T\subseteq\mathcal L(E,F)$ rodziną operatorów.
Jeżeli $\mathcal T$ jest punktowo ograniczona, to jest
równociągła.
\end{twierdzenie}
\begin{proof}
Ustalmy otoczenie $V$ zera w $F$. Z lokalnej wypukłości
możemy wybrać ciągłe seminormy $q_1,\ldots,q_m$ na $F$
i $\varepsilon>0$ takie, aby
$W=\{y:\max_jq_j(y)\le\varepsilon\}\subseteq V$.
Zbiór $W$ jest domknięty, absolutnie wypukły i stanowi
otoczenie zera w $F$. Określmy
\[
 B=\bigcap_{T\in\mathcal T}T^{-1}(W).
\]
Jest to zbiór domknięty i absolutnie wypukły w $E$.
Dla ustalonego $x\in E$ zbiór $\{Tx:T\in\mathcal T\}$
jest ograniczony w $F$. Istnieje więc $t_x>0$ takie,
że $Tx\in t_x W$ dla każdego $T\in\mathcal T$,
a zatem $x\in t_xB$. Zbiór $B$ jest beczką.
Beczkowość $E$ daje otoczenie zera $U\subseteq B$,
więc $T(U)\subseteq W\subseteq V$ dla wszystkich
$T\in\mathcal T$.
\end{proof}

\begin{wniosek}[Klasyczna zasada jednostajnej ograniczoności]
\label{cor:uniform-boundedness}
Niech $E$ będzie przestrzenią Banacha, a $F$ przestrzenią
unormowaną nad $\K$. Jeżeli rodzina
$\mathcal T\subseteq\mathcal L(E,F)$ spełnia
$\sup_{T\in\mathcal T}\|Tx\|_F<\infty$ dla każdego
$x\in E$, to $\sup_{T\in\mathcal T}\|T\|<\infty$.
Nie zakładamy zupełności $F$.
\end{wniosek}
\begin{proof}
Przestrzeń $E$ jest beczkowa. Stosujemy
Twierdzenie~\ref{thm:banach-steinhaus-lcs} do kuli
$V=\{y:\|y\|_F<1\}$. Istnieje $\delta>0$ takie, że
$\|Tx\|_F<1$ dla wszystkich $T\in\mathcal T$
i $\|x\|_E<\delta$. Przez skalowanie otrzymujemy
$\|T\|\le 2/\delta$ dla każdego $T\in\mathcal T$.
\end{proof}

\begin{proof}[Drugi dowód wniosku: bezpośrednio z twierdzenia Baire'a]
Dla $n\in\N$ niech
$A_n=\{x\in E:\sup_{T\in\mathcal T}\|Tx\|_F\le n\}$.
Są to zbiory domknięte, ponieważ każdy $T$ jest ciągły.
Ograniczoność punktowa oznacza
$E=\bigcup_{n=1}^\infty A_n$. Na mocy twierdzenia Baire'a
pewien $A_N$ ma niepuste wnętrze. Istnieją więc
$x_0\in E$ i $r>0$ takie, że
$x_0+\{h:\|h\|_E<r\}\subseteq A_N$.
Dla $\|h\|_E<r$ i każdego $T\in\mathcal T$ mamy
$\|Th\|_F\le\|T(x_0+h)\|_F+\|Tx_0\|_F\le2N$.
Skalowanie daje $\|T\|\le4N/r$, niezależnie od $T$.
W dowodzie nie wystąpiła zupełność przestrzeni wartości $F$.
\end{proof}

\begin{przyklad}[Dlaczego zupełność dziedziny jest istotna?]
\label{ex:c00-not-barrelled}
Niech $E=c_{00}$ z normą $\|x\|_\infty=\sup_n|x_n|$.
Dla $n\ge1$ określmy $f_n(x)=nx_n$. Każdy $f_n\in E'$,
ponieważ $|f_n(x)|\le n\|x\|_\infty$.
Dla ustalonego $x\in c_{00}$ ciąg $(f_n(x))$ jest
ostatecznie zerowy, więc rodzina $H=\{f_n:n\ge1\}$ jest
punktowo ograniczona. Jednak $\|f_n\|=n\to\infty$,
zatem rodzina nie jest równociągła.

Wprost widać też beczkę niebędącą otoczeniem zera:
\[
 B=\{x\in c_{00}: n|x_n|\le1\ \text{dla każdego }n\ge1\}.
\]
Jest ona domknięta, absolutnie wypukła i pochłaniająca,
gdyż każdy $x$ ma skończony nośnik. Nie jest otoczeniem
zera: dla $x^{(n)}=n^{-1/2}e_n$ zachodzi
$\|x^{(n)}\|_\infty\to0$, podczas gdy
$n|x_n^{(n)}|=\sqrt n>1$ dla $n>1$. Tym samym
$E$ nie jest beczkowa.
\end{przyklad}

\begin{przyklad}[Topologia słaba na przestrzeni Banacha]
\label{ex:weak-not-barrelled}
Niech $E$ będzie nieskończenie wymiarową przestrzenią Banacha
i wyposażmy ją w topologię $\sigma(E,E')$. Domknięta kula
normowa $B=\{x:\|x\|\le1\}$ jest słabo domknięta
(z twierdzenia Hahna--Banacha), absolutnie wypukła i pochłaniająca.
Nie jest jednak słabym otoczeniem zera, gdyż każde słabe
otoczenie zawiera niezerową prostą, która nie mieści się
w $B$. Zatem $(E,\sigma(E,E'))$ nie jest beczkowa,
mimo że $E$ jest zupełna w swojej oryginalnej normie.
O beczkowości decyduje \emph{wybrana topologia}, a nie tylko
struktura algebraiczna przestrzeni.
\end{przyklad}

\subsection{Trzy konsekwencje zasady jednostajnej ograniczoności}

\begin{twierdzenie}[Punktowa zbieżność operatorów]\label{thm:pointwise-limit-cont}
Niech $E$ będzie beczkową lokalnie wypukłą TVS Hausdorffa,
a $F$ lokalnie wypukłą TVS Hausdorffa. Załóżmy, że
$(T_i)_{i\in I}$ jest siecią w $\mathcal L(E,F)$,
rodzina $\{T_i:i\in I\}$ jest punktowo ograniczona,
a dla każdego $x\in E$ istnieje granica
$T(x)=\lim_iT_i(x)$ w $F$. Wówczas $T:E\to F$
jest ciągłym operatorem liniowym, a rodzina
$\{T_i:i\in I\}$ jest równociągła.
\end{twierdzenie}
\begin{proof}
Liniowość $T$ wynika z ciągłości dodawania i mnożenia
przez skalar w $F$. Z ograniczoności punktowej i
Twierdzenia~\ref{thm:banach-steinhaus-lcs} otrzymujemy
równociągłość rodziny. Dla danego otoczenia $V$ zera
w $F$ wybieramy domknięte otoczenie $W\subseteq V$.
Istnieje otoczenie $U$ zera w $E$ takie, że
$T_i(U)\subseteq W$ dla wszystkich $i$.
Przechodząc do granicy przy ustalonym $x\in U$, otrzymujemy
$T(x)\in W\subseteq V$. Wobec dowolności $V$ operator
$T$ jest ciągły w zerze, a z liniowości także na całym $E$.
\end{proof}

\begin{wniosek}[Granica punktowa ciągu operatorów]
\label{cor:sequence-operators}
Niech $E$ i $F$ spełniają założenia
Twierdzenia~\ref{thm:pointwise-limit-cont}, a
$T_n\in\mathcal L(E,F)$ dla $n\ge1$. Jeśli
$T_n(x)\to T(x)$ dla każdego $x\in E$, to $T$ jest
ciągłym operatorem liniowym, a rodzina $\{T_n:n\ge1\}$
jest równociągła.
\end{wniosek}
\begin{proof}
Dla każdego $x\in E$ zbiór wartości zbieżnego ciągu
$(T_n(x))$ jest ograniczony w $F$ (jego ogon zbiega do zera
po odjęciu granicy, a skończony początek jest ograniczony).
Stosujemy Twierdzenie~\ref{thm:pointwise-limit-cont}.
\end{proof}

\begin{uwaga}[Dlaczego nie wystarcza sama zbieżność sieci?]
Dla sieci zbieżnej zbiór wszystkich jej wartości nie musi
być ograniczony. Niech $I$ oznacza zbiór skończonych
podzbiorów $\N$, uporządkowany inkluzją, i niech
$a_F=|F|$, gdy $1\notin F$, oraz $a_F=0$, gdy $1\in F$.
Wtedy $a_F\to0$ jako sieć w $\R$, ponieważ na całym
ogonie zaczynającym się od $\{1\}$ mamy $a_F=0$.
Jednocześnie zbiór wszystkich $a_F$ jest nieograniczony.
Przyjmując $E=F=\R$ i $T_F(x)=a_Fx$, otrzymujemy sieć
ciągłych operatorów liniowych zbieżną punktowo do zera,
która nie jest równociągła. Dlatego w twierdzeniu dla
ogólnych sieci jawnie zakładamy ograniczoność punktową.
\end{uwaga}

\begin{twierdzenie}[Jednostajność zbieżności na zbiorach zwartych]
\label{thm:equicont-compact-conv}
Niech $E$ i $F$ będą TVS Hausdorffa,
$\mathcal T\subseteq\mathcal L(E,F)$ rodziną równociągłą,
a $(T_i)_{i\in I}$ siecią z $\mathcal T$ zbieżną punktowo
do operatora $T\in\mathcal L(E,F)$.
Wówczas $T_i\to T$ jednostajnie na każdym zwartym
zbiorze $K\subseteq E$: dla każdego otoczenia $V$ zera
w $F$ istnieje $i_0\in I$ takie, że
$(T_i-T)(K)\subseteq V$ dla $i\ge i_0$.
\end{twierdzenie}
\begin{proof}
Ustalmy $K$ i otoczenie $V$. Wybierzmy zrównoważone
otoczenie $W$ z $W+W+W\subseteq V$. Równociągłość
rodziny $\mathcal T$ i ciągłość $T$ dają wspólne
otoczenie zera $U$ w $E$, które możemy wybrać
zrównoważone, tak że $T_i(U)\subseteq W$ dla wszystkich
$i$ oraz $T(U)\subseteq W$.
Wybieramy $U$ dodatkowo otwarte. Zwartość $K$
pozwala pokryć go skończenie wieloma zbiorami
$x_1+U,\ldots,x_m+U$ o środkach $x_j\in K$.
Punktowa zbieżność zapewnia indeks $i_0$ z
$(T_i-T)x_j\in W$ dla wszystkich $j$ i $i\ge i_0$.
Jeśli $x\in K\cap(x_j+U)$, to
$(T_i-T)x=(T_i-T)x_j+T_i(x-x_j)-T(x-x_j)\in W+W+W
\subseteq V$, co kończy dowód.
\end{proof}

\begin{stwierdzenie}[Ograniczoność w topologiach operatorowych]
\label{prop:sot-bot-bounded}
Niech $E$ będzie beczkową lokalnie wypukłą TVS Hausdorffa,
a $F$ lokalnie wypukłą TVS Hausdorffa. Na $\mathcal L(E,F)$
rozważmy topologię $\tau_s$ zbieżności punktowej i
$\tau_b$ zbieżności jednostajnej na zbiorach ograniczonych.
Wtedy rodzina $\mathcal T\subseteq\mathcal L(E,F)$
jest ograniczona w $\tau_s$ wtedy i tylko wtedy,
gdy jest ograniczona w $\tau_b$.
\end{stwierdzenie}
\begin{proof}
Ponieważ singletony są ograniczone, $\tau_b$ jest
silniejsza od $\tau_s$, więc ograniczoność w $\tau_b$
implikuje ograniczoność w $\tau_s$.

Załóżmy odwrotnie, że $\mathcal T$ jest ograniczona
w $\tau_s$. Zbieżność punktowa oznacza, że dla każdego
$x\in E$ zbiór $\{Tx:T\in\mathcal T\}$ jest ograniczony
w $F$. Z twierdzenia Banacha--Steinhausa rodzina jest
równociągła. Dla ograniczonego $B\subseteq E$ i otoczenia
$V$ zera w $F$ istnieje otoczenie $U$ zera w $E$
takie, że $T(U)\subseteq V$ dla wszystkich $T$.
Ponieważ $B$ jest ograniczony, istnieje $t>0$
z $B\subseteq tU$. Stąd $T(B)\subseteq tV$ dla wszystkich
$T\in\mathcal T$. Oznacza to ograniczoność w $\tau_b$.
\end{proof}

\subsection{Przestrzenie bornologiczne}

\begin{definicja}[Zbiór pochłaniający zbiory ograniczone]
\label{def:bornivorous}
Niech $E$ będzie TVS. Zbiór $C\subseteq E$ nazywamy
\emph{pochłaniającym zbiory ograniczone}, jeśli dla każdego
ograniczonego $B\subseteq E$ istnieje $t>0$ takie, że
$B\subseteq tC$. W szczególności $C$ jest pochłaniający,
ponieważ wszystkie singletony są ograniczone.
\end{definicja}

\begin{definicja}[Przestrzeń bornologiczna]\label{def:bornological}
Lokalnie wypukłą TVS Hausdorffa $E$ nazywamy
\emph{bornologiczną}, jeżeli każdy absolutnie wypukły
zbiór pochłaniający zbiory ograniczone w $E$ jest
otoczeniem zera. Nie wymagamy jego domkniętości.
\end{definicja}

\begin{lemat}[Funkcjonał Minkowskiego i pochłanianie ograniczonych]
\label{lem:minkowski-bornivorous}
Niech $C\subseteq E$ będzie absolutnie wypukły i
pochłaniający. Jego funkcjonał Minkowskiego
$p_C(x)=\inf\{t>0:x\in tC\}$ jest skończoną seminormą.
Zbiór $C$ pochłania wszystkie zbiory ograniczone wtedy
i tylko wtedy, gdy $p_C$ jest ograniczona na każdym
zbiorze ograniczonym w $E$.
\end{lemat}
\begin{proof}
Własności seminormy wynikają z absolutnej wypukłości i
pochłaniania. Jeżeli $B\subseteq tC$, to $p_C(x)\le t$
dla $x\in B$. W drugą stronę, jeśli
$M=\sup_{x\in B}p_C(x)<\infty$, wybieramy $t>M$.
Dla każdego $x\in B$ definicja kresu dolnego zapewnia
$x\in s_x C$ dla pewnego $s_x<t$. Ponieważ $C$ jest
zrównoważony, $s_xC\subseteq tC$, stąd $B\subseteq tC$.
\end{proof}

\begin{twierdzenie}[Kryteria bornologiczności]\label{thm:bornological-equiv}
Niech $E$ będzie lokalnie wypukłą TVS Hausdorffa.
Następujące warunki są równoważne:
\begin{enumerate}[label=\textup{(\roman*)}]
\item $E$ jest bornologiczna;
\item każda seminorma $p:E\to[0,\infty)$, która jest
ograniczona na każdym ograniczonym zbiorze $B\subseteq E$,
jest ciągła;
\item dla każdej lokalnie wypukłej TVS Hausdorffa $F$
każdy liniowy operator $T:E\to F$, który odwzorowuje
zbiory ograniczone w $E$ na zbiory ograniczone w $F$,
jest ciągły.
\end{enumerate}
\end{twierdzenie}
\begin{proof}
(i)$\Rightarrow$(ii). Dla seminormy $p$ jak w (ii)
zbiór $C=\{x:p(x)\le1\}$ jest absolutnie wypukły.
Dla każdego ograniczonego $B$ mamy $M_B=\sup_Bp<\infty$,
więc $B\subseteq tC$ przy dowolnym $t>\max\{1,M_B\}$.
Bornologiczność oznacza, że $C$ jest otoczeniem zera,
co daje ciągłość $p$.

(ii)$\Rightarrow$(iii). Niech $T$ będzie operatorem
odwzorowującym zbiory ograniczone na ograniczone.
Każda ciągła seminorma $q$ na $F$ jest ograniczona
na ograniczonych podzbiorach $F$. Zatem $q\circ T$
jest seminormą ograniczoną na wszystkich ograniczonych
podzbiorach $E$. Z (ii) jest ciągła. Ponieważ topologię
$F$ wyznaczają ciągłe seminormy, operator $T$ jest ciągły.

(iii)$\Rightarrow$(i). Niech $C\subseteq E$ będzie
absolutnie wypukły i pochłaniający wszystkie zbiory
ograniczone. Jego funkcjonał Minkowskiego $p_C$
jest seminormą ograniczoną na każdym takim zbiorze.
Rozważmy przestrzeń $F$, której zbiorem punktów jest $E$,
a topologia jest generowana przez seminormy ciągłe
pierwotnej topologii oraz dodatkowo przez $p_C$.
Jest to lokalnie wypukła TVS Hausdorffa. Identyczność
$T:(E,\text{topologia pierwotna})\to F$ jest liniowa
i odwzorowuje zbiory ograniczone na ograniczone:
wszystkie seminormy wyznaczające topologię $F$
pozostają ograniczone na każdym ograniczonym $B\subseteq E$.
Z (iii) operator $T$ jest ciągły, czyli $p_C$ jest
ciągły w pierwotnej topologii. Z definicji funkcjonału
Minkowskiego wynika $\{x:p_C(x)<1\}\subseteq C$,
zatem $C$ jest otoczeniem zera. Dowiedliśmy (i).
\end{proof}

\begin{twierdzenie}[Przestrzenie unormowane są bornologiczne]
\label{thm:normed-bornological}
Każda przestrzeń unormowana nad $\K$, także niezupełna,
jest bornologiczna.
\end{twierdzenie}
\begin{proof}
Niech $E$ będzie przestrzenią unormowaną, a $C\subseteq E$
zbiorem absolutnie wypukłym, który pochłania wszystkie
zbiory ograniczone. W szczególności pochłania normową
kulę jednostkową $B=\{x:\|x\|\le1\}$, to znaczy
$B\subseteq tC$ dla pewnego $t>0$. Wtedy
$\{x:\|x\|<1/t\}\subseteq C$, więc $C$ jest
otoczeniem zera.
\end{proof}

\begin{przyklad}[Bornologiczna, ale niebeczkowa]
Przestrzeń $c_{00}$ z normą supremum jest bornologiczna
na mocy Twierdzenia~\ref{thm:normed-bornological}, lecz
nie jest beczkowa, jak pokazaliśmy w
Przykładzie~\ref{ex:c00-not-barrelled}.
Zatem bornologiczność nie wystarcza do zastosowania
twierdzenia Banacha--Steinhausa.
\end{przyklad}

\begin{uwaga}[Dwie różne zasady ciągłości]
Beczkowość pozwala przejść od \emph{punktowej ograniczoności
rodziny ciągłych operatorów} do jej równociągłości.
Bornologiczność pozwala zaś przejść od \emph{ograniczoności
pojedynczego operatora liniowego} (przenoszenia zbiorów
ograniczonych na ograniczone) do jego ciągłości.
Te dwa pojęcia mają różne cele i nie są zamienne.
\end{uwaga}

\subsection{Powrót do topologii Mackeya i przestrzeni dystrybucji}

\begin{twierdzenie}[Przestrzenie beczkowe są przestrzeniami Mackeya]
\label{thm:barrelled-mackey}
Niech $E$ będzie beczkową lokalnie wypukłą TVS Hausdorffa
z topologią $\rho$, a $E'$ jej przestrzenią dualną topologiczną.
Wówczas $\rho=\tau(E,E')$, to znaczy pierwotna topologia
jest topologią Mackeya odpowiadającą parze dualnej $(E,E')$.
\end{twierdzenie}
\begin{proof}
Topologia $\rho$ jest zgodna z $(E,E')$, więc z Twierdzenia~\ref{thm:mackey-arens} mamy
$\rho\subseteq\tau(E,E')$.
Niech $K\subseteq E'$ będzie absolutnie wypukłym zbiorem
zwartym w $\sigma(E',E)$. Dla każdego $x\in E$
zbiór $\{f(x):f\in K\}$ jest zwarty w $\K$, a zatem
ograniczony. Twierdzenie~\ref{thm:barrelled-char} daje
równociągłość $K$. Istnieje więc otoczenie $U$ zera,
na którym $p_K(x)=\sup_{f\in K}|f(x)|\le1$.
To dowodzi ciągłości $p_K$ względem $\rho$.
Wszystkie seminormy generujące $\tau(E,E')$ są więc
$\rho$-ciągłe, skąd $\tau(E,E')\subseteq\rho$.
\end{proof}

\begin{wniosek}[Topologia Mackeya przestrzeni Banacha]\label{cor:banach-mackey}
Niech $E$ będzie przestrzenią Banacha, a $E'$ jej
przestrzenią dualną topologiczną. Wówczas
$\tau(E,E')$ jest topologią normową na $E$.
\end{wniosek}
\begin{proof}
Przestrzeń Banacha jest beczkowa, zatem stosujemy
Twierdzenie~\ref{thm:barrelled-mackey}.
\end{proof}

\begin{przyklad}[Aproksymacja dystrybucji Diraca]\label{ex:delta-approx-barrel}
Niech $E=C^\infty([0,1])$ z seminormami
$p_m(\varphi)=\max_{0\le j\le m}\sup_{x\in[0,1]}|\varphi^{(j)}(x)|$,
$m\in\N_0$. Jest to przestrzeń Frécheta, a więc beczkowa.
Dla $n\ge1$ określmy
$T_n(\varphi)=n\int_{[0,1/n]}\varphi(x)\,\mathrm d\lambda(x)$.
Każdy $T_n$ jest ciągłym funkcjonałem liniowym i
$|T_n(\varphi)|\le p_0(\varphi)$, zatem rodzina
$\{T_n:n\ge1\}$ jest równociągła. Ponieważ
$|T_n(\varphi)-\varphi(0)|\le
\sup_{x\in[0,1/n]}|\varphi(x)-\varphi(0)|\to0$,
mamy $T_n(\varphi)\to\varphi(0)$. W późniejszej teorii
dystrybucji funkcjonał $\delta_0(\varphi)=\varphi(0)$
będzie przykładem dystrybucji Diraca; granica zachodzi
punktowo na $E$, a dzięki
Twierdzeniu~\ref{thm:equicont-compact-conv} również
jednostajnie na każdym zwartym zbiorze funkcji testowych
w topologii $E$.
\end{przyklad}

\begin{uwaga}[Granica zakresu obecnego rozdziału]
Nie należy utożsamiać zwartości i ograniczoności w
$C^\infty([0,1])$ ani przenosić automatycznie jej własności
na przestrzeń $\mathcal D(\Omega)$ z topologią LF.
Precyzyjne omówienie zbiorów ograniczonych w przestrzeniach
funkcji testowych nastąpi po konstrukcji granic induktywnych.
\end{uwaga}

% Wersja źródłowa: audyt II, jednostka 11
\setcounter{section}{10}
\section{Przestrzenie Frécheta. Twierdzenia o odwzorowaniu otwartym i wykresie domkniętym}\label{sec:frechet-open-graph}

\subsection{Dlaczego zupełność i metryzowalność występują razem?}

W \hyperref[sec:beczkowe-bornologiczne]{rozdziale 10} udowodniliśmy twierdzenie
Baire'a dla zupełnych przestrzeni metrycznych i wynikającą z niego
beczkowość przestrzeni Frécheta. W tym rozdziale wyprowadzimy dwa dalsze
fundamentalne rezultaty: ciągły liniowy operator surjektywny między
przestrzeniami Frécheta jest otwarty, a operator liniowy określony na
całej przestrzeni Frécheta jest ciągły, jeśli jego wykres jest domknięty
w iloczynie dwóch przestrzeni Frécheta.

W całym rozdziale $\K\in\{\R,\C\}$. Wszystkie przestrzenie
liniowo-topologiczne są Hausdorffa, chyba że podano inaczej.
Nie utożsamiamy zupełności względem kanonicznej struktury jednostajnej
z wymaganiem zbieżności jedynie ciągów Cauchy'ego bez wcześniejszego
sprawdzenia metryzowalności.

\begin{definicja}[Przestrzeń Frécheta]\label{def:frechet11}
Przestrzeń liniowo-topologiczną $E$ nad $\K$ nazywamy
\emph{przestrzenią Frécheta}, jeżeli $E$ jest lokalnie wypukła,
Hausdorffa, metryzowalna i zupełna względem swej kanonicznej
struktury jednostajnej. Oznaczenie $\mathcal L(E,F)$ oznacza
przestrzeń ciągłych operatorów liniowych $E\to F$.
\end{definicja}

\begin{stwierdzenie}[Metryka wyznaczona przez seminormy]\label{prop:frechet-metric11}
Niech $E$ będzie lokalnie wypukłą TVS Hausdorffa, której topologię
wyznacza przeliczalna rodzina ciągłych seminorm $(p_n)_{n\ge1}$.
Wtedy wzór
\[
 d(x,y)=\sum_{n=1}^{\infty}2^{-n}\min\{1,p_n(x-y)\},\qquad x,y\in E,
\]
określa metrykę niezmienniczą względem przesunięć, generującą
początkową topologię $E$. Metryka $d$ jest zupełna wtedy i tylko
wtedy, gdy $E$ jest zupełna jako TVS.
\end{stwierdzenie}
\begin{proof}
Każdy składnik jest nieujemny i niezmienniczy względem przesunięć.
Z nierówności $\min\{1,a+b\}\le \min\{1,a\}+\min\{1,b\}$
wynika nierówność trójkąta. Ponieważ seminormy rozdzielają punkty,
$d(x,y)=0$ implikuje $x=y$.

Jeśli $x_i\to x$ we wszystkich seminormach, to dla dowolnego
$\varepsilon>0$ wybieramy $N$ z $\sum_{n>N}2^{-n}<\varepsilon/2$.
Następnie pierwsze $N$ składników sumy definiującej $d(x_i,x)$
są dostatecznie małe, więc $d(x_i,x)<\varepsilon$.
Odwrotnie, z $d(x_i,x)\to0$ wynika dla ustalonego $n$,
że $\min\{1,p_n(x_i-x)\}\to0$, a więc $p_n(x_i-x)\to0$.
To dowodzi zgodności topologii, również dla sieci.

Z tego samego oszacowania wynika równoważność warunku Cauchy'ego
w metryce $d$ i warunku Cauchy'ego dla każdej seminormy $p_n$.
Metryka $d$ wyznacza więc kanoniczną strukturę jednostajną TVS,
a jej zupełność jest równoważna zupełności przestrzeni $E$.
\end{proof}

\begin{przyklad}[Podstawowe przestrzenie Frécheta]\label{ex:frechet-list11}
Każda przestrzeń Banacha jest przestrzenią Frécheta. Dalsze
przykłady to $\K^{\N}$ z topologią produktową oraz
$C^\infty([0,1])$ z seminormami
$p_n(f)=\max_{0\le j\le n}\sup_{t\in[0,1]}|f^{(j)}(t)|$,
gdzie pochodne na $(0,1)$ mają ciągłe przedłużenia na $[0,1]$.
W przestrzeni $\K^{\N}$ można wziąć seminormy
$p_n(x)=\max_{1\le j\le n}|x_j|$. Każda współrzędna ciągu
Cauchy'ego jest wtedy ciągiem Cauchy'ego w $\K$, a granice
współrzędnych dają granicę w topologii produktowej.
Wykażmy również zupełność $C^\infty([0,1])$. Jeśli $(f_n)$ jest ciągiem
Cauchy'ego we wszystkich seminormach $p_m$, to dla każdego $j\ge0$
ciąg $(f_n^{(j)})$ zbiega jednostajnie na $[0,1]$ do pewnej funkcji
ciągłej $g_j$. Z tożsamości
\[
 f_n^{(j)}(t)-f_n^{(j)}(0)
 =\int_{[0,t]}f_n^{(j+1)}(s)\,\mathrm d\lambda(s)
\]
i jednostajnej zbieżności funkcji podcałkowych otrzymujemy
$g_j(t)-g_j(0)=\int_{[0,t]}g_{j+1}(s)\,\mathrm d\lambda(s)$.
Twierdzenie podstawowe rachunku całkowego daje zatem
$g_j'=g_{j+1}$. Stąd $g_0\in C^\infty([0,1])$ oraz
$p_m(f_n-g_0)\to0$ dla każdego $m$.
\end{przyklad}

\subsection{Podprzestrzenie, iloczyny i przestrzenie ilorazowe}

\begin{twierdzenie}[Stabilność klasy przestrzeni Frécheta]\label{thm:frechet-stability11}
Niech $E$ będzie przestrzenią Frécheta, a $M\subseteq E$
domkniętą podprzestrzenią liniową. Wtedy $M$ oraz $E/M$
ze swoimi naturalnymi topologiami są przestrzeniami Frécheta.
Ponadto przeliczalny produkt przestrzeni Frécheta z topologią
produktową jest przestrzenią Frécheta.
\end{twierdzenie}
\begin{proof}
Niech $d$ będzie zgodną, zupełną i niezmienniczą względem
przesunięć metryką na $E$, jak w Stwierdzeniu~\ref{prop:frechet-metric11}.
Domknięta podprzestrzeń $M$ jest zupełna względem metryki
indukowanej z $d$, a lokalna wypukłość i metryzowalność są
dziedziczne dla podprzestrzeni.

Niech $q:E\to E/M$ będzie rzutowaniem ilorazowym.
Z domkniętości $M$ wynika, że
\[
 \bar d\bigl(q(x),q(y)\bigr)=\inf_{m\in M}d(x-y+m,0)
\]
jest metryką: nie zależy od przedstawicieli klas, spełnia
nierówność trójkąta, a $\bar d(q(x),q(y))=0$ oznacza
$x-y\in\overline M=M$. Ponadto
$q(B_d(x,r))=B_{\bar d}(q(x),r)$, więc $\bar d$ wyznacza
topologię ilorazową (rzutowanie $q$ jest otwarte).

Wykażemy zupełność. Niech $(u_n)$ będzie ciągiem Cauchy'ego
w $E/M$. Wybierzmy rosnący podciąg indeksów $(n_k)$ tak,
aby $\bar d(u_{n_{k+1}},u_{n_k})<2^{-k-1}$.
Wybierzmy przedstawiciela $z_1\in E$ klasy $u_{n_1}$.
Indukcyjnie wybieramy przedstawiciela $z_{k+1}$ klasy
$u_{n_{k+1}}$ z $d(z_{k+1},z_k)<2^{-k}$;
istnienie wynika z definicji $\bar d$. Ciąg $(z_k)$
jest Cauchy'ego w $E$, a więc zbiega do pewnego $z\in E$.
W konsekwencji $u_{n_k}=q(z_k)\to q(z)$.
Ponieważ całe $(u_n)$ jest ciągiem Cauchy'ego,
również $u_n\to q(z)$. Zatem $E/M$ jest zupełna.
Jej lokalna wypukłość wynika z obrazów otwartych,
wypukłych otoczeń zera przez $q$.

Niech teraz $(E_j)_{j\ge1}$ będzie ciągiem przestrzeni Frécheta.
Wybierzmy w każdym $E_j$ zgodną zupełną metrykę
$d_j$, niezmienniczą względem przesunięć. Metryka
\[
 D(x,y)=\sum_{j=1}^\infty 2^{-j}\min\{1,d_j(x_j,y_j)\}
\]
wyznacza topologię produktową na $\prod_{j\ge1}E_j$.
Ciąg $D$-Cauchy'ego jest Cauchy'ego na każdej współrzędnej;
po przejściu do granic współrzędnych i oszacowaniu ogona szeregu
otrzymujemy zbieżność względem $D$. Lokalna wypukłość
wynika z lokalnej wypukłości czynników.
\end{proof}

\begin{uwaga}[Znaczenie domkniętości podprzestrzeni]
Jeśli $M$ nie jest domknięta, iloraz $E/M$ z topologią
ilorazową nie jest Hausdorffa. Dlatego nie jest przestrzenią
Frécheta. Warunek domkniętości w Twierdzeniu~\ref{thm:frechet-stability11}
nie może zostać pominięty.
\end{uwaga}

\subsection{Twierdzenie o odwzorowaniu otwartym}

\begin{lemat}[Krok Baire'a]\label{lem:baire-open11}
Niech $E$ będzie TVS, $F$ lokalnie wypukłą TVS Hausdorffa
będącą przestrzenią Baire'a, a $T:E\to F$ ciągłym,
surjektywnym operatorem liniowym. Jeśli $U\subseteq E$
jest otwartym, absolutnie wypukłym otoczeniem zera, to
$\overline{T(U)}$ jest otoczeniem zera w $F$.
\end{lemat}
\begin{proof}
Ponieważ $U$ jest pochłaniający, a $T$ jest surjektywny,
$F=\bigcup_{n=1}^\infty nT(U)$. Zbiór $nT(U)$ jest
wypukły i zrównoważony, a jego domknięcie również.
Z własności Baire'a wynika, że pewien zbiór
$\overline{nT(U)}=n\overline{T(U)}$ ma niepuste wnętrze.
Stąd także $\overline{T(U)}$ ma niepuste wnętrze.
Jeżeli wypukły, zrównoważony zbiór $C$ ma niepuste wnętrze,
to zawiera otoczenie zera: dla $a\in\operatorname{int}C$
mamy także $-a\in\operatorname{int}C$, a wypukłość
daje otoczenie zera zawarte w $C$, jako zbiór
$\tfrac12((a+V)+(-a+V))=V$ dla dostatecznie małego,
otwartego, zrównoważonego otoczenia $V$ zera.
Stosujemy to do $C=\overline{T(U)}$.
\end{proof}

\begin{twierdzenie}[Twierdzenie o odwzorowaniu otwartym]\label{thm:open-mapping11}
Niech $E$ i $F$ będą przestrzeniami Frécheta nad $\K$.
Jeżeli $T:E\to F$ jest ciągłym i surjektywnym
operatorem liniowym, to $T$ jest odwzorowaniem otwartym:
dla każdego otwartego $O\subseteq E$ zbiór $T(O)$
jest otwarty w $F$.
\end{twierdzenie}
\begin{proof}
Wystarczy pokazać, że obraz dowolnego otoczenia zera
w $E$ zawiera otoczenie zera w $F$. Ustalmy takie
otoczenie $U$ i zmniejszmy je, jeśli trzeba, aby było
otwarte i absolutnie wypukłe. Wybierzmy otwarte,
absolutnie wypukłe otoczenie $U_0$ zera takie, że
$\overline{U_0}\subseteq U$. Istnienie wynika z lokalnej
wypukłości i ciągłości seminorm wyznaczających topologię.

Indukcyjnie konstruujemy otwarte, absolutnie wypukłe
otoczenia $(U_n)_{n\ge1}$ zera w $E$, takie że
$U_n+U_n\subseteq U_{n-1}$ oraz $U_n\subseteq A_n$,
gdzie $(A_n)$ jest ustaloną malejącą przeliczalną bazą
otoczeń zera w $E$. Możemy to zrobić, zmniejszając
kolejne otoczenia. Wówczas $(U_n)$ jest bazą otoczeń zera.

Z Lematu~\ref{lem:baire-open11}, zastosowanego osobno
do każdego $U_n$, wynika, że
$\overline{T(U_n)}$ zawiera otoczenie zera w $F$.
Wybieramy zatem otwarte, absolutnie wypukłe otoczenia
$V_n$ zera w $F$ takie, że
$V_n\subseteq\overline{T(U_n)}\cap B_n$, gdzie $(B_n)$
jest malejącą przeliczalną bazą otoczeń zera w $F$.

Niech $y\in V_1$. Wobec $y\in\overline{T(U_1)}$
i otwartości $V_2$ istnieje $x_1\in U_1$ takie, że
$y-Tx_1\in V_2$. Ponieważ $V_2\subseteq\overline{T(U_2)}$,
możemy wybrać $x_2\in U_2$ tak, aby
$y-T(x_1+x_2)\in V_3$. Kontynuując konstrukcję,
otrzymujemy $x_n\in U_n$ spełniające
\[
 r_n:=y-T\Bigl(\sum_{k=1}^n x_k\Bigr)\in V_{n+1}
 \quad (n\ge1).
\]
Ponieważ $V_{n+1}\subseteq B_{n+1}$, mamy $r_n\to0$ w $F$.

Pokażemy, że szereg $\sum_{n=1}^\infty x_n$ zbiega w $E$.
Z inkluzji $U_n+U_n\subseteq U_{n-1}$ wynika także
$U_n\subseteq U_{n-1}$. Indukcja wsteczna dowodzi, że
$\sum_{k=m}^N x_k\in U_{m-1}$ dla dowolnych $N\ge m\ge1$.
Ponieważ $(U_n)$ jest bazą otoczeń zera, ciąg sum częściowych
jest Cauchy'ego. Zupełność $E$ daje
$x=\lim_{N\to\infty}\sum_{k=1}^N x_k$.
Dla $m=1$ sumy częściowe należą do $U_0$, więc
$x\in\overline{U_0}\subseteq U$. Z ciągłości $T$ i
$r_n\to0$ otrzymujemy $Tx=y$. Udowodniliśmy
$V_1\subseteq T(U)$.

Dla dowolnego zbioru otwartego $O\subseteq E$ i punktu
$y=Tx\in T(O)$ wybieramy otoczenie zera $U_x$ z
$x+U_x\subseteq O$. Wówczas $T(O)$ zawiera otoczenie
$y$ postaci $y+V$, gdzie $V\subseteq T(U_x)$.
Zatem $T(O)$ jest otwarty.
\end{proof}

\begin{uwaga}[Co jest rzeczywiście potrzebne?]\label{rem:open-assumptions11}
Dowód wykorzystuje zupełność i metryzowalność dziedziny,
własność Baire'a oraz metryzowalność przestrzeni wartości,
a także lokalną wypukłość obu przestrzeni. W konsekwencji
ten sam argument działa dla $E$ będącej przestrzenią Frécheta
oraz $F$ będącej metryzowalną lokalnie wypukłą TVS Hausdorffa
o własności Baire'a, nawet jeśli zupełności $F$ nie założono.
Nie twierdzimy natomiast, że dowolny ciągły operator
liniowy między dwiema przestrzeniami Frécheta jest otwarty:
konieczna jest surjektywność.
\end{uwaga}

\begin{wniosek}[Twierdzenie o ciągłości odwrotności]\label{cor:inverse-mapping11}
Niech $E,F$ będą przestrzeniami Frécheta, a
$T\in\mathcal L(E,F)$ liniową bijekcją.
Wtedy $T^{-1}:F\to E$ jest ciągłym operatorem liniowym.
Zatem $T$ jest izomorfizmem liniowo-topologicznym.
\end{wniosek}
\begin{proof}
Z Twierdzenia~\ref{thm:open-mapping11} operator $T$ jest otwarty.
Dla każdego otwartego $U\subseteq E$ zbiór
$(T^{-1})^{-1}(U)=T(U)$ jest zatem otwarty w $F$.
\end{proof}

\begin{wniosek}[Oszacowanie normy operatora odwrotnego]\label{cor:bounded-inverse11}
Niech $E,F$ będą przestrzeniami Banacha i niech
$T:E\to F$ będzie liniowym, ciągłym i bijektywnym
operatorem. Istnieje stała $C>0$ taka, że
$\|T^{-1}y\|_E\le C\|y\|_F$ dla każdego $y\in F$.
Równoważnie, $\|Tx\|_F\ge C^{-1}\|x\|_E$ dla $x\in E$.
\end{wniosek}
\begin{proof}
Z Wniosku~\ref{cor:inverse-mapping11} odwrotność jest ciągła.
Liniowy operator między przestrzeniami unormowanymi jest
ciągły wtedy i tylko wtedy, gdy jest ograniczony; przyjmujemy
$C=\max\{1,\|T^{-1}\|\}$.
\end{proof}

\begin{wniosek}[Izomorfizm z przestrzenią ilorazową]\label{cor:quotient-open11}
Niech $E,F$ będą przestrzeniami Frécheta oraz
$T\in\mathcal L(E,F)$ operatorem surjektywnym.
Wtedy $M=\ker T$ jest domkniętą podprzestrzenią $E$,
a operator
$\widetilde T:E/M\to F$, $\widetilde T(x+M)=Tx$,
jest izomorfizmem liniowo-topologicznym.
\end{wniosek}
\begin{proof}
Ciągłość $T$ i domkniętość $\{0\}\subset F$ dają
domkniętość $M$. Własność uniwersalna ilorazu zapewnia
ciągłość $\widetilde T$. Jest ona liniową bijekcją.
Ponieważ $E/M$ i $F$ są przestrzeniami Frécheta,
Wniosek~\ref{cor:inverse-mapping11} daje ciągłość odwrotności.
\end{proof}

\subsection{Twierdzenie o wykresie domkniętym}

\begin{definicja}[Wykres operatora]\label{def:graph11}
Jeżeli $E,F$ są TVS, a $T:D(T)\subseteq E\to F$
odwzorowaniem, to jego \emph{wykresem} nazywamy
$G(T)=\{(x,Tx):x\in D(T)\}\subseteq E\times F$.
Dla operatora liniowego $G(T)$ jest podprzestrzenią
liniową produktu.
\end{definicja}

\begin{stwierdzenie}[Kryterium ciągowe domknięcia wykresu]
\label{prop:graph-sequential11}
Niech $E,F$ będą metryzowalnymi TVS Hausdorffa, a
$T:E\to F$ operatorem liniowym określonym na całym $E$.
Wykres $G(T)$ jest domknięty w $E\times F$ wtedy i tylko
wtedy, gdy dla każdego ciągu $(x_n)$ w $E$ i każdego
$x\in E$, $y\in F$ spełniających $x_n\to x$ i
$Tx_n\to y$ zachodzi $Tx=y$.
\end{stwierdzenie}
\begin{proof}
Produkt $E\times F$ jest metryzowalny. W przestrzeni
metrycznej zbiór jest domknięty dokładnie wtedy, gdy
zawiera granicę każdego zbieżnego ciągu swoich punktów.
Zastosowanie do ciągu $(x_n,Tx_n)$ daje tezę.
\end{proof}

\begin{twierdzenie}[Twierdzenie o wykresie domkniętym]\label{thm:closed-graph11}
Niech $E,F$ będą przestrzeniami Frécheta, a
$T:E\to F$ operatorem liniowym określonym na całym $E$.
Wówczas $T$ jest ciągły wtedy i tylko wtedy,
gdy jego wykres $G(T)$ jest domknięty w $E\times F$.
\end{twierdzenie}
\begin{proof}
Jeżeli $T$ jest ciągły, to odwzorowanie
$(x,y)\mapsto y-Tx$ jest ciągłe z $E\times F$ do $F$.
Jego przeciwobraz zbioru domkniętego $\{0\}$ to $G(T)$,
a więc wykres jest domknięty.

Odwrotnie, załóżmy, że $G(T)$ jest domknięty.
Iloczyn $E\times F$ jest przestrzenią Frécheta
na mocy Twierdzenia~\ref{thm:frechet-stability11},
a jego domknięta podprzestrzeń $G(T)$ również.
Rzutowanie
$P:G(T)\to E$, $P(x,Tx)=x$, jest ciągłą liniową
bijekcją między przestrzeniami Frécheta. Z Wniosku
\ref{cor:inverse-mapping11} wynika ciągłość $P^{-1}$.
Stąd $T=\pi_F\circ P^{-1}$ jest ciągły, gdzie
$\pi_F:E\times F\to F$ oznacza drugie rzutowanie.
\end{proof}

\begin{center}
\begin{tikzpicture}[>=Stealth, every node/.style={font=\small}]
  \node (g) at (0,1.65) {$G(T)\subseteq E\times F$};
  \node (e) at (-2.75,0) {$E$};
  \node (f) at (2.75,0) {$F$};
  \draw[->] (g) -- node[above left] {$P$} (e);
  \draw[->] (g) -- node[above right] {$\pi_F|_{G(T)}$} (f);
  \draw[->,dashed] (e) -- node[below] {$T=\pi_F\circ P^{-1}$} (f);
\end{tikzpicture}

\smallskip
\small Schemat dowodu twierdzenia o wykresie domkniętym:
$P$ jest ciągłą bijekcją, a jego odwrotność jest ciągła
na mocy twierdzenia o odwzorowaniu otwartym.
\end{center}

\begin{uwaga}[Pełna dziedzina operatora jest niezbędna]
Twierdzenie~\ref{thm:closed-graph11} wymaga, by $T$ był
określony na całym $E$. Operatory różniczkowe jako
operatory nieograniczone w przestrzeniach Banacha mają
zwykle właściwe dziedziny $D(T)\subsetneq E$ i mogą
mieć domknięte wykresy bez ciągłości w normie $E$.
Przykład podamy poniżej.
\end{uwaga}

\subsection{Kontrprzykłady i zastosowania}

\begin{przyklad}[Brak zupełności przeciwdziedziny: bijekcja nieotwarta]
\label{ex:incomplete-codomain11}
Na zbiorze $\ell^1$ rozważmy dwie normy
$\|x\|_1=\sum_{n\ge1}|x_n|$ oraz
$\|x\|_2=(\sum_{n\ge1}|x_n|^2)^{1/2}$.
Operator identycznościowy
$J:(\ell^1,\|\cdot\|_1)\to(\ell^1,\|\cdot\|_2)$
jest ciągłą liniową bijekcją, bo $\|x\|_2\le\|x\|_1$.
Jego odwrotność nie jest ciągła. Dla wektora $u^{(n)}$
o pierwszych $n$ współrzędnych równych $1/n$
i pozostałych równych zeru mamy
$\|u^{(n)}\|_1=1$ i $\|u^{(n)}\|_2=n^{-1/2}\to0$.
Zatem $J$ nie jest otwarty. Przeciwdziedzina
$(\ell^1,\|\cdot\|_2)$ nie jest zupełna: jej
uzupełnieniem jest $\ell^2$, ponieważ $c_{00}$ jest
gęste w $\ell^2$.
\end{przyklad}

\begin{przyklad}[Wykres domknięty nie zapewnia ciągłości bez zupełności]
\label{ex:incomplete-domain-graph11}
Niech $E=(\ell^1,\|\cdot\|_2)$, $F=(\ell^1,\|\cdot\|_1)$,
a $T:E\to F$ będzie odwzorowaniem identycznościowym.
Nie jest ono ciągłe, jak pokazuje ciąg $u^{(n)}$ z
Przykładu~\ref{ex:incomplete-codomain11}. Ma jednak
wykres domknięty w $E\times F$. Istotnie, jeśli
$x_n\to x$ w normie $\ell^2$ oraz
$Tx_n=x_n\to y$ w normie $\ell^1$, to $x_n\to y$
także w normie $\ell^2$. Jednoznaczność granicy
w $\ell^2$ daje $x=y$, czyli $(x,y)\in G(T)$.
Kryterium ciągowe domknięcia wykresu kończy dowód.
Tutaj $E$ jest niezupełna, mimo że $F$ jest przestrzenią Banacha.
\end{przyklad}

\begin{przyklad}[Domknięty nieograniczony operator różniczkowania]
\label{ex:closed-derivative11}
Niech $E=F=C([0,1])$ z normą supremum i
$D(T)=C^1([0,1])$. Określmy $Tf=f'$. Jest to liniowy
operator nieograniczony względem normy $E$:
$f_n(t)=n^{-1}\sin(nt)$ spełnia $\|f_n\|_\infty\to0$,
podczas gdy $\|Tf_n\|_\infty=1$.
Wykres $G(T)$ jest jednak domknięty w $E\times F$.
Jeżeli $f_n\to f$ i $f_n'\to g$ jednostajnie,
to dla $t\in[0,1]$ z twierdzenia podstawowego rachunku
całkowego mamy
\[
 f_n(t)-f_n(0)=\int_{[0,t]}f_n'(s)\,\mathrm d\lambda(s).
\]
Przechodząc do granicy, otrzymujemy
$f(t)-f(0)=\int_{[0,t]}g(s)\,\mathrm d\lambda(s)$.
Ponieważ $g$ jest ciągła, funkcja $f$ jest klasy $C^1$
i $f'=g$. Nie przeczy to Twierdzeniu~\ref{thm:closed-graph11},
ponieważ $D(T)\ne E$.
\end{przyklad}

\begin{przyklad}[Operator pochodnej w przestrzeni Frécheta]
\label{ex:frechet-derivative11}
Niech $E=F=C^\infty([0,1])$ z seminormami $p_n$
z Przykładu~\ref{ex:frechet-list11} i niech $Df=f'$.
Operator $D:E\to F$ jest ciągły, ponieważ
$p_n(Df)\le p_{n+1}(f)$ dla każdego $n$.
Jest również surjektywny: dla każdego $g\in F$ funkcja
$Rg(t)=\int_{[0,t]}g(s)\,\mathrm d\lambda(s)$ należy do $E$
i spełnia $D(Rg)=g$. Jądro $D$ składa się z funkcji stałych.
Zatem Twierdzenie~\ref{cor:quotient-open11} daje
izomorfizm liniowo-topologiczny
$C^\infty([0,1])/\K\cong C^\infty([0,1])$,
gdzie $\K$ oznacza podprzestrzeń funkcji stałych.
W istocie $R$ jest ciągły, ponieważ
$p_n(Rg)\le 2p_n(g)$ dla wszystkich $n\ge0$.
\end{przyklad}

\begin{przyklad}[Gęsty, niedomknięty obraz operatora]
\label{ex:dense-range11}
Naturalne włożenie $\iota:\ell^1\to\ell^2$ jest
ciągłe i różnowartościowe. Jego obraz jest gęsty, gdyż
zawiera $c_{00}$, ale nie jest cały: ciąg
$(n^{-3/4})_{n\ge1}$ należy do $\ell^2\setminus\ell^1$.
Odwrotność $\iota^{-1}:\iota(\ell^1)\to\ell^1$
nie jest ciągła przy topologii indukowanej z $\ell^2$.
Twierdzenie o odwzorowaniu otwartym nie ma zastosowania
do operatora $\ell^1\to\iota(\ell^1)$, ponieważ
przeciwdziedzina z topologią indukowaną nie jest zupełna.
\end{przyklad}

\begin{wniosek}[Równoważność zupełnych norm porównywalnych]
\label{cor:two-banach-norms11}
Niech $V$ będzie przestrzenią wektorową nad $\K$,
a $\|\cdot\|_a$ oraz $\|\cdot\|_b$ normami na $V$
takimi, że obie przestrzenie $(V,\|\cdot\|_a)$ i
$(V,\|\cdot\|_b)$ są Banacha. Jeśli istnieje $C>0$
z $\|x\|_b\le C\|x\|_a$ dla wszystkich $x\in V$,
to istnieje $D>0$ takie, że
$\|x\|_a\le D\|x\|_b$ dla wszystkich $x\in V$.
\end{wniosek}
\begin{proof}
Operator identycznościowy między podanymi przestrzeniami
Banacha jest ciągłą bijekcją. Jego odwrotność jest ciągła
na mocy Wniosku~\ref{cor:inverse-mapping11}, co jest
równoważne żądanemu oszacowaniu.
\end{proof}

\begin{uwaga}[Pomost do rachunku różniczkowego]
W dalszych rozdziałach będziemy badać pochodne Gâteaux
i Frécheta odwzorowań między przestrzeniami Banacha.
W twierdzeniu o lokalnej funkcji odwrotnej zakłada się
odwracalność pochodnej Frécheta, która jest ograniczonym
operatorem liniowym. Wniosek~\ref{cor:bounded-inverse11}
pokazuje, dlaczego algebraiczna odwracalność takiej
pochodnej między przestrzeniami Banacha zapewnia
ciągłość jej odwrotności. Nie należy mylić tego z
odwracalnością samego odwzorowania nieliniowego.
\end{uwaga}

% Wersja źródłowa: audyt II, jednostka 12
\setcounter{section}{11}
\section{Klasyczne przestrzenie Banacha i ich przestrzenie sprzężone}
\label{sec:klasyczne-banach12}

\providecommand{\atlasref}[2]{#2}


\subsection{Przejście od ogólnych przestrzeni liniowo-topologicznych do norm}

Dotychczas prowadziliśmy rozważania dla przestrzeni liniowo-topologicznych,
przestrzeni lokalnie wypukłych i przestrzeni Frécheta. Teraz zawężamy
założenia, aby otrzymać konkretne wzory na normy, ciągłe funkcjonały
i przestrzenie sprzężone. Będziemy pracować nad $\K\in\{\R,\C\}$.
Niech $E'$ oznacza przestrzeń ciągłych funkcjonałów $\K$-liniowych na $E$.
Jeżeli $E$ jest unormowana, oznaczamy przez $\|\cdot\|_{E'}$ normę
operatorową na $E'$. Dla $1<p<\infty$ liczbę $q=p/(p-1)$ nazywamy
\emph{wykładnikiem sprzężonym} do $p$; dla $p=1$ przyjmujemy $q=\infty$.

\begin{definicja}[Przestrzeń unormowana i przestrzeń Banacha]
\label{def:banach12}
Przestrzenią unormowaną nazywamy parę $(E,\|\cdot\|)$, gdzie $E$
jest przestrzenią wektorową nad $\K$, a norma spełnia
$\|x\|=0\iff x=0$, $\|\alpha x\|=|\alpha|\|x\|$ oraz
$\|x+y\|\le\|x\|+\|y\|$. Przestrzeń unormowaną nazywamy
\emph{przestrzenią Banacha}, jeżeli jest zupełna względem metryki
$d(x,y)=\|x-y\|$.
\end{definicja}

\begin{stwierdzenie}[Norma operatora i przestrzeń sprzężona]
\label{prop:operator-norm12}
Niech $E$ będzie przestrzenią unormowaną i $F$ przestrzenią Banacha.
Wtedy przestrzeń $\mathcal L(E,F)$ wszystkich ciągłych operatorów
liniowych $E\to F$, z normą
$\|T\|=\sup_{\|x\|_E\le1}\|Tx\|_F$, jest przestrzenią Banacha.
W szczególności $E'=\mathcal L(E,\K)$ jest zawsze przestrzenią
Banacha, nawet gdy $E$ nie jest zupełna.
\end{stwierdzenie}
\begin{proof}
Ciągłość liniowego $T$ w zerze jest równoważna istnieniu $C\ge0$
takiego, że $\|Tx\|_F\le C\|x\|_E$: dla implikacji nietrywialnej
wystarczy przeskalować niezerowe $x$ do kuli, na której $T$ jest
ograniczony. Norma operatorowa jest zatem skończona i spełnia
aksjomaty normy.

Jeśli $(T_n)$ jest ciągiem Cauchy'ego w normie operatorowej,
to dla każdego $x\in E$ ciąg $(T_nx)$ jest Cauchy'ego w $F$, ponieważ
$\|T_nx-T_mx\|\le\|T_n-T_m\|\|x\|$.
Niech $Tx=\lim_nT_nx$. Granice zachowują liniowość. Ustalmy
$\varepsilon>0$ i wybierzmy $N$ tak, że
$\|T_n-T_m\|<\varepsilon$ dla $m,n\ge N$.
Po przejściu do granicy $m\to\infty$ otrzymujemy
$\|(T_n-T)x\|\le\varepsilon\|x\|$ dla $n\ge N$.
Stąd $T$ jest ciągły i $\|T_n-T\|\to0$.
\end{proof}

\subsection{Przestrzenie ciągowe}

\begin{definicja}[Przestrzenie $\ell^p$, $c_0$ i $\ell^\infty$]
\label{def:sequence-spaces12}
Dla $1\le p<\infty$ określamy
$\ell^p=\{x=(x_n)_{n\ge1}\in\K^{\N}:\sum_{n\ge1}|x_n|^p<\infty\}$
i $\|x\|_p=(\sum_{n\ge1}|x_n|^p)^{1/p}$.
Ponadto $\ell^\infty$ jest przestrzenią ciągów ograniczonych
z normą $\|x\|_\infty=\sup_{n\ge1}|x_n|$, a $c_0$ oznacza
podprzestrzeń ciągów zbieżnych do zera, z tą samą normą.
Przez $c_{00}$ oznaczamy przestrzeń ciągów o skończonym nośniku.
\end{definicja}

\begin{lemat}[Nierówność H\"oldera dla ciągów]
\label{lem:holder-seq12}
Dla $1<p<\infty$, $q=p/(p-1)$, $x\in\ell^p$ i $y\in\ell^q$
zachodzi
$\sum_{n\ge1}|x_ny_n|\le\|x\|_p\|y\|_q$.
Dla $p=1$, $q=\infty$, ta sama nierówność jest prawdziwa.
\end{lemat}
\begin{proof}
Dla $a,b\ge0$ nierówność Younga
$ab\le a^p/p+b^q/q$ wynika z wypukłości funkcji
$t\mapsto e^t$ (równoważnie, z minimalizacji
$t^p/p-tb$ względem $t\ge0$). Jeżeli obie normy są niezerowe,
stosujemy ją do $a=|x_n|/\|x\|_p$ i
$b=|y_n|/\|y\|_q$, sumujemy po $n\le N$,
a następnie przechodzimy do granicy $N\to\infty$.
Przypadki z normą zerową są oczywiste. Dla $p=1$
oszacowujemy bezpośrednio każdy wyraz przez
$|x_n|\|y\|_\infty$.
\end{proof}

\begin{twierdzenie}[Zupełność klasycznych przestrzeni ciągowych]
\label{thm:seq-banach12}
Dla każdego $1\le p\le\infty$ przestrzeń $\ell^p$ jest Banacha.
Przestrzeń $c_0$ jest domkniętą podprzestrzenią $\ell^\infty$,
a zatem także przestrzenią Banacha. Ponadto $c_{00}$ jest gęste
w $\ell^p$ dla $1\le p<\infty$ i w $c_0$, lecz nie jest gęste
w $\ell^\infty$.
\end{twierdzenie}
\begin{proof}
Dla $p=\infty$ ciąg Cauchy'ego $(x^{(k)})$ jest Cauchy'ego
jednostajnie względem współrzędnych. Istnieją granice
$x_n=\lim_kx_n^{(k)}$, a przejście do granicy w oszacowaniu
$\|x^{(k)}-x^{(m)}\|_\infty\le\varepsilon$ daje
$\|x^{(k)}-x\|_\infty\le\varepsilon$. Ciąg $x$ jest ograniczony.

Nierówność Minkowskiego dla skończonych sum ciągów
wynika z nierówności H\"oldera tak samo jak dla całek:
$\|u+v\|_p^p\le(\|u\|_p+\|v\|_p)\|u+v\|_p^{p-1}$
dla $p>1$, zaś dla $p=1$ jest natychmiastowa.

Dla $1\le p<\infty$ wybierzmy z ciągu Cauchy'ego
podciąg $(x^{(k_j)})$ spełniający
$\|x^{(k_{j+1})}-x^{(k_j)}\|_p\le2^{-j}$.
Niech $h_j=x^{(k_{j+1})}-x^{(k_j)}$.
Nierówność Minkowskiego dla sum skończonych daje
$\|\sum_{j=1}^N|h_j|\|_p\le\sum_{j=1}^N\|h_j\|_p\le1$.
Przez monotoniczność sum współrzędnych i ich granic
otrzymujemy $s=\sum_{j\ge1}|h_j|\in\ell^p$.
Szereg $\sum_jh_j$ jest więc bezwzględnie zbieżny
na każdej współrzędnej i określa $h\in\ell^p$;
ponadto norma jego ogona nie przekracza $\sum_{j\ge N}2^{-j}$.
Wobec tego $x^{(k_j)}\to x^{(k_1)}+h$ w $\ell^p$.
Cały ciąg Cauchy'ego ma tę samą granicę.

Jeśli $x^{(k)}\in c_0$ i $x^{(k)}\to x$ jednostajnie,
to dla $\varepsilon>0$ wybieramy $k$ z
$\|x-x^{(k)}\|_\infty<\varepsilon/2$, a następnie $N$
z $|x_n^{(k)}|<\varepsilon/2$ dla $n\ge N$.
Wtedy $x\in c_0$. Obcięcie ciągu do pierwszych $N$
współrzędnych daje zbieżność w $\ell^p$ dla $p<\infty$
i w $c_0$. W $\ell^\infty$ odległość stałego ciągu
$(1,1,\ldots)$ od $c_{00}$ wynosi co najmniej $1$.
\end{proof}

\begin{twierdzenie}[Przestrzenie sprzężone $\ell^p$ i $c_0$]
\label{thm:dual-sequences12}
Niech $1\le p<\infty$ i niech $q$ będzie wykładnikiem sprzężonym.
Każdy $F\in(\ell^p)'$ ma dokładnie jedną reprezentację
\[
 F(x)=\sum_{n=1}^\infty a_nx_n,\qquad x\in\ell^p,
 \quad a\in\ell^q,
\]
a odwzorowanie $a\mapsto F$ jest izometrycznym izomorfizmem
$\ell^q\cong(\ell^p)'$ (dla $p=1$ przyjmujemy $q=\infty$).
Ponadto każdy $F\in(c_0)'$ ma dokładnie jedną reprezentację
tego samego typu z $a\in\ell^1$ i $\|F\|=\|a\|_1$.
\end{twierdzenie}
\begin{proof}
Niech $e_n$ będzie $n$-tym wektorem standardowym i niech
$a_n=F(e_n)$. Dla $1<p<\infty$ i $N\ge1$ oznaczmy
$S_N=\sum_{n=1}^N|a_n|^q$. Jeśli $S_N>0$, wybieramy
$x_n=\overline{a_n}|a_n|^{q-2}S_N^{-1/p}$
dla $n\le N$ (a dla $a_n=0$ przyjmujemy $x_n=0$)
i $x_n=0$ dla $n>N$. Wtedy
$\|x\|_p=1$ i $F(x)=S_N^{1/q}$, ponieważ
$(q-1)p=q$. Zatem $S_N^{1/q}\le\|F\|$ dla każdego $N$.
Stąd $a\in\ell^q$ i $\|a\|_q\le\|F\|$.
Dla $p=1$ mamy bezpośrednio
$|a_n|\le\|F\|\|e_n\|_1=\|F\|$.

Ponieważ $c_{00}$ jest gęste w $\ell^p$ dla $p<\infty$,
na mocy ciągłości i Lematu~\ref{lem:holder-seq12}
$F(x)=\sum_na_nx_n$ dla każdego $x\in\ell^p$.
Nierówność H\"oldera daje $\|F\|\le\|a\|_q$,
a powyższe oszacowania dają równość.

Dla $F\in(c_0)'$ zdefiniujmy ponownie $a_n=F(e_n)$.
Dla każdego $N$ wektor $x\in c_{00}$ o współrzędnych
$x_n=\overline{a_n}/|a_n|$ przy $a_n\ne0$ i
$x_n=0$ przy $a_n=0$, $n\le N$, ma normę co najwyżej $1$.
Zatem $\sum_{n=1}^N|a_n|\le\|F\|$.
Wobec tego $a\in\ell^1$, a gęstość $c_{00}$ w $c_0$
oraz nierówność $|\sum_na_nx_n|\le\|a\|_1\|x\|_\infty$
kończą dowód wraz z identycznością norm.
Jednoznaczność w obu przypadkach wynika z podstawienia $x=e_n$.
\end{proof}

\begin{przyklad}[Dlaczego $(\ell^\infty)'$ jest większa niż $\ell^1$?]
\label{ex:linfty-dual-seq12}
Niech $u=(1,1,\ldots)\in\ell^\infty$. Na podprzestrzeni $c=c_0+\K u\subseteq\ell^\infty$
ciągów zbieżnych określmy $L(x)=\lim_{n\to\infty}x_n$.
Jest to ciągły funkcjonał liniowy o normie $1$.
Z twierdzenia Hahna--Banacha przedłużamy go do
$\widetilde L\in(\ell^\infty)'$ z tą samą normą.
Mamy $\widetilde L(e_n)=0$ dla każdego $n$, ale
$\widetilde L(u)=1$. Gdyby $\widetilde L$
był reprezentowany przez pewne $a\in\ell^1$ wzorem
$\widetilde L(x)=\sum_na_nx_n$, to z podstawienia $e_n$
wynikałoby $a_n=0$ dla wszystkich $n$, co przeczy
$\widetilde L(u)=1$.
\end{przyklad}

\subsection{Przestrzenie \texorpdfstring{$L^p$}{Lp} na przestrzeni z miarą skończoną}

\noindent\emph{Powiązania tematyczne w Bibliotece matematycznej:} \atlasref{ME-14::thm:riesz-fischer-14}{ME-14: twierdzenie Riesza--Fischera}; \atlasref{ME-15::thm:dual-lp-15}{ME-15: reprezentacja funkcjonałów na przestrzeniach funkcji całkowalnych}.


W tej podsekcji $(\Omega,\mathcal A,\mu)$ jest przestrzenią
z miarą dodatnią i skończoną, tj. $\mu(\Omega)<\infty$.
Rozważamy funkcje mierzalne o wartościach w $\K$, utożsamiane,
gdy są równe $\mu$-prawie wszędzie. Zapis
$\int_\Omega f(x)\,\mathrm d\mu(x)$ zawsze oznacza całkę
względem wskazanej miary. Jeżeli $\mu(\Omega)=0$, wszystkie
przestrzenie $L^p(\mu)$ są zerowe i twierdzenia pozostają prawdziwe.

\begin{definicja}[Przestrzenie $L^p(\mu)$]
\label{def:lp12}
Dla $1\le p<\infty$ niech $L^p(\mu)$ oznacza przestrzeń
klas funkcji $f$ spełniających
$\int_\Omega|f(x)|^p\,\mathrm d\mu(x)<\infty$, z normą
$\|f\|_p=(\int_\Omega|f(x)|^p\,\mathrm d\mu(x))^{1/p}$.
Przestrzeń $L^\infty(\mu)$ składa się z klas funkcji istotnie
ograniczonych z normą
$\|f\|_\infty=\operatorname*{ess\,sup}_{x\in\Omega}|f(x)|$.
\end{definicja}

\begin{twierdzenie}[H\"older, Minkowski i zupełność $L^p$]
\label{thm:lp-banach12}
Dla $1\le p\le\infty$ przestrzeń $L^p(\mu)$ jest Banacha.
Dla $1<p<\infty$ i $q=p/(p-1)$ zachodzi nierówność H\"oldera
$\int_\Omega|f(x)g(x)|\,\mathrm d\mu(x)\le\|f\|_p\|g\|_q$,
a dla $p=1$ analogiczny wzór z $q=\infty$.
Dla każdego $1\le p\le\infty$ zachodzi nierówność Minkowskiego
$\|f+g\|_p\le\|f\|_p+\|g\|_p$.
\end{twierdzenie}
\begin{proof}
Dowód H\"oldera powtarza argument z nierównością Younga
z Lematu~\ref{lem:holder-seq12}, po zastąpieniu sumy całką.
Dla $1<p<\infty$ najpierw $|f+g|^p\le2^{p-1}(|f|^p+|g|^p)$
zapewnia $f+g\in L^p$. Następnie nierówność Minkowskiego
otrzymujemy z $|f+g|^p\le(|f|+|g|)|f+g|^{p-1}$
i nierówności H\"oldera
i podzielenia przez $\|f+g\|_p^{p-1}$, o ile norma ta jest
niezerowa. Dla $p=1$ i $p=\infty$ jest to bezpośrednia
nierówność trójkąta.

Niech $(f_n)$ będzie ciągiem Cauchy'ego w $L^p$ dla
$1\le p<\infty$. Wybierzmy podciąg $(f_{n_k})$ taki, że
$\|f_{n_{k+1}}-f_{n_k}\|_p\le2^{-k}$.
Dla reprezentantów $h_k=f_{n_{k+1}}-f_{n_k}$,
Minkowski i twierdzenie o zbieżności monotonicznej dają
\[
 s(x)=\sum_{k\ge1}|h_k(x)|<\infty\quad\text{dla prawie każdego }x,
 \qquad \|s\|_p\le\sum_{k\ge1}2^{-k}.
\]
Szereg $\sum_kh_k$ zbiega więc prawie wszędzie do mierzalnej
funkcji $h\in L^p$, a normy ogonów są oszacowane przez
$\sum_{k\ge N}2^{-k}$. Stąd $f_{n_k}\to f_{n_1}+h$ w $L^p$.
Cauchy'ego całego ciągu zapewnia tę samą granicę.

Dla $p=\infty$ z warunku Cauchy'ego wybieramy podciąg
$(f_{n_k})$ z $\|f_{n_{k+1}}-f_{n_k}\|_\infty\le2^{-k}$.
Po usunięciu przeliczalnej sumy zbiorów zerowych
reprezentantów mamy punktową jednostajną zbieżność
szeregu różnic do istotnie ograniczonej funkcji.
Reszta dowodu przebiega jak dla $p<\infty$.
\end{proof}

\begin{lemat}[Gęstość funkcji prostych]
\label{lem:simple-dense12}
Dla $1\le p<\infty$ funkcje proste są gęste w $L^p(\mu)$.
Każda funkcja prosta jest kombinacją liniową funkcji charakterystycznych
zbiorów mierzalnych.
\end{lemat}
\begin{proof}
Dla rzeczywistej funkcji nieujemnej wybieramy rosnące funkcje
proste $s_n$ z $0\le s_n\uparrow f$ prawie wszędzie
(np. przez obcinanie wartości do $n$ i zaokrąglanie w dół
do wielokrotności $2^{-n}$). Wtedy
$|f-s_n|^p\le|f|^p$ i $|f-s_n|^p\to0$ prawie wszędzie.
Twierdzenie o zbieżności zmajoryzowanej daje
$\|f-s_n\|_p\to0$. Dla funkcji rzeczywistych ogólnych
i zespolonych rozdzielamy części dodatnie i ujemne
oraz części rzeczywiste i urojone.
\end{proof}

\begin{twierdzenie}[Przestrzeń sprzężona $L^p$ dla $1\le p<\infty$]
\label{thm:lp-dual12}
Niech $(\Omega,\mathcal A,\mu)$ będzie przestrzenią z miarą
skończoną, niech $1\le p<\infty$ i $q$ będzie wykładnikiem
sprzężonym. Dla każdego $F\in(L^p(\mu))'$ istnieje dokładnie
jedna klasa $g\in L^q(\mu)$ taka, że
\[
 F(f)=\int_\Omega f(x)g(x)\,\mathrm d\mu(x)
 \quad\text{dla wszystkich }f\in L^p(\mu).
\]
Ponadto $\|F\|=\|g\|_q$, więc
$(L^p(\mu))'\cong L^q(\mu)$ izometrycznie. W przypadku
zespolonym jest to identyfikacja liniowa przy przyjętej konwencji
wzoru bez sprzężenia $g$.
\end{twierdzenie}
\begin{proof}
Nierówność H\"oldera pokazuje, że dowolne $g\in L^q$
definiuje ciągły funkcjonał o normie co najwyżej $\|g\|_q$.
Dla danego $F\in(L^p)'$ definiujemy miarę zespoloną
(lub znakowaną w przypadku rzeczywistym)
$\nu(A)=F(\chi_A)$, $A\in\mathcal A$.
Funkcje charakterystyczne należą do $L^p$, ponieważ $\mu$ jest skończona.
Jeżeli $A_j$ są parami rozłączne, to
$\chi_{\bigcup_{j=1}^N A_j}\to\chi_{\bigcup_{j\ge1}A_j}$
w $L^p$ na mocy zbieżności monotonicznej; ciągłość $F$
daje przeliczalną addytywność $\nu$.
Dla dowolnego skończonego rozkładu $\Omega$ na zbiory
mierzalne $A_j$ wybierzmy liczby $\theta_j\in\K$ o module
nie większym niż $1$ tak, aby
$\theta_j\nu(A_j)=|\nu(A_j)|$. Wówczas
\[
 \sum_j|\nu(A_j)|
 =F\Bigl(\sum_j\theta_j\chi_{A_j}\Bigr)
 \le\|F\|\,\mu(\Omega)^{1/p}.
\]
Zatem $\nu$ ma skończoną wariację i jest absolutnie ciągła
względem $\mu$. Z twierdzenia Radona--Nikodyma dla miar
skończonych otrzymujemy $g\in L^1(\mu)$ takie, że
$\nu(A)=\int_Ag(x)\,\mathrm d\mu(x)$ dla $A\in\mathcal A$.
Przez liniowość i Lemat~\ref{lem:simple-dense12} trzeba jeszcze
wykazać, że $g\in L^q$.

Jeżeli $p=1$, to dla dowolnego $A\in\mathcal A$ stosujemy
nierówność normową do funkcji
$f=\chi_A\overline g/|g|$ na $\{g\ne0\}$,
a poza tym zbiorem przyjmujemy $f=0$. Z rozszerzenia wzoru
z funkcji prostych na funkcje ograniczone (przez jednostajną
aproksymację funkcjami prostymi) otrzymujemy
$\int_A|g|\,\mathrm d\mu\le\|F\|\mu(A)$.
Wybór $A=\{|g|>\|F\|+\varepsilon\}$ pokazuje,
że $\|g\|_\infty\le\|F\|$.

Jeżeli $1<p<\infty$, niech
$g_N=g\chi_{\{|g|\le N\}}$ i niech
$f_N=\overline{g_N}|g_N|^{q-2}$, z wartością $0$ w
punktach, gdzie $g_N=0$. Ponieważ $\mu$ jest skończona,
$f_N$ jest ograniczona, a więc należy do $L^p$.
Wzór reprezentacyjny obowiązuje dla niej, gdyż
jest ona granicą jednostajną funkcji prostych; wobec tego
\[
 \int_\Omega|g_N(x)|^q\,\mathrm d\mu(x)
 =F(f_N)\le\|F\|\left(\int_\Omega|g_N(x)|^q\,\mathrm d\mu(x)\right)^{1/p}.
\]
Jeśli całka po lewej stronie jest dodatnia, dzielimy przez
jej potęgę $1/p$ i otrzymujemy $\|g_N\|_q\le\|F\|$.
Zbieżność monotoniczna daje $g\in L^q$ i $\|g\|_q\le\|F\|$.
Dla $p=1$ wykazaliśmy analogiczną nierówność.
Teraz równość $F(f)=\int fg\,\mathrm d\mu$ wynika z gęstości
funkcji prostych i ciągłości obu stron na $L^p$.
Nierówność H\"oldera daje $\|F\|\le\|g\|_q$, czyli
równość norm. Jednoznaczność $g$ wynika z podstawiania
funkcji charakterystycznych zbiorów mierzalnych.
\end{proof}

\begin{uwaga}[Dlaczego ograniczenie $p<\infty$ jest istotne?]
\label{rem:linfty-dual12}
Przy $p=\infty$ przestrzeń sprzężona $(L^\infty(\mu))'$ jest
na ogół istotnie większa od $L^1(\mu)$. W przypadku miary
Lebesgue'a na $[0,1]$ wybierzmy parami rozłączne zbiory mierzalne
$A_n$ o dodatniej mierze, a $A=\bigcup_nA_n$. Operator
$P:L^\infty([0,1])\to\ell^\infty$ dany wzorem
$P(f)_n=\lambda(A_n)^{-1}\int_{A_n}f(x)\,\mathrm d\lambda(x)$
jest ciągły, liniowy i $\|P\|\le1$.
Jeżeli $\widetilde L$ jest przedłużeniem funkcjonału granicy
z Przykładu~\ref{ex:linfty-dual-seq12}, to $\Phi=\widetilde L\circ P$
spełnia $\Phi(\chi_{A_n})=0$ dla wszystkich $n$, lecz
$\Phi(\chi_A)=1$.
Gdyby $\Phi(f)=\int_{[0,1]}f(x)g(x)\,\mathrm d\lambda(x)$
dla pewnego $g\in L^1$, przeliczalna addytywność całki dałaby
$\Phi(\chi_A)=\sum_n\Phi(\chi_{A_n})=0$, sprzeczność.
\end{uwaga}

\subsection{Przestrzeń \texorpdfstring{$C(K)$}{C(K)} i twierdzenie Riesza--Markowa}

\noindent\emph{Powiązania tematyczne w Bibliotece matematycznej:} \atlasref{ME-21::thm:compact-riesz-21}{ME-21: twierdzenie Riesza--Markowa--Kakutaniego na przestrzeni zwartej}.


\begin{definicja}[Przestrzeń funkcji ciągłych i miary regularne]
\label{def:ck-measures12}
Niech $K$ będzie niepustą zwartą przestrzenią Hausdorffa.
Przez $C(K;\K)$ oznaczamy przestrzeń ciągłych funkcji
$K\to\K$ z normą $\|f\|_\infty=\max_{x\in K}|f(x)|$.
Skończoną miarę borelowską $\nu$ na $K$ nazywamy
\emph{regularną}, jeśli dla każdego borelowskiego $B\subseteq K$
i $\varepsilon>0$ istnieje zbiór domknięty $F$ i zbiór otwarty
$U$ takie, że $F\subseteq B\subseteq U$ oraz
$|\nu|(U\setminus F)<\varepsilon$, gdzie $|\nu|$ oznacza
miarę wariacji całkowitej. Zbiór takich miar ze znakiem
lub zespolonych oznaczamy przez $\mathcal M(K;\K)$.
\end{definicja}

\begin{stwierdzenie}[Zupełność $C(K)$]
\label{prop:ck-banach12}
Niech $\K\in\{\R,\C\}$. Dla każdej niepustej zwartej przestrzeni Hausdorffa $K$
przestrzeń $C(K;\K)$ jest przestrzenią Banacha.
\end{stwierdzenie}
\begin{proof}
Każda funkcja z $C(K;\K)$ jest ograniczona dzięki zwartości. Dla ciągu Cauchy'ego $(f_n)$ w normie supremum istnieje punktowa granica $f$, dzięki zupełności $\K$. Jeśli $|f_n(x)-f_m(x)|<\varepsilon/2$ dla $n,m\ge N$ i każdego $x$, to przechodząc do granicy przy $m\to\infty$ dostajemy $\|f_n-f\|_\infty\le\varepsilon/2$ dla $n\ge N$. Dla ustalonego $x$ wybieramy $n$ z $\|f_n-f\|_\infty<\varepsilon/3$, a następnie otoczenie $x$, na którym $|f_n(y)-f_n(x)|<\varepsilon/3$. Nierówność trójkąta daje $|f(y)-f(x)|<\varepsilon$. Stąd $f\in C(K;\K)$ i zbieżność w normie supremum.
\end{proof}

\begin{lemat}[Reprezentacja dodatniego funkcjonału]
\label{lem:riesz-positive12}
Niech $K$ będzie niepustą zwartą przestrzenią metryczną,
a $T:C(K;\R)\to\R$ dodatnim funkcjonałem liniowym,
tj. $T(f)\ge0$ dla każdej $f\ge0$.
Wówczas istnieje dokładnie jedna skończona regularna dodatnia miara
borelowska $\mu_T$ na $K$, dla której
$T(f)=\int_Kf(x)\,\mathrm d\mu_T(x)$ dla wszystkich $f\in C(K;\R)$.
\end{lemat}
\begin{proof}
Jest to dokładnie przypadek dodatniego funkcjonału na niepustej
zwartej przestrzeni metrycznej z twierdzenia
\atlasref{ME-21::thm:compact-riesz-21}{ME-21: twierdzenie Riesza--Markowa--Kakutaniego}.
W tym miejscu korzystamy z pełnego dowodu zawartego w teorii miary,
łącznie z konstrukcją miary, jej regularnością i jednoznacznością.
\end{proof}

\begin{twierdzenie}[Twierdzenie Riesza--Markowa na zwartej przestrzeni metrycznej]
\label{thm:riesz-markov12}
Niech $K$ będzie niepustą zwartą przestrzenią metryczną.
Dla każdego $F\in C(K;\K)'$ istnieje dokładnie jedna
regularna skończona miara borelowska $\nu\in\mathcal M(K;\K)$
taka, że
\[
 F(f)=\int_K f(x)\,\mathrm d\nu(x)
 \quad\text{dla wszystkich }f\in C(K;\K).
\]
Odwzorowanie $\nu\mapsto F$ jest izometrycznym izomorfizmem
$\mathcal M(K;\K)\cong C(K;\K)'$, przy normie
$\|\nu\|=|\nu|(K)$.
\end{twierdzenie}
\begin{proof}
Najpierw $\K=\R$. Dla $f\ge0$ definiujemy
$F^+(f)=\sup\{F(g):g\in C(K;\R),\ 0\le g\le f\}$.
Wartość supremum należy do $[0,\|F\|\|f\|_\infty]$.
Funkcja $F^+$ jest addytywna na stożku dodatnim.
Jedna nierówność wynika z wyboru $g_1\le f_1$ i
$g_2\le f_2$. Dla przeciwnej, jeśli
$0\le h\le f_1+f_2$, rozkładamy
$h=\min\{h,f_1\}+(h-\min\{h,f_1\})$;
oba składniki są ciągłe, nieujemne i ograniczone odpowiednio
przez $f_1$ i $f_2$. Po przejściu do supremum mamy
$F^+(f_1+f_2)=F^+(f_1)+F^+(f_2)$.
Rozszerzamy zatem $F^+$ na rzeczywisty funkcjonał liniowy
przez $F^+(f)=F^+(f_+)-F^+(f_-)$, gdzie
$f_+=\max\{f,0\}$, $f_-=\max\{-f,0\}$.
Jest on dodatni i ciągły; również $F^-=F^+-F$ jest dodatni,
ponieważ $F^+(f)\ge F(f)$ dla $f\ge0$.
Lemat~\ref{lem:riesz-positive12} daje dodatnie regularne miary
$\nu^+$ i $\nu^-$ reprezentujące odpowiednio $F^+$ i $F^-$.
Zatem $F(f)=\int_Kf\,\mathrm d(\nu^+-\nu^-)$.

Dla $\K=\C$ bierzemy części rzeczywistą i urojoną
ograniczenia $F$ do rzeczywistych funkcji ciągłych,
reprezentujemy je dwiema miarami rzeczywistymi i składamy
$\nu=\nu_1+i\nu_2$. Zespolona liniowość $F$ zapewnia,
że ten sam wzór działa na wszystkich $f\in C(K;\C)$.
Jednoznaczność wynika z jednoznaczności reprezentacji
części rzeczywistych i urojonych.

Dla dowolnej miary o skończonej wariacji mamy
$|\int f\,\mathrm d\nu|\le\|f\|_\infty|\nu|(K)$,
czyli $\|F\|\le|\nu|(K)$. W drugą stronę niech
$u=d\nu/d|\nu|$ będzie pochodną Radona--Nikodyma;
$|u|=1$ $|\nu|$-prawie wszędzie. Regularność miary
$|\nu|$ i własność aproksymacji funkcji charakterystycznych zbiorów
borelowskich funkcjami ciągłymi (przez zbiory domknięte
i otwarte oraz funkcje Urysohna) dają gęstość $C(K;\K)$
w $L^1(|\nu|)$. Aproksymując $\overline u$ funkcjami
ciągłymi o module nie większym niż $1$, otrzymujemy
$\|F\|\ge\lim|\int f_nu\,\mathrm d|\nu||=|\nu|(K)$.
To dowodzi izometrii.
\end{proof}

\begin{przyklad}[Miary Diraca i funkcjonały ewaluacji]
\label{ex:dirac12}
Dla $x_0\in K$ miara Diraca $\delta_{x_0}$ jest zdefiniowana
przez $\delta_{x_0}(B)=1$, gdy $x_0\in B$, i $0$ w przeciwnym
przypadku. Funkcjonał $F_{x_0}(f)=f(x_0)$ spełnia
$F_{x_0}(f)=\int_Kf\,\mathrm d\delta_{x_0}$ i ma normę $1$.
Jeśli $x_0\ne x_1$, funkcjonały ewaluacji są różne,
ponieważ metryka $K$ pozwala skonstruować funkcję ciągłą
rozdzielającą $x_0$ i $x_1$.
\end{przyklad}

\begin{uwaga}[Przygotowanie do przestrzeni dystrybucji]
\label{rem:bridge-distribution12}
Miara $\nu$ działa liniowo na funkcjach testowych przez całkowanie,
a miara Diraca działa przez ewaluację. Te same mechanizmy wystąpią
w przestrzeni dystrybucji $\mathcal D'(\Omega)$, jednak tam
przestrzeń funkcji testowych ma topologię granicy induktywnej,
a nie normę supremum. W szczególności nie wolno bez dodatkowych
argumentów utożsamiać wszystkich dystrybucji z miarami.
\end{uwaga}

% Wersja źródłowa: audyt II, jednostka 13
\setcounter{section}{12}
\section{Przestrzenie Hilberta: geometria, reprezentacja Riesza i formy półtoraliniowe}
\label{sec:hilbert13}

\subsection{Iloczyn skalarny i zupełność}

\begin{definicja}[Iloczyn skalarny i przestrzeń Hilberta]
\label{def:iloczyn-skalarny13}
Niech $H$ będzie przestrzenią wektorową nad $\K\in\{\R,\C\}$.
\emph{Iloczynem skalarnym} na $H$ nazywamy odwzorowanie
$\langle\cdot,\cdot\rangle:H\times H\to\K$ spełniające, dla
$x,y,z\in H$ i $\alpha,\beta\in\K$:
$\langle\alpha x+\beta y,z\rangle=\alpha\langle x,z\rangle+
\beta\langle y,z\rangle$,
$\langle x,y\rangle=\overline{\langle y,x\rangle}$ oraz
$\langle x,x\rangle>0$ dla $x\ne0$.
\emph{W całym rozdziale iloczyn skalarny jest liniowy w pierwszym
argumencie} i antyliniowy w drugim. Przestrzenią Hilberta nazywamy
przestrzeń z iloczynem skalarnym zupełną w normie
$\|x\|=\sqrt{\langle x,x\rangle}$.
\end{definicja}

\begin{lemat}[Nierówność Cauchy'ego--Schwarza]
\label{lem:cs13}
Dla dowolnych $x,y$ w przestrzeni z iloczynem skalarnym zachodzi
$|\langle x,y\rangle|\le\|x\|\|y\|$.
Równość dla $y\ne0$ zachodzi dokładnie wtedy, gdy $x$ jest
wielokrotnością $y$.
\end{lemat}
\begin{proof}
Dla $y\ne0$ przyjmijmy $\alpha=\langle x,y\rangle/\|y\|^2$.
Wówczas $\langle x-\alpha y,y\rangle=0$ oraz
$0\le\|x-\alpha y\|^2=\|x\|^2-|\langle x,y\rangle|^2/\|y\|^2$.
Równość oznacza $x=\alpha y$. Dla $y=0$ teza jest oczywista.
Nierówność trójkąta dla normy wynika teraz z rozwinięcia
$\|x+y\|^2$ i nierówności Cauchy'ego--Schwarza.
\end{proof}

\begin{stwierdzenie}[Tożsamość równoległoboku i polaryzacja]
\label{prop:parallelogram13}
W każdej przestrzeni z iloczynem skalarnym zachodzi
$\|x+y\|^2+\|x-y\|^2=2\|x\|^2+2\|y\|^2$.
Dla $\K=\C$ iloczyn skalarny jest odtworzony z normy wzorem
\[
\langle x,y\rangle=\frac14\bigl(\|x+y\|^2-\|x-y\|^2
+i\|x+iy\|^2-i\|x-iy\|^2\bigr).
\]
Dla $\K=\R$ wystarczą pierwsze dwa wyrazy.
\end{stwierdzenie}
\begin{proof}
Rozwijamy kwadraty norm. Przy naszej konwencji
$\|x+iy\|^2-\|x-iy\|^2=4\operatorname{Im}\langle x,y\rangle$;
różnica dla $x\pm y$ to $4\operatorname{Re}\langle x,y\rangle$.
\end{proof}

\begin{przyklad}[Przestrzenie $\ell^2$ i $L^2$]
\label{ex:hilbert-l2-13}
Na $\ell^2$ przyjmujemy $\langle x,y\rangle=\sum_{n\ge1}x_n\overline{y_n}$.
Na $L^2(\Omega,\mu)$ przyjmujemy
$\langle f,g\rangle=\int_\Omega f(x)\overline{g(x)}\,\mathrm d\mu(x)$.
Nierówność Cauchy'ego--Schwarza zapewnia zbieżność i ciągłość tych
wyrażeń, zaś zupełność wynika z
twierdzeń o zupełności przestrzeni $\ell^2$ i $L^2$
(drugi przykład przy założeniach o mierze z \hyperref[sec:klasyczne-banach12]{rozdziału~12}).
Przestrzeń $C([0,1])$ z normą supremum nie jest przestrzenią Hilberta w tej normie:
dla $f(t)=1$ i $g(t)=t$ mamy
$\|f+g\|_\infty^2+\|f-g\|_\infty^2=5$,
podczas gdy $2\|f\|_\infty^2+2\|g\|_\infty^2=4$.
\end{przyklad}

\subsection{Rzut na zbiór wypukły i ortogonalność}

\begin{twierdzenie}[Twierdzenie o rzucie metrycznym]
\label{thm:projection-convex13}
Niech $H$ będzie rzeczywistą lub zespoloną przestrzenią Hilberta,
a $C\subseteq H$ niepustym, domkniętym i wypukłym zbiorem.
Dla każdego $x\in H$ istnieje dokładnie jeden punkt $p=P_Cx\in C$
spełniający $\|x-p\|=\inf_{z\in C}\|x-z\|$.
Punkt $p\in C$ jest tym minimizatorem wtedy i tylko wtedy, gdy
\[
\operatorname{Re}\langle x-p,z-p\rangle\le0\qquad(z\in C).
\]
Ponadto odwzorowanie $P_C:H\to C$ jest $1$-lipschitzowskie.
\end{twierdzenie}
\begin{proof}
Niech $d=\inf_{z\in C}\|x-z\|$ i wybierzmy $z_n\in C$ z
$\|x-z_n\|^2\to d^2$. Ponieważ $(z_n+z_m)/2\in C$,
tożsamość równoległoboku zastosowana do $x-z_n$ i $x-z_m$ daje
\[
\|z_n-z_m\|^2
=2\|x-z_n\|^2+2\|x-z_m\|^2
-4\left\|x-\frac{z_n+z_m}{2}\right\|^2
\le2\|x-z_n\|^2+2\|x-z_m\|^2-4d^2\longrightarrow0.
\]
Zatem $(z_n)$ jest ciągiem Cauchy'ego. Zupełność $H$ i domkniętość
$C$ dają granicę $p\in C$, a ciągłość normy daje $\|x-p\|=d$.
Jeśli $p,q$ są minimizatorami, ten sam wzór daje $p=q$.

Niech $z\in C$ i $t\in[0,1]$. Minimalność $p$ oznacza
$\|x-p\|^2\le\|x-(p+t(z-p))\|^2$,
a więc
$0\le-2t\operatorname{Re}\langle x-p,z-p\rangle+t^2\|z-p\|^2$.
Dzieląc przez $t>0$ i przechodząc z $t$ do zera otrzymujemy
warunek wariacyjny. Odwrotnie, warunek ten i rozwinięcie
$\|x-z\|^2=\|x-p\|^2+\|z-p\|^2-
2\operatorname{Re}\langle x-p,z-p\rangle$ implikują minimalność.

Dla $p=P_Cx$, $q=P_Cy$ zastosujmy warunek kolejno z $z=q$
i $z=p$. Otrzymujemy
$\operatorname{Re}\langle x-p,q-p\rangle\le0$ oraz
$\operatorname{Re}\langle y-q,p-q\rangle\le0$. Po dodaniu
$\|p-q\|^2\le\operatorname{Re}\langle x-y,p-q\rangle
\le\|x-y\|\|p-q\|$, skąd $\|p-q\|\le\|x-y\|$.
\end{proof}

\begin{definicja}[Dopełnienie ortogonalne]
\label{def:orth-complement13}
Dla $M\subseteq H$ definiujemy
$M^\perp=\{x\in H:\langle x,m\rangle=0\text{ dla wszystkich }m\in M\}$.
Jest to domknięta podprzestrzeń liniowa $H$; istotnie, to
przecięcie jąder ciągłych funkcjonałów $x\mapsto\langle x,m\rangle$.
\end{definicja}

\begin{twierdzenie}[Rozkład ortogonalny]
\label{thm:orth-decomposition13}
Niech $H$ będzie przestrzenią Hilberta, a $M\subseteq H$ domkniętą
podprzestrzenią liniową. Wówczas dla każdego $x\in H$ istnieje
jedyny rozkład $x=m+n$, $m\in M$, $n\in M^\perp$.
Odwzorowanie $P_Mx=m$ jest liniowym operatorem o normie co najwyżej
$1$, spełnia $P_M^2=P_M=P_M^*$ oraz
$\|x\|^2=\|P_Mx\|^2+\|x-P_Mx\|^2$.
Ponadto $M^{\perp\perp}=M$.
Jeśli $M\ne\{0\}$, to $\|P_M\|=1$.
\end{twierdzenie}
\begin{proof}
Z Twierdzenia~\ref{thm:projection-convex13} istnieje $m=P_Mx$.
Warunek wariacyjny zastosowany do $m+tv$ dla każdego $v\in M$
i $t\in\R$, a dla przestrzeni zespolonej także do $m+itv$,
daje $\langle x-m,v\rangle=0$. Stąd $n=x-m\in M^\perp$.
Jeśli $x=m_1+n_1=m_2+n_2$, to $m_1-m_2=n_2-n_1$ należy
do $M\cap M^\perp=\{0\}$. Jednoznaczność rozkładu zapewnia
liniowość $P_M$; twierdzenie Pitagorasa daje jego ograniczoność.

Dla $x=m+n$ i $y=u+v$ mamy
$\langle P_Mx,y\rangle=\langle m,u\rangle
=\langle x,P_My\rangle$; to oznacza $P_M=P_M^*$,
przy czym istnienie operatora sprzężonego wykażemy formalnie
w Twierdzeniu~\ref{thm:adjoint13} poniżej.
Z rozkładu $H=M\oplus M^\perp$ wynika, że
$M^{\perp\perp}=M$. Dla $m\in M\setminus\{0\}$ zachodzi
$P_Mm=m$, zatem norma operatorowa nie może być mniejsza niż $1$.
\end{proof}

\begin{przyklad}[Rzut na funkcje stałe]
\label{ex:projection-mean13}
W przestrzeni $H=L^2([0,1],\lambda)$ niech
$M=\{\text{klasy funkcji stałych}\}$. Wtedy
$P_Mf=(\int_{[0,1]}f(t)\,\mathrm d\lambda(t))\chi_{[0,1]}$.
Rzeczywiście, $\langle f-P_Mf,\chi_{[0,1]}\rangle=0$.
Ogólniej, rzut na linię $\K u$, $u\ne0$, ma wzór
$P_{\K u}x=\langle x,u\rangle u/\|u\|^2$.
\end{przyklad}

\subsection{Twierdzenie Riesza o reprezentacji funkcjonałów}

\begin{twierdzenie}[Twierdzenie Riesza dla przestrzeni Hilberta]
\label{thm:riesz-hilbert13}
Niech $H$ będzie przestrzenią Hilberta nad $\K\in\{\R,\C\}$.
Dla każdego ciągłego $\K$-liniowego funkcjonału $F:H\to\K$
istnieje dokładnie jeden $y\in H$ taki, że
$F(x)=\langle x,y\rangle$ dla wszystkich $x\in H$.
Ponadto $\|F\|_{H'}=\|y\|$.
Odwzorowanie $J:H\to H'$, $J(y)(x)=\langle x,y\rangle$,
jest izometrycznym izomorfizmem \emph{antyliniowym} dla $\K=\C$
i liniowym dla $\K=\R$.
\end{twierdzenie}
\begin{proof}
Jeżeli $F=0$, wybieramy $y=0$. Załóżmy $F\ne0$.
Podprzestrzeń $M=\ker F$ jest domknięta, a $M^\perp\ne\{0\}$:
gdyby $M^\perp=\{0\}$, to z Twierdzenia~\ref{thm:orth-decomposition13}
wynikałoby $M=H$.
Niech $z\in M^\perp$, $z\ne0$. Wtedy $F(z)\ne0$, gdyż
$z\in M\cap M^\perp$ implikowałoby $z=0$.
Dla dowolnego $x\in H$ mamy
$x-F(x)z/F(z)\in M$, więc
\[
0=\left\langle x-\frac{F(x)}{F(z)}z,z\right\rangle
=\langle x,z\rangle-\frac{F(x)}{F(z)}\|z\|^2.
\]
Zatem $F(x)=F(z)\langle x,z\rangle/\|z\|^2$.
Przy naszej konwencji należy wziąć
$y=\overline{F(z)}z/\|z\|^2$, ponieważ
$\langle x,cy\rangle=\overline c\langle x,y\rangle$.
Nierówność Cauchy'ego--Schwarza daje $\|F\|\le\|y\|$,
a podstawienie $x=y/\|y\|$ dla $y\ne0$ daje równość.
Jeśli $\langle x,y_1-y_2\rangle=0$ dla wszystkich $x$, to
$x=y_1-y_2$ zapewnia $y_1=y_2$. Antyliniowość $J$
wynika bezpośrednio z $J(\alpha y+\beta z)
=\overline\alpha J(y)+\overline\beta J(z)$, ponieważ
iloczyn skalarny jest antyliniowy w drugim argumencie.
\end{proof}

\begin{uwaga}[Dwa różne twierdzenia Riesza]
\label{rem:two-riesz13}
Twierdzenie~\ref{thm:riesz-hilbert13} identyfikuje przestrzeń sprzężoną
przestrzeni Hilberta z nią samą (antyliniowo w przypadku zespolonym).
Natomiast twierdzenie Riesza--Markowa z teorii przestrzeni $C(K)$ identyfikuje
przestrzeń sprzężoną $C(K)$ z przestrzenią miar. Nie są to
odmienne wersje tej samej reprezentacji w sensie konstrukcyjnym.
\end{uwaga}

\subsection{Rodziny ortonormalne i bazy przestrzeni Hilberta}

\begin{definicja}[Rodzina ortonormalna i baza przestrzeni Hilberta]
\label{def:on-basis13}
Niech $H$ będzie przestrzenią Hilberta, a $I$ zbiorem indeksów.
Rodzina $(e_i)_{i\in I}\subseteq H$ jest \emph{ortonormalna}, jeśli
$\langle e_i,e_j\rangle=\delta_{ij}$.
Nazywamy ją \emph{bazą ortonormalną przestrzeni Hilberta}, jeśli
$\overline{\lin\{e_i:i\in I\}}=H$.
Nie należy utożsamiać bazy ortonormalnej przestrzeni Hilberta z algebraiczną bazą Hamela.
\end{definicja}

\begin{lemat}[Bessel i ortogonalne sumy skończone]
\label{lem:bessel13}
Niech $H$ będzie przestrzenią Hilberta, $(e_i)_{i\in I}$ rodziną
ortonormalną, a $x\in H$. Dla każdego skończonego $A\subseteq I$
zachodzi
\[
\left\|x-\sum_{i\in A}\langle x,e_i\rangle e_i\right\|^2
=\|x\|^2-\sum_{i\in A}|\langle x,e_i\rangle|^2.
\]
W szczególności $\sum_{i\in A}|\langle x,e_i\rangle|^2\le\|x\|^2$
dla każdego takiego $A$ (nierówność Bessela).
\end{lemat}
\begin{proof}
Wektor $x-\sum_{i\in A}\langle x,e_i\rangle e_i$ jest
prostopadły do każdego $e_j$, $j\in A$. Z twierdzenia Pitagorasa
dla jego rozkładu z sumą ortogonalną wynika wzór.
\end{proof}


\begin{twierdzenie}[Przeliczalność nośnika współczynników Fouriera]
\label{thm:countable-fourier-support13}
Niech $H$ będzie przestrzenią Hilberta nad $\K\in\{\R,\C\}$,
$(e_i)_{i\in I}$ dowolną rodziną ortonormalną (zbiór indeksów $I$
może być nieprzeliczalny), a $x\in H$. Wówczas zbiór
$ I_x:=\{i\in I:\langle x,e_i\rangle\ne0\}$
jest co najwyżej przeliczalny. Dokładniej, dla każdego $n\in\N$,
$n\ge1$, zbiór
$ I_{x,n}:=\{i\in I:|\langle x,e_i\rangle|\ge1/n\}$
jest skończony i ma co najwyżej $n^2\|x\|^2$ elementów.
\end{twierdzenie}
\begin{proof}
Niech $F\subseteq I_{x,n}$ będzie zbiorem skończonym.
Z nierówności Bessela (Lemat~\ref{lem:bessel13}) otrzymujemy
$|F|/n^2\le\sum_{i\in F}|\langle x,e_i\rangle|^2\le\|x\|^2$.
Gdyby $I_{x,n}$ miał więcej niż $n^2\|x\|^2$ elementów,
można byłoby wybrać skończony podzbiór o liczności
przekraczającej tę liczbę, co przeczyłoby otrzymanemu oszacowaniu.
Każdy $I_{x,n}$ jest zatem skończony. Dla każdego
$i\in I_x$ liczba $|\langle x,e_i\rangle|$ jest dodatnia, więc
przekracza $1/n$ dla pewnego $n\ge1$. Stąd
$I_x=\bigcup_{n\ge1}I_{x,n}$ jest sumą przeliczalnej rodziny
zbiorów skończonych.
\end{proof}

\begin{uwaga}[Dlaczego warunek jest silniejszy od samego Bessela?]
\label{rem:fourier-countability13}
Nie zakładaliśmy ośrodkowości $H$ ani przeliczalności $I$.
Twierdzenie dotyczy \emph{każdego pojedynczego wektora}:
przy nieprzeliczalnej bazie ortonormalnej różne wektory mogą
mieć różne przeliczalne nośniki współczynników. Nie wynika
stąd, że całą bazę można ponumerować jednym ciągiem.
\end{uwaga}

\begin{twierdzenie}[Rozwinięcie Fouriera, tożsamość Parsevala i zbieżność bezwarunkowa]
\label{thm:parseval13}
\label{cor:fourier-enumeration13}
Niech $H$ będzie przestrzenią Hilberta nad $\K\in\{\R,\C\}$,
$(e_i)_{i\in I}$ rodziną ortonormalną, a
$M=\overline{\lin\{e_i:i\in I\}}$. Dla każdego $x\in H$
sieć sum $s_A(x)=\sum_{i\in A}\langle x,e_i\rangle e_i$,
indeksowana skończonymi podzbiorami $A\subseteq I$ uporządkowanymi
przez inkluzję, zbiega w normie do $P_Mx$. Ponadto
\[
 \|P_Mx\|^2=\sum_{i\in I}|\langle x,e_i\rangle|^2,
 \qquad
 \|x\|^2=\|x-P_Mx\|^2+
               \sum_{i\in I}|\langle x,e_i\rangle|^2,
\]
gdzie suma po dowolnym zbiorze indeksów oznacza supremum
sum skończonych. Rodzina $(e_i)_{i\in I}$ jest bazą ortonormalną
przestrzeni Hilberta wtedy i tylko wtedy, gdy dla każdego $x\in H$
zachodzi tożsamość Parsevala $\|x\|^2=\sum_{i\in I}
|\langle x,e_i\rangle|^2$; równoważnie, gdy
$M^\perp=\{0\}$.

Jeśli $(e_i)_{i\in I}$ jest bazą ortonormalną, to zbiór
$I_x=\{i:\langle x,e_i\rangle\ne0\}$ jest co najwyżej
przeliczalny. Gdy $I_x$ jest nieskończony, dla każdej bijekcji
$\nu:\N\to I_x$ zachodzi
\[
 x=\sum_{n=1}^{\infty}\langle x,e_{\nu(n)}\rangle e_{\nu(n)}
 \quad\text{w normie $H$},\qquad
 \|x\|^2=\sum_{n=1}^{\infty}|\langle x,e_{\nu(n)}\rangle|^2.
\]
Suma nie zależy od numeracji $I_x$ (zbieżność bezwarunkowa).
Jeśli $I_x$ jest skończony, obowiązuje analogiczna suma skończona.
Dla każdego skończonego $A\subseteq I$ mamy dokładny wzór na błąd
\[
 \left\|x-\sum_{i\in A}\langle x,e_i\rangle e_i\right\|^2
 =\sum_{i\in I\setminus A}|\langle x,e_i\rangle|^2.
\]
\end{twierdzenie}
\begin{proof}
Dla skończonych $A\subseteq B\subseteq I$ ortonormalność daje
$\|s_B-s_A\|^2=\sum_{i\in B\setminus A}|\langle x,e_i\rangle|^2$.
Niech $S=\sup_A\sum_{i\in A}|\langle x,e_i\rangle|^2
\le\|x\|^2$ (nierówność Bessela). Dla $\varepsilon>0$ wybieramy
$A_0$ o sumie większej niż $S-\varepsilon^2$.
Dla $A,B\supseteq A_0$ porównanie z $A\cup B$ daje
$\|s_A-s_B\|<2\varepsilon$, zatem $(s_A)$ jest siecią
Cauchy'ego. Z zupełności $H$ wynika jej zbieżność do pewnego
$s\in M$. Dla każdego $j\in I$ od chwili, gdy $j\in A$,
$\langle x-s_A,e_j\rangle=0$. Po przejściu do granicy
$x-s\in M^\perp$, więc $s=P_Mx$. Przejście do granicy
w $\|s_A\|^2=\sum_{i\in A}|\langle x,e_i\rangle|^2$
oraz twierdzenie Pitagorasa dają obie równości.
Charakterystyka bazy wynika z równoważności $M^\perp=\{0\}$
i $M=H$.

Przeliczalność $I_x$ została już niezależnie udowodniona
w Twierdzeniu~\ref{thm:countable-fourier-support13}.
Jeśli $I_x$ jest nieskończony, to zbiory
$A_N=\{\nu(1),\ldots,\nu(N)\}$ są współkońcowe w rodzinie
skończonych podzbiorów $I_x$: każdy skończony
$A\subseteq I_x$ zawiera się w pewnym $A_N$. Współczynniki
spoza $I_x$ są zerowe, więc sieć sum po podzbiorach $I$ i
sieć sum po podzbiorach $I_x$ mają tę samą granicę.
Stąd $s_{A_N}\to x$ (dla bazy $M=H$), niezależnie od $\nu$.
Suma kwadratów wynika z Parsevala. Wreszcie wektor
$x-s_A(x)$ jest prostopadły do $s_A(x)$, więc
$\|x-s_A(x)\|^2=\|x\|^2-\sum_{i\in A}|\langle x,e_i\rangle|^2$;
ponowne użycie Parsevala daje wzór na błąd.
\end{proof}

\begin{przyklad}[Różnica między rozpięciem algebraicznym i domkniętym]
\label{ex:hilbert-series-not-hamel13}
Niech $(e_n)_{n\ge1}$ będzie ciągiem ortonormalnym w $H$.
Szereg $x=\sum_{n\ge1}2^{-n}e_n$ zbiega w $H$, ponieważ
$\sum_{n\ge1}4^{-n}<\infty$. Dla każdego $n$ mamy
$\langle x,e_n\rangle=2^{-n}\ne0$, więc
$x\notin\lin\{e_n:n\ge1\}$, choć należy do domknięcia
tej podprzestrzeni. Baza ortonormalna przestrzeni Hilberta pozwala zatem na rozwinięcia
\emph{normowo zbieżne}, a baza Hamela dopuszcza wyłącznie
\emph{skończone} kombinacje liniowe.
\end{przyklad}

\begin{twierdzenie}[Istnienie baz ortonormalnych]
\label{thm:hilbert-basis13}
Każda przestrzeń Hilberta ma bazę ortonormalną (dopuszczając pustą
bazę dla $H=\{0\}$). Przestrzeń Hilberta jest ośrodkowa wtedy i tylko
wtedy, gdy posiada bazę ortonormalną co najwyżej przeliczalną.
\end{twierdzenie}
\begin{proof}
W zbiorze rodzin ortonormalnych, uporządkowanym przez rozszerzenie,
każdy łańcuch ma ograniczenie górne dane sumą rodzin.
Lemat Zorna zapewnia rodzinę maksymalną $(e_i)_{i\in I}$.
Jeśli $x\perp e_i$ dla każdego $i$ i $x\ne0$, można ją powiększyć
o $x/\|x\|$, sprzeczność. Zatem jest bazą.

Jeśli $H$ ma przeliczalną bazę, skończone kombinacje liniowe z
współczynnikami w $\mathbb Q$ (dla $\R$) lub
$\mathbb Q+i\mathbb Q$ (dla $\C$) tworzą zbiór przeliczalny i gęsty.
Odwrotnie, jeśli $H$ jest ośrodkowa, nie może zawierać
nieprzeliczalnej rodziny ortonormalnej: kule
$B(e_i,1/3)$ są rozłączne, a każdy taki otwarty zbiór zawiera
punkt ustalonego przeliczalnego zbioru gęstego.
\end{proof}


\begin{twierdzenie}[Nieskończenie wymiarowe przestrzenie Hilberta a wymiar Hamela]
\label{thm:hilbert-hamel-uncountable13}
Niech $H$ będzie nieskończenie wymiarową przestrzenią Hilberta
nad $\K\in\{\R,\C\}$. Wtedy każda baza Hamela przestrzeni $H$
jest nieprzeliczalna. Co więcej, $H$ zawiera liniowo niezależną
rodzinę mocy continuum, zatem
$\dim_{\mathrm{Hamel}}H\ge 2^{\aleph_0}$.
W szczególności twierdzenie dotyczy również ośrodkowych,
nieskończenie wymiarowych przestrzeni Hilberta, które mają
przeliczalne bazy ortonormalne.
\end{twierdzenie}
\begin{proof}
Wybierzmy nieskończony ciąg ortonormalny $(e_n)_{n\ge1}$
w $H$; istnieje on, ponieważ $H$ jest nieskończenie wymiarowa
(ewentualnie wybieramy go z bazy ortonormalnej otrzymanej
w Twierdzeniu~\ref{thm:hilbert-basis13}). Dla $t\in(0,1)$
określmy
\[
 v_t:=\sum_{n=1}^{\infty}t^{n-1}e_n.
\]
Szereg ten zbiega w normie dzięki twierdzeniu Pitagorasa i
$\sum_{n\ge1}t^{2(n-1)}=(1-t^2)^{-1}<\infty$.
Pokażemy, że rodzina $(v_t)_{t\in(0,1)}$ jest liniowo niezależna.
Niech $t_1,\ldots,t_m\in(0,1)$ będą parami różne, a
$\sum_{j=1}^m a_jv_{t_j}=0$. Biorąc iloczyn skalarny tej
równości z $e_n$ dla $n=1,\ldots,m$, dostajemy
$\sum_{j=1}^m a_jt_j^{n-1}=0$. Współczynniki $a_j$
spełniają jednorodny układ o macierzy Vandermonde'a
$(t_j^{n-1})_{1\le n,j\le m}$. Jej wyznacznik
$\prod_{1\le j<k\le m}(t_k-t_j)$ jest niezerowy, więc
$a_1=\cdots=a_m=0$. Rodzina $(v_t)_{t\in(0,1)}$ ma moc
$2^{\aleph_0}$ i jest liniowo niezależna. Każda baza Hamela
musi mieć moc co najmniej tak dużą.
\end{proof}

\begin{uwaga}[Drugi argument: twierdzenie Baire'a]
\label{rem:hilbert-hamel-baire13}
Sama nieprzeliczalność bazy Hamela wynika nawet dla każdej
nieskończenie wymiarowej przestrzeni \emph{Banacha} $X$.
Gdyby $X$ miała przeliczalną bazę Hamela $(h_n)_{n\ge1}$,
to $X=\bigcup_{n\ge1}F_n$, gdzie
$F_n=\lin\{h_1,\ldots,h_n\}$. Każda skończenie wymiarowa
podprzestrzeń $F_n$ jest domknięta (w przestrzeni unormowanej)
i ma puste wnętrze: podprzestrzeń liniowa o niepustym wnętrzu
jest całą przestrzenią, a $X$ nie jest skończenie wymiarowa.
Twierdzenie Baire'a wyklucza pokrycie zupełnej przestrzeni
metrycznej przeliczalną sumą takich zbiorów.
Argument ten dowodzi nieprzeliczalności, natomiast poprzednia
konstrukcja daje również dolne oszacowanie mocy bazy Hamela.
\end{uwaga}

\begin{uwaga}[Dwie wielkości, których nie należy mylić]
\label{rem:hamel-vs-hilbert13}
\emph{Wymiar hilbertowski} to moc bazy ortonormalnej,
a \emph{wymiar Hamela} to moc bazy algebraicznej.
Dla $H=\ell^2(\N)$ pierwszy jest przeliczalny, drugi
co najmniej continuum. Jeśli baza ortonormalna jest
nieprzeliczalna, każdy $x\in H$ nadal ma najwyżej
przeliczalnie wiele niezerowych współczynników.
\end{uwaga}

\begin{stwierdzenie}[Izometryczne współrzędne w $\ell^2(I)$]
\label{prop:hilbert-l2I13}
Niech $(e_i)_{i\in I}$ będzie bazą ortonormalną przestrzeni
Hilberta $H$. Wówczas $U:H\to\ell^2(I)$,
$Ux=(\langle x,e_i\rangle)_{i\in I}$, jest izometrycznym
izomorfizmem liniowym, gdzie
$\ell^2(I)=\{a:I\to\K:\sup_{A\subseteq I\text{ skończony}}
\sum_{i\in A}|a_i|^2<\infty\}$.
\end{stwierdzenie}
\begin{proof}
Liniowość i izometria wynikają z Parsevala. Dla dowolnego
$a\in\ell^2(I)$ sieć $\sum_{i\in A}a_ie_i$ jest Cauchy'ego:
tak jak w dowodzie Twierdzenia~\ref{thm:parseval13},
odrzucenie dostatecznie dużego skończonego zbioru współrzędnych
daje dowolnie mały ogon sumy kwadratów. Niech $x$ będzie granicą;
wtedy $\langle x,e_j\rangle=a_j$ dla każdego $j\in I$.
To dowodzi surjektywności.
\end{proof}

\begin{przyklad}[Gram--Schmidt i szeregi Fouriera]
\label{ex:gram-schmidt13}
Dla liniowo niezależnych $v_1,\ldots,v_n\in H$ konstrukcja
$u_k=v_k-\sum_{j<k}\langle v_k,e_j\rangle e_j$,
$e_k=u_k/\|u_k\|$, tworzy rodzinę ortonormalną o tym samym
rozpięciu liniowym. Dla nieskończonego ciągu liniowo niezależnego
ortogonalizujemy kolejne wyrazy i otrzymujemy bazę ortonormalną domkniętego rozpięcia tych wektorów. Dla $H=L^2([0,2\pi],\lambda)$
układ $e_n(t)=(2\pi)^{-1/2}e^{int}$, $n\in\mathbb Z$, jest
ortonormalny; jego zupełność wymaga dodatkowego twierdzenia
o gęstości wielomianów trygonometrycznych, które zostanie
omówione przy transformacji Fouriera. Nie używamy tej zupełności
bez dowodu.
\end{przyklad}

\subsection{Operatory sprzężone i reprezentacja form}

\begin{twierdzenie}[Operator sprzężony w przestrzeni Hilberta]
\label{thm:adjoint13}
Niech $H,G$ będą przestrzeniami Hilberta nad tym samym ciałem $\K$,
a $T\in\mathcal L(H,G)$. Istnieje dokładnie jeden operator
$T^*\in\mathcal L(G,H)$ taki, że
$\langle Tx,y\rangle_G=\langle x,T^*y\rangle_H$
dla wszystkich $x\in H$, $y\in G$.
Ponadto $\|T^*\|=\|T\|$, $(T^*)^*=T$ i
$(ST)^*=T^*S^*$ dla każdego $S\in\mathcal L(G,W)$,
gdzie $W$ jest kolejną przestrzenią Hilberta.
\end{twierdzenie}
\begin{proof}
Dla ustalonego $y\in G$ funkcja $x\mapsto\langle Tx,y\rangle_G$
jest ciągłym liniowym funkcjonałem na $H$. Z
Twierdzenia~\ref{thm:riesz-hilbert13} istnieje jedyny $z_y\in H$
z $\langle Tx,y\rangle_G=\langle x,z_y\rangle_H$.
Definiujemy $T^*y=z_y$. Z antyliniowości w drugim argumencie
po obu stronach, z uwzględnieniem sprzężenia skalarów, wynika
$T^*(\alpha y_1+\beta y_2)=\alpha T^*y_1+\beta T^*y_2$.
Równość norm funkcjonałów reprezentujących daje
$\|T^*y\|\le\|T\|\|y\|$, więc $T^*$ jest ograniczony.
Tożsamość dla $(T^*)^*$ wynika z hermitowskości iloczynu
skalarnego i jednoznaczności operatora sprzężonego.
Stąd $\|T\|=\|(T^*)^*\|\le\|T^*\|\le\|T\|$.
Wreszcie
$\langle STx,w\rangle=\langle Tx,S^*w\rangle
=\langle x,T^*S^*w\rangle$, co daje wzór dla iloczynu.
\end{proof}

\begin{definicja}[Formy półtoraliniowe]
\label{def:sesquilinear13}
Dla przestrzeni Hilberta $H$ nad $\K\in\{\R,\C\}$
\emph{formą półtoraliniową} nazywamy odwzorowanie
$a:H\times H\to\K$ liniowe w pierwszym i antyliniowe
w drugim argumencie. W przypadku rzeczywistym są to formy
dwuliniowe. Dla $\K=\C$ warunek antyliniowości oznacza
$a(x,\alpha y+\beta z)=\overline\alpha a(x,y)+\overline\beta a(x,z)$. Forma jest \emph{ograniczona}, jeśli
$|a(x,y)|\le M\|x\|\|y\|$ dla pewnego $M\ge0$ i wszystkich
$x,y\in H$. Najmniejsze takie $M$ oznaczamy $\|a\|$.
Formę nazywamy \emph{koercywną}, jeśli istnieje $\alpha>0$
spełniające $\operatorname{Re}a(x,x)\ge\alpha\|x\|^2$
dla wszystkich $x\in H$.
\end{definicja}

\begin{twierdzenie}[Reprezentacja ograniczonych form]
\label{thm:form-repr13}
Niech $H$ będzie przestrzenią Hilberta nad $\K$,
a $a:H\times H\to\K$ formą półtoraliniową ograniczoną.
Wówczas istnieje dokładnie jeden $A\in\mathcal L(H)$ taki, że
$a(x,y)=\langle Ax,y\rangle$ dla wszystkich $x,y\in H$.
Ponadto $\|A\|=\|a\|$.
Forma jest hermitowska ($a(x,y)=\overline{a(y,x)}$) wtedy
i tylko wtedy, gdy $A=A^*$.
\end{twierdzenie}
\begin{proof}
Dla ustalonego $x\in H$ funkcja
$y\mapsto\overline{a(x,y)}$ jest ciągłym liniowym funkcjonałem,
zatem istnieje dokładnie jeden $Ax\in H$ spełniający
$\overline{a(x,y)}=\langle y,Ax\rangle$ dla każdego $y$.
Po sprzężeniu otrzymujemy $a(x,y)=\langle Ax,y\rangle$.
Jednoznaczność daje liniowość $x\mapsto Ax$.
Równość $\|Ax\|=\sup_{\|y\|\le1}|a(x,y)|\le\|a\|\|x\|$
dowodzi ograniczoności i $\|A\|\le\|a\|$.
Z drugiej strony
$|a(x,y)|=|\langle Ax,y\rangle|\le\|A\|\|x\|\|y\|$,
stąd równość norm. Równoważność hermitowskości i $A=A^*$
wynika bezpośrednio z definicji operatora sprzężonego.
\end{proof}

\begin{twierdzenie}[Twierdzenie Laxa--Milgrama]
\label{thm:lax-milgram13}
Niech $H$ będzie przestrzenią Hilberta nad $\K\in\{\R,\C\}$,
a $a:H\times H\to\K$ formą półtoraliniową ograniczoną
i koercywną ze stałą $\alpha>0$. Dla każdego ciągłego
\emph{antyliniowego} funkcjonału $\Lambda:H\to\K$ istnieje
dokładnie jeden $u\in H$ taki, że
$a(u,v)=\Lambda(v)$ dla wszystkich $v\in H$.
Ponadto $\|u\|\le\|\Lambda\|/\alpha$, gdzie
$\|\Lambda\|=\sup_{\|v\|\le1}|\Lambda(v)|$.
\end{twierdzenie}
\begin{proof}
Niech $A\in\mathcal L(H)$ reprezentuje $a$ jak w
Twierdzeniu~\ref{thm:form-repr13}. Z koercywności i nierówności
Cauchy'ego--Schwarza mamy
$\alpha\|x\|^2\le\operatorname{Re}\langle Ax,x\rangle
\le\|Ax\|\|x\|$, zatem $\|Ax\|\ge\alpha\|x\|$.
Stąd $A$ jest iniektywny i ma domknięty obraz:
$Ax_n\to y$ implikuje, że $x_n$ jest ciągiem Cauchy'ego,
a więc $x_n\to x$ i $Ax=y$.

Jeżeli $y\perp\operatorname{ran}A$, to $A^*y=0$.
Wtedy
$\alpha\|y\|^2\le\operatorname{Re}\langle Ay,y\rangle
=\operatorname{Re}\langle y,A^*y\rangle=0$,
więc $y=0$. Zatem obraz $A$ jest gęsty i jednocześnie
domknięty, czyli $\operatorname{ran}A=H$.
Ponieważ $\Lambda$ jest antyliniowy, funkcja
$v\mapsto\overline{\Lambda(v)}$ jest ciągłym liniowym
funkcjonałem. Twierdzenie~\ref{thm:riesz-hilbert13} daje
jedyny $f\in H$ spełniający $\Lambda(v)=\langle f,v\rangle$.
Bierzemy $u=A^{-1}f$; wtedy $a(u,v)=\langle Au,v\rangle
=\Lambda(v)$. Jednoznaczność wynika z iniektywności $A$,
a oszacowanie z nierówności $\|Au\|\ge\alpha\|u\|$
i $\|f\|=\|\Lambda\|$.
\end{proof}

\begin{przyklad}[Jednowymiarowy problem brzegowy]
\label{ex:weak-poisson13}
Niech $H=\{u:[0,1]\to\C:\ u\text{ jest absolutnie ciągła},
u(0)=u(1)=0,\ u'\in L^2([0,1],\lambda)\}$,
przy czym pochodną rozumiemy prawie wszędzie.
Definiujemy $\langle u,v\rangle_H=
\int_{[0,1]}u'(t)\overline{v'(t)}\,\mathrm d\lambda(t)$.
Warunki brzegowe zapewniają dodatnią określoność. Przestrzeń $H$
jest zupełna: odwzorowanie $u\mapsto u'$ jest izometryczną bijekcją
$H\to L^2_0:=\{g\in L^2([0,1],\lambda):
\int_{[0,1]}g(t)\,\mathrm d\lambda(t)=0\}$,
którego obraz jest domknięty. Odwrotność ma wzór
$g\mapsto(x\mapsto\int_{[0,x]}g(t)\,\mathrm d\lambda(t))$.
Dla $u\in H$ nierówność Cauchy'ego--Schwarza daje
$|u(x)|\le\sqrt{x}\,\|u'\|_2$, a więc
$\|u\|_2^2\le\int_{[0,1]}t\|u'\|_2^2\,\mathrm d\lambda(t)
=\frac12\|u'\|_2^2$, a więc
$\|u\|_2\le\|u'\|_2/\sqrt2$.

Dla $f\in L^2([0,1],\lambda)$ forma
$a(u,v)=\int_{[0,1]}u'(t)\overline{v'(t)}\,\mathrm d\lambda(t)$
jest ograniczona i koercywna ze stałą $1$, zaś
$\Lambda(v)=\int_{[0,1]}f(t)\overline{v(t)}\,\mathrm d\lambda(t)$
jest antyliniowa i ciągła na $H$, ponieważ
$|\Lambda(v)|\le\|f\|_2\|v\|_2
\le\|f\|_2\|v\|_H/\sqrt2$.
Twierdzenie~\ref{thm:lax-milgram13} daje jedyne $u\in H$
spełniające $a(u,v)=\Lambda(v)$ dla wszystkich $v\in H$.
Jest to słaba postać zagadnienia $-u''=f$, $u(0)=u(1)=0$.
Aby wyznaczyć rozwiązanie jawnie, określmy
$c=\int_{[0,1]}(1-s)f(s)\,\mathrm d\lambda(s)$ oraz
\[
 g(t)=c-\int_{[0,t]}f(s)\,\mathrm d\lambda(s),\qquad
 w(x)=\int_{[0,x]}g(t)\,\mathrm d\lambda(t).
\]
Ponieważ $f\in L^2([0,1])\subseteq L^1([0,1])$,
funkcja $g$ jest absolutnie ciągła, $g'=-f$ prawie wszędzie
i $g\in L^2$. Funkcja $w$ jest absolutnie ciągła,
$w'=g$, $w(0)=0$, a z twierdzenia Fubiniego
\[
 w(1)=c-\int_{[0,1]}(1-s)f(s)\,\mathrm d\lambda(s)=0.
\]
Zatem $w\in H$ i $w''=-f$ prawie wszędzie.
Dla gładkich funkcji $v$ o zerowych wartościach na brzegu
całkowanie przez części daje
$\int_{[0,1]}w'\overline{v'}\,\mathrm d\lambda
=\int_{[0,1]}f\overline v\,\mathrm d\lambda$.
Równość rozciąga się na całe $H$ przez gęstość.
Istotnie, dla $h=v'\in L^2_0$ wybieramy
$h_n\in C^\infty([0,1])$ takie, że $h_n\to h$ w $L^2$,
i kładziemy
$\widetilde h_n=h_n-\int_{[0,1]}h_n(t)\,\mathrm d\lambda(t)$.
Wówczas $\widetilde h_n\to h$ w $L^2_0$, zaś funkcje
$v_n(x)=\int_{[0,x]}\widetilde h_n(t)\,\mathrm d\lambda(t)$
należą do $C^\infty([0,1])$, znikają na brzegu i zbiegają
do $v$ w normie $H$. Ciągłość obu stron słabej tożsamości
pozwala zatem przejść do granicy $v_n\to v$.
Jednoznaczność rozwiązania z twierdzenia Laxa--Milgrama
daje $u=w$, a więc
$u'(t)=c-\int_{[0,t]}f(s)\,\mathrm d\lambda(s)$.
\end{przyklad}

\begin{uwaga}[Droga do pochodnych Gâteaux i Frécheta]
\label{rem:gateaux-bridge13}
Jeśli $H$ jest rzeczywistą przestrzenią Hilberta, to funkcjonał
$\Phi(x)=\frac12\|x\|^2$ ma pochodną Frécheta
$D\Phi(x)h=\langle h,x\rangle$. Formalnie wynika to z
$\Phi(x+h)-\Phi(x)-\langle h,x\rangle=\frac12\|h\|^2$.
Definicję pochodnej Frécheta i jej relację do pochodnej Gâteaux
wprowadzimy w następnym rozdziale. W przestrzeni zespolonej
$\Phi$ jest funkcją rzeczywistą i rozważamy jej pochodną
\emph{rzeczywistą}, równą $h\mapsto\operatorname{Re}\langle h,x\rangle$.
\end{uwaga}

% Wersja źródłowa: audyt II, jednostka 14
\setcounter{section}{13}
\section{Pochodne Gâteaux i Frécheta w przestrzeniach Banacha i Hilberta}
\label{sec:gateaux-frechet14}

\subsection{Przygotowanie: kierunki, przyrosty i liniowe przybliżenia}

W analizie funkcji jednej zmiennej pochodna opisuje najlepsze liniowe przybliżenie funkcji. Dla odwzorowania między przestrzeniami unormowanymi możliwe są dwie naturalne interpretacje: przybliżenie \emph{osobno w każdym ustalonym kierunku} oraz przybliżenie \emph{jednostajne względem wszystkich małych przyrostów}. Pierwsza prowadzi do pochodnej Gâteaux, druga do pochodnej Frécheta. Różnica jest istotna już w $\R^2$.

Niech $E,F$ będą rzeczywistymi przestrzeniami Banacha, $U\subseteq E$ zbiorem otwartym i $f:U\to F$. Przez $\mathcal L(E,F)$ oznaczamy przestrzeń ograniczonych operatorów liniowych, wyposażoną w normę operatorową. Otwartość $U$ zapewnia, że dla $x\in U$ i $h\in E$ punkty $x+th$ należą do $U$ dla wszystkich dostatecznie małych $|t|$.

\begin{definicja}[Pochodna kierunkowa]
\label{def:directional14}
Dla $x\in U$ i $h\in E$ mówimy, że $f$ ma \emph{dwustronną pochodną kierunkową} w punkcie $x$ w kierunku $h$, jeśli w $F$ istnieje granica
$\partial_h f(x)=\lim_{t\to0,\,t\in\R}(f(x+th)-f(x))/t$.
Nie wymagamy, aby zależność $h\mapsto\partial_h f(x)$ była liniowa ani ciągła. Istnienie wszystkich pochodnych kierunkowych samo w sobie nie oznacza różniczkowalności Gâteaux.
\end{definicja}

\begin{definicja}[Różniczkowalność Gâteaux]
\label{def:gateaux14}
Odwzorowanie $f$ jest \emph{różniczkowalne w sensie Gâteaux} w $x\in U$, jeśli istnieje $A\in\mathcal L(E,F)$ takie, że dla każdego $h\in E$ zachodzi $\partial_h f(x)=Ah$. Operator $A$, jednoznacznie wyznaczony przez te granice, oznaczamy $D_Gf(x)$. Przyjęliśmy definicję wymagającą ciągłości pochodnej jako operatora liniowego; w niektórych źródłach spotyka się słabszą konwencję, dopuszczającą nieograniczone odwzorowanie liniowe.
\end{definicja}

\begin{definicja}[Różniczkowalność Frécheta]
\label{def:frechet14}
Odwzorowanie $f$ jest \emph{różniczkowalne w sensie Frécheta} w $x\in U$, jeśli istnieje $A\in\mathcal L(E,F)$ oraz funkcja $r$ określona dla małych $h\in E$ taka, że
\[
f(x+h)=f(x)+Ah+r(h),\qquad
\lim_{\|h\|_E\to0,\,h\ne0}\frac{\|r(h)\|_F}{\|h\|_E}=0.
\]
Operator $A$ oznaczamy $Df(x)$ albo $D_Ff(x)$; piszemy również $r(h)=o(\|h\|_E)$. Symbol $Df(x)h$ oznacza wartość operatora $Df(x)$ na wektorze $h$, nie jego iloczyn skalarny z $h$.
\end{definicja}

\begin{uwaga}[Założenia o zupełności]
\label{rem:completeness14}
Do samych definicji wystarczą przestrzenie unormowane i otwarta dziedzina. Zupełność potrzebna będzie między innymi do twierdzenia o całkowaniu funkcji o wartościach w przestrzeni Banacha oraz w późniejszych twierdzeniach o funkcji odwrotnej. Stąd w niniejszym rozdziale pracujemy przede wszystkim w przestrzeniach Banacha.
\end{uwaga}

\subsection{Implikacje między pochodnymi}

\begin{twierdzenie}[Fréchet implikuje Gâteaux]
\label{thm:frechet-implies-gateaux14}
Niech $E,F$ będą rzeczywistymi przestrzeniami unormowanymi, $U\subseteq E$ zbiorem otwartym, a $f:U\to F$ odwzorowaniem różniczkowalnym w sensie Frécheta w punkcie $x\in U$. Wówczas $f$ jest ciągłe w $x$, jest różniczkowalne w sensie Gâteaux w $x$ oraz $D_Gf(x)=Df(x)$. Pochodna Frécheta, jeśli istnieje, jest jednoznaczna.
\end{twierdzenie}
\begin{proof}
Niech $A=Df(x)$ i $f(x+h)-f(x)=Ah+r(h)$, gdzie $\|r(h)\|=o(\|h\|)$. Ponieważ $A$ jest ograniczony, $\|Ah+r(h)\|\le\|A\|\|h\|+o(\|h\|)\to0$, co daje ciągłość. Dla ustalonego $h\ne0$ i małego rzeczywistego $t$ mamy
$\|(f(x+th)-f(x))/t-Ah\|=\|r(th)\|/|t|=\|h\|\,\|r(th)\|/\|th\|\to0$.
Dla $h=0$ pochodna kierunkowa wynosi zero. Jeżeli $A$ i $B$ spełniają definicję Frécheta, powyższy argument daje $Ah=Bh$ dla każdego $h$, a więc $A=B$.
\end{proof}

\begin{stwierdzenie}[Kryterium jednostajnego przybliżenia kierunkowego]
\label{prop:uniform-directions14}
Niech $E,F$ będą rzeczywistymi przestrzeniami unormowanymi, $U\subseteq E$ otwarty, $x\in U$ i $A\in\mathcal L(E,F)$. Odwzorowanie $f:U\to F$ ma pochodną Frécheta $Df(x)=A$ wtedy i tylko wtedy, gdy
\[
\lim_{t\to0,\,t\in\R}\ \sup_{\|h\|_E=1}
\left\|\frac{f(x+th)-f(x)}{t}-Ah\right\|_F=0.
\]
W wyrażeniu supremum rozpatrujemy wyłącznie $t\ne0$ o module na tyle małym, że $x+tB_E\subseteq U$; zakładamy też $E\ne\{0\}$.
\end{stwierdzenie}
\begin{proof}
Dla reszty $r(k)=f(x+k)-f(x)-Ak$ mamy
$\sup_{\|h\|=1}\|r(th)\|/|t|=\sup_{\|k\|=|t|}\|r(k)\|/\|k\|$.
Warunek Frécheta zapewnia, że te suprema dążą do zera. Odwrotnie, dla $k\ne0$ połóżmy $t=\|k\|$ i $h=k/\|k\|$; uzyskana jednostajna granica daje $\|r(k)\|/\|k\|\to0$.
\end{proof}

\begin{uwaga}
\label{rem:quantifiers14}
W definicji Gâteaux kwantyfikatory mają kolejność: \emph{dla każdego ustalonego $h$ granica przy $t\to0$}. W definicji Frécheta wymagamy wspólnej kontroli błędu dla wszystkich kierunków. Zamiana kolejności kwantyfikatorów nie jest uprawniona, nawet w skończonym wymiarze.
\end{uwaga}

\subsection{Kontrprzykłady: czego pochodna Gâteaux nie gwarantuje}

\begin{przyklad}[Wszystkie pochodne kierunkowe i brak ciągłości]
\label{ex:directional-discontinuous14}
Dla $(x,y)\ne(0,0)$ połóżmy
$g(x,y)=x^3y/(x^6+y^2)$, zaś $g(0,0)=0$. Dla każdego $(a,b)\in\R^2$ granica $\partial_{(a,b)}g(0,0)$ istnieje i jest równa $0$: jeśli $b\ne0$, to $g(ta,tb)/t=a^3b\,t/(a^6t^4+b^2)\to0$; jeśli $b=0$, iloraz jest zerowy. Zatem $D_Gg(0,0)=0\in\mathcal L(\R^2,\R)$. Jednak na krzywej $y=x^3$, $x\ne0$, mamy $g(x,x^3)=1/2$. Funkcja nie jest ciągła w $(0,0)$, więc nie jest tam różniczkowalna w sensie Frécheta.
\end{przyklad}

\begin{przyklad}[Ciągłość i pochodna Gâteaux nie implikują pochodnej Frécheta]
\label{ex:continuous-gateaux-not-frechet14}
Niech $f(0,0)=0$, zaś dla $(x,y)\ne(0,0)$ niech
\[
f(x,y)=\frac{x^3y}{x^4+y^2}.
\]
Z nierówności $2|x^2y|\le x^4+y^2$ wynika $|f(x,y)|\le|x|/2$; zatem $f$ jest ciągła w $(0,0)$. Dla ustalonego $(a,b)$, gdy $b\ne0$, otrzymujemy
$f(ta,tb)/t=a^3b\,t/(a^4t^2+b^2)\to0$, a dla $b=0$ iloraz jest zerowy. Stąd $D_Gf(0,0)=0$. Natomiast $f(x,x^2)=x/2$ dla $x\ne0$ oraz
$|f(x,x^2)|/\sqrt{x^2+x^4}=1/(2\sqrt{1+x^2})\to1/2$ przy $x\to0^+$. Reszta względem zerowego operatora nie jest $o(\|(x,y)\|_2)$: pochodna Frécheta nie istnieje.
\end{przyklad}

\begin{przyklad}[Sama pochodna kierunkowa nie musi być liniowa]
\label{ex:nonlinear-directional14}
Dla $f(x,y)=|x|+|y|$ w zerze istnieje tylko jednostronna pochodna kierunkowa, równa $|a|+|b|$ w kierunku $(a,b)$; dwustronna pochodna kierunkowa przy $h\ne0$ nie istnieje. Dla przykładu z \emph{dwustronnymi} pochodnymi kierunkowymi, lecz nieliniową zależnością od kierunku, rozważmy $q(0,0)=0$ oraz $q(x,y)=x^3/(x^2+y^2)$ dla $(x,y)\ne0$. Wówczas $\partial_{(a,b)}q(0,0)=a^3/(a^2+b^2)$ dla $(a,b)\ne(0,0)$, a w kierunku zerowym pochodna wynosi $0$. Odwzorowanie to jest jednorodne względem mnożenia kierunku przez rzeczywiste skalary, lecz nie jest addytywne: $\partial_{(1,0)}q+\partial_{(0,1)}q=1$, podczas gdy $\partial_{(1,1)}q=1/2$. Dlatego $q$ nie ma pochodnej Gâteaux według Definicji~\ref{def:gateaux14}.
\end{przyklad}

\begin{stwierdzenie}[Brak reguły łańcuchowej dla samej pochodnej Gâteaux]
\label{prop:chain-fails14}
Istnieją ciągłe odwzorowanie $f:\R^2\to\R$, różniczkowalne w sensie Gâteaux w zerze, oraz gładka krzywa $\gamma:\R\to\R^2$ z $\gamma(0)=0$ takie, że $\frac{\mathrm d}{\mathrm dt}(f\circ\gamma)(0)\ne D_Gf(0)\gamma'(0)$.
\end{stwierdzenie}
\begin{proof}
Niech $f$ będzie funkcją z Przykładu~\ref{ex:continuous-gateaux-not-frechet14}, a $\gamma(t)=(t,t^2)$. Wówczas $f(\gamma(t))=t/2$ dla $t\ne0$, a także dla $t=0$ z ciągłości. Zatem $(f\circ\gamma)'(0)=1/2$, podczas gdy $D_Gf(0)=0$. Niezgodność wynika z niestabilności kierunkowego przybliżenia przy zmiennym kierunku przyrostu.
\end{proof}

\begin{center}
\begin{tikzpicture}[x=2.15cm,y=1.65cm,>=Stealth]
  \draw[->,thin] (-0.17,0)--(1.7,0) node[below] {$x$};
  \draw[->,thin] (0,-0.42)--(0,1.48) node[left] {$y$};
  \draw[gray,densely dashed] (0,0)--(1.32,0.62);
  \draw[gray,densely dashed] (0,0)--(0.79,1.22);
  \draw[gray,densely dashed] (0,0)--(1.42,-0.29);
  \draw[very thick,blue!55!black,domain=0:1.15,samples=70,smooth]
       plot (\x,{\x*\x});
  \fill (0,0) circle (1.05pt);
  \node[anchor=west,blue!55!black] at (1.17,1.22) {\footnotesize $y=x^2$};
  \node[anchor=west] at (0.83,-0.37) {\scriptsize kierunki stałe};
\end{tikzpicture}

{\small Proste przez zero odpowiadają ustalonym kierunkom w definicji
Gâteaux. Krzywa $y=x^2$ prowadzi do zera po zmieniającym się
kierunku; na niej ujawnia się brak różniczkowalności Frécheta.}
\end{center}

\subsection{Pochodna Gâteaux normy w \texorpdfstring{$\ell^1$}{ell1} bez pochodnej Frécheta}

\begin{twierdzenie}[Kontrprzykład w nieskończonym wymiarze]
\label{thm:l1-gateaux14}
Niech $E=\ell^1(\N;\R)$ i $N:E\to\R$, $N(x)=\|x\|_1$. Jeśli $x=(x_n)\in E$ spełnia $x_n\ne0$ dla wszystkich $n$, to $N$ jest różniczkowalne w sensie Gâteaux w $x$ i
$D_GN(x)h=\sum_{n=1}^\infty \operatorname{sgn}(x_n)h_n$.
Funkcja $N$ nie jest jednak różniczkowalna w sensie Frécheta w żadnym takim punkcie $x$.
\end{twierdzenie}
\begin{proof}
Dla ustalonego $h\in\ell^1$ i $t\ne0$ oznaczmy $a_n(t)=(|x_n+th_n|-|x_n|)/t$. Ponieważ $x_n\ne0$, mamy $a_n(t)\to\operatorname{sgn}(x_n)h_n$ dla każdego $n$, przy czym $|a_n(t)|\le|h_n|$. Zbieżność zmajoryzowana dla szeregu (albo argument z obcięciem sumy i oszacowaniem ogona przez $\sum_{n>N}|h_n|$) daje
$\lim_{t\to0}(N(x+th)-N(x))/t=\sum_n \operatorname{sgn}(x_n)h_n$.
Jest to funkcjonał liniowy i ciągły o normie $1$.

Jeżeli istniałaby pochodna Frécheta, musiałaby być równa $D_GN(x)$. Dla $h^{(n)}=-2x_ne_n$ mamy $\|h^{(n)}\|_1=2|x_n|\to0$, ponieważ $x\in\ell^1$, a zarazem $N(x+h^{(n)})=N(x)$ i $D_GN(x)h^{(n)}=-2|x_n|$. Wobec tego
\[
\frac{|N(x+h^{(n)})-N(x)-D_GN(x)h^{(n)}|}{\|h^{(n)}\|_1}=1
\]
dla każdego $n$. Jest to sprzeczne z warunkiem Frécheta.
\end{proof}

\begin{uwaga}
Jeśli pewna współrzędna $x_n$ znika, norma $\ell^1$ nie ma w tym punkcie nawet dwustronnej pochodnej kierunkowej w kierunku $e_n$. Powyższy kontrprzykład pokazuje więc ścisłą różnicę między Gâteaux a Fréchetem również w punktach, w których wszystkie pochodne kierunkowe tworzą ciągły funkcjonał liniowy.
\end{uwaga}

\subsection{Reguły różniczkowania w przestrzeniach Banacha}

\begin{stwierdzenie}[Suma i operator liniowy]
\label{prop:sum-linear14}
Niech $E,F,G$ będą rzeczywistymi przestrzeniami Banacha, $U\subseteq E$ otwarty i $x\in U$. Jeżeli $f,g:U\to F$ są różniczkowalne w sensie Frécheta w $x$, a $\alpha,\beta\in\R$, to $D(\alpha f+\beta g)(x)=\alpha Df(x)+\beta Dg(x)$. Jeśli $T\in\mathcal L(F,G)$, to $D(T\circ f)(x)=T\circ Df(x)$. Ponadto dla $L\in\mathcal L(E,F)$ i $b\in F$ funkcja $f(y)=Ly+b$ ma $Df(y)=L$ w każdym $y\in E$.
\end{stwierdzenie}
\begin{proof}
Wszystkie równości wynikają przez podstawienie rozwinięć $f(x+h)=f(x)+Df(x)h+r_f(h)$ oraz $g(x+h)=g(x)+Dg(x)h+r_g(h)$ i użycie $\|Tr_f(h)\|\le\|T\|\|r_f(h)\|$. W przypadku funkcji afinicznej reszta wynosi zero.
\end{proof}

\begin{twierdzenie}[Reguła łańcuchowa Frécheta]
\label{thm:chain14}
Niech $E,F,G$ będą rzeczywistymi przestrzeniami Banacha, $U\subseteq E$, $V\subseteq F$ otwarte, $f:U\to F$ oraz $g:V\to G$ i $f(U)\subseteq V$. Jeśli $f$ jest różniczkowalne w sensie Frécheta w $x\in U$, a $g$ w $f(x)$, to $g\circ f$ jest różniczkowalne w sensie Frécheta w $x$, a
$D(g\circ f)(x)=Dg(f(x))\circ Df(x)$.
\end{twierdzenie}
\begin{proof}
Oznaczmy $A=Df(x)$, $B=Dg(f(x))$. Wówczas $f(x+h)=f(x)+Ah+r(h)$, $\|r(h)\|=o(\|h\|)$, i dla $k\to0$ mamy $g(f(x)+k)=g(f(x))+Bk+s(k)$, gdzie $\|s(k)\|=o(\|k\|)$. Ponieważ $f$ jest ciągłe w $x$, dla małego $h$ można przyjąć $k=Ah+r(h)$, a $\|k\|\le(\|A\|+1)\|h\|$. Otrzymujemy
$g(f(x+h))-g(f(x))-BAh=Br(h)+s(Ah+r(h))$.
Pierwszy składnik jest $o(\|h\|)$ przez ograniczoność $B$. Drugi także: jeśli $k\ne0$, to
$\|s(k)\|/\|h\|=(\|s(k)\|/\|k\|)(\|k\|/\|h\|)\to0$;
jeśli $k=0$, to $s(k)=0$. Dowód nie wymaga ciągłości pochodnych w otoczeniu rozważanych punktów.
\end{proof}

\begin{twierdzenie}[Reguła dla odwzorowania dwuliniowego]
\label{thm:bilinear14}
Niech $E_1,E_2,G$ będą rzeczywistymi przestrzeniami Banacha i $B:E_1\times E_2\to G$ ciągłym odwzorowaniem dwuliniowym, tj. $\|B(u,v)\|\le C\|u\|\|v\|$. Wówczas $B$ jest różniczkowalne w sensie Frécheta na przestrzeni produktowej $E_1\times E_2$ i
$DB(u,v)(h,k)=B(h,v)+B(u,k)$.
Jeżeli $f:U\to E_1$, $g:U\to E_2$ są różniczkowalne w $x\in U\subseteq E$ otwartym, to
$D(B(f,g))(x)h=B(Df(x)h,g(x))+B(f(x),Dg(x)h)$.
\end{twierdzenie}
\begin{proof}
Dokładna tożsamość
$B(u+h,v+k)-B(u,v)-B(h,v)-B(u,k)=B(h,k)$
daje $\|B(h,k)\|\le C\|h\|\|k\|\le C\max(\|h\|,\|k\|)^2$. Reszta jest więc $o(\max(\|h\|,\|k\|))$, a operator liniowy wskazany w tezie jest ograniczony. Ostatni wzór wynika z reguły łańcuchowej zastosowanej do odwzorowania $z\mapsto(f(z),g(z))$ i do $B$.
\end{proof}

\begin{przyklad}[Operatory i funkcje kwadratowe]
\label{ex:quadratic14}
Niech $E$ będzie rzeczywistą przestrzenią Banacha, $F$ rzeczywistą przestrzenią Banacha i $B:E\times E\to F$ ciągłym odwzorowaniem dwuliniowym. Dla $Q(x)=B(x,x)$ zachodzi
$DQ(x)h=B(h,x)+B(x,h)$.
W szczególności, jeśli $H$ jest rzeczywistą przestrzenią Hilberta i $\Phi(x)=\frac12\|x\|^2$, to
$D\Phi(x)h=\langle h,x\rangle$, a reszta ma postać $\frac12\|h\|^2$.
\end{przyklad}

\subsection{Różniczkowalność klasy \texorpdfstring{$C^1$}{C1} i twierdzenie o przyrostach}

\begin{definicja}[Klasa $C^1$]
\label{def:c1-14}
Niech $E,F$ będą rzeczywistymi przestrzeniami Banacha, $U\subseteq E$ otwarty i $f:U\to F$. Piszemy $f\in C^1(U,F)$, jeśli $f$ jest różniczkowalne w sensie Frécheta w każdym punkcie $U$ oraz odwzorowanie $Df:U\to\mathcal L(E,F)$ jest ciągłe w normie operatorowej.
\end{definicja}

\begin{lemat}[Podstawowe twierdzenie rachunku dla krzywych Banacha]
\label{lem:ftc-banach14}
Niech $F$ będzie rzeczywistą przestrzenią Banacha, $a<b$ i $\gamma:[a,b]\to F$ funkcją klasy $C^1$ (pochodna w końcach jednostronna i ciągle przedłużona). Wówczas $\gamma'$ jest całkowalna w sensie Riemanna z wartościami w $F$ oraz
$\gamma(b)-\gamma(a)=\int_a^b\gamma'(t)\,\mathrm dt$.
\end{lemat}
\begin{proof}
Ciągłość $\gamma'$ na zwartym przedziale daje jej jednostajną ciągłość, a przez kryterium Cauchy'ego sum Riemanna i zupełność $F$ zapewnia istnienie całki. Dla podziału $a=t_0<\cdots<t_m=b$ wybierzmy dowolne $\xi_j\in[t_{j-1},t_j]$. Z definicji pochodnej i jednostajnej ciągłości $\gamma'$ wynika, że przy dostatecznie małej średnicy podziału
$\|\gamma(t_j)-\gamma(t_{j-1})-(t_j-t_{j-1})\gamma'(\xi_j)\|\le\varepsilon(t_j-t_{j-1})$.
Ściślej, dla dowolnego rzeczywistego funkcjonału $\ell\in F'$ o normie co najwyżej $1$ stosujemy skalarne twierdzenie o wartości średniej do $t\mapsto\ell(\gamma(t)-t\gamma'(\xi_j))$ na $[t_{j-1},t_j]$. Otrzymujemy górne oszacowanie modułu przyrostu przez $(t_j-t_{j-1})\sup_{t\in[t_{j-1},t_j]}\|\gamma'(t)-\gamma'(\xi_j)\|$. Z twierdzenia Hahna--Banacha, po supremum po takich $\ell$, dostajemy to samo oszacowanie normy wektora. Jednostajna ciągłość $\gamma'$ zapewnia żądaną nierówność dla dostatecznie drobnych podziałów. Sumowanie oraz przejście z średnicą podziału do zera kończy dowód.
\end{proof}

\begin{twierdzenie}[Wzór całkowy i oszacowanie przyrostu]
\label{thm:mean-value14}
Niech $E,F$ będą rzeczywistymi przestrzeniami Banacha, $U\subseteq E$ zbiorem otwartym, $f\in C^1(U,F)$, a $x,y\in U$ takimi punktami, że $[x,y]=\{x+t(y-x):t\in[0,1]\}\subseteq U$. Wówczas
\[
f(y)-f(x)=\int_0^1 Df(x+t(y-x))(y-x)\,\mathrm dt,
\]
gdzie całka ma wartości w $F$. W szczególności, jeśli $\sup_{z\in[x,y]}\|Df(z)\|\le M$, to $\|f(y)-f(x)\|\le M\|y-x\|$.
\end{twierdzenie}
\begin{proof}
Definiujemy $\gamma(t)=f(x+t(y-x))$. Reguła łańcuchowa daje $\gamma'(t)=Df(x+t(y-x))(y-x)$; ciągłość $Df$ zapewnia ciągłość $\gamma'$. Lemat~\ref{lem:ftc-banach14} daje wzór. Oszacowanie normy całki przez całkę normy wynika z oszacowania sum Riemanna i przejścia do granicy.
\end{proof}

\begin{twierdzenie}[Ciągła pochodna Gâteaux daje pochodną Frécheta]
\label{thm:continuous-gateaux14}
Niech $E,F$ będą rzeczywistymi przestrzeniami Banacha, $U\subseteq E$ otwarty, $f:U\to F$. Załóżmy, że $f$ jest różniczkowalne w sensie Gâteaux w każdym punkcie $U$ i odwzorowanie $D_Gf:U\to\mathcal L(E,F)$ jest ciągłe w normie operatorowej. Wtedy $f\in C^1(U,F)$ oraz $Df(x)=D_Gf(x)$ dla wszystkich $x\in U$.
\end{twierdzenie}
\begin{proof}
Ustalmy $x\in U$ i weźmy małe $h$ takie, że $[x,x+h]\subseteq U$. Funkcja $\gamma(t)=f(x+th)$ jest różniczkowalna dla $t\in(0,1)$ i $\gamma'(t)=D_Gf(x+th)h$, co wynika z definicji kierunkowej. Jej pochodna jest ciągła dzięki założeniu, a $\gamma$ jest ciągła również w końcach przedziału, ponieważ istnienie pochodnej kierunkowej w tych punktach zapewnia ciągłość wzdłuż tej prostej. Z Lematu~\ref{lem:ftc-banach14} mamy
$f(x+h)-f(x)=\int_0^1 D_Gf(x+th)h\,\mathrm dt$.
Po odjęciu $D_Gf(x)h$ i oszacowaniu normy dostajemy
\[
\frac{\|f(x+h)-f(x)-D_Gf(x)h\|}{\|h\|}
\le\sup_{0\le t\le1}\|D_Gf(x+th)-D_Gf(x)\|\longrightarrow0.
\]
To dowodzi różniczkowalności Frécheta; ciągłość pochodnej otrzymujemy z założenia.
\end{proof}

\begin{przyklad}[Pochodna istniejąca, ale nieciągła]
\label{ex:derivative-noncontinuous14}
Na $\R$ niech $f(0)=0$ oraz $f(x)=x^2\sin(1/x^2)$ dla $x\ne0$. Wtedy $f'(0)=0$, natomiast dla $x\ne0$ mamy $f'(x)=2x\sin(1/x^2)-2x^{-1}\cos(1/x^2)$. Zatem $f$ jest różniczkowalna Frécheta w każdym punkcie, lecz nie należy do $C^1(\R,\R)$. Nie należy utożsamiać istnienia pochodnej z jej ciągłością.
\end{przyklad}

\subsection{Norma i gradient w rzeczywistych i zespolonych przestrzeniach Hilberta}

\begin{twierdzenie}[Różniczkowanie kwadratu normy i normy]
\label{thm:hilbert-norm14}
Niech $H$ będzie przestrzenią Hilberta nad $\K\in\{\R,\C\}$, traktowaną jako \emph{rzeczywista} przestrzeń Banacha, i niech $\Phi(x)=\frac12\|x\|^2$. W każdym $x\in H$ istnieje rzeczywista pochodna Frécheta $D_{\R}\Phi(x)h=\operatorname{Re}\langle h,x\rangle$. Dla $x\ne0$ norma $N(x)=\|x\|$ ma rzeczywistą pochodną Frécheta
$D_{\R}N(x)h=\operatorname{Re}\langle h,x\rangle/\|x\|$.
Jeżeli $H\ne\{0\}$, norma nie jest różniczkowalna Gâteaux w $0$ jako funkcja na rzeczywistej przestrzeni Banacha.
\end{twierdzenie}
\begin{proof}
Z rozwinięcia iloczynu skalarnego wynika dokładna tożsamość
$\Phi(x+h)-\Phi(x)-\operatorname{Re}\langle h,x\rangle=\frac12\|h\|^2$, a funkcjonał $h\mapsto\operatorname{Re}\langle h,x\rangle$ jest rzeczywisto-liniowy i ograniczony. To dowodzi pierwszej tezy. Dla $x\ne0$ stosujemy rzeczywistą regułę łańcuchową do $N(x)=\sqrt{2\Phi(x)}$, ponieważ pierwiastek jest gładki na $(0,\infty)$. W zerze iloraz kierunkowy wynosi $\|th\|/t=(|t|/t)\|h\|$; dla $h\ne0$ nie istnieje granica dwustronna.
\end{proof}

\begin{definicja}[Gradient na rzeczywistej przestrzeni Hilberta]
\label{def:gradient14}
Niech $H$ będzie rzeczywistą przestrzenią Hilberta, $U\subseteq H$ otwarty i $f:U\to\R$ różniczkowalna w sensie Frécheta w $x\in U$. \emph{Gradientem} $f$ w $x$ nazywamy jedyny wektor $\nabla f(x)\in H$ spełniający $Df(x)h=\langle h,\nabla f(x)\rangle$ dla wszystkich $h\in H$. Istnienie i jednoznaczność daje twierdzenie Riesza. Dla zespolonej przestrzeni Hilberta i funkcjonału rzeczywistego używamy struktury rzeczywistej i iloczynu skalarnego $(h,k)_{\R}=\operatorname{Re}\langle h,k\rangle$.
\end{definicja}

\begin{przyklad}[Gradient energii kwadratowej]
\label{ex:quadratic-energy14}
Niech $H$ będzie rzeczywistą lub zespoloną przestrzenią Hilberta, $A\in\mathcal L(H)$ operatorem samosprzężonym, $b\in H$, a $J(x)=\frac12\langle Ax,x\rangle-\operatorname{Re}\langle x,b\rangle$. Z samosprzężoności wynika, że $J$ ma wartości rzeczywiste. Dla $h\in H$ otrzymujemy
$D_{\R}J(x)h=\operatorname{Re}\langle h,Ax-b\rangle$, więc $\nabla J(x)=Ax-b$ w rzeczywistej strukturze Hilberta. Jeżeli $A$ jest dodatnio określony w sensie $\langle Ax,x\rangle\ge\alpha\|x\|^2$ dla pewnego $\alpha>0$, to $J$ ma dokładnie jeden minimizator: rozwiązanie $Ax=b$; jego istnienie zapewnia przypadek samosprzężony twierdzenia Laxa--Milgrama.
\end{przyklad}

\begin{przyklad}[Funkcjonał całkowy i pochodna na $L^2$]
\label{ex:integral-functional14}
Niech $(\Omega,\mathcal A,\mu)$ będzie przestrzenią z miarą skończoną i $H=L^2(\Omega,\mu;\R)$. Dla $g\in H$ określmy $J(u)=\frac12\int_\Omega |u(t)-g(t)|^2\,\mathrm d\mu(t)$. Wtedy
$DJ(u)h=\int_\Omega (u(t)-g(t))h(t)\,\mathrm d\mu(t)$,
a $\nabla J(u)=u-g$. Wynik można otrzymać przez rozwinięcie kwadratu oraz nierówność Cauchy'ego--Schwarza; reszta ma dokładnie postać $\frac12\|h\|_2^2$.
\end{przyklad}

\subsection{Rzeczywista a zespolona różniczkowalność}

\begin{definicja}[Pochodna zespolona Frécheta i Gâteaux]
\label{def:complex-differentiability14}
Niech $E,F$ będą zespolonymi przestrzeniami Banacha, $U\subseteq E$ zbiorem otwartym i $f:U\to F$. \emph{Zespoloną pochodną Frécheta} $D_{\C}f(x)$ nazywamy operator $A\in\mathcal L_{\C}(E,F)$ spełniający $f(x+h)=f(x)+Ah+o(\|h\|)$. \emph{Zespolona pochodna Gâteaux} w $x$ oznacza istnienie operatora $A\in\mathcal L_{\C}(E,F)$ takiego, że dla każdego $h\in E$
$\lim_{z\to0,\,z\in\C}(f(x+zh)-f(x))/z=Ah$.
W obu definicjach liniowość wymagana jest nad $\C$, a dla pochodnej Gâteaux parametr $z$ przebiega płaszczyznę zespoloną, nie tylko oś rzeczywistą.
\end{definicja}

\begin{stwierdzenie}[Pochodna zespolona a rzeczywista]
\label{prop:complex-real14}
Niech $E,F$ będą zespolonymi przestrzeniami Banacha, $U\subseteq E$ otwarty i $f:U\to F$. Jeśli $f$ jest zespolenie różniczkowalne w sensie Frécheta w $x\in U$, to jest też rzeczywiście różniczkowalne w sensie Frécheta w $x$, a $D_{\R}f(x)$ jest operatorem rzeczywistym odpowiadającym $D_{\C}f(x)$. Odwrotnie, jeśli $D_{\R}f(x)$ istnieje, to $D_{\C}f(x)$ istnieje wtedy i tylko wtedy, gdy $D_{\R}f(x)(ih)=iD_{\R}f(x)h$ dla wszystkich $h\in E$.
\end{stwierdzenie}
\begin{proof}
Pochodna zespolona spełnia oszacowanie reszty wymagane dla pochodnej rzeczywistej, ponieważ normy nie zmieniają się po zapomnieniu struktury zespolonej. Jeśli istnieje $D_{\R}f(x)=A$, to identyczna reszta $o(\|h\|)$ daje zespoloną pochodną dokładnie wtedy, gdy rzeczywisto-liniowy operator $A$ jest zespolenie liniowy. Dla operatora rzeczywisto-liniowego $A$ warunek $A(ih)=iA(h)$ jest równoważny liniowości zespolonej.
\end{proof}

\begin{przyklad}[Sprzężenie zespolone i kwadrat modułu]
\label{ex:complex-nonholomorphic14}
Funkcja $f:\C\to\C$, $f(z)=\overline z$, jest rzeczywiście różniczkowalna Frécheta w każdym punkcie, z $D_{\R}f(z)h=\overline h$, ale nie jest zespolenie różniczkowalna w żadnym punkcie. Podobnie $z\mapsto|z|^2$ ma rzeczywistą pochodną $h\mapsto2\operatorname{Re}(h\overline z)$, nie ma zaś zespolonej pochodnej Frécheta w punktach $z\ne0$. W punkcie $z=0$ zespolona pochodna Frécheta funkcji $z\mapsto|z|^2$ istnieje i wynosi zero, gdyż $|h|^2=o(|h|)$.
\end{przyklad}

\begin{uwaga}[Poziomy ogólności]
Związki między kierunkową a pełną różniczkowalnością, normą, topologią operatorową i różniczkowalnością zespoloną wykorzystują normowalność dziedziny. Istnieją pojęcia różniczkowania na ogólnych przestrzeniach lokalnie wypukłych (m.in. różniczkowalność Michala--Bastiani i rachunek wygodny), ale wymagają odrębnego aparatu. Nie utożsamiamy ich bez dodatkowych definicji z rachunkiem Frécheta w przestrzeniach Banacha.
\end{uwaga}

% Wersja źródłowa: audyt II, jednostka 15
\setcounter{section}{14}
\section{Pochodne wyższych rzędów i lokalne twierdzenia analizy nieliniowej}
\label{sec:wyzsze-pochodne15}

\subsection{Odwzorowania wieloliniowe i definicja wyższych pochodnych}

W podstawowym rachunku różniczkowym różniczkę $Df(x)$ traktowaliśmy jako ograniczony operator liniowy. Przy dalszym różniczkowaniu otrzymujemy odwzorowania wieloliniowe. Ta perspektywa pozwala sformułować twierdzenie Taylora i twierdzenia o funkcji odwrotnej oraz uwikłanej bez redukowania przestrzeni Banacha do współrzędnych.

\begin{definicja}[Ciągłe odwzorowania wieloliniowe]
\label{def:multilinear15}
Niech $E_1,\ldots,E_m,F$ będą rzeczywistymi przestrzeniami unormowanymi, $m\ge1$. Przez $\mathcal L^m(E_1,\ldots,E_m;F)$ oznaczamy przestrzeń odwzorowań $A:E_1\times\cdots\times E_m\to F$, liniowych w każdym argumencie i ciągłych, wyposażoną w normę
$\|A\|=\sup\{\|A(h_1,\ldots,h_m)\|:\|h_j\|\le1,\ 1\le j\le m\}$. Jeśli wszystkie dziedziny są równe $E$, piszemy $\mathcal L^m(E;F)$, a dla $m=1$ używamy dotychczasowego oznaczenia $\mathcal L(E,F)$. Przez $\mathcal L^m_s(E;F)$ oznaczamy podprzestrzeń odwzorowań symetrycznych, tj. niezmienniczych przy permutacjach argumentów.
\end{definicja}

\begin{stwierdzenie}[Kryterium ograniczoności i zupełność]
\label{prop:multi-banach15}
Niech $E_1,\ldots,E_m$ będą przestrzeniami unormowanymi, a $F$ przestrzenią Banacha. Odwzorowanie $m$-liniowe $A:E_1\times\cdots\times E_m\to F$ jest ciągłe wtedy i tylko wtedy, gdy istnieje $C\ge0$ takie, że
$\|A(h_1,\ldots,h_m)\|\le C\prod_{j=1}^m\|h_j\|$ dla wszystkich $h_j\in E_j$. Przestrzeń $\mathcal L^m(E_1,\ldots,E_m;F)$ jest przestrzenią Banacha.
\end{stwierdzenie}
\begin{proof}
Jeśli $A$ jest ciągłe w zerze, istnieje $\delta>0$ takie, że $\|A(v_1,\ldots,v_m)\|\le1$ dla $\|v_j\|<\delta$. Dla niezerowych $h_j$ podstawiamy $v_j=\delta h_j/(2\|h_j\|)$, co po wykorzystaniu wieloliniowości daje oszacowanie z $C=(2/\delta)^m$. Gdy któryś argument jest zerowy, nierówność jest oczywista. W drugą stronę nierówność zapewnia ciągłość w zerze, a wieloliniowość i rozwinięcie $A(h_1+k_1,\ldots,h_m+k_m)$ w skończoną sumę zapewniają ciągłość w dowolnym punkcie.

Niech $(A_n)$ będzie ciągiem Cauchy'ego w normie operatorowej. Dla ustalonych $h_1,\ldots,h_m$ ciąg $A_n(h_1,\ldots,h_m)$ jest Cauchy'ego w $F$, więc definiuje $A(h_1,\ldots,h_m)$. Granica zachowuje wieloliniowość. Ponieważ normy $\|A_n\|$ są ograniczone przez wspólną stałą, $A$ spełnia oszacowanie wieloliniowe. Przechodząc do granicy w nierówności $\|(A_n-A_p)(h_1,\ldots,h_m)\|\le\varepsilon\prod\|h_j\|$ dla $n,p\ge N$, otrzymujemy $\|A_n-A\|\le\varepsilon$. To dowodzi zupełności.
\end{proof}

\begin{lemat}[Utożsamienie operatorów wieloliniowych]
\label{lem:currying15}
Niech $E,F$ będą rzeczywistymi przestrzeniami unormowanymi, a $m\ge1$. Naturalne odwzorowanie
$\mathcal L(E,\mathcal L^m(E;F))\to\mathcal L^{m+1}(E;F)$, zadane przez $T\mapsto\widetilde T$ oraz $\widetilde T(h_0,h_1,\ldots,h_m)=(Th_0)(h_1,\ldots,h_m)$, jest izometrycznym izomorfizmem liniowym.
\end{lemat}
\begin{proof}
Liniowość, wieloliniowość i ciągłość wynikają z oszacowania $\|\widetilde T(h_0,\ldots,h_m)\|\le\|T\|\prod_{j=0}^m\|h_j\|$. Ponadto, z definicji norm operatorowych,
\[
\begin{aligned}
\|\widetilde T\|&=\sup_{\|h_0\|\le1}\sup_{\substack{\|h_j\|\le1\\1\le j\le m}}\|(Th_0)(h_1,\ldots,h_m)\|\\
&=\|T\|.
\end{aligned}
\]
Dla $A\in\mathcal L^{m+1}(E;F)$ definiujemy $(T_Ah_0)(h_1,\ldots,h_m)=A(h_0,h_1,\ldots,h_m)$. Wówczas $T_A\in\mathcal L(E,\mathcal L^m(E;F))$ oraz $\widetilde T_A=A$, co daje surjektywność.
\end{proof}

\begin{definicja}[Klasa $C^m$ i pochodne wyższych rzędów]
\label{def:ck15}
Niech $E,F$ będą rzeczywistymi przestrzeniami Banacha, $U\subseteq E$ zbiorem otwartym, $m\ge1$. Mówimy, że $f:U\to F$ jest klasy $C^m$, gdy $f\in C^1$ oraz rekurencyjnie $D^{j-1}f:U\to\mathcal L^{j-1}(E;F)$ jest klasy $C^1$ dla $2\le j\le m$. Pochodną $D^jf(x)$ utożsamiamy przez Lemat~\ref{lem:currying15} z elementem $\mathcal L^j(E;F)$. Przyjmujemy $D^0f=f$ oraz $C^\infty(U,F)=\bigcap_{m\ge1}C^m(U,F)$. Ciągłość $D^jf$ rozumiemy w normie $\mathcal L^j(E;F)$.
\end{definicja}

\begin{przyklad}[Odwzorowanie wieloliniowe i wielomiany]
\label{ex:multi-poly15}
Niech $B\in\mathcal L^m(E;F)$ i $P(x)=B(x,\ldots,x)$. Wówczas $P\in C^\infty(E,F)$ oraz
$DP(x)h=\sum_{j=1}^m B(x,\ldots,x,\underset{j}{h},x,\ldots,x)$.
Istotnie, różnica $P(x+h)-P(x)$ składa się z wyrazów, w których $h$ występuje co najmniej raz; wyrazy z dokładnie jednym $h$ tworzą wskazany operator, a suma pozostałych jest $O(\|h\|^2)$ przy ustalonym $x$. Powtarzanie argumentu daje wyższe pochodne, które zerują się od rzędu $m+1$. Dla symetrycznego $B$ zachodzi $DP(x)h=mB(x,\ldots,x,h)$, a $D^mP(x)=m!B$.
\end{przyklad}

\subsection{Symetria drugiej pochodnej i reguły rachunku}

\begin{twierdzenie}[Symetria drugiej pochodnej]
\label{thm:schwarz15}
Niech $E,F$ będą rzeczywistymi przestrzeniami Banacha, $U\subseteq E$ otwarty i $f\in C^2(U,F)$. Dla każdego $x\in U$ i $h,k\in E$ zachodzi
$D^2f(x)[h,k]=D^2f(x)[k,h]$. W szczególności $D^2f(x)\in\mathcal L_s^2(E;F)$.
\end{twierdzenie}
\begin{proof}
Dla małych rzeczywistych $s,t$ połóżmy
$\Delta(s,t)=f(x+sh+tk)-f(x+sh)-f(x+tk)+f(x)$. Z dwukrotnego zastosowania wzoru całkowego na odcinku dla funkcji klasy $C^1$ otrzymujemy
\[
\Delta(s,t)=st\int_0^1\int_0^1 D^2f(x+\alpha sh+\beta tk)[h,k] \,\mathrm d\beta\,\mathrm d\alpha.
\]
Można bowiem najpierw scałkować po kierunku $h$, a następnie różnicę wyrażeń $Df$ po kierunku $k$. Odwrócenie kolejności tych dwóch kroków daje ten sam przyrost $\Delta(s,t)$ oraz wyrażenie z $D^2f(\cdot)[k,h]$. Dla $s,t\ne0$ dzielimy przez $st$ i przechodzimy z $(s,t)$ do $(0,0)$. Ciągłość $D^2f$ w normie operatorowej zapewnia jednostajną zbieżność funkcji podcałkowych na kwadracie $[0,1]^2$ odpowiednio do $D^2f(x)[h,k]$ i $D^2f(x)[k,h]$. Obie granice są równe, co kończy dowód.
\end{proof}


\begin{wniosek}[Symetria wszystkich pochodnych wyższych rzędów]
\label{cor:symmetric-higher15}
Niech $E,F$ będą rzeczywistymi przestrzeniami Banacha, $U\subseteq E$ otwarty, $m\ge2$ i $f\in C^m(U,F)$. Wówczas $D^mf(x)\in\mathcal L_s^m(E;F)$ dla każdego $x\in U$.
\end{wniosek}
\begin{proof}
Stosujemy indukcję po $m$. Dla $m=2$ jest to Twierdzenie~\ref{thm:schwarz15}. Jeśli $m\ge3$, to odwzorowanie $D^{m-2}f:U\to\mathcal L^{m-2}(E;F)$ jest klasy $C^2$, a więc zastosowanie tego twierdzenia do $D^{m-2}f$ daje symetrię $D^mf(x)$ względem pierwszych dwóch argumentów. Z założenia indukcyjnego $D^{m-1}f(z)$ jest symetryczne dla każdego $z\in U$. Różniczkując tożsamości wyrażające tę symetrię w kierunku pierwszego argumentu, stwierdzamy, że $D^mf(x)$ jest symetryczne także względem ostatnich $m-1$ argumentów. Łącznie otrzymujemy niezmienniczość przy wszystkich zamianach sąsiednich argumentów, a te generują wszystkie permutacje.
\end{proof}

\begin{stwierdzenie}[Druga pochodna złożenia]
\label{prop:second-chain15}
Niech $E,F,G$ będą rzeczywistymi przestrzeniami Banacha, $U\subseteq E$ i $V\subseteq F$ otwarte, $f\in C^2(U,F)$, $g\in C^2(V,G)$ oraz $f(U)\subseteq V$. Wówczas $g\circ f\in C^2(U,G)$ i dla $x\in U$ oraz $h,k\in E$ zachodzi
\[
\begin{aligned}
D^2(g\circ f)(x)[h,k]
={}&D^2g(f(x))[Df(x)h,Df(x)k]\\
 &+Dg(f(x))\bigl(D^2f(x)[h,k]\bigr).
\end{aligned}
\]
\end{stwierdzenie}
\begin{proof}
Z reguły łańcuchowej $D(g\circ f)(x)=Dg(f(x))\circ Df(x)$. Kompozycja operatorów jest ciągłym odwzorowaniem dwuliniowym $\mathcal L(F,G)\times\mathcal L(E,F)\to\mathcal L(E,G)$. Stosując regułę różniczkowania ciągłego odwzorowania dwuliniowego do odwzorowań $x\mapsto Dg(f(x))$ i $x\mapsto Df(x)$, otrzymujemy po zadziałaniu na $k$ dwa składniki. Pierwszy, dzięki regule łańcuchowej zastosowanej do $Dg$, ma postać $D^2g(f(x))[Df(x)h,Df(x)k]$, a drugi jest równy $Dg(f(x))(D^2f(x)[h,k])$. Ciągłość wszystkich występujących pochodnych dowodzi przynależności złożenia do $C^2$.
\end{proof}

\begin{przyklad}[Hesjan funkcjonału na Hilbercie]
\label{ex:hessian15}
Niech $H$ będzie rzeczywistą przestrzenią Hilberta, $A\in\mathcal L(H)$ i $b\in H$. Dla $J(x)=\frac12\langle Ax,x\rangle-\langle x,b\rangle$ mamy
$DJ(x)h=\langle h,(A+A^*)x/2-b\rangle$ oraz
$D^2J(x)[h,k]=\langle h,(A+A^*)k/2\rangle$.
Przy konwencji liniowości iloczynu skalarnego w pierwszym argumencie wzór jest symetryczny dla rzeczywistego $H$, ponieważ $(A+A^*)/2$ jest samosprzężony. Jeśli $A=A^*$, to $D^2J(x)[h,k]=\langle h,Ak\rangle$.
\end{przyklad}

\subsection{Twierdzenie Taylora w przestrzeniach Banacha}

\begin{twierdzenie}[Wzór Taylora z resztą całkową]
\label{thm:taylor15}
Niech $E,F$ będą rzeczywistymi przestrzeniami Banacha, $U\subseteq E$ otwarty, $m\ge1$, $f\in C^m(U,F)$, a $x\in U$ i $h\in E$ takie, że odcinek $\{x+th:0\le t\le1\}$ zawiera się w $U$. Wtedy
\[
 f(x+h)=\sum_{j=0}^{m-1}\frac1{j!}D^jf(x)[h^j]
 +\frac1{(m-1)!}\int_0^1(1-t)^{m-1}D^mf(x+th)[h^m] \,\mathrm dt,
\]
gdzie $D^0f(x)[h^0]=f(x)$, a $D^jf(x)[h^j]$ oznacza $D^jf(x)[h,\ldots,h]$. Ponadto
$f(x+h)=\sum_{j=0}^{m}\frac1{j!}D^jf(x)[h^j]+r_m(h)$, gdzie
\[
\|r_m(h)\|\le \frac{\|h\|^m}{m!}\sup_{0\le t\le1}\|D^mf(x+th)-D^mf(x)\|.
\]
W szczególności $\|r_m(h)\|=o(\|h\|^m)$ dla $h\to0$.
\end{twierdzenie}
\begin{proof}
Ustalmy $x,h$ i zdefiniujmy $\gamma(t)=f(x+th)$. Z reguły łańcuchowej i indukcji mamy $\gamma^{(j)}(t)=D^jf(x+th)[h^j]$ dla $1\le j\le m$. Pierwsza pochodna jest ciągła, zatem stosujemy podstawowe twierdzenie rachunku całkowego dla funkcji o wartościach w $F$. Iterowanie wzoru $\gamma(t)-\gamma(0)=\int_0^t\gamma'(s)\,\mathrm ds$ prowadzi do jednokrotnej całki z jądrem $(1-t)^{m-1}/(m-1)!$: krok indukcyjny wynika z zamiany kolejności całkowania na trójkącie $0\le s\le t\le1$, dopuszczalnej dla ciągłej funkcji o wartościach w Banachu (można ją uzasadnić najpierw dla sum Riemanna). Otrzymujemy pierwszy wzór.

Odejmujemy od wyrażenia pod całką stałą $D^mf(x)[h^m]$, której całka z wagą $(1-t)^{m-1}/(m-1)!$ jest równa $D^mf(x)[h^m]/m!$. Następnie stosujemy normę operatorową, oszacowanie normy całki przez całkę normy oraz $\int_0^1(1-t)^{m-1}\,\mathrm dt=1/m$. Ciągłość $D^mf$ w $x$ zapewnia, że supremum w oszacowaniu reszty dąży do zera przy $h\to0$.
\end{proof}

\begin{wniosek}[Wzór drugiego rzędu]
\label{cor:taylor2-15}
Przy założeniach Twierdzenia~\ref{thm:taylor15} dla $m=2$ zachodzi
\[
 f(x+h)=f(x)+Df(x)h+\tfrac12D^2f(x)[h,h]
 +\int_0^1(1-t)\bigl(D^2f(x+th)-D^2f(x)\bigr)[h,h] \,\mathrm dt.
\]
Jeżeli na odcinku $[x,x+h]$ zachodzi $\|D^2f(z)\|\le M$, to reszta po wyrazach rzędu pierwszego ma normę co najwyżej $M\|h\|^2/2$.
\end{wniosek}

\begin{stwierdzenie}[Warunki konieczne i wystarczające minimum]
\label{prop:extremum15}
Niech $E$ będzie rzeczywistą przestrzenią Banacha, $U\subseteq E$ otwarty, $x_0\in U$, a $f\in C^2(U,\R)$. Jeżeli $x_0$ jest lokalnym minimum $f$, to $Df(x_0)=0$ i $D^2f(x_0)[h,h]\ge0$ dla każdego $h\in E$. Odwrotnie, jeżeli $Df(x_0)=0$ oraz istnieje $c>0$ takie, że $D^2f(x_0)[h,h]\ge c\|h\|^2$ dla wszystkich $h\in E$, to $x_0$ jest ścisłym lokalnym minimum $f$; dokładniej, dla dostatecznie małych $h$ zachodzi $f(x_0+h)\ge f(x_0)+(c/4)\|h\|^2$.
\end{stwierdzenie}
\begin{proof}
Dla dowolnego $h$ rozważamy rzeczywistą funkcję $\varphi(t)=f(x_0+th)$. Ma ona minimum w $0$, więc $\varphi'(0)=Df(x_0)h=0$ i $\varphi''(0)=D^2f(x_0)[h,h]\ge0$. W drugim kierunku ciągłość $D^2f$ daje promień $r>0$ taki, że $\|D^2f(x_0+u)-D^2f(x_0)\|\le c/2$ dla $\|u\|<r$. Stąd $D^2f(x_0+u)[h,h]\ge(c/2)\|h\|^2$. Z całkowego wzoru Taylora drugiego rzędu, przy $Df(x_0)=0$, otrzymujemy $f(x_0+h)-f(x_0)\ge(c/2)\|h\|^2\int_0^1(1-t)\,\mathrm dt=c\|h\|^2/4$ dla $\|h\|<r$.
\end{proof}

\begin{przyklad}[Dodatniość bez jednostajności nie wystarcza]
\label{ex:positive-hessian15}
W rzeczywistej przestrzeni Hilberta $H=\ell^2$ niech
$Q(x)=\frac12\sum_{n=1}^\infty x_n^2/n$, zaś $f(x)=Q(x)-\|x\|_2^3$. Funkcja $f$ jest klasy $C^2$ w otoczeniu zera. Mamy $Df(0)=0$ i $D^2f(0)[h,h]=\sum_{n\ge1}h_n^2/n>0$ dla $h\ne0$, ale $0$ nie jest minimum lokalnym. Istotnie, $f(te_n)=t^2/(2n)-t^3<0=f(0)$, gdy $t>1/(2n)$, a takie dodatnie $t$ można wybierać dowolnie małe przez zwiększanie $n$. Funkcja $G(x)=\|x\|_2^3$ jest klasy $C^2$ również w zerze. Dla $x\ne0$ mamy
$DG(x)h=3\|x\|\langle x,h\rangle$ oraz
$D^2G(x)[h,k]=3\|x\|\langle h,k\rangle+3\langle x,h\rangle\langle x,k\rangle/\|x\|$.
Stąd $\|D^2G(x)\|\le6\|x\|$. Ponadto $DG(0)=0$ i $\|DG(x)\|\le3\|x\|^2$, co dowodzi $D^2G(0)=0$ i ciągłości drugiej pochodnej w zerze.
\end{przyklad}

\subsection{Szereg Neumanna i odwracanie operatorów}

\begin{lemat}[Szereg Neumanna]
\label{lem:neumann15}
Niech $E$ będzie rzeczywistą przestrzenią Banacha i $T\in\mathcal L(E)$ spełnia $\|T\|<1$. Wtedy $I-T$ jest odwracalny oraz
$(I-T)^{-1}=\sum_{n=0}^\infty T^n$ w normie operatorowej, przy czym $\|(I-T)^{-1}\|\le(1-\|T\|)^{-1}$.
\end{lemat}
\begin{proof}
Szereg jest absolutnie zbieżny w przestrzeni Banacha $\mathcal L(E)$. Dla sum częściowych $S_N=\sum_{n=0}^N T^n$ zachodzi $(I-T)S_N=S_N(I-T)=I-T^{N+1}$; przechodząc do granicy, otrzymujemy obustronną odwrotność. Oszacowanie normy wynika z szeregu geometrycznego.
\end{proof}

\begin{stwierdzenie}[Otwartość grupy operatorów odwracalnych]
\label{prop:inverse-operator15}
Niech $E$ będzie rzeczywistą przestrzenią Banacha i $\mathrm{GL}(E)$ zbiorem odwracalnych operatorów w $\mathcal L(E)$. Wówczas $\mathrm{GL}(E)$ jest otwarty, a $\iota:\mathrm{GL}(E)\to\mathcal L(E)$, $\iota(A)=A^{-1}$, jest klasy $C^\infty$. Jej pochodna ma postać
$D\iota(A)H=-A^{-1}HA^{-1}$.
\end{stwierdzenie}
\begin{proof}
Dla $A\in\mathrm{GL}(E)$ i $\|A^{-1}H\|<1$ mamy $A+H=A(I+A^{-1}H)$, więc Lemat~\ref{lem:neumann15} daje odwracalność oraz lokalną jednostajną ograniczoność odwrotności. Z tożsamości
$(A+H)^{-1}-A^{-1}=-(A+H)^{-1}HA^{-1}$ wynika
\[
(A+H)^{-1}-A^{-1}+A^{-1}HA^{-1}
=A^{-1}H(A+H)^{-1}HA^{-1}.
\]
Jej prawą stronę ograniczamy przez $C_A\|H\|^2$ dla małych $H$, skąd wynika pochodna Frécheta. Jej ciągłość otrzymujemy z ciągłości odwracania i ciągłości kompozycji operatorów. Powtarzając różniczkowanie wzoru $D\iota(A)H=-A^{-1}HA^{-1}$ z regułą wieloliniową, otrzymujemy rekurencyjnie pochodne wszystkich rzędów i ich ciągłość.
\end{proof}

\subsection{Twierdzenie o funkcji odwrotnej}

\begin{twierdzenie}[Twierdzenie o funkcji odwrotnej dla przestrzeni Banacha]
\label{thm:inverse15}
Niech $E,F$ będą rzeczywistymi przestrzeniami Banacha, $U\subseteq E$ otwarty, $f\in C^1(U,F)$ i $x_0\in U$. Załóżmy, że $A=Df(x_0):E\to F$ jest izomorfizmem liniowym i topologicznym, tj. ma ograniczoną odwrotność $A^{-1}:F\to E$. Wówczas istnieją otoczenia otwarte $U_0\subseteq U$ punktu $x_0$ i $V_0\subseteq F$ punktu $y_0=f(x_0)$ takie, że $f|_{U_0}:U_0\to V_0$ jest bijekcją, a odwrotność $g:V_0\to U_0$ należy do $C^1(V_0,E)$ i spełnia
$Dg(y)=(Df(g(y)))^{-1}$ dla każdego $y\in V_0$.
\end{twierdzenie}
\begin{proof}
\emph{Krok 1: jednostajne przybliżenie liniowe.} Z ciągłości $Df$ w $x_0$ wybieramy $r>0$ tak małe, aby domknięta kula $\overline B(x_0,r)$ zawierała się w $U$ i
$\|I-A^{-1}Df(x)\|\le q<1$ dla $x\in\overline B(x_0,r)$; można przyjąć $q=1/2$. Dla $y\in F$ określamy $T_y(x)=x-A^{-1}(f(x)-y)$. Dla $x,z$ w tej kuli wzór całkowy na odcinku daje
\[
\|T_y(x)-T_y(z)\|\le q\|x-z\|.
\]
Ponadto $T_y(x_0)-x_0=A^{-1}(y-y_0)$. Dla $\|A^{-1}(y-y_0)\|\le(1-q)r$ operator $T_y$ przeprowadza $\overline B(x_0,r)$ w siebie, bo
$\|T_y(x)-x_0\|\le q\|x-x_0\|+\|A^{-1}(y-y_0)\|\le r$.

\emph{Krok 2: punkt stały.} Przypomnijmy zasadę kontrakcji: jeżeli $T$ jest odwzorowaniem $q$-lipschitzowskim pełnej przestrzeni metrycznej w siebie, $q<1$, to ciąg $x_{n+1}=T(x_n)$ jest Cauchy'ego, ponieważ $d(x_{n+1},x_n)\le q^n d(x_1,x_0)$; jego granica jest jedynym punktem stałym. Stosujemy tę zasadę do $T_y$ na domkniętej kuli w $E$. Otrzymujemy jedyny $x$ z $T_y(x)=x$, co jest równoważne $f(x)=y$. Jeżeli $y$ spełnia nierówność ostrą $\|A^{-1}(y-y_0)\|<(1-q)r$, punkt stały leży we wnętrzu kuli. Dodatkowo każdy $x\in B(x_0,r)$ ma na tym obszarze co najwyżej jeden punkt o tej samej wartości $f(x)$: jeżeli $f(x)=f(z)$, to $x,z$ są punktami stałymi tej samej kontrakcji $T_{f(x)}$.

\emph{Krok 3: otwarte otoczenia i ciągłość odwrotności.} Ustalmy $0<r_0<r$ i połóżmy $U_0=B(x_0,r_0)$ oraz $V_0=f(U_0)$. Z oszacowania
$\|(x-z)-A^{-1}(f(x)-f(z))\|\le q\|x-z\|$ wynika
\[
\|x-z\|\le\frac{\|A^{-1}\|}{1-q}\,\|f(x)-f(z)\| \qquad(x,z\in B(x_0,r)).
\]
Zatem $f|_{U_0}$ jest różnowartościowe, a odwrotność $g:V_0\to U_0$ jest lokalnie Lipschitzowska. Zbiór $V_0$ jest otwarty: dla dowolnego $x\in U_0$ wybieramy kulę $\overline B(x,\delta)\subset U_0$; ta sama kontrakcja $T_y$ zachowuje $\overline B(x,\delta)$ przy $y$ dostatecznie bliskim $f(x)$, ponieważ oszacowanie stałej $q$ obowiązuje w całej kuli $B(x_0,r)$. W konsekwencji otoczenie $f(x)$ zawiera się w $V_0$.

\emph{Krok 4: różniczkowalność odwrotności.} Zmniejszając $r$, jeśli trzeba, mamy $\|I-A^{-1}Df(x)\|<1$ dla $x\in U_0$, więc każdy $Df(x)$ jest odwracalny z Lematu~\ref{lem:neumann15}. Ustalmy $y=f(x)\in V_0$ i rozważmy $y+k=f(x+h)$, gdzie $h=g(y+k)-g(y)\to0$. Z różniczkowalności $f$ w $x$:
$k=Df(x)h+\rho(h)$ z $\|\rho(h)\|=o(\|h\|)$. Po pomnożeniu przez $(Df(x))^{-1}$ dostajemy
$h-(Df(x))^{-1}k=-(Df(x))^{-1}\rho(h)=o(\|h\|)=o(\|k\|)$,
gdyż otrzymana w kroku 3 nierówność zapewnia $\|h\|\le C\|k\|$. Stąd $Dg(y)=(Df(x))^{-1}$. Ciągłość $Dg$ wynika z ciągłości $Df$, $g$ oraz odwracania w $\mathrm{GL}(E,F)$ (można zastosować ten sam argument Neumanna do $(Df(x))^{-1}Df(x')$). Tak więc $g\in C^1$.
\end{proof}

\begin{wniosek}[Wersja $C^m$]
\label{cor:inverse-ck15}
Jeżeli przy założeniach Twierdzenia~\ref{thm:inverse15} odwzorowanie $f$ jest klasy $C^m$ dla $m\ge2$ albo klasy $C^\infty$, to lokalna odwrotność $g$ ma tę samą regularność.
\end{wniosek}
\begin{proof}
Wzór $Dg=(Df\circ g)^{-1}$ oraz Stwierdzenie~\ref{prop:inverse-operator15} pozwalają prowadzić indukcję: jeśli $g\in C^j$ i $f\in C^{j+1}$, to $Df\circ g\in C^j$, zatem także $Dg\in C^j$, a więc $g\in C^{j+1}$. Zaczynamy od $g\in C^1$.
\end{proof}

\begin{przyklad}[Nieliniowa mapa na przestrzeni Hilberta]
\label{ex:inverse-hilbert15}
Niech $H$ będzie rzeczywistą przestrzenią Hilberta, $v\in H$ ustalony i $f:H\to H$ dane przez $f(x)=x+\|x\|^2v$. Wtedy $Df(x)h=h+2\langle x,h\rangle v$, a $Df(0)=I$. Twierdzenie~\ref{thm:inverse15} zapewnia lokalną odwrotność w otoczeniu zera. Jeżeli $v\ne0$, nie wnioskujemy stąd, że $f$ ma odwrotność na całym $H$.
\end{przyklad}

\subsection{Twierdzenie o funkcji uwikłanej}

\begin{twierdzenie}[Twierdzenie o funkcji uwikłanej dla przestrzeni Banacha]
\label{thm:implicit15}
Niech $X,Y,Z$ będą rzeczywistymi przestrzeniami Banacha, $\mathcal U\subseteq X\times Y$ zbiorem otwartym, $F\in C^1(\mathcal U,Z)$, $(x_0,y_0)\in\mathcal U$ i $F(x_0,y_0)=0$. Załóżmy, że pochodna częściowa $D_yF(x_0,y_0):Y\to Z$ jest izomorfizmem liniowym i topologicznym. Wtedy istnieją otwarte otoczenia $U_0\subseteq X$ punktu $x_0$, $V_0\subseteq Y$ punktu $y_0$ oraz jedyna funkcja $\varphi:U_0\to V_0$ klasy $C^1$ taka, że dla $(x,y)\in U_0\times V_0$ zachodzi
$F(x,y)=0\ \Longleftrightarrow\ y=\varphi(x)$. Ponadto
\[
D\varphi(x)=-\bigl(D_yF(x,\varphi(x))\bigr)^{-1}D_xF(x,\varphi(x))\qquad(x\in U_0).
\]
Jeżeli $F\in C^m$ dla $m\ge2$ lub $C^\infty$, to również $\varphi$ ma odpowiednią regularność.
\end{twierdzenie}
\begin{proof}
Definiujemy $\Psi:\mathcal U\to X\times Z$ wzorem $\Psi(x,y)=(x,F(x,y))$. Jest to odwzorowanie klasy $C^1$ i
$D\Psi(x_0,y_0)(h,k)=(h,Bh+Ak)$, gdzie $A=D_yF(x_0,y_0)$ i $B=D_xF(x_0,y_0)$. Operator ten jest odwracalny: odwrotny przeprowadza $(h,z)$ w $(h,A^{-1}(z-Bh))$ i jest ograniczony. Twierdzenie~\ref{thm:inverse15} daje lokalny dyfeomorfizm $\Psi$ między otoczeniami $(x_0,y_0)$ oraz $(x_0,0)$. Po zmniejszeniu otoczeń możemy wybrać prostokąt $U_0\times W_0$ zawarty w obrazie $\Psi$ oraz $U_0\times V_0$ zawarty w dziedzinie lokalnego dyfeomorfizmu. Zapiszmy odwrotność $\Psi^{-1}(x,z)=(x,G(x,z))$; pierwsza współrzędna ma tę postać, ponieważ $\Psi$ zachowuje $x$. Połóżmy $\varphi(x)=G(x,0)$. Wówczas $F(x,\varphi(x))=0$.

Dla zapewnienia równoważności na całym $U_0\times V_0$ zwężamy $V_0$ i $U_0$ tak, by $\Psi$ była różnowartościowa na tym prostokącie i $\varphi(U_0)\subset V_0$. Jeśli $F(x,y)=0$, to $\Psi(x,y)=(x,0)=\Psi(x,\varphi(x))$, zatem $y=\varphi(x)$. Różniczkując tożsamość $F(x,\varphi(x))=0$ i używając reguły łańcuchowej, otrzymujemy
$D_xF(x,\varphi(x))+D_yF(x,\varphi(x))D\varphi(x)=0$.
Po ewentualnym zmniejszeniu otoczeń $D_yF$ jest odwracalne na grafie $\varphi$, bo zbiór operatorów odwracalnych jest otwarty. Stąd wzór na $D\varphi$. Wersja wyższej regularności wynika z Wniosku~\ref{cor:inverse-ck15}.
\end{proof}

\begin{przyklad}[Rodzina równań operatorowych]
\label{ex:implicit-operator15}
Niech $E$ będzie rzeczywistą przestrzenią Banacha, $K\in\mathcal L(E)$ i $g\in E$. Na $\R\times E$ definiujemy $F(\lambda,u)=u-\lambda Ku-g$. W punkcie $(0,g)$ mamy $F(0,g)=0$ oraz $D_uF(0,g)=I$, zatem istnieje lokalnie jednoznaczne rozwiązanie $u(\lambda)$ klasy $C^\infty$ równania $(I-\lambda K)u(\lambda)=g$. Z wzoru na pochodną funkcji uwikłanej wynika
$u'(\lambda)=(I-\lambda K)^{-1}Ku(\lambda)$, w szczególności $u'(0)=Kg$. Różniczkując raz jeszcze, otrzymujemy $u''(0)=2K^2g$. Dla $|\lambda|\,\|K\|<1$ rozwiązanie można również zapisać w postaci zbieżnego szeregu Neumanna $u(\lambda)=\sum_{n\ge0}\lambda^nK^ng$.
\end{przyklad}

\begin{uwaga}[Dlaczego przestrzenie Banacha i dlaczego pochodna Frécheta?]
\label{rem:limits-inverse15}
Dowód twierdzenia o funkcji odwrotnej wykorzystuje zupełność domkniętej kuli (zasada kontrakcji) oraz jednostajną kontrolę operatora $Df(x)$ w normie. Sama odwracalność pochodnej Gâteaux w punkcie, bez stosownych założeń o ciągłości i pełnej różniczkowalności, nie daje takiego wniosku. W ogólnych przestrzeniach lokalnie wypukłych istnieją odrębne twierdzenia o funkcji odwrotnej, lecz nie wolno przenosić powyższego dowodu na dowolną TVS przez formalne zastąpienie normy seminormami.
\end{uwaga}

% Wersja źródłowa: audyt II, jednostka 15A
\setcounter{section}{15}
\noindent\textbf{Wymagane wyniki.} Korzystamy z twierdzenia Hahna--Banacha, twierdzenia Banacha--Steinhausa, podstawowego twierdzenia rachunku dla krzywych o wartościach w przestrzeni Banacha i faktu, że ciągłe funkcjonały liniowe rozdzielają punkty lokalnie wypukłej przestrzeni Hausdorffa. Pełną teorię całek na takich przestrzeniach rozwijamy poniżej.
\medskip
% \hyperref[sec:integral-lcs15a]{Moduł 15A}. Celowo nie zmieniamy numeracji \hyperref[sec:wariacyjny16]{rozdziałów 16}--33.
\begingroup
\renewcommand{\thesection}{15A}
\renewcommand{\theHsection}{15A}
\section{Całkowanie funkcji o wartościach w przestrzeniach lokalnie wypukłych}
\label{sec:integral-lcs15a}

\providecommand{\atlasref}[2]{#2}


W \hyperref[sec:gateaux-frechet14]{rozdziale~14} udowodniliśmy podstawowe twierdzenie
Newtona--Leibniza dla krzywych w przestrzeni Banacha, potrzebne w rachunku
różniczkowym. Teraz budujemy szerszą teorię: obiektem całkowanym może być
funkcja przyjmująca wartości w przestrzeni lokalnie wypukłej, której topologia
nie musi pochodzić od jednej normy. Nie należy utożsamiać tego zagadnienia
z całkowaniem funkcji mierzalnych: całki Bochnera, Pettisa, Dunforda i Gelfanda
będą przedmiotem \hyperref[sec:integral-vector15b]{modułu~15B}.

W całym module $\K\in\{\R,\C\}$, a $E$ jest przestrzenią liniowo-topologiczną
Hausdorffa nad $\K$, lokalnie wypukłą, z topologią wyznaczoną przez rodzinę
ciągłych seminorm $\mathcal P$. \emph{Zupełność ciągowa} oznacza, że każdy
ciąg Cauchy'ego w $E$ ma granicę w $E$; nie mylimy jej z ogólną zupełnością
względem sieci Cauchy'ego. Jeśli nie zaznaczono inaczej, $a<b$ są liczbami
rzeczywistymi. Funkcja na odcinku jest rozumiana z topologią podprzestrzeni.

\subsection{Sumy Riemanna i twierdzenie o istnieniu}

\noindent\emph{Powiązanie tematyczne w Bibliotece matematycznej:} \atlasref{AM1-29::ma1:def:riemann-sum}{AM1-29: sumy Riemanna dla funkcji skalarnych}.


\begin{definicja}[Całka Riemanna o wartościach w przestrzeni liniowo-topologicznej]
\label{def:riemann-lcs15a}
Niech $a<b$ będą liczbami rzeczywistymi, $\K\in\{\R,\C\}$,
$E$ przestrzenią lokalnie wypukłą Hausdorffa nad $\K$,
$\mathcal P$ rodziną ciągłych seminorm wyznaczających jej
topologię i $f:[a,b]\to E$ funkcją. \emph{Podziałem z punktami pośrednimi} nazywamy
parę $P=(a=t_0<\cdots<t_N=b;\,\xi_j\in[t_{j-1},t_j])$.
Piszemy $|P|=\max_j(t_j-t_{j-1})$ i
$S(f,P)=\sum_{j=1}^N(t_j-t_{j-1})f(\xi_j)\in E$,
gdzie $N\ge1$ i $1\le j\le N$ w warunku na punkty pośrednie.
Funkcja $f$ jest całkowalna w sensie Riemanna w $E$, jeżeli istnieje $I\in E$
takie, że dla każdej seminormy $p\in\mathcal P$ i $\varepsilon>0$ istnieje
$\delta>0$ taka, iż $p(S(f,P)-I)<\varepsilon$ dla \emph{każdego} podziału
z punktami pośrednimi spełniającego $|P|<\delta$.
Warunek jest równoważny zbieżności sum Riemanna
w topologii $E$, ponieważ skończone przecięcia kul seminorm
tworzą bazę otoczeń zera. W szczególności zbieżność zachodzi
względem każdej ciągłej seminormy, nie tylko seminorm z $\mathcal P$.
Granica jest jednoznaczna w przestrzeni Hausdorffa.
Wówczas oznaczamy $I=\int_a^b f(t)\,\mathrm dt$.
Dla funkcji określonej na odcinku o końcach $a,b$ przy
$a=b$ przyjmujemy całkę równą zeru, a przy $b<a$ definiujemy
$\int_a^b f(t)\,\mathrm dt=-\int_b^a f(t)\,\mathrm dt$,
jeżeli całka po prawej stronie istnieje.
\end{definicja}

\begin{uwaga}[Znaczenie topologii]
Definicja wymaga zbieżności \emph{w topologii $E$}, a nie tylko po
zastosowaniu każdego funkcjonału $\ell\in E'$. Tę drugą własność nazywamy
słabą całkowalnością skalarną i sama nie gwarantuje istnienia wektora całki.
Powyższa granica jest granicą rodziny wszystkich sum Riemanna przy
$|P|\to0$; nie sprowadzamy jej z definicji do jednego wybranego ciągu podziałów.
\end{uwaga}

\begin{twierdzenie}[Istnienie całki przy zupełności ciągowej]
\label{thm:riemann-existence15a}
Niech $a<b$ będą liczbami rzeczywistymi, $\K\in\{\R,\C\}$,
a $E$ zupełną ciągowo przestrzenią lokalnie wypukłą Hausdorffa
nad $\K$. Wtedy każda funkcja ciągła $f:[a,b]\to E$
jest całkowalna w sensie Riemanna. W szczególności wynik
zachodzi w każdej zupełnej przestrzeni lokalnie wypukłej Hausdorffa.
\end{twierdzenie}
\begin{proof}
Dla $p\in\mathcal P$ definiujemy moduł ciągłości
$\omega_p(r)=\sup\{p(f(t)-f(s)):s,t\in[a,b],\ |s-t|\le r\}$.
Ciągłość $f$ na zwartym odcinku oznacza jednostajną ciągłość
w topologii $E$, zatem $\omega_p(r)\to0$ dla $r\downarrow0$.

Niech $P,Q$ będą dwoma podziałami z punktami pośrednimi, o średnicach
co najwyżej $r$. Rozważmy wspólne rozdrobnienie ich podziałów
przedziału i wybierzmy na nim dowolne punkty pośrednie, otrzymując $R$.
Każdy przedział $R$ mieści się w jednym przedziale $P$, a odpowiadające
mu punkty pośrednie są odległe o najwyżej $r$; analogicznie dla $Q$.
Po rozpisaniu sum względem wspólnego rozdrobnienia dostajemy
\[
 p(S(f,P)-S(f,R))\le(b-a)\omega_p(r),\qquad
 p(S(f,Q)-S(f,R))\le(b-a)\omega_p(r).
\]
Zatem $p(S(f,P)-S(f,Q))\le2(b-a)\omega_p(r)$.

Wybierzmy ciąg podziałów dyadycznych $P_n$, na przykład ze środkami
przedziałów jako punktami pośrednimi, o średnicach dążących do zera.
Powyższa nierówność pokazuje, że $(S(f,P_n))$ jest ciągiem Cauchy'ego
w $E$. Z zupełności ciągowej istnieje $I\in E$, do którego zbiega.
Dla dowolnego $P$ o średnicy mniejszej od $r$ stosujemy ostatnie
oszacowanie do $P$ i $P_n$ przy dostatecznie dużym $n$, a następnie
przechodzimy do granicy $n\to\infty$:
$p(S(f,P)-I)\le2(b-a)\omega_p(r)$.
Ponieważ dotyczy to każdej seminormy, zachodzi warunek z definicji.
Jeżeli $I,J\in E$ spełniają warunek z definicji, to dla
każdego $p\in\mathcal P$ i $\varepsilon>0$ wybieramy podział
$P$ tak drobny, że $p(I-S(f,P))<\varepsilon/2$ i
$p(S(f,P)-J)<\varepsilon/2$. Z nierówności trójkąta
$p(I-J)<\varepsilon$, więc $p(I-J)=0$. Ponieważ $\mathcal P$
rozdziela punkty, mamy $I=J$.
\end{proof}

\begin{stwierdzenie}[Liniowość i oszacowania seminormowe]
\label{prop:riemann-seminorm15a}
Niech $a<b$ będą rzeczywiste, $\K\in\{\R,\C\}$,
$E$ przestrzenią lokalnie wypukłą Hausdorffa nad $\K$,
a $f,g:[a,b]\to E$ funkcjami ciągłymi i całkowalnymi w $E$.
Dla $\alpha,\beta\in\K$ funkcja $\alpha f+\beta g$ jest
całkowalna i zachodzi pierwsza z poniższych równości.
Druga zachodzi dla $c\in[a,b]$, jeśli ograniczenia $f$ do $[a,c]$
i $[c,b]$ są całkowalne w $E$ (w szczególności gdy $E$ jest zupełna ciągowo):
\[
 \int_a^b(\alpha f(t)+\beta g(t))\,\mathrm dt
 =\alpha\int_a^b f(t)\,\mathrm dt+\beta\int_a^b g(t)\,\mathrm dt,
 \qquad \int_a^b f(t)\,\mathrm dt=\int_a^c f(t)\,\mathrm dt
 +\int_c^b f(t)\,\mathrm dt.
\]
Dla stałej funkcji o wartości $v\in E$ zachodzi
$\int_a^b v\,\mathrm dt=(b-a)v$.
Jeżeli $E$ jest zupełna ciągowo, to dla wszystkich
$u,v,w\in[a,b]$ zachodzi ponadto addytywność całek z orientacją:
$\int_u^v f(t)\,\mathrm dt+\int_v^w f(t)\,\mathrm dt
=\int_u^w f(t)\,\mathrm dt$.
Dla każdej ciągłej seminormy $p$ na $E$ mamy
\[
 p\left(\int_a^b f(t)\,\mathrm dt\right)
 \le\int_a^b p(f(t))\,\mathrm dt
 \le (b-a)\sup_{t\in[a,b]}p(f(t)).
\]
\end{stwierdzenie}
\begin{proof}
Oznaczmy całki $f,g$ przez $I_f,I_g$. Z liniowości sum
Riemanna otrzymujemy
$p(S(\alpha f+\beta g,P)-(\alpha I_f+\beta I_g))
\le|\alpha|p(S(f,P)-I_f)+|\beta|p(S(g,P)-I_g)$.
Dla każdej seminormy $p$ prawa strona dąży do zera
wraz ze średnicą podziału, co daje istnienie całki
kombinacji liniowej i wskazaną równość.
Jeżeli $a<c<b$, wybieramy ciągi podziałów odcinków $[a,c]$
i $[c,b]$ o średnicach dążących do zera. Ich połączenia są
podziałami $[a,b]$ o średnicach dążących do zera, a odpowiadające
sumy Riemanna są sumami dwóch sum na pododcinkach.
Przechodząc do granicy, uzyskujemy addytywność.
Dla $c=a$ lub $c=b$ równość wynika z konwencji dla całki
na odcinku zdegenerowanym. Dla funkcji stałej każda suma
Riemanna wynosi $(b-a)v$, co dowodzi wzoru.
Przy zupełności ciągowej z twierdzenia o istnieniu mamy
całki na wszystkich pododcinkach. Oznaczmy
$H(x)=\int_a^x f(t)\,\mathrm dt$. Addytywność dla
$a\le u\le v\le b$ daje $\int_u^v f(t)\,\mathrm dt=H(v)-H(u)$.
Na mocy konwencji orientacji tożsamość zachodzi również
dla $v<u$ oraz $u=v$. Dla dowolnych $u,v,w$ sumujemy
dwie takie różnice, otrzymując żądaną addytywność.
Nierówność trójkąta i dodatnia jednorodność seminormy dają
$p(S(f,P))\le\sum_j(t_j-t_{j-1})p(f(\xi_j))$.
Po przejściu do granicy lewa strona zbiega do
$p(\int_a^b f(t)\,\mathrm dt)$, a prawa do skalarnej
całki funkcji ciągłej $p\circ f$. Ostatnia nierówność jest skalarna.
\end{proof}

\begin{twierdzenie}[Przemienność z odwzorowaniami liniowymi ciągłymi]
\label{thm:riemann-functorial15a}
Niech $a<b$ będą rzeczywiste, $\K\in\{\R,\C\}$,
$E,F$ przestrzeniami lokalnie wypukłymi Hausdorffa nad $\K$,
$f:[a,b]\to E$ funkcją ciągłą i całkowalną w $E$, a
$T:E\to F$ odwzorowaniem liniowym ciągłym. Wtedy $T\circ f$
jest całkowalna w $F$ i
\[
 T\left(\int_a^b f(t)\,\mathrm dt\right)
 =\int_a^b T(f(t))\,\mathrm dt.
\]
W szczególności równość zachodzi dla każdego $\ell\in E'$, a
całkę w $E$ wyznaczają jednoznacznie jej wartości skalarne
$\ell(\int_a^b f(t)\,\mathrm dt)$ dla $\ell\in E'$.
Nie zakładamy zupełności $F$.
\end{twierdzenie}
\begin{proof}
Dla każdego $P$ mamy $S(T\circ f,P)=T(S(f,P))$.
Ciągłość $T$ przenosi granicę sum na $T(I)$; ten wektor jest
całką w $F$, nawet jeśli $F$ nie jest zupełna.
Funkcjonały z $E'$ rozdzielają punkty $E$ na mocy
twierdzenia Hahna--Banacha i hausdorffowskości $E$,
więc równości wszystkich obrazów skalarnych wyznaczają $I$.
\end{proof}

\begin{wniosek}[Środek ciężkości krzywej]
\label{cor:barycenter15a}
Niech $a<b$ będą rzeczywiste, $\K\in\{\R,\C\}$,
$E$ zupełną ciągowo przestrzenią lokalnie wypukłą Hausdorffa
nad $\K$, a $f:[a,b]\to E$ funkcją ciągłą. Wówczas
\[
 (b-a)^{-1}\int_a^b f(t)\,\mathrm dt
 \in\overline{\operatorname{conv}}\bigl(f([a,b])\bigr),
\]
gdzie domknięcie i otoczkę wypukłą rozumiemy w topologii $E$.
\end{wniosek}
\begin{proof}
Dla każdego podziału $P$ wektor $(b-a)^{-1}S(f,P)$ jest
kombinacją wypukłą wartości $f(\xi_j)$, ponieważ jego
współczynniki $(t_j-t_{j-1})/(b-a)$ są nieujemne i sumują się do $1$.
Zamkniętość wskazanego zbioru i zbieżność sum dowodzą tezy.
\end{proof}

\subsection{Rola zupełności i całka niewłaściwa}

\begin{przyklad}[Funkcja ciągła bez całki w przestrzeni niezupełnej]
\label{ex:noncomplete-integral15a}
Rozważmy $E=c_{00}(\N;\R)$, przestrzeń ciągów rzeczywistych
o skończonym nośniku z normą odziedziczoną z $\ell^2(\N;\R)$.
Niech $e_n$ oznacza $n$-ty wektor jednostkowy.
Obcięcia każdego ciągu z $\ell^2$ do pierwszych $N$ współrzędnych
zbiegają w normie $\ell^2$ do tego ciągu, ponieważ ogon
szeregu kwadratów modułów dąży do zera. Stąd $c_{00}$ jest
gęsta w $\ell^2$, lecz nie jest domknięta (np. obcięcia
$(2^{-n})_{n\ge1}$ zbiegają poza $c_{00}$).
Jest zatem niezupełna, a jej uzupełnieniem jest $\ell^2$.
Wybierzmy nieujemną, niezerową funkcję ciągłą $\eta:\R\to\R$
o zwartym nośniku zawartym w $(0,1)$.
Dla $t\in[0,1]$ niech
\[
 f(t)=\sum_{n=1}^{\infty}2^{-n}\eta(2^{n+1}t-1)e_n.
\]
Przedziały, na których kolejne składniki są niezerowe,
leżą w rozłącznych zbiorach $(2^{-(n+1)},2^{-n})$.
Dlatego dla każdego $t$ suma ma co najwyżej jeden niezerowy
wyraz i $f(t)\in c_{00}$. W otoczeniu każdego $t_0>0$ tylko skończenie wiele składników
może być niezerowych, gdyż prawe końce ich przedziałów dążą
do zera. Funkcja jest tam skończoną sumą funkcji ciągłych.
Mamy $f(0)=0$; dla $0<t<2^{-N}$ jedyny ewentualny
niezerowy składnik ma indeks $n\ge N$, więc
$\|f(t)\|_2\le2^{-N}\|\eta\|_\infty\to0$ przy $N\to\infty$.
To dowodzi ciągłości w zerze.
Całka istnieje w $\ell^2$, lecz jej $n$-ta współrzędna to
\[
 \int_0^1 2^{-n}\eta(2^{n+1}t-1)\,\mathrm dt
 =2^{-2n-1}\int_0^1\eta(s)\,\mathrm ds>0.
\]
Ma ona więc nieskończenie wiele niezerowych współrzędnych i nie
należy do $c_{00}$. Gdyby $f$ była całkowalna w $E$, przemienność
z ciągłymi funkcjonałami współrzędnych wymusiłaby tę samą
wartość o nieskończonym nośniku; sprzeczność.

Co więcej, definiując $g:[0,2]\to c_{00}$ przez $g(t)=f(t)$ dla
$t\in[0,1]$ oraz $g(t)=-f(2-t)$ dla $t\in[1,2]$, otrzymujemy
funkcję ciągłą, ponieważ $f(1)=0$. Jej całka w $\ell^2$ jest zerem: dla każdego ciągłego
funkcjonału $\ell$ podstawienie $s=2-t$ daje
$\int_1^2\ell(g(t))\,\mathrm dt=-\int_0^1\ell(f(s))\,\mathrm ds$.
Przemienność całek z funkcjonałami i rozdzielanie punktów
dają tę samą równość całek wektorowych, a addytywność
daje całkę zero. Sumy Riemanna $g$ zbiegają
więc do zera także w normie odziedziczonej przez $c_{00}$.
Jednak ograniczenie $g$ do $[0,1]$ nie jest całkowalne w $c_{00}$.
Pokazuje to konieczność dodatkowego założenia w addytywności całki
względem przedziału.
\end{przyklad}

\begin{stwierdzenie}[Całka w uzupełnieniu]
\label{prop:completion-integral15a}
Niech $a<b$ będą rzeczywiste, $\K\in\{\R,\C\}$,
$E$ przestrzenią lokalnie wypukłą Hausdorffa nad $\K$,
a $j:E\to\widehat E$ jej kanonicznym zanurzeniem do uzupełnienia.
Dla ciągłej $f:[a,b]\to E$ funkcja $j\circ f$ jest
całkowalna w $\widehat E$. Całka $f$ w $E$ istnieje wtedy
i tylko wtedy, gdy $\int_a^b j(f(t))\,\mathrm dt\in j(E)$;
wówczas $j(\int_a^b f(t)\,\mathrm dt)
=\int_a^b j(f(t))\,\mathrm dt$.
\end{stwierdzenie}
\begin{proof}
Uzupełnienie $\widehat E$ jest zupełne, zatem pierwsza teza
wynika z twierdzenia o istnieniu. Jeżeli całka w $E$ istnieje,
przemienność z $j$ daje konieczność. W drugą stronę załóżmy,
że całka w $\widehat E$ jest postaci $j(v)$ dla $v\in E$.
Rodzina sum Riemanna należy do $j(E)$ i zbiega w $\widehat E$
do $j(v)$. Topologia $E$ jest topologią indukowaną przez
zanurzenie $j$, więc sumy te zbiegają do $v$ w $E$.
\end{proof}

\begin{twierdzenie}[Kryterium całki niewłaściwej]
\label{thm:improper-lcs15a}
Niech $a\in\R$, $\K\in\{\R,\C\}$, $E$ będzie przestrzenią
lokalnie wypukłą Hausdorffa nad $\K$, zupełną ciągowo,
a $f:[a,\infty)\to E$ funkcją ciągłą.
Załóżmy, że dla każdej ciągłej seminormy $p$ zachodzi
$\int_a^\infty p(f(t))\,\mathrm dt<\infty$ w sensie
niewłaściwej całki skalarnej. Wówczas istnieje
$\lim_{R\to\infty}\int_a^R f(t)\,\mathrm dt\in E$, oznaczana przez
$\int_a^\infty f(t)\,\mathrm dt$, oraz
$p(\int_a^\infty f(t)\,\mathrm dt)
\le\int_a^\infty p(f(t))\,\mathrm dt$.
\end{twierdzenie}
\begin{proof}
Dla $R'>R\ge a$ z oszacowania seminormowego i addytywności
całki wynika
$p(\int_a^{R'}f(t)\,\mathrm dt-\int_a^R f(t)\,\mathrm dt)
\le\int_R^{R'}p(f(t))\,\mathrm dt$.
Ogon ostatniej całki dąży do zera, dlatego ciąg
$\bigl(\int_a^{a+n}f(t)\,\mathrm dt\bigr)_{n\ge1}$
jest ciągiem Cauchy'ego w $E$ i ma granicę $I$.
Dla ustalonych $p$ i $\varepsilon>0$ wybieramy $n$ tak duże,
że $p(\int_a^{a+n}f(t)\,\mathrm dt-I)<\varepsilon/2$
oraz $\int_{a+n}^\infty p(f(t))\,\mathrm dt<\varepsilon/2$.
Dla każdego rzeczywistego $R\ge a+n$ mamy zatem
$p(\int_a^R f(t)\,\mathrm dt-I)<\varepsilon$ z nierówności
trójkąta i oszacowania ogona. To daje granicę dla rzeczywistego
parametru $R\to\infty$. Dla każdego $R$ zachodzi
$p(\int_a^R f(t)\,\mathrm dt)\le\int_a^R p(f(t))\,\mathrm dt$;
ciągłość seminormy i granica całki skalarnej dają oszacowanie końcowe.
\end{proof}

\subsection{Twierdzenie Newtona--Leibniza i rachunek z parametrem}

\begin{twierdzenie}[Podstawowe twierdzenie rachunku całkowego]
\label{thm:ftc-lcs15a}
Niech $a<b$ będą rzeczywiste, a $E$ zupełną ciągowo
przestrzenią lokalnie wypukłą Hausdorffa nad $\K\in\{\R,\C\}$.
Krzywą klasy $C^1$ na $[a,b]$ rozumiemy jako krzywą ciągłą
z pochodną w topologii $E$ wewnątrz przedziału i pochodnymi
jednostronnymi na końcach, tworzącymi ciągłą funkcję
na $[a,b]$. Niech $f:[a,b]\to E$ będzie ciągła. Wtedy funkcja
$F:[a,b]\to E$, $F(x)=\int_a^x f(t)\,\mathrm dt$, jest klasy
$C^1$ na $[a,b]$ w powyższym sensie oraz
$F'(x)=f(x)$ dla każdego $x\in[a,b]$ (z odpowiednimi
pochodnymi jednostronnymi).
Jeżeli $G:[a,b]\to E$ jest klasy $C^1$, to
\[
 G(b)-G(a)=\int_a^b G'(t)\,\mathrm dt.
\]
W szczególności jest to uogólnienie wyniku z \hyperref[sec:gateaux-frechet14]{rozdziału~14},
a nie jego ponowny dowód w przypadku przestrzeni Banacha.
\end{twierdzenie}
\begin{proof}
Dla $x,x+h\in[a,b]$ i $h\ne0$ z addytywności całki mamy
$h^{-1}(F(x+h)-F(x))-f(x)
=h^{-1}\int_x^{x+h}(f(t)-f(x))\,\mathrm dt$.
Oszacowanie seminormowe (również dla $h<0$ po zmianie
orientacji całki) daje
\[
 p\left(h^{-1}(F(x+h)-F(x))-f(x)\right)
 \le\sup_{\substack{t\in[a,b]\\|t-x|\le|h|}}p(f(t)-f(x))
 \longrightarrow0.
\]
Stąd pochodna istnieje w topologii $E$ i jest równa $f(x)$.
Istnienie skończonej pochodnej (również jednostronnej)
zapewnia ciągłość $F$ w $x$: przyrost jest iloczynem $h$
i ilorazu różnicowego zbieżnego do $f(x)$.
Pochodna jest ciągła na $[a,b]$, bo $f$ jest ciągła.

Dla drugiej tezy nie musimy odtwarzać dowodu z przestrzeni Banacha.
Dla każdego $\ell\in E'$ złożenie $\ell\circ G$ jest skalarne
klasy $C^1$ z pochodną $\ell(G'(t))$. Skalarne twierdzenie
Newtona--Leibniza i przemienność całki z $\ell$ dają
$\ell(G(b)-G(a))=\ell(\int_a^b G'(t)\,\mathrm dt)$.
Ponieważ $E'$ rozdziela punkty, zachodzi żądana równość w $E$.
\end{proof}

\begin{twierdzenie}[Różniczkowanie całki zależnej od parametru]
\label{thm:parametric-integral15a}
Niech $a<b$ będą rzeczywiste, $\K\in\{\R,\C\}$,
$E$ zupełną ciągowo przestrzenią lokalnie wypukłą Hausdorffa
nad $\K$, $J\subseteq\R$ otwartym przedziałem i
$f:[a,b]\times J\to E$ funkcją ciągłą. Załóżmy, że dla
każdego $(t,s)\in[a,b]\times J$ istnieje pochodna
$\partial_s f(t,s)$ w topologii $E$ i że funkcja
$(t,s)\mapsto\partial_s f(t,s)$ jest ciągła.
Wtedy $F(s)=\int_a^b f(t,s)\,\mathrm dt$ jest funkcją klasy $C^1$
o wartościach w $E$ oraz
$F'(s)=\int_a^b\partial_s f(t,s)\,\mathrm dt$ dla $s\in J$.
\end{twierdzenie}
\begin{proof}
Ustalmy $s\in J$ i wybierzmy $\delta>0$ tak, aby
$[s-\delta,s+\delta]\subset J$. Dla każdej ciągłej seminormy $p$
mamy $p(F(s+h)-F(s))\le(b-a)\sup_{t\in[a,b]}
p(f(t,s+h)-f(t,s))\to0$ dzięki jednostajnej ciągłości
na zwartym prostokącie. Dla $0<|h|<\delta$ stosujemy
wektorowe twierdzenie Newtona--Leibniza do krzywej
$u\mapsto f(t,u)$, otrzymując
$f(t,s+h)-f(t,s)=\int_s^{s+h}\partial_s f(t,u)\,\mathrm du$.
Po odjęciu $h\partial_s f(t,s)$, podzieleniu przez $h$
i użyciu oszacowania seminormowego dostajemy
\[
 p\left(\frac{f(t,s+h)-f(t,s)}h-\partial_s f(t,s)\right)
 \le\sup_{u\text{ między }s\text{ i }s+h}
 p(\partial_s f(t,u)-\partial_s f(t,s)).
\]
Z liniowości całki i oszacowania seminormowego mamy
\[
 p\left(\frac{F(s+h)-F(s)}h
       -\int_a^b\partial_s f(t,s)\,\mathrm dt\right)
 \le(b-a)\sup_{\substack{t\in[a,b]\\u\text{ między }s\text{ i }s+h}}
 p(\partial_s f(t,u)-\partial_s f(t,s)).
\]
Jednostajna ciągłość $\partial_s f$ na wybranym zwartym
prostokącie daje zbieżność prawej strony do zera, więc
pochodna ma wskazaną postać. Dla jej ciągłości stosujemy
$p(F'(s+h)-F'(s))\le(b-a)\sup_{t\in[a,b]}
p(\partial_s f(t,s+h)-\partial_s f(t,s))\to0$.
\end{proof}


\begin{twierdzenie}[Reguła Leibniza z ruchomymi granicami]
\label{thm:leibniz-moving15a}
Niech $\K\in\{\R,\C\}$, $E$ będzie zupełną ciągowo
przestrzenią lokalnie wypukłą Hausdorffa nad $\K$,
$I,J$ otwartymi przedziałami rzeczywistymi, $f:I\times J\to E$ ciągłą funkcją
z pochodną $\partial_s f$ istniejącą w topologii $E$
w każdym punkcie i ciągłą na $I\times J$, a $a,b:J\to I$
funkcjami klasy $C^1$. Całkę z odwróconymi granicami
rozumiemy jako ujemną całkę w porządku rosnącym.
Wtedy $F(s)=\int_{a(s)}^{b(s)}f(t,s)\,\mathrm dt$
jest klasy $C^1$ w topologii $E$ i
\[
 F'(s)=f(b(s),s)b'(s)-f(a(s),s)a'(s)
       +\int_{a(s)}^{b(s)}\partial_s f(t,s)\,\mathrm dt.
\]
\end{twierdzenie}
\begin{proof}
Ustalmy $s\in J$. Wybieramy zwarty przedział $K\subset I$
zawierający $a(s),b(s)$ we wnętrzu i $\delta>0$ tak, że
$[s-\delta,s+\delta]\subset J$ i dla $|h|\le\delta$ obie
granice $a(s+h),b(s+h)$ należą do $K$.
Jest to możliwe dzięki ciągłości $a,b$ i otwartości $I,J$. Wszystkie całki
istnieją; addytywność całek z orientacją daje
\[
 \begin{split}
 F(s+h)-F(s)={}&
 \int_{a(s)}^{b(s)}(f(t,s+h)-f(t,s))\,\mathrm dt\\
 &+\int_{b(s)}^{b(s+h)} f(t,s+h)\,\mathrm dt
 -\int_{a(s)}^{a(s+h)} f(t,s+h)\,\mathrm dt.
 \end{split}
\]
Pierwszy składnik po podzieleniu przez $h$ dąży
do całki pochodnej na mocy twierdzenia o stałych
granicach (po odwróceniu granic zmieniamy znak,
a przy $a(s)=b(s)$ obie całki są zerowe).
Dla dowolnej ciągłej seminormy $p$
różnica drugiego ilorazu od
$((b(s+h)-b(s))/h)f(b(s),s)$ ma seminormę co najwyżej
\[
 \left|\frac{b(s+h)-b(s)}h\right|
 \sup_{t\text{ między }b(s),b(s+h)}
 p(f(t,s+h)-f(b(s),s)).
\]
Pierwszy czynnik pozostaje ograniczony, a drugi
dąży do zera z ciągłości $f$ w punkcie $(b(s),s)$.
Analogiczne oszacowanie dla $a$ daje wzór pochodnej,
również gdy granice są równe lub zmieniają kolejność.
Ciągłość odwzorowań $u\mapsto f(b(u),u)b'(u)$
i $u\mapsto f(a(u),u)a'(u)$ wynika z ciągłości $f,a,b,a',b'$
i ciągłości mnożenia przez skalar w $E$. Dla ostatniego oznaczmy $g(t,u)=\partial_s f(t,u)$ oraz
$G(u)=\int_{a(u)}^{b(u)}g(t,u)\,\mathrm dt$.
Rozbicie przyrostu jak wyżej i oszacowanie seminormowe dają
\[
\begin{split}
 p(G(s+h)-G(s))\le{}&
 |b(s)-a(s)|\sup_{t\in K}p(g(t,s+h)-g(t,s))\\
 &+\bigl(|b(s+h)-b(s)|+|a(s+h)-a(s)|\bigr)M_p,
\end{split}
\]
gdzie $M_p=\sup_{(t,u)\in K\times[s-\delta,s+\delta]}p(g(t,u))$
jest skończone przez zwartość. Pierwszy składnik dąży
do zera dzięki jednostajnej ciągłości w seminormie $p$,
drugi dzięki ciągłości $a,b$. To daje ciągłość $G$ w $E$.
\end{proof}

\begin{twierdzenie}[Wielokrotne różniczkowanie całki z parametrem]
\label{thm:higher-parametric15a}
Niech $a<b$ będą rzeczywiste, $m\ge1$, $\K\in\{\R,\C\}$,
$E$ zupełną ciągowo przestrzenią lokalnie wypukłą Hausdorffa
nad $\K$ i $J\subseteq\R$ otwartym przedziałem. Jeżeli $f:[a,b]\times J\to E$ jest ciągła, a każda
pochodna $\partial_s^j f$, $1\le j\le m$, istnieje
w topologii $E$ i jest wspólnie ciągła względem $(t,s)$,
to
$F(s)=\int_a^b f(t,s)\,\mathrm dt$ jest klasy $C^m$ i
$F^{(j)}(s)=\int_a^b\partial_s^j f(t,s)\,\mathrm dt$
dla $1\le j\le m$.
\end{twierdzenie}
\begin{proof}
Przypadek $j=1$ jest twierdzeniem o różniczkowaniu
całki z parametrem. Jeżeli wzór zachodzi dla $j<m$,
stosujemy to samo twierdzenie do funkcji $\partial_s^j f$,
której pochodna względem $s$ jest ciągła na mocy
założeń. Otrzymujemy wzór dla $j+1$ i ciągłość tej
pochodnej. Indukcja kończy dowód.
\end{proof}


\begin{twierdzenie}[Twierdzenie Fubiniego dla funkcji ciągłych]
\label{thm:fubini-continuous15a}
Niech $a<b$ i $c<d$ będą liczbami rzeczywistymi,
$\K\in\{\R,\C\}$, $E$ zupełną ciągowo przestrzenią
lokalnie wypukłą Hausdorffa nad $\K$ i
$f:[a,b]\times[c,d]\to E$ funkcją ciągłą. Wówczas obie funkcje
$s\mapsto\int_c^d f(s,t)\,\mathrm dt$ oraz
$t\mapsto\int_a^b f(s,t)\,\mathrm ds$ są ciągłe i
\[
 \int_a^b\left(\int_c^d f(s,t)\,\mathrm dt\right)\mathrm ds
 =\int_c^d\left(\int_a^b f(s,t)\,\mathrm ds\right)\mathrm dt.
\]
Nie zakładamy mierzalności Bochnerowskiej ani twierdzenia Fubiniego
dla ogólnych funkcji mierzalnych.
\end{twierdzenie}
\begin{proof}
Dla każdej seminormy $p$ jednostajna ciągłość $f$
na zwartym prostokącie daje ciągłość obydwu funkcji
całek iterowanych: na przykład
$p(\int_c^d(f(s,t)-f(s_0,t))\,\mathrm dt)
\le(d-c)\sup_{t\in[c,d]}p(f(s,t)-f(s_0,t))$.

Niech $P=(s_j;\xi_j)$ i $Q=(t_k;\eta_k)$ będą podziałami
z punktami pośrednimi odpowiednio odcinków $[a,b]$ i $[c,d]$;
piszemy $\Delta s_j=s_j-s_{j-1}$ i $\Delta t_k=t_k-t_{k-1}$.
Dla $p$ oznaczmy przez $\omega_p^{(1)}(r)$ i $\omega_p^{(2)}(r)$
suprema seminormy różnicy wartości $f$ przy zmianie tylko
odpowiedniej zmiennej o co najwyżej $r$.
Oba dążą do zera dla $r\downarrow0$ dzięki jednostajnej
ciągłości na prostokącie. Niech
$D(P,Q)=\sum_{j,k}\Delta s_j\Delta t_k f(\xi_j,\eta_k)$.
Dla $A(s)=\int_c^d f(s,t)\,\mathrm dt$ mamy
\[
\begin{split}
 p(S(A,P)-D(P,Q))
 &\le\sum_{j,k}\Delta s_j
 \int_{t_{k-1}}^{t_k}p(f(\xi_j,t)-f(\xi_j,\eta_k))\,\mathrm dt\\
 &\le(b-a)(d-c)\omega_p^{(2)}(|Q|).
\end{split}
\]
Istotnie, na każdym podprzedziale całkujemy różnicę
od funkcji stałej o wartości $f(\xi_j,\eta_k)$, po czym
stosujemy addytywność i oszacowanie seminormowe.
Dla $B(t)=\int_a^b f(s,t)\,\mathrm ds$ analogicznie
$p(S(B,Q)-D(P,Q))\le(b-a)(d-c)\omega_p^{(1)}(|P|)$.
Wybierzmy ciągi podziałów $P_n,Q_n$ o średnicach dążących
do zera. Sumy $S(A,P_n)$ i $S(B,Q_n)$ zbiegają do
odpowiednich całek iterowanych z twierdzenia o istnieniu
i ciągłości $A,B$. Powyższe oszacowania dają
$p(S(A,P_n)-S(B,Q_n))\to0$ dla każdej $p$.
Ciągłość seminorm i rozdzielanie punktów dają równość
całek iterowanych.
\end{proof}

\subsection{Przykłady funkcyjne i całkowanie rodzin operatorów}

\begin{przyklad}[Całkowanie w $C^\infty(\Omega)$]
\label{ex:c-infty-integral15a}
Niech $d\ge1$, $\K\in\{\R,\C\}$, $a<b$ będą rzeczywiste
i $\Omega\subset\R^d$ będzie niepustym zbiorem otwartym.
Przestrzeń $C^\infty(\Omega;\K)$ funkcji skalarnych
wyposażamy w seminormy
$p_{K,m}(u)=\max_{|\alpha|\le m}\sup_{x\in K}|\partial^\alpha u(x)|$,
gdzie $K\subseteq\Omega$ jest niepustym zbiorem zwartym,
a $m\in\N_0$.
Jest to zupełna przestrzeń lokalnie wypukła, co szczegółowo
udowodnimy w \hyperref[sec:gladkie-schwartz17]{rozdziale~17}. Jeżeli $t\mapsto u_t$ jest
ciągłe z $[a,b]$ do tej przestrzeni, to
$u=\int_a^b u_t\,\mathrm dt\in C^\infty(\Omega)$ i
\[
 \partial^\alpha u(x)
 =\int_a^b\partial^\alpha u_t(x)\,\mathrm dt
 \qquad(x\in\Omega,\ \alpha\in\N_0^d).
\]
Istotnie, $p_{K,m}(\partial^\alpha v)\le p_{K,m+|\alpha|}(v)$,
więc $\partial^\alpha$ jest ciągłym operatorem liniowym.
Ewaluacja $v\mapsto v(x)$ jest ciągła, ponieważ
$|v(x)|=p_{\{x\},0}(v)$. Z dwukrotnego zastosowania
przemienności całki z odwzorowaniem liniowym ciągłym
otrzymujemy wskazany wzór, bez zakładania różniczkowania
pod całką skalarną.
\end{przyklad}

\begin{przyklad}[Całkowanie w przestrzeni Schwartza]
\label{ex:schwartz-integral15a}
Niech $d\ge1$, $\K\in\{\R,\C\}$ i $a<b$ będą rzeczywiste.
Niech $\mathcal S(\R^d)=\mathcal S(\R^d;\K)$ oznacza
przestrzeń funkcji gładkich o wartościach w $\K$,
dla których $q_{m,k}(u)=\max_{|\alpha|\le k}\sup_{x\in\R^d}
(1+|x|)^m|\partial^\alpha u(x)|$ jest skończone
przy wszystkich $m,k\in\N_0$; tę topologię również
omówimy w \hyperref[sec:gladkie-schwartz17]{rozdziale~17}. Przestrzeń ta jest zupełna.
Dla ciągłego $t\mapsto u_t\in\mathcal S(\R^d)$ całka
$u=\int_a^b u_t\,\mathrm dt$ należy do $\mathcal S(\R^d)$
i spełnia $q_{m,k}(u)\le\int_a^b q_{m,k}(u_t)\,\mathrm dt$.
Różniczkowanie jest na tej przestrzeni ciągłym
odwzorowaniem liniowym: dla każdego wielowskaźnika $\alpha$
seminorma $q_{m,k}(\partial^\alpha u)$ jest ograniczona
przez $q_{m,k+|\alpha|}(u)$. Stąd
$\partial^\alpha(\int_a^b u_t\,\mathrm dt)
=\int_a^b\partial^\alpha u_t\,\mathrm dt$
(jako równość w przestrzeni Schwartza).
\end{przyklad}

\begin{twierdzenie}[Silna całka rodziny operatorów]
\label{thm:strong-operator-integral15a}
Niech $a<b$ będą rzeczywiste, $X,Y$ przestrzeniami Banacha
nad wspólnym ciałem $\K\in\{\R,\C\}$ i
$T:[a,b]\to\mathcal L(X,Y)$ rodziną operatorów taką, że
$t\mapsto T(t)x$ jest ciągła w $Y$ dla każdego $x\in X$.
Wtedy wzór
$Ax=\int_a^b T(t)x\,\mathrm dt$ określa operator
$A\in\mathcal L(X,Y)$, a
\[
 \|A\|\le(b-a)\sup_{t\in[a,b]}\|T(t)\|<\infty.
\]
Nie wynika stąd ciągłość $t\mapsto T(t)$ w normie operatorowej.
\end{twierdzenie}
\begin{proof}
Dla każdego $x$ zbiór $\{T(t)x:t\in[a,b]\}$ jest ograniczony,
gdyż jest ciągłym obrazem zwartego odcinka. Twierdzenie
Banacha--Steinhausa zapewnia
$M=\sup_{t\in[a,b]}\|T(t)\|<\infty$.
Istnienie całki wektorowej w przestrzeni Banacha $Y$
wynika z twierdzenia o istnieniu, a liniowość z
liniowości całki. Dla $x\in X$ oszacowanie seminormowe
daje $\|Ax\|\le\int_a^b\|T(t)x\|\,\mathrm dt
\le(b-a)M\|x\|$, co dowodzi ciągłości $A$.
\end{proof}

\begin{przyklad}[Ciągłość silna nie oznacza ciągłości normowej]
\label{ex:strong-not-norm-cont15a}
W zespolonej przestrzeni Hilberta $H=L^2(\R,\lambda_1;\C)$ niech
$(U_t f)(x)=e^{itx}f(x)$ dla $f\in H$ i $t\in\R$.
Każdy $U_t$ jest liniową izometrią, ponieważ $|e^{itx}|=1$,
i $U_{t+s}=U_tU_s$, $U_0=I_H$. Dla każdego $f\in H$
i $t\to0$ mamy
$\|U_tf-f\|_2^2=\int_\R |e^{itx}-1|^2|f(x)|^2\,\mathrm d\lambda_1(x)
\to0$ z twierdzenia o zbieżności zmajoryzowanej, gdyż
funkcja podcałkowa jest ograniczona przez $4|f|^2$.
Dla dowolnego $s\in\R$ równość
$\|U_{s+t}f-U_sf\|_2=\|U_tf-f\|_2$ daje silną ciągłość w $s$.
Jednak dla $t\ne0$ mamy $\|U_t-I_H\|\le2$ z nierówności
trójkąta. Dla dowolnego $0<\varepsilon<2$ wybieramy
ograniczony otwarty przedział $A$ zawierający $\pi/t$,
na którym $|e^{itx}-1|>2-\varepsilon$; istnieje on z
ciągłości mnożnika i jego wartości w $\pi/t$.
Wektor $v=\chi_A/\sqrt{\lambda_1(A)}$ ma normę $1$, a
$\|(U_t-I_H)v\|_2^2\ge(2-\varepsilon)^2$.
Przechodząc z $\varepsilon\downarrow0$, otrzymujemy
$\|U_t-I_H\|=2$ dla każdego $t\ne0$.
Można więc całkować wektory $U_tf$ względem $t$ w $H$,
ale nie wolno bez dodatkowego dowodu zakładać ciągłości
$t\mapsto U_t$ w normie $\mathcal L(H)$.
\end{przyklad}

\begin{uwaga}[Granice zakresu]
Definicja sum Riemanna wymaga jedynie topologii i działań liniowych,
ale powyższe twierdzenia o istnieniu i oszacowaniach wykorzystują
lokalną wypukłość, seminormy i zupełność ciągową. W przestrzeni
niezupełnej ciągłość nie gwarantuje istnienia całki; w ogólnej
przestrzeni liniowo-topologicznej nielokalnie wypukłej nie można
bez uzasadnienia stosować oszacowań przez rodzinę seminorm.
Dla funkcji mierzalnych konieczne są dodatkowe pojęcia mierzalności;
teorię tę rozwiniemy osobno w \hyperref[sec:integral-vector15b]{module~15B}.
\end{uwaga}

\addtocounter{section}{-1}
\endgroup

% Wersja źródłowa: audyt II, jednostka 15B
\setcounter{section}{15}
\noindent\textbf{Wymagane wyniki.} Wykorzystujemy podstawy skalarnej całki Lebesgue'a, twierdzenia o zbieżności zmajoryzowanej i Fubiniego, twierdzenia Hahna--Banacha i o wykresie domkniętym oraz definicje przestrzeni Banacha i jej przestrzeni sprzężonej.\medskip
% \hyperref[sec:integral-vector15b]{Moduł 15B}; umieszczony po 15A, bez przesuwania \hyperref[sec:wariacyjny16]{rozdziałów 16}--33.
\begingroup
\renewcommand{\thesection}{15B}
\renewcommand{\theHsection}{15B}
\section{Całki Bochnera, Pettisa, Dunforda i Gelfanda oraz miary wektorowe}
\label{sec:integral-vector15b}

W \hyperref[sec:integral-lcs15a]{module 15A} całkowaliśmy \emph{funkcje ciągłe na odcinku} w topologii
przestrzeni lokalnie wypukłej. Teraz dopuszczamy funkcje mierzalne;
różne pojęcia całki wynikają z tego, czy granica istnieje w normie,
czy najpierw całkujemy wszystkie obrazy skalarne, i w jakiej przestrzeni
musi leżeć wektor reprezentujący te całki. Nie powtarzamy konstrukcji
skalarnej całki Lebesgue'a.

W całym module $(S,\Sigma,\mu)$ jest \emph{zupełną przestrzenią z miarą
skończoną}, $\K\in\{\R,\C\}$, zaś $X$ jest przestrzenią Banacha nad
$\K$, chyba że jawnie przyjmiemy inne założenia. Przez $X'$ i $X''$
oznaczamy odpowiednio przestrzeń sprzężoną i podwójnie sprzężoną; $J_X:X\to X''$
jest zanurzeniem kanonicznym danym przez $J_X(x)(x')=x'(x)$.
W przestrzeniach $L^1$ utożsamiamy funkcje równe prawie wszędzie.
W wypowiedziach o wartościach punktowych i całce Riemanna
rozpatrujemy konkretne funkcje, a nie klasy równoważności. Dla $A\in\Sigma$ zapis
$\int_A f(s)\,\mathrm d\mu(s)$ oznacza $\int_S\chi_A(s)f(s)\,\mathrm d\mu(s)$.

\subsection{Mierzalność silna i całka Bochnera}

\begin{definicja}[Mierzalność silna i słaba]
\label{def:measurability15b}
Niech $(S,\Sigma,\mu)$ będzie zupełną przestrzenią z miarą
skończoną, a $X$ przestrzenią Banacha nad $\K\in\{\R,\C\}$.
Funkcję $f:S\to X$ nazywamy \emph{silnie mierzalną}, jeżeli istnieją
funkcje proste $s_n=\sum_{j=1}^{m_n}\chi_{A_{n,j}}x_{n,j}$,
$m_n\ge1$, $A_{n,j}\in\Sigma$, $x_{n,j}\in X$, takie że $s_n(s)\to f(s)$
w normie dla prawie każdego $s\in S$. Nazywamy ją
\emph{słabo mierzalną}, jeżeli $x'\circ f:S\to\K$ jest
mierzalna dla każdego $x'\in X'$. Funkcję $f$ nazywamy
\emph{istotnie ośrodkowo-wartościową}, jeśli poza pewnym zbiorem
miary zero jej wartości leżą w ośrodkowej podprzestrzeni domkniętej
przestrzeni $X$.
\end{definicja}

\begin{twierdzenie}[Twierdzenie Pettisa o mierzalności]
\label{thm:pettis-meas15b}
Niech $(S,\Sigma,\mu)$ będzie zupełną przestrzenią z miarą
skończoną, a $X$ przestrzenią Banacha nad $\K\in\{\R,\C\}$.
Funkcja $f:S\to X$ jest silnie mierzalna wtedy i tylko wtedy, gdy
jest słabo mierzalna oraz istotnie ośrodkowo-wartościowa.
\end{twierdzenie}
\begin{proof}
Jeżeli $s_n\to f$ prawie wszędzie, to dla każdego $x'\in X'$
mamy $x'(s_n)\to x'(f)$, co daje mierzalność słabą.
Wartości wszystkich $s_n$ tworzą zbiór przeliczalny;
jego domknięta otoczka liniowa jest ośrodkową podprzestrzenią,
w której leży $f(s)$ poza wspólnym zbiorem miary zero.

Odwrotnie, po usunięciu zbioru miary zero można przyjąć, że $f$
przyjmuje wartości w ośrodkowej przestrzeni Banacha $Y\subset X$.
Jeżeli $Y=\{0\}$, funkcja jest równa zeru prawie wszędzie,
więc jest silnie mierzalna. Dalej zakładamy $Y\ne0$.
Wybierzmy gęsty ciąg $(y_j)$ w $Y$. Istnieje przeliczalna
rodzina $(\ell_k)$ elementów $X'$ o normie co najwyżej $1$
taka, że dla każdego $y\in Y$ zachodzi
$\|y\|=\sup_k|\ell_k(y)|$: otrzymujemy ją, dobierając
funkcjonały normujące gęsty zbiór sfery jednostkowej $Y$
i przedłużając je za pomocą twierdzenia Hahna--Banacha.
Dokładniej, dla gęstego ciągu $(u_k)$ w sferze jednostkowej
$Y$ wybieramy $\ell_k\in X'$ o normie $1$ i
$\ell_k(u_k)=1$. Dla dowolnego $y\ne0$ oraz $\varepsilon>0$
istnieje $k$ takie, że $\|u_k-y/\|y\|\|<\varepsilon$,
więc $|\ell_k(y)|\ge(1-\varepsilon)\|y\|$.
Przeciwne oszacowanie wynika z $\|\ell_k\|=1$;
dla $y=0$ równość jest natychmiastowa.
W konsekwencji $s\mapsto\|f(s)-y_j\|$ jest mierzalna,
ponieważ jest supremum przeliczalnej rodziny mierzalnych
funkcji skalarnych.

Dla $n\ge1$ definiujemy rozłączne, mierzalne zbiory $B_{n,j}$,
przypisując $s$ pierwszemu indeksowi $j$, dla którego
$\|f(s)-y_j\|<1/n$. Zbiory te pokrywają $S$ z dokładnością
do zbioru miary zero, a funkcja przyjmująca wartość $y_j$
na $B_{n,j}$ przybliża $f$ z błędem mniejszym od $1/n$.
Ponieważ $\mu(S)<\infty$, można wybrać $m_n$ tak, aby
$\mu(\bigcup_{j>m_n}B_{n,j})<2^{-n}$, i przyjąć wartość
$0$ na tym ogonie. Powstaje funkcja prosta $s_n$;
z lematu Borela--Cantellego otrzymujemy $s_n(s)\to f(s)$
prawie wszędzie. Jest to mierzalność silna.
\end{proof}

\begin{definicja}[Całka Bochnera]
\label{def:bochner15b}
Niech $(S,\Sigma,\mu)$ będzie zupełną przestrzenią z miarą
skończoną, a $X$ przestrzenią Banacha nad $\K\in\{\R,\C\}$.
Funkcję silnie mierzalną $f:S\to X$ nazywamy
\emph{całkowalną w sensie Bochnera}, jeśli istnieje ciąg
funkcji prostych $(s_n)$ taki, że
$\int_S\|f(s)-s_n(s)\|\,\mathrm d\mu(s)\to0$.
Wtedy definiujemy
$\int_S f(s)\,\mathrm d\mu(s)
=\lim_{n\to\infty}\int_S s_n(s)\,\mathrm d\mu(s)$,
gdzie całka funkcji prostej $\sum_{j=1}^m\chi_{A_j}x_j$
jest równa $\sum_{j=1}^m\mu(A_j)x_j$ dla parami rozłącznych
$A_j\in\Sigma$ i $x_j\in X$. Każda funkcja prosta ma
takie przedstawienie; przez wspólne rozdrobnienie dwóch
skończonych podziałów (z uwzględnieniem zbioru, na którym
funkcja ma wartość zero) i skończoną addytywność $\mu$
otrzymujemy niezależność tej sumy od przedstawienia.
Przez $L^1(S,\mu;X)$, również $L^1(S;X)$ przy ustalonej
mierze $\mu$, oznaczamy przestrzeń klas funkcji całkowalnych
w sensie Bochnera z normą
$\|f\|_{L^1}=\int_S\|f(s)\|\,\mathrm d\mu(s)$.
Istnienie całki i niezależność od wyboru aproksymacji
uzasadniamy w następnym twierdzeniu.
\end{definicja}

\begin{twierdzenie}[Kryterium całkowalności i oszacowanie]
\label{thm:bochner-criterion15b}
Niech $(S,\Sigma,\mu)$ będzie zupełną przestrzenią z miarą
skończoną, a $X$ przestrzenią Banacha nad $\K\in\{\R,\C\}$.
Funkcja $f:S\to X$ jest całkowalna w sensie Bochnera wtedy
i tylko wtedy, gdy jest silnie mierzalna i
$\int_S\|f(s)\|\,\mathrm d\mu(s)<\infty$.
Całka nie zależy od ciągu przybliżeń, jest liniowa i spełnia
\[
 \left\|\int_A f(s)\,\mathrm d\mu(s)\right\|
 \le\int_A\|f(s)\|\,\mathrm d\mu(s)\qquad(A\in\Sigma).
\]
Jeżeli $f$ jest całkowalna w sensie Bochnera, to dla każdej
przestrzeni Banacha $Y$ nad $\K$ i operatora
$T\in\mathcal L(X,Y)$ funkcja $Tf$ jest całkowalna w sensie
Bochnera oraz dla każdego $A\in\Sigma$ zachodzi
$T(\int_A f(s)\,\mathrm d\mu(s))=\int_A T(f(s))\,\mathrm d\mu(s)$.
\end{twierdzenie}
\begin{proof}
Dla funkcji prostych $h$ mamy
$\|\int_A h(s)\,\mathrm d\mu(s)\|\le\int_A\|h(s)\|\,\mathrm d\mu(s)$.
Jeżeli $\int_S\|f-s_n\|\,\mathrm d\mu\to0$, to
$\int_S\|s_n-s_m\|\,\mathrm d\mu\le
\int_S\|s_n-f\|\,\mathrm d\mu+\int_S\|f-s_m\|\,\mathrm d\mu$.
Stąd całki $s_n$ po $A$ są ciągiem Cauchy'ego w zupełnej
przestrzeni $X$. Takie samo oszacowanie dla dwóch różnych
ciągów aproksymujących $f$ daje niezależność granicy od
aproksymacji. Ograniczenie funkcji prostych do $A$ przez
mnożenie przez $\chi_A$ zachowuje zbieżność błędu w całce,
więc uzasadnia również istnienie całki po $A$.
Jeżeli $s_n$ i $t_n$ aproksymują odpowiednio $f$ i $h$,
to $\alpha s_n+\beta t_n$ aproksymują $\alpha f+\beta h$
dla $\alpha,\beta\in\K$; liniowość wynika z liniowości
całki funkcji prostych, wynikającej ze wspólnego
rozdrobnienia ich podziałów. Ponadto
$|\|f\|-\|s_n\||\le\|f-s_n\|$, stąd całkowalność normy
i zbieżność całek norm, co uzasadnia przejście do granicy
w oszacowaniu całki po $A$.

Jeśli z kolei $f$ jest silnie mierzalna i ma całkowalną normę,
wybierzmy funkcje proste $u_n\to f$ prawie wszędzie.
Funkcja $s_n$ jest równa $u_n$ tam, gdzie
$\|u_n(s)\|\le2\|f(s)\|$, a zeru w pozostałych punktach.
Oba zbiory są mierzalne, zatem $s_n$ jest funkcją prostą,
$\|s_n\|\le2\|f\|$, a $s_n\to f$ prawie wszędzie:
dla $f(s)\ne0$ nierówność definiująca $s_n$ jest
spełniona od pewnego miejsca, a dla $f(s)=0$ mamy
$s_n(s)=0$. Stąd $\|s_n-f\|\le3\|f\|$ i skalarne
twierdzenie o zbieżności zmajoryzowanej zapewnia
$\int_S\|f(s)-s_n(s)\|\,\mathrm d\mu(s)\to0$.
Dla operatora $T$ każda funkcja $Ts_n$ jest prosta,
a silna mierzalność $Tf$ wynika z ciągłości $T$ i silnej
mierzalności $f$. Ponadto
$\int\|Tf-Ts_n\|\le\|T\|\int\|f-s_n\|$, a równość
$T\int s_n=\int Ts_n$ przechodzi do granicy.
\end{proof}

\begin{twierdzenie}[Zbieżność zmajoryzowana Bochnera]
\label{thm:bochner-dct15b}
Niech $(S,\Sigma,\mu)$ będzie zupełną przestrzenią z miarą
skończoną, a $X$ przestrzenią Banacha nad $\K\in\{\R,\C\}$.
Niech $f:S\to X$ będzie funkcją, a $(f_n)_{n\ge1}$ ciągiem
silnie mierzalnych funkcji $f_n:S\to X$ takim, że
$f_n(s)\to f(s)$ w normie dla prawie każdego $s\in S$.
Jeżeli istnieje $g\in L^1(S,\mu)$, $g\ge0$, taka że
$\|f_n(s)\|\le g(s)$ prawie wszędzie dla wszystkich $n$,
to $f$ i $f_n$ są całkowalne w sensie Bochnera oraz
$\int_S\|f_n(s)-f(s)\|\,\mathrm d\mu(s)\to0$, a więc
$\int_S f_n(s)\,\mathrm d\mu(s)\to\int_S f(s)\,\mathrm d\mu(s)$ w normie.
\end{twierdzenie}
\begin{proof}
Granica prawie wszędzie funkcji silnie mierzalnych jest
silnie mierzalna: dla każdego $n$ wybieramy funkcję prostą $v_n$
taką, że $\mu(\{s:\|v_n(s)-f_n(s)\|>1/n\})<2^{-n}$.
Jest to możliwe dzięki zbieżności prawie wszędzie aproksymacji
prostych i skończoności miary. Lemat Borela--Cantellego daje
$\|v_n(s)-f_n(s)\|\to0$ prawie wszędzie, więc $v_n(s)\to f(s)$
prawie wszędzie. Ponadto $\|f\|\le g$ i
$\|f_n-f\|\le2g$ prawie wszędzie. Skalarne twierdzenie
o zbieżności zmajoryzowanej i kryterium całkowalności
Bochnera kończą dowód.
\end{proof}

\begin{twierdzenie}[Twierdzenie Fubiniego dla całki Bochnera]
\label{thm:bochner-fubini15b}
Niech $(S,\Sigma,\mu)$ i $(T,\mathcal T,\nu)$ będą
zupełnymi przestrzeniami z miarami skończonymi, a $X$
przestrzenią Banacha nad $\K\in\{\R,\C\}$.
Niech $f:S\times T\to X$ będzie silnie mierzalna względem
uzupełnienia miary produktowej i spełnia
$\int_{S\times T}\|f(s,t)\|\,\mathrm d(\mu\otimes\nu)(s,t)<\infty$.
Całki na $S\times T$ rozumiemy względem uzupełnienia
miary $\mu\otimes\nu$; ich wartości nie zależą od wyboru
przedstawiciela funkcji z dokładnością do zbioru miary zero.
Wtedy dla $\mu$-prawie każdego $s$ funkcja $f(s,\cdot)$
oraz dla $\nu$-prawie każdego $t$ funkcja $f(\cdot,t)$
są całkowalne w sensie Bochnera. Po przyjęciu wartości $0$
dla całek po przekrojach wyjątkowych obie funkcje otrzymane
przez całkowanie po jednej zmiennej są całkowalne w sensie
Bochnera i
\[
 \int_{S\times T}f\,\mathrm d(\mu\otimes\nu)
 =\int_S\left(\int_Tf(s,t)\,\mathrm d\nu(t)\right)\mathrm d\mu(s)
 =\int_T\left(\int_Sf(s,t)\,\mathrm d\mu(s)\right)\mathrm d\nu(t).
\]
\end{twierdzenie}
\begin{proof}
W pierwszej kolejności redukujemy zagadnienie do funkcji mierzalnej
względem nieuzupełnionej produktowej $\sigma$-algebry. Istotnie,
funkcja $f$ jest prawie wszędzie granicą funkcji prostych mierzalnych
względem uzupełnienia. Każdy z przeliczalnie wielu zbiorów
występujących w tych funkcjach różni się o zbiór zawarty w mierzalnym
zbiorze miary produktowej zero od pewnego zbioru produktowo
mierzalnego. Zastępując te zbiory odpowiednimi przedstawicielami,
otrzymujemy ciąg produktowo mierzalnych funkcji prostych $(u_n)$.
Poza jednym produktowo mierzalnym zbiorem $N$ miary zero
ciąg ten zbiega do $f$. Zbiór punktów, w których $(u_n)$
jest ciągiem Cauchy'ego, jest produktowo mierzalny, gdyż
można go zapisać jako
$\bigcap_{k\ge1}\bigcup_{M\ge1}\bigcap_{i,j\ge M}
\{(s,t):\|u_i(s,t)-u_j(s,t)\|<1/k\}$.
Z zupełności $X$ na tym zbiorze istnieje granica; definiujemy
$\widetilde f$ jako tę granicę, a poza tym zbiorem jako $0$.
Mnożąc $u_n$ przez funkcję charakterystyczną zbioru zbieżności,
otrzymujemy produktowo mierzalne funkcje proste zbieżne
punktowo do $\widetilde f$. W szczególności jej norma jest
produktowo mierzalna. Funkcje $f$ i $\widetilde f$ są równe
poza $N$. Na mocy skalarnego twierdzenia Fubiniego przekroje
$N$ mają miarę zero dla prawie wszystkich wartości każdego
z parametrów. Wystarczy więc wykazać tezę dla $\widetilde f$,
którą dalej oznaczamy przez $f$.

Funkcje postaci $\chi_{A\times B}(s,t)x$, gdzie $A\in\Sigma$,
$B\in\mathcal T$ i $x\in X$, spełniają żądane tożsamości.
Ich skończone kombinacje liniowe są gęste w
$L^1(S\times T;X)$: prostokąty generują produktową
$\sigma$-algebrę, a skończoność miar pozwala, za pomocą
argumentu klasy monotonicznej, aproksymować funkcję charakterystyczną
każdego zbioru produktowo mierzalnego w skalarnej normie $L^1$.
Istotnie, klasa zbiorów, których funkcje charakterystyczne należą
do domknięcia liniowej otoczki funkcji charakterystycznych prostokątów,
zawiera algebrę skończonych sum prostokątów (przez rozbicie na
rozłączne prostokąty). Jest zamknięta na rosnące sumy i malejące
przecięcia ciągów zbiorów, ponieważ przy skończonej mierze ich
funkcje charakterystyczne zbiegają w $L^1$. Twierdzenie o klasie
monotonicznej daje zatem całą produktową $\sigma$-algebrę.
Aproksymacja funkcji prostych daje ten sam wynik dla
funkcji o wartościach w $X$.

Wybierzmy zatem skończone kombinacje prostokątne $g_n$
takie, że
$\int_{S\times T}\|f-g_n\|\,\mathrm d(\mu\otimes\nu)<4^{-n}$.
Skalarne twierdzenie Tonellego daje
\[
 \int_S\sum_{n=1}^\infty
 \left(\int_T\|f(s,t)-g_n(s,t)\|\,\mathrm d\nu(t)\right)
 \mathrm d\mu(s)<\infty.
\]
Silna mierzalność przekrojów dla prawie każdego $s$ wynika
z przekrojów produktowo mierzalnych funkcji prostych
aproksymujących $f$ punktowo poza produktowym zbiorem miary
zero: jego przekroje są zerowe dla prawie każdego $s$.
Skalarne twierdzenie Tonellego dla $\|f\|$ daje całkowalność
normy przekroju dla prawie każdego $s$. Powyższe oszacowanie
sumy błędów daje ponadto $g_n(s,\cdot)\to f(s,\cdot)$
w $L^1(T;X)$.
Zgodnie z kryterium Bochnera przekrój jest całkowalny.
Dla takiego $s$ oznaczmy
$F(s)=\int_T f(s,t)\,\mathrm d\nu(t)$ oraz
$F_n(s)=\int_T g_n(s,t)\,\mathrm d\nu(t)$.
Przyjmujemy $F=0$ na wyjątkowym zbiorze miary zero.
Funkcje $F_n$ (określone na całym $S$ przez całkowanie
$g_n$ po drugiej zmiennej) są proste i $F_n(s)\to F(s)$
prawie wszędzie
na mocy oszacowania całki Bochnera po przekroju. Zatem $F$
jest silnie mierzalna i możemy zastosować Tonellego, otrzymując
\[
 \int_S\|F(s)-F_n(s)\|\,\mathrm d\mu(s)
 \le\int_{S\times T}\|f-g_n\|\,\mathrm d(\mu\otimes\nu)
 <4^{-n}.
\]
Z tego oszacowania i całkowalności $F_n$ wynika całkowalność
$F$ w sensie Bochnera.
Tożsamość dla $g_n$, oszacowanie Bochnera i przejście
do granicy dają
$\int_S F\,\mathrm d\mu
 =\int_{S\times T}f\,\mathrm d(\mu\otimes\nu)$.
Zamiana ról $S$ i $T$ prowadzi do drugiej całki iterowanej
oraz kończy dowód. Wszystkie wnioski dla pierwotnej $f$
są prawdziwe z dokładnością do odpowiednich zbiorów
miary zero.
\end{proof}

\begin{stwierdzenie}[Zgodność z całką Riemanna]
\label{prop:riemann-bochner15b}
Niech $a,b\in\R$, $a<b$, a $X$ będzie przestrzenią Banacha
nad $\K\in\{\R,\C\}$. Jeśli $f:[a,b]\to X$ jest ciągła, to jej całka Riemanna
w przestrzeni Banacha z \hyperref[sec:integral-lcs15a]{modułu 15A} jest równa całce Bochnera
względem miary Lebesgue'a na $[a,b]$. Jeżeli $X\ne\{0\}$, całkowalność Bochnera nie implikuje
całkowalności Riemanna, nawet na odcinku $[0,1]$.
\end{stwierdzenie}
\begin{proof}
Ciągłość na zwartym odcinku daje jednostajną ciągłość.
Funkcje proste, stałe na odcinkach drobnego podziału,
przybliżają $f$ jednostajnie; ich całki Bochnera są
odpowiednimi sumami Riemanna, a zbieżność zachodzi w normie.
Kontrprzykładem jest $f=\chi_{\mathbb Q\cap[0,1]}x_0$
dla niezerowego $x_0\in X$: funkcja jest równa zeru prawie
wszędzie, więc ma zerową całkę Bochnera, lecz nie jest
całkowalna w sensie Riemanna: wybierając w każdym
podprzedziale dodatniej długości punkt wymierny, otrzymujemy
sumę Riemanna równą $x_0$, a wybierając wyłącznie punkty
niewymierne, otrzymujemy sumę równą $0$. Oba wybory są
możliwe przy dowolnie małej średnicy podziału, więc nie
istnieje wspólna granica sum. W przypadku $X=\{0\}$
każda funkcja na odcinku jest całkowalna w obu sensach.
\end{proof}

\subsection{Całki wyznaczane przez wartości skalarne}

\begin{definicja}[Całka Dunforda i całka Pettisa]
\label{def:dunford-pettis15b}
Niech $(S,\Sigma,\mu)$ będzie zupełną przestrzenią z miarą
skończoną, a $X$ przestrzenią Banacha nad $\K\in\{\R,\C\}$.
Niech $f:S\to X$ będzie słabo mierzalna i niech
$x'\circ f\in L^1(S,\mu)$ dla każdego $x'\in X'$.
\emph{Całką Dunforda} funkcji $f$ na $A\in\Sigma$ nazywamy
element $D_A\in X''$ określony równością
$D_A(x')=\int_Ax'(f(s))\,\mathrm d\mu(s)$ dla każdego
$x'\in X'$. Funkcję spełniającą powyższe założenia
nazywamy całkowalną w sensie Dunforda; istnienie $D_A$
uzasadnia następne twierdzenie.
Funkcja $f$ jest \emph{całkowalna w sensie Pettisa}, gdy
dla każdego $A\in\Sigma$ istnieje $x_A\in X$ spełniający
$x'(x_A)=\int_Ax'(f(s))\,\mathrm d\mu(s)$ dla każdego
$x'\in X'$. Funkcjonały z $X'$ rozdzielają punkty $X$
na mocy twierdzenia Hahna--Banacha, więc $x_A$ jest
jednoznacznie wyznaczony. Wtedy
$x_A=\int_A^{\mathrm P} f(s)\,\mathrm d\mu(s)$.
Nie wystarcza wymagać reprezentacji tylko dla $A=S$:
wyjaśni to modyfikacja poniższego przykładu z $c_0$.
\end{definicja}

\begin{twierdzenie}[Istnienie całki Dunforda]
\label{thm:dunford-closedgraph15b}
Niech $(S,\Sigma,\mu)$ będzie zupełną przestrzenią
z miarą skończoną, $\K\in\{\R,\C\}$, a $X$ przestrzenią
Banacha nad $\K$. Niech funkcja $f:S\to X$ będzie słabo
mierzalna i niech $x'\circ f\in L^1(S,\mu)$ dla każdego
$x'\in X'$. Wówczas operator $Q_f:X'\to L^1(S,\mu)$,
$Q_f(x')=x'\circ f$, jest liniowy i ograniczony.
Dla każdego $A\in\Sigma$ wzór
$D_A(x')=\int_Ax'(f(s))\,\mathrm d\mu(s)$
określa element $D_A\in X''$, nawet gdy $f$ nie jest
całkowalna w sensie Bochnera, oraz
$\|D_A\|\le\|Q_f\|$.
\end{twierdzenie}
\begin{proof}
Jeżeli $x_n'\to x'$ w normie $X'$ i
$Q_f(x_n')\to h$ w $L^1$, to pewien podciąg
$Q_f(x_{n_k}')$ zbiega do $h$ prawie wszędzie.
Równocześnie dla każdego $s\in S$ mamy
$x_{n_k}'(f(s))\to x'(f(s))$, więc $h=Q_f(x')$ prawie
wszędzie. Wykres operatora $Q_f$ jest zatem domknięty.
Twierdzenie o wykresie domkniętym dla przestrzeni Banacha
daje $\|Q_f\|<\infty$. Stąd
$|D_A(x')|\le\int_A|x'(f(s))|\,\mathrm d\mu(s)
\le\|Q_f\|\|x'\|$, co definiuje element przestrzeni podwójnie sprzężonej.
\end{proof}

\begin{twierdzenie}[Hierarchia trzech całek]
\label{thm:hierarchy-integrals15b}
Niech $(S,\Sigma,\mu)$ będzie zupełną przestrzenią z miarą
skończoną, a $X$ przestrzenią Banacha nad $\K\in\{\R,\C\}$.
Jeżeli $f:S\to X$ jest całkowalna w sensie Bochnera,
to jest całkowalna w sensie Pettisa i Dunforda, a dla
każdego $A\in\Sigma$ zachodzi
$D_A=J_X(\int_A f(s)\,\mathrm d\mu(s))$.
Funkcja $f:S\to X$ słabo mierzalna i skalarnie całkowalna
(tzn. $x'\circ f\in L^1(S,\mu)$ dla każdego $x'\in X'$) jest
całkowalna w sensie Pettisa wtedy i tylko wtedy, gdy
$D_A\in J_X(X)$ dla każdego $A\in\Sigma$.
Całkowalność w sensie Dunforda nie implikuje w ogólności
całkowalności w sensie Pettisa, a całkowalność w sensie
Pettisa nie implikuje całkowalności w sensie Bochnera.
\end{twierdzenie}
\begin{proof}
Dla funkcji prostej równość
$x'(\int_A f(s)\,\mathrm d\mu(s))=\int_Ax'(f(s))\,\mathrm d\mu(s)$
jest bezpośrednia. Przybliżenie w $L^1(S;X)$
i $|x'(v)|\le\|x'\|\|v\|$ pozwalają przejść do granicy
w przypadku Bochnera. Równoważność w drugiej tezie wynika bezpośrednio z definicji:
$D_A=J_X(x_A)$ oznacza, że $D_A(x')=x'(x_A)$ dla każdego
$x'\in X'$, a istnienie takiego $x_A$ dla wszystkich
$A\in\Sigma$ jest warunkiem całkowalności w sensie Pettisa. Zapowiedziane
kontrprzykłady podajemy niżej.
\end{proof}

\begin{przyklad}[Całka Dunforda w $X''$ poza $J_X(X)$]
\label{ex:dunford-not-pettis15b}
Niech $S=\N=\{1,2,\ldots\}$, $\Sigma=\mathcal P(\N)$ i
$\mu(\{n\})=2^{-n}$. W przestrzeni Banacha $X=c_0$
przyjmijmy $f(n)=2^ne_n$, gdzie $(e_n)$ jest standardową
bazą jednostkową $c_0$. Ponieważ $X'=\ell^1$, dla
$a=(a_n)\in\ell^1$ mamy
$\int_S|a(f(n))|\,\mathrm d\mu(n)
=\sum_{n\ge1}|a_n|<\infty$.
Całka Dunforda ma zatem współrzędne
$D_S=(1,1,\ldots)\in\ell^\infty=X''$ i nie leży w
$J_X(c_0)$. Funkcja nie jest całkowalna w sensie Pettisa
ani Bochnera; co więcej,
$\int_S\|f(n)\|\,\mathrm d\mu(n)=\sum_{n\ge1}1=\infty$.

Rozważmy teraz $S_*=\N\times\{-1,1\}$ z $\sigma$-algebrą
wszystkich podzbiorów i miarą
$\mu_*(\{(n,\eta)\})=2^{-n-1}$ dla $\eta\in\{-1,1\}$.
Niech $h(n,\eta)=\eta 2^ne_n\in c_0$.
Dla $a\in\ell^1$ całka modułu skalaryzacji jest równa
$\sum_n|a_n|$, zatem $h$ jest całkowalna w sensie Dunforda.
Na całym $S_*$ wkłady przeciwnych znaków znoszą się i
$D_{S_*}=0\in J_{c_0}(c_0)$ (szereg jest bezwzględnie
zbieżny po każdej skalaryzacji). Natomiast dla
$A=\N\times\{1\}$ otrzymujemy
$D_A=(1/2,1/2,\ldots)\notin J_{c_0}(c_0)$.
Stąd $h$ nie jest całkowalna w sensie Pettisa, mimo że
całka Dunforda po całym $S_*$ jest reprezentowana
przez wektor z $c_0$.
\end{przyklad}

\begin{przyklad}[Całka Pettisa bez całki Bochnera]
\label{ex:pettis-not-bochner15b}
Niech $S=[0,1]$ z miarą Lebesgue'a, zaś
$X=\ell^2([0,1])$. Dla $t\in[0,1]$ definiujemy $f(t)=e_t$.
Każdy $a\in\ell^2([0,1])$ ma co najwyżej przeliczalnie
wiele niezerowych współrzędnych: dla każdego $n\ge1$
zbiór $\{t:|a_t|\ge 1/n\}$ jest skończony, gdyż w przeciwnym
razie sumy kwadratów po skończonych podzbiorach byłyby
nieograniczone. Suma tych zbiorów zawiera wszystkie
niezerowe współrzędne. Na mocy twierdzenia Riesza o
reprezentacji funkcjonałów w przestrzeni Hilberta każda
skalaryzacja ma postać $t\mapsto\langle f(t),a\rangle$
(przy konwencji liniowości względem pierwszej zmiennej).
Jest ona równa zeru prawie wszędzie i jest mierzalna,
ponieważ miara Lebesgue'a jest zupełna. Całki skalarne po każdym $A\in\Sigma$ znikają,
więc $f$ ma całkę Pettisa równą $0$ na każdym $A$.
Nie jest jednak silnie mierzalna: zbiór jej wartości poza
jakimkolwiek zbiorem miary zero jest nieprzeliczalną rodziną
punktów odległych o $\sqrt 2$, a więc nie może leżeć
w ośrodkowej podprzestrzeni $X$. Całka Bochnera nie istnieje,
choć $\|f(t)\|=1$ dla każdego $t$.
\end{przyklad}

\begin{przyklad}[Silna mierzalność i całka Pettisa bez całki Bochnera]
\label{ex:pettis-strong-not-bochner15b}
Niech $S=\N=\{1,2,\ldots\}$, $\Sigma=\mathcal P(\N)$,
$\mu(A)=\sum_{n\in A}2^{-n}$ dla $A\in\Sigma$ oraz
$X=\ell^2(\N;\K)$, gdzie $\K\in\{\R,\C\}$. Miara $\mu$ jest
zupełna i $\mu(S)=1$. Przestrzeń $X$ jest ośrodkową
przestrzenią Hilberta. Oznaczmy przez $(e_n)$ jej standardową
bazę ortonormalną i zdefiniujmy
$f:S\to X$ wzorem $f(n)=\frac{2^n}{n}e_n$.

\emph{Silna mierzalność.} Dla $N\ge1$ funkcja
$s_N=\sum_{n=1}^N\chi_{\{n\}}\frac{2^n}{n}e_n$ jest prosta.
Dla każdego $k\in S$ i każdego $N\ge k$ mamy $s_N(k)=f(k)$,
zatem $s_N(k)\to f(k)$ w normie w każdym punkcie.
W szczególności $f$ jest silnie mierzalna, a na mocy
twierdzenia Pettisa o mierzalności również słabo mierzalna.

\emph{Całkowalność skalarna.} Dla dowolnego $x'\in X'$
twierdzenie Riesza o reprezentacji funkcjonałów daje
wektor $a=(a_n)\in\ell^2(\N;\K)$ taki, że
$x'(v)=\langle v,a\rangle$ dla $v\in X$.
Przyjmujemy iloczyn skalarny liniowy względem pierwszej
zmiennej, więc $x'(e_n)=\overline{a_n}$ (w przypadku
rzeczywistym sprzężenie niczego nie zmienia). Z definicji
miary atomowej i nierówności Cauchy'ego--Schwarza otrzymujemy
\[
 \int_S|x'(f(n))|\,\mathrm d\mu(n)
 =\sum_{n=1}^{\infty}\frac{|a_n|}{n}
 \le \|a\|_{\ell^2}
       \left(\sum_{n=1}^{\infty}\frac1{n^2}\right)^{1/2}
 <\infty.
\]
Zatem $f$ jest całkowalna w sensie Dunforda.

\emph{Reprezentacja całek Pettisa.} Dla każdego $A\in\Sigma$
definiujemy wektor $v_A=(\chi_A(n)/n)_{n\ge1}$.
Ponieważ $\sum_{n\ge1}|\chi_A(n)/n|^2
\le\sum_{n\ge1}1/n^2<\infty$, mamy $v_A\in X$.
Dla dowolnego $x'\in X'$ reprezentowanego przez $a$ jak wyżej
zachodzi
\[
 \int_A x'(f(n))\,\mathrm d\mu(n)
 =\sum_{n\in A}\frac{\overline{a_n}}{n}
 =\langle v_A,a\rangle=x'(v_A).
\]
Szereg skalarny jest bezwzględnie zbieżny na mocy
poprzedniego oszacowania. Ponieważ równość zachodzi dla
każdego $A$ i każdego $x'$, funkcja $f$ jest całkowalna
w sensie Pettisa, a
$\int_A^{\mathrm P} f(n)\,\mathrm d\mu(n)=v_A$.
W szczególności całka po $S$ jest wektorem
$(1,1/2,1/3,\ldots)\in\ell^2$.

\emph{Brak całkowalności Bochnera.} Mamy $\|e_n\|=1$,
więc
\[
 \int_S\|f(n)\|\,\mathrm d\mu(n)
 =\sum_{n=1}^{\infty}2^{-n}\frac{2^n}{n}
 =\sum_{n=1}^{\infty}\frac1n=\infty.
\]
Z kryterium całkowalności Bochnera wynika, że $f$ nie
jest całkowalna w tym sensie, pomimo silnej mierzalności.
W odróżnieniu od \hyperref[ex:pettis-not-bochner15b]{poprzedniego przykładu}
przeszkodą nie jest brak silnej mierzalności, lecz brak
całkowalności normy. Pokazuje to, że nawet w ośrodkowej
przestrzeni Hilberta całkowalność Pettisa jest słabsza
od całkowalności Bochnera.
\end{przyklad}

\begin{uwaga}[Refleksywność i całka Dunforda]
Jeżeli $X$ jest refleksywna, to $J_X(X)=X''$ i każdy
$D_A$ jest reprezentowany przez wektor z $X$. Zatem każda
słabo mierzalna, skalarnie całkowalna funkcja do refleksywnej
przestrzeni Banacha jest całkowalna w sensie Pettisa.
Nie oznacza to całkowalności Bochnera; przykład z
przestrzeni Hilberta $\ell^2([0,1])$ pozostaje ważny.
\end{uwaga}

\begin{twierdzenie}[Zbieżność zmajoryzowana dla całki Pettisa]
\label{thm:pettis-convergence15b}
Niech $(S,\Sigma,\mu)$ będzie zupełną przestrzenią z miarą
skończoną, a $X$ refleksywną przestrzenią Banacha nad
$\K\in\{\R,\C\}$. Niech $f:S\to X$ będzie funkcją, a
$(f_n)_{n\ge1}$ ciągiem słabo mierzalnych funkcji $f_n:S\to X$
takim, że
$f_n(s)\to f(s)$ w normie dla prawie każdego $s$.
Załóżmy, że $\|f_n(s)\|\le g(s)$ dla wszystkich $n$
i prawie wszystkich $s$, gdzie $g\in L^1(S,\mu)$,
$g\ge0$, oraz że funkcje
$s\mapsto\|f_n(s)-f(s)\|$ są mierzalne.
Wtedy $f_n$ i $f$ są całkowalne w sensie Pettisa i dla
każdego $A\in\Sigma$ zachodzi
$\int_A^{\mathrm P}f_n(s)\,\mathrm d\mu(s)
\to\int_A^{\mathrm P}f(s)\,\mathrm d\mu(s)$ w normie $X$.
\end{twierdzenie}
\begin{proof}
Funkcja $f$ jest słabo mierzalna, bo dla każdego $x'\in X'$
jej skalaryzacja jest granicą mierzalnych $x'\circ f_n$.
Dominacja przez $g$ daje całkowalność skalarną $f_n$ i $f$.
Z refleksywności $X$ oraz twierdzenia o istnieniu całki
Dunforda wynika, że wszystkie te funkcje mają całkę Pettisa.
Całka Pettisa jest liniowa na przestrzeni funkcji
całkowalnych w tym sensie: liniowość całek skalarnych
pokazuje, że kombinacja liniowa wektorów reprezentujących
jest wektorem reprezentującym odpowiednią kombinację funkcji,
a jednoznaczność wynika z twierdzenia Hahna--Banacha.
Dla każdego $A\in\Sigma$ i $x'\in X'$ o normie co najwyżej
$1$ zachodzi
$|x'(\int_A^{\mathrm P}(f_n(s)-f(s))\,\mathrm d\mu(s))|
\le\int_A\|f_n(s)-f(s)\|\,\mathrm d\mu(s)$.
Biorąc supremum po takich $x'$ otrzymujemy oszacowanie
normy różnicy całek. Skalarne twierdzenie o zbieżności
zmajoryzowanej stosuje się do mierzalnych funkcji
$\|f_n-f\|\le2g$; jego prawa strona dąży do zera.
\end{proof}

\subsection{Całka Gelfanda i miary wektorowe}

\begin{definicja}[Całka Gelfanda]
\label{def:gelfand15b}
Niech $(S,\Sigma,\mu)$ będzie zupełną przestrzenią z miarą
skończoną, a $Y$ przestrzenią Banacha nad $\K\in\{\R,\C\}$.
Niech $f:S\to Y'$ będzie funkcją \emph{słabo-* mierzalną}, tzn.
$s\mapsto f(s)(y)$ jest mierzalna dla każdego $y\in Y$.
Zakładamy, że $f(\cdot)(y)\in L^1(S,\mu)$ dla każdego
$y\in Y$. \emph{Całką Gelfanda} na $A\in\Sigma$ nazywamy
jedyny $g_A\in Y'$ spełniający
$g_A(y)=\int_Af(s)(y)\,\mathrm d\mu(s)$ dla wszystkich
$y\in Y$. Stosujemy oznaczenie
$g_A=\int_A^{\mathrm G}f(s)\,\mathrm d\mu(s)$.
Istnienie tego funkcjonału uzasadnia następne twierdzenie.
\end{definicja}

\begin{twierdzenie}[Istnienie całki Gelfanda i jej addytywność]
\label{thm:gelfand-exists15b}
Niech $(S,\Sigma,\mu)$ będzie zupełną przestrzenią
z miarą skończoną, $\K\in\{\R,\C\}$, a $Y$
przestrzenią Banacha nad $\K$. Niech $f:S\to Y'$
będzie słabo-* mierzalna oraz
$f(\cdot)(y)\in L^1(S,\mu)$ dla każdego $y\in Y$.
Wtedy dla każdego $A\in\Sigma$ istnieje dokładnie jeden
$g_A\in Y'$ taki, że
$g_A(y)=\int_A f(s)(y)\,\mathrm d\mu(s)$
dla wszystkich $y\in Y$. Odwzorowanie
$A\mapsto g_A$ jest przeliczalnie addytywne
w topologii słabej-* $\sigma(Y',Y)$; bez dodatkowych
założeń nie musi być przeliczalnie addytywne w normie $Y'$.
\end{twierdzenie}
\begin{proof}
Operator $Q:Y\to L^1(S,\mu)$, $(Qy)(s)=f(s)(y)$,
jest domknięty. Istotnie, z $y_n\to y$ i $Qy_n\to h$
w $L^1$ wybieramy podciąg zbieżny prawie wszędzie;
punktowa ciągłość każdego $f(s)\in Y'$ daje $h=Qy$.
Z twierdzenia o wykresie domkniętym $Q$ jest ograniczony.
Zatem $|\int_Af(s)(y)\,\mathrm d\mu(s)|\le\|Q\|\|y\|$,
a wzór definiuje $g_A\in Y'$.
Dla parami rozłącznych $A_n$ i $y\in Y$ skalarna
addytywność całki Lebesgue'a daje
$g_{\bigcup_nA_n}(y)=\sum_ng_{A_n}(y)$.
Jest to dokładnie zbieżność w topologii $\sigma(Y',Y)$.
\end{proof}

\begin{przyklad}[Słaba-* addytywność bez addytywności normowej]
\label{ex:gelfand-weakstar15b}
Niech $Y=\ell^1$, $Y'=\ell^\infty$, a na $S=\N$ przyjmijmy
$\mu(\{n\})=2^{-n}$ oraz $f(n)=2^ne_n\in\ell^\infty$.
Dla $y=(y_n)\in\ell^1$ zachodzi
$\int_S|f(n)(y)|\,\mathrm d\mu(n)=\sum_n|y_n|<\infty$.
W konsekwencji $g_A=\chi_A\in\ell^\infty$, gdzie
$\chi_A=(\chi_A(n))_{n\ge1}$. Dla $A_m=\{m,m+1,\ldots\}$
mamy $\|g_{A_m}\|_\infty=1$ przy każdym $m$, mimo że
$g_{A_m}\to0$ w topologii $\sigma(\ell^\infty,\ell^1)$:
dla $y\in\ell^1$ zachodzi
$|g_{A_m}(y)|\le\sum_{n\ge m}|y_n|\to0$.
Ponieważ $g_S-\sum_{n=1}^N g_{\{n\}}=g_{A_{N+1}}$,
szereg wartości na singletonach nie zbiega w normie do $g_S$.
\end{przyklad}

\begin{definicja}[Miara wektorowa i całkowita wariacja]
\label{def:vector-measure15b}
Niech $(S,\Sigma)$ będzie przestrzenią mierzalną, a $X$
przestrzenią Banacha nad $\K\in\{\R,\C\}$.
Odwzorowanie $m:\Sigma\to X$ nazywamy \emph{miarą wektorową},
jeśli $m(\emptyset)=0$ i dla każdego ciągu parami rozłącznych
$A_n\in\Sigma$ zachodzi
$m(\bigcup_{n\ge1}A_n)=\lim_{N\to\infty}\sum_{n=1}^Nm(A_n)$
w normie $X$. Jej \emph{całkowitą wariacją} nazywamy funkcję
$|m|:\Sigma\to[0,\infty]$ określoną dla $A\in\Sigma$ przez
$|m|(A)=\sup\{\sum_{j=1}^k\|m(A_j)\|:\ A_j\text{ są
parami rozłącznymi zbiorami mierzalnymi, }
\bigcup_{j=1}^kA_j=A,\ k\ge1\}$.
Dopuszczamy puste składniki podziału. W następnym stwierdzeniu
wykażemy, że $|m|$ jest miarą dodatnią.
\end{definicja}

\begin{stwierdzenie}[Przeliczalna addytywność całkowitej wariacji]
\label{prop:variation-measure15b}
Niech $(S,\Sigma)$ będzie przestrzenią mierzalną, a $X$
przestrzenią Banacha nad $\K\in\{\R,\C\}$.
Jeżeli $m:\Sigma\to X$ jest miarą wektorową, to
$|m|:\Sigma\to[0,\infty]$ jest miarą dodatnią (niekoniecznie
skończoną). W szczególności dla parami rozłącznych
$A_n\in\Sigma$ zachodzi
$|m|(\bigcup_{n\ge1}A_n)=\sum_{n\ge1}|m|(A_n)$.
\end{stwierdzenie}
\begin{proof}
Z definicji $|m|(\emptyset)=0$. Niech
$A=\bigcup_{n\ge1}A_n$ będzie sumą parami rozłącznych
zbiorów mierzalnych. Dla każdego $N$ można zestawić
skończone partycje zbiorów $A_1,\ldots,A_N$ w jedną
partycję $A$, dodając ewentualnie
$A\setminus\bigcup_{n=1}^N A_n$.
Po przejściu do supremów i następnie z $N\to\infty$
dostajemy $|m|(A)\ge\sum_{n\ge1}|m|(A_n)$,
również gdy prawa strona jest nieskończona.

Dla nierówności odwrotnej wybierzmy dowolną skończoną
partycję $B_1,\ldots,B_k$ zbioru $A$.
Z przeliczalnej addytywności $m$ w normie i nierówności
trójkąta wynika
\[
 \sum_{j=1}^k\|m(B_j)\|
 \le\sum_{j=1}^k\sum_{n=1}^\infty
        \|m(B_j\cap A_n)\|
 \le\sum_{n=1}^\infty |m|(A_n).
\]
Po wzięciu supremum po wszystkich partycjach $A$
otrzymujemy nierówność odwrotną i tezę.
\end{proof}

\begin{twierdzenie}[Miara generowana przez całkę Bochnera]
\label{thm:bochner-measure15b}
Niech $(S,\Sigma,\mu)$ będzie zupełną przestrzenią z miarą
skończoną, a $X$ przestrzenią Banacha nad $\K\in\{\R,\C\}$.
Jeśli $f:S\to X$ jest całkowalna w sensie Bochnera, to
odwzorowanie $m_f:\Sigma\to X$ określone wzorem
$m_f(A)=\int_Af(s)\,\mathrm d\mu(s)$ jest miarą wektorową,
a jej całkowita wariacja spełnia
\[
 |m_f|(A)=\int_A\|f(s)\|\,\mathrm d\mu(s)
 \qquad(A\in\Sigma).
\]
W szczególności $|m_f|$ jest skończoną miarą bezwzględnie
ciągłą względem $\mu$.
\end{twierdzenie}
\begin{proof}
Dla parami rozłącznych $A_n$ i $A=\bigcup_nA_n$
z oszacowania całki otrzymujemy
$\|m_f(A)-\sum_{n=1}^Nm_f(A_n)\|
\le\int_{\bigcup_{n>N}A_n}\|f(s)\|\,\mathrm d\mu(s)\to0$.
Dla każdej skończonej partycji $A$ suma norm całek
jest ograniczona przez $\int_A\|f\|$, czyli
$|m_f|(A)\le\int_A\|f\|$.
Gdy $f=\sum_j\chi_{B_j}v_j$ jest prosta o rozłącznych
zbiorach poziomicowych, partycja $A$ przez $A\cap B_j$
daje równość. Dla ogólnej $f$ wybierzmy funkcje proste
$f_n\to f$ w $L^1$. Z nierówności
$\bigl||m_f|(A)-|m_{f_n}|(A)\bigr|
\le\int_A\|f(s)-f_n(s)\|\,\mathrm d\mu(s)$
oraz zbieżności norm w $L^1$ otrzymujemy równość.
Istotnie, dla każdej skończonej partycji $A=\bigcup_j A_j$
mamy
$\sum_j\|m_f(A_j)\|\le
\sum_j\|m_{f_n}(A_j)\|+
\sum_j\|m_{f-f_n}(A_j)\|
\le |m_{f_n}|(A)+\int_A\|f-f_n\|\,\mathrm d\mu$.
Biorąc supremum i zamieniając role $f$ oraz $f_n$,
otrzymujemy przytoczone oszacowanie.
\end{proof}

\begin{przyklad}[Miara wektorowa o nieskończonej całkowitej wariacji]
\label{ex:pettis-infinite-variation15b}
Niech $S=\N=\{1,2,\ldots\}$, $\Sigma=\mathcal P(\N)$,
$\mu(\{n\})=2^{-n}$ oraz $X=\ell^2(\N;\K)$,
gdzie $\K\in\{\R,\C\}$. W
\hyperref[ex:pettis-strong-not-bochner15b]{przykładzie dotyczącym silnej mierzalności}
wykazaliśmy, że funkcja $f(n)=\frac{2^n}{n}e_n$ jest
całkowalna w sensie Pettisa, a dla każdego $A\in\Sigma$
jej całka ma postać
\[
 m(A)=\int_A^{\mathrm P}f(n)\,\mathrm d\mu(n)
 =\left(\frac{\chi_A(n)}{n}\right)_{n\ge1}.
\]
Wykażemy, że $m:\Sigma\to X$ jest miarą wektorową
bezwzględnie ciągłą względem $\mu$, lecz jej całkowita
wariacja nie jest skończona.

\emph{Przeliczalna addytywność w normie.}
Oczywiście $m(\emptyset)=0$. Niech $(A_j)_{j\ge1}$ będzie
ciągiem parami rozłącznych zbiorów z $\Sigma$ i niech
$A=\bigcup_{j\ge1}A_j$, $B_N=\bigcup_{j>N}A_j$.
Ponieważ zbiory $A_j$ są rozłączne, dla każdego $n\in\N$
mamy $\chi_A(n)-\sum_{j=1}^N\chi_{A_j}(n)=\chi_{B_N}(n)$.
Zatem
\[
 \left\|m(A)-\sum_{j=1}^Nm(A_j)\right\|_{\ell^2}^2
 =\sum_{n=1}^{\infty}\frac{\chi_{B_N}(n)}{n^2}.
\]
Zbiory $B_N$ maleją do zbioru pustego: każdy element $A$
należy do dokładnie jednego $A_j$, a element spoza $A$
nie należy do żadnego $B_N$. Stąd dla każdego $n$
zachodzi $\chi_{B_N}(n)\to0$. Ponadto
$0\le\chi_{B_N}(n)/n^2\le1/n^2$ i szereg
$\sum_{n\ge1}1/n^2$ jest zbieżny. Twierdzenie o zbieżności
zmajoryzowanej dla miary liczącej daje zbieżność prawej
strony do zera. Otrzymujemy przeliczalną addytywność
$m$ w normie $X$. Jeżeli $\mu(A)=0$, to $A=\emptyset$,
gdyż każdy singleton ma dodatnią miarę; w szczególności
$m(A)=0$.

\emph{Całkowita wariacja.}
Dla dowolnego $A\in\Sigma$ i każdego skończonego podziału
$A=\bigcup_{j=1}^k C_j$ zachodzi
\[
 \sum_{j=1}^k\|m(C_j)\|_{\ell^2}
 =\sum_{j=1}^k\left(\sum_{n\in C_j}\frac1{n^2}\right)^{1/2}
 \le\sum_{j=1}^k\sum_{n\in C_j}\frac1n
 =\sum_{n\in A}\frac1n.
\]
Nierówność wynika z oszacowania normy $\ell^2$ przez
sumę modułów współrzędnych, które dla szeregu nieskończonego
otrzymujemy przez przejście do granicy w oszacowaniu
skończonych sum; dopuszczamy wartość $+\infty$.
Biorąc supremum po podziałach, otrzymujemy
$|m|(A)\le\sum_{n\in A}1/n$.
Z drugiej strony, dla każdego $N\ge1$ możemy podzielić $A$
na singletony $\{n\}$ dla $n\in A\cap\{1,\ldots,N\}$
oraz pozostały zbiór $A\setminus\{1,\ldots,N\}$.
Ponieważ $m(\{n\})=e_n/n$, definicja wariacji daje
$|m|(A)\ge\sum_{n\in A,\ n\le N}1/n$.
Przechodząc z $N\to\infty$, otrzymujemy
\[
 |m|(A)=\sum_{n\in A}\frac1n
 \qquad(A\in\Sigma).
\]
W szczególności $|m|(S)=\sum_{n\ge1}1/n=\infty$.

Zauważmy, że wartości samej miary są jednostajnie ograniczone:
$\|m(A)\|_{\ell^2}\le(\sum_{n\ge1}1/n^2)^{1/2}$
dla każdego $A$. Mimo to jej całkowita wariacja jest
nieskończona. Przykład pokazuje, że miara generowana przez
całkę Pettisa może być przeliczalnie addytywna w normie,
lecz nie spełniać warunku skończonej wariacji, który
zawsze zachodzi dla miary generowanej przez całkę Bochnera.
\end{przyklad}

\subsection{Własności całek wektorowych i zależność od zbioru całkowania}

W dalszym ciągu miara jest skończona i zupełna. Zasadnicza
różnica pomiędzy całkami nie polega na liniowości, która
zachodzi dla wszystkich czterech całek, lecz na przestrzeni
wartości, rodzinie funkcjonałów użytych do skalaryzacji
oraz topologii, w której można sumować całki po rozłącznych
zbiorach. Wprowadzamy oznaczenia $P_A(f)=\int_A^{\mathrm P}f\,\mathrm d\mu$,
$D_A(f)$ dla całki Dunforda i $G_A(h)=\int_A^{\mathrm G}h\,\mathrm d\mu$.

\begin{twierdzenie}[Liniowość, ograniczenie i mnożenie przez funkcje skalarne]
\label{thm:vector-integral-multipliers15b}
Niech $(S,\Sigma,\mu)$ będzie zupełną przestrzenią z miarą
skończoną, $X,Y$ przestrzeniami Banacha nad $\K\in\{\R,\C\}$,
$f:S\to X$ funkcją całkowalną w sensie Dunforda, a
$h:S\to Y'$ funkcją słabo-* mierzalną, dla której
$h(\cdot)(y)\in L^1(S,\mu)$ dla każdego $y\in Y$.
Wprowadźmy dla $A\in\Sigma$
\[
 \begin{split}
 q_f(A)&=\sup_{\|x'\|\le1,\ x'\in X'}
       \int_A|x'(f(s))|\,\mathrm d\mu(s),\\
 r_h(A)&=\sup_{\|y\|\le1,\ y\in Y}
       \int_A|h(s)(y)|\,\mathrm d\mu(s).
 \end{split}
\]
Wówczas $q_f(S)$ i $r_h(S)$ są skończone. Dla każdej
$\varphi\in L^\infty(S,\mu;\K)$ (wybieramy mierzalnego
przedstawiciela klasy) funkcja $\varphi f$ jest
całkowalna w sensie Dunforda, $\varphi h$ w sensie Gelfanda,
a gdy $f$ jest całkowalna w sensie Pettisa, to także
$\varphi f$ jest całkowalna w sensie Pettisa. Dla każdego
$A\in\Sigma$ zachodzą oszacowania
\[
 \begin{split}
 \|D_A(\varphi f)\|&\le\|\varphi\|_\infty q_f(A),\\
 \|G_A(\varphi h)\|&\le\|\varphi\|_\infty r_h(A),\\
 \|P_A(\varphi f)\|&\le\|\varphi\|_\infty q_f(A)
       \quad\text{(w przypadku Pettisa).}
 \end{split}
\]
Każda z tych całek jest liniowa względem funkcji
całkowanej w odpowiedniej klasie całkowalności i nie zmienia
się po zmianie funkcji na zbiorze miary zero. Dla ustalonego
$A$ jest też liniowa względem mnożnika $\varphi$.
Jeżeli $f$ jest całkowalna w sensie Bochnera, to
$\varphi f$ jest całkowalna w tym sensie i
$q_f(A)\le\int_A\|f\|\,\mathrm d\mu$.
\end{twierdzenie}
\begin{proof}
Z \hyperref[thm:dunford-closedgraph15b]{twierdzenia o istnieniu całki Dunforda}
mamy $q_f(S)=\|Q_f\|$, gdzie $Q_f(x')=x'\circ f$.
Analogicznie dowód
\hyperref[thm:gelfand-exists15b]{twierdzenia o istnieniu całki Gelfanda}
daje $r_h(S)=\|Q\|$, gdzie $(Qy)(s)=h(s)(y)$.
Mnożenie funkcji skalarnej z $L^1$ przez funkcję z
$L^\infty$ zachowuje całkowalność, a wszystkie
odpowiednie skalaryzacje pozostają mierzalne. Definicje
całek Dunforda i Gelfanda dają więc istnienie oraz
oszacowania, gdyż dla $\|x'\|\le1$ mamy
$|\int_A\varphi(s)x'(f(s))\,\mathrm d\mu(s)|
\le\|\varphi\|_\infty\int_A|x'(f(s))|\,\mathrm d\mu(s)$,
a dla Gelfanda stosujemy ten sam argument do $h(s)(y)$.
Bierzemy supremum po odpowiedniej kuli jednostkowej.

Niech teraz $f$ będzie całkowalna w sensie Pettisa.
Dla skalarnej funkcji prostej
$\psi=\sum_{j=1}^k c_j\chi_{B_j}$ wektor
$\sum_j c_j P_{A\cap B_j}(f)$ reprezentuje całki skalarne
funkcji $\psi f$ po $A$, co wynika z liniowości
całki skalarnej. Zatem $\psi f$ jest całkowalna w sensie
Pettisa. Z twierdzenia Hahna--Banacha i poprzedniego
oszacowania wynika
$\|P_A(\psi f)\|\le\|\psi\|_\infty q_f(A)$.
Wybierzmy skalarne funkcje proste $\psi_k$ zbieżne do
$\varphi$ w normie $L^\infty$ (po wyborze ograniczonego
przedstawiciela wystarczy podzielić jego ograniczony
zbiór wartości na zbiory o coraz mniejszej średnicy).
Z liniowości dla funkcji prostych dostajemy
\[
 \|P_A(\psi_k f)-P_A(\psi_l f)\|
 \le\|\psi_k-\psi_l\|_\infty q_f(A).
\]
Zupełność $X$ daje granicę $v_A$ w normie. Dla każdego
$x'\in X'$ różnica całek skalarnych odpowiadających
$\psi_k f$ i $\varphi f$ ma moduł nie większy niż
$\|\psi_k-\varphi\|_\infty\int_A|x'(f)|\,\mathrm d\mu$,
a więc dąży do zera. Ciągłość $x'$ daje
$x'(v_A)=\int_A\varphi x'(f)\,\mathrm d\mu$.
Ponieważ $A$ było dowolne, jest to całkowalność Pettisa.
Przejście do granicy daje także oszacowanie normy.

Liniowość i niezmienniczość względem zmian na zbiorach
miary zero wynikają z odpowiednich własności całki
skalarnej oraz jednoznaczności reprezentacji: w $X$ przez
funkcjonały z $X'$, w $X''$ przez wartości na $X'$, a
w $Y'$ przez wartości na $Y$. W przypadku Bochnera
$\varphi f$ jest silnie mierzalna (iloczyny prostych
aproksymacji punktowo przybliżają iloczyn poza zbiorem
zerowym), a $\|\varphi f\|\le\|\varphi\|_\infty\|f\|$
prawie wszędzie. Kryterium Bochnera kończy ten przypadek;
nierówność dla $q_f$ wynika z $|x'(f)|\le\|x'\|\|f\|$.
\end{proof}

\begin{stwierdzenie}[Przechodzenie przez operatory liniowe]
\label{prop:vector-integrals-operators15b}
Niech $(S,\Sigma,\mu)$ będzie zupełną przestrzenią z miarą
skończoną, $X,Z,Y,W$ przestrzeniami Banacha nad
$\K\in\{\R,\C\}$, $T\in\mathcal L(X,Z)$ i
$R\in\mathcal L(W,Y)$. Jeżeli $f:S\to X$ jest
całkowalna w sensie Dunforda, to $Tf$ jest całkowalna
w sensie Dunforda i dla każdego $A\in\Sigma$
$D_A(Tf)=T''D_A(f)$.
Jeżeli $f$ jest całkowalna w sensie Pettisa, to $Tf$
jest całkowalna w sensie Pettisa i
$P_A(Tf)=T P_A(f)$.
Jeżeli $h:S\to Y'$ jest całkowalna w sensie Gelfanda,
to $R'h:S\to W'$ jest całkowalna w sensie Gelfanda i
$G_A(R'h)=R'G_A(h)$.
\end{stwierdzenie}
\begin{proof}
Przez $T':Z'\to X'$ rozumiemy operator
$T'z'=z'\circ T$, a przez $T'':X''\to Z''$ jego
operator sprzężony. Dla $z'\in Z'$ skalaryzacja
$z'\circ Tf=(T'z')\circ f$ jest mierzalna i całkowalna.
Zatem
$D_A(Tf)(z')=D_A(f)(T'z')=(T''D_A(f))(z')$.
Jeżeli istnieje wektor Pettisa $P_A(f)$, to
$z'(TP_A(f))=(T'z')(P_A(f))
=\int_Az'(Tf(s))\,\mathrm d\mu(s)$.
To dowodzi zarówno całkowalności, jak i żądanej równości.
Dla Gelfanda, $w\in W$ i $R':Y'\to W'$ mamy
$(R'h(s))(w)=h(s)(Rw)$. Ta funkcja skalarna jest
mierzalna i całkowalna na mocy założeń, a jej całka
jest równa $G_A(h)(Rw)=(R'G_A(h))(w)$.
Funkcjonały są równe, ponieważ przyjmują te same wartości
na każdym $w\in W$.
\end{proof}

\begin{stwierdzenie}[Addytywność Dunforda i Gelfanda]
\label{prop:dunford-gelfand-additivity15b}
Niech $(S,\Sigma,\mu)$ będzie zupełną przestrzenią z miarą
skończoną, $X,Y$ przestrzeniami Banacha nad
$\K\in\{\R,\C\}$, $f:S\to X$ funkcją całkowalną
w sensie Dunforda i $h:S\to Y'$ funkcją całkowalną
w sensie Gelfanda. Odwzorowania $A\mapsto D_A(f)$
i $A\mapsto G_A(h)$ są przeliczalnie addytywne odpowiednio
w topologiach $\sigma(X'',X')$ i $\sigma(Y',Y)$,
lecz w ogólności nie w normie. Są skończenie addytywne
w normie, jednostajnie ograniczone oraz zerują się na
zbiorach miary zero.
Jeżeli istnieje $g\in L^1(S,\mu)$, $g\ge0$, taka że
$\|f(s)\|\le g(s)$ prawie wszędzie, to
$A\mapsto D_A(f)$ jest miarą wektorową o wariacji
$|D|(A)\le\int_A g\,\mathrm d\mu$.
Analogicznie, jeżeli $\|h(s)\|\le g(s)$ prawie wszędzie,
to $A\mapsto G_A(h)$ jest miarą wektorową
o wariacji $|G|(A)\le\int_A g\,\mathrm d\mu$.
Nie wymaga się mierzalności samych norm.
\end{stwierdzenie}
\begin{proof}
Dla ciągu parami rozłącznych $A_j\in\Sigma$ oraz
$A=\bigcup_j A_j$ skalarna addytywność daje
$D_A(f)(x')=\sum_j D_{A_j}(f)(x')$ dla każdego $x'\in X'$,
a dla Gelfanda $G_A(h)(y)=\sum_jG_{A_j}(h)(y)$ dla każdego
$y\in Y$. Są to dokładnie deklarowane zbieżności słabe-*.
Skończona addytywność jest równością wektorów, gdyż
w obu przypadkach punkty są rozdzielane przez wskazane
skalaryzacje. Oszacowania z
\hyperref[thm:vector-integral-multipliers15b]{twierdzenia o mnożnikach}
dają jednostajne ograniczenie, a zerowanie się na zbiorach
miary zero wynika z własności całek skalarnych.
Brak addytywności w normie Dunforda pokazuje pierwszy
\hyperref[ex:dunford-not-pettis15b]{przykład z $c_0$}:
$D_S=(1,1,\ldots)$, a $\sum_{n=1}^ND_{\{n\}}$ ma jedynki
tylko na pierwszych $N$ pozycjach, zatem norma różnicy
wynosi $1$. Dla Gelfanda analogiczny argument podano
w \hyperref[ex:gelfand-weakstar15b]{przykładzie z $\ell^\infty$}.

Przy dominacji przez $g$ mamy
$\|D_B(f)\|\le\int_Bg\,\mathrm d\mu$ dla każdego $B$,
ponieważ $|x'(f)|\le g$ dla $\|x'\|\le1$ prawie wszędzie.
Stosując skończoną addytywność do $A$ i reszty
$B_N=\bigcup_{j>N}A_j$, dostajemy
$\|D_A(f)-\sum_{j=1}^ND_{A_j}(f)\|\le\int_{B_N}g\,\mathrm d\mu$.
Prawa strona dąży do zera z twierdzenia o zbieżności
zmajoryzowanej, ponieważ $B_N\downarrow\emptyset$.
Dla każdego skończonego podziału $A$ suma norm całek
jest ograniczona przez $\int_Ag\,\mathrm d\mu$;
supremum daje oszacowanie wariacji. Dowód dla Gelfanda
jest identyczny po zastąpieniu $x'(f)$ przez $h(\cdot)(y)$.
\end{proof}

\begin{twierdzenie}[Twierdzenie Orlicza--Pettisa]
\label{thm:orlicz-pettis15b}
Niech $X$ będzie przestrzenią Banacha nad $\K\in\{\R,\C\}$
i niech $(u_n)_{n\ge1}$ będzie ciągiem w $X$.
Załóżmy, że dla każdego $M\subseteq\N$ istnieje $v_M\in X$
taki, że $\sum_{n\in M,\ n\le N}u_n\to v_M$ w topologii
$\sigma(X,X')$ przy $N\to\infty$. Wtedy szereg
$\sum_{n\ge1}u_n$ jest bezwarunkowo zbieżny w normie,
a dla każdego $M$ suma podszeregu w normie jest równa $v_M$.
\end{twierdzenie}
\begin{proof}
Dla każdego $x'\in X'$ każdy podszereg
$\sum_{n\in M}x'(u_n)$ jest zbieżny. W konsekwencji
$\sum_n|x'(u_n)|<\infty$: w przypadku rzeczywistym
w przeciwnym razie podszereg wyrazów nieujemnych
lub niedodatnich byłby rozbieżny. W przypadku zespolonym
stosujemy to do części rzeczywistych i urojonych
(ich podszeregi są zbieżne), a następnie
$|z|\le|\operatorname{Re}z|+|\operatorname{Im}z|$.

Niech $E=\{0,1\}^{\N}$ z topologią produktową, a dla
$\alpha=(\alpha_n)\in E$ niech
$V(\alpha)=v_{\{n:\alpha_n=1\}}$. Przestrzeń $E$ jest
zupełna względem metryki
$d(\alpha,\beta)=\sum_n2^{-n}|\alpha_n-\beta_n|$:
ciąg Cauchy'ego stabilizuje się na każdej współrzędnej,
a zbieżność do ciągu tych ustalonych wartości wynika
z oszacowania ogona przez $\sum_{n>N}2^{-n}$.
Dla każdego $x'\in X'$ mamy
$x'(V(\alpha))=\sum_n\alpha_n x'(u_n)$. Szereg ten
zbiega jednostajnie na $E$, ponieważ jego ogon jest
ograniczony przez $\sum_{n>N}|x'(u_n)|$.
Sumy skończone są ciągłe, więc $x'\circ V$ jest ciągła.
To oznacza ciągłość $V$ z $E$ do $(X,\sigma(X,X'))$.

Niech $Z=\overline{\lin\{u_n:n\ge1\}}$ (domknięcie w normie).
Jest to ośrodkowa podprzestrzeń: kombinacje skończone
o współczynnikach wymiernych (w przypadku zespolonym
o wymiernych częściach rzeczywistych i urojonych)
tworzą w niej zbiór przeliczalny i gęsty. Podprzestrzeń
$Z$ jest słabo domknięta z twierdzenia Hahna--Banacha:
dla punktu spoza $Z$ istnieje ciągły funkcjonał zerujący
się na $Z$, lecz nie w tym punkcie. Wszystkie sumy
skończone podszeregów leżą w $Z$, a więc również ich
słabe granice $V(\alpha)$ leżą w $Z$.

Ustalmy $\varepsilon>0$ i wybierzmy gęsty ciąg $(z_j)$
w $Z$. Zbiory
$F_j=\{\alpha\in E:\|V(\alpha)-z_j\|\le\varepsilon/3\}$
pokrywają $E$. Są domknięte w $E$: domknięta kula
w $X$ jest słabo domknięta, ponieważ normę wyznacza
supremum modułów funkcjonałów z kuli jednostkowej $X'$,
a $V$ jest słabo ciągłe. Twierdzenie Baire'a daje indeks
$j$, dla którego $F_j$ ma niepuste wnętrze.
Bazę topologii produktowej tworzą zbiory określone przez
ustalenie skończonej liczby współrzędnych. Istnieją zatem
$N_0$ i wartości $\alpha_1,\ldots,\alpha_{N_0}$ takie,
że wszystkie ciągi o tych początkowych współrzędnych
należą do $F_j$.

Dla dowolnego skończonego
$H\subseteq\{N_0+1,N_0+2,\ldots\}$ wybierzmy dwa ciągi
z tego zbioru bazowego: pierwszy po $N_0$ ma jedynki
dokładnie na $H$, drugi same zera. Odpowiednie wartości
$V$ różnią się o $\sum_{n\in H}u_n$, co wynika
bezpośrednio z definicji sum skończonych i jednoznaczności
słabej granicy. Obie wartości leżą w tej samej kuli,
stąd
$\|\sum_{n\in H}u_n\|\le2\varepsilon/3<\varepsilon$.
Otrzymaliśmy kryterium Cauchy'ego dla sum po dowolnych
skończonych podzbiorach ogona szeregu. W szczególności
dla każdego $M$ sumy $\sum_{n\in M,\ n\le N}u_n$
są ciągiem Cauchy'ego w normie, więc mają granicę w $X$.
Musi być ona równa $v_M$, ponieważ zbieżność w normie
implikuje słabą i słaba granica jest jednoznaczna.

Dla dowolnej permutacji indeksów sumy początkowe po
dostatecznie wielu wyrazach zawierają wszystkie indeksy
do $N_0$; pozostałe składniki mają normę sumy mniejszą
od $\varepsilon$. Suma szeregu w naturalnej kolejności
różni się od sumy do $N_0$ o wektor normy co najwyżej
$2\varepsilon/3$, przez przejście do granicy.
W istocie każda dostatecznie długa suma po permutacji
różni się od sumy do $N_0$ o wektor normy co najwyżej
$2\varepsilon/3$, na mocy tego samego oszacowania dla $H$.
Jej odległość od sumy szeregu jest zatem co najwyżej
$4\varepsilon/3$. Ponieważ $\varepsilon>0$ było dowolne,
każda permutacja daje tę samą granicę w normie.
To jest bezwarunkowa zbieżność.
\end{proof}

\begin{twierdzenie}[Miara generowana przez całkę Pettisa]
\label{thm:pettis-vector-measure15b}
Niech $(S,\Sigma,\mu)$ będzie zupełną przestrzenią z miarą
skończoną, $X$ przestrzenią Banacha nad $\K\in\{\R,\C\}$,
a $f:S\to X$ funkcją całkowalną w sensie Pettisa.
Wówczas $m_f:\Sigma\to X$, $m_f(A)=P_A(f)$,
jest przeliczalnie addytywną w normie miarą wektorową,
bezwzględnie ciągłą względem $\mu$.
Dla każdego ciągu parami rozłącznych $A_j\in\Sigma$ szereg
$\sum_jP_{A_j}(f)$ jest bezwarunkowo zbieżny w normie
i ma sumę $P_{\bigcup_j A_j}(f)$.
Wariacja $|m_f|(S)$ nie musi być skończona.
\end{twierdzenie}
\begin{proof}
Ustalmy dowolny ciąg parami rozłącznych zbiorów
$(A_j)_{j\ge1}$ z $\Sigma$ i niech $u_j=P_{A_j}(f)$.
Dla dowolnego $M\subseteq\N$
zbiór $B_M=\bigcup_{j\in M}A_j$ jest mierzalny, a
$v_M=P_{B_M}(f)$ istnieje na mocy całkowalności Pettisa.
Dla każdego $x'\in X'$ addytywność całki skalarnej daje
\[
 \lim_{N\to\infty}\sum_{\substack{j\in M\\j\le N}}x'(u_j)
 =\int_{B_M}x'(f(s))\,\mathrm d\mu(s)=x'(v_M).
\]
Wszystkie podszeregi mają zatem słabe granice w $X$.
Z twierdzenia Orlicza--Pettisa otrzymujemy bezwarunkową
zbieżność w normie i dla $M=\N$ żądaną addytywność.
Całka po zbiorze miary zero jest równa zeru, ponieważ
wszystkie jej skalaryzacje są zerowe; dotyczy to również
zbioru pustego. Brak skończonej wariacji wykazaliśmy
w \hyperref[ex:pettis-infinite-variation15b]{przykładzie z $\ell^2$}.
\end{proof}

\begin{twierdzenie}[Półwariacja miary Pettisa i ciągłość całek]
\label{thm:pettis-semivariation15b}
Niech $(S,\Sigma,\mu)$ będzie zupełną przestrzenią z miarą
skończoną, $X$ przestrzenią Banacha nad $\K\in\{\R,\C\}$,
a $f:S\to X$ funkcją całkowalną w sensie Pettisa.
Niech $m_f(A)=P_A(f)$ dla $A\in\Sigma$. Definiujemy
\emph{półwariację} miary $m_f$ na $A$ jako
$\|m_f\|_{\mathrm{sv}}(A)=
\sup_{\|x'\|\le1}|x'\circ m_f|(A)$, gdzie
$|x'\circ m_f|$ oznacza wariację miary skalarnej
$B\mapsto x'(m_f(B))$.
Wówczas
\[
 \|m_f\|_{\mathrm{sv}}(A)=q_f(A)
 =\sup_{\|\varphi\|_\infty\le1}\|P_A(\varphi f)\|,
\]
gdzie supremum jest brane po $\varphi\in L^\infty(S,\mu;\K)$,
a $q_f$ określono w
\hyperref[thm:vector-integral-multipliers15b]{twierdzeniu o mnożnikach}.
Ponadto
\[
 \lim_{\delta\downarrow0}
 \sup_{\substack{A\in\Sigma\\\mu(A)\le\delta}}q_f(A)=0.
\]
W szczególności dla $A_n,A\in\Sigma$ warunek
$\mu(A_n\triangle A)\to0$ implikuje
$\|P_{A_n}(f)-P_A(f)\|\to0$.
\end{twierdzenie}
\begin{proof}
Miara skalarna $x'\circ m_f$ ma gęstość $x'\circ f$
względem $\mu$. Jej wariacja na $A$ wynosi
$\int_A|x'(f)|\,\mathrm d\mu$: górne oszacowanie wynika
z nierówności trójkąta dla całki skalarnej, dolne
otrzymujemy przez aproksymację funkcji fazowej
$\overline{x'(f)}/|x'(f)|$ na zbiorze, na którym mianownik
jest niezerowy (przyjmujemy zero w pozostałych punktach),
prostymi funkcjami o module nie większym niż $1$.
Dla takiej funkcji prostej moduł całki iloczynu jest
nie większy od sumy modułów całek po składnikach podziału,
czyli od wariacji. Jednostajna aproksymacja fazy daje
w granicy całkę modułu. Biorąc supremum po $x'$,
otrzymujemy pierwszą równość.
Oszacowanie z twierdzenia o mnożnikach daje
$\sup_{\|\varphi\|_\infty\le1}\|P_A(\varphi f)\|\le q_f(A)$.
Dla ustalonego $x'$ z kuli jednostkowej wybieramy
$\varphi$ równą opisanej wyżej fazie; wtedy
$x'(P_A(\varphi f))=\int_A|x'(f)|\,\mathrm d\mu$.
Supremum po $x'$ daje nierówność przeciwną.

Dla $A\in\Sigma$ połóżmy
$p(A)=\sup\{\|m_f(B)\|:B\in\Sigma,\ B\subseteq A\}$.
Mamy $p(A)\le q_f(A)\le4p(A)$.
Pierwsza nierówność wynika z oszacowania całki.
Dla drugiej ustalmy $x'$ o normie co najwyżej $1$ i
oznaczmy $u=x'\circ f$. W przypadku zespolonym
$\int_A|u|\le\int_A|\operatorname{Re}u|
+\int_A|\operatorname{Im}u|$.
Całkę modułu części rzeczywistej rozbijamy na zbiory,
na których jest ona nieujemna i ujemna. Każda z dwóch
otrzymanych całek jest modułem części rzeczywistej
$x'(m_f(B))$ dla pewnego mierzalnego $B\subseteq A$,
zatem jest nie większa niż $p(A)$. Tak samo postępujemy
z częścią urojoną. W przypadku rzeczywistym wystarczają
dwa składniki. Biorąc supremum po $x'$, uzyskujemy
żądane oszacowanie.

Pokażemy, że jeżeli $E_j\downarrow\emptyset$, to
$p(E_j)\to0$. Z przeliczalnej addytywności $m_f$ w normie
wynika najpierw $m_f(C_j)\to0$ dla każdego
$C_j\downarrow\emptyset$: rozkład $C_1$ na parami
rozłączne różnice $C_j\setminus C_{j+1}$ daje szereg
zbieżny w normie, którego ogon od $j$ jest równy $m_f(C_j)$.
Gdyby $p(E_j)$ nie dążyło do zera, to z monotoniczności
$p$ istniałoby $\varepsilon>0$ takie, że dla każdego $j$
można wybrać $B\subseteq E_j$ mierzalny z $\|m_f(B)\|>2\varepsilon$.
Rekurencyjnie wybieramy indeksy $j_1<j_2<\cdots$ i
$B_k\subseteq E_{j_k}$ z $\|m_f(B_k)\|>2\varepsilon$.
Po wyborze $B_k$ dobieramy $j_{k+1}$ tak, aby
$\|m_f(B_k\cap E_{j_{k+1}})\|<\varepsilon$, co jest
możliwe dzięki ciągłości dla malejących zbiorów.
Zbiory $D_k=B_k\setminus E_{j_{k+1}}$ są parami rozłączne,
a ze skończonej addytywności i nierówności trójkąta
wynika $\|m_f(D_k)\|>\varepsilon$. Przeczy to
przeliczalnej addytywności w normie, bo wyrazy szeregu
$\sum_km_f(D_k)$ zbieżnego w normie muszą dążyć do zera.
Zatem $p(E_j)\to0$ i również $q_f(E_j)\to0$.

Jeżeli $E_j\downarrow E$ oraz $\mu(E)=0$, to
$q_f(E_j)=q_f(E_j\setminus E)$, ponieważ całki skalarne
po $E$ znikają. Stosujemy wynik do $E_j\setminus E$.
Gdyby żądana jednostajna ciągłość względem $\mu$ nie
zachodziła, istniałoby $\varepsilon>0$ i zbiory $A_j$
takie, że $\mu(A_j)\le2^{-j}$, lecz $q_f(A_j)\ge\varepsilon$.
Niech $E_j=\bigcup_{k\ge j}A_k$. Wtedy $E_j$ maleją,
$\mu(E_j)\le\sum_{k\ge j}2^{-k}\to0$, więc ich
przecięcie ma miarę zero. Zatem $q_f(E_j)\to0$,
wbrew $q_f(E_j)\ge q_f(A_j)\ge\varepsilon$.
Na koniec skończona addytywność daje
$P_{A_n}(f)-P_A(f)=P_{A_n\setminus A}(f)-P_{A\setminus A_n}(f)$.
Norma tej różnicy jest nie większa niż
$2q_f(A_n\triangle A)$, co kończy dowód.
\end{proof}

\begin{uwaga}[Półwariacja w przykładzie atomowym]
Dla miary $m(A)=(\chi_A(n)/n)_{n\ge1}$ z
\hyperref[ex:pettis-infinite-variation15b]{przykładu nieskończonej wariacji}
nierówność Cauchy'ego--Schwarza daje
\[
 \|m\|_{\mathrm{sv}}(A)
 =\sup_{\|a\|_{\ell^2}\le1}\sum_{n\in A}\frac{|a_n|}{n}
 =\left(\sum_{n\in A}\frac1{n^2}\right)^{1/2}.
\]
Równość zachodzi, gdy $a$ jest wektorem
$(\chi_A(n)/n)_{n\ge1}$ podzielonym przez jego normę
(dla $A=\emptyset$ obie strony są zerowe). W szczególności
półwariacja na $S$ jest skończona, choć całkowita wariacja
na $S$ jest nieskończona.
\end{uwaga}

\begin{twierdzenie}[Zbieżność skalarna i zbieżność w normie całek]
\label{thm:vector-integral-convergence15b}
Niech $(S,\Sigma,\mu)$ będzie zupełną przestrzenią z miarą
skończoną, $X,Y$ przestrzeniami Banacha nad
$\K\in\{\R,\C\}$, $(f_n)$ ciągiem funkcji całkowalnych
w sensie Dunforda z $S$ do $X$, a $(h_n)$ ciągiem funkcji
całkowalnych w sensie Gelfanda z $S$ do $Y'$.
Niech $f:S\to X$ i $h:S\to Y'$ będą funkcjami.
Jeżeli dla każdego $x'\in X'$ funkcje $x'\circ f_n$
zbiegają prawie wszędzie do $x'\circ f$ i istnieje
$g_{x'}\in L^1(S,\mu)$, $g_{x'}\ge0$, taka że
$|x'\circ f_n|\le g_{x'}$ prawie wszędzie dla każdego $n$,
to $f$ jest całkowalna w sensie Dunforda i dla każdego
$A\in\Sigma$ zachodzi
$D_A(f_n)\to D_A(f)$ w topologii $\sigma(X'',X')$.
Analogicznie, jeśli dla każdego $y\in Y$ funkcje
$h_n(\cdot)(y)$ zbiegają prawie wszędzie do $h(\cdot)(y)$
i są zmajoryzowane przez pewną nieujemną funkcję
$g_y\in L^1(S,\mu)$, to $h$ jest całkowalna w sensie
Gelfanda i $G_A(h_n)\to G_A(h)$ w $\sigma(Y',Y)$.
Jeżeli w pierwszym przypadku wszystkie $f_n$ oraz $f$
są całkowalne w sensie Pettisa, to zachodzi również
$P_A(f_n)\to P_A(f)$ w $\sigma(X,X')$.

Jeżeli ponadto $q_{f_n-f}(S)\to0$, to całki Dunforda
zbiegają w normie jednostajnie względem $A\in\Sigma$;
w przypadku Pettisa ten sam warunek daje zbieżność
w normie całek Pettisa. Dla Gelfanda analogicznym
warunkiem jest $r_{h_n-h}(S)\to0$.
Wystarczający warunek dla takiej zbieżności to
$\|f_n-f\|\to0$ prawie wszędzie i mierzalność tych norm
oraz oszacowanie $\|f_n-f\|\le g$ prawie wszędzie dla
wspólnej nieujemnej $g\in L^1(S,\mu)$; analogicznie dla
$h_n-h$.
\end{twierdzenie}
\begin{proof}
Dla ustalonego $x'$ zupełność miary zapewnia mierzalność
granicy prawie wszędzie $x'\circ f$, a skalarne twierdzenie
o zbieżności zmajoryzowanej daje jej całkowalność i
$\int_Ax'(f_n)\,\mathrm d\mu\to\int_Ax'(f)\,\mathrm d\mu$.
Z definicji całki Dunforda oznacza to zadeklarowaną
zbieżność słabą-*. Jeżeli całki Pettisa istnieją,
ta sama równość jest dokładnie zbieżnością słabą
ich wartości. Dla Gelfanda stosujemy ten argument
do każdego $y\in Y$. Zbiory wyjątkowe mogą zależeć
od testującego funkcjonału lub wektora, ponieważ
argument prowadzimy oddzielnie dla każdej skalaryzacji.

Wszystkie różnice należą do odpowiednich klas
całkowalności na mocy liniowości. Dla każdego $A$
z oszacowań całek i monotoniczności $q$ oraz $r$ mamy
$\|D_A(f_n)-D_A(f)\|\le q_{f_n-f}(S)$ oraz
$\|G_A(h_n)-G_A(h)\|\le r_{h_n-h}(S)$.
Dla Pettisa zachodzi identyczne oszacowanie z $q$,
ponieważ normę w $X$ wyznaczają funkcjonały z kuli
jednostkowej $X'$. Prawa strona nie zależy od $A$.
Przy dodatkowych założeniach o mierzalności i dominacji
normy różnicy zachodzi
$q_{f_n-f}(S)\le\int_S\|f_n-f\|\,\mathrm d\mu\to0$
z twierdzenia o zbieżności zmajoryzowanej; analogicznie
dla $r_{h_n-h}$.
\end{proof}

\begin{uwaga}[Zbieżność skalarna nie zapewnia zbieżności w normie]
Na jednopunktowej przestrzeni z miarą masy $1$ niech
$f_n(s)=e_n\in\ell^2$. Wszystkie funkcje są całkowalne
nawet w sensie Bochnera; dla każdego funkcjonału
$x'\in(\ell^2)'$ mamy $x'(e_n)\to0$ z reprezentacji
Riesza, a $|x'(e_n)|\le\|x'\|$. Ich całki Pettisa są
równe $e_n$ i zbiegają słabo do zera, lecz
$\|e_n\|=1$. Całki Dunforda są równe $J_{\ell^2}(e_n)$,
więc również nie zbiegają w normie do zera.
Dla Gelfanda przyjmujemy na tej samej przestrzeni
$h_n(s)=e_n\in(\ell^1)'=\ell^\infty$. Dla $y\in\ell^1$
mamy $h_n(s)(y)=y_n\to0$ i $|y_n|\le\|y\|_{\ell^1}$,
więc całki Gelfanda zbiegają słabo-* do zera, ale
ich normy są stale równe $1$. Pokazuje to konieczność
odróżnienia dominacji poszczególnych skalaryzacji od
warunku jednostajnego względem kuli jednostkowej.
\end{uwaga}

\begin{stwierdzenie}[Zależność całki Gelfanda od całki Dunforda i Pettisa]
\label{prop:gelfand-dunford-restriction15b}
Niech $(S,\Sigma,\mu)$ będzie zupełną przestrzenią
z miarą skończoną, $Y$ przestrzenią Banacha nad
$\K\in\{\R,\C\}$, a $h:S\to Y'$ funkcją całkowalną
w sensie Gelfanda. Jeśli dodatkowo $h$, traktowana jako
funkcja o wartościach w przestrzeni Banacha $Y'$, jest
całkowalna w sensie Dunforda, to dla każdego $A\in\Sigma$
\[
 G_A(h)=J_Y' D_A(h),
 \qquad J_Y':Y'''\to Y',\quad
 (J_Y'F)(y)=F(J_Yy).
\]
Jeżeli $h$ jest całkowalna w sensie Pettisa w $Y'$, to
$G_A(h)=P_A(h)$ dla wszystkich $A\in\Sigma$.
Odwrotne implikacje dotyczące całkowalności Pettisa
nie zachodzą w ogólności.
\end{stwierdzenie}
\begin{proof}
Zanurzenie kanoniczne $J_Y:Y\to Y''$ spełnia
$(J_Yy)(y')=y'(y)$. Dla każdego $y\in Y$ mamy
$D_A(h)(J_Yy)=\int_A(J_Yy)(h(s))\,\mathrm d\mu(s)
=\int_Ah(s)(y)\,\mathrm d\mu(s)=G_A(h)(y)$.
To daje pierwszą równość. Jeżeli istnieje całka Pettisa
$P_A(h)\in Y'$, to zastosowanie definicji do funkcjonału
$J_Yy\in Y''$ daje $P_A(h)(y)=G_A(h)(y)$
dla każdego $y$, czyli drugą równość. Kontrprzykład
poniżej pokazuje, że istnienie całki Gelfanda (nawet
wraz z całkowalnością w sensie Dunforda) nie zapewnia
całkowalności Pettisa.
\end{proof}

\begin{przyklad}[Całka Gelfanda niebędąca całką Pettisa]
\label{ex:gelfand-not-pettis15b}
Niech $S=\N=\{1,2,\ldots\}$, $\Sigma=\mathcal P(\N)$,
$\mu(\{n\})=2^{-n}$, $Y=\ell^1(\N;\K)$ i
$Y'=\ell^\infty(\N;\K)$, $\K\in\{\R,\C\}$.
W \hyperref[ex:gelfand-weakstar15b]{przykładzie słabej-* addytywności}
rozważaliśmy $h(n)=2^ne_n$ i otrzymaliśmy $G_A(h)=\chi_A$.
Funkcja $h$ jest nawet całkowalna w sensie Dunforda
w przestrzeni $\ell^\infty$. Istotnie, jest silnie
mierzalna przez skończone obcięcia, a dla każdego
$\Phi\in(\ell^\infty)'$ oraz $N\ge1$ wybieramy skalary
$c_n$ o module co najwyżej $1$, dla których
$c_n\Phi(e_n)=|\Phi(e_n)|$. Wówczas
$\|\sum_{n=1}^Nc_ne_n\|_\infty\le1$, stąd
$\sum_{n=1}^N|\Phi(e_n)|
=|\Phi(\sum_{n=1}^Nc_ne_n)|\le\|\Phi\|$. Przejście do granicy
daje $\int_S|\Phi(h(n))|\,\mathrm d\mu(n)
=\sum_n|\Phi(e_n)|<\infty$.

Gdyby $h$ była całkowalna w sensie Pettisa, to jej
całka po $S$ musiałaby być równa $G_S(h)=(1,1,\ldots)$
na mocy poprzedniego stwierdzenia. Niech $c\subset\ell^\infty$
oznacza podprzestrzeń ciągów zbieżnych. Funkcjonał
$L:c\to\K$, $L((b_n))=\lim_n b_n$, jest liniowy
i ograniczony normą $1$. Twierdzenie Hahna--Banacha
daje jego ciągłe przedłużenie $\Phi$ na $\ell^\infty$.
Mamy $\Phi(e_n)=0$ dla każdego $n$, ponieważ
$e_n\in c$ ma granicę zero, natomiast
$\Phi((1,1,\ldots))=1$. Zatem
$\int_S\Phi(h(n))\,\mathrm d\mu(n)=0\ne1
=\Phi(G_S(h))$, co przeczy definicji całki Pettisa.
\end{przyklad}

\subsection{Twierdzenia Fubiniego dla całek słabych}

W twierdzeniu Fubiniego dla całki Bochnera silna mierzalność
pozwala wybrać jeden zbiór wyjątkowy dla mierzalności
przekrojów, a całkowalność normy na produkcie zapewnia
całkowalność normy prawie wszystkich przekrojów dzięki
skalarnemu twierdzeniu Tonellego.
Przy całkach słabych zbiory wyjątkowe ze skalarnego
twierdzenia Fubiniego mogą zależeć od funkcjonału.
Nie wolno bez dodatkowego argumentu zastąpić ich
nieprzeliczalnej rodziny jednym zbiorem miary zero.

\begin{twierdzenie}[Twierdzenie Fubiniego dla całki Gelfanda]
\label{thm:gelfand-fubini15b}
Niech $(S,\Sigma,\mu)$ i $(T,\mathcal T,\nu)$ będą
zupełnymi przestrzeniami z miarami skończonymi,
$\rho$ uzupełnieniem miary $\mu\otimes\nu$, a $Y$
przestrzenią Banacha nad $\K\in\{\R,\C\}$.
Niech $h:S\times T\to Y'$ będzie całkowalna w sensie
Gelfanda względem $\rho$. Załóżmy, że istnieją
mierzalne zbiory zerowe $N_S\subseteq S$, $N_T\subseteq T$
takie, że dla każdego $s\notin N_S$ funkcja $h(s,\cdot)$
jest całkowalna w sensie Gelfanda względem $\nu$,
a dla każdego $t\notin N_T$ funkcja $h(\cdot,t)$ jest
całkowalna w sensie Gelfanda względem $\mu$.
Dla $B\in\mathcal T$ i $A\in\Sigma$ określmy
\[
 H_B(s)=\int_B^{\mathrm G}h(s,t)\,\mathrm d\nu(t),
 \qquad K_A(t)=\int_A^{\mathrm G}h(s,t)\,\mathrm d\mu(s),
\]
przyjmując wartości zero odpowiednio na $N_S$ i $N_T$.
Wówczas $H_B:S\to Y'$ i $K_A:T\to Y'$ są całkowalne
w sensie Gelfanda oraz
\[
 \int_{A\times B}^{\mathrm G}h\,\mathrm d\rho
 =\int_A^{\mathrm G}H_B\,\mathrm d\mu
 =\int_B^{\mathrm G}K_A\,\mathrm d\nu .
\]
\end{twierdzenie}
\begin{proof}
Ustalmy $y\in Y$ oraz $B\in\mathcal T$.
Funkcja skalarna $u_y(s,t)=h(s,t)(y)$ jest całkowalna
względem $\rho$. Skalarne twierdzenie Fubiniego
dla uzupełnionej miary produktowej daje całkowalność
przekroju i mierzalność funkcji jego całek dla prawie
każdego $s$. Można stosować to twierdzenie, ponieważ
$u_y$ ma produktowo mierzalnego przedstawiciela
równego jej prawie wszędzie. Zbiór, na którym funkcja
i przedstawiciel się różnią, zawiera się w produktowo
mierzalnym zbiorze miary zero. Twierdzenie Tonellego
dla jego funkcji charakterystycznej daje przekroje
miary zero dla prawie każdego parametru; zupełność miar
czynników pozwala wówczas zastąpić przekroje przedstawiciela
przekrojami pierwotnej funkcji.
Na $S\setminus N_S$ z definicji całki Gelfanda zachodzi
$H_B(s)(y)=\int_B u_y(s,t)\,\mathrm d\nu(t)$.
Funkcja $H_B(\cdot)(y)$ jest więc prawie wszędzie
równa mierzalnej funkcji skalarnej z twierdzenia Fubiniego.
Zupełność $\mu$ daje jej mierzalność, a ponadto
\[
 \int_S|H_B(s)(y)|\,\mathrm d\mu(s)
 \le\int_{S\times B}|u_y|\,\mathrm d\rho<\infty.
\]
Ponieważ $y$ było dowolne, $H_B$ jest całkowalna
w sensie Gelfanda. Dla każdego $A\in\Sigma$ mamy
\[
 \left(\int_A^{\mathrm G}H_B\,\mathrm d\mu\right)(y)
 =\int_A\int_B u_y(s,t)\,\mathrm d\nu(t)\,\mathrm d\mu(s)
 =\int_{A\times B}u_y\,\mathrm d\rho.
\]
Prawa strona jest wartością całki Gelfanda funkcji $h$
po $A\times B$ na wektorze $y$. Równość dla wszystkich
$y\in Y$ daje równość funkcjonałów. Zamiana ról
przestrzeni $S$ i $T$ daje drugi porządek całkowania.
\end{proof}

\begin{twierdzenie}[Twierdzenie Fubiniego dla całki Dunforda]
\label{thm:dunford-fubini15b}
Niech $(S,\Sigma,\mu)$ i $(T,\mathcal T,\nu)$ będą
zupełnymi przestrzeniami z miarami skończonymi,
$\rho$ uzupełnieniem $\mu\otimes\nu$, a $X$ przestrzenią
Banacha nad $\K\in\{\R,\C\}$.
Niech $f:S\times T\to X$ będzie całkowalna w sensie
Dunforda względem $\rho$. Załóżmy, że poza mierzalnymi
zbiorami miary zero parametrów jej przekroje
$f(s,\cdot)$ i $f(\cdot,t)$ są całkowalne w sensie
Dunforda względem odpowiednich miar.
Dla $A\in\Sigma$ i $B\in\mathcal T$ połóżmy
$U_B(s)=D_B(f(s,\cdot))\in X''$ oraz
$V_A(t)=D_A(f(\cdot,t))\in X''$, z wartością zero
na odpowiednich zbiorach wyjątkowych.
Wówczas $U_B$ i $V_A$ są całkowalne w sensie Gelfanda
w przestrzeni sprzężonej $(X')'=X''$ oraz
\[
 D_{A\times B}(f)=
 \int_A^{\mathrm G}U_B\,\mathrm d\mu
 =\int_B^{\mathrm G}V_A\,\mathrm d\nu.
\]
Nie stwierdzamy tu całkowalności Dunforda funkcji
$U_B$ i $V_A$ względem wszystkich funkcjonałów z $X'''$.
\end{twierdzenie}
\begin{proof}
Zanurzenie kanoniczne $J_X:X\to X''$ daje funkcję
$h=J_X\circ f:S\times T\to (X')'$. Dla każdego $x'\in X'$
mamy $h(s,t)(x')=x'(f(s,t))$. Mierzalność i całkowalność
tych funkcji skalarnych są dokładnie założeniami
o całkowalności Dunforda funkcji $f$. Stąd $h$ jest
całkowalna w sensie Gelfanda z przestrzenią testującą
$Y=X'$, a dla dowolnego mierzalnego $C$ w uzupełnieniu
produktu jej całka Gelfanda po $C$ jest równa $D_C(f)$.
Ten sam argument dla przekrojów identyfikuje ich całki
Gelfanda z $U_B$ i $V_A$. Twierdzenie Fubiniego
dla Gelfanda zastosowane do $h$ daje wszystkie
stwierdzone wnioski.
\end{proof}

\begin{twierdzenie}[Twierdzenie Fubiniego dla całki Pettisa]
\label{thm:pettis-fubini15b}
Niech $(S,\Sigma,\mu)$ i $(T,\mathcal T,\nu)$ będą
zupełnymi przestrzeniami z miarami skończonymi,
$\rho$ uzupełnieniem $\mu\otimes\nu$, a $X$ przestrzenią
Banacha nad $\K\in\{\R,\C\}$.
Niech $f:S\times T\to X$ będzie całkowalna w sensie
Pettisa względem $\rho$. Załóżmy ponadto, że poza
mierzalnymi zbiorami miary zero parametrów przekroje
$f(s,\cdot)$ i $f(\cdot,t)$ są całkowalne w sensie
Pettisa względem odpowiednich miar.
Dla $A\in\Sigma$ i $B\in\mathcal T$ definiujemy
\[
 F_B(s)=\int_B^{\mathrm P}f(s,t)\,\mathrm d\nu(t),
 \qquad L_A(t)=\int_A^{\mathrm P}f(s,t)\,\mathrm d\mu(s),
\]
z wartością zero na odpowiednich zbiorach wyjątkowych.
Wówczas $F_B$ i $L_A$ są całkowalne w sensie Pettisa
i zachodzi
\[
 \int_{A\times B}^{\mathrm P}f\,\mathrm d\rho
 =\int_A^{\mathrm P}F_B\,\mathrm d\mu
 =\int_B^{\mathrm P}L_A\,\mathrm d\nu.
\]
\end{twierdzenie}
\begin{proof}
Całkowalność Pettisa implikuje całkowalność Dunforda,
a dla przekrojów mamy $U_B=J_X\circ F_B$ i
$V_A=J_X\circ L_A$ w oznaczeniach poprzedniego twierdzenia.
Zatem z całkowalności Gelfanda funkcji $U_B$ wynika,
że dla każdego $x'\in X'$ funkcja
$x'(F_B)=U_B(\cdot)(x')$ jest mierzalna i całkowalna
względem $\mu$. Nie wystarcza to jeszcze do całkowalności
Pettisa: potrzebna jest reprezentacja całek po wszystkich
mierzalnych podzbiorach $S$.

Niech więc $C\in\Sigma$ będzie dowolny. Z twierdzenia
Fubiniego dla Dunforda i całkowalności Pettisa na
produkcie otrzymujemy
\[
 D_C(F_B)=D_{C\times B}(f)
 =J_X\left(\int_{C\times B}^{\mathrm P}f\,\mathrm d\rho\right).
\]
Pierwsza równość wynika z testowania obu stron
na każdym $x'\in X'$; druga z definicji Pettisa.
W szczególności $D_C(F_B)\in J_X(X)$ dla każdego $C$.
Kryterium reprezentowalności całki Dunforda daje
całkowalność Pettisa funkcji $F_B$ i żądaną równość
dla $C=A$. Dla $L_A$ postępujemy analogicznie,
wykorzystując wszystkie prostokąty $A\times D$,
$D\in\mathcal T$.
\end{proof}

\begin{stwierdzenie}[Warunki zapewniające całkowalność przekrojów]
\label{prop:weak-fubini-separable15b}
Niech $(S,\Sigma,\mu)$, $(T,\mathcal T,\nu)$ będą
zupełnymi przestrzeniami z miarami skończonymi,
$\rho$ uzupełnieniem $\mu\otimes\nu$, a $Y$ ośrodkową
przestrzenią Banacha nad $\K\in\{\R,\C\}$.
Niech $h:S\times T\to Y'$ będzie słabo-* mierzalna
względem $\rho$. Załóżmy, że istnieje nieujemna
$\rho$-mierzalna funkcja $g\in L^1(\rho)$ taka, że
$\|h(s,t)\|\le g(s,t)$ dla $\rho$-prawie wszystkich
$(s,t)$. Wtedy $h$ i prawie wszystkie przekroje
w obu kierunkach są całkowalne w sensie Gelfanda,
więc zachodzi twierdzenie Fubiniego dla Gelfanda.
Nie trzeba osobno zakładać mierzalności normy $h$:
wynika ona z ośrodkowości $Y$ i słabej-* mierzalności $h$.

Jeżeli $X$ jest przestrzenią Banacha nad $\K$
o ośrodkowej przestrzeni sprzężonej $X'$,
$f:S\times T\to X$ jest słabo mierzalna względem $\rho$
i $\|f\|\le g$ prawie wszędzie dla takiego $g$,
to zachodzi analogiczny wniosek dla Dunforda.
Jeżeli dodatkowo $X$ jest refleksywna, to zachodzi
również wniosek dla Pettisa.
\end{stwierdzenie}
\begin{proof}
Dla każdego $y\in Y$ nierówność
$|h(s,t)(y)|\le\|y\|g(s,t)$ prawie wszędzie zapewnia
całkowalność skalarną na produkcie. Wybierzmy gęsty
ciąg $(y_k)$ w domkniętej kuli jednostkowej $Y$.
Dla każdego $(s,t)$ ciągłość funkcjonału $h(s,t)$ daje
$\|h(s,t)\|=\sup_k|h(s,t)(y_k)|$, więc norma jest
mierzalna jako supremum przeliczalnej rodziny funkcji
mierzalnych. Dla każdego $k$ skalarna
mierzalność na produkcie zapewnia mierzalność
przekrojów $h(s,\cdot)(y_k)$ dla prawie każdego $s$
i $h(\cdot,t)(y_k)$ dla prawie każdego $t$.
Można to uzasadnić przez wybór produktowo mierzalnego
przedstawiciela i zastosowanie skalarnego Fubiniego
do funkcji charakterystycznej mierzalnego zbioru
zerowego, zawierającego zbiór, na którym przedstawiciele
się różnią.

Ponieważ rodzina $(y_k)$ jest przeliczalna, można
usunąć jeden zbiór miary zero w $S$, poza którym
wszystkie te przekroje są mierzalne. Powiększamy
ten zbiór o zbiór wyjątkowy ze skalarnego Fubiniego
dla $g$ oraz o zbiór parametrów, dla których przekrój
zbioru, na którym nie zachodzi dominacja, nie jest
zerowy. Ściślej, ten ostatni zbiór w produkcie
zawieramy w produktowo mierzalnym zbiorze miary zero
i stosujemy do niego skalarne twierdzenie Fubiniego.
W konsekwencji dla każdego pozostałego $s$ funkcja
$g(s,\cdot)$ jest mierzalna i całkowalna, a
$\|h(s,t)\|\le g(s,t)$ dla $\nu$-prawie każdego $t$.

Ustalmy taki $s$ i dowolny $y\in Y$. Dla $y=0$
skalaryzacja jest zerowa. Dla $y\ne0$ wybieramy
dla każdego $j\ge1$ indeks $k_j$ taki, że
$\|y_{k_j}-y/\|y\|\|<1/j$.
Dla każdego $t\in T$ ciągłość funkcjonału $h(s,t)$
daje $\|y\|h(s,t)(y_{k_j})\to h(s,t)(y)$.
Funkcja ta jest zatem mierzalna jako granica funkcji
mierzalnych. Jej moduł jest prawie wszędzie ograniczony
przez $\|y\|g(s,\cdot)$, więc jest całkowalna.
To dowodzi całkowalności Gelfanda przekroju dla każdego
$s$ poza jednym zbiorem zerowym, niezależnym od $y$.
Zamiana ról $S$ i $T$ daje drugi kierunek.

Dla Dunforda stosujemy powyższy argument do
$Y=X'$ i $h=J_X\circ f$, zauważając, że
$\|J_X(f(s,t))\|=\|f(s,t)\|$. Jeżeli $X$ jest
refleksywna, to $J_X(X)=X''$, zatem każda całka
Dunforda na produkcie i w przekrojach jest reprezentowana
przez wektor z $X$. Wtedy stosujemy twierdzenie
Fubiniego dla Pettisa.
\end{proof}

\begin{przyklad}[Całkowalność na produkcie nie gwarantuje całkowalności przekrojów]
\label{ex:weak-fubini-failure15b}
Niech $S=T=[0,1]$ z zupełną miarą Lebesgue'a
$\lambda$, $\rho$ uzupełnieniem $\lambda\otimes\lambda$,
a $X=\ell^2([0,1];\K)$, $\K\in\{\R,\C\}$.
Dla $s\in[0,1]$ oznaczmy przez $e_s$ wektor jednostkowy
o jedynej niezerowej współrzędnej na indeksie $s$.
Definiujemy $k(0)=0$, $k(t)=1/t$ dla $t>0$, oraz
$f(s,t)=k(t)e_s$.

Dla każdego $a\in X$ zbiór $C_a=\{s:a_s\ne0\}$
jest co najwyżej przeliczalny (jak wykazaliśmy
w \hyperref[ex:pettis-not-bochner15b]{przykładzie nieośrodkowej przestrzeni Hilberta}).
Przez twierdzenie Riesza wszystkie skalaryzacje $f$
mają postać $k(t)\overline{a_s}$ i znikają poza
$C_a\times T$. Ten zbiór jest produktowo mierzalny
i ma miarę $\rho$ równą zero. Zupełność $\rho$
daje mierzalność skalaryzacji, a ponieważ jest ona
równa zeru prawie wszędzie, jej całka modułu wynosi zero.
Dla dowolnego mierzalnego $C$ w uzupełnieniu produktu
wektor zero reprezentuje wszystkie całki skalarne
po $C$. Zatem $f$ jest całkowalna w sensie Pettisa,
a więc także Dunforda, na produkcie, i wszystkie
jej całki Pettisa są zerowe.

Ustalmy jednak dowolne $s\in S$. Funkcjonał
$x_s'(v)=\langle v,e_s\rangle$ spełnia
$x_s'(f(s,t))=k(t)$, a $\int_T k\,\mathrm d\lambda=\infty$.
Dlatego żaden przekrój $f(s,\cdot)$ nie jest całkowalny
nawet w sensie Dunforda. Zbiór złych parametrów
jest całym $S$, chociaż dla każdego ustalonego
funkcjonału jego zbiory wyjątkowe mogą być zerowe.

Dla ustalonego $t$ przekrój $f(\cdot,t)$ jest natomiast
całkowalny w sensie Pettisa z całką zero na każdym
mierzalnym podzbiorze $S$: wszystkie skalaryzacje
znikają poza przeliczalnymi zbiorami. Jeden porządek
całkowania jest więc możliwy, drugi nie istnieje.
Przyjmując $Y=X'$ oraz $h=J_X\circ f:S\times T\to Y'$,
otrzymujemy również kontrprzykład dla Gelfanda.
Istotnie, dla każdego $x'\in Y=X'$ zachodzi
$h(s,t)(x')=x'(f(s,t))$. Wszystkie skalaryzacje
funkcji $h$ na produkcie są więc całkowalne,
lecz dla każdego $s$ skalaryzacja przekroju przez
$x_s'\in Y$ nie jest całkowalna. Konstrukcja działa
nad obydwoma ciałami skalarów i nie wymaga identyfikacji
przestrzeni Hilberta z jej przestrzenią sprzężoną.
\end{przyklad}

\subsection{Własność Radona--Nikodyma}

\begin{definicja}[Własność Radona--Nikodyma]
\label{def:rnp15b}
Przestrzeń Banacha $X$ nad $\K\in\{\R,\C\}$ ma
\emph{własność Radona--Nikodyma}
(RNP), jeżeli dla każdej skończonej przestrzeni z miarą
$(S,\Sigma,\mu)$ oraz każdej miary wektorowej
$m:\Sigma\to X$ o skończonej całkowitej wariacji,
bezwzględnie ciągłej względem $\mu$ (tzn. $m(A)=0$ dla
każdego $A\in\Sigma$ spełniającego $\mu(A)=0$), istnieje funkcja
$f\in L^1(S,\mu;X)$ taka, że
$m(A)=\int_Af(s)\,\mathrm d\mu(s)$ dla wszystkich $A\in\Sigma$.
Przy mierze niezupełnej całkowalność w sensie Bochnera
rozumiemy względem jej uzupełnienia. Jak uzasadnimy poniżej,
każda taka funkcja ma przedstawiciela będącego punktową
granicą funkcji prostych mierzalnych względem $\Sigma$;
dla tego przedstawiciela całki norm można obliczać
względem pierwotnej miary $\mu$.
Można też ograniczyć definicję RNP do miar zupełnych:
miara wektorowa $m$ rozszerza się na uzupełnienie przez
$\bar m(A)=m(B)$, gdy $B\in\Sigma$ i różnica symetryczna
$A\triangle B$ jest zawarta w mierzalnym zbiorze miary zero.
Jest to niezależne od $B$: dla dwóch przedstawicieli
$B,C\in\Sigma$ ich różnica symetryczna ma miarę zero,
więc $m(B\setminus C)=m(C\setminus B)=0$.
Niech $(A_j)$ będzie ciągiem parami rozłącznych zbiorów
z uzupełnienia, a $B_j\in\Sigma$ ich przedstawicielami.
Zbiory $C_j=B_j\setminus\bigcup_{i<j}B_i$ są parami
rozłączne, a $C_j$ jest nadal przedstawicielem $A_j$,
gdyż wszystkie przecięcia $B_j\cap B_i$ dla $i\ne j$
mają miarę zero. Ponadto $\bigcup_j C_j$ jest przedstawicielem
$\bigcup_j A_j$, zatem przeliczalna addytywność $m$ daje
$\bar m(\bigcup_j A_j)=\sum_j m(C_j)=\sum_j\bar m(A_j)$
w normie. Dla skończonych podziałów $S$ taki sam wybór
rozłącznych przedstawicieli daje
$\sum_j\|\bar m(A_j)\|=\sum_j\|m(C_j)\|\le |m|(S)$.
Ponieważ $\bar m$ rozszerza $m$, zachodzi
$|\bar m|(S)=|m|(S)<\infty$. Rozszerzenie jest bezwzględnie
ciągłe względem uzupełnionej miary, bo przedstawiciel
zbioru miary zero również ma miarę zero.
Gęstość względem uzupełnienia ma przedstawiciela
silnie mierzalnego względem $\Sigma$: zastępujemy zbiory
w ciągu funkcji prostych mierzalnymi przedstawicielami,
a granicę określamy na mierzalnym zbiorze punktów,
w których ciąg jest Cauchy'ego, i jako zero poza nim.
Mnożąc nowe funkcje proste przez funkcję charakterystyczną
tego zbioru, uzyskujemy zbieżność punktową do przedstawiciela
na całym $S$. Przedstawiciel jest równy pierwotnej gęstości
prawie wszędzie i ma tę samą całkę normy. Przybliżenia
w normie $L^1$ można otrzymać przez przycięcie funkcji
prostych jak w dowodzie kryterium całkowalności Bochnera.
Jeżeli $X$ ma RNP dla miar zupełnych, zastosowanie tej
własności do $\bar m$ daje więc gęstość dla pierwotnej
miary $m$; implikacja odwrotna jest natychmiastowa.
\end{definicja}

\begin{przyklad}[Brak RNP w przestrzeni funkcji ciągłych]
\label{ex:rnp-counterexample15b}
Niech $X=C([0,1])$ z normą supremum i niech $\mu$ będzie
miarą Lebesgue'a na $[0,1]$. Definiujemy
$m(A)(t)=\mu(A\cap[0,t])$ dla $t\in[0,1]$.
Dla $0\le u\le t\le1$ mamy
$|m(A)(t)-m(A)(u)|=\mu(A\cap(u,t])\le t-u$,
więc $m(A)$ jest ciągła. Skończona addytywność $m$ i
równość $m(\emptyset)=0$ wynikają ze skalarnej addytywności
$\mu$. Dla parami rozłącznych $A_n$ zachodzi
$\|m(\bigcup_{n>N}A_n)\|_\infty
=\mu(\bigcup_{n>N}A_n)\to0$,
a skończona addytywność zastosowana do sumy wszystkich
$A_n$ rozbitej na pierwszych $N$ zbiorów i ogon daje
addytywność przeliczalną w normie. Zatem $m$ jest miarą
wektorową. Co więcej, $\|m(A)\|_\infty=\mu(A)$, więc
$|m|(A)=\mu(A)$; wariacja jest skończona i
bezwzględnie ciągła względem $\mu$.

Gdyby $m(A)=\int_Af(s)\,\mathrm d\mu(s)$ dla pewnej funkcji
$f\in L^1([0,1];C([0,1]))$, to dla każdego wymiernego
$q\in[0,1]$ i wszystkich zbiorów mierzalnych $A$
mielibyśmy
$\int_A f(s)(q)\,\mathrm d\mu(s)
=\int_A\chi_{[0,q]}(s)\,\mathrm d\mu(s)$.
Stąd $f(s)(q)=\chi_{[0,q]}(s)$ dla prawie wszystkich $s$.
Po usunięciu jednego wspólnego zbioru miary zero równość
zachodzi dla wszystkich wymiernych $q$ jednocześnie.
Dla takiego $s\in(0,1)$ wybieramy ciągi wymiernych
$q_n<s<r_n$ zbieżne do $s$. Wówczas $f(s)(q_n)=0$
i $f(s)(r_n)=1$ dla każdego $n$. Ciągłość $f(s)$
w punkcie $s$ wymusza jednocześnie $f(s)(s)=0$ i $f(s)(s)=1$.
Otrzymaliśmy sprzeczność: przestrzeń $C([0,1])$ nie ma RNP.
\end{przyklad}

\begin{uwaga}[Rozszerzenie na przestrzenie lokalnie wypukłe]
\label{rem:lcs-pettis15b}
Dla przestrzeni lokalnie wypukłej Hausdorffa $E$ sensowna
pozostaje \emph{całka Pettisa w $E$}: dla funkcji $f:S\to E$
wymagamy całkowalności wszystkich skalaryzacji $\ell\circ f$,
$\ell\in E'$, i dla każdego $A\in\Sigma$ istnienia $v_A\in E$
takiego, że $\ell(v_A)=\int_A\ell(f(s))\,\mathrm d\mu(s)$.
Rozdzielanie punktów przez $E'$ zapewnia jednoznaczność
$v_A$. Jeżeli $T:E\to F$ jest ciągłym odwzorowaniem liniowym
między przestrzeniami lokalnie wypukłymi Hausdorffa,
to $Tf$ jest całkowalna w sensie Pettisa w $F$ i
$T(v_A)=\int_A^{\mathrm P} T(f(s))\,\mathrm d\mu(s)$:
dla każdego $\varphi\in F'$ wystarczy zastosować
$\varphi\circ T\in E'$. Nie zakładamy jednak, że sama
zupełność $E$ wymusza istnienie całek Pettisa wszystkich
funkcji słabo mierzalnych: jest to własność mocniejsza.
\end{uwaga}

\begin{uwaga}[Zakres teorii]
W osobnym skrypcie teorii miary rozwijamy skalarne podstawy
całkowania, natomiast tu wyjaśniliśmy różnice pomiędzy
całkowaniem w normie, całkowaniem skalarnym z wartościami
w przestrzeni podwójnie sprzężonej i całkowaniem słabym-* do przestrzeni sprzężonej.
Równania różniczkowe wykorzystujące te konstrukcje należą
do osobnego skryptu PDE. Ujęcie całek dla ogólnych przestrzeni
liniowo-topologicznych nie musi być oparte na seminormach;
bez lokalnej wypukłości argumenty Hahna--Banacha nie są
automatycznie dostępne.
\end{uwaga}

\addtocounter{section}{-1}
\endgroup

% Wersja źródłowa: audyt II, jednostka 16
\setcounter{section}{15}
\section{Funkcjonały wypukłe i podstawy rachunku wariacyjnego}
\label{sec:wariacyjny16}

\subsection{Od pochodnej Frécheta do problemu minimum}

W poprzednich rozdziałach badaliśmy lokalną strukturę funkcji
różniczkowalnych na przestrzeniach Banacha. Rachunek wariacyjny pyta
o istnienie i charakter minimizatorów funkcjonału $J:E\to\R$,
często na nieskończenie wymiarowej przestrzeni funkcji.
W tym rozdziale $H$ oznacza rzeczywistą przestrzeń Hilberta,
a $E$ rzeczywistą przestrzeń Banacha. Pisząc $\langle\cdot,\cdot\rangle$
w $H$, przyjmujemy liniowość w pierwszym argumencie.

\begin{definicja}[Wypukłość funkcjonału]
\label{def:convex-fun16}
Niech $E$ będzie rzeczywistą przestrzenią liniową,
a $C\subseteq E$ zbiorem wypukłym.
Funkcjonał $J:C\to(-\infty,+\infty]$ nazywamy
\emph{wypukłym}, jeżeli
$J((1-t)x+ty)\le(1-t)J(x)+tJ(y)$ dla wszystkich $x,y\in C$
i $0<t<1$. Dla $t=0,1$ przyjmujemy równość, bez
mnożenia zera przez nieskończoność.
Funkcjonał nazywamy \emph{właściwym}, gdy jego dziedzina
efektywna $\operatorname{dom}J=\{x\in C:J(x)<+\infty\}$
jest niepusta. Z wypukłości $J$ wynika wypukłość
$\operatorname{dom}J$, ponieważ kombinacja wypukła
dwóch skończonych wartości jest skończona.
Funkcjonał wypukły jest \emph{ściśle wypukły}, jeśli nierówność
jest ostra dla różnych $x,y\in\operatorname{dom}J$
i $0<t<1$. W przypadku funkcjonału rzeczywistego
$\operatorname{dom}J=C$.
\end{definicja}

\begin{twierdzenie}[Nierówność styczna]
\label{thm:convex-derivative16}
Niech $E$ będzie rzeczywistą przestrzenią Banacha,
$U\subseteq E$ niepustym otwartym zbiorem wypukłym, a
$J:U\to\R$ funkcjonałem różniczkowalnym w sensie
Frécheta w każdym punkcie $U$.
Następujące warunki są równoważne:
\begin{enumerate}
\item $J$ jest wypukły na $U$;
\item $J(y)\ge J(x)+DJ(x)(y-x)$ dla wszystkich $x,y\in U$.
\end{enumerate}
Jeżeli $DJ(x_0)=0$, gdzie $x_0\in U$, a $J$ jest wypukły,
to $x_0$ jest globalnym minimizatorem $J$ na $U$.
Odwrotnie, jeśli $x_0$ jest lokalnym minimizatorem
różniczkowalnego $J$ i $x_0\in U$, to $DJ(x_0)=0$.
\end{twierdzenie}
\begin{proof}
Dla ustalonych $x,y\in U$ niech
$\phi(t)=J(x+t(y-x))$, $0\le t\le1$.
Jeżeli $J$ jest wypukły, to dla $0<t\le1$
z definicji wypukłości
$[\phi(t)-\phi(0)]/t\le\phi(1)-\phi(0)$.
Z reguły łańcuchowej (Twierdzenie~\ref{thm:chain14})
pochodna $\phi$ w zerze wynosi $DJ(x)(y-x)$
(otwartość $U$ pozwala przedłużyć parametr nieco poza zero).
Przechodząc z $t\downarrow0$, otrzymujemy
$DJ(x)(y-x)\le J(y)-J(x)$.

Odwrotnie, ustalmy $x,y\in U$, $0<t<1$ i
$z=(1-t)x+ty\in U$.
Nierówność z punktu 2 daje
$J(x)\ge J(z)+DJ(z)(x-z)$ oraz
$J(y)\ge J(z)+DJ(z)(y-z)$.
Mnożąc je odpowiednio przez $1-t$ i $t$ oraz dodając,
otrzymujemy
\[
 (1-t)J(x)+tJ(y)\ge J(z)+DJ(z)((1-t)x+ty-z)=J(z).
\]
Dla $t=0,1$ zachodzi równość, więc $J$ jest wypukły.
Jeżeli $DJ(x_0)=0$, nierówność styczna zapewnia
$J(y)\ge J(x_0)$. Jeśli zaś $x_0$ jest lokalnym minimum,
dla każdego $h\in E$ i dostatecznie małych $|t|$
punkt $x_0+th$ należy do $U$ i do otoczenia lokalnego
minimum. Funkcja $t\mapsto J(x_0+th)$ ma zatem
minimum w zerze, a jej pochodna w zerze, równa
$DJ(x_0)h$, musi znikać. Ponieważ dotyczy to
każdego $h\in E$, mamy $DJ(x_0)=0$.
\end{proof}

\begin{definicja}[Silna wypukłość]
\label{def:strong-convex16}
Niech $H$ będzie rzeczywistą przestrzenią Hilberta
i $\mu>0$. Funkcjonał $J:H\to\R$ jest
$\mu$-silnie wypukły, jeśli $x\mapsto J(x)-\frac\mu2\|x\|^2$
jest wypukły.
\end{definicja}

\begin{stwierdzenie}[Koercywność z silnej wypukłości]
\label{prop:strong-coercive16}
Niech $H$ będzie rzeczywistą przestrzenią Hilberta,
$\mu>0$, a $J:H\to\R$ funkcjonałem $\mu$-silnie wypukłym
i różniczkowalnym w sensie Frécheta na $H$. Wtedy dla $x,y\in H$
zachodzi
\[
J(y)\ge J(x)+DJ(x)(y-x)+\frac\mu2\|y-x\|^2.
\]
W szczególności $J$ jest koercywny, tj.
$J(y)\to+\infty$, gdy $\|y\|\to\infty$,
i ma co najwyżej jeden minimizator.
\end{stwierdzenie}
\begin{proof}
Stosujemy Twierdzenie~\ref{thm:convex-derivative16}
do funkcjonału $G(z)=J(z)-\mu\|z\|^2/2$.
Z rozwinięcia
$\|x+h\|^2=\|x\|^2+2\langle x,h\rangle+\|h\|^2$
wynika $DG(x)h=DJ(x)h-\mu\langle x,h\rangle$.
Nierówność styczna dla $G$ daje
\[
 J(y)-\tfrac\mu2\|y\|^2
 \ge J(x)-\tfrac\mu2\|x\|^2
       +DJ(x)(y-x)-\mu\langle x,y-x\rangle .
\]
Po przeniesieniu składników kwadratowych na prawą stronę
i wykorzystaniu tożsamości
$\|y\|^2-\|x\|^2-2\langle x,y-x\rangle=\|y-x\|^2$
otrzymujemy żądaną nierówność.
Dla $x=0$ mamy
$J(y)\ge J(0)-\|DJ(0)\|\|y\|+\mu\|y\|^2/2$,
co zapewnia koercywność. Jeżeli $x\ne y$ są minimizatorami,
to wypukłość $G$ przy $t=1/2$ i tożsamość
$\tfrac12(\|x\|^2+\|y\|^2)-\|(x+y)/2\|^2
=\|x-y\|^2/4$ dają
$J((x+y)/2)\le [J(x)+J(y)]/2-\mu\|x-y\|^2/8$,
sprzecznie z minimalnością.
\end{proof}

\begin{przyklad}[Niezbędne jednostajne oszacowanie hesjanu]
\label{ex:variational-hessian16}
Niech $H$ będzie rzeczywistą przestrzenią Hilberta
i $J\in C^2(H,\R)$. Jeżeli istnieje $\mu>0$ takie, że
$D^2J(x)[h,h]\ge\mu\|h\|^2$ dla wszystkich $x,h\in H$,
to $J$ jest $\mu$-silnie wypukły.
Istotnie, dla $x,y\in H$ i $h=y-x$ funkcja
$\phi(t)=J(x+th)-\frac\mu2\|x+th\|^2$
ma drugą pochodną
$\phi''(t)=D^2J(x+th)[h,h]-\mu\|h\|^2\ge0$.
Zatem $\phi'$ jest niemalejąca. Dla $0<t<1$
z twierdzenia o wartości średniej na $[0,t]$ i $[t,1]$
wynika
$(\phi(t)-\phi(0))/t\le(\phi(1)-\phi(t))/(1-t)$.
Po pomnożeniu przez $t(1-t)$ i uporządkowaniu
otrzymujemy $\phi(t)\le(1-t)\phi(0)+t\phi(1)$,
czyli wypukłość funkcjonału $J-\mu\|\cdot\|^2/2$.
Sama dodatnia określoność $D^2J(0)$, nawet wraz
z $DJ(0)=0$, nie gwarantuje lokalnego minimum:
\hyperref[ex:positive-hessian15]{kontrprzykład w rozdziale 15}
podaje dokładnie taki funkcjonał na $\ell^2$.
\end{przyklad}

\subsection{Zbieżność słaba i metoda bezpośrednia}

\begin{definicja}[Półciągłość z dołu i koercywność]
\label{def:weak-lsc16}
Niech $H$ będzie rzeczywistą przestrzenią Hilberta,
$C\subseteq H$ niepustym zbiorem i $J:C\to(-\infty,+\infty]$.
Funkcjonał $J$ jest \emph{słabo półciągły z dołu}, jeśli dla każdej sieci
$(x_\alpha)$ w $C$ zbieżnej do $x\in C$ w topologii
słabej $\sigma(H,H')$ zachodzi
$J(x)\le\liminf_\alpha J(x_\alpha)$, gdzie
$\liminf_\alpha J(x_\alpha)=
\sup_\alpha\inf_{\beta\ge\alpha}J(x_\beta)$.
Równoważnie, każdy zbiór podpoziomowy
$\{x\in C:J(x)\le a\}$, $a\in\R$, jest domknięty
w topologii słabej indukowanej na $C$.
Istotnie, warunek dla sieci zapewnia, że granica sieci
z podpoziomu nadal należy do tego podpoziomu.
Odwrotnie, gdyby $\liminf_\alpha J(x_\alpha)<J(x)$,
można byłoby wybrać rzeczywiste $a$ ściśle pomiędzy
tymi wartościami. Otwarte dopełnienie podpoziomu
zawiera $x$, więc od pewnego indeksu $J(x_\alpha)>a$,
co wymusza $\liminf_\alpha J(x_\alpha)\ge a$,
sprzecznie z wyborem $a$.
Analogicznie definiuje się półciągłość z dołu
w topologii normowej.
Funkcjonał jest \emph{koercywny na $C$}, jeśli dla
każdego $a\in\R$ istnieje $R>0$ takie, że
$J(x)>a$ dla każdego $x\in C$ o $\|x\|>R$.
Równoważnie, wszystkie podpoziomy są ograniczone.
\end{definicja}

\begin{lemat}[Podpoziomowe zbiory wypukłe]
\label{lem:convex-weaklsc16}
Niech $H$ będzie rzeczywistą przestrzenią Hilberta,
a $J:H\to(-\infty,+\infty]$ funkcjonałem wypukłym
i półciągłym z dołu w topologii normowej.
Wtedy $J$ jest słabo półciągły z dołu.
\end{lemat}
\begin{proof}
Dla każdego $a\in\R$ zbiór $C_a=\{x:J(x)\le a\}$
jest normowo domknięty przez półciągłość z dołu
i wypukły przez wypukłość $J$. Z Wniosku~\ref{cor:weak-convex-closure}
o równości domknięć zbioru wypukłego wynika,
że każdy normowo domknięty zbiór wypukły jest
domknięty także w topologii słabej. Zatem wszystkie $C_a$
są słabo domknięte, co daje tezę.
\end{proof}

\begin{lemat}[Słaba zbieżność podciągu w przestrzeni Hilberta]
\label{lem:hilbert-weakseq16}
Niech $H$ będzie rzeczywistą przestrzenią Hilberta,
a $(x_n)$ ciągiem ograniczonym w $H$.
Istnieje podciąg $(x_{n_k})$ i $x\in H$ takie,
że $x_{n_k}\to x$ słabo w $H$.
Nie zakładamy ośrodkowości całej przestrzeni $H$.
\end{lemat}
\begin{proof}
Niech $M=\sup_n\|x_n\|<\infty$ i
$H_0=\overline{\operatorname{lin}\{x_n:n\in\N\}}$.
Przestrzeń $H_0$ jest ośrodkowa, ponieważ skończone
kombinacje liniowe $x_n$ o współczynnikach wymiernych
tworzą w niej przeliczalny zbiór gęsty. Jest również
przestrzenią Hilberta jako domknięta podprzestrzeń $H$.
Jeśli $H_0=\{0\}$, cały ciąg jest zerowy i teza zachodzi.
W przeciwnym razie wybieramy bazę ortonormalną $(e_j)$
skończoną lub przeliczalnie nieskończoną.
Dla każdego $j$ ciąg $\langle x_n,e_j\rangle$ jest ograniczony.
Z twierdzenia Bolzana--Weierstrassa wybieramy podciąg,
na którym pierwsza współrzędna ma granicę. Z niego
wybieramy podciąg dla drugiej współrzędnej i tak dalej.
W przypadku skończonej bazy ostatni podciąg ma
wszystkie wymagane granice. W przypadku nieskończonej
bazy wybieramy $k$-ty wyraz $k$-tego podciągu.
Indeksy tak wybranego ciągu rosną, a jego ogon
od $k=j$ jest podciągiem $j$-tego podciągu, więc
$j$-ta współrzędna jest zbieżna dla każdego $j$.

Na przestrzeni skończonych kombinacji wektorów $e_j$
definiujemy $L(v)=\lim_k\langle x_{n_k},v\rangle$.
Istnienie granicy wynika z jej istnienia na każdym $e_j$,
a $|L(v)|\le M\|v\|$. Funkcjonał $L$ przedłuża się
jednoznacznie na $H_0$ jako ciągły funkcjonał liniowy.
Z twierdzenia Riesza istnieje $x\in H_0$ takie,
że $L(v)=\langle x,v\rangle$ dla $v\in H_0$.
Ponieważ $\|x\|=\|L\|\le M$, dla dowolnego $v\in H_0$
i skończonej kombinacji $w$ wektorów $e_j$ zachodzi
\[
 |\langle x_{n_k}-x,v\rangle|
 \le 2M\|v-w\|+|\langle x_{n_k}-x,w\rangle|.
\]
Dla $\varepsilon>0$ wybieramy $w$ tak, aby pierwszy
składnik był mniejszy od $\varepsilon/2$, a następnie
$k$ dostatecznie duże, aby drugi składnik był mniejszy
od $\varepsilon/2$. Gdy $M=0$, ciąg jest zerowy.
Otrzymujemy zbieżność $\langle x_{n_k},v\rangle
\to\langle x,v\rangle$ dla każdego $v\in H_0$.
Dla $v\in H_0^\perp$ obie strony są zerowe.
Dla dowolnego $v\in H$ stosujemy rozkład
$H=H_0\oplus H_0^\perp$ z Twierdzenia~\ref{thm:orth-decomposition13}.
Z twierdzenia Riesza każdy ciągły funkcjonał liniowy
na $H$ jest dany przez iloczyn skalarny z takim $v$.
Otrzymana zbieżność jest więc zbieżnością słabą w $H$.
\end{proof}

\begin{twierdzenie}[Metoda bezpośrednia w przestrzeni Hilberta]
\label{thm:direct-method16}
Niech $H$ będzie rzeczywistą przestrzenią Hilberta,
$C\subseteq H$ niepustym, normowo domkniętym i wypukłym
zbiorem, a $J:C\to(-\infty,+\infty]$ funkcjonałem,
który jest właściwy (tzn. przyjmuje przynajmniej jedną
wartość skończoną), ograniczony z dołu, koercywny na $C$
i półciągły z dołu dla topologii słabej indukowanej na $C$.
Wtedy istnieje $u\in C$ takie, że
$J(u)=\inf_{v\in C}J(v)$.
Jeżeli dodatkowo $J$ jest ściśle wypukły,
minimizator jest jedyny.
\end{twierdzenie}
\begin{proof}
Niech $m=\inf_C J$, który jest skończony na mocy
właściwości i ograniczoności z dołu. Dla każdego $n\ge1$ z definicji infimum wybieramy
$u_n\in C$ takie, że $m\le J(u_n)<m+1/n$.
Wszystkie $u_n$ należą do podpoziomu $\{v\in C:J(v)\le m+1\}$,
który jest ograniczony z koercywności. Zatem
$(u_n)$ jest ograniczonym ciągiem w $H$.
Na mocy Lematu~\ref{lem:hilbert-weakseq16}
pewien podciąg $u_{n_k}\to u$ słabo w $H$.
Ponieważ $C$ jest normowo domknięty i wypukły,
jest również słabo domknięty na mocy
Wniosku~\ref{cor:weak-convex-closure}, więc $u\in C$.
Słaba półciągłość z dołu daje
$J(u)\le\liminf_k J(u_{n_k})=m$; z definicji infimum
zachodzi też $J(u)\ge m$. Jeśli $u\ne v$ są minimizatorami,
ścisła wypukłość daje $J((u+v)/2)<m$, sprzeczność.
\end{proof}

\begin{uwaga}[Czego nie daje sama koercywność?]
\label{rem:coercive16}
W nieskończenie wymiarowej przestrzeni Hilberta
domknięta kula jednostkowa nie jest zwarta w normie:
zawiera ciąg ortonormalny, którego różne wyrazy
są odległe o $\sqrt2$, więc żaden jego podciąg
nie jest zbieżny w normie.
Sama ciągłość normowa i koercywność funkcjonału
nie dają zatem zwartości podpoziomu potrzebnej
do bezpośredniego zastosowania twierdzenia Weierstrassa. W twierdzeniu wykorzystaliśmy
słabą zbieżność podciągów i słabą półciągłość z dołu.
W ogólnej przestrzeni Banacha naturalną rolę przestrzeni
Hilberta przejmują przestrzenie refleksywne.
W rozdziale 26 udowodnimy ogólną wersję twierdzenia
o słabo zbieżnym podciągu ciągu ograniczonego.
Powyższy elementarny dowód w przestrzeni Hilberta
nie wymaga tej późniejszej teorii.
\end{uwaga}

\subsection{Subgradienty i rzutowanie proksymalne}

\begin{definicja}[Subgradient i subróżniczka]
\label{def:subgrad16}
Niech $J:H\to(-\infty,+\infty]$ będzie właściwym
funkcjonałem wypukłym, a $u\in H$ punktem takim,
że $J(u)<+\infty$. Wektor $p\in H$ nazywamy
\emph{subgradientem} $J$ w $u$, jeśli
$J(v)\ge J(u)+\langle p,v-u\rangle$ dla każdego $v\in H$.
Zbiór wszystkich subgradientów oznaczamy $\partial J(u)$
i nazywamy \emph{subróżniczką} $J$ w $u$.
Dla $J(u)=+\infty$ przyjmujemy $\partial J(u)=\varnothing$.
\end{definicja}

\begin{stwierdzenie}[Warunek optymalności i monotoniczność]
\label{prop:subgrad16}
Jeżeli $J:H\to(-\infty,+\infty]$ jest właściwy i wypukły,
to $u$ minimalizuje $J$ wtedy i tylko wtedy, gdy
$0\in\partial J(u)$. Jeśli $p\in\partial J(u)$ oraz
$q\in\partial J(v)$, to
$\langle p-q,u-v\rangle\ge0$.
Jeżeli $J$ jest różniczkowalny Frécheta w $u$ i skończony
na otoczeniu $u$, to $\partial J(u)=\{\nabla J(u)\}$.
\end{stwierdzenie}
\begin{proof}
Pierwsza równoważność wynika wprost z definicji.
Z nierówności subgradientowych
$J(v)\ge J(u)+\langle p,v-u\rangle$ i
$J(u)\ge J(v)+\langle q,u-v\rangle$,
po dodaniu dostajemy monotoniczność.
Jeśli $J$ jest różniczkowalny, to dla $p\in\partial J(u)$
podstawiamy $v=u+th$ z dodatnim i ujemnym $t$,
po czym przechodzimy do zera. Otrzymujemy
$\langle p,h\rangle=DJ(u)h$ dla każdego $h\in H$.
Odwrotnie, dla $v$ z dziedziny efektywnej $J$ i $0<t<1$
wypukłość daje
$(J(u+t(v-u))-J(u))/t\le J(v)-J(u)$.
Przechodząc z $t$ do zera, otrzymujemy
$DJ(u)(v-u)\le J(v)-J(u)$.
Dla $J(v)=+\infty$ nierówność zachodzi automatycznie.
Reprezentacja Riesza daje zatem $\nabla J(u)\in\partial J(u)$.
\end{proof}

\begin{przyklad}[Subróżniczka normy w Hilbercie]
\label{ex:norm-subgrad16}
Dla $J(x)=\|x\|$ zachodzi
$\partial J(0)=\{p:\|p\|\le1\}$ i
$\partial J(x)=\{x/\|x\|\}$ dla $x\ne0$.
Jeżeli $\|p\|\le1$, to $\langle p,v\rangle\le\|v\|$
dla każdego $v\in H$, zatem $p\in\partial J(0)$.
Odwrotnie, dla $p\in\partial J(0)$, podstawiając $v=p$,
otrzymujemy $\|p\|^2\le\|p\|$, czyli $\|p\|\le1$.
Dla $x\ne0$ reguła łańcuchowa zastosowana do
$\|x\|=\sqrt{\langle x,x\rangle}$ daje
$DJ(x)h=\langle x,h\rangle/\|x\|$.
Stwierdzenie~\ref{prop:subgrad16} daje drugi wzór.
W ogólnej przestrzeni Banacha analogiczny subgradient normy
jest funkcjonałem z przestrzeni dualnej, a jego istnienie
dla $x\ne0$ zapewnia twierdzenie Hahna--Banacha.
\end{przyklad}

\begin{twierdzenie}[Odwzorowanie proksymalne]
\label{thm:prox16}
Niech $J:H\to(-\infty,+\infty]$ będzie właściwy,
wypukły, normowo półciągły z dołu i ograniczony z dołu,
a $\lambda>0$. Dla każdego $y\in H$ istnieje dokładnie
jeden punkt $u=\operatorname{prox}_{\lambda J}(y)$
minimalizujący funkcjonał
$v\mapsto J(v)+\|v-y\|^2/(2\lambda)$.
Jest on scharakteryzowany równoważnością
$y-u\in\lambda\partial J(u)$.
Ponadto odwzorowanie
$\operatorname{prox}_{\lambda J}:H\to H$ jest
nieekspansywne:
$\|\operatorname{prox}_{\lambda J}(y)-
\operatorname{prox}_{\lambda J}(z)\|\le\|y-z\|$.
\end{twierdzenie}
\begin{proof}
Funkcjonał $\Phi_y(v)=J(v)+\|v-y\|^2/(2\lambda)$
jest właściwy, wypukły i normowo półciągły z dołu,
a więc słabo półciągły z dołu na mocy
Lematu~\ref{lem:convex-weaklsc16}.
Jest koercywny, gdyż $J$ jest ograniczony z dołu,
i jest ściśle wypukły dzięki składnikowi kwadratowemu.
Twierdzenie~\ref{thm:direct-method16} dla $C=H$
daje jednoznaczny minimizator $u$.

Jeśli $u$ minimalizuje $\Phi_y$, to dla $v$ z
$J(v)<+\infty$ i $0<t<1$, z wypukłości $J$
oraz minimalności mamy
\[
0\le\Phi_y(u+t(v-u))-\Phi_y(u)
\le t[J(v)-J(u)]+\frac t\lambda\langle u-y,v-u\rangle
+\frac{t^2}{2\lambda}\|v-u\|^2.
\]
Po podzieleniu przez $t$ i przejściu z $t$ do zera
otrzymujemy $J(v)\ge J(u)+\langle(y-u)/\lambda,v-u\rangle$,
a więc $(y-u)/\lambda\in\partial J(u)$.
Odwrotnie, ta nierówność oraz rozwinięcie
$\|v-y\|^2$ wokół $u$ dają
$\Phi_y(v)-\Phi_y(u)\ge\|v-u\|^2/(2\lambda)$,
co zapewnia minimalność.

Niech $u=\operatorname{prox}_{\lambda J}(y)$ oraz
$v=\operatorname{prox}_{\lambda J}(z)$.
Monotoniczność subróżniczki daje
$\langle (y-u)-(z-v),u-v\rangle\ge0$.
Stąd $\|u-v\|^2\le\langle y-z,u-v\rangle
\le\|y-z\|\,\|u-v\|$,
co po podzieleniu przez $\|u-v\|$ (gdy $u\ne v$)
daje tezę.
\end{proof}

\subsection{Przykład: słabe zagadnienie brzegowe}

Niech $\lambda$ oznacza miarę Lebesgue'a na $[0,1]$.
Wprowadzamy przestrzeń $H_0^1(0,1)$ przez opis elementarny:
\[
H_0^1(0,1)=\Bigl\{u:[0,1]\to\R:\ 
 u(x)=\int_0^x v(t)\,\mathrm dt\ (0\le x\le1),\ 
 v\in L^2(0,1),\ \int_0^1 v(t)\,\mathrm dt=0\Bigr\}.
\]
Funkcje $u$ są absolutnie ciągłe, ich pochodna $u'=v$
istnieje prawie wszędzie, a warunki brzegowe
$u(0)=u(1)=0$ są spełnione. Przyjmujemy
$\langle u,w\rangle_H=\int_0^1u'(t)w'(t)\,\mathrm dt$.
Odwzorowanie $u\mapsto u'$ jest liniową izometrią
$H_0^1(0,1)$ na domkniętą podprzestrzeń
$\{v\in L^2(0,1):\int_0^1v(t)\,\mathrm dt=0\}$,
przy czym podprzestrzeń ta jest jądrem ciągłego funkcjonału
$v\mapsto\int_0^1v(t)\,\mathrm dt$ (jego norma jest co najwyżej $1$
na mocy nierówności Cauchy'ego--Schwarza).
Surjektywność izometrii wynika z definicji $H_0^1(0,1)$,
a iniektywność z warunku $u(0)=0$ i wzoru całkowego.
Jako domknięta podprzestrzeń przestrzeni Hilberta $L^2(0,1)$
jej obraz jest zupełny, więc $H_0^1(0,1)$ jest przestrzenią Hilberta.

\begin{lemat}[Nierówność Poincarégo na odcinku]
\label{lem:poincare16}
Dla każdego $u\in H_0^1(0,1)$ zachodzi
$\|u\|_{L^2}\le2^{-1/2}\|u'\|_{L^2}$.
\end{lemat}
\begin{proof}
Ponieważ $u(x)=\int_0^xu'(t)\,\mathrm dt$,
nierówność Cauchy'ego--Schwarza daje
$|u(x)|^2\le x\int_0^x|u'(t)|^2\,\mathrm dt
\le x\|u'\|_{L^2}^2$.
Całkując po $x\in[0,1]$, otrzymujemy
$\|u\|_{L^2}^2\le\frac12\|u'\|_{L^2}^2$.
\end{proof}

\begin{twierdzenie}[Zasada minimalizacji dla problemu Dirichleta]
\label{thm:dirichlet16}
Niech $f\in L^2(0,1;\R)$. Na $H=H_0^1(0,1)$
określmy funkcjonał
\[
J(u)=\frac12\int_0^1|u'(t)|^2\,\mathrm dt
-\int_0^1 f(t)u(t)\,\mathrm dt.
\]
Istnieje dokładnie jeden minimizator $u\in H$ funkcjonału $J$.
Jest on dokładnie jedynym elementem $H$ spełniającym
równanie słabe
\[
\int_0^1u'(t)v'(t)\,\mathrm dt
=\int_0^1f(t)v(t)\,\mathrm dt\qquad(v\in H).
\]
Dla $f(t)=1$ rozwiązaniem jest $u(t)=t(1-t)/2$.
\end{twierdzenie}
\begin{proof}
Funkcjonał $J$ jest ciągły i ściśle wypukły na $H$.
Z Lematu~\ref{lem:poincare16} i nierówności
Cauchy'ego--Schwarza wynika
$J(u)\ge\frac12\|u\|_H^2-
2^{-1/2}\|f\|_{L^2}\|u\|_H$, więc jest koercywny
i ograniczony z dołu. Ponieważ jest wypukły
i normowo ciągły, jest słabo półciągły z dołu.
Twierdzenie~\ref{thm:direct-method16} daje minimizator,
a ścisła wypukłość zapewnia jego jednoznaczność.

Dla $h\in H$ rozwinięcie kwadratu daje
\[
J(u+h)-J(u)=\langle u,h\rangle_H-\int_0^1 f(t)h(t)\,\mathrm dt
+\tfrac12\|h\|_H^2.
\]
Część liniowa jest ograniczona przez
$(\|u\|_H+2^{-1/2}\|f\|_{L^2})\|h\|_H$,
a reszta podzielona przez $\|h\|_H$ dąży do zera.
Zatem $J$ jest różniczkowalny Frécheta, a jego pochodna ma postać
$DJ(u)v=\int_0^1u'(t)v'(t)\,\mathrm dt-
\int_0^1 f(t)v(t)\,\mathrm dt$.
Twierdzenie~\ref{thm:convex-derivative16} daje
równoważność $DJ(u)=0$ i globalnej minimalności.
Jeśli $f\equiv1$, funkcja $u(t)=t(1-t)/2$ spełnia
$u(0)=u(1)=0$ oraz $u''(t)=-1$; całkowanie przez części
dla absolutnie ciągłego $v\in H$ daje
$\int_0^1u'v'\,\mathrm dt=\int_0^1v\,\mathrm dt$.
Przez jednoznaczność jest to szukane rozwiązanie.
\end{proof}

\subsection{Wypukłe maksima i twierdzenie Kreina--Milmana}

Twierdzenie Kreina--Milmana i zasada maksimum Bauera orzekają,
że niepusty zwarty wypukły podzbiór przestrzeni lokalnie wypukłej
Hausdorffa jest domkniętą otoczką wypukłą swoich punktów
ekstremalnych, a każda wypukła i półciągła z góry funkcja
rzeczywista na tym zbiorze osiąga maksimum w punkcie ekstremalnym. Ich rola w optymalizacji jest odmienna od metody bezpośredniej:
\emph{zwartość wypukłego zbioru} daje możliwość badania maksimum
funkcjonału wypukłego na punktach ekstremalnych,
natomiast \emph{koercywność i słaba półciągłość z dołu}
służą do znajdowania minimum.

\begin{stwierdzenie}[Maksimum funkcjonału afinicznego]
\label{prop:affine-extreme16}
Niech $E$ będzie lokalnie wypukłą TVS Hausdorffa nad $\K$,
$K\subseteq E$ niepustym, zwartym i wypukłym zbiorem,
a $A:K\to\R$ ciągłą funkcją afiniczną,
czyli $A((1-t)x+ty)=(1-t)A(x)+tA(y)$.
Wtedy istnieje punkt ekstremalny $e\in K$, w którym
$A$ osiąga maksimum, oraz (być może inny) punkt ekstremalny,
w którym osiąga minimum.
\end{stwierdzenie}
\begin{proof}
Funkcje $A$ i $-A$ są wypukłe i ciągłe,
więc zastosowanie zasady maksimum Bauera do obu daje tezę.
\end{proof}

\begin{przyklad}[Miary prawdopodobieństwa i ich punkty ekstremalne]
\label{ex:prob-measures16}
Niech $K$ będzie niepustą zwartą przestrzenią metryczną,
a $\mathcal P(K)$ zbiorem borelowskich miar probabilistycznych na $K$.
Na mocy twierdzenia Riesza--Markowa identyfikujemy
$\mathcal P(K)$ z dodatnimi funkcjonałami na $C(K;\R)$,
przyjmującymi wartość $1$ na funkcji stałej $1$.
W topologii słabej-* $\sigma(C(K)',C(K))$ jest to zbiór
zwarty i wypukły: jako podzbiór kuli jednostkowej
jest domknięty przez warunki
$L(1)=1$ i $L(\varphi)\ge0$ dla każdej $\varphi\ge0$.

Punktami ekstremalnymi $\mathcal P(K)$ są dokładnie
miary Diraca $\delta_x$, $x\in K$.
Dla $\delta_x=t\mu+(1-t)\nu$, $0<t<1$, dodatniość
miar i równość ich mas wymuszają
$\mu(K\setminus\{x\})=\nu(K\setminus\{x\})=0$,
więc $\mu=\nu=\delta_x$.
Jeśli natomiast $\mu$ nie jest miarą Diraca,
istnieje zbiór borelowski $B$ o
$0<\mu(B)<1$: w przeciwnym razie nośnik miary
musiałby być jednopunktowy (dla zwartej metrycznej $K$
wynika to z rozdzielania różnych punktów nośnika
rozłącznymi otoczeniami o dodatniej mierze).
Wówczas
$\mu=\mu(B)\,\mu_B+(1-\mu(B))\,\mu_{K\setminus B}$,
gdzie $\mu_B$ i $\mu_{K\setminus B}$ są znormalizowanymi
ograniczeniami $\mu$, a obie miary są różne.
Twierdzenie Kreina--Milmana stwierdza więc, że
$\mathcal P(K)$ jest słabo-* domkniętą otoczką wypukłą
rodziny $\{\delta_x:x\in K\}$.
W szczególności dla $\varphi\in C(K;\R)$
maksimum funkcjonału $\mu\mapsto\int_K\varphi(x)\,\mathrm d\mu(x)$
można osiągnąć na pewnej mierze Diraca.
\end{przyklad}

\begin{uwaga}[Dalsza perspektywa]
Twierdzenie Kreina--Milmana jest rezultatem o \emph{zwartych
zbiorach wypukłych} w ogólnych przestrzeniach lokalnie wypukłych,
a nie jedynie o kulach jednostkowych przestrzeni Banacha.
Z kolei subróżniczki i operatory proksymalne prowadzą
do teorii operatorów monotonicznych i metod minimalizacji,
co wykracza poza podstawowy zakres obecnego rozdziału.
\end{uwaga}

% Wersja źródłowa: audyt II, jednostka 26
\setcounter{section}{25}
\section{Geometria przestrzeni Banacha: refleksywność, zwartość słaba i wypukłość jednostajna}
\label{sec:geometria26}

\providecommand{\atlasref}[2]{#2}


\noindent\textbf{Wyniki wykorzystywane.} Korzystamy z twierdzenia
Hahna--Banacha dla przestrzeni unormowanych
(Twierdzenie~\ref{thm:hb-normed}), separacji wypukłej
(Wniosek~\ref{cor:point-closed-convex}), podstaw topologii
słabych z rozdziału~8, twierdzenia Banacha--Alaoglu
(Wniosek~\ref{cor:alaoglu-norm}) oraz zasady jednostajnej
ograniczoności (Wniosek~\ref{cor:uniform-boundedness}).
Twierdzenia Goldstine'a, Eberleina--\v Smuliana i Milmana--Pettisa
udowodnimy tutaj.
\par\medskip

W tym rozdziale $X$ oznacza przestrzeń Banacha nad $\K\in\{\R,\C\}$,
$X'$ jej przestrzeń dualną topologiczną z normą operatorową,
a $X''=(X')'$ i $X'''=(X'')'$ kolejne przestrzenie dualne
z tymi samymi konwencjami. Funkcjonały są liniowe nad $\K$,
również w przypadku zespolonym.
Przez $B_X=\{x\in X:\|x\|\le1\}$ rozumiemy domkniętą kulę
jednostkową, natomiast $\sigma(X,X')$ oznacza topologię słabą.
Domknięcia z indeksem $w$ dotyczą topologii słabej,
a z indeksem $w^\ast$ w $X''$ topologii $\sigma(X'',X')$.
Celem jest ustalenie, kiedy ograniczoność daje zwartość słabą,
co wolno wówczas wnioskować o ciągach oraz jak geometria normy
wpływa na refleksywność.

\subsection{Kanoniczne zanurzenie i twierdzenie Goldstine'a}

\noindent\emph{Powiązania tematyczne w Bibliotece matematycznej:} \atlasref{LA-07::thm:lin7-bidualnosc}{LA-07: kanoniczne odwzorowanie w podwójną przestrzeń sprzężoną w wymiarze skończonym}.


\begin{definicja}[Refleksywność]
\label{def:refleksywnosc26}
Przestrzeń Banacha $X$ nazywamy \emph{refleksywną}, jeżeli
kanoniczne odwzorowanie liniowe
$J_X:X\to X''$, określone wzorem
$(J_Xx)(f)=f(x)$ dla $x\in X$ i $f\in X'$, jest surjektywne.
\end{definicja}

\begin{stwierdzenie}[Kanoniczne zanurzenie]
\label{prop:canonical26}
Dla każdej przestrzeni unormowanej $X$ odwzorowanie
$J_X:X\to X''$ jest liniową izometrią. Jeżeli $X$ jest zupełna,
to $J_X(X)$ jest normowo domkniętą podprzestrzenią $X''$.
Ponadto $J_X$ jest homeomorfizmem przestrzeni
$(X,\sigma(X,X'))$ na swój obraz z topologią podprzestrzeni
$(X'',\sigma(X'',X'))$.
\end{stwierdzenie}
\begin{proof}
Dla ustalonego $x\in X$ odwzorowanie $f\mapsto f(x)$
jest liniowe oraz $|f(x)|\le\|f\|\|x\|$, więc należy
do $X''$. Liniowość $J_X$ wynika z liniowości każdego $f$.
Dla $x\in X$ mamy
$\|J_Xx\|=\sup_{\|f\|\le1}|f(x)|\le\|x\|$.
Dla $x\ne0$ Wniosek~\ref{cor:norming} zapewnia funkcjonał
$f\in X'$ o normie $1$ i wartości $f(x)=\|x\|$.
Daje to odwrotną nierówność; dla $x=0$ mamy równość.
Izometria jest różnowartościowa. Obraz przestrzeni zupełnej
przez izometrię jest zupełny, a zupełna podprzestrzeń
przestrzeni metrycznej jest domknięta.

Topologia słaba na $X$ jest wyznaczona przez seminormy
$x\mapsto|f(x)|$ dla $f\in X'$, natomiast topologia
słaba-* na $X''$ przez $\Phi\mapsto|\Phi(f)|$ dla
tych samych $f$. Ponieważ $|(J_Xx)(f)|=|f(x)|$,
przeciwobrazy tych seminorm są dokładnie rodziną
wyznaczającą topologię słabą. Stąd $J_X$ i jego
odwrotność na obrazie są ciągłe w podanych topologiach.
\end{proof}

\begin{twierdzenie}[Goldstine]
\label{thm:goldstine26}
Niech $X$ będzie przestrzenią unormowaną nad $\K$.
Wówczas $J_X(B_X)$ jest gęsty w $B_{X''}$ względem topologii
$\sigma(X'',X')$.
\end{twierdzenie}
\begin{proof}
Z izometryczności $J_X$ mamy $J_X(B_X)\subseteq B_{X''}$.
Kula $B_{X''}$ jest słabo-* domknięta, ponieważ jest
przecięciem zbiorów $\{\Psi:|\Psi(f)|\le1\}$ po
$f\in B_{X'}$. Przypuśćmy, że
$\Phi\in B_{X''}\setminus\overline{J_X(B_X)}^{\,w^\ast}$.
Z bazy otoczeń słabych-* (Twierdzenie~\ref{thm:slaba-baza})
istnieją $m\ge1$, $f_1,\ldots,f_m\in X'$ i $\eta>0$,
dla których żaden $x\in B_X$ nie spełnia wszystkich
nierówności $|f_j(x)-\Phi(f_j)|<\eta$.
Wektor $(\Phi(f_1),\ldots,\Phi(f_m))$ leży więc poza
domknięciem w $\K^m$ wypukłego zbioru
$\{(f_1(x),\ldots,f_m(x)):x\in B_X\}$.
Stosując Wniosek~\ref{cor:point-closed-convex} w
skończonym wymiarze, oddzielamy ten wektor od domknięcia.
Nad $\C$ rzeczywisty funkcjonał liniowy na $\C^m$
jest częścią rzeczywistą funkcjonału zespolonego
(Lemat~\ref{lem:realification}); nad $\R$ przyjmujemy
$\operatorname{Re}t=t$. Otrzymujemy skalary
$c_1,\ldots,c_m\in\K$ i liczbę rzeczywistą $a$ spełniające
\[
\sup_{x\in B_X}\operatorname{Re}\Big(\sum_{j=1}^m c_jf_j(x)\Big)
<a<\operatorname{Re}\Big(\sum_{j=1}^m c_j\Phi(f_j)\Big).
\]
Dla $f=\sum_jc_jf_j$ mamy
$\sup_{x\in B_X}\operatorname{Re}f(x)=\|f\|$:
nierówność $\le$ wynika z definicji normy, a dla
każdego $x$ z $f(x)\ne0$ można pomnożyć $x$ przez
skalar modułu $1$, aby otrzymać wartość $|f(x)|$.
W przypadku $f(x)=0$ nie trzeba zmieniać $x$.
Z liniowości $\Phi$ prawa strona nierówności to
$\operatorname{Re}\Phi(f)\le|\Phi(f)|\le\|\Phi\|\|f\|
\le\|f\|$. Otrzymana sprzeczność kończy dowód.
\end{proof}

\begin{twierdzenie}[Refleksywność a zwartość słaba kuli]
\label{thm:reflexive-ball26}
Niech $X$ będzie przestrzenią Banacha. Następujące warunki
są równoważne:
\begin{enumerate}[label=\textup{(\roman*)}]
\item $X$ jest refleksywna;
\item $B_X$ jest zwarta w topologii $\sigma(X,X')$;
\item każdy niepusty, ograniczony, normowo domknięty i wypukły
podzbiór $X$ jest słabo zwarty.
\end{enumerate}
\end{twierdzenie}
\begin{proof}
Jeżeli $J_X$ jest surjektywne, to z izometryczności i
Stwierdzenia~\ref{prop:canonical26} jego ograniczenie
do $B_X$ jest homeomorfizmem na $B_{X''}$ w odpowiednich
topologiach słabej i słabej-*. Stosujemy Wniosek~
\ref{cor:alaoglu-norm} do przestrzeni unormowanej $X'$:
kula jej przestrzeni dualnej $X''$ jest słabo-* zwarta.
Zatem (i) pociąga (ii).
Jeżeli zachodzi (ii), to dla każdego ograniczonego normowo,
zamkniętego i wypukłego $C\subseteq X$ istnieje $R>0$ z
$C\subseteq RB_X$. Z Wniosku~\ref{cor:weak-convex-closure} wynika, że
normowo domknięty zbiór wypukły jest słabo domknięty.
Zatem $C$ jest domkniętym podzbiorem słabo zwartej kuli $RB_X$,
co daje (iii). Przypadek $C=\varnothing$ nie wpływa na argument.
Warunek (iii) zastosowany do $B_X$ daje (ii).
Wreszcie jeśli $B_X$ jest słabo zwarta, to $J_X(B_X)$ jest
słabo-* zwarta, a zatem słabo-* domknięta w $X''$,
gdyż funkcjonały ewaluacji na $X'$ rozdzielają punkty
$X''$. Twierdzenie~\ref{thm:goldstine26} daje
$J_X(B_X)=B_{X''}$. Dla $\Phi\ne0$ wybieramy
$x\in B_X$ takie, że $J_Xx=\Phi/\|\Phi\|$; wtedy
$J_X(\|\Phi\|x)=\Phi$. Element zerowy jest obrazem zera.
Otrzymujemy surjektywność $J_X$, czyli (i).
\end{proof}

\begin{wniosek}[Podprzestrzenie i przestrzeń dualna]
\label{cor:reflexivity-stability26}
Jeżeli $X$ jest refleksywną przestrzenią Banacha, to każda
jej domknięta podprzestrzeń liniowa jest refleksywna.
Ponadto $X$ jest refleksywna wtedy i tylko wtedy, gdy $X'$
jest refleksywna.
\end{wniosek}
\begin{proof}
Niech $Y\subseteq X$ będzie domkniętą podprzestrzenią.
Jako zbiór domknięty i wypukły $Y$ jest słabo domknięta.
Zatem $B_Y=B_X\cap Y$ jest słabo zwarta w $X$.
Ograniczenia funkcjonałów z $X'$ należą do $Y'$,
a z Twierdzenia~\ref{thm:hb-normed} każdy $g\in Y'$
przedłuża się do $X'$ z zachowaniem normy.
Obie rodziny seminorm na $Y$ są więc identyczne,
co daje równość topologii słabej $Y$ i topologii
indukowanej z $X$. Ponieważ $Y$ jest normowo domknięta
w przestrzeni Banacha, sama jest przestrzenią Banacha. Twierdzenie~\ref{thm:reflexive-ball26}
daje refleksywność $Y$.

Jeżeli $X$ jest refleksywna, to $J_X:X\to X''$ jest
izomorfizmem izometrycznym. Dla $G\in X'''$ określamy
$f(x)=G(J_Xx)$. Liniowość oraz oszacowanie
$|f(x)|\le\|G\|\|x\|$ dają $f\in X'$.
Dla $\Phi=J_Xx\in X''$
mamy $(J_{X'}f)(\Phi)=\Phi(f)=f(x)=G(\Phi)$.
Surjektywność $J_X$ daje $J_{X'}f=G$, więc $X'$ jest refleksywna.
Odwrotnie, gdy $X'$ jest refleksywna i $J_X(X)\ne X''$,
obraz $J_X(X)$ jest normowo domknięty z zupełności $X$.
Twierdzenie~\ref{thm:point-closed-subspace}, zastosowane
do tej właściwej podprzestrzeni przestrzeni unormowanej $X''$,
daje niezerowy $G\in X'''$ zerujący się na $J_X(X)$.
Z refleksywności $X'$ wynika $G=J_{X'}f$ dla pewnego $f\in X'$.
Wtedy $0=G(J_Xx)=f(x)$ dla wszystkich $x\in X$, więc
$f=0$ i $G=0$, sprzeczność.
\end{proof}

\subsection{Twierdzenie Eberleina--\v Smuliana}

W dowolnej przestrzeni topologicznej zwartości nie można
charakteryzować wyłącznie ciągami. Dla słabej topologii
przestrzeni Banacha zachodzi jednak ważny wyjątek.

\begin{lemat}[Skończona rodzina normująca]
\label{lem:finite-norming26}
Niech $X$ będzie przestrzenią unormowaną,
a $M\subseteq X''$ skończenie wymiarową podprzestrzenią.
Istnieje skończony niepusty zbiór $F\subseteq B_{X'}$ taki, że
$\|\Psi\|\le 2\max_{f\in F}|\Psi(f)|$ dla wszystkich $\Psi\in M$.
\end{lemat}
\begin{proof}
Dla $M=\{0\}$ wystarczy przyjąć $F=\{0\}$.
W pozostałym przypadku dla każdego $\Psi\in M$ o normie $1$
wybieramy $f_\Psi\in B_{X'}$ z
$|\Psi(f_\Psi)|>3/4$, co jest możliwe z definicji
normy operatorowej $\Psi$. Sfera jednostkowa
skończenie wymiarowej przestrzeni $M$ jest zwarta
w topologii normowej.
Pokrywamy ją kulami $\{\Theta\in M:\|\Theta-\Psi\|<1/4\}$
i wybieramy skończone podpokrycie o środkach
$\Psi_1,\ldots,\Psi_r$. Dla $\Theta$ na sferze i odpowiedniego
$j$ mamy
\[
 |\Theta(f_{\Psi_j})|
 \ge|\Psi_j(f_{\Psi_j})|-\|\Theta-\Psi_j\|\|f_{\Psi_j}\|
 >3/4-1/4=1/2.
\]
Kładziemy $F=\{f_{\Psi_1},\ldots,f_{\Psi_r}\}$.
Dla niezerowego $\Theta\in M$ stosujemy nierówność
do $\Theta/\|\Theta\|$, a dla zera teza jest oczywista.
\end{proof}

\begin{twierdzenie}[Eberlein--\v Smulian]
\label{thm:eberlein-smulian26}
Niech $X$ będzie przestrzenią Banacha, a $A\subseteq X$
dowolnym podzbiorem. Następujące warunki są równoważne:
\begin{enumerate}[label=\textup{(\roman*)}]
\item domknięcie $\overline A^{\,w}$ jest słabo zwarte;
\item z każdego ciągu punktów $A$ można wybrać podciąg
zbieżny słabo do pewnego punktu $X$;
\item każdy ciąg w $A$ ma słaby punkt skupienia w $X$
(w każdym jego słabym otoczeniu leżą wyrazy ciągu o dowolnie dużych indeksach).
\end{enumerate}
\end{twierdzenie}
\begin{proof}
Gdy $A=\varnothing$, wszystkie warunki zachodzą.
Zakładamy więc $A\ne\varnothing$; przypadek $X=\{0\}$
jest natychmiastowy.

\emph{Krok pierwszy: (i) $\Rightarrow$ (ii).}
Niech $(x_n)\subseteq A$ i
$Y=\overline{\lin\{x_n:n\ge1\}}^{\,\|\cdot\|}$.
Jest to ośrodkowa podprzestrzeń: skończone kombinacje
liniowe $x_n$ o współczynnikach wymiernych (nad $\C$
o wymiernych częściach rzeczywistych i urojonych)
tworzą przeliczalny zbiór gęsty w $Y$.
Podprzestrzeń $Y$ jest normowo domknięta i wypukła,
więc jest słabo domknięta
(Wniosek~\ref{cor:weak-convex-closure}). Zbiór
$K=\overline A^{\,w}\cap Y$ jest słabo zwarty i zawiera $(x_n)$.
Wybieramy ciąg gęsty $(y_j)$ w $Y$ i funkcjonały
$f_j\in B_{X'}$ takie, że $f_j(y_j)=\|y_j\|$
(Wniosek~\ref{cor:norming}; dla $y_j=0$ bierzemy $f_j=0$).
Taka rodzina rozdziela punkty $Y$: dla $y\ne0$ wystarczy
wybrać $y_j$ w odległości mniejszej niż $\|y\|/3$ od $y$.
Wówczas $|f_j(y)|\ge\|y_j\|-\|y-y_j\|
\ge\|y\|-2\|y-y_j\|>\|y\|/3$.
Na $K$ funkcja
\[
d(u,v)=\sum_{j=1}^\infty 2^{-j}\min\{1,|f_j(u-v)|\}
\]
jest metryką: szereg jest zbieżny, symetria wynika
z modułu, a rozdzielanie punktów z powyższej własności
rodziny $(f_j)$. Nierówność trójkąta wynika z
podaddytywności $t\mapsto\min\{1,t\}$ dla $t\ge0$
oraz nierówności trójkąta dla modułu.
Dla ustalonego $v\in K$ każda suma częściowa
$d_N(\cdot,v)$ jest słabo ciągła i
$0\le d(u,v)-d_N(u,v)\le2^{-N}$ dla wszystkich $u,v\in K$.
Stąd $d(\cdot,v)$ jest słabo ciągła jako jednostajna
granica funkcji ciągłych; kule metryczne są słabo
otwarte w $K$. Topologia metryczna jest więc słabsza
od topologii słabej i jest Hausdorffa. Ponieważ $K$ jest słabo zwarty,
identyczność z $K$ słabo zwartą do $K$ metryczną jest
homeomorfizmem: ciągła bijekcja z przestrzeni zwartej
do Hausdorffa ma ciągłą odwrotność. W konsekwencji $K$
jest zwartą przestrzenią metryczną, więc $(x_n)$ ma w $K$
podciąg zbieżny słabo.

\emph{Krok drugi: (ii) $\Rightarrow$ (iii).}
Granica podciągu słabo zbieżnego jest słabym punktem skupienia
całego ciągu, gdyż indeksy podciągu rosną do nieskończoności.

\emph{Krok trzeci: (iii) $\Rightarrow$ (i).}
Najpierw $A$ jest ograniczony w normie. W przeciwnym razie
rodzina $\{J_Xx:x\in A\}\subseteq X''$, rozpatrywana jako
rodzina funkcjonałów na przestrzeni Banacha $X'$, byłaby
nieograniczona jednostajnie w normie, gdyż
$\|J_Xx\|=\|x\|$. Przestrzeń $X'$ jest przestrzenią Banacha nawet
dla unormowanej przestrzeni $X$. Z zasady jednostajnej
ograniczoności (Wniosek~\ref{cor:uniform-boundedness})
istniałby $f\in X'$ z
$\sup_{x\in A}|f(x)|=\infty$. Wybierając $x_n\in A$ z
$|f(x_n)|\ge n$, dostalibyśmy ciąg bez słabego punktu
skupienia: dla każdego $x\in X$ słabe otoczenie
$\{y:|f(y)-f(x)|<1\}$ nie zawiera żadnego $x_n$
dla $n>|f(x)|+1$.

Niech $\Phi\in\overline{J_X(A)}^{\,w^\ast}\subseteq X''$.
Skonstruujemy ciąg $(x_n)\subseteq A$ oraz skończone
zbiory $F_n\subseteq B_{X'}$. Wybieramy $x_1\in A$ i zbiór
$F_1$ z Lematu~\ref{lem:finite-norming26}, normujący z dokładnością
do czynnika $2$ przestrzeń
$M_1=\lin\{\Phi,J_Xx_1\}\subseteq X''$.
Gdy $x_1,\ldots,x_n$ i $F_1,\ldots,F_n$ są ustalone,
wybieramy $x_{n+1}\in A$ tak, by
$|\Phi(f)-f(x_{n+1})|<1/(n+1)$ dla każdego
$f\in F_1\cup\cdots\cup F_n$; jest to możliwe, ponieważ
$\Phi$ należy do słabo-* domknięcia zbioru $J_X(A)$,
a skończona rodzina warunków wyznacza
podstawowe otoczenie słabe-* punktu $\Phi$.
Następnie wybieramy $F_{n+1}$ normujący z dokładnością
do czynnika $2$ podprzestrzeń
$M_{n+1}=\lin\{\Phi,J_Xx_1,\ldots,J_Xx_{n+1}\}$.

Na mocy (iii) istnieje słaby punkt skupienia $x\in X$
ciągu $(x_n)$. Dla każdego $f\in\bigcup_nF_n$ mamy
$f(x_n)\to\Phi(f)$: jeśli $f\in F_j$, to
$|f(x_n)-\Phi(f)|<1/n$ dla $n>j$.
Punkt $f(x)$ jest punktem skupienia ciągu $(f(x_n))$,
bo przeciwobrazy otoczeń przez $f$ są słabo otwarte.
Jednoznaczność granicy skalarnej daje $f(x)=\Phi(f)$.
Ponadto $x$ należy do słabego domknięcia zbioru
$\{x_n\}$, a stąd do normowo domkniętej podprzestrzeni
$Y=\overline{\lin\{x_n:n\ge1\}}^{\,\|\cdot\|}$,
która jest słabo domknięta na mocy twierdzenia Hahna--Banacha.
W szczególności $J_Xx$ należy do domknięcia normowego
przestrzeni $\lin\{J_Xx_n:n\ge1\}$, zawartej w
$\bigcup_nM_n$. Przestrzenie $M_n$ są rosnące,
więc ich suma mnogościowa jest podprzestrzenią liniową.
Ponieważ $\Phi\in M_1$, również
$\Psi=\Phi-J_Xx$ należy do normowego domknięcia
$\bigcup_nM_n$. Dla dowolnego $\eta>0$ wybieramy
$n$ i $G\in M_n$ z $\|G-\Psi\|<\eta$. Z konstrukcji dla wszystkich
$f\in F_n$ zachodzi $\Psi(f)=0$, a więc
$G(f)=(G-\Psi)(f)$.
Ponieważ $G\in M_n$ oraz $F_n\subseteq B_{X'}$,
z własności rodziny normującej mamy
$\|G\|\le2\max_{f\in F_n}|G(f)|\le2\|G-\Psi\|$.
Stąd $\|\Psi\|\le3\|\Psi-G\|<3\eta$.
Dowolność $\eta>0$ daje $\Psi=0$, a zatem $\Phi=J_Xx$.

Pokazaliśmy, że słabo-* domknięcie $J_X(A)$ zawiera się
w $J_X(X)$. Z pierwszej części kroku trzeciego zbiór $A$
jest normowo ograniczony, więc istnieje $R>0$ takie,
że $J_X(A)\subseteq RB_{X''}$. Ta kula jest słabo-*
domknięta i zwarta (Wniosek~\ref{cor:alaoglu-norm}).
Zatem $C=\overline{J_X(A)}^{\,w^\ast}$ jest słabo-*
zwarty i, jak wykazaliśmy, $C\subseteq J_X(X)$.
Z własności domknięcia w podprzestrzeni wynika
$\overline{J_X(A)}^{\,w^\ast|_{J_X(X)}}=C\cap J_X(X)=C$.
Ponieważ $J_X$ jest homeomorfizmem na obraz
(Stwierdzenie~\ref{prop:canonical26}), otrzymujemy
$J_X(\overline A^{\,w})=C$.
Przenosząc zwartość przez odwrotność tego
homeomorfizmu, dostajemy słabą zwartość
$\overline A^{\,w}$.
\end{proof}

\begin{wniosek}[Ograniczone ciągi w przestrzeni refleksywnej]
\label{cor:bounded-subsequence26}
Jeżeli $X$ jest refleksywną przestrzenią Banacha,
to każdy ograniczony ciąg w $X$ ma podciąg słabo zbieżny.
\end{wniosek}
\begin{proof}
Wszystkie wyrazy ciągu należą do pewnej słabo zwartej kuli
$RB_X$ na mocy Twierdzenia~\ref{thm:reflexive-ball26}.
Stosujemy Twierdzenie~\ref{thm:eberlein-smulian26}.
\end{proof}

\subsection{Własność Schura i geometria przestrzeni \texorpdfstring{$\ell^1$}{ell1}}

\begin{definicja}[Własność Schura]
\label{def:schur26}
Przestrzeń Banacha $X$ ma \emph{własność Schura}, jeżeli dla
każdego ciągu $(x_n)\subseteq X$ i $x\in X$ ze zbieżności
$x_n\to x$ słabo wynika $\|x_n-x\|\to0$.
\end{definicja}

\begin{twierdzenie}[Własność Schura w $\ell^1$]
\label{thm:schur-l1-26}
Przestrzeń Banacha $\ell^1(\N,\K)$ ma własność Schura.
\end{twierdzenie}
\begin{proof}
Jeżeli $x_n\to x$ słabo, to $x_n-x\to0$ słabo,
więc wystarczy rozważyć ciąg $(x_n)$ słabo zbieżny
do zera. Dla każdej współrzędnej $j$ odwzorowanie
$x\mapsto x(j)$ jest liniowym funkcjonałem ciągłym,
gdyż $|x(j)|\le\|x\|_1$. Wobec tego $x_n(j)\to0$.
Przypuśćmy, że normy nie dążą do zera. Po przejściu do
podciągu możemy założyć, że $\|x_n\|_1\ge\varepsilon>0$
dla wszystkich $n$. Rekurencyjnie wybieramy podciąg
$y_k=x_{n_k}$ oraz liczby całkowite
$0=m_0<m_1<m_2<\cdots$ takie, że
\[
\sum_{j\le m_{k-1}}|y_k(j)|<\varepsilon/8,
\qquad
\sum_{j>m_k}|y_k(j)|<\varepsilon/8.
\]
Rozpoczynamy od $m_0=n_0=0$. Mając $m_{k-1}$ i
$n_{k-1}$, wybieramy $n_k>n_{k-1}$ tak duże, aby
$\sum_{j\le m_{k-1}}|x_{n_k}(j)|<\varepsilon/8$:
jest to skończona suma wyrazów dążących do zera
(dla $m_{k-1}=0$ suma jest zerowa).
Ustaliwszy $y_k=x_{n_k}$, wybieramy
$m_k>m_{k-1}$, korzystając z
$\sum_{j\ge1}|y_k(j)|<\infty$.
Na bloku $I_k=\{m_{k-1}+1,\ldots,m_k\}$ otrzymujemy
\[
 \sum_{j\in I_k}|y_k(j)|
 =\|y_k\|_1-\sum_{j\notin I_k}|y_k(j)|
 >3\varepsilon/4,\qquad
 \sum_{j\notin I_k}|y_k(j)|<\varepsilon/4.
\]
Bloki są parami rozłączne, a ponieważ $m_k\ge k$,
ich suma mnogościowa jest równa $\N$.
Definiujemy $a=(a_j)\in\ell^\infty$ tak, że dla $j\in I_k$
wybieramy $a_j=\overline{y_k(j)}/|y_k(j)|$, gdy $y_k(j)\ne0$,
i $a_j=0$ w przeciwnym wypadku.
Nad $\R$ sprzężenie jest identycznością.
Mamy $|a_j|\le1$ dla wszystkich $j$, a na $I_k$
zachodzi $a_jy_k(j)=|y_k(j)|$, także gdy $y_k(j)=0$.
Definiujemy jeden funkcjonał (niezależny od $k$)
$f(x)=\sum_{j\ge1}a_jx(j)$. Szereg jest bezwzględnie
zbieżny, gdyż $\sum_j|a_jx(j)|\le\|x\|_1$.
Liniowość wynika z liniowości sum częściowych,
a oszacowanie $|f(x)|\le\|x\|_1$ daje ciągłość,
zgodnie z Twierdzeniem~\ref{thm:dual-sequences12}.
Dla każdego $k$ z rozbicia bezwzględnie zbieżnego
szeregu i nierówności $\operatorname{Re}z\ge-|z|$
wynika
\[
\operatorname{Re}f(y_k)\ge
\sum_{j\in I_k}|y_k(j)|-\sum_{j\notin I_k}|y_k(j)|
>\varepsilon/2.
\]
Tymczasem $(y_k)$ jest podciągiem ciągu słabo
zbieżnego do zera, więc $f(y_k)\to0$ dla tego
ustalonego $f\in(\ell^1)'$. Otrzymujemy sprzeczność.
\end{proof}

\begin{wniosek}[Nierefleksywność przestrzeni $\ell^1$]
\label{cor:l1-nonreflexive26}
Nieskończenie wymiarowa przestrzeń Banacha $\ell^1(\N,\K)$
nie jest refleksywna.
\end{wniosek}
\begin{proof}
Niech $e_n(j)=1$ dla $j=n$ oraz $e_n(j)=0$
dla $j\ne n$. Mamy $\|e_n\|_1=1$ i
$\|e_n-e_m\|_1=2$ dla $n\ne m$, więc ciąg jest
ograniczony, lecz żaden jego podciąg nie jest
ciągiem Cauchy'ego w normie.
Gdyby $\ell^1$ była refleksywna,
Wniosek~\ref{cor:bounded-subsequence26} zapewniałby podciąg
słabo zbieżny do pewnego $x\in\ell^1$.
Twierdzenie~\ref{thm:schur-l1-26} dawałoby jego
zbieżność normową, a więc własność Cauchy'ego.
Sprzeczność dowodzi nierefleksywności $\ell^1$.
\end{proof}

\begin{przyklad}[Słabo zbieżny ciąg bez zbieżności normowej]
\label{ex:en-hilbert26}
Niech $H$ będzie nieskończenie wymiarową przestrzenią Hilberta.
Istnieje w niej nieskończony ciąg ortonormalny $(e_n)$:
po wyborze $e_1,\ldots,e_n$ ich liniowa powłoka jest
właściwą domkniętą podprzestrzenią, więc z rozkładu
ortogonalnego (Twierdzenie~\ref{thm:orth-decomposition13})
można wybrać wektor jednostkowy $e_{n+1}$
w jej dopełnieniu ortogonalnym. Nie zakładamy
ośrodkowości $H$ ani zupełności tego układu ortonormalnego.
Dla każdego $x\in H$ nierówność Bessela
(Lemat~\ref{lem:bessel13}) daje
$\sum_{n\ge1}|\langle e_n,x\rangle|^2\le\|x\|^2$,
więc $\langle e_n,x\rangle\to0$.
Z twierdzenia Riesza
(Twierdzenie~\ref{thm:riesz-hilbert13}) każdy $f\in H'$
ma postać $f(y)=\langle y,x\rangle$ dla pewnego $x\in H$.
Zatem $f(e_n)\to0$ dla każdego $f\in H'$, czyli
$e_n\to0$ słabo w $H$. Jednak $\|e_n\|=1$
dla wszystkich $n$.
Dlatego nieskończenie wymiarowa przestrzeń Hilberta
nie ma własności Schura.
\end{przyklad}

\subsection{Jednostajna wypukłość i twierdzenie Milmana--Pettisa}

\noindent\emph{Powiązania tematyczne w Bibliotece matematycznej:} \atlasref{ME-15::thm:uniform-convexity-15}{ME-15: jednostajna wypukłość przestrzeni funkcji całkowalnych}; \atlasref{ME-15::thm:reflexive-lp-15}{ME-15: refleksywność przestrzeni funkcji całkowalnych}.


\begin{definicja}[Ścisła i jednostajna wypukłość]
\label{def:uniform-convex26}
Niech $X$ będzie przestrzenią unormowaną.
Mówimy, że norma jest \emph{ściśle wypukła}, jeżeli
dla różnych $x,y\in X$ o normach $1$ zachodzi
$\|(x+y)/2\|<1$.
Przestrzeń $X$ jest \emph{jednostajnie wypukła}, jeżeli
dla każdego $\varepsilon\in(0,2]$ istnieje
$\delta(\varepsilon)>0$ takie, że dla wszystkich
$x,y\in X$ o normach $1$ i $\|x-y\|\ge\varepsilon$
zachodzi $\|(x+y)/2\|\le1-\delta(\varepsilon)$.
Przez $\delta_X(\varepsilon)$ oznaczamy dowolną
dodatnią stałą spełniającą ten warunek, nie większą
niż $1$. Jednostajna wypukłość implikuje ścisłą
wypukłość: dla różnych wektorów jednostkowych
wystarczy przyjąć $\varepsilon=\|x-y\|>0$.
\end{definicja}

\begin{przyklad}[Moduł wypukłości przestrzeni Hilberta]
\label{ex:modulus-hilbert26}
Niech $H$ będzie rzeczywistą lub zespoloną przestrzenią
Hilberta. Dla $x,y\in H$ o normach $1$ tożsamość
równoległoboku (Stwierdzenie~\ref{prop:parallelogram13}) daje
$\|x+y\|^2+\|x-y\|^2=4$, czyli
$\|(x+y)/2\|^2=1-\|x-y\|^2/4$.
Jeżeli $\|x-y\|\ge\varepsilon$, to
$\|(x+y)/2\|\le\sqrt{1-\varepsilon^2/4}$.
Zatem $H$ jest jednostajnie wypukła, a można przyjąć
$\delta_H(\varepsilon)=1-\sqrt{1-\varepsilon^2/4}>0$
dla $0<\varepsilon\le2$.
Jeżeli wymiar rzeczywisty $H$ wynosi co najmniej $2$,
jest to optymalna stała: dla dwóch wektorów
ortonormalnych względem rzeczywistego iloczynu skalarnego
$u,v$ bierzemy
$x=\sqrt{1-\varepsilon^2/4}\,u+(\varepsilon/2)v$
i $y=\sqrt{1-\varepsilon^2/4}\,u-(\varepsilon/2)v$.
Oba są jednostkowe i $\|x-y\|=\varepsilon$.
W przypadku zespolonym używamy rzeczywistej części
iloczynu skalarnego, która wyznacza tę samą normę.
Natomiast przestrzeń $\ell^1$ nie jest ściśle wypukła:
$\|e_1\|_1=\|e_2\|_1=\|(e_1+e_2)/2\|_1=1$.
\end{przyklad}

\begin{twierdzenie}[Milman--Pettis]
\label{thm:milman-pettis26}
Każda jednostajnie wypukła przestrzeń Banacha jest refleksywna.
\end{twierdzenie}
\begin{proof}
Wystarczy wykazać, że każdy $\Phi\in X''$ o normie $1$
należy do $J_X(X)$: element zerowy jest obrazem zera,
a dla $\Phi\ne0$ stosujemy wynik do $\Phi/\|\Phi\|$
i mnożymy przeciwobraz przez $\|\Phi\|$.
Obejmuje to również przypadek $X=\{0\}$.
Ustalmy więc $\Phi$ o normie $1$.
Z Twierdzenia~\ref{thm:goldstine26} i charakteryzacji
domknięcia za pomocą sieci istnieje sieć
$(x_\alpha)\subseteq B_X$ taka, że
$J_Xx_\alpha\to\Phi$ słabo-*.
Ustalmy $\varepsilon\in(0,2]$ i wybierzmy
$\delta>0$ mniejsze od $\min\{\delta_X(\varepsilon/2),\varepsilon/4\}$.
Z definicji normy $\Phi$ wybieramy $g\in B_{X'}$
tak, że $|\Phi(g)|>1-\delta/2>0$.
Dzieląc przez $\|g\|$ i mnożąc przez skalar modułu $1$,
otrzymujemy $f\in X'$ o normie $1$ spełniający
$\operatorname{Re}\Phi(f)>1-\delta/2$.
Zbieżność słaba-* daje indeks $\alpha_0$ taki, że
$|f(x_\alpha)-\Phi(f)|<\delta/2$ dla wszystkich
$\alpha\ge\alpha_0$. Wówczas
$\operatorname{Re}f(x_\alpha)>1-\delta>0$.
Ponieważ $\|f\|=1$ i $x_\alpha\in B_X$, mamy
$1-\delta<\|x_\alpha\|\le1$ i dla tych indeksów
możemy położyć $u_\alpha=x_\alpha/\|x_\alpha\|$.
Wektory $u_\alpha$ mają normę $1$, a dzielenie
dodatniej części rzeczywistej przez liczbę z $(0,1]$
daje $\operatorname{Re}f(u_\alpha)>1-\delta$.
Dla $\alpha,\beta\ge\alpha_0$ otrzymujemy
\[
 \left\|\frac{u_\alpha+u_\beta}{2}\right\|
 \ge\operatorname{Re}f\left(\frac{u_\alpha+u_\beta}{2}\right)
 >1-\delta.
\]
Dzięki wybranemu $\delta$ jednostajna wypukłość daje
$\|u_\alpha-u_\beta\|<\varepsilon/2$.
Ponadto $\|u_\alpha-x_\alpha\|=1-\|x_\alpha\|<\delta$,
zatem $\|x_\alpha-x_\beta\|<\varepsilon/2+2\delta<\varepsilon$.
Stąd $(x_\alpha)$ jest siecią Cauchy'ego w normie.
Zupełność $X$ daje $x_\alpha\to x$ w normie:
w przestrzeni metrycznej zupełność dla ciągów
implikuje zupełność dla sieci Cauchy'ego.
Istotnie, dla każdego $n$ wybieramy indeks,
od którego odległości dwóch wyrazów są mniejsze
niż $2^{-n}$; korzystając z kierunkowości, możemy
wybrać rosnące indeksy $\alpha_n$ o tej własności.
Ciąg $(x_{\alpha_n})$ jest Cauchy'ego i zbiega
do pewnego $x$. Dla $\alpha\ge\alpha_n$
mamy $\|x_\alpha-x_{\alpha_n}\|<2^{-n}$,
więc cała sieć zbiega do $x$.
Izometryczność $J_X$ daje zbieżność normową,
a zatem słabą-* sieci $J_Xx_\alpha$ do $J_Xx$.
Z jednoznaczności granicy w topologii słabej-*
Hausdorffa wynika $\Phi=J_Xx$.
Wraz ze skalowaniem z początku dowodu daje to
surjektywność $J_X$.
\end{proof}

\begin{twierdzenie}[Własność Radona--Riesza]
\label{thm:radon-riesz26}
Niech $X$ będzie jednostajnie wypukłą przestrzenią unormowaną,
a $(x_n)\subseteq X$ i $x\in X$. Jeżeli $x_n\to x$
słabo oraz $\|x_n\|\to\|x\|$, to $\|x_n-x\|\to0$.
\end{twierdzenie}
\begin{proof}
Gdy $x=0$, teza wynika wprost ze zbieżności norm.
W przeciwnym razie, ponieważ $\|x_n\|\to\|x\|>0$,
dla dostatecznie dużych $n$ możemy położyć
$u_n=x_n/\|x_n\|$ i $u=x/\|x\|$.
Mamy $\|u_n\|=\|u\|=1$ i dla każdego $f\in X'$
zachodzi $f(u_n)=f(x_n)/\|x_n\|\to f(x)/\|x\|=f(u)$.
Zatem $u_n\to u$ słabo.
Gdyby $\|u_n-u\|\not\to0$, istniałyby
$\varepsilon\in(0,2]$ i podciąg z $\|u_{n_k}-u\|\ge\varepsilon$.
Z jednostajnej wypukłości wynika
$\|(u_{n_k}+u)/2\|\le1-\delta_X(\varepsilon)$.
Tymczasem $(u_{n_k}+u)/2\to u$ słabo, a norma jest
słabo półciągła z dołu, ponieważ
$\|z\|=\sup_{f\in B_{X'}}|f(z)|$
(Wniosek~\ref{cor:norming}) jest supremum
funkcji słabo ciągłych. Otrzymalibyśmy
\[
 1=\|u\|\le\liminf_k\|(u_{n_k}+u)/2\|
 \le1-\delta_X(\varepsilon),
\]
sprzeczność. Stąd $\|u_n-u\|\to0$.
Ostatecznie
$\|x_n-x\|\le\|x_n\|\|u_n-u\|
+|\|x_n\|-\|x\||\to0$, gdyż ciąg norm jest ograniczony.
Zupełność $X$ nie jest potrzebna w tym dowodzie.
\end{proof}

\subsection{Granice i kontrprzykłady}

\begin{twierdzenie}[Słabe domknięcie sfery jednostkowej]
\label{thm:sphere-weak-dense26}
Niech $X$ będzie nieskończenie wymiarową przestrzenią
unormowaną. Wówczas słabe domknięcie sfery jednostkowej
$S_X=\{x\in X:\|x\|=1\}$ jest równe $B_X$.
\end{twierdzenie}
\begin{proof}
Jest to Twierdzenie~\ref{thm:sfera-slaba}; przypomnijmy
argument w kontekście geometrii kuli.
Kula $B_X$ jest wypukła i normowo domknięta,
a więc słabo domknięta
(Wniosek~\ref{cor:weak-convex-closure}).
Ponieważ $S_X\subseteq B_X$, mamy
$\overline{S_X}^{\,w}\subseteq B_X$.
Punkty sfery należą do jej domknięcia.
Niech więc $x\in B_X$ i $\|x\|<1$.
W dowolnym słabym otoczeniu $x$ istnieje
otoczenie podstawowe
\[
 U=\{y\in X:|f_j(y-x)|<\eta,\ j=1,\ldots,m\},
\]
dla pewnych $f_1,\ldots,f_m\in X'$ i $\eta>0$.
Liniowe odwzorowanie $y\mapsto(f_1(y),\ldots,f_m(y))$
do $\K^m$ nie może być różnowartościowe, gdyż
$X$ jest nieskończenie wymiarowa. Istnieje zatem
niezerowy $v\in\bigcap_{j=1}^m\ker f_j$.
Dla każdego $t\ge0$ mamy $x+tv\in U$, gdyż
$f_j(x+tv-x)=0$.
Funkcja $t\mapsto\|x+tv\|$ jest ciągła
(z odwrotnej nierówności trójkąta), w zerze
przyjmuje wartość $\|x\|<1$, a ponadto
$\|x+tv\|\ge t\|v\|-\|x\|\to+\infty$.
Z twierdzenia o wartości pośredniej otrzymujemy
$t_0>0$ z $\|x+t_0v\|=1$.
Zatem $x+t_0v\in U\cap S_X$.
Każde słabe otoczenie $x$ przecina sferę,
co daje $B_X\subseteq\overline{S_X}^{\,w}$
i kończy dowód.
\end{proof}

\begin{uwaga}[Co wynika z twierdzenia Eberleina--\v Smuliana]
\label{rem:weak-compact-what26}
Twierdzenie~\ref{thm:eberlein-smulian26} charakteryzuje
względną słabą zwartość: $\overline A^{\,w}$ jest
słabo zwarte wtedy i tylko wtedy, gdy każdy ciąg
w $A$ ma podciąg słabo zbieżny do punktu $X$.
Nie wymaga się tutaj, aby granica należała do $A$.
Na przykład $A=(0,1)\subseteq\R$ spełnia ten warunek,
ale nie jest zwarty: ciąg $1/(n+1)$ nie ma podciągu
z granicą w $A$.

Dla samego zbioru $A$ słaba zwartość jest równoważna
temu, że każdy ciąg w $A$ ma podciąg słabo zbieżny
do punktu $A$. Jeżeli $A$ jest słabo zwarty, jest
słabo domknięty (topologia słaba jest Hausdorffa);
twierdzenie Eberleina--\v Smuliana daje wymagany
podciąg z granicą w $\overline A^{\,w}=A$.
Odwrotnie, założenie o podciągach z granicą w $A$
zapewnia względną słabą zwartość. Dla dowolnego
$x\in\overline A^{\,w}$ zbiór $A$ jest zawarty
w słabo zwartej przestrzeni $\overline A^{\,w}$.
Samo twierdzenie o względnej zwartości nie pozwala
jednak zastąpić domknięcia domknięciem ciągowym
w dowolnej przestrzeni topologicznej.
W naszym przypadku stosujemy konstrukcję z kroku
trzeciego dowodu Twierdzenia~\ref{thm:eberlein-smulian26}
do $\Phi=J_Xx$: skonstruowany ciąg w $A$ ma
podciąg słabo zbieżny do $y\in A$. Punkt $y$ jest
słabym punktem skupienia tego ciągu, więc ten sam
argument normujący daje $\Phi=J_Xy$.
Z różnowartościowości $J_X$ mamy $x=y\in A$.
Stąd $A=\overline A^{\,w}$ i $A$ jest słabo zwarty.

Słaba zwartość nie oznacza zwartości normowej.
W nieskończenie wymiarowej przestrzeni Hilberta
kula $B_H$ jest słabo zwarta na mocy
Twierdzeń~\ref{thm:milman-pettis26}
i~\ref{thm:reflexive-ball26}, ale zawiera ciąg
ortonormalny z Przykładu~\ref{ex:en-hilbert26}.
Dla $n\ne m$ mamy $\|e_n-e_m\|^2=2$, więc żaden
jego podciąg nie jest normowo zbieżny.
Z kolei w $\ell^1$ sfera ma zero w słabym
domknięciu na mocy Twierdzenia~\ref{thm:sphere-weak-dense26},
ale żaden ciąg jej punktów nie zbiega słabo do zera,
gdyż własność Schura wymuszałaby zbieżność norm do zera.
W przestrzeniach innych niż przestrzenie Banacha
równoważności zwartości słabej i słabej zwartości
ciągowej nie wolno zakładać bez dodatkowego dowodu.
\end{uwaga}

\begin{przyklad}[Refleksywność bez jednostajnej wypukłości danej normy]
\label{ex:reflexive-not-uniform26}
Niech $X=\K^2$ z normą $\|(a,b)\|_\infty=\max\{|a|,|b|\}$.
Jest to przestrzeń Banacha, a jej kula jednostkowa
jest normowo zwarta jako domknięty i ograniczony
zbiór w skończonym wymiarze. Jest więc również
słabo zwarta, ponieważ identyczność z topologii
normowej do słabej jest ciągła.
Twierdzenie~\ref{thm:reflexive-ball26} daje refleksywność.
Jednak dla $u=(1,1)$ i $v=(1,-1)$ zachodzi
$\|u\|_\infty=\|v\|_\infty=1$,
$\|u-v\|_\infty=2$ oraz $\|(u+v)/2\|_\infty=1$.
Norma nie jest nawet ściśle wypukła, a więc nie
jest jednostajnie wypukła.
Przykład dotyczy ustalonej normy: ta sama przestrzeń
z normą euklidesową jest jednostajnie wypukła.
\end{przyklad}

% Wersja źródłowa: audyt II, jednostka 27
\setcounter{section}{26}
\noindent\textbf{Wyniki wykorzystywane.}
Przyjmujemy znane twierdzenia Hahna--Banacha,
Banacha--Steinhausa, o odwzorowaniu otwartym,
twierdzenie Eberleina--\v Smuliana oraz podstawy
ortogonalności i bazy ortonormalnej przestrzeni Hilberta.
\par\medskip
\section{Bazy Schaudera, układy biortogonalne i zbieżność bezwarunkowa}
\label{sec:bazy27}

W tym rozdziale $X$ oznacza przestrzeń Banacha nad $\K\in\{\R,\C\}$,
$X'$ jej przestrzeń dualną topologiczną, a $\mathcal L(X,Y)$
przestrzeń operatorów liniowych ciągłych. Bazy ortonormalne przestrzeni
Hilberta omówiono we wcześniejszym rozdziale o przestrzeniach Hilberta; tutaj badamy bazy
\emph{uporządkowane}, właściwe ogólnym przestrzeniom Banacha.
Nie należy mylić baz Schaudera z bazami Hamela: w przypadku tych
ostatnich dopuszcza się wyłącznie \emph{skończone} kombinacje liniowe.

\subsection{Definicja, współczynniki i ciągłość rzutów częściowych}

\begin{definicja}[Baza Schaudera i współczynniki]
\label{def:schauder27}
Niech $X$ będzie przestrzenią Banacha, a $(e_n)_{n\ge1}\subset X$
ciągiem wektorów niezerowych. Nazywamy go \emph{bazą Schaudera}
przestrzeni $X$, jeśli dla każdego $x\in X$ istnieje dokładnie jeden
ciąg skalarów $(a_n)_{n\ge1}$, dla którego
$x=\lim_{N\to\infty}\sum_{n=1}^N a_ne_n$ w normie $X$.
Skalar $a_n$ oznaczamy $e_n^*(x)$ i nazywamy $n$-tym
\emph{współczynnikiem} wektora $x$. Odwzorowanie
$P_Nx=\sum_{n=1}^Ne_n^*(x)e_n$ nazywamy
\emph{rzutem częściowym}; kładziemy też $P_0=0$.
\end{definicja}

Jednoznaczność implikuje liniowość $e_n^*$ i $P_N$; nie zakładaliśmy
natomiast jeszcze ich ciągłości. Jest to istotne: ciągłość
współczynników stanowi treść następnego twierdzenia.

\begin{lemat}[Przestrzeń ciągów współczynników]
\label{lem:coefficient-space27}
Niech $X$ będzie przestrzenią Banacha, a $(e_n)_{n\ge1}\subset X$
ciągiem wektorów niezerowych. Zbiór
\[
E_e=\bigl\{a=(a_n)\in\K^{\N}:\sum_{n=1}^{\infty}a_ne_n
\text{ jest zbieżny w }X\bigr\}
\]
z normą $\|a\|_{E_e}=\sup_{N\ge1}\|\sum_{n=1}^Na_ne_n\|_X$
jest przestrzenią Banacha. Ponadto odwzorowanie
$S:E_e\to X$, $S(a)=\sum_{n\ge1}a_ne_n$, jest liniowe
ciągłe i spełnia $\|S\|\le1$.
\end{lemat}
\begin{proof}
Norma jest skończona dzięki zbieżności sum częściowych.
Jeżeli wszystkie sumy częściowe znikają, to z różnic kolejnych
sum i z $e_n\ne0$ wynika $a_n=0$ dla wszystkich $n$.
Niech $(a^{(k)})$ będzie ciągiem Cauchy'ego w $E_e$.
Dla każdego $n$ mamy
$|a_n^{(k)}-a_n^{(\ell)}|\,\|e_n\|
\le2\|a^{(k)}-a^{(\ell)}\|_{E_e}$,
więc istnieje granica $a_n=\lim_k a_n^{(k)}$.
Przechodząc do granicy przy ustalonym $N$, otrzymujemy
$\|\sum_{n=1}^N(a_n^{(k)}-a_n)e_n\|\le\eta$
jednostajnie w $N$, gdy $k$ jest dostatecznie duże,
bo analogiczne oszacowanie z $a^{(\ell)}$ jest
prawdziwe dla każdego $\ell$.
Wybieramy takie $k$. Dla $M>N$ mamy
\[
\Big\|\sum_{n=N+1}^Ma_ne_n\Big\|
\le2\eta+\Big\|\sum_{n=N+1}^Ma_n^{(k)}e_n\Big\|.
\]
W ostatnim wyrażeniu, dla ustalonego $k$, sumy ogonowe
są dowolnie małe przy dużym $N$. Ponieważ $\eta>0$
można przyjąć dowolnie małe, szereg dla $a$ jest zbieżny.
Ponadto $\|a^{(k)}-a\|_{E_e}\to0$.
Dla każdego $a\in E_e$ mamy
$\|S(a)\|\le\|a\|_{E_e}$, co dowodzi tezy o $S$.
\end{proof}

\begin{twierdzenie}[Ciągłość współczynników i stała bazy]
\label{thm:basis-continuity27}
Jeżeli $(e_n)_{n\ge1}$ jest bazą Schaudera przestrzeni
Banacha $X$, to $e_n^*\in X'$ dla każdego $n$, a każdy
$P_N\in\mathcal L(X,X)$. Ponadto
\[
K_e:=\sup_{N\ge1}\|P_N\|<\infty,\qquad K_e\ge1.
\]
Liczbę $K_e$ nazywamy \emph{stałą bazy}.
Dla każdego $n\ge1$ zachodzi $e_n^*(e_m)=\delta_{nm}$
oraz $\|e_n^*\|\,\|e_n\|\le2K_e$.
\end{twierdzenie}
\begin{proof}
Dla bazy Schaudera odwzorowanie $S:E_e\to X$
z Lematu~\ref{lem:coefficient-space27} jest liniową ciągłą
bijekcją. Twierdzenie o odwzorowaniu otwartym zapewnia
ciągłość $S^{-1}$. Dla każdego $N$ zachodzi
$\|P_Nx\|=\|\sum_{n=1}^Ne_n^*(x)e_n\|
\le\|S^{-1}x\|_{E_e}$,
stąd $K_e\le\|S^{-1}\|<\infty$.
Dla $n\ge1$ tożsamość
$(P_n-P_{n-1})x=e_n^*(x)e_n$
daje $\|e_n^*\|\,\|e_n\|\le\|P_n\|+\|P_{n-1}\|
\le2K_e$. Zatem wszystkie współczynniki są ciągłe.
Biortogonalność wynika z jednoznaczności reprezentacji $e_m$.
Każdy $P_N$ jest niezerowym rzutem (pozostawia $e_1$),
a więc $\|P_N\|\ge1$.
\end{proof}

\begin{wniosek}[Normalizacja i ośrodkowość]
\label{cor:separable27}
Każda przestrzeń Banacha z bazą Schaudera jest ośrodkowa.
Jeżeli $(e_n)$ jest bazą Schaudera, to
$(e_n/\|e_n\|)$ również jest bazą Schaudera o tej samej stałej
bazy; jej współczynniki są równe $\|e_n\|e_n^*(x)$.
\end{wniosek}
\begin{proof}
Skończone kombinacje $e_n$ o współczynnikach w $\mathbb Q$
(dla $\K=\C$ w $\mathbb Q+i\mathbb Q$) tworzą zbiór przeliczalny
i gęsty w $X$, bo sumy częściowe bazy przybliżają każdy $x$.
Zmiana normalizacji nie zmienia ani sum częściowych,
ani ich granic, zatem nie zmienia operatorów $P_N$.
\end{proof}

\begin{definicja}[Układ biortogonalny i ciąg bazowy]
\label{def:biorthogonal27}
Niech $X$ będzie przestrzenią Banacha.
Układ $(e_n,f_n)_{n\ge1}\subset X\times X'$ nazywamy
\emph{biortogonalnym}, jeżeli $f_n(e_m)=\delta_{nm}$
dla wszystkich $m,n$. Ciąg $(e_n)$ nazywamy
\emph{ciągiem bazowym}, jeżeli jest bazą Schaudera swojej
normowo domkniętej otoczki liniowej $[e_n]
=\overline{\operatorname{span}\{e_n:n\ge1\}}^{\,\|\cdot\|}$.
\end{definicja}

\begin{twierdzenie}[Kryterium bazy przez rzuty]
\label{thm:biorthogonal-criterion27}
Niech $X$ będzie przestrzenią Banacha, zaś
$(e_n,f_n)_{n\ge1}\subset X\times X'$ układem biortogonalnym,
dla którego $\operatorname{span}\{e_n:n\ge1\}$ jest gęsta w $X$.
Oznaczmy $P_Nx=\sum_{n=1}^Nf_n(x)e_n$.
Wówczas $(e_n)$ jest bazą Schaudera $X$, której
funkcjonałami współrzędnych są $f_n$, wtedy i tylko wtedy,
gdy $\sup_N\|P_N\|<\infty$.
\end{twierdzenie}
\begin{proof}
Warunek jest konieczny na mocy
Twierdzenia~\ref{thm:basis-continuity27}.
W przeciwną stronę niech $K=\sup_N\|P_N\|<\infty$.
Jeśli $y$ jest skończoną kombinacją $e_n$, to $P_Ny=y$
dla dostatecznie dużych $N$. Dla dowolnego $x\in X$
wybieramy takie $y$ bliskie $x$; wówczas
$\|x-P_Nx\|\le(1+K)\|x-y\|$ dla dużych $N$.
Zatem $P_Nx\to x$ w normie.
Jeżeli również $x=\sum_na_ne_n$, to przez ciągłość $f_j$
i biortogonalność dostajemy $f_j(x)=a_j$.
To dowodzi jednoznaczności rozwinięcia.
\end{proof}

\begin{wniosek}[Kryterium ciągu bazowego]
\label{cor:basic-sequence27}
Niech $(e_n)_{n\ge1}$ będzie ciągiem wektorów niezerowych
w przestrzeni Banacha $X$. Jest to ciąg bazowy wtedy i tylko wtedy,
gdy jest liniowo niezależny i istnieje $K<\infty$ takie, że
$\|\sum_{n=1}^Na_ne_n\|
\le K\|\sum_{n=1}^Ma_ne_n\|$
dla wszystkich $M\ge N\ge1$ i wszystkich $a_1,\ldots,a_M\in\K$.
\end{wniosek}
\begin{proof}
Dla ciągu bazowego nierówność wynika z ograniczoności
rzutów częściowych na $[e_n]$.
Odwrotnie, na algebraicznej otoczce liniowej $e_n$
jednoznacznie definiujemy $P_N$ przez ucięcie współczynników.
Założona nierówność daje $\|P_N\|\le K$.
Każdy $P_N$ przedłuża się do $[e_n]$ przez ciągłość;
wówczas $P_Nx\to x$ przez gęstość skończonych kombinacji.
Współczynniki są ciągłe, ponieważ
$(P_n-P_{n-1})x$ należy do jednowymiarowej przestrzeni
$\K e_n$, a ciągłość funkcjonałów współrzędnych
wynika z podzielenia przez stały niezerowy wektor $e_n$.
Stosujemy Twierdzenie~\ref{thm:biorthogonal-criterion27}
do przestrzeni $[e_n]$.
\end{proof}

\subsection{Bazy monotoniczne i klasyczne przykłady}

\begin{definicja}[Baza monotoniczna]
\label{def:monotonic27}
Bazę Schaudera $(e_n)$ nazywamy \emph{monotoniczną}, gdy
$\|P_N\|=1$ dla wszystkich $N\ge1$.
\end{definicja}

\begin{przyklad}[Przestrzenie ciągowe]
\label{ex:sequence-bases27}
W przestrzeniach Banacha $c_0(\N,\K)$ i $\ell^p(\N,\K)$
dla $1\le p<\infty$ wektory standardowe
$e_n=(\delta_{nj})_{j\ge1}$ tworzą bazę Schaudera.
Współczynnik $e_n^*(x)$ to $n$-ta współrzędna,
a $P_N(x)$ zeruje wszystkie współrzędne po $N$.
Dla $\ell^p$ zachodzi
$\|x-P_Nx\|_p^p=\sum_{j>N}|x_j|^p\to0$,
a dla $c_0$ mamy $\|x-P_Nx\|_\infty=\sup_{j>N}|x_j|\to0$.
Stąd $\|P_N\|=1$ i bazy są monotoniczne.
W $\ell^\infty$ ciąg standardowy nie jest bazą Schaudera:
jego domknięta otoczka liniowa równa jest $c_0\ne\ell^\infty$.
\end{przyklad}

\begin{przyklad}[Baza sumująca w $c_0$]
\label{ex:summing-basis27}
Niech $s_n=\sum_{j=1}^ne_j\in c_0$.
Dla $x=(x_j)\in c_0$ przyjmijmy $a_n=x_n-x_{n+1}$.
Wtedy $\sum_{n=1}^Na_ns_n$ ma na pozycjach $j\le N$
wartość $x_j-x_{N+1}$, a po pozycjach $j>N$ wartość $0$.
Zatem
$\|x-\sum_{n=1}^Na_ns_n\|_\infty
\le2\sup_{j>N}|x_j|\to0$.
Jednoznaczność współczynników wynika z tożsamości
$a_n=x_n-x_{n+1}$; układ $(s_n)$ jest więc bazą Schaudera.
Jej rzuty częściowe mają normę $2$ dla każdego $N\ge1$,
ponieważ $\|P_Nx\|_\infty
\le\|x\|_\infty+|x_{N+1}|\le2\|x\|_\infty$,
a równość otrzymujemy przez wybór $x_1=1$, $x_{N+1}=-1$
i pozostałych współrzędnych równych zeru (dla $N\ge1$).
Ta baza nie jest monotoniczna.
\end{przyklad}

\subsection{Bazy bezwarunkowe i zmiana kolejności sumowania}

\begin{definicja}[Baza bezwarunkowa]
\label{def:unconditional27}
Bazę Schaudera $(e_n)$ przestrzeni Banacha $X$ nazywamy
\emph{bezwarunkową}, jeżeli dla każdego $x\in X$ i każdej
permutacji $\pi:\N\to\N$ szereg
$\sum_{n=1}^\infty e_{\pi(n)}^*(x)e_{\pi(n)}$
jest zbieżny w normie $X$ do $x$.
Dla skończonego $F\subseteq\N$ piszemy
$P_Fx=\sum_{n\in F}e_n^*(x)e_n$.
\end{definicja}

\begin{lemat}[Permutacje a sumy po zbiorach skończonych]
\label{lem:net-unconditional27}
Niech $(u_n)$ będzie ciągiem w przestrzeni Banacha $X$
i załóżmy, że szereg $\sum_nu_n$ jest zbieżny w zadanej
kolejności do $u\in X$. Następujące warunki są równoważne:
\begin{enumerate}[label=\textup{(\roman*)}]
\item wszystkie permutacje szeregu są zbieżne do $u$;
\item sieć $(\sum_{n\in F}u_n)_{F\Subset\N}$, skierowana przez inkluzję,
 zbiega do $u$;
\item dla każdego $\varepsilon>0$ istnieje skończony $F_0\subseteq\N$
 taki, że $\|\sum_{n\in F}u_n\|<\varepsilon$
 dla każdego skończonego $F\subseteq\N\setminus F_0$.
\end{enumerate}
\end{lemat}
\begin{proof}
Jeśli (iii) zachodzi, to dla skończonych
$F,G\supseteq F_0$ różnica ich sum jest różnicą sum po
$F\setminus G$ i $G\setminus F$. Oba zbiory są rozłączne
z $F_0$, zatem nierówność trójkąta daje normę różnicy
mniejszą niż $2\varepsilon$. Zatem sieć jest
siecią Cauchy'ego. W przestrzeni Banacha ma granicę,
która musi być równa $u$, ponieważ zawiera ona jako podsieć
sumy po $\{1,\ldots,N\}$.
Implikacja (ii) $\Rightarrow$ (i) wynika stąd, że
zbiory $\{\pi(1),\ldots,\pi(N)\}$ wyczerpują $\N$.
Załóżmy, że (iii) nie zachodzi. Istnieje $\varepsilon>0$
takie, że po usunięciu dowolnego zbioru skończonego
można znaleźć zbiór skończony $F$ z normą sumy
przynajmniej $\varepsilon$. Rekurencyjnie konstruujemy
rozłączne zbiory skończone $F_k$, których elementy są coraz
większe, i których sumy mają normę przynajmniej $\varepsilon$.
Każde $F_k$ wybieramy poza poprzednimi zbiorami i poza
$\{1,\ldots,k\}$. Permutację $\N$ budujemy blokami,
wyliczając kolejno elementy $F_k$ i wstawiając
między blokami najmniejszy dotąd niewykorzystany indeks.
Wyrazy tak przestawionego szeregu mają sumy częściowe,
które na końcach każdego bloku $F_k$ różnią się o co najmniej
$\varepsilon$. Szereg po tej permutacji nie jest więc
zbieżny, co przeczy (i).
\end{proof}

\begin{twierdzenie}[Kryterium bezwarunkowości przez rzuty]
\label{thm:unconditional-projections27}
Niech $(e_n)$ będzie bazą Schaudera przestrzeni Banacha $X$.
Wówczas następujące warunki są równoważne:
\begin{enumerate}[label=\textup{(\roman*)}]
\item $(e_n)$ jest bazą bezwarunkową;
\item istnieje $U<\infty$ takie, że $\|P_F\|\le U$
 dla każdego skończonego $F\subseteq\N$;
\item dla każdego $x\in X$ rodzina
$\{P_Fx:F\Subset\N\}$ jest ograniczona w normie.
\end{enumerate}
Najmniejsze $U$ z (ii) nazywamy \emph{stałą bezwarunkowości
przez opuszczanie współczynników}.
\end{twierdzenie}
\begin{proof}
Z (i) i Lematu~\ref{lem:net-unconditional27} wynika, że
$P_Fx\to x$ w sieci zbiorów skończonych, zatem rodzina
$\{P_Fx\}$ jest ograniczona. To daje (iii).
Twierdzenie Banacha--Steinhausa zastosowane do operatorów
$P_F\in\mathcal L(X,X)$ daje (iii) $\Rightarrow$ (ii).
Jeżeli zachodzi (ii), to dla $x\in X$ i dla przybliżenia
$y\in\operatorname{span}\{e_n\}$ mamy
$\|P_F(x-y)\|\le U\|x-y\|$.
Gdy $F$ zawiera wszystkie indeksy występujące w $y$,
to $P_Fy=y$, skąd $\|P_Fx-x\|\le(U+1)\|x-y\|$.
Gęstość skończonych kombinacji implikuje zbieżność
sieci $P_Fx\to x$, a Lemat~\ref{lem:net-unconditional27}
daje (i).
\end{proof}

\begin{wniosek}[Bazy ortonormalne i bezwarunkowość]
\label{cor:hilbert-unconditional27}
Niech $H$ będzie ośrodkową, nieskończenie wymiarową
przestrzenią Hilberta, a $(e_n)_{n\ge1}$ dowolną numeracją
jej bazy ortonormalnej. Wtedy $(e_n)$ jest monotoniczną
bazą Schaudera i bezwarunkową o stałej $U=1$.
\end{wniosek}
\begin{proof}
Twierdzenie Parsevala daje zbieżność rozwinięcia
$\sum_n\langle x,e_n\rangle e_n$ do $x$.
Dla każdego skończonego $F$ twierdzenie Pitagorasa
daje $\|P_Fx\|^2=\sum_{n\in F}|\langle x,e_n\rangle|^2
\le\|x\|^2$. Zarówno $P_F$, jak i $P_N$ są
niezerowymi rzutami, więc mają normę dokładnie $1$.
Twierdzenie~\ref{thm:unconditional-projections27} kończy dowód.
\end{proof}

\begin{przyklad}[Baza warunkowa w przestrzeni $c_0$]
\label{ex:conditional-summing27}
Baza sumująca $(s_n)$ z Przykładu~\ref{ex:summing-basis27}
nie jest bezwarunkowa. Weźmy $x\in c_0$ o współrzędnych
$x_n=(-1)^n/\sqrt n$. Wtedy
$a_n=s_n^*(x)=x_n-x_{n+1}=(-1)^n(1/\sqrt n+1/\sqrt{n+1})$.
Szereg skalarów $\sum_n a_n$ jest zbieżny (jego suma to
$x_1$), ale nie jest bezwzględnie zbieżny.
Gdyby szereg wektorowy $\sum_n a_ns_n$ był bezwarunkowo
zbieżny w $c_0$, to po zastosowaniu ciągłego
funkcjonału $\delta_1(y)=y_1$ uzyskalibyśmy bezwarunkową
zbieżność $\sum_n a_n$ w $\K$, co jest równoważne
zbieżności bezwzględnej. Otrzymujemy sprzeczność.
\end{przyklad}

\begin{przyklad}[Bezwarunkowość nie oznacza absolutnej zbieżności]
\label{ex:not-absolute27}
W przestrzeni Hilberta $\ell^2$ szereg
$\sum_{n\ge1}n^{-1}e_n$ jest zbieżny bezwarunkowo
(każdy wybór skończonej rodziny wyrazów ma kwadrat
normy równy sumie $n^{-2}$), ale
$\sum_{n\ge1}\|n^{-1}e_n\|=\sum_{n\ge1}n^{-1}=\infty$.
\end{przyklad}

\subsection{Bazy ściągające i zupełne w sensie ograniczonym}

\begin{definicja}[Dwa warunki na bazę]
\label{def:shrinking-bcomplete27}
Niech $(e_n)$ będzie bazą Schaudera przestrzeni Banacha $X$,
a $(e_n^*)\subset X'$ jej funkcjonałami współrzędnych.
Bazę nazywamy \emph{ściągającą} (ang. \emph{shrinking}), gdy
$\overline{\operatorname{span}\{e_n^*:n\ge1\}}^{\,\|\cdot\|}=X'$.
Nazywamy ją \emph{zupełną w sensie ograniczonym}
(ang. \emph{boundedly complete}), jeżeli dla każdego ciągu
skalarów $(a_n)$ warunek
$\sup_N\|\sum_{n=1}^Na_ne_n\|<\infty$
implikuje zbieżność szeregu $\sum_na_ne_n$ w normie $X$.
\end{definicja}

\begin{twierdzenie}[Kryterium refleksywności przez bazę]
\label{thm:james-basis27}
Niech $X$ będzie przestrzenią Banacha z bazą Schaudera $(e_n)$.
Wtedy $X$ jest refleksywna wtedy i tylko wtedy, gdy
$(e_n)$ jest jednocześnie ściągająca i zupełna w sensie
ograniczonym.
\end{twierdzenie}
\begin{proof}
Załóżmy najpierw refleksywność $X$.
Aby dowieść zupełności w sensie ograniczonym,
niech $s_N=\sum_{n=1}^Na_ne_n$ będzie ograniczony.
Kula zawierająca wszystkie $s_N$ jest słabo zwarta.
Istnieje zatem podciąg słabo zbieżny do pewnego $x\in X$
na mocy twierdzenia Eberleina--\v Smuliana.
Dla każdego $j$ ciąg $e_j^*(s_N)$ jest od pewnego miejsca
stały i równy $a_j$, więc $e_j^*(x)=a_j$.
Stąd $P_Nx=s_N$ dla wszystkich $N$, a definicja bazy
daje $s_N\to x$ w normie.

Niech teraz $f\in X'$. Oznaczmy
$Y_N=\overline{\operatorname{span}\{e_j:j>N\}}^{\,\|\cdot\|}$.
Pokażemy, że $\|f|_{Y_N}\|\to0$, gdzie norma oznacza
normę ograniczenia $f$ do domkniętej podprzestrzeni $Y_N$.
Gdyby tak nie było, istniałyby $\varepsilon>0$, liczby
$N_k\nearrow\infty$ i wektory $y_k\in Y_{N_k}$
o normie $1$ takie, że $|f(y_k)|\ge\varepsilon$.
Wektory te można wybierać rekurencyjnie, gdyż $Y_{N+1}\subseteq Y_N$,
a więc normy ograniczeń tworzą ciąg nierosnący.
Refleksywność i twierdzenie Eberleina--\v Smuliana
dają podciąg $(y_{k_\ell})$ słabo zbieżny do $y\in X$.
Dla każdego ustalonego $j$ mamy
$e_j^*(y_{k_\ell})=0$ dla dużych $\ell$, więc $e_j^*(y)=0$.
Jednoznaczność rozwinięcia wymusza $y=0$,
co jest sprzeczne z $|f(y_{k_\ell})|\ge\varepsilon$.
Ponieważ $(I-P_N)X\subseteq Y_N$ i $\|I-P_N\|\le1+K_e$,
dostajemy $\|f-P_N'f\|\le(1+K_e)\|f|_{Y_N}\|\to0$.
Funkcjonał $P_N'f=f\circ P_N$ jest skończoną kombinacją
$e_1^*,\ldots,e_N^*$; baza jest zatem ściągająca.

Odwrotnie, niech baza ma obie własności i niech
$\Phi\in X''$. Kładziemy $a_n=\Phi(e_n^*)$ oraz
$s_N=\sum_{n=1}^Na_ne_n$. Dla $f\in X'$ mamy
$f(s_N)=\Phi(P_N'f)$, a więc
$|f(s_N)|\le K_e\|\Phi\|\,\|f\|$.
Z definicji normy wynika $\|s_N\|\le K_e\|\Phi\|$.
Zupełność w sensie ograniczonym daje
$x=\lim_Ns_N\in X$. Dla każdego $n$ otrzymujemy
$e_n^*(x)=\Phi(e_n^*)$.
Zatem $J_Xx$ i $\Phi$ pokrywają się na liniowej otoczce
$\{e_n^*\}$, która jest gęsta w $X'$ dzięki ściąganiu bazy.
Ciągłość obu funkcjonałów na $X'$ daje $J_Xx=\Phi$.
Kanoniczne odwzorowanie $J_X:X\to X''$ jest surjektywne,
czyli $X$ jest refleksywna.
\end{proof}

\begin{przyklad}[Przykłady rozdzielające oba warunki]
\label{ex:james-examples27}
Standardowa baza przestrzeni $c_0$ jest ściągająca,
gdyż $c_0'=\ell^1$, a jej funkcjonały współrzędnych
odpowiadają wektorom standardowym w $\ell^1$.
Nie jest zupełna w sensie ograniczonym: sumy
$\sum_{n=1}^Ne_n$ mają normę $1$ w $c_0$, ale
nie tworzą ciągu Cauchy'ego.
Standardowa baza przestrzeni $\ell^1$ jest zupełna w sensie
ograniczonym, ponieważ ograniczoność sum częściowych oznacza
$\sup_N\sum_{n=1}^N|a_n|<\infty$.
Nie jest ściągająca, gdyż $(\ell^1)'=\ell^\infty$,
a domknięcie liniowej otoczki funkcjonałów współrzędnych
to właściwa podprzestrzeń $c_0$ przestrzeni $\ell^\infty$.
\end{przyklad}

\subsection{Czego istnienie bazy nie gwarantuje?}

\begin{uwaga}[Baza ortonormalna a baza Hamela]
\label{rem:hamel-vs-schauder27}
Jeżeli przestrzeń Hilberta jest ośrodkowa i nieskończenie
wymiarowa, to przeliczalna baza ortonormalna jest bazą
Schaudera, lecz nigdy bazą Hamela całej przestrzeni:
wektory z bazy Hamela dają skończone kombinacje liniowe.
Co więcej, dla nieskończenie wymiarowej przestrzeni Hilberta
każda baza Hamela jest nieprzeliczalna. Dowody tych faktów,
a także przeliczalność niezerowych współczynników
względem dowolnej, nawet nieprzeliczalnej rodziny
ortonormalnej, znajdują się we wcześniejszym rozdziale o przestrzeniach Hilberta.
\end{uwaga}

\begin{uwaga}[Ośrodkowość nie wystarcza]
\label{rem:enflo27}
Istnieją ośrodkowe przestrzenie Banacha bez bazy Schaudera.
Wynik ten ma charakter głęboki i nie wynika z samej
nierówności Banacha--Steinhausa. Konstrukcje tego typu
wiążą się z niepowodzeniem własności aproksymacji,
którą omówimy w następnym rozdziale. Istnienie
\emph{ciągu bazowego} w każdej nieskończenie wymiarowej
przestrzeni Banacha nie oznacza, że jest on bazą całej
przestrzeni.
\end{uwaga}

% Wersja źródłowa: audyt II, jednostka 28
\setcounter{section}{27}
\noindent\textbf{Wyniki wykorzystywane.}
Przyjmujemy twierdzenia Hahna--Banacha i Banacha--Steinhausa,
teorię baz Schaudera i norm operatorowych, podstawy przestrzeni
Hilberta oraz kryterium ośrodkowości. Wszystkie nowe konstrukcje
i ich zależności są formułowane i dowodzone w tym rozdziale.
\par\medskip
\section{Własności aproksymacji przestrzeni Banacha}
\label{sec:aproksymacja28}

W całym rozdziale $X,Y,Z$ są przestrzeniami Banacha nad
$\K\in\{\R,\C\}$. Przez $\mathcal L(X,Y)$ oznaczamy
przestrzeń operatorów liniowych ciągłych z normą operatorową, zaś
$\mathcal F(X,Y)$ -- zbiór tych $T\in\mathcal L(X,Y)$, których
obraz ma skończony wymiar. Wyrażenie \emph{operator skończonego rzędu}
nie oznacza operatora będącego jedynie odwzorowaniem algebraicznym:
zawsze wymagana jest jego ciągłość. W rozdziale o bazach Schaudera
wykazano, że rzuty częściowe takiej bazy są jednostajnie ograniczonymi
operatorami skończonego rzędu. Tutaj wyjaśnimy, jaka własność
przestrzeni za tym stoi i dlaczego nie każda przestrzeń Banacha ją ma.

\subsection{Aproksymowanie jednostajnie na zbiorach zwartych}

\begin{definicja}[AP, BAP i MAP]
\label{def:ap-bap-map28}
Niech $X$ będzie przestrzenią Banacha, a $\lambda\ge1$.
Mówimy, że $X$ ma \emph{własność aproksymacji} (AP), jeżeli dla
każdego zwartego $K\subseteq X$ i każdego $\varepsilon>0$
istnieje $T\in\mathcal F(X,X)$ spełniający
$\sup_{x\in K}\|Tx-x\|<\varepsilon$.
Mówimy, że $X$ ma \emph{$\lambda$-ograniczoną własność aproksymacji}
($\lambda$-BAP), jeżeli można ponadto zawsze wybrać taki $T$
ze $\|T\|\le\lambda$. Własność BAP oznacza istnienie
$\lambda<\infty$ spełniającego ten warunek, natomiast
\emph{metryczna własność aproksymacji} (MAP) jest równoznaczna
z $1$-BAP.
\end{definicja}

Przestrzenie skończenie wymiarowe mają MAP: wystarczy wziąć
$T=I_X$. W przestrzeniach nieskończenie wymiarowych wymaganie
zbieżności w normie operatorowej byłoby zbyt silne.

\begin{stwierdzenie}[Brak aproksymacji identyczności w normie operatorowej]
\label{prop:no-norm-identity28}
Jeżeli $X$ jest nieskończenie wymiarową przestrzenią Banacha,
to dla każdego $T\in\mathcal F(X,X)$ zachodzi
$\|I_X-T\|\ge1$.
\end{stwierdzenie}
\begin{proof}
Ponieważ $T(X)$ ma wymiar skończony, $T$ nie może być
różnowartościowy na nieskończenie wymiarowej przestrzeni $X$.
Wybieramy $x\in\ker T$ o $\|x\|=1$. Wtedy
$\|(I_X-T)x\|=1$, co daje tezę.
\end{proof}

\begin{lemat}[Kryterium BAP na zbiorach skończonych]
\label{lem:finite-bap28}
Niech $X$ będzie przestrzenią Banacha, $\lambda\ge1$.
Następujące warunki są równoważne:
\begin{enumerate}[label=\textup{(\roman*)}]
\item $X$ ma $\lambda$-BAP;
\item dla wszystkich $m\ge1$, $x_1,\ldots,x_m\in X$ i
$\varepsilon>0$ istnieje $T\in\mathcal F(X,X)$ takie,
że $\|T\|\le\lambda$ oraz
$\max_{1\le j\le m}\|Tx_j-x_j\|<\varepsilon$.
\end{enumerate}
\end{lemat}
\begin{proof}
Pierwszy warunek implikuje drugi, ponieważ zbiór skończony
jest zwarty. W przeciwną stronę ustalmy zwarty $K\subset X$
i $\varepsilon>0$. Dla $K=\varnothing$ teza jest oczywista.
W przeciwnym razie istnieją $x_1,\ldots,x_m\in K$
takie, że dla każdego $x\in K$ pewien $j$ spełnia
$\|x-x_j\|<\varepsilon/[3(1+\lambda)]$.
Wybierzmy $T$ jak w (ii) z błędem mniejszym niż
$\varepsilon/3$. Wówczas
\[
\|Tx-x\|\le\|T(x-x_j)\|+\|Tx_j-x_j\|+\|x_j-x\|
<\varepsilon.
\]
W dowodzie korzystamy jednocześnie ze skończoności siatki
$\varepsilon$ i ze \emph{wspólnej granicy norm} aproksymantów.
\end{proof}

\begin{wniosek}[Hierarchia własności]
\label{cor:hierarchy28}
Dla każdej przestrzeni Banacha zachodzi
$\mathrm{MAP}\Rightarrow\mathrm{BAP}\Rightarrow\mathrm{AP}$.
Jeśli $X$ ma $\lambda$-BAP, to ma też $\mu$-BAP
dla każdego $\mu\ge\lambda$.
\end{wniosek}
\begin{proof}
Dwie pierwsze implikacje i monotoniczność stałej wynikają
bezpośrednio z Definicji~\ref{def:ap-bap-map28}.
\end{proof}

\begin{uwaga}[Dlaczego AP nie definiuje się na zbiorach skończonych?]
\label{rem:finite-exact28}
Dla \emph{dowolnej} przestrzeni Banacha $X$ i skończonego
zbioru $\{x_1,\ldots,x_m\}\subset X$ można znaleźć
$T\in\mathcal F(X,X)$ takie, że $Tx_j=x_j$ dla każdego $j$.
Istotnie, niech $E=\operatorname{span}\{x_1,\ldots,x_m\}$
i wybierzmy bazę $u_1,\ldots,u_r$ przestrzeni $E$.
Funkcjonały współrzędnych na $E$ są ciągłe, a twierdzenie
Hahna--Banacha przedłuża je do $f_1,\ldots,f_r\in X'$.
Operator $Tx=\sum_{j=1}^rf_j(x)u_j$ jest ciągły,
skończonego rzędu i jest identycznością na $E$.
Jego norma może jednak zależeć w sposób niekontrolowany
od wyboru $E$. To właśnie odróżnia powyższy fakt od BAP.
\end{uwaga}

\subsection{Topologia zbieżności na zbiorach zwartych i granice operatorów}

\begin{definicja}[Zbieżność zwarta i silna operatorów]
\label{def:topologies-ap28}
Niech $X,Y$ będą przestrzeniami Banacha. Topologia
\emph{zbieżności jednostajnej na zbiorach zwartych}
na $\mathcal L(X,Y)$ jest wyznaczona przez seminormy
$p_K(T)=\sup_{x\in K}\|Tx\|$, gdzie $K$ przebiega
zbiory zwarte w $X$. \emph{Silna topologia operatorowa}
(SOT) jest wyznaczona przez seminormy $p_x(T)=\|Tx\|$
indeksowane przez $x\in X$. W obu przypadkach
zbieżność rozumiemy także dla sieci, a nie tylko ciągów.
\end{definicja}

\begin{twierdzenie}[Operatorowa charakterystyka AP i BAP]
\label{thm:operator-character-ap28}
Niech $X$ będzie przestrzenią Banacha i $\lambda\ge1$.
\begin{enumerate}[label=\textup{(\roman*)}]
\item $X$ ma AP wtedy i tylko wtedy, gdy
$I_X$ należy do domknięcia $\mathcal F(X,X)$
w topologii zbieżności jednostajnej na zbiorach zwartych.
\item $X$ ma $\lambda$-BAP wtedy i tylko wtedy, gdy
$I_X$ należy do domknięcia zbioru
$\{T\in\mathcal F(X,X):\|T\|\le\lambda\}$ w SOT.
Dla tego ograniczonego zbioru aproksymantów można
równoważnie użyć topologii zbieżności zwartej.
\end{enumerate}
\end{twierdzenie}
\begin{proof}
Pierwszy punkt jest dokładnym przepisaniem definicji AP
na język otoczeń w topologii $p_K$; skończone przecięcia
otoczeń odpowiadają sumie skończonej liczby zbiorów zwartych.
Pierwsza część (ii) wynika z Lematu~\ref{lem:finite-bap28},
ponieważ otoczenia w SOT wyznaczają warunki na skończenie
wielu wektorach. Na zbiorze operatorów o normach
nieprzekraczających $\lambda$ zbieżność punktowa do
$I_X$ implikuje jednostajną aproksymację na każdym
ustalonym zbiorze zwartym: stosujemy argument ze
skończoną siatką z dowodu Lematu~\ref{lem:finite-bap28}.
Implikacja odwrotna jest oczywista.
\end{proof}

\begin{stwierdzenie}[Silna granica sieci skończonego rzędu istnieje zawsze]
\label{prop:universal-sot28}
Dla \emph{każdej} przestrzeni Banacha $X$ istnieje sieć
$(T_E)_{E\in\mathfrak F(X)}$ operatorów skończonego rzędu,
która zbiega w SOT do $I_X$. Tutaj $\mathfrak F(X)$
oznacza zbiór wszystkich skończenie wymiarowych podprzestrzeni
$E\subseteq X$, skierowany inkluzją.
\end{stwierdzenie}
\begin{proof}
Dla każdego $E\in\mathfrak F(X)$ wybierzmy
ciągły rzut $T_E:X\to E$, konstruowany jak
w Uwadze~\ref{rem:finite-exact28}. Jeżeli $x\in X$,
to dla wszystkich $E\supseteq\operatorname{span}\{x\}$
zachodzi $T_Ex=x$. Oznacza to $T_Ex\to x$
w normie dla każdego $x\in X$, czyli zbieżność w SOT.
\end{proof}

\begin{uwaga}[Niebezpieczne utożsamienie AP ze zbieżnością punktową]
\label{rem:net-vs-ap28}
Twierdzenie~\ref{prop:universal-sot28} nie daje AP:
konstruowana sieć może nie mieć \emph{wspólnego ograniczenia}
na normy operatorowe, a nawet nie musi być punktowo
ograniczona jako cała rodzina. Zbieżność sieci nie implikuje
ograniczoności całego jej zbioru wartości. Dlatego
zastosowanie do takiej sieci twierdzenia Banacha--Steinhausa
bez dodatkowych przesłanek byłoby błędem.
Istnieją przestrzenie Banacha bez AP, mimo że również
w nich zachodzi Twierdzenie~\ref{prop:universal-sot28}.
\end{uwaga}

\begin{twierdzenie}[Ciągi aproksymantów i ośrodkowość]
\label{thm:sequence-bap28}
Niech $X$ będzie ośrodkową przestrzenią Banacha.
Wówczas $X$ ma BAP wtedy i tylko wtedy, gdy istnieje
\emph{ciąg} $(T_n)_{n\ge1}\subset\mathcal F(X,X)$
taki, że $T_nx\to x$ dla każdego $x\in X$.
Każda przestrzeń Banacha, dla której istnieje taki ciąg,
jest ośrodkowa, bez żadnego założenia ośrodkowości z góry.
\end{twierdzenie}
\begin{proof}
Załóżmy $\lambda$-BAP i ośrodkowość.
Niech $(x_j)_{j\ge1}$ będzie gęstym ciągiem w $X$.
Na mocy Lematu~\ref{lem:finite-bap28} dla każdego $n$
istnieje $T_n\in\mathcal F(X,X)$ o normie nie większej
niż $\lambda$ i takim, że $\|T_nx_j-x_j\|<1/n$
dla $1\le j\le n$. Zatem $T_nx_j\to x_j$ dla każdego $j$.
Dla dowolnego $x\in X$ i $x_j$ bliskiego $x$ mamy
$\|T_nx-x\|\le(\lambda+1)\|x-x_j\|+\|T_nx_j-x_j\|$,
co daje $T_nx\to x$.
Odwrotnie, jeżeli $T_nx\to x$ dla każdego $x$, to rodzina
$\{T_n:n\ge1\}$ jest punktowo ograniczona. Twierdzenie
Banacha--Steinhausa daje $\sup_n\|T_n\|<\infty$.
Przybliżenie na skończonych zbiorach, a następnie
Lemat~\ref{lem:finite-bap28}, zapewniają BAP.
Ponadto $\bigcup_nT_n(X)$ leży w przeliczalnej sumie
skończenie wymiarowych podprzestrzeni. Jej liniowa otoczka
ma przeliczalny zbiór gęsty, a zawiera gęsty zbiór
$\{\lim_nT_nx:x\in X\}=X$ w swoim domknięciu.
Wobec tego $X$ jest ośrodkowa.
\end{proof}

\subsection{Bazy Schaudera a metryczna własność aproksymacji}

\begin{twierdzenie}[Baza Schaudera implikuje BAP]
\label{thm:basis-implies-bap28}
Niech $X$ będzie przestrzenią Banacha z bazą Schaudera
$(e_n)_{n\ge1}$ oraz rzutami częściowymi $P_N$.
Jeżeli $K_e=\sup_N\|P_N\|$, to $X$ ma $K_e$-BAP.
W szczególności baza monotoniczna daje MAP.
\end{twierdzenie}
\begin{proof}
Każdy $P_N$ ma skończony rząd i spełnia
$\|P_N\|\le K_e$, a definicja bazy daje
$P_Nx\to x$ w normie dla każdego $x\in X$.
Lemat~\ref{lem:finite-bap28} wraz z oszacowaniem
ze skończoną siatką dowodzi $K_e$-BAP.
Jeżeli baza jest monotoniczna, $K_e=1$.
\end{proof}

\begin{przyklad}[Przestrzenie ciągowe]
\label{ex:sequence-map28}
Przestrzenie Banacha $c_0$ i $\ell^p$ dla
$1\le p<\infty$ mają MAP: rzuty obcinające ciąg
po $N$ wyrazach mają normę $1$ i zbieżność punktową
do identyczności udowodniono już w rozdziale o bazach Schaudera.
Nie powtarzamy tamtego dowodu. Inaczej wygląda sytuacja
w przestrzeni $\ell^\infty$, mimo że i ona ma MAP.
\end{przyklad}

\begin{twierdzenie}[MAP w dowolnej przestrzeni Hilberta]
\label{thm:hilbert-map28}
Każda rzeczywista lub zespolona przestrzeń Hilberta $H$,
bez założenia ośrodkowości, ma MAP.
\end{twierdzenie}
\begin{proof}
Niech $K\subseteq H$ będzie zwarty i $\varepsilon>0$.
Dla $K=\varnothing$ nie ma nic do dowodzenia.
W przeciwnym razie wybieramy skończoną $\varepsilon/2$-siatkę
$\{x_1,\ldots,x_m\}\subset K$ i oznaczamy
$E=\operatorname{span}\{x_1,\ldots,x_m\}$.
Niech $P_E$ będzie rzutem ortogonalnym na $E$.
Jeśli $E\ne\{0\}$, to $\|P_E\|=1$, a w przeciwnym
razie $\|P_E\|=0$. Ponieważ $P_Ex_j=x_j$,
dla $x\in K$ i odpowiedniego $j$ zachodzi
$\|x-P_Ex\|\le\|x-x_j\|+\|P_E(x_j-x)\|
<\varepsilon$. Operator $P_E$ ma skończony rząd.
\end{proof}

\begin{twierdzenie}[MAP przestrzeni funkcji ciągłych na odcinku]
\label{thm:c01-map28}
Przestrzeń Banacha $C([0,1],\K)$ z normą supremum
ma MAP. Istnieje ciąg operatorów skończonego rzędu
$L_n$ o normach równych $1$, zbieżny punktowo do
identyczności.
\end{twierdzenie}
\begin{proof}
Dla $n\ge1$ przyjmijmy $t_k=k/n$ ($0\le k\le n$).
W szczególności operator $L_n$ zależy liniowo od wartości
$f(t_0),\ldots,f(t_n)$, co jawnie zapewnia skończoność jego rzędu.
Niech $L_nf$ będzie funkcją ciągłą, liniową na każdym
$[t_k,t_{k+1}]$ i spełniającą $L_nf(t_k)=f(t_k)$.
Operator $L_n$ jest liniowy; jego obraz leży w
$(n+1)$-wymiarowej przestrzeni funkcji odcinkami liniowych.
Ponieważ $L_nf(t)$ jest kombinacją wypukłą
$f(t_k)$ i $f(t_{k+1})$, mamy
$\|L_nf\|_\infty\le\|f\|_\infty$. Dla funkcji stałej
$1$ zachodzi $L_n1=1$, więc $\|L_n\|=1$.
Z jednostajnej ciągłości $f$ na $[0,1]$ wynika
$\|L_nf-f\|_\infty\to0$. Stosujemy
Twierdzenie~\ref{thm:sequence-bap28} lub bezpośrednio
Lemat~\ref{lem:finite-bap28} z $\lambda=1$.
\end{proof}

\begin{twierdzenie}[MAP bez bazy Schaudera]
\label{thm:linfty-map28}
Przestrzeń Banacha $\ell^\infty(\N,\K)$ ma MAP,
lecz nie ma żadnej bazy Schaudera.
\end{twierdzenie}
\begin{proof}
Ustalmy $x^{(1)},\ldots,x^{(m)}\in\ell^\infty$ oraz
$\varepsilon>0$. Zbiór
$A=\{(x^{(1)}_n,\ldots,x^{(m)}_n):n\in\N\}\subset\K^m$
jest ograniczony, a więc ma skończoną
$\varepsilon/2$-siatkę. Przypisujemy każdemu indeksowi $n$ pierwszy z wybranych
punktów siatki leżący w odległości mniejszej niż $\varepsilon/2$
od odpowiadającego mu wektora. Otrzymujemy podział $\N$
na skończenie wiele niepustych zbiorów $A_1,\ldots,A_r$ tak,
że dla $n,k\in A_j$ zachodzi
$|x_n^{(s)}-x_k^{(s)}|<\varepsilon$
dla każdego $1\le s\le m$ (na mocy nierówności trójkąta).
Wybierzmy po jednym $n_j\in A_j$ i połóżmy
$(Ty)_n=y_{n_j}$, jeżeli $n\in A_j$.
Jest to liniowy operator skończonego rzędu,
$\|T\|=1$ i
$\|Tx^{(s)}-x^{(s)}\|_\infty<\varepsilon$ dla każdego $s$.
Lemat~\ref{lem:finite-bap28} dowodzi MAP.

Przestrzeń $\ell^\infty$ nie jest ośrodkowa:
nieprzeliczalna rodzina ciągów $\chi_A=(\chi_A(n))_{n\ge1}$,
gdzie $A$ przebiega wszystkie podzbiory $\N$,
ma parami odległość $1$ w normie supremum.
Tymczasem każda przestrzeń Banacha z bazą Schaudera
jest ośrodkowa. Zatem $\ell^\infty$ nie ma takiej bazy.
\end{proof}

\begin{uwaga}[Brak uniwersalnego ciągu w przestrzeni nieośrodkowej]
\label{rem:linfty-sequence28}
Twierdzenia~\ref{thm:linfty-map28} oraz
\ref{thm:sequence-bap28} pokazują jednocześnie, że
w $\ell^\infty$ nie istnieje \emph{ciąg} operatorów
skończonego rzędu zbieżny silnie do $I$.
Istnieje natomiast jednostajnie ograniczona \emph{sieć}
takich operatorów zbieżna na wszystkich zbiorach zwartych.
Jest to powód, dla którego definicje AP i BAP nie powinny
opierać się wyłącznie na ciągach.
\end{uwaga}

\subsection{Aproksymowanie operatorów zwartych}

\begin{definicja}[Operatory zwarte i aproksymowalne]
\label{def:compact-approximable28}
Niech $X,Y$ będą przestrzeniami Banacha.
Operator $T\in\mathcal L(X,Y)$ jest \emph{zwarty},
jeżeli $T(B_X)$ ma zwarte domknięcie w $Y$, gdzie
$B_X=\{x\in X:\|x\|\le1\}$.
Oznaczamy zbiór takich operatorów przez
$\mathcal K(X,Y)$. Operator jest \emph{aproksymowalny},
jeżeli należy do domknięcia $\mathcal F(X,Y)$
w normie operatorowej; ich zbiór oznaczamy
$\mathcal A(X,Y)$.
\end{definicja}

\begin{twierdzenie}[Operatory aproksymowalne są zwarte]
\label{thm:approximable-compact28}
Dla wszystkich przestrzeni Banacha $X,Y$ zachodzi
$\mathcal F(X,Y)\subseteq\mathcal A(X,Y)\subseteq
\mathcal K(X,Y)\subseteq\mathcal L(X,Y)$.
Ponadto $\mathcal K(X,Y)$ jest domkniętą podprzestrzenią
liniową $\mathcal L(X,Y)$.
\end{twierdzenie}
\begin{proof}
Każdy operator skończonego rzędu jest zwarty, ponieważ
jego ograniczony obraz kuli jednostkowej ma zwarte
domknięcie w skończenie wymiarowej przestrzeni.
Jeżeli $T_n\to T$ w normie operatorowej i każdy $T_n$
jest zwarty, to dla $\varepsilon>0$ wybieramy $n$ z
$\|T-T_n\|<\varepsilon/3$. Zbiór $T_n(B_X)$ jest
całkowicie ograniczony, więc pokrywają go skończenie
wiele kul promienia $\varepsilon/3$.
Otrzymujemy skończone pokrycie $T(B_X)$ kulami
promienia $\varepsilon$, a zatem całkowitą ograniczoność.
Ponieważ $Y$ jest zupełna, $\overline{T(B_X)}$ jest zwarty.
To pokazuje domkniętość zbioru operatorów zwartych.
Liniowość wynika ze zwartości sum i skalarnych
wielokrotności zbiorów względnie zwartych.
\end{proof}

\begin{twierdzenie}[AP w przestrzeni wartości i aproksymacja operatorów]
\label{thm:ap-compact-operators28}
Niech $X,Y$ będą przestrzeniami Banacha, przy czym
$Y$ ma AP. Wówczas każdy operator zwarty
$T\in\mathcal K(X,Y)$ jest aproksymowalny, tzn.
$\mathcal K(X,Y)=\mathcal A(X,Y)$.
Jeżeli $Y$ ma $\lambda$-BAP, to dla każdego
$T\in\mathcal K(X,Y)$ i każdego $\varepsilon>0$
istnieje $F\in\mathcal F(X,Y)$ z
$\|F-T\|<\varepsilon$ oraz $\|F\|\le\lambda\|T\|$.
\end{twierdzenie}
\begin{proof}
Niech $K=\overline{T(B_X)}$; z założenia jest zwarty.
Z AP przestrzeni $Y$ otrzymujemy dla każdego
$\varepsilon>0$ operator skończonego rzędu
$S:Y\to Y$ taki, że
$\sup_{y\in K}\|Sy-y\|<\varepsilon$.
Operator $F=S\circ T$ ma skończony rząd i
$\|F-T\|=\sup_{x\in B_X}\|S(Tx)-Tx\|<\varepsilon$.
Przy $\lambda$-BAP można ponadto wybrać $\|S\|\le\lambda$,
stąd $\|F\|\le\lambda\|T\|$.
\end{proof}

\begin{uwaga}[Właściwy kierunek w Twierdzeniu~\ref{thm:ap-compact-operators28}]
\label{rem:range-not-domain28}
Założenie AP dotyczy tutaj \emph{przestrzeni wartości $Y$},
a nie automatycznie przestrzeni wyjściowej $X$:
aproksymujemy zwarty zbiór $T(B_X)\subseteq Y$.
Utożsamienie tych dwóch ról bez dodatkowego twierdzenia
o dualności byłoby nieuzasadnione.
\end{uwaga}

\subsection{Trwałość własności i ograniczenia twierdzeń}

\begin{twierdzenie}[Trwałość dla podprzestrzeni dopełnionych]
\label{thm:complemented28}
Niech $X$ będzie przestrzenią Banacha i $Y\subseteq X$
domkniętą podprzestrzenią. Załóżmy, że istnieje
ciągły rzut $Q:X\to Y$, tj. $Q|_Y=I_Y$.
Jeżeli $X$ ma AP, to $Y$ ma AP.
Jeżeli $X$ ma $\lambda$-BAP, to $Y$ ma
$\mu$-BAP dla $\mu=\max\{1,\lambda\|Q\|\}$.
W szczególności, gdy $X$ ma MAP i $\|Q\|=1$,
to $Y$ ma MAP.
\end{twierdzenie}
\begin{proof}
Ustalmy zwarty $K\subseteq Y$. Jest on również
zwarty w $X$. Wybierzmy operator skończonego rzędu
$T:X\to X$ przybliżający identyczność jednostajnie
na $K$ i połóżmy $S=Q\circ T|_Y:Y\to Y$.
Operator $S$ ma skończony rząd, a dla $y\in K$
$\|Sy-y\|=\|Q(Ty-y)\|\le\|Q\|\,\|Ty-y\|$.
Ponieważ dokładność aproksymacji $T$ możemy wybrać
dowolnie dużą, otrzymujemy AP na $Y$.
Przy $\lambda$-BAP zachodzi dodatkowo
$\|S\|\le\|Q\|\lambda$, co dowodzi drugiej części.
\end{proof}

\begin{twierdzenie}[Niezmienniczość AP i BAP przy izomorfizmach]
\label{thm:isomorphic28}
Niech $X,Y$ będą przestrzeniami Banacha i niech
$U:X\to Y$ będzie liniowym izomorfizmem topologicznym.
Wtedy $X$ ma AP wtedy i tylko wtedy, gdy $Y$ ma AP.
Jeżeli $X$ ma $\lambda$-BAP, to $Y$ ma
$\mu$-BAP z $\mu=\max\{1,\lambda\|U\|\|U^{-1}\|\}$.
W szczególności istnienie BAP jest niezmiennicze
względem równoważnych norm, natomiast liczba $\lambda$
może się zmienić.
\end{twierdzenie}
\begin{proof}
Dla zwartego $K\subseteq Y$ jego przeciwobraz
$U^{-1}(K)$ jest zwarty w $X$. Niech
$T\in\mathcal F(X,X)$ przybliża identyczność na
$U^{-1}(K)$. Operator $S=UTU^{-1}\in\mathcal F(Y,Y)$
spełnia $\|Sy-y\|\le\|U\|\,\|TU^{-1}y-U^{-1}y\|$.
W przypadku BAP mamy dodatkowo
$\|S\|\le\|U\|\|U^{-1}\|\lambda$.
Argument jest symetryczny względem $X,Y$.
\end{proof}

\begin{uwaga}[Granice implikacji: kontrprzykłady głębokie]
\label{rem:enflo-figiel28}
Nie każda ośrodkowa, a nawet nie każda \emph{refleksywna}
przestrzeń Banacha ma AP: istnieje taki kontrprzykład
Enflo. Istnieją też przestrzenie Banacha mające AP,
ale niemające BAP (konstrukcja Figiela--Johnsona),
oraz mające BAP, lecz nie MAP. Klasyczne źródła
dwóch pierwszych kontrprzykładów to prace Enflo
(\emph{Acta Mathematica} 130, 1973, s.~309--317) oraz
Figiela i Johnsona (\emph{Proceedings of the AMS} 41,
1973, s.~197--200). Wyniki te \emph{przytaczamy bez dowodu}, ponieważ szczegółowe
konstrukcje wykraczają poza zakres rozdziału.
W przeciwieństwie do tych rezultatów kontrprzykład
``MAP nie implikuje istnienia bazy Schaudera'' z
Twierdzenia~\ref{thm:linfty-map28} jest całkowicie
elementarny i został udowodniony powyżej.
\end{uwaga}

\begin{uwaga}[Co naprawdę daje istnienie bazy?]
\label{rem:recap28}
Zestawiając Twierdzenia~\ref{thm:basis-implies-bap28}
i~\ref{thm:linfty-map28}, otrzymujemy ściśle
jednokierunkową zależność między własnościami bazowymi
i aproksymacyjnymi. Z drugiej strony nie wolno
zastępować \emph{jednostajnej} aproksymacji na
zbiorach zwartych dowolną aproksymacją punktową przez
sieć operatorów skończonego rzędu: ta ostatnia zachodzi
w każdej przestrzeni Banacha na mocy
Twierdzenia~\ref{prop:universal-sot28}.
\end{uwaga}

% Wersja źródłowa: audyt II, jednostka 29
\setcounter{section}{28}
\section{Operatory ograniczone i zwarte: kryteria, transpozycja i ideały}
\label{sec:operatory29}
\noindent\textbf{Wyniki wykorzystywane.}
Korzystamy z normującej postaci twierdzenia Hahna--Banacha
(Wniosek~\ref{cor:norming}), zasady jednostajnej ograniczoności
(Wniosek~\ref{cor:uniform-boundedness}), twierdzenia o odwzorowaniu
otwartym (Twierdzenie~\ref{thm:open-mapping11}), twierdzenia
Eberleina--\v Smuliana (Twierdzenie~\ref{thm:eberlein-smulian26})
oraz Twierdzenia~\ref{thm:approximable-compact28}.
Korzystamy również z klasycznych twierdzeń Arzeli--Ascolego
i Fubiniego oraz regularności miary Lebesgue'a.
\par\medskip

W rozdziale przyjmujemy, że $X,Y,Z$ są przestrzeniami Banacha nad
$\K\in\{\R,\C\}$. Zapis $\mathcal L(X,Y)$ oznacza przestrzeń
ograniczonych operatorów liniowych, $\mathcal F(X,Y)$ -- operatory
ciągłe skończonego rzędu, a $\mathcal K(X,Y)$ -- operatory zwarte.
Te pojęcia, a także domkniętość $\mathcal K(X,Y)$ w normie operatorowej,
ustanowiono w rozdziale o własnościach aproksymacji. Nie powtarzamy
tamtych dowodów. Dla operatora $T:X\to Y$ przez $T':Y'\to X'$
oznaczamy jego operator sprzężony (transpozycję), dany wzorem
$(T'y')(x)=y'(Tx)$. Normy operatora i jego transpozycji są równe:
$\|T'\|=\|T\|$. Istotnie, z Wniosku~\ref{cor:norming}
i definicji normy operatorowej wynika
\[
\|T'\|=\sup_{\|y'\|\le1}\sup_{\|x\|\le1}|y'(Tx)|
=\sup_{\|x\|\le1}\|Tx\|=\|T\|.
\]

\subsection{Zwartość i kryteria ciągowe}

\begin{twierdzenie}[Trzy kryteria zwartości]
\label{thm:compact-criteria29}
Dla $T\in\mathcal L(X,Y)$ następujące warunki są równoważne:
\begin{enumerate}[label=\textup{(\roman*)}]
\item $T$ jest zwarty, tj. $\overline{T(B_X)}$ jest zwarty w $Y$;
\item dla każdego ograniczonego ciągu $(x_n)$ w $X$ ciąg $(Tx_n)$
      ma podciąg zbieżny w normie $Y$;
\item dla każdego $\varepsilon>0$ zbiór $T(B_X)$ ma skończoną
      $\varepsilon$-siatkę w $Y$.
\end{enumerate}
\end{twierdzenie}
\begin{proof}
(i) $\Rightarrow$ (ii): dowolny ograniczony ciąg $(x_n)$
mieści się w $rB_X$ dla pewnego $r>0$. Zbiór
$r\overline{T(B_X)}$ jest zwarty, więc ciąg $(Tx_n)$ ma
podciąg zbieżny w normie.
(ii) $\Rightarrow$ (iii): gdyby $T(B_X)$ nie miał skończonej
$\varepsilon$-siatki dla pewnego $\varepsilon>0$, wybieralibyśmy
indukcyjnie $x_n\in B_X$ tak, aby $Tx_{n+1}$ leżał poza sumą
kul otwartych promienia $\varepsilon$ wokół $Tx_1,\ldots,Tx_n$.
Wtedy $\|Tx_n-Tx_m\|\ge\varepsilon$ dla $n\ne m$,
więc żaden podciąg nie byłby ciągiem Cauchy'ego.
(iii) $\Rightarrow$ (i): warunek (iii) oznacza całkowitą
ograniczoność $T(B_X)$. Domknięcie tego zbioru jest również
całkowicie ograniczone: dla zadanego promienia wybieramy
siatkę o połowie tego promienia, a każdy punkt domknięcia
przybliżamy punktem zbioru z błędem mniejszym niż połowa
promienia. Domknięty podzbiór zupełnej przestrzeni $Y$ jest
zupełny. Z metrycznego kryterium zwartości otrzymujemy (i).
\end{proof}

\begin{lemat}[Lemat Riesza o prawie prostopadłym wektorze]
\label{lem:riesz29}
Niech $M$ będzie właściwą domkniętą podprzestrzenią przestrzeni
unormowanej $X$, a $0<\theta<1$. Istnieje $x\in X$ takie, że
$\|x\|=1$ oraz $\operatorname{dist}(x,M)>\theta$.
\end{lemat}
\begin{proof}
Wybieramy $z\in X\setminus M$. Ponieważ $M$ jest domknięta,
$d:=\operatorname{dist}(z,M)>0$. Wybieramy $y\in M$ z
$\|z-y\|<d/\theta$ i kładziemy $x=(z-y)/\|z-y\|$.
Dla każdego $m\in M$ zachodzi
$\|x-m\|=\|z-y-\|z-y\|m\|/\|z-y\|\ge d/\|z-y\|$,
ponieważ $y+\|z-y\|m\in M$. Biorąc infimum po $m\in M$,
dostajemy $\operatorname{dist}(x,M)\ge d/\|z-y\|>\theta$.
\end{proof}

\begin{twierdzenie}[Zwartość identyczności]
\label{thm:identity-compact29}
Dla przestrzeni Banacha $X$ operator $I_X$ jest zwarty wtedy
i tylko wtedy, gdy $X$ ma skończony wymiar.
\end{twierdzenie}
\begin{proof}
W skończonym wymiarze ograniczone zbiory mają zwarte domknięcie.
Gdy $X$ ma nieskończony wymiar, wybieramy indukcyjnie wektory
$x_n\in B_X$ o normie $1$ w taki sposób, że
$\operatorname{dist}(x_{n+1},\operatorname{span}\{x_1,\ldots,x_n\})>1/2$.
Lemat~\ref{lem:riesz29} pozwala wykonać każdy krok, gdyż
podprzestrzeń skończenie wymiarowa jest domknięta i właściwa.
Wówczas $\|x_n-x_m\|>1/2$ dla $n\ne m$, więc kula jednostkowa
nie jest względnie zwarta.
\end{proof}

\begin{twierdzenie}[Zwartość a składanie operatorów]
\label{thm:ideal-compact29}
Jeżeli $S\in\mathcal L(Y,Z)$ i $T\in\mathcal L(X,Y)$, to:
\begin{enumerate}[label=\textup{(\roman*)}]
\item $T\in\mathcal K(X,Y)$ implikuje
      $ST\in\mathcal K(X,Z)$;
\item $S\in\mathcal K(Y,Z)$ implikuje
      $ST\in\mathcal K(X,Z)$.
\end{enumerate}
Ponadto $\mathcal K(X,Y)$ jest podprzestrzenią liniową zamkniętą
w normie operatorowej.
\end{twierdzenie}
\begin{proof}
W (i) zbiór $T(B_X)$ jest względnie zwarty, a ciągłość $S$
zapewnia względną zwartość jego obrazu. W (ii) zbiór $T(B_X)$
jest ograniczony, więc zawiera się w $rB_Y$ dla pewnego $r>0$;
$S(rB_Y)$ jest względnie zwarty. Ostatnie zdanie jest treścią
Twierdzenia~\ref{thm:approximable-compact28}.
\end{proof}

\begin{przyklad}[Zwartość nie wymaga skończonego rzędu]
\label{ex:volterra29}
W przestrzeni Banacha $C([0,1],\K)$ definiujemy operator Volterry
$(Vf)(t)=\int_{0}^{t}f(s)\,\mathrm d\lambda_1(s)$.
Liniowość wynika z liniowości całki, a oszacowanie
$|(Vf)(t)|\le t\|f\|_\infty$ daje $\|V\|\le1$.
Wartości funkcji $Vf$ zależą od $t$ w sposób ciągły, gdyż
$|(Vf)(t)-(Vf)(u)|\le\|f\|_\infty|t-u|$.
Dla $\|f\|_\infty\le1$ mamy $\|Vf\|_\infty\le1$ i
$|(Vf)(t)-(Vf)(u)|\le|t-u|$. Zatem rodzina $V(B_X)$
jest jednostajnie ograniczona i jednakowo ciągła;
z twierdzenia Arzeli--Ascolego ma zwarte domknięcie.
Operator $V$ jest zatem zwarty. Jednocześnie obrazy wielomianów
$f_n(t)=t^n$, równe $t^{n+1}/(n+1)$, są liniowo niezależne,
więc rząd $V$ jest nieskończony. Ponadto z podstawowego twierdzenia rachunku całkowego
$(Vf)'=f$ na $(0,1)$, więc $Vf=0$ implikuje $f=0$
także na końcach odcinka przez ciągłość. Operator jest
różnowartościowy, ale jego obraz nie jest domknięty: gdyby był domknięty, odwrotność
$V^{-1}:V(X)\to X$ byłaby ciągła z twierdzenia o odwzorowaniu
otwartym, a stąd $I_X=V^{-1}V$ byłaby zwarta, wbrew
Twierdzeniu~\ref{thm:identity-compact29}.
\end{przyklad}

\subsection{Twierdzenie Schaudera o operatorze sprzężonym}

\begin{lemat}[Jednostajna aproksymacja ograniczonych funkcjonałów]
\label{lem:restrictions29}
Niech $K$ będzie zwartym podzbiorem przestrzeni Banacha $Y$.
Rodzina $\{y'|_K:\|y'\|\le1\}$ jest względnie zwarta
w $C(K,\K)$ z normą supremum.
\end{lemat}
\begin{proof}
Dla $k\in K$ mamy $|y'(k)|\le\|k\|$, więc rodzina jest
jednostajnie ograniczona. Ponadto
$|y'(k)-y'(l)|\le\|k-l\|$ dla dowolnych $k,l\in K$:
rodzina jest jednakowo ciągła. Twierdzenie Arzeli--Ascolego
(dla funkcji zespolonych stosujemy je do części rzeczywistych
i urojonych) daje względną zwartość w normie supremum.
\end{proof}

\begin{twierdzenie}[Schauder: zwartość transpozycji]
\label{thm:schauder-adjoint29}
Niech $X,Y$ będą przestrzeniami Banacha nad tym samym ciałem
$\K$, a $T\in\mathcal L(X,Y)$. Wówczas
\[
T\in\mathcal K(X,Y)\quad\Longleftrightarrow\quad
T'\in\mathcal K(Y',X').
\]
\end{twierdzenie}
\begin{proof}
Załóżmy najpierw, że $T$ jest zwarty, i niech
$K=\overline{T(B_X)}\subseteq Y$. Dla $y'_1,y'_2\in B_{Y'}$
z definicji transpozycji dostajemy
\[
\|T'y'_1-T'y'_2\|_{X'}
=\sup_{x\in B_X}|(y'_1-y'_2)(Tx)|
\le\sup_{k\in K}|(y'_1-y'_2)(k)|.
\]
Lemat~\ref{lem:restrictions29} pokazuje, że rodzina
restrykcji $y'|_K$ ma zwarte domknięcie w $C(K,\K)$.
Dla każdego $\varepsilon>0$ wybieramy skończoną
$\varepsilon/3$-siatkę tej rodziny. Z każdej kuli siatki,
która przecina rodzinę, wybieramy jedną restrykcję
$y'_j|_K$ dla $y'_j\in B_{Y'}$. Każda restrykcja jest
od wybranej restrykcji odległa o mniej niż $2\varepsilon/3$.
Z powyższego oszacowania wynika, że
$T'y'_1,\ldots,T'y'_m$ tworzą skończoną $\varepsilon$-siatkę
$T'(B_{Y'})$. Twierdzenie~\ref{thm:compact-criteria29}
daje zwartość $T'$.

Jeżeli natomiast $T'$ jest zwarty, to właśnie udowodniona
implikacja, zastosowana do $T':Y'\to X'$, daje zwartość
$T'':X''\to Y''$. Oznaczmy przez $J_X:X\to X''$ i
$J_Y:Y\to Y''$ kanoniczne izometryczne zanurzenia.
Dla każdego $x\in X$ i $y'\in Y'$ mamy
$(T''J_Xx)(y')=(J_Xx)(T'y')=(T'y')(x)=y'(Tx)=(J_YTx)(y')$,
co dowodzi tożsamości $T''J_Xx=J_YTx$. Ponieważ $J_X(B_X)\subseteq B_{X''}$,
zbiór $J_YT(B_X)$ jest względnie zwarty w $Y''$.
Obraz $J_Y(Y)$ jest domknięty (ponieważ $Y$ jest zupełna),
a $J_Y:Y\to J_Y(Y)$ jest izometrią. Stąd $T(B_X)$
jest względnie zwarty w $Y$, co dowodzi tezy.
\end{proof}

\begin{wniosek}[Sprzężenie zachowuje zwartość]
\label{cor:schauder-iterated29}
Dla każdego $T\in\mathcal L(X,Y)$ zwartość $T$ jest
równoważna zwartości $T'$ i $T''$; nie wymagamy
refleksywności żadnej z przestrzeni.
\end{wniosek}
\begin{proof}
Pierwsza równoważność jest treścią
Twierdzenia~\ref{thm:schauder-adjoint29}; druga wynika
z zastosowania tego twierdzenia do $T'$ oraz z tożsamości
$T''J_X=J_YT$.
\end{proof}

\subsection{Zbieżność słaba, operatory całkowicie ciągłe i słabo zwarte}

\begin{twierdzenie}[Operator zwarty na ciągu słabo zbieżnym]
\label{thm:weak-norm29}
Niech $T\in\mathcal K(X,Y)$. Jeżeli ciąg $(x_n)$ zbiega
słabo do $x\in X$, to $\|Tx_n-Tx\|\to0$.
To samo zachodzi dla \emph{ograniczonej} sieci słabo zbieżnej.
\end{twierdzenie}
\begin{proof}
Dla ciągu słabo zbieżnego twierdzenie Banacha--Steinhausa
zastosowane do obrazów $J_Xx_n\in X''$ zapewnia
$\sup_n\|x_n\|<\infty$. Jeżeli $\|Tx_n-Tx\|$ nie zbiega
do zera, można wybrać podciąg, którego odległość od $Tx$
jest ograniczona od dołu. Ze zwartości $T$ mamy podpodciąg
zbieżny w normie do pewnego $y\in Y$. Ponieważ $T$
jest ciągły względem topologii słabych (dla $y'\in Y'$
mamy $y'(Tx_n)=(T'y')(x_n)\to (T'y')(x)$),
granica słaba tego podpodciągu
musi być równa $Tx$, więc $y=Tx$ -- sprzeczność.
Dla ograniczonej sieci jej obraz jest względnie zwarty
w normie. Gdyby obrazy nie zbiegały do $Tx$, istniałaby
kula $B(Tx,\varepsilon)$, poza którą obrazy leżą
na współkońcowym zbiorze indeksów. Możemy wówczas
wybrać podsieć o wszystkich wyrazach poza tą kulą.
Ze zwartości można wybrać dalszą podsieć normowo zbieżną,
a jej granica słaba musiałaby być równa $Tx$.
Jednak dopełnienie kuli $B(Tx,\varepsilon)$ jest normowo
domknięte, więc granica normowa należy do tego dopełnienia.
Otrzymujemy sprzeczność.
\end{proof}

\begin{definicja}[Operator całkowicie ciągły i słabo zwarty]
\label{def:completely-continuous29}
Operator $T\in\mathcal L(X,Y)$ nazywamy
\emph{całkowicie ciągłym}, jeżeli odwzorowuje każdy ciąg
słabo zbieżny w $X$ na ciąg zbieżny w normie $Y$.
Operator $T$ jest \emph{słabo zwarty}, jeżeli
$T(B_X)$ ma słabo zwarte domknięcie w $Y$.
\end{definicja}

\begin{przyklad}[Implikacje są ostre]
\label{ex:weak-vs-compact29}
Każdy operator zwarty jest całkowicie ciągły na mocy
Twierdzenia~\ref{thm:weak-norm29} i słabo zwarty, ponieważ
normowo zwarte domknięcie obrazu kuli jest słabo zwarte
oraz słabo domknięte (topologia słaba jest Hausdorffa).
Jednak operator identyczności na $\ell^1$ jest całkowicie
ciągły na mocy własności Schura, choć nie jest zwarty,
ponieważ $\ell^1$ ma nieskończony wymiar. Nie jest też
słabo zwarty: ciąg $(e_n)$ nie ma podciągu słabo zbieżnego,
gdyż własność Schura (Twierdzenie~\ref{thm:schur-l1-26})
wymuszałaby zbieżność normową, a $\|e_n-e_m\|_1=2$
dla $n\ne m$. Stosujemy twierdzenie Eberleina--\v Smuliana.
Z drugiej strony naturalne zanurzenie
$J:\ell^1\to\ell^2$, $J(x)=x$, jest ograniczone
($\|x\|_2\le\|x\|_1$) i słabo zwarte, gdyż obraz kuli
leży w słabo zwartej kuli refleksywnej przestrzeni $\ell^2$.
Jest także całkowicie ciągłe dzięki własności Schura,
ale nie jest zwarte: $Je_n=e_n$ i $\|e_n-e_m\|_2=\sqrt2$
dla $n\ne m$. Natomiast $I_{\ell^2}$ jest słabo zwarty,
lecz nie jest całkowicie ciągły, ponieważ $e_n\rightharpoonup0$
(współczynniki każdego wektora z $\ell^2$ dążą do zera),
ale $\|e_n\|_2=1$.
\end{przyklad}

\begin{twierdzenie}[Kryterium w refleksywnej przestrzeni Banacha]
\label{thm:reflexive-complete-cont29}
Jeżeli $X$ jest refleksywną przestrzenią Banacha, to
$T\in\mathcal L(X,Y)$ jest zwarty wtedy i tylko wtedy,
gdy jest całkowicie ciągły.
\end{twierdzenie}
\begin{proof}
Jedna implikacja wynika z
Twierdzenia~\ref{thm:weak-norm29}. Odwrotnie,
niech $(x_n)$ będzie dowolnym ciągiem z $B_X$.
Refleksywność daje słabą zwartość $B_X$, a twierdzenie
Eberleina--\v Smuliana zapewnia podciąg $(x_{n_k})$
słabo zbieżny do pewnego $x\in B_X$.
Całkowita ciągłość $T$ daje normową zbieżność
$Tx_{n_k}\to Tx$. Stosujemy
Twierdzenie~\ref{thm:compact-criteria29}.
\end{proof}

\subsection{Aproksymacja operatorów zwartych i przykłady funkcyjne}

\begin{twierdzenie}[Kryterium ogonów w przestrzeniach ciągowych]
\label{thm:tails29}
Niech $Y=c_0$ lub $Y=\ell^p$ dla $1\le p<\infty$,
a $P_N:Y\to Y$ oznacza obcięcie ciągu po $N$ współrzędnych.
Dla każdego $T\in\mathcal L(X,Y)$ następujące warunki są
równoważne:
\begin{enumerate}[label=\textup{(\roman*)}]
\item $T$ jest zwarty;
\item $\lim_{N\to\infty}\|(I_Y-P_N)T\|=0$;
\item $\lim_{N\to\infty}\sup_{x\in B_X}
      \|(I_Y-P_N)Tx\|=0$.
\end{enumerate}
\end{twierdzenie}
\begin{proof}
Równoważność (ii) i (iii) jest konsekwencją definicji normy
operatorowej, gdyż dla każdego $N$ norma operatora
$(I_Y-P_N)T$ jest równa supremum w (iii).
Niech $T$ będzie zwarty i $K=\overline{T(B_X)}$.
Ponieważ $P_Ny\to y$ dla każdego $y\in Y$, a
$\|I_Y-P_N\|\le1$, zbieżność jest jednostajna na $K$:
wybieramy skończoną siatkę w $K$ i stosujemy oszacowanie
$\|(I-P_N)y\|\le\|(I-P_N)(y-y_j)\|+
\|(I-P_N)y_j\|$. Dla zadanej $\varepsilon>0$ wybieramy
siatkę promienia $\varepsilon/3$; dla dostatecznie dużych $N$
wszystkie skończenie liczne drugie składniki są mniejsze
niż $\varepsilon/3$. Supremum po $K$ jest więc nie większe
niż $2\varepsilon/3$. Stąd (ii).
W drugą stronę $P_NT$ jest operatorem skończonego rzędu,
a więc jest zwarty. Z założenia $P_NT\to T$ w normie
operatorowej, a domkniętość klasy operatorów zwartych
daje (i).
\end{proof}

\begin{przyklad}[Operatory diagonalne]
\label{ex:diagonal29}
Niech $Y=c_0$ albo $Y=\ell^p$, $1\le p<\infty$,
oraz $a=(a_n)\in\ell^\infty$. Operator
$D_a:Y\to Y$, $(D_ax)_n=a_nx_n$, jest ograniczony
i $\|D_a\|\le\|a\|_\infty$ z oszacowania współrzędnych.
Ponieważ $\|D_ae_n\|=|a_n|$ i $\|e_n\|=1$, nierówność
odwrotna wynika po wzięciu supremum po $n$. Twierdzenie~\ref{thm:tails29}
pokazuje, że $D_a$ jest zwarty dokładnie wtedy, gdy
$a_n\to0$, ponieważ
$\|(I-P_N)D_a\|=\sup_{n>N}|a_n|$.
Na przykład $D_a$ dla $a_n=1/n$ jest zwartym
operatorem różnowartościowym o niedomkniętym obrazie:
gdyby obraz był domknięty, twierdzenie o odwzorowaniu
otwartym dawałoby ciągłość odwrotności na obrazie, a zatem
$\|e_n\|\le C\|D_ae_n\|=C/n$, co jest niemożliwe.
\end{przyklad}

\begin{twierdzenie}[Operatory całkowe o jądrze kwadratowalnym]
\label{thm:hilbert-schmidt29}
Niech $k\in L^2([0,1]^2,\lambda_2)$ i niech
$H=L^2([0,1],\lambda_1)$ będzie przestrzenią Hilberta.
Wówczas wzór
\[
(T_kf)(t)=\int_{[0,1]}k(t,s)f(s)\,\mathrm d\lambda_1(s)
\]
definiuje dla każdego $f\in H$ funkcję prawie wszędzie
określoną, której klasa w $H$ nie zależy od reprezentantów
$k$ i $f$. Otrzymujemy operator liniowy ciągły $T_k:H\to H$. Jest on zwarty i spełnia
$\|T_k\|\le\|k\|_{L^2(\lambda_2)}$.
\end{twierdzenie}
\begin{proof}
Z twierdzenia Fubiniego dla $|k|^2$ wynika, że
$k(t,\cdot)\in L^2$ dla prawie każdego $t$. Dla $f\in H$
iloczyn $k(t,s)f(s)$ jest całkowalny na kwadracie dzięki
nierówności Cauchy'ego--Schwarza w $L^2(\lambda_2)$,
ponieważ funkcja $(t,s)\mapsto f(s)$ ma normę $\|f\|_2$.
Fubini zapewnia mierzalność całki względem $s$ oraz
niezależność jej klasy od zmian reprezentantów na zbiorach
miary zero. Liniowość wynika z liniowości całki.
Dla prawie każdego $t$ nierówność Cauchy'ego--Schwarza daje:
$|(T_kf)(t)|^2\le
(\int_{[0,1]}|k(t,s)|^2\,\mathrm d\lambda_1(s))\|f\|_2^2$.
Całkując po $t$ względem $\lambda_1$ otrzymujemy
$\|T_kf\|_2\le\|k\|_2\|f\|_2$.
Aby uzasadnić potrzebną aproksymację jądra, zauważmy, że
funkcje charakterystyczne prostokątów $A\times B$ mają postać
$\chi_A(t)\chi_B(s)$, a ich skończone kombinacje
liniowe są gęste w $L^2([0,1]^2,\lambda_2)$.
Wynika to z gęstości funkcji prostych oraz aproksymacji
zbiorów mierzalnych przez skończone sumy prostokątów
w pseudometryce $A,B\mapsto\lambda_2(A\triangle B)$
(identyfikującej zbiory różniące się o zbiór miary zero).
Istotnie, z regularności miary Lebesgue'a każdy mierzalny
zbiór $E$ w kwadracie można zawrzeć pomiędzy zbiorem
zwartym $C$ i względnie otwartym $O$ z
$\lambda_2(O\setminus C)<\delta$. Każdy punkt $C$ leży
w względnie otwartym prostokącie zawartym w $O$.
Ze zwartości wybieramy skończone pokrycie $C$ takimi
prostokątami; ich suma $R$ spełnia
$\lambda_2(E\triangle R)<\delta$.
Funkcję $\chi_R$ zapisujemy jako skończoną kombinację
funkcji charakterystycznych prostokątów przez zasadę
włączeń i wyłączeń. Ponieważ
$\|\chi_E-\chi_R\|_2^2=\lambda_2(E\triangle R)$,
gęstość funkcji prostych kończy uzasadnienie.
Wybieramy zatem jądra $k_n$ będące skończonymi sumami
$g_j(t)h_j(s)$ i zbiegające do $k$ w $L^2$.
Jeżeli $k_n(t,s)=\sum_{j=1}^{m_n}g_j(t)h_j(s)$,
gdzie $g_j,h_j\in L^2([0,1],\lambda_1)$, to
\[
(T_{k_n}f)(t)=\sum_{j=1}^{m_n}g_j(t)
\int_{[0,1]}h_j(s)f(s)\,\mathrm d\lambda_1(s).
\]
Każdy współczynnik jest dobrze określony na mocy nierówności
Cauchy'ego--Schwarza. Obraz $T_{k_n}$ zawiera się więc
w $\operatorname{span}\{g_1,\ldots,g_{m_n}\}$. Z oszacowania normy wynika
$\|T_k-T_{k_n}\|\le\|k-k_n\|_2\to0$.
Zatem $T_k$ jest zwarty jako granica normowa operatorów
skończonego rzędu.
\end{proof}

\begin{uwaga}[Nie wszystkie operatory zwarte są aproksymowalne]
\label{rem:nonapproximable29}
W Twierdzeniu~\ref{thm:ap-compact-operators28} wykazano,
że jeżeli przestrzeń wartości $Y$ ma AP, to $\mathcal K(X,Y)=\mathcal A(X,Y)$.
Bez założenia AP w przestrzeni wartości takie twierdzenie
nie zachodzi automatycznie. Dlatego aproksymacja przez
operator skończonego rzędu nie jest równoważną
\emph{definicją} zwartości w dowolnych przestrzeniach Banacha.
\end{uwaga}

\subsection{Ideały operatorowe i przygotowanie teorii Fredholma}

\begin{definicja}[Ideał operatorowy]
\label{def:operator-ideal29}
Przez \emph{ideał operatorowy} w klasie przestrzeni Banacha
rozumiemy przyporządkowanie każdej parze $(X,Y)$
podprzestrzeni liniowej $\mathcal I(X,Y)\subseteq\mathcal L(X,Y)$,
zawierającej wszystkie operatory skończonego rzędu,
takie, że dla wszystkich przestrzeni Banacha $W,Z$
i operatorów $R\in\mathcal L(W,X)$, $T\in\mathcal I(X,Y)$,
$S\in\mathcal L(Y,Z)$ zachodzi
$STR\in\mathcal I(W,Z)$.
\end{definicja}

\begin{twierdzenie}[Trzy ideały i domkniętość]
\label{thm:ideals29}
Rodziny $\mathcal F$, $\mathcal A$ i $\mathcal K$ są
ideałami operatorowymi. Przy tym $\mathcal A(X,Y)$
i $\mathcal K(X,Y)$ są domknięte w normie operatorowej,
a $\mathcal F(X,Y)$ nie jest domknięta, jeżeli
$X=Y=\ell^2$. Dla wszystkich $X,Y$ zachodzi
$\mathcal F(X,Y)\subseteq\mathcal A(X,Y)\subseteq
\mathcal K(X,Y)$.
\end{twierdzenie}
\begin{proof}
Suma dwóch operatorów skończonego rzędu ma obraz zawarty
w sumie ich skończenie wymiarowych obrazów, a mnożenie
przez skalar zachowuje skończony rząd. Zatem $\mathcal F(X,Y)$
jest podprzestrzenią liniową, podobnie jej normowe domknięcie
$\mathcal A(X,Y)$. Składanie operatora skończonego rzędu z operatorami
ograniczonymi z lewej i prawej strony zachowuje
skończoność wymiaru obrazu; dlatego $\mathcal F$ jest ideałem.
Jeżeli $T_n\to T$ w normie i każdy $T_n$ ma skończony
rząd, to $ST_nR\to STR$ w normie na mocy
$\|ST_nR-STR\|\le\|S\|\|T_n-T\|\|R\|$.
To daje własność ideału $\mathcal A$ i jego domkniętość
z definicji. Twierdzenie~\ref{thm:ideal-compact29}
daje własność ideału $\mathcal K$; jego domkniętość
udowodniono wcześniej. Operator diagonalny
$D_{(1/n)}:\ell^2\to\ell^2$ ma nieskończony rząd,
a jego obcięcia skończonego rzędu zbiegają do niego
w normie. To pokazuje, że $\mathcal F(\ell^2,\ell^2)$
nie jest domknięta. Zawierania między trzema ideałami
wynikają z definicji i zwartości operatorów
skończonego rzędu.
\end{proof}

\begin{definicja}[Norma istotna operatora]
\label{def:essential-norm29}
Dla $T\in\mathcal L(X,Y)$ definiujemy
$\|T\|_{\mathrm{ess}}=
\inf\{\|T-K\|:K\in\mathcal K(X,Y)\}$.
Jest to odległość operatora $T$ od domkniętej
podprzestrzeni operatorów zwartych; może wynosić zero
dla operatora niezerowego. Nie jest więc normą
na całej przestrzeni $\mathcal L(X,Y)$, ale wyznacza normę
na przestrzeni ilorazowej
$\mathcal L(X,Y)/\mathcal K(X,Y)$. Gdy $X=Y$,
przestrzeń ta jest dodatkowo algebrą ilorazową.
\end{definicja}

\begin{twierdzenie}[Norma istotna identyczności]
\label{thm:essential-identity29}
Jeżeli $X$ jest nieskończenie wymiarową przestrzenią
Banacha, to $\|I_X\|_{\mathrm{ess}}=1$.
\end{twierdzenie}
\begin{proof}
Wybór $K=0$ pokazuje, że odległość jest nie większa
niż $1$. Załóżmy, że istnieje $K\in\mathcal K(X,X)$
z $\|I_X-K\|<1$. Stosując Lemat~\ref{lem:neumann15}
do operatora $T=I_X-K$, otrzymujemy odwracalność
$I_X-T=K$. W szczególności
$I_X=K^{-1}K$ byłby zwarty na mocy
Twierdzenia~\ref{thm:ideal-compact29}, wbrew
Twierdzeniu~\ref{thm:identity-compact29}.
\end{proof}

\begin{lemat}[Pierwszy krok teorii Fredholma]
\label{lem:compact-eigenvalue29}
Niech $T\in\mathcal K(X,X)$ i $\lambda\in\K\setminus\{0\}$.
Jeśli operator $\lambda I_X-T$ jest różnowartościowy,
to istnieje $c>0$ takie, że
$\|(\lambda I_X-T)x\|\ge c\|x\|$ dla wszystkich $x\in X$.
W szczególności obraz $\operatorname{ran}(\lambda I_X-T)$
jest domknięty.
\end{lemat}
\begin{proof}
Gdyby nie było takiego $c$, istniałby ciąg $(x_n)$
o normach równych $1$, spełniający
$(\lambda I_X-T)x_n\to0$. Ze zwartości $T$ możemy
wybrać podciąg, dla którego $Tx_{n_k}\to y$ w normie.
Wówczas $\lambda x_{n_k}=Tx_{n_k}+
(\lambda I_X-T)x_{n_k}\to y$, więc $x_{n_k}\to x=y/\lambda$
i $\|x\|=1$. Przez ciągłość operatora otrzymujemy
$(\lambda I_X-T)x=0$, co przeczy jego różnowartościowości.
Jeżeli obrazy $(\lambda I_X-T)x_n$ tworzą ciąg Cauchy'ego,
to uzyskane oszacowanie implikuje, że $(x_n)$ jest
ciągiem Cauchy'ego. Zupełność $X$ i ciągłość operatora
pokazują, że obraz jest domknięty.
\end{proof}

\begin{uwaga}[Granica zakresu rozdziału]
\label{rem:fredholm-next29}
Ostatni lemat nie stwierdza jeszcze surjektywności
$\lambda I_X-T$ ani nie opisuje krotności wartości
własnych. W następnym rozdziale dowiedziemy twierdzenia
Riesza--Schaudera, alternatywy Fredholma i stabilności
indeksu przy zaburzeniach zwartych, nie powtarzając
kryteriów zwartości i twierdzenia o transpozycji.
\end{uwaga}

% Wersja źródłowa: audyt II, jednostka 30
\setcounter{section}{29}
\noindent\textbf{Wyniki wykorzystywane.} Przyjmujemy znane twierdzenie Hahna--Banacha,
lemat Riesza, twierdzenie o odwzorowaniu otwartym, własności operatorów
zwartych i twierdzenie Schaudera o zwartości transpozycji.
\par\medskip

\section{Teoria Fredholma i twierdzenie Riesza--Schaudera}
\label{sec:fredholm30}

W tym rozdziale $X,Y$ są przestrzeniami Banacha nad
$\K\in\{\R,\C\}$, a $\mathcal L(X,Y)$ i $\mathcal K(X,Y)$
oznaczają odpowiednio operatory liniowe ciągłe i zwarte.
Kryteria zwartości i twierdzenie Schaudera o zwartości transpozycji
zostały już udowodnione. Najpierw przedstawimy alternatywę
Fredholma w przestrzeni Hilberta, w wersji korzystającej z
przybliżeń przez rzuty ortogonalne, a następnie argument
działający w dowolnych przestrzeniach Banacha.

\subsection{Operatory Fredholma i ich indeks}

\begin{definicja}[Operator Fredholma]
\label{def:fredholm30}
Niech $X,Y$ będą przestrzeniami Banacha i $T\in\mathcal L(X,Y)$.
Operator $T$ nazywamy \emph{operatorem Fredholma}, jeżeli
$\operatorname{ran}T$ jest domknięty, $\dim\ker T<\infty$ i
$\operatorname{codim}\operatorname{ran}T:=\dim(Y/\operatorname{ran}T)<\infty$.
Jego \emph{indeks Fredholma} wynosi
\[\operatorname{ind}T=\dim\ker T-\operatorname{codim}\operatorname{ran}T.\]
Symbol $X'$ oznacza przestrzeń dualną topologiczną $X$.
\end{definicja}

\begin{lemat}[Uzupełniające podprzestrzenie]
\label{lem:fredholm-complements30}
Każda skończenie wymiarowa podprzestrzeń $N$ przestrzeni Banacha
$X$ jest domknięta i ma domknięte dopełnienie topologiczne
$X=N\oplus X_1$, z ciągłymi rzutami. Jeżeli $R\subseteq Y$
jest domkniętą podprzestrzenią Banacha o skończonym kowymiarze,
to $Y=R\oplus C$, gdzie $C$ ma skończony wymiar i oba rzuty
są ciągłe.
\end{lemat}
\begin{proof}
Dla bazy $u_1,\ldots,u_n$ przestrzeni $N$ funkcjonały
współrzędnych są ciągłe na $N$ i przedłużają się na $X$
przez twierdzenie Hahna--Banacha do $f_1,\ldots,f_n\in X'$.
Wówczas $P(x)=\sum_{j=1}^n f_j(x)u_j$ jest ciągłym rzutem
na $N$, a $X_1=\ker P$. W drugim przypadku wybieramy bazę
$[v_1],\ldots,[v_m]$ skończenie wymiarowego ilorazu $Y/R$
i definiujemy $C=\operatorname{span}\{v_1,\ldots,v_m\}$.
Odwzorowanie przypisujące $y\in Y$ jego składową w $C$
przez ciągłe funkcjonały współrzędnych $Y/R$ jest ciągłym
rzutem o jądrze $R$. Obie dekompozycje są topologiczne.
\end{proof}

\begin{twierdzenie}[Obraz domknięty i warunek zgodności]
\label{thm:fredholm-compatibility30}
Jeżeli $T\in\mathcal L(X,Y)$ ma domknięty obraz, to
\[\operatorname{ran}T=\{y\in Y:\ f(y)=0\ \text{dla każdego }f\in\ker T'\}.\]
Jeżeli ponadto $T$ jest operatorem Fredholma, to
$\dim\ker T'=\operatorname{codim}\operatorname{ran}T$.
Zatem $Tx=y$ ma rozwiązanie wtedy i tylko wtedy, gdy
$f(y)=0$ dla wszystkich $f\in\ker T'$, a zbiór
rozwiązań ma postać $x_0+\ker T$.
\end{twierdzenie}
\begin{proof}
Z definicji $T'$ mamy $\ker T'=(\operatorname{ran}T)^\perp$,
gdzie symbol $R^\perp=\{f\in Y':f|_R=0\}$ oznacza
anihilator. Gdy $R=\operatorname{ran}T$ jest domknięty,
twierdzenie Hahna--Banacha oddziela każdy $y\notin R$
od $R$ ciągłym funkcjonałem zerującym się na $R$.
Stąd $R=(R^\perp)_\perp$. Ponadto $R^\perp$ jest
izomorficzna z $(Y/R)'$, które przy skończonym wymiarze
$Y/R$ ma ten sam wymiar. Reszta wynika z liniowości.
\end{proof}

\subsection{Alternatywa Fredholma w przestrzeni Hilberta}

\begin{twierdzenie}[Alternatywa Fredholma dla operatora zwartego]
\label{thm:fredholm-hilbert30}
Niech $H$ będzie zespoloną przestrzenią Hilberta, a
$K\in\mathcal K(H,H)$. Wówczas $T=I_H-K$ ma domknięty
obraz oraz skończenie wymiarowe jądra $\ker T$ i $\ker T^*$,
przy czym $\dim\ker T=\dim\ker T^*$.
Dla $f\in H$ równanie $Tu=f$ ma rozwiązanie dokładnie wtedy,
gdy $f\perp\ker T^*$; wszystkie rozwiązania tworzą
klasę $u_0+\ker T$.
\end{twierdzenie}
\begin{proof}
\emph{Krok 1: skończony wymiar jądra.}
Gdyby $N=\ker T$ był nieskończenie wymiarowy, zawierałby
nieskończony układ ortonormalny $(e_n)$. Równość $Ke_n=e_n$
przeczyłaby zwartości $K$.

\emph{Krok 2: domkniętość obrazu.}
Niech $M=N^\perp$. Istnieje $c>0$ takie, że
$\|Tx\|\ge c\|x\|$ dla $x\in M$: w przeciwnym przypadku
istniałby ciąg $x_n\in M$ o normach $1$, dla którego
$Tx_n\to0$. Ze zwartości $K$ wybralibyśmy podciąg
$(x_{n_j})$ z $Kx_{n_j}\to z$, a wtedy
$x_{n_j}=Tx_{n_j}+Kx_{n_j}\to z\in M\cap N$, przy $\|z\|=1$.
To sprzeczność. Jeżeli $Tx_n\to y$, zastępujemy $x_n$
jego rzutem ortogonalnym na $M$; oszacowanie dolne
czyni ten ciąg ciągiem Cauchy'ego. Jego granica $x\in M$
spełnia $Tx=y$.

\emph{Krok 3: indeks zerowy.}
Zwartość $K(B_H)$ daje skończenie wymiarową podprzestrzeń
$W\subseteq H$, dla której
$\|K-P_WK\|<1$: wystarczy wziąć podprzestrzeń
rozpiętą na punktach odpowiednio dokładnej skończonej siatki.
Piszemy $F=P_WK$, $E=K-F$, $\|E\|<1$ i
$G=(I-E)^{-1}F$. Szereg Neumanna dowodzi, że $I-E$
jest odwracalny, a
\[I-K=(I-E)(I-G).\]
Operator $G$ ma skończony rząd. Niech $W'=\operatorname{ran}G$
i rozłóżmy $H=W'\oplus(W')^\perp$.
W tych współrzędnych
\[(I-G)(w,z)=((I_{W'}-G|_{W'})w-Gz,\ z).\]
Składowa druga identyczności pozwala eliminować $-Gz$
i redukuje obliczenie wymiaru jądra i kowymiaru obrazu
do skończenie wymiarowego operatora $I_{W'}-G|_{W'}$.
Na mocy twierdzenia o wymiarze jądra i obrazu obie liczby są
równe. Mnożenie z lewej przez $I-E$ nie zmienia
indeksu, więc $\operatorname{ind}T=0$.

\emph{Krok 4: rozwiązywalność.}
Ponieważ $\operatorname{ran}T$ jest domknięty,
$(\operatorname{ran}T)^\perp=\ker T^*$ daje
$\operatorname{ran}T=(\ker T^*)^\perp$. Z kroku 3 wynika
$\dim\ker T=\dim\ker T^*$, zaś opis rozwiązań
wynika z liniowości.
\end{proof}

\begin{uwaga}[Zakres powyższego dowodu]
\label{rem:fredholm-hilbert-scope30}
W dowodzie wykorzystano rzut ortogonalny na skończenie
wymiarową podprzestrzeń. Dowolna przestrzeń Banacha
nie musi mieć własności aproksymacji, dlatego nie wolno
przenosić tego argumentu bez zmian. Alternatywę Fredholma
w przestrzeniach Banacha udowodnimy bez takiej aproksymacji.
\end{uwaga}

\subsection{Twierdzenie Atkinsona i indeks w przestrzeniach Banacha}

\begin{twierdzenie}[Twierdzenie Atkinsona]
\label{thm:atkinson30}
Dla $T\in\mathcal L(X,Y)$ następujące warunki są równoważne:
\begin{enumerate}[label=\textup{(\roman*)}]
\item $T$ jest operatorem Fredholma;
\item istnieje $S\in\mathcal L(Y,X)$, dla którego
$ST-I_X\in\mathcal K(X,X)$ i $TS-I_Y\in\mathcal K(Y,Y)$.
\end{enumerate}
Gdy (i) zachodzi, obie różnice można wybrać skończonego rzędu.
\end{twierdzenie}
\begin{proof}
Załóżmy (i) i wybierzmy rozkłady $X=N\oplus X_1$,
$Y=R\oplus C$, gdzie $N=\ker T$, $R=\operatorname{ran}T$.
Ograniczenie $A=T|_{X_1}:X_1\to R$ jest bijekcją między
przestrzeniami Banacha, więc $A^{-1}$ jest ograniczony.
Dla $y=r+c$ definiujemy $S(y)=A^{-1}r$.
Wtedy $ST=I_X-P_N$ i $TS=I_Y-P_C$, gdzie oba rzuty
mają skończony rząd.

Odwrotnie niech $ST=I_X-K_1$ i $TS=I_Y-K_2$,
gdzie $K_1,K_2$ są zwarte. Na $N=\ker T$ operator
$K_1$ działa jak identyczność, więc $N$ jest skończenie
wymiarowe (inaczej lemat Riesza przeczyłby zwartości).
Wybierzmy dopełnienie $X=N\oplus X_1$.
Jeśli $T|_{X_1}$ nie byłby ograniczony z dołu, istniałby
ciąg jednostkowy $x_n\in X_1$ z $Tx_n\to0$.
Równość $x_n=STx_n+K_1x_n$ oraz zwartość $K_1$
dawałyby normowo zbieżny podciąg do niezerowego wektora
w $N\cap X_1$, co jest niemożliwe. Zatem obraz $T$
jest domknięty.
Z równości $S'T'=I_{Y'}-K_2'$ (otrzymanej przez
transpozycję $TS=I_Y-K_2$) wynika, że
$\ker T'\subseteq\ker(I_{Y'}-K_2')$.
Operator $K_2'$ jest zwarty na mocy twierdzenia Schaudera,
więc drugie jądro ma skończony wymiar.
Twierdzenie~\ref{thm:fredholm-compatibility30} daje
$\operatorname{codim}\operatorname{ran}T=\dim\ker T'<\infty$.
\end{proof}

\begin{twierdzenie}[Lokalna stałość indeksu]
\label{thm:index-local30}
Niech $T\in\mathcal L(X,Y)$ będzie operatorem Fredholma.
Istnieje $\delta>0$ takie, że każdy $U\in\mathcal L(X,Y)$
z $\|U-T\|<\delta$ jest Fredholma i
$\operatorname{ind}U=\operatorname{ind}T$.
\end{twierdzenie}
\begin{proof}
Wybierzmy rozkłady $X=N\oplus X_1$, $Y=R\oplus C$
z dowodu twierdzenia Atkinsona. W ich współrzędnych
$T(n,x)=(Ax,0)$, gdzie $A:X_1\to R$ jest izomorfizmem.
Dla $U$ bliskiego $T$ mamy
\[U(n,x)=(Bn+Dx,\ En+Fx),\]
gdzie $B:N\to R$, $D:X_1\to R$, $E:N\to C$,
$F:X_1\to C$. Operator $D$ jest odwracalny, gdy
$\|A^{-1}(D-A)\|<1$, co zachodzi dla $U$ dostatecznie
bliskiego $T$.
Dokonujemy dwóch \emph{odwracalnych} zmian współrzędnych.
Najpierw w dziedzinie podstawiamy
$(n,x)=(n,x'-D^{-1}Bn)$. Otrzymujemy
\[U(n,x)=(Dx',\ (E-FD^{-1}B)n+Fx').\]
Następnie w przeciwdziedzinie stosujemy
$(r,c)\mapsto(r,c-FD^{-1}r)$. Wynik ma postać
\[(n,x')\mapsto(Dx',\ (E-FD^{-1}B)n).\]
Obraz jest domknięty, ponieważ $D$ jest izomorfizmem,
a $E-FD^{-1}B:N\to C$ działa między przestrzeniami
skończenie wymiarowymi. Indeks drugiego bloku to
$\dim N-\dim C$ niezależnie od jego rzędu, gdyż
\[\dim\ker M-\dim(C/\operatorname{ran}M)=\dim N-\dim C\]
dla każdego $M:N\to C$ w skończonym wymiarze.
Stąd $\operatorname{ind}U=\dim N-\dim C=\operatorname{ind}T$.
\end{proof}

\begin{twierdzenie}[Zaburzenia zwarte zachowują indeks]
\label{thm:index-compact-perturb30}
Niech $T\in\mathcal L(X,Y)$ będzie operatorem Fredholma,
a $K\in\mathcal K(X,Y)$. Wówczas $T+K$ jest Fredholma i
$\operatorname{ind}(T+K)=\operatorname{ind}T$.
\end{twierdzenie}
\begin{proof}
Z twierdzenia Atkinsona istnieje $S:Y\to X$ takie, że
$ST-I_X$ i $TS-I_Y$ są zwarte. Dla każdego $t\in[0,1]$
obie różnice $S(T+tK)-I_X$ i $(T+tK)S-I_Y$ są zwarte.
Zatem $T+tK$ jest Fredholma dla wszystkich $t$.
Jego indeks zależy lokalnie stale od $t$ na mocy
Twierdzenia~\ref{thm:index-local30}; ponieważ $[0,1]$
jest spójny, jest stały.
\end{proof}

\begin{wniosek}[Alternatywa Fredholma w przestrzeni Banacha]
\label{cor:fredholm-banach30}
Niech $X$ będzie przestrzenią Banacha i
$K\in\mathcal K(X,X)$. Wówczas $I_X-K$ jest operatorem
Fredholma indeksu zero. Równanie $(I_X-K)x=y$ jest
rozwiązywalne dokładnie wtedy, gdy $f(y)=0$ dla każdego
$f\in\ker(I_{X'}-K')$. Jeśli $\ker(I_X-K)=\{0\}$,
to $I_X-K$ jest izomorfizmem liniowo-topologicznym.
\end{wniosek}
\begin{proof}
Operator $I_X$ jest Fredholma indeksu zero; zaburzenie
$-K$ zachowuje indeks. Warunek zgodności wynika z
Twierdzenia~\ref{thm:fredholm-compatibility30}. Gdy jądro
jest zerowe, indeks zero i domkniętość obrazu dają
surjektywność. Ciągłość odwrotności wynika z twierdzenia
o odwzorowaniu otwartym.
\end{proof}

\begin{przyklad}[Indeks może być niezerowy]
\label{ex:shift-index30}
W $X=\ell^2(\N)$ przesunięcie w prawo
$R(x_1,x_2,\ldots)=(0,x_1,x_2,\ldots)$ ma zerowe jądro
i obraz kowymiaru jeden, a więc indeks $-1$.
Przesunięcie w lewo $L(x_1,x_2,\ldots)=(x_2,x_3,\ldots)$
ma indeks $1$. Nie każdy operator Fredholma ma więc
indeks zero.
\end{przyklad}

\subsection{Twierdzenie Riesza--Schaudera i widmo operatora zwartego}

\begin{definicja}[Widmo i rezolwenta]
\label{def:spectrum30}
Niech $X$ będzie zespoloną przestrzenią Banacha i
$K\in\mathcal L(X,X)$. Definiujemy
$\rho(K)=\{\lambda\in\C:\lambda I_X-K\text{ jest odwracalny}\}$
oraz $\sigma(K)=\C\setminus\rho(K)$.
Liczbę $\lambda$ nazywamy wartością własną, gdy
$\ker(\lambda I_X-K)\ne\{0\}$.
\end{definicja}

\begin{twierdzenie}[Riesz--Schauder]
\label{thm:riesz-schauder30}
Niech $X$ będzie nieskończenie wymiarową zespoloną
przestrzenią Banacha, a $K\in\mathcal K(X,X)$.
Wtedy:
\begin{enumerate}[label=\textup{(\roman*)}]
\item $0\in\sigma(K)$, a $\sigma(K)$ jest niepustym
zwartym podzbiorem $\C$;
\item każdy $\lambda\in\sigma(K)\setminus\{0\}$ jest
wartością własną o skończenie wymiarowym jądrze;
\item dla każdego $\varepsilon>0$ zbiór
$\{\lambda\in\sigma(K):|\lambda|\ge\varepsilon\}$
jest skończony;
\item dla $\lambda\ne0$ uogólniona przestrzeń własna
$\bigcup_{m\ge1}\ker(\lambda I_X-K)^m$ jest
skończenie wymiarowa i równa
$\ker(\lambda I_X-K)^{m_0}$ dla pewnego $m_0\ge1$.
\end{enumerate}
\end{twierdzenie}
\begin{proof}
Dla $|\lambda|>\|K\|$ szereg Neumanna dowodzi
odwracalności $\lambda I_X-K$. Zbiór rezolwenty jest otwarty:
jeśli $\lambda_0 I_X-K$ ma odwrotność, to dla małego
$|\lambda-\lambda_0|$ stosujemy szereg Neumanna do
$I+(\lambda-\lambda_0)(\lambda_0 I_X-K)^{-1}$.
Widmo jest więc domknięte i ograniczone. Ponieważ
$K$ nie może być odwracalny (inaczej
$I_X=K^{-1}K$ byłby zwarty), $0\in\sigma(K)$.

Dla $\lambda\ne0$ operator $\lambda I_X-K$ jest Fredholma
indeksu zero jako niezerowa wielokrotność $I_X-\lambda^{-1}K$.
Jeżeli nie jest odwracalny, nie może być różnowartościowy:
zerowe jądro i indeks zero wymuszałyby surjektywność.
Jądro jest skończenie wymiarowe. To dowodzi (ii).

Gdyby (iii) zawiodło, wybralibyśmy parami różne
niezerowe wartości własne $\lambda_n$ o
$|\lambda_n|\ge\varepsilon$ i odpowiadające im wektory
własne $v_n\ne0$. Wektory te są liniowo niezależne.
Niech $M_n=\operatorname{span}\{v_1,\ldots,v_n\}$.
Podprzestrzenie $M_n$ są niezmiennicze względem $K$.
Lemat Riesza dostarcza $x_n\in M_n$ o normie $1$ i
$\operatorname{dist}(x_n,M_{n-1})>1/2$ dla $n\ge2$.
Na ilorazie $M_n/M_{n-1}$ operator $K$ działa przez
mnożenie przez $\lambda_n$, więc
$Kx_n-\lambda_nx_n\in M_{n-1}$. Jeśli $m<n$,
to $Kx_m\in M_{n-1}$, a zatem
\[\|Kx_n-Kx_m\|\ge
|\lambda_n|\operatorname{dist}(x_n,M_{n-1})>\varepsilon/2.\]
To przeczy zwartości $K$.

Dla (iv) połóżmy $N_m=\ker(\lambda I_X-K)^m$.
Wielomian $\lambda^{-m}(\lambda I_X-K)^m$ ma postać
$I_X-Q_m$ z $Q_m$ zwartym (każdy niezerowy składnik
w $Q_m$ zawiera dodatnią potęgę $K$). Na mocy
alternatywy Fredholma każde $N_m$ ma skończony wymiar.
Gdyby rosnący ciąg $N_m$ nie stabilizował się, można
byłoby wybrać rosnące indeksy $n_j$ z
$N_{n_j-1}\subsetneq N_{n_j}$ oraz wektory
$x_j\in N_{n_j}$ o normie $1$ takie, że
$\operatorname{dist}(x_j,N_{n_j-1})>1/2$.
Dla $i<j$ zachodzą $Kx_i\in N_{n_j-1}$ i
$Kx_j-\lambda x_j\in N_{n_j-1}$; stąd
$\|Kx_j-Kx_i\|>|\lambda|/2$, co przeczy zwartości.
Zatem ciąg jąder stabilizuje się.
\end{proof}

\begin{przyklad}[Widmo operatora diagonalnego]
\label{ex:diagonal-spectrum30}
W $\ell^2(\N)$ operator
$K(x_1,x_2,\ldots)=(x_1,x_2/2,x_3/3,\ldots)$
jest zwarty jako granica operatorów skończonego rzędu.
Ma widmo $\sigma(K)=\{0\}\cup\{1/n:n\ge1\}$;
każda niezerowa wartość własna jest prosta.
\end{przyklad}

\begin{przyklad}[Operator Volterry: widmo bez niezerowych wartości własnych]
\label{ex:volterra-spectrum30}
W $X=C([0,1],\C)$ definiujemy
$(Vf)(t)=\int_{[0,t]}f(s)\,\mathrm d\lambda_1(s)$.
Operator $V$ jest zwarty. Iterując wzór na całkę,
dostajemy $\|V^n\|\le 1/n!$, ponieważ jądro $V^n$
jest równe $(t-s)^{n-1}/(n-1)!$ dla $0\le s\le t$.
Dla $\lambda\ne0$ szereg
$\lambda^{-1}\sum_{n\ge0}(V/\lambda)^n$
jest zbieżny w normie operatorowej i określa
$(\lambda I_X-V)^{-1}$. Zatem $\sigma(V)=\{0\}$,
mimo że $V\ne0$.
\end{przyklad}

\subsection{Zastosowanie w teorii równań różniczkowych cząstkowych}

\begin{uwaga}[Od operatora eliptycznego do operatora zwartego]
\label{rem:elliptic-fredholm30}
W skrypcie \emph{Równania różniczkowe cząstkowe, tom I}
rozpatruje się faktoryzację $L=L_\mu(I_V-K)$,
w której $L_\mu:V\to V^*$ jest izomorfizmem,
a $K:V\to V$ operatorem zwartym otrzymanym dzięki
zwartemu zanurzeniu Sobolewa. Wówczas $\operatorname{ind}L=0$,
a warunek rozwiązywalności wynika z anihilatora obrazu.
Wyprowadzenie odpowiednich nierówności koercywności
i zagadnienia brzegowe pozostają w skrypcie PDE.
\end{uwaga}

% Wersja źródłowa: audyt II, jednostka 31
\setcounter{section}{30}
\section{Algebry Banacha, widmo i holomorficzny rachunek funkcyjny}
\label{sec:banach-algebras31}

\providecommand{\atlasref}[2]{#2}

\noindent\textbf{Wyniki wykorzystywane.}
Korzystamy z twierdzeń Hahna--Banacha, Banacha--Alaoglu,
lematu Neumanna (Lemat~\ref{lem:neumann15}) oraz podstawowych
twierdzeń skalarnej analizy zespolonej (Cauchy'ego, Liouville'a,
o usuwaniu osobliwości i zasadzie maksimum). Twierdzenie
Stone'a--Weierstrassa udowodnimy w tym rozdziale.
\par\medskip

W poprzednim rozdziale wprowadziliśmy widmo ograniczonego operatora
na zespolonej przestrzeni Banacha, przede wszystkim w związku z operatorami
zwartymi. Teraz rozpatrzymy widmo elementu \emph{dowolnej} zespolonej algebry
Banacha. Otrzymamy w ten sposób wspólne podstawy teorii operatorów, algebr
funkcyjnych oraz algebr $C^*$. Nie powtarzamy twierdzenia Riesza--Schaudera:
jego twierdzenia będą przykładami ogólnego rachunku spektralnego.

\subsection{Algebra Banacha i elementy odwracalne}

\noindent\emph{Powiązania tematyczne w Bibliotece matematycznej:} \atlasref{LA-03::prop:lin3-end-algebra}{LA-03: algebra endomorfizmów}.


\begin{definicja}[Zespolona algebra Banacha]
\label{def:banach-algebra31}
\emph{Zespoloną algebrą Banacha} nazywamy zespoloną przestrzeń Banacha
$A$ z łącznym, dwuliniowym mnożeniem $A\times A\to A$ spełniającym
$\|ab\|\le\|a\|\|b\|$ dla wszystkich $a,b\in A$.
W tym rozdziale, o ile nie zaznaczono inaczej, zakładamy, że istnieje
element jednostkowy $1_A$ i $\|1_A\|=1$. Algebra nie musi być przemienna.
Element $a\in A$ nazywamy odwracalnym, gdy istnieje $b\in A$ takie,
że $ab=ba=1_A$; oznaczamy wtedy $b=a^{-1}$.
\end{definicja}

\begin{przyklad}[Trzy podstawowe algebry]
\label{ex:banach-algebras31}
Jeśli $X$ jest niezerową zespoloną przestrzenią Banacha, to $\mathcal L(X,X)$
z normą operatorową i składaniem jest algebrą Banacha.
Jeśli $K$ jest niepustą zwartą przestrzenią Hausdorffa, to $C(K,\C)$
z normą supremum i mnożeniem punktowym jest przemienną algebrą Banacha.
Przestrzeń $\ell^1(\N_0)$ z iloczynem splotowym
$(a*b)_n=\sum_{j=0}^n a_jb_{n-j}$ jest przemienną algebrą Banacha,
której jednostką jest $(1,0,0,\ldots)$.
W ostatnim przypadku
\[
\sum_{n\ge0}|(a*b)_n|\le
\sum_{n\ge0}\sum_{j=0}^n|a_j||b_{n-j}|
=\sum_{j,k\ge0}|a_j||b_k|=\|a\|_1\|b\|_1.
\]
Zmiana kolejności w sumach skończonych dla każdej
współrzędnej daje łączność splotu, a zamiana $j$ na
$n-j$ daje przemienność. Zupełność wynika z zupełności
przestrzeni $\ell^1$.
\end{przyklad}

\begin{lemat}[Szereg Neumanna i otwartość grupy odwracalnych]
\label{lem:neumann31}
Niech $A$ będzie zespoloną algebrą Banacha z jedynką
$1_A$ taką, że $\|1_A\|=1$.
Jeżeli $u\in A$ i $\|u\|<1$, to $1_A-u$ jest odwracalny i
$(1_A-u)^{-1}=\sum_{n=0}^\infty u^n$ (zbieżność w normie).
Zbiór $A^\times$ wszystkich elementów odwracalnych jest otwarty,
a odwzorowanie $a\mapsto a^{-1}$ jest na nim ciągłe.
\end{lemat}
\begin{proof}
Wersja operatorowa szeregu Neumanna została już udowodniona
w Lemacie~\ref{lem:neumann15}. Sprowadzimy do niej przypadek
dowolnej algebry Banacha z jedynką. Dla $v\in A$ określmy
operator lewego mnożenia $L_v:A\to A$, $L_v(x)=vx$.
Z podmultiplikatywności normy wynika $\|L_v\|\le\|v\|$,
a z łączności mnożenia $L_vL_w=L_{vw}$ i $L_{1_A}=I_A$,
gdzie $I_A$ jest operatorem identycznościowym na przestrzeni
Banacha $A$. Ponadto $L_v(1_A)=v$, więc odwzorowanie
$v\mapsto L_v$ jest iniektywne.

Jeżeli $\|u\|<1$, to $\|L_u\|<1$. Z Lematu~\ref{lem:neumann15}
otrzymujemy $B=(I_A-L_u)^{-1}=\sum_{n=0}^\infty L_u^n$
w normie operatorowej. Ewaluacja operatora na $1_A$ jest
ciągła, ponieważ $\|T(1_A)\|\le\|T\|\|1_A\|$.
Stąd $v=B(1_A)=\sum_{n=0}^\infty u^n$ w normie $A$.
Ciągłość $w\mapsto L_w$ i równość $L_{u^n}=L_u^n$
dają $L_v=B$. Wobec tego
$L_{(1_A-u)v}=I_A=L_{v(1_A-u)}$, a iniektywność
reprezentacji daje $(1_A-u)v=v(1_A-u)=1_A$.
Jest więc $v=(1_A-u)^{-1}$, jak należało wykazać.

Dla $a\in A^\times$ i $b\in A$ spełniającego
$\|a^{-1}(b-a)\|<1$ rozkład $b=a(1_A+a^{-1}(b-a))$
oraz udowodniona właśnie wersja szeregu dają odwracalność
$b$. To dowodzi otwartości $A^\times$, ponieważ wystarczy
$\|b-a\|<1/\|a^{-1}\|$. Jeśli
$\|a^{-1}(b-a)\|\le1/2$, oszacowanie szeregu daje
$\|b^{-1}\|\le2\|a^{-1}\|$ (przy $\|1_A\|=1$).
Tożsamość $b^{-1}-a^{-1}=b^{-1}(a-b)a^{-1}$ prowadzi
więc do $\|b^{-1}-a^{-1}\|\le2\|a^{-1}\|^2\|b-a\|$,
co dowodzi ciągłości odwracania.
\end{proof}

\subsection{Widmo i rezolwenta}

\begin{definicja}[Widmo elementu algebry]
\label{def:spectrum31}
Dla $a\in A$ definiujemy \emph{zbiór rezolwenty}
$\rho_A(a)=\{z\in\C:z1_A-a\in A^\times\}$ oraz
\emph{widmo} $\sigma_A(a)=\C\setminus\rho_A(a)$.
Na $\rho_A(a)$ definiujemy \emph{rezolwentę}
$R_a(z)=(z1_A-a)^{-1}$.
Indeks dolny $A$ pomijamy, gdy algebra jest ustalona. W przypadku
$A=\mathcal L(X,X)$ jest to definicja z poprzedniego rozdziału.
\end{definicja}

\begin{twierdzenie}[Zwartość i niepustość widma]
\label{thm:spectrum-nonempty31}
Niech $A$ będzie zespoloną algebrą Banacha z jedynką i $a\in A$.
Wtedy $\sigma_A(a)$ jest niepustym zwartym podzbiorem
$\{z\in\C:|z|\le\|a\|\}$. Rezolwenta $R_a$ jest holomorficzna
na $\rho_A(a)$ jako funkcja o wartościach w przestrzeni Banacha $A$.
Dla $z,w\in\rho_A(a)$ zachodzi tożsamość rezolwenty
\[
 R_a(z)-R_a(w)=(w-z)R_a(z)R_a(w).
\]
\end{twierdzenie}
\begin{proof}
Dla $|z|>\|a\|$ lemat~\ref{lem:neumann31} daje
$R_a(z)=\sum_{n\ge0}a^n z^{-n-1}$. Zatem widmo jest ograniczone,
a jego dopełnienie otwarte jako przeciwobraz $A^\times$
przez ciągłe odwzorowanie $z\mapsto z1_A-a$.
Dla $z_0\in\rho(a)$ mamy faktoryzację
$z1_A-a=(z_0 1_A-a)(1_A+(z-z_0)R_a(z_0))$.
Stąd dla $|z-z_0|\|R_a(z_0)\|<1$ otrzymujemy
\[
R_a(z)=\sum_{n\ge0}(-1)^n(z-z_0)^n R_a(z_0)^{n+1}.
\]
Szereg jest jednostajnie zbieżny w normie na każdym
mniejszym domkniętym dysku i dowodzi holomorficzności.
Dla tożsamości rezolwenty obliczamy bez założenia przemienności
algebry:
\[
R_a(z)[(w1_A-a)-(z1_A-a)]R_a(w)
=R_a(z)-R_a(w).
\]
Lewa strona jest równa $(w-z)R_a(z)R_a(w)$.

Gdyby widmo było puste, $R_a$ byłaby całkowitą funkcją holomorficzną
$\C\to A$. Dla $|z|>2\|a\|$ mielibyśmy
$\|R_a(z)\|\le 2/|z|$. Na pozostałym domkniętym dysku rezolwenta jest ograniczona
z ciągłości (możemy użyć dysku o promieniu $1+2\|a\|$).
Jest więc ograniczona na całej płaszczyźnie.
Dla dowolnego $\varphi\in A'$ funkcja skalarna
$\varphi\circ R_a$ jest całkowita i ograniczona, a z twierdzenia
Liouville'a stała. Granica przy $|z|\to\infty$ daje
$\varphi(R_a(z))=0$. Funkcjonały ciągłe rozdzielają punkty $A$,
więc $R_a(z)=0$, co przeczy $(z1_A-a)R_a(z)=1_A$.
\end{proof}

\begin{przyklad}[Widmo operatora mnożenia]
\label{ex:multiplication-spectrum31}
Dla $f\in C(K,\C)$ zachodzi $\sigma_{C(K)}(f)=f(K)$.
Istotnie, $z1-f$ ma odwrotność w $C(K)$ dokładnie wtedy,
gdy nigdzie nie znika; odwrotnością jest wówczas $1/(z-f)$.
W szczególności $\sigma(f)$ może być dowolnym niepustym zwartym
podzbiorem $\C$ (weźmy $K$ równy temu zbiorowi i $f(z)=z$).
\end{przyklad}

\begin{uwaga}[Widmo zależy od algebry]
Jeżeli $B\subseteq A$ jest domkniętą podalgebrą z tą samą jednostką,
to $\sigma_A(b)\subseteq\sigma_B(b)$ dla $b\in B$:
odwrotność w $B$ jest również odwrotnością w $A$.
Równość nie musi zachodzić. Niech
$\mathbb D=\{z\in\C:|z|<1\}$ i
$\mathbb T=\{z\in\C:|z|=1\}$.
Rozważmy algebrę dyskową
$\mathcal A(\overline{\mathbb D})$, tj. funkcje ciągłe
na $\overline{\mathbb D}$ i holomorficzne w $\mathbb D$.
Ograniczanie do brzegu daje izometryczną podalgebrę
$B\subseteq C(\mathbb T)$: izometrię zapewnia zasada maksimum,
a domkniętość wynika z zupełności algebry dyskowej:
ciąg Cauchy'ego w normie supremum na brzegu jest, na mocy
zasady maksimum, ciągiem Cauchy'ego na całym domkniętym
dysku. Jego granica jest ciągła i holomorficzna wewnątrz
dysku dzięki zbieżności jednostajnej na zbiorach zwartych.
Dla funkcji współrzędnej $u(z)=z$ na $\mathbb T$ zachodzi
$\sigma_{C(\mathbb T)}(u)=\mathbb T$.
Jednak $\sigma_B(u)=\overline{\mathbb D}$: dla
$|\lambda|>1$ funkcja $1/(u-\lambda)$ jest w $B$,
a dla $|\lambda|<1$ nie mogłaby być odwrotnością
$u-\lambda$ w $B$, gdyż jej holomorficzne przedłużenie
musiałoby spełniać $(z-\lambda)g(z)=1$ w dysku:
funkcja $(z-\lambda)g(z)-1$ znika na brzegu, więc
na mocy zasady maksimum znika w całym dysku.
Podstawienie $z=\lambda$ daje sprzeczność.
Dla $|\lambda|=1$ odwrotność nie istnieje już w
$C(\mathbb T)$.
\end{uwaga}

\begin{twierdzenie}[Wzór na promień spektralny]
\label{thm:spectral-radius31}
Niech $A$ będzie zespoloną algebrą Banacha z jedynką, $a\in A$.
Wówczas granica $\lim_{n\to\infty}\|a^n\|^{1/n}$ istnieje
i jest równa promieniowi spektralnemu
\[
 r(a):=\max\{|z|:z\in\sigma_A(a)\}
 =\lim_{n\to\infty}\|a^n\|^{1/n}
 =\inf_{n\ge1}\|a^n\|^{1/n}.
\]
\end{twierdzenie}
\begin{proof}
Nierówność $\|a^{n+m}\|\le\|a^n\|\|a^m\|$ daje istnienie granicy
$L$ i równość z infimum. Rzeczywiście, jeśli $a^m=0$ dla
pewnego $m$, teza o granicy jest natychmiastowa.
W przeciwnym razie, ustalając $m\ge1$ i zapisując $n=qm+r$,
$0\le r<m$, otrzymujemy
$\|a^n\|\le\|a^m\|^q\max_{0\le j<m}\|a^j\|$.
Po wyciągnięciu pierwiastka stopnia $n$ i przejściu do granicy
z góry dostajemy
$\limsup_n\|a^n\|^{1/n}\le\|a^m\|^{1/m}$.
Ponieważ $m$ jest dowolne, a każdy wyraz ciągu pierwiastków
jest nie mniejszy niż ich infimum, otrzymujemy żądaną granicę.
Jeżeli $|z|>L$, wybieramy $q$ z $L/|z|<q<1$.
Dla dostatecznie dużych $n$ mamy $\|a^n z^{-n}\|\le q^n$.
Szereg $\sum_{n\ge0}a^n z^{-n-1}$ jest więc normowo
zbieżny. Pomnożenie jego sumy częściowej z dowolnej strony
przez $z1_A-a$ daje $1_A-a^{N+1}z^{-N-1}$,
którego granicą jest $1_A$. Dowodzi to odwracalności.
Zatem $r(a)\le L$.

Dla $R>\|a\|$ z szeregu Neumanna i skalarnego wzoru Cauchy'ego
(stosowanego po złożeniu z dowolnym $\varphi\in A'$) otrzymujemy
\[
 a^n=\frac{1}{2\pi i}\int_{|z|=R}z^nR_a(z)\,\mathrm dz.
\]
Całka jest granicą sum Riemanna w zupełnej przestrzeni Banacha $A$;
krzywa ma dodatnią orientację. Holomorficzność rezolwenty pozwala
przesunąć okrąg do dowolnego promienia $R>r(a)$, gdyż pierścień
między okręgami nie przecina widma. W konsekwencji
$\|a^n\|\le R^{n+1}\sup_{|z|=R}\|R_a(z)\|$.
Po wyciągnięciu pierwiastka stopnia $n$ uzyskujemy $L\le R$.
Przechodząc z $R$ do $r(a)$ dostajemy tezę.
\end{proof}

\begin{przyklad}[Element quasi-nilpotentny]
Operator Volterry $V$ na $C([0,1],\C)$ z
poprzedniego rozdziału ma $\sigma(V)=\{0\}$,
mimo że $V\ne0$. Zatem $r(V)=0$ i
$\lim_n\|V^n\|^{1/n}=0$. Jednak dla funkcji stałej $1$ mamy
$(V^m1)(t)=t^m/m!$ z indukcyjnego wzoru na całkę
(Przykład~\ref{ex:volterra-spectrum30}), więc $V^m\ne0$
dla każdego $m\ge1$. Operator nie jest nilpotentny.
\end{przyklad}

\subsection{Charaktery, ideały maksymalne i transformacja Gelfanda}

\noindent\emph{Powiązania tematyczne w Bibliotece matematycznej:} \atlasref{AA-09::aa:def:prime-max-ideal}{AA-09: ideały maksymalne}.


\begin{definicja}[Charakter i przestrzeń charakterów]
\label{def:character31}
Niech $A$ będzie \emph{przemienną} zespoloną algebrą Banacha
z jedynką. Charakterem nazywamy niezerowy homomorfizm
zespolonych algebr $\chi:A\to\C$ (a priori bez założenia ciągłości).
Zbiór wszystkich charakterów oznaczamy $\Delta(A)$.
\emph{Transformacja Gelfanda} przyporządkowuje elementowi $a\in A$
funkcję $\widehat a:\Delta(A)\to\C$, $\widehat a(\chi)=\chi(a)$.
\end{definicja}

\begin{lemat}[Twierdzenie Gelfanda--Mazura]
\label{lem:gelfand-mazur31}
Każda zespolona algebra Banacha z jedynką, w której każdy niezerowy
element jest odwracalny, jest izomorficzna jako unitalna algebra z $\C$.
\end{lemat}
\begin{proof}
Dla $a\in A$ widmo $\sigma(a)$ jest niepuste na mocy
twierdzenia~\ref{thm:spectrum-nonempty31}. Wybierzmy
$z\in\sigma(a)$. Element $a-z1_A$ nie jest odwracalny,
więc przy założeniu o $A$ musi być równy zeru.
Zatem każdy $a$ jest skalarną wielokrotnością $1_A$.
\end{proof}

\begin{twierdzenie}[Charaktery i widmo elementu]
\label{thm:gelfand-spectrum31}
Niech $A$ będzie przemienną zespoloną algebrą Banacha z jedynką.
Wówczas każdy charakter $\chi\in\Delta(A)$ jest ciągły,
$\|\chi\|=1$, a $\Delta(A)$ jest niepustą zwartą przestrzenią
w topologii słabej-* $\sigma(A',A)$. Dla każdego $a\in A$ zachodzi
\[
 \sigma_A(a)=\{\chi(a):\chi\in\Delta(A)\},\qquad
 \|\widehat a\|_\infty=r(a)\le\|a\|.
\]
W szczególności $a\mapsto\widehat a$ jest ciągłym homomorfizmem
unitalnym $A\to C(\Delta(A),\C)$.
\end{twierdzenie}
\begin{proof}
Dla charakteru $\chi$ wybieramy $b\in A$ z $\chi(b)\ne0$.
Równość $\chi(b)=\chi(1_A)\chi(b)$ daje $\chi(1_A)=1$.
Jeżeli $\chi(a)\notin\sigma(a)$, element $a-\chi(a)1_A$
byłby odwracalny, chociaż jego obraz przez $\chi$ jest zerem.
To niemożliwe, więc $|\chi(a)|\le r(a)\le\|a\|$.
Zatem $\chi$ jest ciągły i ma normę $1$.

Jeżeli $z\in\sigma(a)$, element $a-z1_A$ jest nieodwracalny.
Ideał $A(a-z1_A)$ jest właściwy, bo w przemiennej
algebrze zawieranie jednostki oznaczałoby odwracalność
$a-z1_A$. Lemat Zorna stosujemy do właściwych ideałów
zawierających ten ideał: suma łańcucha jest ideałem
właściwym, ponieważ nie zawiera jednostki. Otrzymujemy
ideał maksymalny $M$. Jego domknięcie jest ideałem
z ciągłości mnożenia. Gdyby $\overline M=A$, istniałby
$m\in M$ z $\|m-1_A\|<1$, a więc odwracalny, co przeczy
właściwości $M$. Maksymalność daje $\overline M=M$.
Iloraz $A/M$ jest zupełny w normie ilorazowej; norma
jest podmultiplikatywna, bo można niezależnie przybliżać
infima norm reprezentantów obu klas. Jednostka $[1_A]$
ma normę nie większą niż $1$, a z niezerowości i
$[1_A]^2=[1_A]$ wynika nierówność odwrotna.
Dla $b\notin M$ maksymalność daje $M+Ab=A$, czyli
$1_A=m+cb$ dla pewnych $m\in M,c\in A$.
W ilorazie $[c][b]=[1_A]$; przemienność daje też
odwrotność prawostronną. Zatem $A/M$ jest zespoloną
algebrą Banacha, której każdy niezerowy element jest odwracalny. Z lematu~\ref{lem:gelfand-mazur31}
$A/M\simeq\C$, a odwzorowanie ilorazowe złożone z tym
izomorfizmem daje charakter $\chi$ spełniający $\chi(a)=z$.
To dowodzi formuły na widmo i niepustości $\Delta(A)$.

Topologia słaba-* czyni $\chi\mapsto\chi(a)$ ciągłym.
Zbiór charakterów jest w kuli jednostkowej $A'$ domknięty,
ponieważ równania $\chi(1_A)=1$ i
$\chi(ab)=\chi(a)\chi(b)$ dla wszystkich $a,b\in A$
określają przecięcie zbiorów domkniętych w topologii
słabej-*. Każdy funkcjonał w tym przecięciu jest
niezerowym homomorfizmem. Z twierdzenia Banacha--Alaoglu
kula jest zwarta, a topologia słaba-* jest Hausdorffa,
gdyż różne funkcjonały różnią się na pewnym elemencie $A$.
Norma $\widehat a$ jest maksimum $|\chi(a)|$, czyli $r(a)$.
Liniowość, multiplikatywność i zachowanie jednostki
transformacji wynikają z tych samych własności charakterów;
oszacowanie normy daje jej ciągłość.
\end{proof}

\begin{uwaga}[Algebry niepółproste]
Transformacja Gelfanda nie musi być różnowartościowa.
Na przykład w dwuwymiarowej algebrze
$A=\C[\epsilon]/(\epsilon^2)$, wyposażonej w normę
$\|a+b\epsilon\|=|a|+|b|$, istnieje tylko jeden charakter,
$\chi(a+b\epsilon)=a$. Niezerowy element $\epsilon$ przechodzi na
funkcję zerową. Jądro transformacji Gelfanda jest przecięciem jąder
wszystkich charakterów. Każde takie jądro jest ideałem
maksymalnym, gdyż charakter jest surjekcją na $\C$.
Odwrotnie każdy ideał maksymalny jest jądrem charakteru,
jak wykazano w Twierdzeniu~\ref{thm:gelfand-spectrum31}.
Przecięcie wszystkich ideałów maksymalnych to radykał
Jacobsona algebry przemiennej. Ze wzoru
$\|\widehat a\|_\infty=r(a)$ wynika, że składa się on
dokładnie z elementów o promieniu spektralnym zero.
\end{uwaga}

\subsection{Holomorficzny rachunek funkcyjny}

\noindent\emph{Powiązania tematyczne w Bibliotece matematycznej:} \atlasref{LA-18::thm:lin18-analytic-calculus}{LA-18: analityczny rachunek funkcyjny macierzy}.


\begin{definicja}[Rachunek funkcyjny Dunforda]
\label{def:holomorphic-calculus31}
Niech $A$ będzie zespoloną algebrą Banacha z jedynką, $a\in A$,
a $U\subseteq\C$ otwartym otoczeniem $\sigma_A(a)$.
Dla $f\in\mathcal O(U)$ wybieramy skończoną, zorientowaną
rodzinę zamkniętych krzywych kawałkami $C^1$, oznaczoną $\Gamma$,
leżących w $U\setminus\sigma(a)$ i stanowiących dodatnio
zorientowany brzeg obszaru wielokątnego $D$, dla którego
$\sigma(a)\subset D\subseteq\overline D\subset U$.
Dopuszczamy wiele składowych i dziury; wspólne krawędzie
kwadratów w konstrukcji brzegu znoszą się z przeciwnymi
orientacjami. Indeks $\operatorname{Ind}_\Gamma(w)$
wynosi $1$ w $D$ poza brzegiem i $0$ poza $\overline D$.
Takie kontury istnieją: wybieramy $d>0$ tak, aby zamknięte
$d$-otoczenie zwartego widma mieściło się w $U$.
Na siatce kwadratowej o średnicy oczka mniejszej niż
$d/4$ bierzemy wszystkie domknięte kwadraty przecinające
zamknięte $d/2$-otoczenie widma. Ich suma jest skończona,
zawiera widmo we wnętrzu i mieści się w $U$.
Wewnętrzne krawędzie usuwamy, a pozostałe orientujemy
tak, by wnętrze leżało po lewej stronie. W każdym
wierzchołku liczba wchodzących krawędzi jest równa
liczbie wychodzących (tak było dla każdego kwadratu
i pozostaje po skreśleniu par). Skończony zbiór tych
krawędzi rozkłada się więc na zamknięte drogi wielokątne,
również gdy niektóre kwadraty stykają się tylko wierzchołkami.
Dla punktu poza wszystkimi krawędziami siatki indeks
sumy brzegów wynika z addytywności indeksu: wynosi $1$
w jednym z wybranych kwadratów i $0$ poza nimi.
Dla punktu na usuniętej krawędzi lub wewnętrznym wierzchołku
ten sam wniosek wynika z lokalnej stałości indeksu poza
pozostawionym konturem.
Tę konstrukcję można powtórzyć w dowolnym mniejszym
otoczeniu widma, otrzymując kontury zagnieżdżone.
Definiujemy
\[
 f(a):=\frac{1}{2\pi i}\int_\Gamma f(z)(z1_A-a)^{-1}\,\mathrm dz.
\]
Całka oznacza granicę sum Riemanna po parametryzacjach konturów
w normie przestrzeni Banacha $A$; zapis obejmuje sumę całek
po wszystkich składowych konturu. Ciągłość funkcji podcałkowej na każdym odcinku parametryzacji
zapewnia istnienie całki w zupełnej przestrzeni $A$.
\end{definicja}

\begin{twierdzenie}[Własności holomorficznego rachunku funkcyjnego]
\label{thm:holomorphic-calculus31}
Niech $A,a,U$ będą jak w definicji~\ref{def:holomorphic-calculus31}.
Odwzorowanie $f\mapsto f(a)$ z $\mathcal O(U)$ do $A$ jest unitalnym
homomorfizmem algebr. Dla wielomianu $p$ otrzymujemy zwykły element
$p(a)$, a dla $z\notin\sigma(a)$ rachunek funkcyjny funkcji
$w\mapsto(z-w)^{-1}$ daje $R_a(z)$. Wartość $f(a)$ zależy wyłącznie
od kiełka $f$ w otoczeniu widma i jest zgodna z ograniczaniem $U$.
Ponadto
\[
 \sigma_A(f(a))=f(\sigma_A(a)).
\]
\end{twierdzenie}
\begin{proof}
Dla dwóch dopuszczalnych obszarów $D_1,D_2$ wybieramy
trzeci $D_0$ z $\sigma(a)\subset D_0$ i
$\overline D_0\subset D_1\cap D_2$. Funkcja
$z\mapsto\varphi(f(z)R_a(z))$ jest holomorficzna
na otoczeniu każdego $\overline D_j\setminus D_0$
dla $j=1,2$. Skalarne twierdzenie Cauchy'ego na tych
obszarach (z uwzględnieniem orientacji wewnętrznego
brzegu) daje równość całek po $\partial D_j$ i
$\partial D_0$. Ciągły funkcjonał komutuje z całką,
co wynika z sum Riemanna. Twierdzenie Hahna--Banacha
przenosi te równości na równość w $A$.
Ten sam argument dla wspólnego mniejszego otoczenia
dowodzi zgodności z ograniczaniem $U$ i zależności od kiełka.
Liniowość jest bezpośrednia. Dla funkcji całkowitych
$1$ i $z$ wolno najpierw rozszerzyć dziedzinę do $\C$
(zgodność z ograniczaniem), a następnie użyć brzegu dużego kwadratu zawierającego widmo,
na którym $|z|>\|a\|$, i rozwinięcia rezolwenty. Jednostajna zbieżność szeregu na tym brzegu pozwala
całkować wyraz po wyrazie. Całka z $z^k$ podzielona
przez $2\pi i$ wynosi $1$ dla $k=-1$ oraz $0$
dla pozostałych całkowitych $k$ (istnieje funkcja
pierwotna $z^{k+1}/(k+1)$ dla $k\ne-1$).
Dostajemy odpowiednio $1_A$ i $a$.

Dla mnożenia wybieramy obszary $D_i,D_o$ z
$\sigma(a)\subset D_i$, $\overline D_i\subset D_o$
i $\overline D_o\subset U$. Kontury $\Gamma_i,\Gamma_o$
są rozłączne; $z$ oznacza zmienną na $\Gamma_o$,
a $w$ na $\Gamma_i$. Ciągłość mnożenia i jednostajna
zbieżność sum Riemanna uzasadniają zapis iloczynu całek
jako całki podwójnej oraz zmianę kolejności całkowania.
Używamy tożsamości
$R_a(z)R_a(w)=(R_a(w)-R_a(z))/(z-w)$.
Po podstawieniu do podwójnej całki dla $f(a)g(a)$
i całkowaniu najpierw po konturze zewnętrznym w wyrazie z $R_a(w)$
otrzymujemy przez wzór Cauchy'ego $f(w)g(w)R_a(w)$;
wyraz z $R_a(z)$ znika, ponieważ odpowiadająca mu funkcja
$w\mapsto g(w)/(z-w)$ jest holomorficzna wewnątrz konturu
wewnętrznego. Istotnie, $\operatorname{Ind}_{\Gamma_o}(w)=1$
dla $w\in\Gamma_i$, a $z\in\Gamma_o$ nie należy do
$\overline D_i$. Skalarne wzory Cauchy'ego stosujemy
na wszystkich składowych z ich orientacjami, również
gdy $U$ jest niespójny lub obszary mają dziury.
Daje to $(fg)(a)=f(a)g(a)$, a wraz z liniowością
i obliczeniem $1(a),z(a)$ także zgodność z wielomianami.

Dla funkcji $g(w)=(z-w)^{-1}$, $z\notin\sigma(a)$,
rozważanej na dostatecznie małym otoczeniu widma, zachodzi $(z1_A-a)g(a)=1_A$ na mocy własności homomorfizmu;
Analogicznie $g(a)(z1_A-a)=1_A$ (obrazy funkcji
komutują). Stąd $g(a)=R_a(z)$.
Jeśli $\lambda\notin f(\sigma(a))$, funkcja
$1/(f-\lambda)$ jest holomorficzna na pewnym otoczeniu widma.
Z własności homomorfizmu wynika odwracalność
$f(a)-\lambda1_A$, czyli $\sigma(f(a))\subseteq f(\sigma(a))$.
Jeśli $\lambda=f(\mu)$ dla $\mu\in\sigma(a)$,
funkcja $h(w)=(f(w)-f(\mu))/(w-\mu)$ ma usuwalną
osobliwość w $\mu$ i jest holomorficzna na otoczeniu widma.
Gdyby $f(a)-f(\mu)1_A=(a-\mu1_A)h(a)$ było odwracalne,
to czynniki, które komutują, miałyby odwrotności
(jeżeli $bc=cb$ i $d=(bc)^{-1}$, to $b,c$ komutują
z $bc$, a więc także z $d$; stąd $b(cd)=(cd)b=1_A$
oraz $c(bd)=(bd)c=1_A$). Otrzymalibyśmy sprzeczność
z $\mu\in\sigma(a)$; stąd przeciwne zawieranie.
\end{proof}

\begin{przyklad}[Eksponenta i funkcje wymierne]
Jeśli $a\in A$, szereg
$e^a=\sum_{n\ge0}a^n/n!$ zbiega w normie i odpowiada
holomorficznemu rachunkowi funkcyjnemu funkcji $e^z$:
szereg Taylora zbiega jednostajnie na konturze, a
oszacowanie normy całki przez długość konturu i supremum
normy funkcji podcałkowej pozwala przejść do granicy.
Twierdzenie~\ref{thm:holomorphic-calculus31} daje
$\sigma(e^a)=e^{\sigma(a)}$.
Jeśli $q$ jest wielomianem bez zer na $\sigma(a)$,
a $p$ wielomianem, to $(p/q)(a)=p(a)q(a)^{-1}$.
\end{przyklad}

\subsection{Algebry \texorpdfstring{$C^*$}{C*} i ich przemienny model}

\begin{definicja}[Algebra $C^*$]
\label{def:cstar31}
\emph{Algebrą $C^*$} nazywamy zespoloną algebrę Banacha $A$
z odwzorowaniem $a\mapsto a^*$ spełniającym
$(a+b)^*=a^*+b^*$, $(\lambda a)^*=\bar\lambda a^*$,
$(ab)^*=b^*a^*$, $(a^*)^*=a$ oraz
$\|a^*a\|=\|a\|^2$. W tym rozdziale algebry $C^*$
są unitalne. Element $a$ jest \emph{samosprzężony}, gdy $a=a^*$,
a \emph{normalny}, gdy $a^*a=aa^*$.
\end{definicja}

\begin{przyklad}[Dwie algebry $C^*$]
Dla niezerowej zespolonej przestrzeni Hilberta $H$
algebra $\mathcal L(H,H)$ z operacją operatora
sprzężonego jest algebrą $C^*$.
Równość $\|T^*\|=\|T\|$ wynika z oszacowania
$\|T^*y\|=\sup_{\|x\|\le1}|\langle Tx,y\rangle|
\le\|T\|\|y\|$ i zastosowania tego oszacowania do $T^*$.
Stąd $\|T^*T\|\le\|T\|^2$, zaś dla $\|x\|\le1$
mamy $\|Tx\|^2=|\langle T^*Tx,x\rangle|\le\|T^*T\|$,
co po wzięciu supremum daje odwrotną nierówność.
Dla niepustej przestrzeni zwartej Hausdorffa $K$
algebra $C(K,\C)$ z operacją $f^*(x)=\overline{f(x)}$
jest przemienną algebrą $C^*$, ponieważ
$\|f^*f\|_\infty=\sup_{x\in K}|f(x)|^2=\|f\|_\infty^2$.
\end{przyklad}

\begin{lemat}[Norma i widmo elementu normalnego]
\label{lem:cstar-normal31}
Niech $A$ będzie unitalną algebrą $C^*$, a $a\in A$ elementem
normalnym. Wtedy $r(a)=\|a\|$.
Jeżeli $a=a^*$, to $\sigma(a)\subseteq\R$.
\end{lemat}
\begin{proof}
Inwolucja zachowuje jednostkę: dla każdego $b=a^*$
mamy $1_A^* b=(a1_A)^*=b$ i $b1_A^*=(1_Aa)^*=b$,
więc $1_A^*=1_A$. Z własności $C^*$ wynika
$\|b^*\|=\|b\|$ (dla $b=0$ równość jest oczywista):
$\|b\|^2=\|b^*b\|\le\|b^*\|\|b\|$,
a analogicznie po zamianie $b$ na $b^*$.
Dla samosprzężonego $h$ mamy
$\|h^2\|=\|h^*h\|=\|h\|^2$.
Jeżeli $a$ jest normalny, to
$\|(a^2)^*a^2\|=\|(a^*a)^2\|=\|a^*a\|^2$;
przez własność $C^*$ otrzymujemy $\|a^2\|=\|a\|^2$.
Ponieważ $a$ komutuje z $a^*$, każda jego potęga jest
normalna: $(a^m)^*=(a^*)^m$ komutuje z $a^m$.
Indukcja daje $\|a^{2^n}\|=\|a\|^{2^n}$.
Wzór na promień spektralny implikuje $r(a)=\|a\|$.

Jeśli $h=h^*$, to dla $t\in\R$ i $u_t=e^{ith}$
ciągłość antyliniowej inwolucji pozwala zastosować ją
wyraz po wyrazie w szeregu eksponenty, dając
$u_t^*=e^{-ith}$. Z własności homomorfizmu rachunku
holomorficznego (Twierdzenie~\ref{thm:holomorphic-calculus31})
i tożsamości $e^{itz}e^{-itz}=1$ wynika
$u_t^*u_t=u_tu_t^*=1_A$. Własność $C^*$ daje
$\|u_t\|^2=\|1_A\|=1$, więc $\|u_t\|=1$.
Na mocy twierdzenia o odwzorowaniu widma
$e^{it\lambda}\in\sigma(u_t)$ dla $\lambda\in\sigma(h)$.
Zatem $|e^{it\lambda}|\le1$ dla wszystkich $t\in\R$,
a po uwzględnieniu również $-t$ otrzymujemy
$\operatorname{Im}\lambda=0$.
\end{proof}

\begin{lemat}[Zespolone twierdzenie Stone'a--Weierstrassa]
\label{lem:stone-weierstrass31}
Niech $K$ będzie niepustą zwartą przestrzenią Hausdorffa, a
$B\subseteq C(K,\C)$ podalgebrą zawierającą funkcje stałe,
rozdzielającą punkty $K$ i zamkniętą na sprzężenie zespolone.
Wtedy $B$ jest gęsta w $C(K,\C)$ w normie supremum.
\end{lemat}
\begin{proof}
Najpierw ustalmy, że każdą rzeczywistą funkcję ciągłą na odcinku
można jednostajnie aproksymować wielomianami rzeczywistymi.
Po zmianie zmiennej wystarczy odcinek $[0,1]$. Dla $h\in C([0,1],\R)$
rozważamy wielomiany Bernsteina
$B_nh(t)=\sum_{k=0}^nh(k/n)\binom nk t^k(1-t)^{n-k}$.
Ponieważ wagi są nieujemne, sumują się do $1$ na mocy wzoru dwumianowego, a ich
średnia dla zmiennej $k/n$ wynosi $t$ i wariancja
$t(1-t)/n\le1/(4n)$ (z tożsamości
$k\binom nk=n\binom{n-1}{k-1}$ oraz
$k(k-1)\binom nk=n(n-1)\binom{n-2}{k-2}$ dla $n\ge2$,
z pominięciem wyrazów $k<2$; dla $n=1$ drugi moment
silniowy wynosi zero).
Rozdzielamy sumę na $|k/n-t|\le\delta$ i pozostałe
indeksy. W pierwszej części różnica wartości $h$ jest
ograniczona przez moduł ciągłości, w drugiej przez
$2\|h\|_\infty$, a suma wag drugiej części nie przekracza
$t(1-t)/(n\delta^2)$ na mocy oszacowania drugiego momentu.
Otrzymujemy, przy oznaczeniu
$\omega_h(\delta)=\sup\{|h(s)-h(t)|:s,t\in[0,1],\ |s-t|\le\delta\}$,
oszacowanie
$\|B_nh-h\|_\infty\le\omega_h(\delta)
+\|h\|_\infty/(2n\delta^2)$ dla każdego $\delta>0$.
Wybierając najpierw małe $\delta$, a następnie duże $n$,
otrzymujemy zbieżność jednostajną.

Niech $B_\R=\{f\in B:f(K)\subseteq\R\}$.
Dzięki zamknięciu na sprzężenie, część rzeczywista i urojona
dowolnego elementu $B$ należą do $B_\R$.
Oznaczmy przez $D$ domknięcie $B_\R$ w $C(K,\R)$.
Domknięcie $D$ jest algebrą rzeczywistą z jednostką,
ponieważ dodawanie i mnożenie są ciągłe w normie supremum.
Dla $u\in D$ aproksymujemy funkcję $s\mapsto|s|$
wielomianami na $[-\|u\|_\infty,\|u\|_\infty]$
(przypadek $u=0$ jest natychmiastowy). Ich złożenia
z $u$ należą do $D$ i zbiegają jednostajnie do $|u|$.
Zatem $|u|\in D$.
Stąd $\max(u,v)=(u+v+|u-v|)/2$ i
$\min(u,v)=(u+v-|u-v|)/2$ należą do $D$.

Ustalmy $h\in C(K,\R)$ i $\varepsilon>0$.
Dla $x\ne y$ wybieramy element $B$ o różnych wartościach
w $x,y$. Jego część rzeczywista lub urojona daje
$v\in B_\R$ z $v(x)\ne v(y)$. Kładziemy
$u_{xy}=h(x)+(h(y)-h(x))(v-v(x))/(v(y)-v(x))$.
Dla $x=y$ wybieramy stałą $h(x)$.
W obu przypadkach $u_{xy}(x)=h(x)$ i $u_{xy}(y)=h(y)$.
Dla ustalonego $x$ z ciągłości wynika, że dla każdego $y$
istnieje otoczenie $V_y$ punktu $y$, na którym
$u_{xy}<h+\varepsilon$. Po wybraniu skończonego podpokrycia
odpowiednia funkcja $g_x=\min_j u_{xy_j}\in D$ spełnia
$g_x<h+\varepsilon$ na całym $K$ i $g_x(x)=h(x)$.
W pewnym otoczeniu $U_x$ punktu $x$ zachodzi także
$g_x>h-\varepsilon$. Wybierając skończenie wiele takich
otoczeń pokrywających $K$ i biorąc $g=\max_j g_{x_j}\in D$,
otrzymujemy $|g-h|<\varepsilon$ na $K$.
Zatem $D=C(K,\R)$. Każdą funkcję zespoloną można rozłożyć
na część rzeczywistą i urojoną, więc $B$ jest gęsta w $C(K,\C)$.
\end{proof}

\begin{twierdzenie}[Przemienne twierdzenie Gelfanda--Naimarka]
\label{thm:commutative-cstar31}
Niech $A$ będzie przemienną unitalną algebrą $C^*$.
Wówczas transformacja Gelfanda jest izometrycznym izomorfizmem
algebr z inwolucją
$A\simeq C(\Delta(A),\C)$,
przy czym $\widehat{a^*}=\overline{\widehat a}$.
\end{twierdzenie}
\begin{proof}
Dla elementu samosprzężonego $h$ mamy
$\chi(h)\in\sigma(h)\subseteq\R$ dla każdego charakteru $\chi$.
Każdy $a$ rozkłada się jako $a=h+ik$ z
$h=(a+a^*)/2$, $k=(a-a^*)/(2i)$ samosprzężonymi.
Zatem $\chi(a^*)=\overline{\chi(a)}$.
Ponieważ w przemiennej algebrze każdy element jest normalny,
lemat~\ref{lem:cstar-normal31} i twierdzenie~\ref{thm:gelfand-spectrum31}
dają $\|\widehat a\|_\infty=r(a)=\|a\|$.
Izometria jest różnowartościowa, a jej obraz jest domknięty
przez zupełność $A$. Obraz zawiera stałe, ponieważ
$\widehat{\lambda1_A}=\lambda$. Rozdziela punkty:
dla $\chi_1\ne\chi_2$ istnieje $a$ z
$\chi_1(a)\ne\chi_2(a)$. Jest podalgebrą zamkniętą
na sprzężenie dzięki wykazanej tożsamości dla $a^*$.
Z lematu~\ref{lem:stone-weierstrass31} wynika,
że jest gęsty w $C(\Delta(A))$, a więc równy całej tej przestrzeni.
Zatem odwzorowanie Gelfanda jest na $C(\Delta(A),\C)$.
\end{proof}

\begin{uwaga}[Granica zakresu]
Twierdzenie~\ref{thm:commutative-cstar31} dotyczy algebr $C^*$
\emph{przemiennych}. Pełne twierdzenie Gelfanda--Naimarka o
wiernej reprezentacji dowolnej algebry $C^*$ na przestrzeni Hilberta
wymaga m.in. konstrukcji GNS i będzie osobnym rozszerzeniem.
Nie utożsamiamy też rachunku holomorficznego z pełnym
ciągłym rachunkiem funkcyjnym dla operatorów normalnych;
ten drugi zostanie opracowany w \hyperref[sec:spectral-hilbert32]{rozdziale~32}.
\end{uwaga}


\subsection{Dwa przykłady pokazujące znaczenie struktury algebry}

\begin{przyklad}[Charaktery algebry splotowej $\ell^1(\N_0)$]
\label{ex:l1-characters31}
Niech $A=\ell^1(\N_0)$ z mnożeniem splotowym opisanym
w przykładzie~\ref{ex:banach-algebras31}. Oznaczmy przez $e_n$ ciąg,
którego $n$-ty wyraz jest równy $1$, a wszystkie pozostałe są zerowe.
Wówczas $e_n=e_1^{*n}$ i $e_0=1_A$.
Jeżeli $\chi\in\Delta(A)$, to $z:=\chi(e_1)$ spełnia
$|z|\le\|e_1\|_1=1$. Przez multiplikatywność i ciągłość charakteru
otrzymujemy dla dowolnego $a=(a_n)\in A$
\[
 \chi(a)=\sum_{n=0}^\infty a_nz^n.
\]
Odwrotnie, dla każdego $|z|\le1$ ten szereg określa charakter:
jest absolutnie zbieżny i zachowuje splot dzięki możliwości zamiany
kolejności sum bezwzględnie zbieżnych. Zatem przestrzeń charakterów
$\Delta(A)$ jest homeomorficzna z domkniętym dyskiem jednostkowym:
odwzorowanie $z\mapsto\chi_z$ jest ciągłe w topologii
słabej-* dzięki jednostajnej zbieżności szeregu
$\sum a_nz^n$ na domkniętym dysku dla każdego $a\in A$,
jest bijekcją, a odwrotność $\chi\mapsto\chi(e_1)$
jest ciągła z definicji tej topologii; a transformacja Gelfanda przyporządkowuje
ciągowi $a$ funkcję $\widehat a(z)=\sum_{n\ge0}a_nz^n$.
W szczególności $\sigma_A(e_1)=\overline{\mathbb D}$,
choć $e_1$ nie jest elementem odwracalnym.
\end{przyklad}

\begin{przyklad}[Przesunięcie jednostronne — widmo niebędące zbiorem wartości własnych]
\label{ex:unilateral-shift31}
Na zespolonej przestrzeni Hilberta $H=\ell^2(\N)$ niech
$S(x_1,x_2,\ldots)=(0,x_1,x_2,\ldots)$.
Jest to izometria o normie $1$, więc $\sigma(S)\subseteq\overline{\mathbb D}$.
Dla każdego $|\lambda|<1$ operator sprzężony $S^*$ ma niezerowy
wektor własny $v=(1,\overline\lambda,\overline\lambda^2,\ldots)$
z wartością własną $\overline\lambda$.
Istotnie, $\sum_{n\ge0}|\lambda|^{2n}<\infty$,
a $S^*$ jest przesunięciem w lewo, więc
$(\lambda I-S)^*v=0$. Dla każdego $x\in H$
iloczyn skalarny $\langle(\lambda I-S)x,v\rangle$
jest zerowy. Zatem obraz $\lambda I-S$ jest zawarty
we właściwej domkniętej podprzestrzeni $\{v\}^\perp$,
nie jest gęsty i operator nie jest surjektywny.
Stąd $\lambda\in\sigma(S)$.
Z domkniętości widma dostajemy $\sigma(S)=\overline{\mathbb D}$.
Jednak $S$ nie ma żadnych wartości własnych: równanie $Sx=\lambda x$
dla $\lambda\ne0$ wymusza kolejno $x_1=0,x_2=0,\ldots$,
a dla $\lambda=0$ daje $x=0$ z różnowartościowości $S$.
Pokazuje to, dlaczego w ogólnej teorii operatorów nie należy
utożsamiać widma ze zbiorem wartości własnych, co było możliwe
poza zerem w twierdzeniu Riesza--Schaudera dla operatorów zwartych.
\end{przyklad}

% Wersja źródłowa: audyt II, jednostka 32
\setcounter{section}{31}
\section{Teoria spektralna operatorów ograniczonych na przestrzeniach Hilberta}
\label{sec:spectral-hilbert32}
\noindent\textbf{Wyniki wykorzystywane.}
Korzystamy z przemiennego twierdzenia Gelfanda--Naimarka
(Twierdzenie~\ref{thm:commutative-cstar31}), twierdzenia
Stone'a--Weierstrassa (Lemat~\ref{lem:stone-weierstrass31}),
twierdzenia Riesza--Markowa o reprezentacji dodatnich
funkcjonałów na $C(K)$, własności całki Lebesgue'a
i twierdzenia Riesza--Schaudera (Twierdzenie~\ref{thm:riesz-schauder30}).
\par\medskip

Poprzedni rozdział dostarczył holomorficznego rachunku funkcyjnego
w dowolnej zespolonej algebrze Banacha i twierdzenia Gelfanda--Naimarka
dla przemiennych algebr $C^*$. Dla operatorów normalnych na przestrzeniach
Hilberta otrzymamy znacznie więcej: rachunek wszystkich ograniczonych
funkcji borelowskich oraz reprezentację przez \emph{miarę spektralną}.
Nie zakładamy ośrodkowości przestrzeni Hilberta. Iloczyn skalarny jest
liniowy w pierwszym argumencie, tak jak w \hyperref[sec:hilbert13]{rozdziale~13}.

\subsection{Operatory normalne i ciągły rachunek funkcyjny}

\begin{definicja}[Operator normalny, samosprzężony i unitarny]
\label{def:normal32}
Niech $H$ będzie niezerową zespoloną przestrzenią Hilberta,
a $T\in\mathcal L(H,H)$. Operator $T$ jest \emph{normalny}, gdy
$TT^*=T^*T$, \emph{samosprzężony}, gdy $T=T^*$, oraz
\emph{unitarny}, gdy $T^*T=TT^*=I_H$. Operator samosprzężony
$T$ jest \emph{dodatni}, co zapisujemy $T\geq0$, gdy
$\langle Tx,x\rangle\geq0$ dla każdego $x\in H$.
\end{definicja}

\begin{lemat}[Zgodność widma z podalgebrą $C^*$]
\label{lem:spectral-permanence32}
Niech $H$ będzie niezerową zespoloną przestrzenią Hilberta,
$T\in\mathcal L(H,H)$ operatorem normalnym i
$A=C^*(I_H,T)$ domkniętą w normie podalgebrą generowaną
przez $I_H,T,T^*$. Wówczas $A$ jest przemienną unitalną
algebrą $C^*$ i $\sigma_A(T)=\sigma_{\mathcal L(H,H)}(T)$.
\end{lemat}
\begin{proof}
Normalność daje przemienność $T$ i $T^*$, a więc całej
podalgebry $A$. Zawsze $\sigma_{\mathcal L(H,H)}(T)\subseteq\sigma_A(T)$.
Niech odwrotnie $\lambda\notin\sigma_{\mathcal L(H,H)}(T)$ i
$B=T-\lambda I_H$. Wówczas $\|x\|\le\|B^{-1}\|\|Bx\|$, więc dla
$m=\|B^{-1}\|^{-2}>0$ zachodzi
$\|Bx\|^2\geq m\|x\|^2$ dla wszystkich $x\in H$.
Operator $S=B^*B$ należy do $A$, jest dodatni i spełnia
$mI_H\leq S\leq\|S\|I_H$. Przypomnijmy uzasadnienie wzoru na normę operatora
samosprzężonego $R$. Niech
$c=\sup_{\|x\|=1}|\langle Rx,x\rangle|$.
Dla jednostkowych $x,y$ dobieramy fazę wektora $y$ tak,
aby $\langle Rx,y\rangle$ było nieujemne rzeczywiste.
Z tożsamości
\[
4\operatorname{Re}\langle Rx,y\rangle
=\langle R(x+y),x+y\rangle-\langle R(x-y),x-y\rangle
\]
i prawa równoległoboku dostajemy
$4|\langle Rx,y\rangle|\le c(\|x+y\|^2+\|x-y\|^2)=4c$.
Supremum po $x,y$ daje $\|R\|\le c$, a odwrotna
nierówność wynika z nierówności Cauchy'ego--Schwarza.
Stosując ten wzór do $R=I_H-S/\|S\|$, otrzymujemy
$\|I_H-S/\|S\|\|\leq1-m/\|S\|<1$.
Szereg Neumanna dowodzi, że $S^{-1}\in A$. Ponieważ $B$ jest
normalny, $B$ i $B^*$ komutują z $S^{-1}$; zatem
$S^{-1}B^*$ jest obustronną odwrotnością $B$ w $A$.
W konsekwencji $\lambda\notin\sigma_A(T)$.
\end{proof}

\begin{twierdzenie}[Ciągły rachunek funkcyjny operatora normalnego]
\label{thm:continuous-calculus32}
Niech $H$ i $T$ będą jak w lemacie~\ref{lem:spectral-permanence32},
a $K=\sigma(T)\subset\C$. Istnieje dokładnie jeden izometryczny
unitalny homomorfizm algebr z inwolucją
$\Phi_T:C(K,\C)\to\mathcal L(H,H)$ taki, że
$\Phi_T(\mathrm{id}_K)=T$. Dla $f\in C(K,\C)$ piszemy $f(T)=\Phi_T(f)$.
Zachodzą
$\|f(T)\|=\max_{z\in K}|f(z)|$,
$\overline f(T)=f(T)^*$ i $\sigma(f(T))=f(K)$.
Jeżeli $p(z,\overline z)$ jest wielomianem w dwóch zmiennych,
to $p(T, T^*)=p|_K(T)$.
\end{twierdzenie}
\begin{proof}
Przemienne twierdzenie Gelfanda--Naimarka z
poprzedniego rozdziału daje izometryczny
izomorfizm $A\simeq C(\Delta(A),\C)$, zaś lemat
\ref{lem:spectral-permanence32} identyfikuje $K$ z widmem $T$
w $A$. Odwzorowanie
$\Delta(A)\ni\chi\mapsto\chi(T)\in K$ jest ciągłą surjekcją.
Jest różnowartościowe: charakter jest określony przez wartości
na $T,T^*$ i jednostce, a $\chi(T^*)=\overline{\chi(T)}$.
Jest więc homeomorfizmem zwartych przestrzeni Hausdorffa.
Przeniesienie transformacji Gelfanda daje $\Phi_T$.
Jednoznaczność wynika z gęstości wielomianów w $z,\overline z$
w $C(K)$, czyli z twierdzenia Stone'a--Weierstrassa.
Izometria i zgodność z inwolucją wynikają z twierdzenia
Gelfanda--Naimarka. Dla równości widm rozważmy
$B=C^*(I_H,f(T))\subseteq A$. Operator $f(T)$ jest
normalny. Lemat~\ref{lem:spectral-permanence32} dla
$f(T)$ daje $\sigma_{\mathcal L(H,H)}(f(T))=\sigma_B(f(T))$,
a inkluzje widm przy włączaniu algebr dają
\[
\sigma_{\mathcal L(H,H)}(f(T))\subseteq\sigma_A(f(T))
\subseteq\sigma_B(f(T)).
\]
Wszystkie trzy zbiory są więc równe. W $A\simeq C(K)$
widmo $f(T)$ jest dokładnie $f(K)$.
\end{proof}

\begin{uwaga}[Związek z rachunkiem holomorficznym]
Jeżeli $f$ jest holomorficzna na otoczeniu $K$, to $f(T)$
z twierdzenia~\ref{thm:continuous-calculus32} zgadza się z całką
Dunforda z poprzedniego rozdziału.
Istotnie, dla $w\notin K$ rachunek ciągły daje
$(w-z)^{-1}(T)=(wI_H-T)^{-1}$. Jeżeli $\Gamma$ jest
konturem dopuszczalnym w definicji całki Dunforda,
ciągłość homomorfizmu $\Phi_T$ pozwala przenieść go
pod znak całki (jako granicę sum Riemanna). Wtedy
\[
 \frac1{2\pi i}\int_\Gamma f(w)(wI_H-T)^{-1}\,\mathrm dw
 =\Phi_T\!\left(\frac1{2\pi i}\int_\Gamma
 \frac{f(w)}{w-z}\,\mathrm dw\right)=\Phi_T(f|_K),
\]
gdzie ostatnia równość jest skalarnym wzorem Cauchy'ego
stosowanym punktowo dla $z\in K$.
Dla operatora nienormalnego sam rachunek
ciągły na $C(\sigma(T))$ z własnością $\overline f(T)=f(T)^*$
na ogół nie istnieje.
\end{uwaga}

\subsection{Miary spektralne i twierdzenie spektralne}

\begin{definicja}[Miara spektralna]
\label{def:pvm32}
Niech $H$ będzie zespoloną przestrzenią Hilberta i $K\subset\C$
niepustym zwartym zbiorem. \emph{Miarą spektralną} na
$\mathcal B(K)$ nazywamy odwzorowanie
$E:\mathcal B(K)\to\mathcal L(H,H)$ spełniające:
$E(B)$ jest rzutem ortogonalnym dla każdego zbioru borelowskiego $B$,
$E(\varnothing)=0$, $E(K)=I_H$ oraz
$E(B\cap C)=E(B)E(C)$ dla $B,C\in\mathcal B(K)$.
Ponadto dla parami rozłącznych $(B_n)_{n\geq1}$ wymaga się
\emph{ciągowej addytywności w silnej topologii operatorowej}, tj.
$E(\bigcup_nB_n)x=\lim_{N\to\infty}\sum_{n=1}^NE(B_n)x$
dla każdego $x\in H$.
Dla $x\in H$ związana z $E$ dodatnia skończona miara skalarna
to $\mu_x(B):=\langle E(B)x,x\rangle=\|E(B)x\|^2$.
\end{definicja}

\begin{lemat}[Model cykliczny]
\label{lem:cyclic-model32}
Niech $T\in\mathcal L(H,H)$ będzie normalny, $K=\sigma(T)$,
$a\in H\setminus\{0\}$ i
$H_a=\overline{\{f(T)a:f\in C(K,\C)\}}$.
Wówczas $H_a$ jest podprzestrzenią redukującą $T$
(niezmienniczą względem $T$ i $T^*$).
Istnieje jedyna skończona dodatnia regularna miara borelowska
$\nu_a$ na $K$ i izometria liniowa na
$U_a:L^2(K,\nu_a)\to H_a$ takie, że
$U_a1=a$ i $U_aM_z=(T|_{H_a})U_a$, gdzie
$(M_zf)(w)=wf(w)$.
\end{lemat}
\begin{proof}
Z domknięcia algebry $C(K)$ na sprzężenie otrzymujemy
$T H_a\subseteq H_a$ i $T^*H_a\subseteq H_a$.
Definiujemy $L_a(f)=\langle f(T)a,a\rangle$. Jest to ciągły
dodatni funkcjonał na $C(K)$, ponieważ dla $f\geq0$
możemy zapisać $f=|\sqrt f|^2$, a więc
$f(T)=(\sqrt f(T))^*\sqrt f(T)\geq0$.
Z twierdzenia Riesza--Markowa o reprezentacji dodatnich
funkcjonałów na $C(K)$ otrzymujemy jednoznaczną dodatnią,
skończoną regularną miarę borelowską $\nu_a$ spełniającą
$L_a(f)=\int_K f(z)\,\mathrm d\nu_a(z)$.
Mamy $\nu_a(K)=\|a\|^2$ oraz
\[
 \langle f(T)a,g(T)a\rangle
 =\langle (\overline g f)(T)a,a\rangle
 =\int_K f(z)\overline{g(z)}\,\mathrm d\nu_a(z).
\]
Zatem odwzorowanie $f\mapsto f(T)a$ zależy tylko od
klasy $f$ w $L^2(\nu_a)$ i jest izometrią liniową na
swój obraz w $H$. Nie zakładamy, że półnorma $L^2$
jest normą na wszystkich funkcjach ciągłych.
Aby wykazać gęstość ich klas, dla borelowskiego $B$
i $\varepsilon>0$ wybieramy z regularności miary
zwarty $C\subseteq B$ i względnie otwarty $O\supseteq B$
z $\nu_a(O\setminus C)<\varepsilon^2$.
Lemat Urysohna daje $h\in C(K,[0,1])$ równe $1$ na $C$
i $0$ poza $O$, więc
$\|h-\chi_B\|_{L^2(\nu_a)}^2\le\nu_a(O\setminus C)<\varepsilon^2$.
Gęstość funkcji prostych w $L^2$ kończy argument. Izometria przedłuża się do $U_a$ na całym
$L^2(K,\nu_a)$; jej obraz jest domknięty i równy $H_a$.
Ponieważ $U_a(zf)=T U_a(f)$ dla funkcji ciągłych,
przez gęstość otrzymujemy równość operatorową.
Jeżeli inna miara i izometria na spełniają te warunki,
to izometria jest unitarna, zatem przeplata również
sprzężone operatory mnożenia. Dla wielomianu
$p(z,\overline z)$ obraz funkcji $p$ musi być równy
$p(T,T^*)a$. Gęstość tych wielomianów w $C(K)$
i skończoność miar dają to samo dla każdej funkcji ciągłej.
Parowanie z $a$ identyfikuje całki funkcji ciągłych
z $L_a$, więc jednoznaczność Riesza--Markowa daje
równość miar. Gęstość w $L^2$ daje równość izometrii.
\end{proof}

\begin{lemat}[Całkowanie względem dowolnej miary spektralnej]
\label{lem:pvm-integral32}
Niech $E$ będzie miarą spektralną z Definicji~\ref{def:pvm32}.
Dla każdej ograniczonej borelowskiej funkcji $f$ całka
$\int_K f\,\mathrm dE$ jest dobrze określona jako granica
w normie operatorowej całek funkcji prostych aproksymujących
$f$ jednostajnie. Jest to unitalny homomorfizm algebr
z inwolucją, o normie nie większej niż $1$.
Dla $\mu_{x,y}(B)=\langle E(B)x,y\rangle$ zachodzą
\[
\left\langle\left(\int_K f\,\mathrm dE\right)x,y\right\rangle
=\int_K f\,\mathrm d\mu_{x,y},\qquad
\left\|\left(\int_K f\,\mathrm dE\right)x\right\|^2
=\int_K |f|^2\,\mathrm d\mu_x .
\]
\end{lemat}
\begin{proof}
Dla $s=\sum_{j=1}^m c_j\chi_{B_j}$ na skończonym
borelowskim podziale $K$ definiujemy
$I_E(s)=\sum_{j=1}^m c_j E(B_j)$. Skończona addytywność
i wspólne rozdrobnienie podziałów dają niezależność
od zapisu, liniowość i multiplikatywność; samosprzężoność
rzutów daje $I_E(\overline s)=I_E(s)^*$.
Ortogonalność obrazów rzutów daje
\[
\|I_E(s)x\|^2=\sum_j |c_j|^2\|E(B_j)x\|^2
\le\|s\|_\infty^2\|x\|^2 .
\]
Ograniczoną funkcję borelowską aproksymujemy jednostajnie
funkcjami prostymi, dzieląc ograniczony zbiór jej wartości
na skończenie wiele borelowskich zbiorów o średnicy
dążącej do zera. Oszacowanie normy i zupełność
$\mathcal L(H,H)$ dają granicę niezależną od aproksymacji.
Ciągłość mnożenia i sprzężenia przenosi własności na granicę.
Silna addytywność $E$ daje ciągową addytywność
$\mu_{x,y}$. Dla skończonego rozłącznego układu $(B_j)$
\[
\sum_j|\langle E(B_j)x,y\rangle|
=\sum_j|\langle E(B_j)x,E(B_j)y\rangle|
\le\|x\|\|y\|,
\]
na mocy nierówności Cauchy'ego--Schwarza dla sum i
ortogonalności. Miara zespolona ma zatem skończone
wahanie całkowite. Dla funkcji prostych oba wzory
wynikają z definicji, a jednostajna aproksymacja
i skończoność miar dają oba wzory dla $f$.
\end{proof}

\begin{twierdzenie}[Twierdzenie spektralne dla operatora normalnego]
\label{thm:spectral-normal32}
Niech $H$ będzie niezerową zespoloną przestrzenią Hilberta,
$T\in\mathcal L(H,H)$ operatorem normalnym, a $K=\sigma(T)$.
Istnieje dokładnie jedna miara spektralna
$E:\mathcal B(K)\to\mathcal L(H,H)$ taka, że
$T=\int_Kz\,\mathrm dE(z)$.
Dla każdej ograniczonej borelowskiej funkcji $f:K\to\C$
istnieje operator $f(T)=\int_K f(z)\,\mathrm dE(z)$ spełniający
$\langle f(T)x,y\rangle=\int_K f(z)\,\mathrm d\mu_{x,y}(z)$,
gdzie $\mu_{x,y}(B):=\langle E(B)x,y\rangle$
jest skończoną miarą zespoloną, otrzymaną przez
polaryzację miar dodatnich $\mu_x$.
Zachodzą: $(fg)(T)=f(T)g(T)$,
$f(T)^*=\overline f(T)$ oraz
$\|f(T)x\|^2=\int_K|f(z)|^2\,\mathrm d\mu_x(z)$.
Rachunek ten przedłuża ciągły rachunek funkcyjny z
twierdzenia~\ref{thm:continuous-calculus32}.
\end{twierdzenie}
\begin{proof}
\textbf{Krok 1: rozkład na podprzestrzenie cykliczne.}
Rozważamy rodziny parami ortogonalnych, niezerowych podprzestrzeni
postaci $H_a$ z lematu~\ref{lem:cyclic-model32}. Są one redukujące.
Porządkujemy te rodziny przez włączanie; suma łańcucha
jest rodziną parami ortogonalną. Lemat Zorna daje
rodzinę maksymalną $(H_{a_j})_{j\in J}$.
Ortogonalne dopełnienie domkniętej sumy jest redukujące,
bo suma jest niezmiennicza względem $T$ i $T^*$.
Każdy niezerowy wektor z dopełnienia generowałby kolejną
podprzestrzeń cykliczną: dopełnienie jest niezmiennicze
względem wielomianów w $T,T^*$, a przez aproksymację
także względem każdego $f(T)$ dla $f\in C(K)$.
Przeczy to maksymalności.
Zatem $H=\bigoplus_{j\in J}H_{a_j}$, także dla nieprzeliczalnego $J$.

\textbf{Krok 2: konstrukcja miary.}
Z lematu~\ref{lem:cyclic-model32} każdy $T|_{H_{a_j}}$
jest unitarnie równoważny operatorowi mnożenia $M_z$ na
$L^2(K,\nu_{a_j})$. Oznaczmy tę równoważność przez $U_j$.
Dla borelowskiego $B\subseteq K$ definiujemy
$E(B)|_{H_{a_j}}=U_jM_{\chi_B}U_j^{-1}$.
Są to rzuty ortogonalne, komutują i spełniają
$E(B\cap C)=E(B)E(C)$, $E(K)=I_H$. Jeżeli $(B_n)$ jest
rodziną parami rozłączną, to dla każdego $x_j\in H_{a_j}$
\[
 \bigl\|\bigl(E(\textstyle\bigcup_n B_n)
 -\sum_{n=1}^{N}E(B_n)\bigr)x_j\bigr\|^2
 =\int_{\bigcup_{n>N}B_n}|U_j^{-1}x_j|^2\,\mathrm d\nu_{a_j}
 \longrightarrow0.
\]
Każdy wektor $x\in\bigoplus_{j\in J}H_{a_j}$ ma co najwyżej
przeliczalnie wiele niezerowych składowych; sumując powyższe
oszacowania i stosując zbieżność zmajoryzowaną dla szeregu
nieujemnych składników (każdy jest ograniczony przez
$\|x_j\|^2$, a $\sum_j\|x_j\|^2=\|x\|^2$),
otrzymujemy silną ciągową addytywność.

\textbf{Krok 3: rachunek funkcyjny.}
Dla ograniczonej funkcji borelowskiej $f$ definiujemy
$f(T)|_{H_{a_j}}=U_jM_fU_j^{-1}$.
Normy tych operatorów nie przekraczają $\|f\|_\infty$,
więc określają jeden operator na całej sumie ortogonalnej.
Dla funkcji prostych tak określony operator jest równy
$I_E(f)$ z Lematu~\ref{lem:pvm-integral32}. Dla
ograniczonych funkcji borelowskich jednostajna aproksymacja
i oszacowanie normy operatorów mnożenia dają tę samą
równość. Wszystkie własności całkowe wynikają z tego lematu.
Dla $f(z)=z$ dostajemy $T$. Dla funkcji ciągłych
otrzymujemy wcześniejszy rachunek, gdyż oba homomorfizmy
zgadzają się na $z$ i $\overline z$, a te generują
$C(K)$ jako zamkniętą algebrę.

\textbf{Krok 4: jednoznaczność.}
Jeżeli $F$ jest inną miarą spektralną o podanej własności,
to niezależnie zdefiniowany rachunek z
Lematu~\ref{lem:pvm-integral32} oparty na $F$
przyporządkowuje wielomianom
w $z,\overline z$ operatory $p(T,T^*)$.
Z twierdzenia Stone'a--Weierstrassa to samo zachodzi
dla wszystkich $f\in C(K)$. Dla dowolnych $x,y\in H$
miary zespolone $\langle E(\cdot)x,y\rangle$ i
$\langle F(\cdot)x,y\rangle$ mają więc identyczne całki
z wszystkich funkcji ciągłych. Każda z nich jest regularna:
wynika to przez polaryzację ze skończoności i regularności
miar dodatnich $\mu_x$ na zwartym zbiorze metrycznym $K$.
Jednoznaczność w twierdzeniu Riesza--Markowa daje
równość miar, a więc $E(B)=F(B)$ dla każdego zbioru
borelowskiego $B$.
\end{proof}

\begin{uwaga}[Zakres wykorzystanego twierdzenia miarowego]
Dowód twierdzenia spektralnego sprowadza jego część miarową
do \emph{twierdzenia Riesza--Markowa}: każdy dodatni ciągły
funkcjonał na $C(K)$ dla zwartej przestrzeni Hausdorffa $K$
ma jednoznaczną reprezentację przez skończoną regularną miarę
borelowską. Jest to rezultat skalarnej teorii miary, a nie
wniosek z samego twierdzenia Hahna--Banacha. W pozostałych
krokach dowodu konstrukcja miary projekcyjnej i jej jednoznaczność
zostały przeprowadzone jawnie. Nie używaliśmy założenia
ośrodkowości przestrzeni Hilberta.
\end{uwaga}

\subsection{Konsekwencje miary spektralnej}

\begin{twierdzenie}[Nośnik, widmo i wartości własne]
\label{thm:spectral-support32}
Niech $T$ i $E$ będą jak w twierdzeniu~\ref{thm:spectral-normal32}.
Wtedy $\sigma(T)$ jest najmniejszym domkniętym zbiorem
$F\subseteq K$ takim, że $E(F)=I_H$.
Dla każdego $\lambda\in\C$:
$\lambda$ jest wartością własną $T$ wtedy i tylko wtedy,
gdy $E(\{\lambda\}\cap K)\ne0$, a wówczas
$\ker(T-\lambda I_H)=E(\{\lambda\}\cap K)H$.
Dla $\lambda\in\sigma(T)$ i $\varepsilon>0$ istnieje
$\|x\|=1$ takie, że $\|(T-\lambda I_H)x\|<\varepsilon$.
\end{twierdzenie}
\begin{proof}
Jeśli $U\subseteq\C$ jest otwarty i $U\cap K\ne\varnothing$,
wybieramy niezerową funkcję $f\in C(K)$ o nośniku
zawartym w $U\cap K$ (z odpowiednią interpretacją
nośnika względnego w $K$).
Izometria rachunku ciągłego daje $f(T)\ne0$.
Gdyby $E(U\cap K)=0$, tożsamość
$f(T)=f(T)E(U\cap K)$ dawałaby sprzeczność.
Każdy względnie otwarty niepusty podzbiór $K$
ma zatem niezerową projekcję spektralną. Stąd
$K$ jest nośnikiem $E$.

Jeśli $Tx=\lambda x$, to
$0=\|(T-\lambda I_H)x\|^2
=\int_K|z-\lambda|^2\,\mathrm d\mu_x(z)$,
więc dla każdego $n\ge1$ zbiór
$\{z\in K:|z-\lambda|\ge1/n\}$ ma $\mu_x$-miarę zero.
Ich suma to $K\setminus\{\lambda\}$, stąd
$\|(I_H-E(\{\lambda\}\cap K))x\|^2=0$,
czyli $x=E(\{\lambda\}\cap K)x$.
Odwrotnie, na obrazie tego rzutu operator $T$ działa
jako mnożenie przez $\lambda$.
Wybierając niezerowy wektor z
$E(K\cap B(\lambda,\varepsilon/2))H$ i normalizując go,
otrzymujemy wymagany wektor przybliżony dzięki wzorowi
na normę z twierdzenia~\ref{thm:spectral-normal32}:
$\|(T-\lambda I_H)x\|\le\varepsilon/2<\varepsilon$.
\end{proof}

\begin{twierdzenie}[Borelowski rachunek funkcyjny i zbieżność]
\label{thm:borel-convergence32}
Niech $T\in\mathcal L(H,H)$ będzie normalny z miarą spektralną $E$
na $K=\sigma(T)$. Dla ograniczonej funkcji borelowskiej
$f:K\to\C$ operator $f(T)$ zależy jedynie od wartości $f$
poza zbiorem borelowskim $N$ z $E(N)=0$, a
\[
 \|f(T)\|=\inf\{M\ge0:E(\{z\in K:|f(z)|>M\})=0\}.
\]
Jeżeli $f_n$ są ograniczonymi funkcjami borelowskimi,
$\sup_n\|f_n\|_\infty<\infty$ oraz $f_n(z)\to f(z)$
dla ograniczonej borelowskiej funkcji $f$ poza zbiorem
$E$-zerowym, to $f_n(T)x\to f(T)x$
w normie $H$ dla każdego $x\in H$.
W szczególności $0\le f\le g$ pociąga $f(T)\le g(T)$
dla rzeczywistych ograniczonych funkcji borelowskich $f,g$.
\end{twierdzenie}
\begin{proof}
Jeśli $f=g$ poza zbiorem $N$ z $E(N)=0$, to
$\mu_x(N)=0$ dla wszystkich $x$, a wzór na normę
zastosowany do $f-g$ daje $f(T)=g(T)$.
W twierdzeniu o zbieżności zakładamy, że $f$ jest
ograniczona i borelowska; ewentualnie można przyjąć
ją równą zero na wyjątkowym zbiorze borelowskim.
Dla $x\in H$ wzór na normę daje
$\|(f_n(T)-f(T))x\|^2
=\int_K|f_n(z)-f(z)|^2\,\mathrm d\mu_x(z)$.
Zbieżność zmajoryzowana dowodzi tezy o zbieżności.
Jeżeli $E(\{|f|>M\})=0$, to całki normowe dają
$\|f(T)\|\le M$. Odwrotnie, gdy $E(\{|f|>M\})\ne0$,
możemy wybrać wektor jednostkowy z obrazu tego rzutu;
wtedy $\mu_x$ jest skupiona na $\{|f|>M\}$ i ma masę
$1$. Całka z $|f|^2-M^2$ jest ściśle dodatnia:
zbiór, na którym funkcja ta jest dodatnia, jest sumą
zbiorów, na których jest nie mniejsza niż $1/n$,
a przynajmniej jeden ma dodatnią miarę. Stąd
$\|f(T)x\|>M$. Przechodząc do infimum,
otrzymujemy wzór na normę.
Pozytywność wynika z
$\langle f(T)x,x\rangle=\int_Kf\,\mathrm d\mu_x$.
\end{proof}

\begin{wniosek}[Operatory samosprzężone, dodatnie i unitarne]
\label{cor:spectral-types32}
Niech $T\in\mathcal L(H,H)$ będzie normalny.
Wówczas $T=T^*$ wtedy i tylko wtedy, gdy
$\sigma(T)\subseteq\R$; jest dodatni wtedy i tylko wtedy, gdy
$\sigma(T)\subseteq[0,\infty)$, a unitarny wtedy i tylko wtedy,
gdy $\sigma(T)\subseteq\{z\in\C:|z|=1\}$.
Dla samosprzężonego $T$ projekcje
$E(( -\infty,t]\cap\sigma(T))$, $t\in\R$, tworzą
monotoniczną rodzinę rzutów opisującą rozkład spektralny.
\end{wniosek}
\begin{proof}
Z izometrii rachunku ciągłego wynika, że $T-T^*=0$
dokładnie wtedy, gdy $z-\overline z=0$ na $K$.
Podobnie dla normalnego $T$ oba iloczyny $T^*T$ i $TT^*$
odpowiadają funkcji $|z|^2$, więc unitarność jest
równoważna warunkowi $|z|=1$ na $K$.
Jeśli $K\subseteq[0,\infty)$, to całka z funkcji
$z$ względem każdej $\mu_x$ jest nieujemna, zatem $T\ge0$.
Odwrotnie, gdy $T\ge0$, jest samosprzężony, więc $K\subseteq\R$.
Jeśli $\lambda\in K$ i $\lambda<0$, wybieramy względnie
otwarte otoczenie $V$ punktu $\lambda$ w $K$ zawarte
w $(-\infty,\lambda/2)$. Z Twierdzenia~\ref{thm:spectral-support32}
mamy $E(V)\ne0$. Dla jednostkowego $x\in E(V)H$ miara
$\mu_x$ jest skupiona na $V$ i ma masę $1$, więc
$\langle Tx,x\rangle\le\lambda/2<0$, sprzeczność.
Dla operatora samosprzężonego monotoniczność wynika z
$E(B)E(C)=E(B)$ dla $B\subseteq C$.
\end{proof}

\begin{przyklad}[Widmo ciągłe operatora mnożenia]
\label{ex:mult-spectrum32}
W zespolonej przestrzeni Hilberta $H=L^2([0,1],\lambda_1)$
niech $(Tf)(t)=tf(t)$. Operator jest samosprzężony,
bo mnożnik jest rzeczywisty. Dla $\lambda\notin[0,1]$
odwrotność $\lambda I-T$ jest mnożeniem przez
ograniczoną funkcję $(\lambda-t)^{-1}$.
Dla $\lambda\in[0,1]$ i $\varepsilon>0$ zbiór
$A=[0,1]\cap(\lambda-\varepsilon/2,\lambda+\varepsilon/2)$
ma dodatnią miarę. Funkcja
$f=\chi_A/\sqrt{\lambda_1(A)}$ ma normę $1$ i
$\|(\lambda I-T)f\|\le\varepsilon/2$, co wyklucza
istnienie ograniczonej odwrotności.
Zatem $\sigma(T)=[0,1]$. Miara spektralna ma postać
$(E(B)f)(t)=\chi_{B\cap[0,1]}(t)f(t)$.
Rzuty mnożenia przez funkcje charakterystyczne
spełniają aksjomaty miary spektralnej, w tym silną
addytywność z twierdzenia o zbieżności zmajoryzowanej.
Całka funkcji $z$ jest mnożeniem przez $t$ przez
jednostajną aproksymację funkcjami prostymi, więc
jednoznaczność miary spektralnej uzasadnia podany wzór.
Każda pojedyncza wartość $\lambda\in[0,1]$ ma miarę Lebesgue'a
równą $0$, więc $E(\{\lambda\})=0$ i $T$ nie ma wartości
własnych. Dowodzi to, że nawet dla operatora samosprzężonego
pełny rozkład na wektory własne nie musi istnieć.
\end{przyklad}

\begin{przyklad}[Nieośrodkowa przestrzeń Hilberta]
\label{ex:nonsep-spectrum32}
Niech $I$ będzie dowolnym niepustym zbiorem i $H=\ell^2(I)$.
Dla ograniczonej funkcji $a:I\to\C$ definiujemy
$Te_i=a(i)e_i$ na kanonicznej bazie ortonormalnej.
Współrzędnie $(Tx)_i=a(i)x_i$, więc
$\|Tx\|\le\sup_i|a(i)|\|x\|$ i $T^*$ jest mnożeniem
przez $\overline a$. Stąd normalność.
Poza $\overline{a(I)}$ odwrotność $\lambda I-T$ jest
mnożeniem przez ograniczoną funkcję $i\mapsto(\lambda-a(i))^{-1}$.
Dla $\lambda\in\overline{a(I)}$ istnieją indeksy $i_n$
z $|a(i_n)-\lambda|<1/n$, więc
$\|(\lambda I-T)e_{i_n}\|<1/n$ wyklucza ograniczoną
odwrotność. Zatem $\sigma(T)=\overline{a(I)}$.
Rzut $E(B)$ ma współrzędne $\chi_{a^{-1}(B)}(i)x_i$;
mnożenie tych funkcji charakterystycznych daje aksjomaty
algebraiczne, a zbieżność zmajoryzowana po nośniku
współrzędnych $x$ daje silną addytywność.
Całka funkcji $z$ daje $T$ przez aproksymację prostą,
więc z jednoznaczności miary jest to rzut na domkniętą
rozpiętość $\{e_i:a(i)\in B\}$. Każdy wektor ma co najwyżej przeliczalnie wiele
niezerowych współczynników: dla każdego $n\ge1$ zbiór
$\{i:|x_i|\ge1/n\}$ jest skończony, a suma tych zbiorów
obejmuje wszystkie niezerowe współczynniki.
Sam zbiór indeksów $I$ może być nieprzeliczalny.
\end{przyklad}

\subsection{Rodzina spektralna i rozkład polarny}

\begin{twierdzenie}[Rodzina spektralna operatora samosprzężonego]
\label{thm:spectral-family32}
Niech $T\in\mathcal L(H,H)$ będzie samosprzężony, $E$ jego
miarą spektralną na $K=\sigma(T)\subseteq\R$, a
$P_t=E(K\cap(-\infty,t])$ dla $t\in\R$.
Wtedy $P_sP_t=P_{\min(s,t)}$, $P_s\le P_t$ dla $s\le t$,
$\lim_{t\to-\infty}P_tx=0$ i
$\lim_{t\to+\infty}P_tx=x$ dla wszystkich $x\in H$.
Dla każdego $t\in\R$ i każdego ciągu $t_n\downarrow t$
zachodzi $P_{t_n}x\to P_tx$ w normie; dla $t_n<t$
i $t_n\uparrow t$ granicą jest $E(K\cap(-\infty,t))x$.
Ponadto
\[
 \langle Tx,y\rangle=\int_{K}t\,\mathrm d\mu_{x,y}(t)
 \quad(x,y\in H),
\]
a różnica $P_t-E(K\cap(-\infty,t))=E(K\cap\{t\})$
jest rzutem na $\ker(T-tI_H)$; jej rząd (być może
nieskończony) jest krotnością geometryczną wartości własnej $t$.
\end{twierdzenie}
\begin{proof}
Monotoniczność i równość iloczynów wynikają z własności
$E(B\cap C)=E(B)E(C)$. Ponieważ $K$ jest zwarty,
dla dostatecznie małego $t$ mamy $P_t=0$, a dla dostatecznie
dużego $t$ mamy $P_t=I_H$.
Jeżeli $t_n\downarrow t$, to zbiory $K\cap(-\infty,t_n]$
maleją do $K\cap(-\infty,t]$. Dla każdego $x\in H$
\[
 \|(P_{t_n}-P_t)x\|^2
 =\mu_x(K\cap(t,t_n])\longrightarrow0
\]
na mocy ciągłości miary skończonej z góry.
Analogicznie dla $t_n<t$ i $t_n\uparrow t$ zbiory
$K\cap(-\infty,t_n]$ rosną do $K\cap(-\infty,t)$.
Wzór na parowanie jest szczególnym przypadkiem
rachunku spektralnego, a ostatnia równość wynika
z addytywności miary i z twierdzenia~\ref{thm:spectral-support32}.
\end{proof}

\begin{twierdzenie}[Rozkład polarny operatora ograniczonego]
\label{thm:polar-decomp32}
Niech $H$ będzie zespoloną przestrzenią Hilberta i
$T\in\mathcal L(H,H)$ dowolnym operatorem (niekoniecznie normalnym).
Istnieje dokładnie jedna \emph{częściowa izometria}
$U\in\mathcal L(H,H)$, której przestrzenią początkową jest
$(\ker T)^\perp$ i przestrzenią końcową
$\overline{\operatorname{Ran}T}$, taka że
$T=U|T|$, gdzie $|T|=(T^*T)^{1/2}$ jest jedynym dodatnim
pierwiastkiem z $T^*T$. Co więcej, $\ker U=\ker T$
i $U^*U$ jest rzutem ortogonalnym na $(\ker T)^\perp$.
Dla operatora normalnego $T$ czynniki $U$ i $|T|$ komutują.
\end{twierdzenie}
\begin{proof}
Ponieważ $T^*T\ge0$, ciągły rachunek funkcyjny
dostarcza dodatniego operatora $|T|$ spełniającego
$|T|^2=T^*T$. Uzasadnijmy jednoznaczność dodatniego
pierwiastka dowolnego dodatniego operatora $A$.
Jeśli $B\ge0$ i $B^2=A$, wybieramy wielomiany $p_n$
zbiegające jednostajnie do $\sqrt{\cdot}$ na
$[0,\|A\|]$ (dla $A=0$ mamy $\|B\|^2=\|B^2\|=0$).
Ponieważ $\sigma(B)\subseteq[0,\|B\|]$ oraz
$\|B\|^2=\|A\|$, rachunek ciągły dla $B$ daje
$p_n(B^2)\to B$: funkcje $p_n(t^2)$ zbiegają
jednostajnie do $t$ na $\sigma(B)$.
Rachunek dla $A$ daje jednocześnie $p_n(A)\to A^{1/2}$,
więc $B=A^{1/2}$. Wówczas
$\||T|x\|^2=\langle T^*Tx,x\rangle=\|Tx\|^2$, więc
$\ker|T|=\ker T$ oraz
$\overline{\operatorname{Ran}|T|}=(\ker T)^\perp$.
W ostatniej równości użyliśmy
$(\operatorname{Ran}|T|)^\perp=\ker|T|^*=\ker|T|$.
Na $\operatorname{Ran}|T|$ definiujemy $U_0(|T|x)=Tx$.
Jest to dobrze określona izometria, którą można przedłużyć
przez ciągłość na $(\ker T)^\perp$.
Na $\ker T$ kładziemy $U=0$. Otrzymujemy częściową
izometrię o wskazanych przestrzeniach, normie co najwyżej $1$
i z własnością $U|T|=T$. Jej obraz na przestrzeni
początkowej jest domknięty (izometria z przestrzeni
zupełnej) i równy $\overline{\operatorname{Ran}T}$.
Dla rozkładu $x=x_1+x_0$ z $x_1\perp\ker T$,
$x_0\in\ker T$, mamy $Ux=Ux_1$, zatem $\ker U=\ker T$.
Dla $y=y_1+y_0$ analogicznie
$\langle U^*Ux,y\rangle=\langle x_1,y_1\rangle
=\langle x_1,y\rangle$, więc $U^*U$ jest wskazanym rzutem.
Każda inna częściowa izometria spełniająca warunki tezy
musi zgadzać się z $U$ na gęstym zbiorze
$\operatorname{Ran}|T|$ w przestrzeni początkowej oraz znikać
na $\ker T$, co dowodzi jednoznaczności.
Jeżeli $T$ jest normalny, w modelu spektralnym
$|T|$ odpowiada mnożeniu przez $|z|$, a $U$ przez
funkcję $u(z)=z/|z|$ dla $z\ne0$ i $u(0)=0$:
operator $u(T)$ znika na $\ker T$, jest izometrią na
jej dopełnieniu z wzoru na normę i spełnia
$u(T)|T|=T$; jednoznaczność identyfikuje go z $U$.
Ponieważ funkcje $u(z)$ i $|z|$ komutują punktowo,
operatory $U$ i $|T|$ również komutują.
\end{proof}

\begin{przyklad}[Rozkład polarny macierzy]
Dla $T=\left(\begin{smallmatrix}0&2\\0&0\end{smallmatrix}\right)$
na $\C^2$ mamy
$|T|=\operatorname{diag}(0,2)$ oraz
$U=\left(\begin{smallmatrix}0&1\\0&0\end{smallmatrix}\right)$.
Zatem $T=U|T|$, a $U$ jest częściową izometrią;
nie jest operatorem unitarnym, gdyż $U^*U\ne I$.
\end{przyklad}

\subsection{Operator normalny zwarty jako przypadek dyskretny}

\begin{twierdzenie}[Twierdzenie spektralne dla zwartego operatora normalnego]
\label{thm:compact-normal32}
Niech $H$ będzie niezerową zespoloną przestrzenią Hilberta,
a $T\in\mathcal L(H,H)$ operatorem zwartym i normalnym.
Niezerowe wartości widma tworzą zbiór co najwyżej przeliczalny,
każda z nich jest wartością własną o skończonej krotności,
a jedynym możliwym punktem skupienia tego zbioru jest $0$.
Przestrzeń $H$ jest sumą ortogonalną przestrzeni własnych $T$,
włącznie z $\ker T$. Dla różnych niezerowych wartości
własnych $\lambda$ oznaczmy przez $P_\lambda$ rzut na
$\ker(T-\lambda I_H)$. Wtedy
$T=\sum_{\lambda\in\sigma(T)\setminus\{0\}}
\lambda P_\lambda$, przy czym szereg jest zbieżny w normie
operatorowej (sumy rozumiemy po skończonych podzbiorach indeksów).
\end{twierdzenie}
\begin{proof}
Dla nieskończenie wymiarowego $H$ stosujemy
Twierdzenie~\ref{thm:riesz-schauder30}: niezerowe punkty
widma są izolowanymi wartościami własnymi skończonej
krotności i mają jedyny możliwy punkt skupienia w zerze.
Zbiory punktów o module co najmniej $1/n$ są skończone,
więc ich suma jest co najwyżej przeliczalna.
Dla skończenie wymiarowego $H$ widmo jest skończonym
zbiorem pierwiastków wielomianu charakterystycznego,
a wszystkie przestrzenie własne mają skończony wymiar.
Z twierdzenia~\ref{thm:spectral-support32} każda niepusta
względnie otwarta część widma ma niezerową projekcję spektralną.
Poza zerem widmo jest zatem sumą pojedynczych punktów,
a odpowiadające im projekcje są ortogonalne.
Dopełnienie ortogonalne ich obrazów jest obrazem
$E(K\cap\{0\})$ i równa się $\ker T$.
Dla skończonego zbioru niezerowych wartości własnych $F$
wzór z Twierdzenia~\ref{thm:borel-convergence32} daje
\[
\left\|T-\sum_{\lambda\in F}\lambda P_\lambda\right\|
=\sup\bigl(\{0\}\cup\{|\lambda|:\lambda\in K\setminus(F\cup\{0\})\}\bigr).
\]
Istotnie, każda pozostająca niezerowa wartość własna
ma niezerową projekcję, a reszta funkcji spektralnej
jest zerowa na $F$ i w zerze. Wzór obejmuje też
przypadek, gdy nie pozostaje żadna niezerowa wartość.
Ta wielkość dąży do zera, gdy $F$ obejmuje wszystkie
wartości własne o module większym od ustalonego $\varepsilon>0$.
W rezultacie otrzymujemy zbieżność w normie operatorowej.
\end{proof}

\begin{uwaga}[Granica twierdzenia]
Dla operatorów nienormalnych nie można w ogólności znaleźć
miary spektralnej projekcyjnej reprezentującej operator.
Na przykład niezerowa macierz
$\left(\begin{smallmatrix}0&1\\0&0\end{smallmatrix}\right)$
ma widmo $\{0\}$, lecz nie jest zerowa; gdyby miała reprezentację
$T=\int z\,\mathrm dE(z)$ z miarą spektralną na tym widmie,
byłaby równa $0$. Rozszerzenie rachunku spektralnego na
operatory nieograniczone opracujemy w \hyperref[sec:unbounded33]{rozdziale~33}.
\end{uwaga}

% Wersja źródłowa: audyt II, jednostka 33
\begingroup
\setlength{\parskip}{3pt}
\setcounter{section}{32}
\noindent\textbf{Wyniki wykorzystywane.} Z poprzednich rozdziałów korzystamy z twierdzenia o wykresie domkniętym, teorii operatorów sprzężonych na przestrzeniach Hilberta, twierdzenia Plancherela i twierdzenia spektralnego dla operatorów ograniczonych normalnych wraz z rachunkiem borelowskim. Dowody twierdzenia spektralnego dla operatorów nieograniczonych, twierdzenia von Neumanna o rozszerzeniach i twierdzenia Stone'a podajemy w tym rozdziale.\par\medskip
\section{Operatory nieograniczone na przestrzeniach Hilberta}
\label{sec:unbounded33}

W \hyperref[sec:spectral-hilbert32]{rozdziale~32} skonstruowaliśmy miarę spektralną
operatora \emph{ograniczonego} i normalnego. Operator różniczkowania na
$L^2$, operator mnożenia przez zmienną oraz generatory jednoparametrowych
grup unitarnych są na ogół nieograniczone. Zasadnicza różnica polega na tym,
że \emph{dziedzina operatora jest częścią jego definicji}. W całym rozdziale
$H$ oznacza niezerową zespoloną przestrzeń Hilberta, a iloczyn skalarny jest liniowy
w pierwszym argumencie. Zapis $\mathcal L(H)$ oznacza operatory ograniczone
na całym $H$.

\subsection{Dziedziny, wykresy, operatory sprzężone}

\begin{definicja}[Operator liniowy i jego wykres]
\label{def:unbounded-operator33}
Operatorem liniowym w $H$ nazywamy liniowe odwzorowanie
$A:D(A)\to H$, gdzie $D(A)\subseteq H$ jest podprzestrzenią liniową.
Jego \emph{wykresem} jest $G(A)=\{(x,Ax):x\in D(A)\}\subseteq H\times H$.
Operator jest \emph{domknięty}, gdy $G(A)$ jest domknięty, oraz
\emph{domykalny}, gdy domknięcie $G(A)$ jest wykresem pewnego operatora;
operator ten nazywamy \emph{domknięciem} $\overline A$.
Dla operatorów $A,B$ piszemy $A\subseteq B$, jeżeli
$D(A)\subseteq D(B)$ i $Bx=Ax$ dla $x\in D(A)$.
\end{definicja}

\begin{stwierdzenie}[Kryteria domkniętości i norma wykresowa]
\label{prop:graph33}
Niech $A:D(A)\subseteq H\to H$ będzie liniowy.
Wówczas $A$ jest domknięty wtedy i tylko wtedy, gdy dla każdego
ciągu $x_n\in D(A)$ z $x_n\to x$ i $Ax_n\to y$ w $H$ zachodzi
$x\in D(A)$ i $Ax=y$. Ponadto $A$ jest domknięty wtedy i tylko wtedy,
gdy $D(A)$ z normą $\|x\|_A=(\|x\|^2+\|Ax\|^2)^{1/2}$ jest
przestrzenią Hilberta z iloczynem skalarnym
$\langle x,y\rangle_A=\langle x,y\rangle+\langle Ax,Ay\rangle$.
Jeżeli $D(A)=H$, to operator domknięty jest ograniczony.
\end{stwierdzenie}
\begin{proof}
Pierwsza równoważność jest ciągowym kryterium domkniętości
podprzestrzeni przestrzeni metrycznej $H\times H$.
Odwzorowanie $D(A)\ni x\mapsto(x,Ax)\in G(A)$ jest izometrią
przy wskazanych iloczynach skalarnych, a $G(A)$ jest zupełny dokładnie
wtedy, gdy jest domknięty. Ostatnia teza wynika z twierdzenia
o wykresie domkniętym dla przestrzeni Banacha
(\hyperref[sec:frechet-open-graph]{rozdział~11}).
\end{proof}

\begin{definicja}[Zbiór rezolwenty operatora domkniętego]
Niech $A:D(A)\subseteq H\to H$ będzie operatorem domkniętym.
Do jego \emph{zbioru rezolwenty} $\rho(A)$ należą te $z\in\C$,
dla których $A-zI:D(A)\to H$ jest bijekcją, a jej
odwrotność $(A-zI)^{-1}:H\to H$ jest ograniczona.
\emph{Widmo} to $\sigma(A)=\C\setminus\rho(A)$.
Dla operatora domkniętego ograniczoność odwrotności wynika już
z jej istnienia i twierdzenia o wykresie domkniętym,
ale podajemy ją jawnie dla odróżnienia od operatorów niedomkniętych.
\end{definicja}

\begin{definicja}[Operator sprzężony]
\label{def:adjoint33}
Niech $A:D(A)\subseteq H\to H$ będzie liniowy i gęsto określony,
czyli $\overline{D(A)}=H$. Definiujemy
\[
 D(A^*)=\{y\in H:\text{istnieje }z\in H\text{ takie, że }
 \langle Ax,y\rangle=\langle x,z\rangle\ (x\in D(A))\},
 \qquad A^*y=z.
\]
Wektor $z$ jest jednoznaczny na mocy gęstości $D(A)$.
Operator $A$ jest \emph{symetryczny}, gdy $A\subseteq A^*$,
a \emph{samosprzężony}, gdy $A=A^*$, włącznie z równością dziedzin.
\end{definicja}

\begin{twierdzenie}[Sprzężenie i domykalność]
\label{thm:closable-adjoint33}
Niech $A:D(A)\subseteq H\to H$ będzie gęsto określony.
Wówczas $A^*$ jest operatorem domkniętym (jego dziedzina
nie musi być gęsta). Operator $A$ jest domykalny wtedy i tylko wtedy,
gdy $D(A^*)$ jest gęsta w $H$. W tym przypadku
$\overline A=A^{**}$. Każdy gęsto określony operator symetryczny
jest domykalny, a jego domknięcie również jest symetryczne.
\end{twierdzenie}
\begin{proof}
Jeżeli $y_n\to y$ i $A^*y_n\to z$, to dla każdego $x\in D(A)$
$\langle Ax,y\rangle=\lim_n\langle Ax,y_n\rangle
=\lim_n\langle x,A^*y_n\rangle=\langle x,z\rangle$;
stąd $y\in D(A^*)$ i $A^*y=z$.

Sprawdźmy kryterium domykalności. Dla wektora $(u,v)\in H\times H$
warunek $(u,v)\perp G(A)$ jest równoważny
$\langle u,x\rangle+\langle v,Ax\rangle=0$ dla wszystkich
$x\in D(A)$. Po sprzężeniu tej równości otrzymujemy
$v\in D(A^*)$ oraz $u=-A^*v$. Zatem
$G(A)^\perp=\{(-A^*v,v):v\in D(A^*)\}$.
W szczególności $(0,w)\in\overline{G(A)}$ wtedy i tylko wtedy,
gdy $\langle w,v\rangle=0$ dla wszystkich $v\in D(A^*)$.
Liniowa podprzestrzeń $\overline{G(A)}$ jest wykresem operatora
dokładnie wtedy, gdy nie zawiera niezerowych wektorów $(0,w)$.
Otrzymujemy równoważność z gęstością $D(A^*)$.
Gdy ta gęstość zachodzi, dla $(x,y)\in H\times H$ mamy
\[
 (x,y)\in G(A)^{\perp\perp}
 \ \Longleftrightarrow\
 \langle y,v\rangle=\langle x,A^*v\rangle
 \quad(v\in D(A^*)).
\]
Po sprzężeniu jest to dokładnie warunek $x\in D(A^{**})$
i $A^{**}x=y$. Zatem $G(A^{**})=\overline{G(A)}$.
Jeżeli $A$ jest symetryczny, to $D(A)\subseteq D(A^*)$, a więc
$D(A^*)$ jest gęsta; przejście do granic w tożsamości
$\langle Ax,y\rangle=\langle x,Ay\rangle$ zachowuje symetrię.
\end{proof}

\begin{uwaga}[Operator ograniczony na niewłaściwej dziedzinie]
Nawet operator zerowy na gęstej właściwej podprzestrzeni $D\subsetneq H$
nie jest samosprzężony: jego operator sprzężony jest zerowy na całym
$H$, natomiast jego domknięcie jest operatorem zerowym na $H$.
Nie wystarczy zatem sprawdzać samego wzoru działania operatora.
\end{uwaga}

\subsection{Symetria, samosprzężoność i rezolwenta}

\begin{lemat}[Oszacowanie poza osią rzeczywistą]
\label{lem:symmetric-resolvent33}
Niech $A$ będzie gęsto określonym operatorem symetrycznym w $H$.
Dla $z=a+ib\in\C$ z $b\ne0$ i $x\in D(A)$ zachodzi
$\|(A-zI)x\|\ge |b|\|x\|$. Jeśli $A$ jest ponadto domknięty,
$\operatorname{Ran}(A-zI)$ jest podprzestrzenią domkniętą.
\end{lemat}
\begin{proof}
Ponieważ $\langle Ax,x\rangle$ jest rzeczywiste,
$|\operatorname{Im}\langle(A-zI)x,x\rangle|=|b|\|x\|^2$.
Nierówność Cauchy'ego--Schwarza daje oszacowanie.
Gdy $A$ jest domknięty, oszacowanie implikuje, że
$(A-zI)x_n\to y$ pociąga zbieżność $x_n\to x$;
ponadto $Ax_n=(A-zI)x_n+zx_n\to y+zx$, więc
domkniętość $A$ daje $x\in D(A)$ i $(A-zI)x=y$.
\end{proof}

\begin{twierdzenie}[Kryterium samosprzężoności]
\label{thm:selfadjoint-range33}
Niech $A$ będzie gęsto określonym, symetrycznym operatorem w $H$.
Wówczas następujące warunki są równoważne:
\begin{enumerate}[label=\textup{(\roman*)}]
\item $A$ jest samosprzężony;
\item $\operatorname{Ran}(A+iI)=H$ i $\operatorname{Ran}(A-iI)=H$.
\end{enumerate}
Jeśli warunki te zachodzą, dla każdego $z\in\C\setminus\R$
operator $A-zI:D(A)\to H$ jest bijekcją, a jego odwrotność
$R(z,A)=(A-zI)^{-1}\in\mathcal L(H)$ spełnia
$\|R(z,A)\|\le1/|\operatorname{Im}z|$.
\end{twierdzenie}
\begin{proof}
Jeżeli $A=A^*$, to $A$ jest domknięty. Z lematu
\ref{lem:symmetric-resolvent33} oba obrazy $A\pm iI$ są domknięte.
Ogólnie, dla $z\in\C$, z definicji sprzężenia wynika
\[
 \operatorname{Ran}(A-zI)^\perp=\ker(A^*-\overline z I).
\]
Istotnie, $\langle (A-zI)x,y\rangle=0$ dla każdego $x\in D(A)$
jest równoważne $\langle Ax,y\rangle=\langle x,\overline z y\rangle$.
Ortogonalne dopełnienia obrazów $A\pm iI$ wynoszą więc odpowiednio
$\ker(A^*\mp iI)$, a te są zerowe: dla $y\in D(A)$
równość $(A\mp iI)y=0$ i oszacowanie z lematu wymuszają $y=0$.
Obrazy są więc również gęste, a zatem równe $H$.

Odwrotnie, załóżmy surjektywność obu operatorów.
Wówczas $A$ jest domknięty: jeśli $x_n\to x$ i $Ax_n\to y$,
to $(A+iI)x_n\to y+ix$; wybierając $u\in D(A)$
z $(A+iI)u=y+ix$ i stosując oszacowanie z lematu,
otrzymujemy $x_n\to u$, więc $x=u\in D(A)$ i $Ax=y$.
Dla $y\in D(A^*)$ wybieramy $u\in D(A)$ z
$(A+iI)u=(A^*+iI)y$. Wówczas
$(A^*+iI)(y-u)=0$, a
$\ker(A^*+iI)=\operatorname{Ran}(A-iI)^\perp=\{0\}$.
Stąd $y=u\in D(A)$ i $A=A^*$.

Dla dowolnego $z\notin\R$ ponawiamy argument z domkniętością
obrazu i jego gęstością, używając $A=A^*$ oraz oszacowania
z lematu. Oszacowanie normy odwrotności jest natychmiastowe.
\end{proof}

\begin{przyklad}[Mnożenie przez funkcję rzeczywistą]
\label{ex:mult-unbounded33}
Niech $(X,\Sigma,\mu)$ będzie przestrzenią z miarą, a
$m:X\to\R$ funkcją mierzalną, skończoną prawie wszędzie.
W przestrzeni Hilberta $H=L^2(X,\mu)$ operator
$(M_mf)(x)=m(x)f(x)$ z dziedziną
$D(M_m)=\{f\in L^2(X,\mu):mf\in L^2(X,\mu)\}$
jest gęsto określony i samosprzężony. Istotnie, funkcje
$f\chi_{\{|m|\le n\}}$ należą do $D(M_m)$ i zbiegają do $f$
w $L^2$. Dla każdego $g\in H$ funkcje
$f_\pm(x)=g(x)/(m(x)\pm i)$ należą do $D(M_m)$,
ponieważ $|m/(m\pm i)|\le1$, i spełniają
$(M_m\pm iI)f_\pm=g$. Operator jest symetryczny, więc
poprzednie twierdzenie daje samosprzężoność.
Dla $X=\R$ z miarą Lebesgue'a $\lambda_1$ i $m(x)=x$
operator jest nieograniczony: wektory $f_n=\chi_{[n,n+1]}$
należą do dziedziny, mają normę $1$, a $\|M_m f_n\|\ge n$.
\end{przyklad}

\subsection{Miary spektralne i rachunek funkcji nieograniczonych}

\begin{lemat}[Całkowanie nieograniczonych funkcji względem miary spektralnej]
\label{lem:unbounded-integral33}
Niech $K\subseteq\C$ będzie niepustym zbiorem borelowskim
(niekoniecznie zwartym), $H$ zespoloną przestrzenią Hilberta,
a $E:\mathcal B(K)\to\mathcal L(H)$ miarą spektralną:
$E(\varnothing)=0$, $E(K)=I_H$, każde $E(B)$ jest rzutem ortogonalnym,
$E(B\cap C)=E(B)E(C)$ i $E$ jest przeliczalnie addytywna
w silnej topologii operatorowej. Niech $f:K\to\C$ będzie
funkcją borelowską, niekoniecznie ograniczoną.
Wówczas istnieje gęsto określony operator domknięty $f(E)$
o dziedzinie
\[
 D(f(E))=\left\{x\in H:\int_K|f(z)|^2\,\mathrm d\mu_x(z)<\infty\right\}.
\]
Dla $x\in D(f(E))$ zachodzi
\[
 \|f(E)x\|^2=\int_K|f(z)|^2\,\mathrm d\mu_x(z),
\]
gdzie $\mu_x(B)=\langle E(B)x,x\rangle$.
Zachodzi $f(E)^*=\overline f(E)$; w szczególności dla
$f$ rzeczywistej $E$-prawie wszędzie operator $f(E)$ jest
samosprzężony. Dla ograniczonej borelowskiej funkcji $g$
operator $g(E)$ zachowuje $D(f(E))$ oraz
$f(E)g(E)=(fg)(E)$ z równością dziedzin. Ponadto
$g(E)f(E)x=(fg)(E)x$ dla $x\in D(f(E))$.
Dla ograniczonej $f$ konstrukcja jest zwykłym borelowskim
rachunkiem funkcyjnym.
\end{lemat}
\begin{proof}
\textbf{Krok 1: rachunek funkcji ograniczonych.}
Dla funkcji prostej $s=\sum_{j=1}^m c_j\chi_{B_j}$, gdzie
$B_j$ są parami rozłączne, definiujemy
$s(E)=\sum_{j=1}^m c_jE(B_j)$.
Dopełniając rozłączny układ do podziału $K$ (ze współczynnikiem
zero na pozostałej części) i stosując wspólne rozdrobnienie
dwóch podziałów, otrzymujemy niezależność od reprezentacji
funkcji prostej oraz liniowość. Argument jest taki sam jak
w Lemacie~\ref{lem:pvm-integral32}; nie wymaga zwartości $K$.
Ponadto $1(E)=I_H$. Własności rzutów dają
$\|s(E)x\|^2=\int_K|s|^2\,\mathrm d\mu_x$,
$(st)(E)=s(E)t(E)$ i $s(E)^*=\overline s(E)$
dla funkcji prostych. Każdą ograniczoną funkcję borelowską
można jednostajnie aproksymować funkcjami prostymi.
Z nierówności $\|s(E)\|\le\|s\|_\infty$ otrzymujemy
przez granicę w normie operatorowej rachunek wszystkich
ograniczonych funkcji borelowskich, z tymi samymi
własnościami algebraicznymi i normowymi.

\textbf{Krok 2: dziedzina i domkniętość.}
Funkcja $f$ jest skończona punktowo. Dla ograniczonych
funkcji równych poza zbiorem $N$ z $E(N)=0$ równość operatorów
wynika ze wzoru na normę ich różnicy, ponieważ $\mu_x(N)=0$
dla każdego $x\in H$. Po modyfikacji dowolnej funkcji $f$ na
takim zbiorze nie zmieniają się całki $|f|^2$ ani operatory
jej ograniczonych obcięć, a więc również poniższa konstrukcja
dziedziny i granicy.
Kładziemy $B_n=\{z\in K:|f(z)|\le n\}$,
$P_n=E(B_n)$ i $f_n=f\chi_{B_n}$.
Wtedy $P_n\to I_H$ silnie, a $f_n$ są ograniczone.
Jeżeli $x$ należy do przestrzeni $D$ określonej całką w tezie,
to
\[
 \|(f_n(E)-f_m(E))x\|^2
 =\int_K|f_n-f_m|^2\,\mathrm d\mu_x\longrightarrow0.
\]
Definiujemy $f(E)x=\lim_n f_n(E)x$.
Odwrotnie, gdy $f_n(E)x$ ma granicę, to jego normy są
ograniczone, a monotoniczna zbieżność całek
$\int|f_n|^2\,\mathrm d\mu_x$ dowodzi $x\in D$.
Ponieważ $P_nH\subseteq D$ (dla $x\in H$ miara $\mu_{P_nx}$
jest skupiona na $B_n$ i $\int|f|^2\,\mathrm d\mu_{P_nx}\le n^2\|x\|^2$)
i $P_n\to I_H$ silnie, dziedzina $D$ jest gęsta.
Jest ona podprzestrzenią liniową, gdyż równoważnie składa się
z tych $x$, dla których ciąg liniowych odwzorowań $f_n(E)x$
jest zbieżny. Operator graniczny jest więc liniowy.
Tożsamość na normę wynika
z monotonicznej zbieżności całek.

Aby wykazać domkniętość, przypuśćmy, że $x_j\in D$,
$x_j\to x$ oraz $f(E)x_j\to y$. Dla każdego $n$
z rachunku ograniczonego wynika
$P_nf(E)x_j=f_n(E)x_j$.
Przechodząc do granicy przy ustalonym $n$ otrzymujemy
$P_ny=f_n(E)x$, a więc $\|f_n(E)x\|\le\|y\|$.
Monotoniczna zbieżność całek pokazuje, że $x\in D$.
Następnie $f_n(E)x=P_ny\to y$, czyli $f(E)x=y$.

\textbf{Krok 3: operator sprzężony i mnożenie.}
Dla $x\in D(f(E))$ i $y\in D(\overline f(E))$,
przechodząc do granicy w tożsamości dla $f_n(E)$,
dostajemy $\langle f(E)x,y\rangle
=\langle x,\overline f(E)y\rangle$.
Zatem $\overline f(E)\subseteq f(E)^*$.
Jeśli zaś $y\in D(f(E)^*)$ i $z=f(E)^*y$, to
testując tożsamość sprzężenia na $P_nx$ dla dowolnego
$x\in H$ otrzymujemy
$\overline{f_n}(E)y=P_nz$.
Stąd $\|\overline{f_n}(E)y\|\le\|z\|$.
Monotoniczna zbieżność całek implikuje
$y\in D(\overline f(E))$, a granica daje
$\overline f(E)y=z$. Zatem $f(E)^*=\overline f(E)$.
Jeżeli $g$ jest ograniczona, to z rachunku ograniczonego
\[
 \mu_{g(E)x}(C)=\|(\chi_Cg)(E)x\|^2
 =\int_C |g|^2\,\mathrm d\mu_x
 \quad(C\in\mathcal B(K)).
\]
Stąd $\int|f|^2\,\mathrm d\mu_{g(E)x}=\int|fg|^2\,\mathrm d\mu_x$
(przez aproksymację funkcjami prostymi i zbieżność monotoniczną).
W szczególności $D(f(E)g(E))=D((fg)(E))$, a dziedzina ta zawiera
$D(f(E))$, ponieważ prawa całka jest wtedy nie większa niż
$\|g\|_\infty^2\int|f|^2\,\mathrm d\mu_x$.
Uzasadnijmy przy tym ogólną zasadę obcięć: jeżeli $h$ jest
borelowska, $x\in D(h(E))$, a ograniczone funkcje $h_n$
zbiegają do $h$ punktowo i spełniają $|h_n|\le |h|$, to
\[
 \|h_n(E)x-h(E)x\|^2=\int_K|h_n-h|^2\,\mathrm d\mu_x.
\]
Istotnie, dla ustalonego $n$ stosujemy wzór dla funkcji
ograniczonych $h_n$ oraz $h\chi_{\{|h|\le m\}}$ i przechodzimy
z $m$ do nieskończoności. Po stronie operatorowej korzystamy
z definicji $h(E)$, a po stronie całkowej ze zbieżności
zmajoryzowanej przez $4|h|^2$. Kolejne przejście z $n$ do
nieskończoności daje zbieżność w normie $H$.
Stosując tę zasadę do $h=fg$ i $h_n=f_ng$, otrzymujemy
$(f_ng)(E)x\to(fg)(E)x$ dla $x\in D((fg)(E))$. Z drugiej strony $(f_ng)(E)x=f_n(E)g(E)x$,
co dowodzi równości $f(E)g(E)=(fg)(E)$ z równością dziedzin.
Dla $x\in D(f(E))$ ograniczoność $g(E)$ pozwala przejść do granicy
także w $g(E)f_n(E)x$, dając równość podaną w tezie.
W szczególności dla funkcji $f_1,f_2$ i $x$ należącego do obu
dziedzin, rachunek jest addytywny na tym wspólnym zbiorze:
stosujemy wspólne obcięcia $\{|f_1|\le n,\ |f_2|\le n\}$ i zbieżność
zmajoryzowaną przez $2(|f_1|^2+|f_2|^2)$.
Nie oznacza to na ogół równości dziedzin sumy operatorów
i operatora odpowiadającego sumie funkcji.
\end{proof}

\begin{twierdzenie}[Twierdzenie spektralne dla operatora samosprzężonego]
\label{thm:unbounded-spectral33}
Niech $A:D(A)\subseteq H\to H$ będzie gęsto określonym
operatorem samosprzężonym. Istnieje dokładnie jedna miara spektralna
$E:\mathcal B(\R)\to\mathcal L(H)$ taka, że
\[
 D(A)=\left\{x\in H:\int_\R t^2\,\mathrm d\mu_x(t)<\infty\right\},
 \qquad Ax=\int_\R t\,\mathrm dE(t)\,x\quad(x\in D(A)).
\]
Dla każdego $z\in\C\setminus\R$ zachodzi
$(A-zI)^{-1}=\int_\R(t-z)^{-1}\,\mathrm dE(t)$.
Dla dowolnej borelowskiej funkcji $f:\R\to\C$
definiujemy $f(A)$ według
lematu~\ref{lem:unbounded-integral33}; operator ten jest domknięty,
a jego dziedzina jest wyznaczona przez całkę $\int|f|^2\,\mathrm d\mu_x$.
\end{twierdzenie}
\begin{proof}
\textbf{Krok 1: transformacja Cayleya.}
Z twierdzenia~\ref{thm:selfadjoint-range33} operatory
$A\pm iI$ odwzorowują bijektywnie $D(A)$ na $H$.
Definiujemy ograniczony operator
$U=(A-iI)(A+iI)^{-1}=I_H-2i(A+iI)^{-1}$.
Dla $y=(A+iI)x$ symetria daje
$\|Uy\|=\|(A-iI)x\|=\|(A+iI)x\|=\|y\|$.
Surjektywność $A-iI$ pokazuje, że $U$ jest unitarny.
Ponadto $I_H-U=2i(A+iI)^{-1}$ jest różnowartościowy,
a zatem $\ker(I_H-U)=\{0\}$.

\textbf{Krok 2: przeniesienie miary spektralnej.}
Niech $P$ będzie jedyną miarą spektralną operatora
unitarnego $U$ na $\sigma(U)\subseteq\mathbb T$.
Przedłużamy ją na okrąg jednostkowy, kładąc
$P(C):=P(C\cap\sigma(U))$ dla borelowskich $C\subseteq\mathbb T$.
Ponieważ $\ker(I_H-U)=P(\{1\})H=\{0\}$,
mamy $P(\{1\})=0$. Definiujemy rzeczywistą funkcję borelowską
$h(w)=i(1+w)/(1-w)$ dla $w\ne1$ na okręgu jednostkowym.
Jej wartości są rzeczywiste: dla $w=e^{i\theta}\ne1$
liczba $h(w)=-\cot(\theta/2)\in\R$.
Wartość $h(1)$ ustalamy równą $0$, ponieważ $P(\{1\})=0$.
Kładziemy $E(B)=P(h^{-1}(B))$ dla $B\in\mathcal B(\R)$.
Otrzymujemy miarę spektralną na $\R$.

\textbf{Krok 3: identyfikacja operatora.}
Operator $B=h(P)$ z lematu~\ref{lem:unbounded-integral33}
jest samosprzężony. Dla $w\ne1$ zachodzi
$(h(w)+i)^{-1}=(1-w)/(2i)$.
Oznaczmy tę ograniczoną funkcję po prawej przez $g$.
Ponieważ $hg=1-ig$ poza zbiorem $P$-zerowym, iloczyn $hg$
jest ograniczony $P$-prawie wszędzie. Z dowodu lematu
\ref{lem:unbounded-integral33} wynika $g(P)H\subseteq D(B)$
i $(B+iI)g(P)=I_H$, natomiast dla $x\in D(B)$ mamy
$g(P)(B+iI)x=x$. Rachunek spektralny daje więc obustronną odwrotność
$(B+iI)^{-1}=(I_H-U)/(2i)=(A+iI)^{-1}$.
Odwrotności wyznaczają jednoznacznie dziedziny i działania
operatorów, więc $A=B$. Wzór na dziedzinę i całkę
spektralną wynika z lematu i definicji miary obrazowej:
$\mu_x^E(C)=\mu_x^P(h^{-1}(C))$, zatem
$\int_\R t^2\,\mathrm d\mu_x^E=\int_{\mathbb T}|h|^2\,\mathrm d\mu_x^P$.
Dla funkcji prostych całkowanie względem $E$ odpowiada
całkowaniu ich złożeń z $h$ względem $P$; obcięcia i granica
w normie wektora dają identyczność całek nieograniczonych.

\textbf{Krok 4: jednoznaczność.}
Jeżeli $F$ jest inną miarą spektralną reprezentującą $A$,
to ograniczony rachunek borelowski daje
$U=(A-iI)(A+iI)^{-1}=\int_\R (t-i)/(t+i)\,\mathrm dF(t)$.
Jednoznaczność miary spektralnej operatora normalnego
$U$ z \hyperref[sec:spectral-hilbert32]{rozdziału~32}
stosujemy po sprawdzeniu nośnika miary obrazowej $Q$ na $\mathbb T$.
Jeśli $w\notin\sigma(U)$, niech $\delta>0$ będzie takie, że
$\|(U-wI)y\|\ge\delta\|y\|$ dla wszystkich $y\in H$.
Gdyby $Q(\mathbb T\cap B(w,\delta/2))\ne0$, wektor jednostkowy
z obrazu tego rzutu spełniałby $\|(U-wI)y\|\le\delta/2$,
na mocy wzoru całkowego, co jest sprzecznością.
Otwarty zbiór $\mathbb T\setminus\sigma(U)$ jest pokryty
przeliczalnie wieloma takimi otoczeniami (okrąg ma przeliczalną bazę).
Ich suma jest $Q$-zerowa, gdyż wszystkie miary skalarne $\mu_x^Q$
są przeliczalnie addytywne. Zatem $Q(\sigma(U))=I_H$.
Jednoznaczność na $\sigma(U)$ daje $Q=P$ także na $\mathbb T$.
Funkcja ta jest różnowartościowa z $\R$ na
$\mathbb T\setminus\{1\}$, o odwrotności $h$, więc $F=E$.
Dla $z\notin\R$ funkcje $g_z(t)=(t-z)^{-1}$ i $tg_z(t)=1+zg_z(t)$
są ograniczone. Tak jak w kroku 3, $g_z(E)H\subseteq D(A)$,
$(A-zI)g_z(E)=I_H$ i $g_z(E)(A-zI)x=x$ dla $x\in D(A)$.
To dowodzi wzoru na rezolwentę.
\end{proof}

\begin{wniosek}[Widmo i przestrzenie spektralne]
\label{cor:unbounded-spectrum33}
Niech $A$ będzie samosprzężony z miarą spektralną $E$.
Wówczas $\sigma(A)\subseteq\R$ i jest równe nośnikowi $E$,
tj. najmniejszemu domkniętemu zbiorowi $F\subseteq\R$
z $E(F)=I_H$. Dla $\lambda\in\R$ zachodzi
$\ker(A-\lambda I)=E(\{\lambda\})H$.
Jeżeli $\lambda\in\sigma(A)$, to dla każdego $\varepsilon>0$
istnieje $x\in D(A)$ o normie $1$ taki, że
$\|(A-\lambda I)x\|<\varepsilon$.
\end{wniosek}
\begin{proof}
Najpierw istnienie nośnika: suma wszystkich względnie otwartych
zbiorów $E$-zerowych w $\R$ jest sumą przeliczalnie wielu
takich zbiorów, ponieważ $\R$ ma przeliczalną bazę.
Jest więc $E$-zerowa, co wynika z przeliczalnej addytywności
każdej $\mu_x$. Jej dopełnienie jest najmniejszym domkniętym
zbiorem pełnej miary spektralnej. Poza osią rzeczywistą
rezolwenta istnieje z twierdzenia \ref{thm:selfadjoint-range33}.
Dla rzeczywistego $\lambda\notin\operatorname{supp}E$ istnieje $\delta>0$
takie, że $E((\lambda-\delta,\lambda+\delta))=0$.
Kładziemy $g(t)=(t-\lambda)^{-1}$ dla $|t-\lambda|\ge\delta$
oraz $g(t)=0$ w pozostałych punktach. Funkcje $g$ i $tg$
są ograniczone, a $(t-\lambda)g(t)=1$ $E$-prawie wszędzie.
Argument z kroku 3 dowodu twierdzenia \ref{thm:unbounded-spectral33}
daje obustronną odwrotność $g(E)$ operatora $A-\lambda I$. Odwrotnie, jeśli
$\lambda\in\operatorname{supp}E$, to
$E((\lambda-\varepsilon/2,\lambda+\varepsilon/2))\ne0$.
Wybieramy jednostkowy wektor $x$ z obrazu tego rzutu;
jego miara skalarna jest skupiona na tym ograniczonym przedziale,
a więc $x\in D(A)$ i $\|(A-\lambda I)x\|\le\varepsilon/2<\varepsilon$. Operator nie może mieć wtedy
ograniczonej odwrotności. Tożsamość dotycząca jądra wynika
z $\|(A-\lambda I)x\|^2=\int|t-\lambda|^2\,\mathrm d\mu_x(t)$:
zerowa całka oznacza, że $\mu_x$ jest skupiona na $\{\lambda\}$
(rozważamy zbiory $|t-\lambda|\ge1/n$), czyli $E(\{\lambda\})x=x$.
Odwrotnie, każdy wektor z obrazu tego rzutu należy do $D(A)$
i jest zerowany przez $A-\lambda I$.
\end{proof}

\begin{przyklad}[Różne dziedziny tego samego wzoru]
Niech $H=L^2(\R,\lambda_1)$.
Operator $Af=xf$ z dziedziną $D(A)=\{f:xf\in L^2\}$ jest
samosprzężony; jego miara spektralna to
$(E(C)f)(x)=\chi_C(x)f(x)$ (aksjomaty i reprezentacja całkowa
wynikają ze zbieżności zmajoryzowanej i obcięć). Każdy niepusty
otwarty przedział ma dodatnią miarę Lebesgue'a, więc nośnikiem
$E$ jest $\R$ i $\sigma(A)=\R$. Dla operatora
$Bf=-if'$ z dziedziną $H^1(\R)=\{f\in L^2(\R):f'\in L^2(\R)\}$
(pochodna rozumiana w sensie dystrybucyjnym) transformacja Fouriera
z \hyperref[sec:temperowane22]{rozdziału~22} daje w sensie
dystrybucyjnym $\widehat{Bf}(\xi)=\xi\widehat f(\xi)$.
Twierdzenie Plancherela (Twierdzenie~\ref{thm:plancherel22})
i odwrotna transformacja dają równoważność
$f'\in L^2$ oraz $\xi\widehat f\in L^2$; zatem unitarna
transformacja Fouriera przenosi dokładnie $D(B)$ na $D(M_\xi)$.
Wobec tego $B$ jest
samosprzężony i ma widmo $\R$. Tymczasem ograniczenie
operatora $-if'$ do przestrzeni $C_c^\infty(\R)$
jest tylko symetryczne i domykalne, a nie domknięte.
Symetria wynika z całkowania przez części, a gęstość z gęstości
$C_c^\infty(\R)$ w $L^2$. Dla $f(x)=e^{-x^2}$ wybieramy
$\eta\in C_c^\infty(\R)$ równą $1$ w otoczeniu zera i
$f_n(x)=\eta(x/n)f(x)$. Wtedy $f_n\to f$ i
$f_n'=\eta(\cdot/n)f'+n^{-1}\eta'(\cdot/n)f\to f'$ w $L^2$
przez zbieżność zmajoryzowaną oraz oszacowanie drugiego składnika.
Jednak $f$ nie ma zwartego nośnika, więc wykres ograniczenia
nie jest domknięty.
\end{przyklad}

\subsection{Rozszerzenia samosprzężone i indeksy defektu}

\begin{definicja}[Przestrzenie defektu]
\label{def:deficiency33}
Dla gęsto określonego domkniętego operatora symetrycznego
$S$ definiujemy przestrzenie defektu
$N_+=\ker(S^*-iI)$ i $N_-=\ker(S^*+iI)$
oraz indeksy defektu $n_\pm=\dim N_\pm$ rozumiane jako liczby
kardynalne, tj. moce baz ortonormalnych (w szczególności
dopuszczamy wymiary nieprzeliczalne).
\end{definicja}

\begin{twierdzenie}[Rozkład von Neumanna i rozszerzenia]
\label{thm:von-neumann33}
Niech $S$ będzie gęsto określonym domkniętym operatorem
symetrycznym w zespolonej przestrzeni Hilberta $H$.
W normie wykresowej $D(S^*)$ zachodzi ortogonalny rozkład
\[
 D(S^*)=D(S)\oplus N_+\oplus N_-.
\]
Operator $S$ jest samosprzężony wtedy i tylko wtedy,
gdy $N_+=N_-=\{0\}$. Istnieje rozszerzenie samosprzężone
$S\subseteq\widetilde S=\widetilde S^*$ wtedy i tylko wtedy,
gdy $\dim N_+=\dim N_-$. Wszystkie takie rozszerzenia są
parametryzowane operatorami unitarnymi $V:N_+\to N_-$:
\[
 D(S_V)=D(S)\dotplus\{u+Vu:u\in N_+\},\qquad
 S_V(x+u+Vu)=Sx+iu-iVu.
\]
\end{twierdzenie}
\begin{proof}
Operator $S^*$ jest domknięty, a $D(S)$ jest domkniętą
podprzestrzenią przestrzeni Hilberta $D(S^*)$ z normą wykresową.
Dla $x\in D(S)$ i $u\in N_\pm$ otrzymujemy, korzystając
z $S\subseteq S^*$,
$\langle x,u\rangle_{S^*}=\langle x,u\rangle+
\langle Sx,\pm i u\rangle=0$.
Również $N_+\perp N_-$ w iloczynie wykresowym, ponieważ
$\langle iu,-ip\rangle=-\langle u,p\rangle$.
Każda przestrzeń $N_\pm$ jest domknięta w $H$: jeśli
$u_n\in N_\pm$ i $u_n\to u$, to $S^*u_n=\pm iu_n\to\pm iu$,
więc z domkniętości $S^*$ wynika $u\in N_\pm$.
Na $N_\pm$ norma wykresowa jest równa $\sqrt2$ razy normie $H$,
co uzasadnia użycie zupełności tych przestrzeni poniżej.

Niech $z\in D(S^*)$ będzie prostopadły w normie wykresowej
do $D(S)$. To znaczy
$\langle x,z\rangle+\langle Sx,S^*z\rangle=0$
dla każdego $x\in D(S)$. Z definicji operatora sprzężonego
wynika $S^*z\in D(S^*)$ oraz $S^{*2}z=-z$.
Stąd $z_\pm=(z\mp iS^*z)/2$ należą odpowiednio do $N_\pm$
i $z=z_++z_-$. Otrzymujemy cały rozkład.

Dla wektorów $u,v\in N_+$ i $p,q\in N_-$
forma brzegowa
$\mathfrak b(w,z)=\langle S^*w,z\rangle-
\langle w,S^*z\rangle$ spełnia
$\mathfrak b(u+p,v+q)=2i(\langle u,v\rangle-\langle p,q\rangle)$.
Forma $\mathfrak b$ znika, gdy jeden z argumentów należy do
$D(S)$, na mocy definicji sprzężenia. Dla argumentów z $N_+$ i
$N_-$ wyrazy mieszane znoszą się, ponieważ sprzężenie skalaru
$\pm i$ w drugim argumencie zmienia jego znak.
Każde domknięte rozszerzenie $S\subseteq T\subseteq S^*$
odpowiada domkniętej podprzestrzeni
$L\subseteq N_+\oplus N_-$ przez
$D(T)=D(S)\oplus L$. Operator jest symetryczny dokładnie
wtedy, gdy $\mathfrak b$ znika na $L\times L$.
Ponieważ $S\subseteq T$, mamy $T^*\subseteq S^*$.
Dla $z\in D(S^*)$ warunek $z\in D(T^*)$ jest równoważny
$\langle Tw,z\rangle=\langle w,S^*z\rangle$ dla wszystkich
$w\in D(T)$, czyli $\mathfrak b(w,z)=0$ dla wszystkich $w\in L$.
Zatem $D(T^*)=D(S)\oplus L^{\perp_{\mathfrak b}}$, gdzie
$L^{\perp_{\mathfrak b}}=\{z\in N_+\oplus N_-:\mathfrak b(w,z)=0
\ \text{dla wszystkich }w\in L\}$.
Samosprzężoność oznacza więc $L=L^{\perp_{\mathfrak b}}$.

Niech najpierw $L$ odpowiada operatorowi samosprzężonemu.
Jeżeli $u+p\in L$, $u\in N_+$, $p\in N_-$, to
$\mathfrak b(u+p,u+p)=0$ daje $\|u\|=\|p\|$.
Rzut $L$ na $N_+$ jest więc różnowartościowy, a $L$
ma postać wykresu pewnej izometrii
$V:D(V)\subseteq N_+\to N_-$. Podprzestrzeń $D(V)$
jest domknięta: jeśli $u_n\to u$ w $N_+$, to
$Vu_n$ jest ciągiem Cauchy'ego i zbiega do pewnego $p\in N_-$.
Ponieważ $L$ jest domknięta w $N_+\oplus N_-$,
$u+p\in L$, a więc $u\in D(V)$ i $Vu=p$. Jeśli $D(V)\ne N_+$,
istniałby $0\ne w\in D(V)^\perp\subset N_+$.
Wówczas $\mathfrak b(w,l)=0$ dla każdego $l\in L$, ale
$w\notin L$, wbrew $L=L^{\perp_{\mathfrak b}}$.
A więc $D(V)=N_+$. Analogicznie obraz $V$ jest domknięty,
a gdyby nie był całym $N_-$, wektor
$0\ne q\in (\operatorname{Ran}V)^\perp$ należałby do
$L^{\perp_{\mathfrak b}}\setminus L$. Zatem $V$ jest unitarny.
Odwrotnie, dla dowolnego unitarnego $V:N_+\to N_-$
jego wykres $L=\{u+Vu:u\in N_+\}$ jest domknięty
i spełnia $L=L^{\perp_{\mathfrak b}}$: dla $v+q\in N_+\oplus N_-$
warunek $\mathfrak b(u+Vu,v+q)=0$ dla każdego $u\in N_+$
oznacza $\langle u,v-V^*q\rangle=0$, a więc $v=V^*q$,
równoważnie $q=Vv$. Otrzymujemy samosprzężone rozszerzenie
o wskazanej dziedzinie. Operator unitarny między $N_+$ i $N_-$
istnieje dokładnie wtedy, gdy ich bazy ortonormalne mają równe moce:
w jednym kierunku obraz bazy przez unitarny operator jest bazą,
w drugim bijekcja baz wyznacza izometrię na ich liniowych
rozpiętościach, przedłużającą się na całe przestrzenie.
Gdy $N_+=N_- =\{0\}$, rozkład dziedziny daje
$D(S^*)=D(S)$, a więc $S=S^*$.
\end{proof}

\begin{przyklad}[Pochodna na przedziale i warunki brzegowe]
\label{ex:momentum-interval33}
Na $H=L^2(0,1)$ rozważmy
$Sf=-if'$ z dziedziną
$D(S)=\{f\in H^1(0,1):f(0)=f(1)=0\}$.
Przestrzeń $H^1(0,1)$ oznacza funkcje z $L^2(0,1)$
o słabej pochodnej w $L^2$, identyfikowane z ich
absolutnie ciągłymi reprezentantami.
Gęstość dziedziny wynika z zawierania $C_c^\infty(0,1)$.
Całkowanie przez części daje symetrię. Norma wykresowa jest
normą $H^1$, która jest zupełna: jeżeli $f_n\to f$ i $f_n'\to g$
w $L^2$, testowanie na funkcjach testowych daje $f'=g$ w sensie
dystrybucyjnym. Wartości brzegowe są ciągłe w tej normie, gdyż
dla $t\in[0,1]$ mamy
$|f(0)|\le |f(t)|+\int_0^1|f'(s)|\,\mathrm d\lambda_1(s)$
i analogicznie dla $f(1)$. Całkując względem $t$ oraz stosując
nierówność Cauchy'ego--Schwarza, dostajemy
$|f(0)|,|f(1)|\le\|f\|_{L^2}+\|f'\|_{L^2}$.
Zatem $D(S)$ jest domkniętą podprzestrzenią $H^1$, a operator
$S$ jest domknięty na mocy kryterium normy wykresowej.
Operator sprzężony ma ten sam wzór $S^*f=-if'$, lecz
$D(S^*)=H^1(0,1)$. Istotnie, testowanie warunku
sprzężenia na funkcjach z $C_c^\infty(0,1)$ wymusza istnienie
słabej pochodnej w $L^2$, a zatem przynależność do $H^1$.
Odwrotnie, całkowanie przez części dla dwóch funkcji z $H^1$
pokazuje, że wyrazy brzegowe znikają, ponieważ
wektor z $D(S)$ ma zerowe wartości na obu końcach.
Równania $S^*f=if$ i $S^*f=-if$ mają odpowiednio
rozwiązania $f(x)=ce^{-x}$ i $f(x)=ce^x$.
Zatem $n_+=n_-=1$. Wzór na formę brzegową brzmi
$\mathfrak b(f,g)=-i(f(1)\overline{g(1)}-f(0)\overline{g(0)})$.
Rozszerzenia samosprzężone są dokładnie operatorami
$-if'$ o dziedzinach
$D(S_\theta)=\{f\in H^1(0,1):f(1)=e^{i\theta}f(0)\}$,
$\theta\in[0,2\pi)$. Aby to uzasadnić, odwzorowanie śladu
$\gamma:H^1(0,1)\to\C^2$, $\gamma(f)=(f(0),f(1))$,
jest surjektywne (wystarczą funkcje afiniczne) i ma jądro $D(S)$.
Dziedziny rozszerzeń samosprzężonych odpowiadają więc podprzestrzeniom
$L\subseteq\C^2$ równym własnemu anihilatorowi względem formy
$-i(b\overline d-a\overline c)$ na parach $(a,b),(c,d)$.
Taka podprzestrzeń ma wymiar $1$, gdyż forma jest niezdegenerowana.
Dla jej niezerowego generatora $(a,b)$ izotropowość daje $|a|=|b|$,
więc $a\ne0$ i $b=e^{i\theta}a$. Odwrotnie każda taka prosta
jest izotropowa i ma wymiar $1$, stąd jest równa własnemu
anihilatorowi. Są to różne realizacje \emph{tego samego wzoru}.
\end{przyklad}

\begin{przyklad}[Brak rozszerzenia samosprzężonego]
Na $L^2(0,\infty)$ operator $-if'$ z dziedziną
$H_0^1(0,\infty)=\{f\in H^1(0,\infty):f(0)=0\}$
jest domknięty i symetryczny. Przyjmujemy, że $H^1(0,\infty)$
składa się z funkcji $L^2$ o słabej pochodnej w $L^2$;
każda z nich ma lokalnie absolutnie ciągłego reprezentanta.
Dziedzina jest gęsta, gdyż zawiera $C_c^\infty(0,\infty)$.
Zupełność $H^1$ wynika z tego samego argumentu dystrybucyjnego
co na odcinku, a ślad w zerze jest ciągły dzięki oszacowaniu
na $[0,1]$. To daje domkniętość. Dla $f\in H^1(0,\infty)$
pochodna $|f|^2$ jest równa $2\operatorname{Re}(f'\overline f)$
i należy do $L^1$ przez nierówność Cauchy'ego--Schwarza.
Zatem $|f(t)|^2$ ma skończoną granicę przy $t\to\infty$,
a z całkowalności $|f|^2$ wynika, że granica ta jest zerowa.
Całkowanie przez części na $[0,R]$ i przejście z $R$ do
nieskończoności usuwa więc wyraz brzegowy w nieskończoności.
Dla dwóch funkcji z dziedziny znika również wyraz w zerze,
co daje symetrię. Testowanie sprzężenia na funkcjach testowych
wymusza słabą pochodną w $L^2$, a całkowanie przez części
z dowolną funkcją z $H^1$ daje inkluzję odwrotną.
Jego operator sprzężony ma zatem dziedzinę $H^1(0,\infty)$. Z równań defektu otrzymujemy
$N_+=\lin\{e^{-x}\}$ i $N_- =\{0\}$, ponieważ $e^x\notin L^2$.
Indeksy defektu różnią się, więc operator nie ma rozszerzenia
samosprzężonego w tej samej przestrzeni Hilberta.
\end{przyklad}

\subsection{Jednoparametrowe grupy unitarne i twierdzenie Stone'a}

\begin{definicja}[Silnie ciągła grupa unitarna]
Niech $(U_t)_{t\in\R}$ będzie rodziną operatorów unitarnych
na $H$ z $U_0=I_H$ i $U_{t+s}=U_tU_s$ dla $s,t\in\R$.
Jest to \emph{silnie ciągła jednoparametrowa grupa unitarna}, gdy
$\lim_{t\to0}\|U_tx-x\|=0$ dla każdego $x\in H$.
Jej generator definiujemy na dziedzinie
$D(G)=\{x:\lim_{t\to0}(U_tx-x)/t\text{ istnieje w }H\}$
przez $Gx=\lim_{t\to0}(U_tx-x)/t$.
\end{definicja}

\begin{twierdzenie}[Twierdzenie Stone'a]
\label{thm:stone33}
Dla każdej silnie ciągłej jednoparametrowej grupy unitarnej
$(U_t)_{t\in\R}$ na zespolonej przestrzeni Hilberta $H$
istnieje dokładnie jeden operator samosprzężony $A$ taki, że
$U_t=e^{itA}$ dla wszystkich $t\in\R$, gdzie
$e^{itA}=\int_\R e^{it\lambda}\,\mathrm dE_A(\lambda)$.
Ponadto $D(A)=D(G)$ i $G=iA$.
Odwrotnie, dla każdego samosprzężonego $A$ rachunek spektralny
wyznacza silnie ciągłą grupę unitarną $(e^{itA})_{t\in\R}$,
a jej generator jest równy $iA$.
\end{twierdzenie}
\begin{proof}
\textbf{Krok 1: konstrukcja grupy z operatora.}
Dla samosprzężonego $A$ funkcja $\lambda\mapsto e^{it\lambda}$
jest ograniczona i ma moduł $1$, zatem rachunek borelowski
definiuje operator $U_t=e^{itA}$ spełniający
$U_t^*U_t=U_tU_t^*=1(E_A)=I_H$. Operator jest więc unitarny.
Równości $U_{t+s}=U_tU_s$ i $U_0=I_H$ są punktowymi
tożsamościami funkcji w rachunku spektralnym.
Dla każdego $x\in H$ mamy
$\|(U_t-I_H)x\|^2=\int_\R|e^{it\lambda}-1|^2\,\mathrm d\mu_x(\lambda)
\to0$ na mocy zbieżności zmajoryzowanej.
Dla $x\in D(A)$ nierówność
$|(e^{it\lambda}-1)/t|\le|\lambda|$ (z całkowania po odcinku
funkcji $s\mapsto i\lambda e^{is\lambda}$) i zbieżność punktowa
ilorazu różnicowego do $i\lambda$ dają
\[
 \left\|\frac{U_tx-x}{t}-iAx\right\|^2
 =\int_\R\left|\frac{e^{it\lambda}-1}{t}-i\lambda\right|^2
 \,\mathrm d\mu_x(\lambda)\longrightarrow0.
\]
Równość wynika z addytywności na wspólnej dziedzinie
w dowodzie Lematu~\ref{lem:unbounded-integral33}, a funkcja pod
całką jest ograniczona przez całkowalną funkcję $4\lambda^2$.
Odwrotnie, gdy dla $x$ iloraz ma granicę, lemat Fatou
dla miary $\mu_x$ i
$\lim_{t\to0}|(e^{it\lambda}-1)/t|^2=\lambda^2$
dają, biorąc np. ciąg $t=1/n$,
$\int_\R\lambda^2\,\mathrm d\mu_x\le
\liminf_n\|n(U_{1/n}x-x)\|^2<\infty$.
Zatem $D(G)=D(A)$.

\textbf{Krok 2: gęstość i skośna symetria generatora.}
Niech teraz $(U_t)$ będzie dowolną silnie ciągłą grupą unitarną.
Dla $\phi\in C_c^\infty(\R)$ i $x\in H$ definiujemy
$x_\phi=\int_\R\phi(t)U_tx\,\mathrm d\lambda_1(t)$.
Jest to całka w normie przestrzeni Hilberta, zdefiniowana
przez granice sum Riemanna na zwartym nośniku $\phi$;
nie potrzeba tu ogólnej teorii całki Pettisa. Ciągłość
$t\mapsto U_tx$ w każdym punkcie wynika z
$\|U_{t+s}x-U_tx\|=\|U_sx-x\|$.
Przesunięcie zmiennej pod całką daje dla $s\ne0$
\[
 \frac{U_sx_\phi-x_\phi}{s}
 =\int_\R\frac{\phi(t-s)-\phi(t)}s U_tx\,\mathrm d\lambda_1(t).
\]
Dla $|s|\le1$ ilorazy mają wspólny zwarty nośnik i zbiegają
jednostajnie do $-\phi'$ (z jednostajnej ciągłości $\phi'$).
Oszacowanie normy całki przez całkę norm dowodzi zbieżności
do $-\int_\R\phi'(t)U_tx\,\mathrm d\lambda_1(t)$ w $H$.
W szczególności $x_\phi\in D(G)$. Wybierając
nieujemną aproksymację jedności $\phi_n$ wokół zera,
otrzymujemy $x_{\phi_n}\to x$ z silnej ciągłości grupy.
Dlatego $D(G)$ jest gęsta.
Dla $x,y\in D(G)$ różniczkowanie
$\langle U_tx,U_ty\rangle=\langle x,y\rangle$ w zerze daje
$\langle Gx,y\rangle+\langle x,Gy\rangle=0$.
Stąd operator $A=-iG$ jest symetryczny.

\textbf{Krok 3: surjektywność $A\pm iI$.}
Dla $x\in H$ kładziemy
$R_+x=\int_0^\infty e^{-t}U_tx\,\mathrm d\lambda_1(t)$,
$R_-x=\int_0^\infty e^{-t}U_{-t}x\,\mathrm d\lambda_1(t)$.
Są to granice w normie całek na odcinkach $[0,N]$;
$\|R_\pm\|\le1$, ponieważ $\|U_tx\|=\|x\|$.
Dla $s\in\R$ zmiana zmiennej daje
$U_sR_+x=e^s\int_s^\infty e^{-r}U_rx\,\mathrm d\lambda_1(r)$.
Iloraz różnicowy w zerze zbiega więc w normie do
$R_+x-x$ (z silnej ciągłości $U_rx$ w $r=0$).
Stąd $R_+x\in D(G)$ i $(I_H-G)R_+x=x$.
Podobnie
$U_sR_-x=e^{-s}\int_{-s}^\infty e^{-r}U_{-r}x\,\mathrm d\lambda_1(r)$,
zatem $GR_-x=x-R_-x$ i $(I_H+G)R_-x=x$.
W konsekwencji
$\operatorname{Ran}(I_H\pm G)=H$. Ponieważ $A=-iG$,
mamy $A+iI=i(I_H-G)$ i $A-iI=-i(I_H+G)$ na $D(G)$,
więc $\operatorname{Ran}(A\pm iI)=H$.
Twierdzenie~\ref{thm:selfadjoint-range33} daje samosprzężoność $A$.

\textbf{Krok 4: identyfikacja i jednoznaczność grupy.}
Z własności grupowej wynika
$U_tD(G)\subseteq D(G)$ i $GU_tx=U_tGx$:
istotnie, iloraz różnicowy w punkcie $t$ jest obrazem
ilorazu różnicowego w zerze przez ograniczony operator $U_t$.
Analogiczna własność zachodzi dla $e^{itA}$
na mocy rachunku spektralnego.
Dla $x\in D(A)$ zarówno $t\mapsto U_tx$, jak i
$t\mapsto e^{itA}x$ są więc różniczkowalne i spełniają
równanie $v'(t)=iAv(t)$ oraz warunek $v(0)=x$.
Ich różnica $w$ ma stałą normę, ponieważ
$\frac{\mathrm d}{\mathrm dt}\|w(t)\|^2
=2\operatorname{Re}\langle iAw(t),w(t)\rangle=0$.
Ponieważ $w(0)=0$, mamy $w=0$.
Gęstość $D(A)$ i unitarność obu grup rozszerzają równość
na całe $H$. Jednoznaczność $A$ wynika z definicji generatora.
\end{proof}
\endgroup

\end{document}
